\documentclass[a4paper,12pt]{article}

% --- Pakete ---
\usepackage[utf8]{inputenc}   % Umlaute direkt eingeben
\usepackage[T1]{fontenc}
\usepackage[english]{babel}
\usepackage{lmodern}          % Schöne Schrift
\usepackage{geometry}         % Seitenränder anpassen
\usepackage{booktabs}         % Für schönere Tabellen
\usepackage{array}
\usepackage{setspace}         % Zeilenabstand
\usepackage{graphicx}
\usepackage{caption}
\usepackage{url}
\usepackage{todonotes}
\usepackage{float}
\usepackage{pgfplots}
\usepackage{xfp}      % provides \fpeval for calculations
\usepackage{multirow}
\usepackage[
separate-uncertainty = true,
multi-part-units = repeat
]{siunitx}
\usepackage[hidelinks, bookmarks=true]{hyperref}
\pgfplotsset{compat=1.18}

% --- Einstellungen ---
\geometry{a4paper, margin=2.5cm}
\setstretch{1.25}

% --- Dokument ---
\begin{document}
	
	% Titel
	\begin{center}
		{\LARGE \textbf{Praktikum Hochfrequenz-Schaltungstechnik WS25/26}} \\[1em]
		{\large  \textbf{Bericht}} \\[0.5em]
	\end{center}
	
\begin{center}
	\begin{tabular}{|l|p{4cm}|l|p{4cm}|}
		\hline
		\textbf{Name:} &  Sebastian Grigorevski & \textbf{Matrikelnr.:} &  35690104 \\
		\hline
		\textbf{Versuchsnr.:} & 4 & \textbf{Datum:} & \today \\
		\hline
	\end{tabular}
\end{center}
	
	\vspace{1cm}
	\setcounter{section}{4}  % setzt den Section Zähler auf 1
	\setcounter{subsection}{4}  % setzt den Subsection Zähler auf 5
	\setcounter{subsubsection}{0}  % setzt den Subsubsection Zähler auf 0

	\subsubsection{}
	
	At 1 M$\Omega$ the DSO acts as a high resistance measurment tool and it measures the value of the source
	When changing the input impedance of the DSO from 1 M$\Omega$ to 50 $\Omega$ by putting a terminator in parallel to the DSO following change can be observed: The measured voltage is half the value from before due to the voltage divider in the matched system.
	
	\subsubsection{}

	The results of the FFT of a rectangle, triangle and a sawtooth wave are displayed in following table:
	
	\begin{table}[h]
	\centering
		\caption{Fourier coefficients comparison}
	\begin{tabular}{l|c|c}
		Signal  & $a_n$ & $b_n$ \\ \hline
		Rectangle  & $0$ & $\frac{1-(-1)^n}{\pi n}$ \\
		Triangle  & $\frac{2}{\pi^2}\frac{1-(-1)^n}{n^2}$ & $0$ \\
		Sawtooth  & $0$ & $-\frac{2(-1)^{n}}{n}$
	\end{tabular}
	\end{table}

	The firs difference lays in the symmetry of the signal shapes. Both rectangle and sawtooth signals as displayed in \cite{matheonline_fourier_tutorial} are anti symmetric, thus they only have sine contributions as can be seen from 
	\[
	f(x) = \frac{a_0}{2} + \sum_{n=1}^\infty \left( a_n \cos(nx) + b_n \sin(nx) \right).
	\]
	The triangle function on the other hand is symmetric and only consists of cosines.
	
	The rectangle and the triangle functions only have uneven harmonics since the nominator terminates at even multiples on the base frequency. Therefore their frequency components are at $f, 3f, 5f ...$ . Where as the saw signal is composed from all multiple harmonics of the base frequency. 
	
	Further the rectangle and the sawtooth Fourier coefficients decay by $b_n\sim 1/n$ for higher harmonics. The triangle components are proportional with $a_n\sim 1/n^2$, meaning less influence from higher harmonics than with the other signal shapes.
	
	\subsubsection{}
	
	The amplitudes of the harmonics were determined using an FFT measurement and are
	given in decibels. The fundamental at \(f_1 = 500\,\mathrm{kHz}\) is used as the
	reference.
	
	\begin{table}[h]
		\centering		
		\caption{FFT measurement values of the measured voltages}
		\label{tab:fft}
		\begin{tabular}{c c c c c c}
			\hline
			Harmonic & Frequency & Sine [dB] & Rectangular [dB] & Triangular [dB] & Sawtooth [dB] \\
			\hline
			1 & \(500\,\mathrm{kHz}\) & \(+10\) & \(+12.4\) & \(+8\) & \(+5.6\) \\
			2 & \(1.0\,\mathrm{MHz}\) & \(-47.2\) & \(-45.6\) & \(-65\) & \(-1.2\) \\
			3 & \(1.5\,\mathrm{MHz}\) & \(-64.0\) & \(-2\) & \(-14\) & \(-5.6\) \\
			4 & \(2.0\,\mathrm{MHz}\) & \(-52.4\) & \(-41.6\) & \(-67.2\) & \(-8.4\) \\
			5 & \(2.5\,\mathrm{MHz}\) & \(-56.4\) & \(-3.6\) & \(-23.6\) & \(-9.6\) \\
			6 & \(3.0\,\mathrm{MHz}\) & \(-54.0\) & \(-\) & \(-64\) & \(-12\) \\
			7 & \(3.5\,\mathrm{MHz}\) & \(-60.8\) & \(-6.4\) & \(-29.2\) & \(-13.2\) \\
			\hline
		\end{tabular}

	\end{table}
	
	The harmonic distortion factor is defined as
	\[
	\mathrm{THD}=\frac{\sqrt{U_2^2+U_3^2+U_4^2+\dots}}{U_1} \cite{klirr}
	\]
	Since the harmonic amplitudes are given in decibels, they must first be converted to
	linear amplitude ratios:
	
	\[
	\frac{U_n}{U_1}
	=
	10^{\frac{L_n - L_1}{20}}
	\]
	
	Substituting this expression into the THD definition yields:
	
	\[
	\mathrm{THD}
	=
	\sqrt{
		\sum_{n=2}^{N}
		\left(10^{\frac{L_n - L_1}{20}}\right)^2
	}
	=
	\sqrt{
		\sum_{n=2}^{N}
		10^{\frac{L_n - L_1}{10}}
	}
	\]
	
	Using the measured FFT values, the following result is obtained:
	
	\[
	\mathrm{THD} = 1.79 \times 10^{-3} \equiv 0.179\%
	\]
	
	Compared with the measurements of the spectrum analyzer shown in fig. \ref{fig:four_signals} the distinct peak at 500 kHz is visible while further harmonics either do not appear or are overshadowed by noise.
	
	Further a difference between the DSO and the spectrum analyzer results can be observed in the scaling where the spectrum analyzer defines the ratio with input and output power. Nevertheless both instruments confirm the discussed theoretical characteristics where the rectangle and the triangle signals only contain peaks at uneven harmonics as also can be seen in table \ref{tab:fft}. There both instruments show a quicker decay with the triangle signal as expected. The sawtooth signal also shows matching characteristics and peaks between instruments and theory.
	
	One difference remains at the sine wave measurement where the DSO captures higher order peaks slightly above the noise level while the analyzer lacks those peaks which in an ideal sine wave are not supposed to appear.
	
%	\begin{figure}[H]
%		\centering
%		%% Creator: Matplotlib, PGF backend
%%
%% To include the figure in your LaTeX document, write
%%   \input{<filename>.pgf}
%%
%% Make sure the required packages are loaded in your preamble
%%   \usepackage{pgf}
%%
%% Also ensure that all the required font packages are loaded; for instance,
%% the lmodern package is sometimes necessary when using math font.
%%   \usepackage{lmodern}
%%
%% Figures using additional raster images can only be included by \input if
%% they are in the same directory as the main LaTeX file. For loading figures
%% from other directories you can use the `import` package
%%   \usepackage{import}
%%
%% and then include the figures with
%%   \import{<path to file>}{<filename>.pgf}
%%
%% Matplotlib used the following preamble
%%
\begingroup%
\makeatletter%
\begin{pgfpicture}%
\pgfpathrectangle{\pgfpointorigin}{\pgfqpoint{6.300000in}{3.500000in}}%
\pgfusepath{use as bounding box, clip}%
\begin{pgfscope}%
\pgfsetbuttcap%
\pgfsetmiterjoin%
\definecolor{currentfill}{rgb}{1.000000,1.000000,1.000000}%
\pgfsetfillcolor{currentfill}%
\pgfsetlinewidth{0.000000pt}%
\definecolor{currentstroke}{rgb}{1.000000,1.000000,1.000000}%
\pgfsetstrokecolor{currentstroke}%
\pgfsetdash{}{0pt}%
\pgfpathmoveto{\pgfqpoint{0.000000in}{0.000000in}}%
\pgfpathlineto{\pgfqpoint{6.300000in}{0.000000in}}%
\pgfpathlineto{\pgfqpoint{6.300000in}{3.500000in}}%
\pgfpathlineto{\pgfqpoint{0.000000in}{3.500000in}}%
\pgfpathlineto{\pgfqpoint{0.000000in}{0.000000in}}%
\pgfpathclose%
\pgfusepath{fill}%
\end{pgfscope}%
\begin{pgfscope}%
\pgfsetbuttcap%
\pgfsetmiterjoin%
\definecolor{currentfill}{rgb}{1.000000,1.000000,1.000000}%
\pgfsetfillcolor{currentfill}%
\pgfsetlinewidth{0.000000pt}%
\definecolor{currentstroke}{rgb}{0.000000,0.000000,0.000000}%
\pgfsetstrokecolor{currentstroke}%
\pgfsetstrokeopacity{0.000000}%
\pgfsetdash{}{0pt}%
\pgfpathmoveto{\pgfqpoint{0.758026in}{0.565123in}}%
\pgfpathlineto{\pgfqpoint{6.150000in}{0.565123in}}%
\pgfpathlineto{\pgfqpoint{6.150000in}{3.350000in}}%
\pgfpathlineto{\pgfqpoint{0.758026in}{3.350000in}}%
\pgfpathlineto{\pgfqpoint{0.758026in}{0.565123in}}%
\pgfpathclose%
\pgfusepath{fill}%
\end{pgfscope}%
\begin{pgfscope}%
\pgfpathrectangle{\pgfqpoint{0.758026in}{0.565123in}}{\pgfqpoint{5.391974in}{2.784877in}}%
\pgfusepath{clip}%
\pgfsetbuttcap%
\pgfsetroundjoin%
\pgfsetlinewidth{0.501875pt}%
\definecolor{currentstroke}{rgb}{0.690196,0.690196,0.690196}%
\pgfsetstrokecolor{currentstroke}%
\pgfsetstrokeopacity{0.700000}%
\pgfsetdash{{1.850000pt}{0.800000pt}}{0.000000pt}%
\pgfpathmoveto{\pgfqpoint{1.748388in}{0.565123in}}%
\pgfpathlineto{\pgfqpoint{1.748388in}{3.350000in}}%
\pgfusepath{stroke}%
\end{pgfscope}%
\begin{pgfscope}%
\pgfsetbuttcap%
\pgfsetroundjoin%
\definecolor{currentfill}{rgb}{0.000000,0.000000,0.000000}%
\pgfsetfillcolor{currentfill}%
\pgfsetlinewidth{0.803000pt}%
\definecolor{currentstroke}{rgb}{0.000000,0.000000,0.000000}%
\pgfsetstrokecolor{currentstroke}%
\pgfsetdash{}{0pt}%
\pgfsys@defobject{currentmarker}{\pgfqpoint{0.000000in}{-0.048611in}}{\pgfqpoint{0.000000in}{0.000000in}}{%
\pgfpathmoveto{\pgfqpoint{0.000000in}{0.000000in}}%
\pgfpathlineto{\pgfqpoint{0.000000in}{-0.048611in}}%
\pgfusepath{stroke,fill}%
}%
\begin{pgfscope}%
\pgfsys@transformshift{1.748388in}{0.565123in}%
\pgfsys@useobject{currentmarker}{}%
\end{pgfscope}%
\end{pgfscope}%
\begin{pgfscope}%
\definecolor{textcolor}{rgb}{0.000000,0.000000,0.000000}%
\pgfsetstrokecolor{textcolor}%
\pgfsetfillcolor{textcolor}%
\pgftext[x=1.748388in,y=0.467901in,,top]{\color{textcolor}\rmfamily\fontsize{10.000000}{12.000000}\selectfont \(\displaystyle {1}\)}%
\end{pgfscope}%
\begin{pgfscope}%
\pgfpathrectangle{\pgfqpoint{0.758026in}{0.565123in}}{\pgfqpoint{5.391974in}{2.784877in}}%
\pgfusepath{clip}%
\pgfsetbuttcap%
\pgfsetroundjoin%
\pgfsetlinewidth{0.501875pt}%
\definecolor{currentstroke}{rgb}{0.690196,0.690196,0.690196}%
\pgfsetstrokecolor{currentstroke}%
\pgfsetstrokeopacity{0.700000}%
\pgfsetdash{{1.850000pt}{0.800000pt}}{0.000000pt}%
\pgfpathmoveto{\pgfqpoint{2.848791in}{0.565123in}}%
\pgfpathlineto{\pgfqpoint{2.848791in}{3.350000in}}%
\pgfusepath{stroke}%
\end{pgfscope}%
\begin{pgfscope}%
\pgfsetbuttcap%
\pgfsetroundjoin%
\definecolor{currentfill}{rgb}{0.000000,0.000000,0.000000}%
\pgfsetfillcolor{currentfill}%
\pgfsetlinewidth{0.803000pt}%
\definecolor{currentstroke}{rgb}{0.000000,0.000000,0.000000}%
\pgfsetstrokecolor{currentstroke}%
\pgfsetdash{}{0pt}%
\pgfsys@defobject{currentmarker}{\pgfqpoint{0.000000in}{-0.048611in}}{\pgfqpoint{0.000000in}{0.000000in}}{%
\pgfpathmoveto{\pgfqpoint{0.000000in}{0.000000in}}%
\pgfpathlineto{\pgfqpoint{0.000000in}{-0.048611in}}%
\pgfusepath{stroke,fill}%
}%
\begin{pgfscope}%
\pgfsys@transformshift{2.848791in}{0.565123in}%
\pgfsys@useobject{currentmarker}{}%
\end{pgfscope}%
\end{pgfscope}%
\begin{pgfscope}%
\definecolor{textcolor}{rgb}{0.000000,0.000000,0.000000}%
\pgfsetstrokecolor{textcolor}%
\pgfsetfillcolor{textcolor}%
\pgftext[x=2.848791in,y=0.467901in,,top]{\color{textcolor}\rmfamily\fontsize{10.000000}{12.000000}\selectfont \(\displaystyle {2}\)}%
\end{pgfscope}%
\begin{pgfscope}%
\pgfpathrectangle{\pgfqpoint{0.758026in}{0.565123in}}{\pgfqpoint{5.391974in}{2.784877in}}%
\pgfusepath{clip}%
\pgfsetbuttcap%
\pgfsetroundjoin%
\pgfsetlinewidth{0.501875pt}%
\definecolor{currentstroke}{rgb}{0.690196,0.690196,0.690196}%
\pgfsetstrokecolor{currentstroke}%
\pgfsetstrokeopacity{0.700000}%
\pgfsetdash{{1.850000pt}{0.800000pt}}{0.000000pt}%
\pgfpathmoveto{\pgfqpoint{3.949194in}{0.565123in}}%
\pgfpathlineto{\pgfqpoint{3.949194in}{3.350000in}}%
\pgfusepath{stroke}%
\end{pgfscope}%
\begin{pgfscope}%
\pgfsetbuttcap%
\pgfsetroundjoin%
\definecolor{currentfill}{rgb}{0.000000,0.000000,0.000000}%
\pgfsetfillcolor{currentfill}%
\pgfsetlinewidth{0.803000pt}%
\definecolor{currentstroke}{rgb}{0.000000,0.000000,0.000000}%
\pgfsetstrokecolor{currentstroke}%
\pgfsetdash{}{0pt}%
\pgfsys@defobject{currentmarker}{\pgfqpoint{0.000000in}{-0.048611in}}{\pgfqpoint{0.000000in}{0.000000in}}{%
\pgfpathmoveto{\pgfqpoint{0.000000in}{0.000000in}}%
\pgfpathlineto{\pgfqpoint{0.000000in}{-0.048611in}}%
\pgfusepath{stroke,fill}%
}%
\begin{pgfscope}%
\pgfsys@transformshift{3.949194in}{0.565123in}%
\pgfsys@useobject{currentmarker}{}%
\end{pgfscope}%
\end{pgfscope}%
\begin{pgfscope}%
\definecolor{textcolor}{rgb}{0.000000,0.000000,0.000000}%
\pgfsetstrokecolor{textcolor}%
\pgfsetfillcolor{textcolor}%
\pgftext[x=3.949194in,y=0.467901in,,top]{\color{textcolor}\rmfamily\fontsize{10.000000}{12.000000}\selectfont \(\displaystyle {3}\)}%
\end{pgfscope}%
\begin{pgfscope}%
\pgfpathrectangle{\pgfqpoint{0.758026in}{0.565123in}}{\pgfqpoint{5.391974in}{2.784877in}}%
\pgfusepath{clip}%
\pgfsetbuttcap%
\pgfsetroundjoin%
\pgfsetlinewidth{0.501875pt}%
\definecolor{currentstroke}{rgb}{0.690196,0.690196,0.690196}%
\pgfsetstrokecolor{currentstroke}%
\pgfsetstrokeopacity{0.700000}%
\pgfsetdash{{1.850000pt}{0.800000pt}}{0.000000pt}%
\pgfpathmoveto{\pgfqpoint{5.049597in}{0.565123in}}%
\pgfpathlineto{\pgfqpoint{5.049597in}{3.350000in}}%
\pgfusepath{stroke}%
\end{pgfscope}%
\begin{pgfscope}%
\pgfsetbuttcap%
\pgfsetroundjoin%
\definecolor{currentfill}{rgb}{0.000000,0.000000,0.000000}%
\pgfsetfillcolor{currentfill}%
\pgfsetlinewidth{0.803000pt}%
\definecolor{currentstroke}{rgb}{0.000000,0.000000,0.000000}%
\pgfsetstrokecolor{currentstroke}%
\pgfsetdash{}{0pt}%
\pgfsys@defobject{currentmarker}{\pgfqpoint{0.000000in}{-0.048611in}}{\pgfqpoint{0.000000in}{0.000000in}}{%
\pgfpathmoveto{\pgfqpoint{0.000000in}{0.000000in}}%
\pgfpathlineto{\pgfqpoint{0.000000in}{-0.048611in}}%
\pgfusepath{stroke,fill}%
}%
\begin{pgfscope}%
\pgfsys@transformshift{5.049597in}{0.565123in}%
\pgfsys@useobject{currentmarker}{}%
\end{pgfscope}%
\end{pgfscope}%
\begin{pgfscope}%
\definecolor{textcolor}{rgb}{0.000000,0.000000,0.000000}%
\pgfsetstrokecolor{textcolor}%
\pgfsetfillcolor{textcolor}%
\pgftext[x=5.049597in,y=0.467901in,,top]{\color{textcolor}\rmfamily\fontsize{10.000000}{12.000000}\selectfont \(\displaystyle {4}\)}%
\end{pgfscope}%
\begin{pgfscope}%
\pgfpathrectangle{\pgfqpoint{0.758026in}{0.565123in}}{\pgfqpoint{5.391974in}{2.784877in}}%
\pgfusepath{clip}%
\pgfsetbuttcap%
\pgfsetroundjoin%
\pgfsetlinewidth{0.501875pt}%
\definecolor{currentstroke}{rgb}{0.690196,0.690196,0.690196}%
\pgfsetstrokecolor{currentstroke}%
\pgfsetstrokeopacity{0.700000}%
\pgfsetdash{{1.850000pt}{0.800000pt}}{0.000000pt}%
\pgfpathmoveto{\pgfqpoint{6.150000in}{0.565123in}}%
\pgfpathlineto{\pgfqpoint{6.150000in}{3.350000in}}%
\pgfusepath{stroke}%
\end{pgfscope}%
\begin{pgfscope}%
\pgfsetbuttcap%
\pgfsetroundjoin%
\definecolor{currentfill}{rgb}{0.000000,0.000000,0.000000}%
\pgfsetfillcolor{currentfill}%
\pgfsetlinewidth{0.803000pt}%
\definecolor{currentstroke}{rgb}{0.000000,0.000000,0.000000}%
\pgfsetstrokecolor{currentstroke}%
\pgfsetdash{}{0pt}%
\pgfsys@defobject{currentmarker}{\pgfqpoint{0.000000in}{-0.048611in}}{\pgfqpoint{0.000000in}{0.000000in}}{%
\pgfpathmoveto{\pgfqpoint{0.000000in}{0.000000in}}%
\pgfpathlineto{\pgfqpoint{0.000000in}{-0.048611in}}%
\pgfusepath{stroke,fill}%
}%
\begin{pgfscope}%
\pgfsys@transformshift{6.150000in}{0.565123in}%
\pgfsys@useobject{currentmarker}{}%
\end{pgfscope}%
\end{pgfscope}%
\begin{pgfscope}%
\definecolor{textcolor}{rgb}{0.000000,0.000000,0.000000}%
\pgfsetstrokecolor{textcolor}%
\pgfsetfillcolor{textcolor}%
\pgftext[x=6.150000in,y=0.467901in,,top]{\color{textcolor}\rmfamily\fontsize{10.000000}{12.000000}\selectfont \(\displaystyle {5}\)}%
\end{pgfscope}%
\begin{pgfscope}%
\definecolor{textcolor}{rgb}{0.000000,0.000000,0.000000}%
\pgfsetstrokecolor{textcolor}%
\pgfsetfillcolor{textcolor}%
\pgftext[x=3.454013in,y=0.288889in,,top]{\color{textcolor}\rmfamily\fontsize{10.000000}{12.000000}\selectfont Frequency (Hz)}%
\end{pgfscope}%
\begin{pgfscope}%
\definecolor{textcolor}{rgb}{0.000000,0.000000,0.000000}%
\pgfsetstrokecolor{textcolor}%
\pgfsetfillcolor{textcolor}%
\pgftext[x=6.150000in,y=0.302778in,right,top]{\color{textcolor}\rmfamily\fontsize{10.000000}{12.000000}\selectfont \(\displaystyle \times{10^{6}}{}\)}%
\end{pgfscope}%
\begin{pgfscope}%
\pgfpathrectangle{\pgfqpoint{0.758026in}{0.565123in}}{\pgfqpoint{5.391974in}{2.784877in}}%
\pgfusepath{clip}%
\pgfsetbuttcap%
\pgfsetroundjoin%
\pgfsetlinewidth{0.501875pt}%
\definecolor{currentstroke}{rgb}{0.690196,0.690196,0.690196}%
\pgfsetstrokecolor{currentstroke}%
\pgfsetstrokeopacity{0.700000}%
\pgfsetdash{{1.850000pt}{0.800000pt}}{0.000000pt}%
\pgfpathmoveto{\pgfqpoint{0.758026in}{0.871206in}}%
\pgfpathlineto{\pgfqpoint{6.150000in}{0.871206in}}%
\pgfusepath{stroke}%
\end{pgfscope}%
\begin{pgfscope}%
\pgfsetbuttcap%
\pgfsetroundjoin%
\definecolor{currentfill}{rgb}{0.000000,0.000000,0.000000}%
\pgfsetfillcolor{currentfill}%
\pgfsetlinewidth{0.803000pt}%
\definecolor{currentstroke}{rgb}{0.000000,0.000000,0.000000}%
\pgfsetstrokecolor{currentstroke}%
\pgfsetdash{}{0pt}%
\pgfsys@defobject{currentmarker}{\pgfqpoint{-0.048611in}{0.000000in}}{\pgfqpoint{-0.000000in}{0.000000in}}{%
\pgfpathmoveto{\pgfqpoint{-0.000000in}{0.000000in}}%
\pgfpathlineto{\pgfqpoint{-0.048611in}{0.000000in}}%
\pgfusepath{stroke,fill}%
}%
\begin{pgfscope}%
\pgfsys@transformshift{0.758026in}{0.871206in}%
\pgfsys@useobject{currentmarker}{}%
\end{pgfscope}%
\end{pgfscope}%
\begin{pgfscope}%
\definecolor{textcolor}{rgb}{0.000000,0.000000,0.000000}%
\pgfsetstrokecolor{textcolor}%
\pgfsetfillcolor{textcolor}%
\pgftext[x=0.344444in, y=0.822981in, left, base]{\color{textcolor}\rmfamily\fontsize{10.000000}{12.000000}\selectfont \(\displaystyle {\ensuremath{-}100}\)}%
\end{pgfscope}%
\begin{pgfscope}%
\pgfpathrectangle{\pgfqpoint{0.758026in}{0.565123in}}{\pgfqpoint{5.391974in}{2.784877in}}%
\pgfusepath{clip}%
\pgfsetbuttcap%
\pgfsetroundjoin%
\pgfsetlinewidth{0.501875pt}%
\definecolor{currentstroke}{rgb}{0.690196,0.690196,0.690196}%
\pgfsetstrokecolor{currentstroke}%
\pgfsetstrokeopacity{0.700000}%
\pgfsetdash{{1.850000pt}{0.800000pt}}{0.000000pt}%
\pgfpathmoveto{\pgfqpoint{0.758026in}{1.186807in}}%
\pgfpathlineto{\pgfqpoint{6.150000in}{1.186807in}}%
\pgfusepath{stroke}%
\end{pgfscope}%
\begin{pgfscope}%
\pgfsetbuttcap%
\pgfsetroundjoin%
\definecolor{currentfill}{rgb}{0.000000,0.000000,0.000000}%
\pgfsetfillcolor{currentfill}%
\pgfsetlinewidth{0.803000pt}%
\definecolor{currentstroke}{rgb}{0.000000,0.000000,0.000000}%
\pgfsetstrokecolor{currentstroke}%
\pgfsetdash{}{0pt}%
\pgfsys@defobject{currentmarker}{\pgfqpoint{-0.048611in}{0.000000in}}{\pgfqpoint{-0.000000in}{0.000000in}}{%
\pgfpathmoveto{\pgfqpoint{-0.000000in}{0.000000in}}%
\pgfpathlineto{\pgfqpoint{-0.048611in}{0.000000in}}%
\pgfusepath{stroke,fill}%
}%
\begin{pgfscope}%
\pgfsys@transformshift{0.758026in}{1.186807in}%
\pgfsys@useobject{currentmarker}{}%
\end{pgfscope}%
\end{pgfscope}%
\begin{pgfscope}%
\definecolor{textcolor}{rgb}{0.000000,0.000000,0.000000}%
\pgfsetstrokecolor{textcolor}%
\pgfsetfillcolor{textcolor}%
\pgftext[x=0.413889in, y=1.138581in, left, base]{\color{textcolor}\rmfamily\fontsize{10.000000}{12.000000}\selectfont \(\displaystyle {\ensuremath{-}90}\)}%
\end{pgfscope}%
\begin{pgfscope}%
\pgfpathrectangle{\pgfqpoint{0.758026in}{0.565123in}}{\pgfqpoint{5.391974in}{2.784877in}}%
\pgfusepath{clip}%
\pgfsetbuttcap%
\pgfsetroundjoin%
\pgfsetlinewidth{0.501875pt}%
\definecolor{currentstroke}{rgb}{0.690196,0.690196,0.690196}%
\pgfsetstrokecolor{currentstroke}%
\pgfsetstrokeopacity{0.700000}%
\pgfsetdash{{1.850000pt}{0.800000pt}}{0.000000pt}%
\pgfpathmoveto{\pgfqpoint{0.758026in}{1.502407in}}%
\pgfpathlineto{\pgfqpoint{6.150000in}{1.502407in}}%
\pgfusepath{stroke}%
\end{pgfscope}%
\begin{pgfscope}%
\pgfsetbuttcap%
\pgfsetroundjoin%
\definecolor{currentfill}{rgb}{0.000000,0.000000,0.000000}%
\pgfsetfillcolor{currentfill}%
\pgfsetlinewidth{0.803000pt}%
\definecolor{currentstroke}{rgb}{0.000000,0.000000,0.000000}%
\pgfsetstrokecolor{currentstroke}%
\pgfsetdash{}{0pt}%
\pgfsys@defobject{currentmarker}{\pgfqpoint{-0.048611in}{0.000000in}}{\pgfqpoint{-0.000000in}{0.000000in}}{%
\pgfpathmoveto{\pgfqpoint{-0.000000in}{0.000000in}}%
\pgfpathlineto{\pgfqpoint{-0.048611in}{0.000000in}}%
\pgfusepath{stroke,fill}%
}%
\begin{pgfscope}%
\pgfsys@transformshift{0.758026in}{1.502407in}%
\pgfsys@useobject{currentmarker}{}%
\end{pgfscope}%
\end{pgfscope}%
\begin{pgfscope}%
\definecolor{textcolor}{rgb}{0.000000,0.000000,0.000000}%
\pgfsetstrokecolor{textcolor}%
\pgfsetfillcolor{textcolor}%
\pgftext[x=0.413889in, y=1.454182in, left, base]{\color{textcolor}\rmfamily\fontsize{10.000000}{12.000000}\selectfont \(\displaystyle {\ensuremath{-}80}\)}%
\end{pgfscope}%
\begin{pgfscope}%
\pgfpathrectangle{\pgfqpoint{0.758026in}{0.565123in}}{\pgfqpoint{5.391974in}{2.784877in}}%
\pgfusepath{clip}%
\pgfsetbuttcap%
\pgfsetroundjoin%
\pgfsetlinewidth{0.501875pt}%
\definecolor{currentstroke}{rgb}{0.690196,0.690196,0.690196}%
\pgfsetstrokecolor{currentstroke}%
\pgfsetstrokeopacity{0.700000}%
\pgfsetdash{{1.850000pt}{0.800000pt}}{0.000000pt}%
\pgfpathmoveto{\pgfqpoint{0.758026in}{1.818007in}}%
\pgfpathlineto{\pgfqpoint{6.150000in}{1.818007in}}%
\pgfusepath{stroke}%
\end{pgfscope}%
\begin{pgfscope}%
\pgfsetbuttcap%
\pgfsetroundjoin%
\definecolor{currentfill}{rgb}{0.000000,0.000000,0.000000}%
\pgfsetfillcolor{currentfill}%
\pgfsetlinewidth{0.803000pt}%
\definecolor{currentstroke}{rgb}{0.000000,0.000000,0.000000}%
\pgfsetstrokecolor{currentstroke}%
\pgfsetdash{}{0pt}%
\pgfsys@defobject{currentmarker}{\pgfqpoint{-0.048611in}{0.000000in}}{\pgfqpoint{-0.000000in}{0.000000in}}{%
\pgfpathmoveto{\pgfqpoint{-0.000000in}{0.000000in}}%
\pgfpathlineto{\pgfqpoint{-0.048611in}{0.000000in}}%
\pgfusepath{stroke,fill}%
}%
\begin{pgfscope}%
\pgfsys@transformshift{0.758026in}{1.818007in}%
\pgfsys@useobject{currentmarker}{}%
\end{pgfscope}%
\end{pgfscope}%
\begin{pgfscope}%
\definecolor{textcolor}{rgb}{0.000000,0.000000,0.000000}%
\pgfsetstrokecolor{textcolor}%
\pgfsetfillcolor{textcolor}%
\pgftext[x=0.413889in, y=1.769782in, left, base]{\color{textcolor}\rmfamily\fontsize{10.000000}{12.000000}\selectfont \(\displaystyle {\ensuremath{-}70}\)}%
\end{pgfscope}%
\begin{pgfscope}%
\pgfpathrectangle{\pgfqpoint{0.758026in}{0.565123in}}{\pgfqpoint{5.391974in}{2.784877in}}%
\pgfusepath{clip}%
\pgfsetbuttcap%
\pgfsetroundjoin%
\pgfsetlinewidth{0.501875pt}%
\definecolor{currentstroke}{rgb}{0.690196,0.690196,0.690196}%
\pgfsetstrokecolor{currentstroke}%
\pgfsetstrokeopacity{0.700000}%
\pgfsetdash{{1.850000pt}{0.800000pt}}{0.000000pt}%
\pgfpathmoveto{\pgfqpoint{0.758026in}{2.133607in}}%
\pgfpathlineto{\pgfqpoint{6.150000in}{2.133607in}}%
\pgfusepath{stroke}%
\end{pgfscope}%
\begin{pgfscope}%
\pgfsetbuttcap%
\pgfsetroundjoin%
\definecolor{currentfill}{rgb}{0.000000,0.000000,0.000000}%
\pgfsetfillcolor{currentfill}%
\pgfsetlinewidth{0.803000pt}%
\definecolor{currentstroke}{rgb}{0.000000,0.000000,0.000000}%
\pgfsetstrokecolor{currentstroke}%
\pgfsetdash{}{0pt}%
\pgfsys@defobject{currentmarker}{\pgfqpoint{-0.048611in}{0.000000in}}{\pgfqpoint{-0.000000in}{0.000000in}}{%
\pgfpathmoveto{\pgfqpoint{-0.000000in}{0.000000in}}%
\pgfpathlineto{\pgfqpoint{-0.048611in}{0.000000in}}%
\pgfusepath{stroke,fill}%
}%
\begin{pgfscope}%
\pgfsys@transformshift{0.758026in}{2.133607in}%
\pgfsys@useobject{currentmarker}{}%
\end{pgfscope}%
\end{pgfscope}%
\begin{pgfscope}%
\definecolor{textcolor}{rgb}{0.000000,0.000000,0.000000}%
\pgfsetstrokecolor{textcolor}%
\pgfsetfillcolor{textcolor}%
\pgftext[x=0.413889in, y=2.085382in, left, base]{\color{textcolor}\rmfamily\fontsize{10.000000}{12.000000}\selectfont \(\displaystyle {\ensuremath{-}60}\)}%
\end{pgfscope}%
\begin{pgfscope}%
\pgfpathrectangle{\pgfqpoint{0.758026in}{0.565123in}}{\pgfqpoint{5.391974in}{2.784877in}}%
\pgfusepath{clip}%
\pgfsetbuttcap%
\pgfsetroundjoin%
\pgfsetlinewidth{0.501875pt}%
\definecolor{currentstroke}{rgb}{0.690196,0.690196,0.690196}%
\pgfsetstrokecolor{currentstroke}%
\pgfsetstrokeopacity{0.700000}%
\pgfsetdash{{1.850000pt}{0.800000pt}}{0.000000pt}%
\pgfpathmoveto{\pgfqpoint{0.758026in}{2.449208in}}%
\pgfpathlineto{\pgfqpoint{6.150000in}{2.449208in}}%
\pgfusepath{stroke}%
\end{pgfscope}%
\begin{pgfscope}%
\pgfsetbuttcap%
\pgfsetroundjoin%
\definecolor{currentfill}{rgb}{0.000000,0.000000,0.000000}%
\pgfsetfillcolor{currentfill}%
\pgfsetlinewidth{0.803000pt}%
\definecolor{currentstroke}{rgb}{0.000000,0.000000,0.000000}%
\pgfsetstrokecolor{currentstroke}%
\pgfsetdash{}{0pt}%
\pgfsys@defobject{currentmarker}{\pgfqpoint{-0.048611in}{0.000000in}}{\pgfqpoint{-0.000000in}{0.000000in}}{%
\pgfpathmoveto{\pgfqpoint{-0.000000in}{0.000000in}}%
\pgfpathlineto{\pgfqpoint{-0.048611in}{0.000000in}}%
\pgfusepath{stroke,fill}%
}%
\begin{pgfscope}%
\pgfsys@transformshift{0.758026in}{2.449208in}%
\pgfsys@useobject{currentmarker}{}%
\end{pgfscope}%
\end{pgfscope}%
\begin{pgfscope}%
\definecolor{textcolor}{rgb}{0.000000,0.000000,0.000000}%
\pgfsetstrokecolor{textcolor}%
\pgfsetfillcolor{textcolor}%
\pgftext[x=0.413889in, y=2.400982in, left, base]{\color{textcolor}\rmfamily\fontsize{10.000000}{12.000000}\selectfont \(\displaystyle {\ensuremath{-}50}\)}%
\end{pgfscope}%
\begin{pgfscope}%
\pgfpathrectangle{\pgfqpoint{0.758026in}{0.565123in}}{\pgfqpoint{5.391974in}{2.784877in}}%
\pgfusepath{clip}%
\pgfsetbuttcap%
\pgfsetroundjoin%
\pgfsetlinewidth{0.501875pt}%
\definecolor{currentstroke}{rgb}{0.690196,0.690196,0.690196}%
\pgfsetstrokecolor{currentstroke}%
\pgfsetstrokeopacity{0.700000}%
\pgfsetdash{{1.850000pt}{0.800000pt}}{0.000000pt}%
\pgfpathmoveto{\pgfqpoint{0.758026in}{2.764808in}}%
\pgfpathlineto{\pgfqpoint{6.150000in}{2.764808in}}%
\pgfusepath{stroke}%
\end{pgfscope}%
\begin{pgfscope}%
\pgfsetbuttcap%
\pgfsetroundjoin%
\definecolor{currentfill}{rgb}{0.000000,0.000000,0.000000}%
\pgfsetfillcolor{currentfill}%
\pgfsetlinewidth{0.803000pt}%
\definecolor{currentstroke}{rgb}{0.000000,0.000000,0.000000}%
\pgfsetstrokecolor{currentstroke}%
\pgfsetdash{}{0pt}%
\pgfsys@defobject{currentmarker}{\pgfqpoint{-0.048611in}{0.000000in}}{\pgfqpoint{-0.000000in}{0.000000in}}{%
\pgfpathmoveto{\pgfqpoint{-0.000000in}{0.000000in}}%
\pgfpathlineto{\pgfqpoint{-0.048611in}{0.000000in}}%
\pgfusepath{stroke,fill}%
}%
\begin{pgfscope}%
\pgfsys@transformshift{0.758026in}{2.764808in}%
\pgfsys@useobject{currentmarker}{}%
\end{pgfscope}%
\end{pgfscope}%
\begin{pgfscope}%
\definecolor{textcolor}{rgb}{0.000000,0.000000,0.000000}%
\pgfsetstrokecolor{textcolor}%
\pgfsetfillcolor{textcolor}%
\pgftext[x=0.413889in, y=2.716583in, left, base]{\color{textcolor}\rmfamily\fontsize{10.000000}{12.000000}\selectfont \(\displaystyle {\ensuremath{-}40}\)}%
\end{pgfscope}%
\begin{pgfscope}%
\pgfpathrectangle{\pgfqpoint{0.758026in}{0.565123in}}{\pgfqpoint{5.391974in}{2.784877in}}%
\pgfusepath{clip}%
\pgfsetbuttcap%
\pgfsetroundjoin%
\pgfsetlinewidth{0.501875pt}%
\definecolor{currentstroke}{rgb}{0.690196,0.690196,0.690196}%
\pgfsetstrokecolor{currentstroke}%
\pgfsetstrokeopacity{0.700000}%
\pgfsetdash{{1.850000pt}{0.800000pt}}{0.000000pt}%
\pgfpathmoveto{\pgfqpoint{0.758026in}{3.080408in}}%
\pgfpathlineto{\pgfqpoint{6.150000in}{3.080408in}}%
\pgfusepath{stroke}%
\end{pgfscope}%
\begin{pgfscope}%
\pgfsetbuttcap%
\pgfsetroundjoin%
\definecolor{currentfill}{rgb}{0.000000,0.000000,0.000000}%
\pgfsetfillcolor{currentfill}%
\pgfsetlinewidth{0.803000pt}%
\definecolor{currentstroke}{rgb}{0.000000,0.000000,0.000000}%
\pgfsetstrokecolor{currentstroke}%
\pgfsetdash{}{0pt}%
\pgfsys@defobject{currentmarker}{\pgfqpoint{-0.048611in}{0.000000in}}{\pgfqpoint{-0.000000in}{0.000000in}}{%
\pgfpathmoveto{\pgfqpoint{-0.000000in}{0.000000in}}%
\pgfpathlineto{\pgfqpoint{-0.048611in}{0.000000in}}%
\pgfusepath{stroke,fill}%
}%
\begin{pgfscope}%
\pgfsys@transformshift{0.758026in}{3.080408in}%
\pgfsys@useobject{currentmarker}{}%
\end{pgfscope}%
\end{pgfscope}%
\begin{pgfscope}%
\definecolor{textcolor}{rgb}{0.000000,0.000000,0.000000}%
\pgfsetstrokecolor{textcolor}%
\pgfsetfillcolor{textcolor}%
\pgftext[x=0.413889in, y=3.032183in, left, base]{\color{textcolor}\rmfamily\fontsize{10.000000}{12.000000}\selectfont \(\displaystyle {\ensuremath{-}30}\)}%
\end{pgfscope}%
\begin{pgfscope}%
\definecolor{textcolor}{rgb}{0.000000,0.000000,0.000000}%
\pgfsetstrokecolor{textcolor}%
\pgfsetfillcolor{textcolor}%
\pgftext[x=0.288889in,y=1.957562in,,bottom,rotate=90.000000]{\color{textcolor}\rmfamily\fontsize{10.000000}{12.000000}\selectfont Amplitude (dB)}%
\end{pgfscope}%
\begin{pgfscope}%
\pgfpathrectangle{\pgfqpoint{0.758026in}{0.565123in}}{\pgfqpoint{5.391974in}{2.784877in}}%
\pgfusepath{clip}%
\pgfsetrectcap%
\pgfsetroundjoin%
\pgfsetlinewidth{1.003750pt}%
\definecolor{currentstroke}{rgb}{0.000000,0.000000,0.000000}%
\pgfsetstrokecolor{currentstroke}%
\pgfsetdash{}{0pt}%
\pgfpathmoveto{\pgfqpoint{0.758026in}{1.636537in}}%
\pgfpathlineto{\pgfqpoint{0.763418in}{1.620757in}}%
\pgfpathlineto{\pgfqpoint{0.768810in}{1.477751in}}%
\pgfpathlineto{\pgfqpoint{0.774202in}{1.525091in}}%
\pgfpathlineto{\pgfqpoint{0.779594in}{1.429424in}}%
\pgfpathlineto{\pgfqpoint{0.784986in}{1.538898in}}%
\pgfpathlineto{\pgfqpoint{0.790378in}{1.475778in}}%
\pgfpathlineto{\pgfqpoint{0.795770in}{1.616812in}}%
\pgfpathlineto{\pgfqpoint{0.801161in}{1.489586in}}%
\pgfpathlineto{\pgfqpoint{0.806553in}{1.472819in}}%
\pgfpathlineto{\pgfqpoint{0.811945in}{1.472819in}}%
\pgfpathlineto{\pgfqpoint{0.817337in}{1.582293in}}%
\pgfpathlineto{\pgfqpoint{0.822729in}{1.423507in}}%
\pgfpathlineto{\pgfqpoint{0.828121in}{1.612867in}}%
\pgfpathlineto{\pgfqpoint{0.838905in}{1.468874in}}%
\pgfpathlineto{\pgfqpoint{0.844297in}{1.484654in}}%
\pgfpathlineto{\pgfqpoint{0.849689in}{1.562568in}}%
\pgfpathlineto{\pgfqpoint{0.855081in}{1.466902in}}%
\pgfpathlineto{\pgfqpoint{0.860473in}{1.608922in}}%
\pgfpathlineto{\pgfqpoint{0.865865in}{1.465916in}}%
\pgfpathlineto{\pgfqpoint{0.871257in}{1.717410in}}%
\pgfpathlineto{\pgfqpoint{0.876649in}{1.512269in}}%
\pgfpathlineto{\pgfqpoint{0.882041in}{1.527063in}}%
\pgfpathlineto{\pgfqpoint{0.887433in}{1.462957in}}%
\pgfpathlineto{\pgfqpoint{0.892825in}{1.604977in}}%
\pgfpathlineto{\pgfqpoint{0.898217in}{1.509311in}}%
\pgfpathlineto{\pgfqpoint{0.903609in}{1.460984in}}%
\pgfpathlineto{\pgfqpoint{0.909001in}{1.618784in}}%
\pgfpathlineto{\pgfqpoint{0.914393in}{1.412658in}}%
\pgfpathlineto{\pgfqpoint{0.919785in}{1.522132in}}%
\pgfpathlineto{\pgfqpoint{0.925177in}{1.663166in}}%
\pgfpathlineto{\pgfqpoint{0.930569in}{1.521146in}}%
\pgfpathlineto{\pgfqpoint{0.935961in}{1.520159in}}%
\pgfpathlineto{\pgfqpoint{0.941353in}{1.471833in}}%
\pgfpathlineto{\pgfqpoint{0.946745in}{1.913674in}}%
\pgfpathlineto{\pgfqpoint{0.952137in}{1.502407in}}%
\pgfpathlineto{\pgfqpoint{0.957529in}{1.769681in}}%
\pgfpathlineto{\pgfqpoint{0.962921in}{1.517201in}}%
\pgfpathlineto{\pgfqpoint{0.968313in}{1.500434in}}%
\pgfpathlineto{\pgfqpoint{0.973705in}{1.499448in}}%
\pgfpathlineto{\pgfqpoint{0.979097in}{1.641468in}}%
\pgfpathlineto{\pgfqpoint{0.984489in}{1.514242in}}%
\pgfpathlineto{\pgfqpoint{0.989881in}{1.734176in}}%
\pgfpathlineto{\pgfqpoint{0.995273in}{1.592156in}}%
\pgfpathlineto{\pgfqpoint{1.000665in}{1.654290in}}%
\pgfpathlineto{\pgfqpoint{1.006057in}{1.448163in}}%
\pgfpathlineto{\pgfqpoint{1.011448in}{1.448163in}}%
\pgfpathlineto{\pgfqpoint{1.016840in}{1.399837in}}%
\pgfpathlineto{\pgfqpoint{1.022232in}{1.493531in}}%
\pgfpathlineto{\pgfqpoint{1.027624in}{1.666125in}}%
\pgfpathlineto{\pgfqpoint{1.033016in}{1.697685in}}%
\pgfpathlineto{\pgfqpoint{1.038408in}{1.491558in}}%
\pgfpathlineto{\pgfqpoint{1.043800in}{1.490572in}}%
\pgfpathlineto{\pgfqpoint{1.049192in}{1.774612in}}%
\pgfpathlineto{\pgfqpoint{1.054584in}{1.394906in}}%
\pgfpathlineto{\pgfqpoint{1.059976in}{1.599059in}}%
\pgfpathlineto{\pgfqpoint{1.065368in}{1.599059in}}%
\pgfpathlineto{\pgfqpoint{1.070760in}{2.071474in}}%
\pgfpathlineto{\pgfqpoint{1.076152in}{2.133607in}}%
\pgfpathlineto{\pgfqpoint{1.081544in}{1.818007in}}%
\pgfpathlineto{\pgfqpoint{1.086936in}{1.895921in}}%
\pgfpathlineto{\pgfqpoint{1.092328in}{2.005395in}}%
\pgfpathlineto{\pgfqpoint{1.097720in}{1.658235in}}%
\pgfpathlineto{\pgfqpoint{1.103112in}{1.499448in}}%
\pgfpathlineto{\pgfqpoint{1.108504in}{1.577362in}}%
\pgfpathlineto{\pgfqpoint{1.113896in}{1.719382in}}%
\pgfpathlineto{\pgfqpoint{1.119288in}{1.529036in}}%
\pgfpathlineto{\pgfqpoint{1.124680in}{1.433369in}}%
\pgfpathlineto{\pgfqpoint{1.130072in}{1.591169in}}%
\pgfpathlineto{\pgfqpoint{1.135464in}{1.479723in}}%
\pgfpathlineto{\pgfqpoint{1.140856in}{1.541857in}}%
\pgfpathlineto{\pgfqpoint{1.146248in}{1.526077in}}%
\pgfpathlineto{\pgfqpoint{1.151640in}{1.525091in}}%
\pgfpathlineto{\pgfqpoint{1.157032in}{1.476764in}}%
\pgfpathlineto{\pgfqpoint{1.162424in}{1.666125in}}%
\pgfpathlineto{\pgfqpoint{1.167816in}{1.491558in}}%
\pgfpathlineto{\pgfqpoint{1.173208in}{1.537912in}}%
\pgfpathlineto{\pgfqpoint{1.178600in}{1.474792in}}%
\pgfpathlineto{\pgfqpoint{1.183992in}{2.483727in}}%
\pgfpathlineto{\pgfqpoint{1.189384in}{2.987701in}}%
\pgfpathlineto{\pgfqpoint{1.194776in}{3.191855in}}%
\pgfpathlineto{\pgfqpoint{1.200168in}{3.223415in}}%
\pgfpathlineto{\pgfqpoint{1.205560in}{3.111968in}}%
\pgfpathlineto{\pgfqpoint{1.210952in}{2.748042in}}%
\pgfpathlineto{\pgfqpoint{1.216344in}{1.832801in}}%
\pgfpathlineto{\pgfqpoint{1.221735in}{1.516214in}}%
\pgfpathlineto{\pgfqpoint{1.227127in}{1.720368in}}%
\pgfpathlineto{\pgfqpoint{1.232519in}{1.799268in}}%
\pgfpathlineto{\pgfqpoint{1.237911in}{1.419562in}}%
\pgfpathlineto{\pgfqpoint{1.243303in}{1.481696in}}%
\pgfpathlineto{\pgfqpoint{1.248695in}{1.418576in}}%
\pgfpathlineto{\pgfqpoint{1.254087in}{1.370249in}}%
\pgfpathlineto{\pgfqpoint{1.259479in}{1.984684in}}%
\pgfpathlineto{\pgfqpoint{1.264871in}{1.811103in}}%
\pgfpathlineto{\pgfqpoint{1.270263in}{1.368277in}}%
\pgfpathlineto{\pgfqpoint{1.297223in}{1.364332in}}%
\pgfpathlineto{\pgfqpoint{1.308007in}{1.363345in}}%
\pgfpathlineto{\pgfqpoint{1.329575in}{1.360387in}}%
\pgfpathlineto{\pgfqpoint{1.340359in}{1.359400in}}%
\pgfpathlineto{\pgfqpoint{1.361927in}{1.356442in}}%
\pgfpathlineto{\pgfqpoint{1.372711in}{1.355455in}}%
\pgfpathlineto{\pgfqpoint{1.383495in}{1.354469in}}%
\pgfpathlineto{\pgfqpoint{1.388887in}{1.495503in}}%
\pgfpathlineto{\pgfqpoint{1.394279in}{1.352497in}}%
\pgfpathlineto{\pgfqpoint{1.399671in}{1.352497in}}%
\pgfpathlineto{\pgfqpoint{1.405063in}{1.793351in}}%
\pgfpathlineto{\pgfqpoint{1.410455in}{1.161164in}}%
\pgfpathlineto{\pgfqpoint{1.421239in}{1.160178in}}%
\pgfpathlineto{\pgfqpoint{1.426631in}{1.269652in}}%
\pgfpathlineto{\pgfqpoint{1.432022in}{1.159192in}}%
\pgfpathlineto{\pgfqpoint{1.437414in}{1.205545in}}%
\pgfpathlineto{\pgfqpoint{1.442806in}{1.157219in}}%
\pgfpathlineto{\pgfqpoint{1.448198in}{1.204559in}}%
\pgfpathlineto{\pgfqpoint{1.453590in}{1.219353in}}%
\pgfpathlineto{\pgfqpoint{1.458982in}{1.202587in}}%
\pgfpathlineto{\pgfqpoint{1.464374in}{1.154260in}}%
\pgfpathlineto{\pgfqpoint{1.475158in}{1.153274in}}%
\pgfpathlineto{\pgfqpoint{1.496726in}{1.150315in}}%
\pgfpathlineto{\pgfqpoint{1.507510in}{1.149329in}}%
\pgfpathlineto{\pgfqpoint{1.529078in}{1.146370in}}%
\pgfpathlineto{\pgfqpoint{1.534470in}{1.099030in}}%
\pgfpathlineto{\pgfqpoint{1.539862in}{1.208504in}}%
\pgfpathlineto{\pgfqpoint{1.545254in}{1.160178in}}%
\pgfpathlineto{\pgfqpoint{1.550646in}{1.049718in}}%
\pgfpathlineto{\pgfqpoint{1.556038in}{1.096072in}}%
\pgfpathlineto{\pgfqpoint{1.572214in}{1.094099in}}%
\pgfpathlineto{\pgfqpoint{1.577606in}{1.282473in}}%
\pgfpathlineto{\pgfqpoint{1.582998in}{1.139467in}}%
\pgfpathlineto{\pgfqpoint{1.588390in}{1.139467in}}%
\pgfpathlineto{\pgfqpoint{1.593782in}{1.043800in}}%
\pgfpathlineto{\pgfqpoint{1.599174in}{1.042814in}}%
\pgfpathlineto{\pgfqpoint{1.604566in}{1.090154in}}%
\pgfpathlineto{\pgfqpoint{1.609958in}{1.104948in}}%
\pgfpathlineto{\pgfqpoint{1.615350in}{1.103962in}}%
\pgfpathlineto{\pgfqpoint{1.620742in}{1.040841in}}%
\pgfpathlineto{\pgfqpoint{1.626134in}{1.181875in}}%
\pgfpathlineto{\pgfqpoint{1.631526in}{1.133549in}}%
\pgfpathlineto{\pgfqpoint{1.636918in}{1.101989in}}%
\pgfpathlineto{\pgfqpoint{1.642309in}{1.148343in}}%
\pgfpathlineto{\pgfqpoint{1.647701in}{1.115797in}}%
\pgfpathlineto{\pgfqpoint{1.653093in}{1.068456in}}%
\pgfpathlineto{\pgfqpoint{1.658485in}{1.256830in}}%
\pgfpathlineto{\pgfqpoint{1.663877in}{1.050704in}}%
\pgfpathlineto{\pgfqpoint{1.669269in}{1.050704in}}%
\pgfpathlineto{\pgfqpoint{1.674661in}{1.333758in}}%
\pgfpathlineto{\pgfqpoint{1.680053in}{1.490572in}}%
\pgfpathlineto{\pgfqpoint{1.685445in}{1.110865in}}%
\pgfpathlineto{\pgfqpoint{1.690837in}{1.189765in}}%
\pgfpathlineto{\pgfqpoint{1.696229in}{1.204559in}}%
\pgfpathlineto{\pgfqpoint{1.701621in}{1.093113in}}%
\pgfpathlineto{\pgfqpoint{1.707013in}{0.951093in}}%
\pgfpathlineto{\pgfqpoint{1.717797in}{1.375181in}}%
\pgfpathlineto{\pgfqpoint{1.723189in}{0.996460in}}%
\pgfpathlineto{\pgfqpoint{1.728581in}{1.042814in}}%
\pgfpathlineto{\pgfqpoint{1.733973in}{1.262748in}}%
\pgfpathlineto{\pgfqpoint{1.739365in}{1.183848in}}%
\pgfpathlineto{\pgfqpoint{1.750149in}{1.118755in}}%
\pgfpathlineto{\pgfqpoint{1.755541in}{1.118755in}}%
\pgfpathlineto{\pgfqpoint{1.760933in}{1.071415in}}%
\pgfpathlineto{\pgfqpoint{1.766325in}{1.197655in}}%
\pgfpathlineto{\pgfqpoint{1.771717in}{1.244995in}}%
\pgfpathlineto{\pgfqpoint{1.777109in}{1.276555in}}%
\pgfpathlineto{\pgfqpoint{1.782501in}{1.134535in}}%
\pgfpathlineto{\pgfqpoint{1.787893in}{1.355455in}}%
\pgfpathlineto{\pgfqpoint{1.793285in}{1.260775in}}%
\pgfpathlineto{\pgfqpoint{1.798677in}{1.087195in}}%
\pgfpathlineto{\pgfqpoint{1.804069in}{1.481696in}}%
\pgfpathlineto{\pgfqpoint{1.809461in}{1.134535in}}%
\pgfpathlineto{\pgfqpoint{1.814853in}{1.276555in}}%
\pgfpathlineto{\pgfqpoint{1.820245in}{1.339675in}}%
\pgfpathlineto{\pgfqpoint{1.825637in}{1.355455in}}%
\pgfpathlineto{\pgfqpoint{1.831029in}{1.276555in}}%
\pgfpathlineto{\pgfqpoint{1.836421in}{1.355455in}}%
\pgfpathlineto{\pgfqpoint{1.841813in}{1.371236in}}%
\pgfpathlineto{\pgfqpoint{1.847205in}{1.434356in}}%
\pgfpathlineto{\pgfqpoint{1.852596in}{1.244995in}}%
\pgfpathlineto{\pgfqpoint{1.857988in}{1.450136in}}%
\pgfpathlineto{\pgfqpoint{1.863380in}{1.229215in}}%
\pgfpathlineto{\pgfqpoint{1.868772in}{1.292335in}}%
\pgfpathlineto{\pgfqpoint{1.874164in}{1.292335in}}%
\pgfpathlineto{\pgfqpoint{1.879556in}{1.308115in}}%
\pgfpathlineto{\pgfqpoint{1.884948in}{1.260775in}}%
\pgfpathlineto{\pgfqpoint{1.890340in}{1.292335in}}%
\pgfpathlineto{\pgfqpoint{1.895732in}{1.276555in}}%
\pgfpathlineto{\pgfqpoint{1.901124in}{1.244995in}}%
\pgfpathlineto{\pgfqpoint{1.906516in}{1.102975in}}%
\pgfpathlineto{\pgfqpoint{1.911908in}{1.197655in}}%
\pgfpathlineto{\pgfqpoint{1.917300in}{1.166095in}}%
\pgfpathlineto{\pgfqpoint{1.922692in}{1.150315in}}%
\pgfpathlineto{\pgfqpoint{1.928084in}{1.213435in}}%
\pgfpathlineto{\pgfqpoint{1.933476in}{1.118755in}}%
\pgfpathlineto{\pgfqpoint{1.938868in}{1.134535in}}%
\pgfpathlineto{\pgfqpoint{1.944260in}{1.165109in}}%
\pgfpathlineto{\pgfqpoint{1.949652in}{1.354469in}}%
\pgfpathlineto{\pgfqpoint{1.955044in}{1.117769in}}%
\pgfpathlineto{\pgfqpoint{1.960436in}{1.101989in}}%
\pgfpathlineto{\pgfqpoint{1.965828in}{1.149329in}}%
\pgfpathlineto{\pgfqpoint{1.971220in}{1.133549in}}%
\pgfpathlineto{\pgfqpoint{1.976612in}{0.802169in}}%
\pgfpathlineto{\pgfqpoint{1.982004in}{1.149329in}}%
\pgfpathlineto{\pgfqpoint{1.987396in}{0.975749in}}%
\pgfpathlineto{\pgfqpoint{1.992788in}{0.691709in}}%
\pgfpathlineto{\pgfqpoint{1.998180in}{1.101989in}}%
\pgfpathlineto{\pgfqpoint{2.003572in}{1.101989in}}%
\pgfpathlineto{\pgfqpoint{2.008964in}{1.070429in}}%
\pgfpathlineto{\pgfqpoint{2.014356in}{0.912629in}}%
\pgfpathlineto{\pgfqpoint{2.019748in}{1.196669in}}%
\pgfpathlineto{\pgfqpoint{2.025140in}{1.054649in}}%
\pgfpathlineto{\pgfqpoint{2.030532in}{1.543829in}}%
\pgfpathlineto{\pgfqpoint{2.035924in}{1.196669in}}%
\pgfpathlineto{\pgfqpoint{2.041316in}{1.244009in}}%
\pgfpathlineto{\pgfqpoint{2.046708in}{1.117769in}}%
\pgfpathlineto{\pgfqpoint{2.052100in}{1.228229in}}%
\pgfpathlineto{\pgfqpoint{2.057492in}{1.149329in}}%
\pgfpathlineto{\pgfqpoint{2.062883in}{1.133549in}}%
\pgfpathlineto{\pgfqpoint{2.068275in}{1.070429in}}%
\pgfpathlineto{\pgfqpoint{2.079059in}{1.196669in}}%
\pgfpathlineto{\pgfqpoint{2.084451in}{1.212449in}}%
\pgfpathlineto{\pgfqpoint{2.089843in}{1.165109in}}%
\pgfpathlineto{\pgfqpoint{2.095235in}{1.086209in}}%
\pgfpathlineto{\pgfqpoint{2.100627in}{1.101989in}}%
\pgfpathlineto{\pgfqpoint{2.106019in}{1.196669in}}%
\pgfpathlineto{\pgfqpoint{2.111411in}{1.259789in}}%
\pgfpathlineto{\pgfqpoint{2.116803in}{1.165109in}}%
\pgfpathlineto{\pgfqpoint{2.122195in}{1.307129in}}%
\pgfpathlineto{\pgfqpoint{2.127587in}{1.116783in}}%
\pgfpathlineto{\pgfqpoint{2.132979in}{1.243023in}}%
\pgfpathlineto{\pgfqpoint{2.138371in}{1.306143in}}%
\pgfpathlineto{\pgfqpoint{2.143763in}{1.243023in}}%
\pgfpathlineto{\pgfqpoint{2.149155in}{1.243023in}}%
\pgfpathlineto{\pgfqpoint{2.154547in}{1.179903in}}%
\pgfpathlineto{\pgfqpoint{2.159939in}{1.085223in}}%
\pgfpathlineto{\pgfqpoint{2.165331in}{1.211463in}}%
\pgfpathlineto{\pgfqpoint{2.170723in}{1.164123in}}%
\pgfpathlineto{\pgfqpoint{2.176115in}{1.590183in}}%
\pgfpathlineto{\pgfqpoint{2.181507in}{1.053663in}}%
\pgfpathlineto{\pgfqpoint{2.186899in}{1.006323in}}%
\pgfpathlineto{\pgfqpoint{2.192291in}{1.227243in}}%
\pgfpathlineto{\pgfqpoint{2.197683in}{1.195683in}}%
\pgfpathlineto{\pgfqpoint{2.208467in}{0.990543in}}%
\pgfpathlineto{\pgfqpoint{2.213859in}{1.243023in}}%
\pgfpathlineto{\pgfqpoint{2.219251in}{1.085223in}}%
\pgfpathlineto{\pgfqpoint{2.224643in}{1.211463in}}%
\pgfpathlineto{\pgfqpoint{2.230035in}{1.258803in}}%
\pgfpathlineto{\pgfqpoint{2.235427in}{1.006323in}}%
\pgfpathlineto{\pgfqpoint{2.240819in}{1.227243in}}%
\pgfpathlineto{\pgfqpoint{2.246211in}{1.243023in}}%
\pgfpathlineto{\pgfqpoint{2.251603in}{1.132563in}}%
\pgfpathlineto{\pgfqpoint{2.256995in}{1.069443in}}%
\pgfpathlineto{\pgfqpoint{2.262387in}{1.132563in}}%
\pgfpathlineto{\pgfqpoint{2.267778in}{1.164123in}}%
\pgfpathlineto{\pgfqpoint{2.273170in}{1.037883in}}%
\pgfpathlineto{\pgfqpoint{2.283954in}{1.306143in}}%
\pgfpathlineto{\pgfqpoint{2.294738in}{0.974763in}}%
\pgfpathlineto{\pgfqpoint{2.300130in}{1.227243in}}%
\pgfpathlineto{\pgfqpoint{2.305522in}{1.306143in}}%
\pgfpathlineto{\pgfqpoint{2.310914in}{1.258803in}}%
\pgfpathlineto{\pgfqpoint{2.316306in}{1.164123in}}%
\pgfpathlineto{\pgfqpoint{2.321698in}{1.273597in}}%
\pgfpathlineto{\pgfqpoint{2.327090in}{1.100017in}}%
\pgfpathlineto{\pgfqpoint{2.332482in}{1.100017in}}%
\pgfpathlineto{\pgfqpoint{2.343266in}{1.068456in}}%
\pgfpathlineto{\pgfqpoint{2.348658in}{1.210477in}}%
\pgfpathlineto{\pgfqpoint{2.359442in}{1.210477in}}%
\pgfpathlineto{\pgfqpoint{2.364834in}{1.115797in}}%
\pgfpathlineto{\pgfqpoint{2.370226in}{0.942216in}}%
\pgfpathlineto{\pgfqpoint{2.375618in}{0.973776in}}%
\pgfpathlineto{\pgfqpoint{2.381010in}{1.084236in}}%
\pgfpathlineto{\pgfqpoint{2.386402in}{1.226257in}}%
\pgfpathlineto{\pgfqpoint{2.391794in}{1.305157in}}%
\pgfpathlineto{\pgfqpoint{2.397186in}{0.942216in}}%
\pgfpathlineto{\pgfqpoint{2.402578in}{1.257817in}}%
\pgfpathlineto{\pgfqpoint{2.407970in}{1.100017in}}%
\pgfpathlineto{\pgfqpoint{2.413362in}{1.163137in}}%
\pgfpathlineto{\pgfqpoint{2.418754in}{1.257817in}}%
\pgfpathlineto{\pgfqpoint{2.424146in}{1.210477in}}%
\pgfpathlineto{\pgfqpoint{2.429538in}{1.131577in}}%
\pgfpathlineto{\pgfqpoint{2.434930in}{1.210477in}}%
\pgfpathlineto{\pgfqpoint{2.440322in}{1.036896in}}%
\pgfpathlineto{\pgfqpoint{2.445714in}{1.068456in}}%
\pgfpathlineto{\pgfqpoint{2.451106in}{1.052676in}}%
\pgfpathlineto{\pgfqpoint{2.456498in}{1.226257in}}%
\pgfpathlineto{\pgfqpoint{2.461890in}{1.052676in}}%
\pgfpathlineto{\pgfqpoint{2.467282in}{1.068456in}}%
\pgfpathlineto{\pgfqpoint{2.472674in}{1.226257in}}%
\pgfpathlineto{\pgfqpoint{2.478065in}{1.320937in}}%
\pgfpathlineto{\pgfqpoint{2.483457in}{1.178917in}}%
\pgfpathlineto{\pgfqpoint{2.488849in}{1.226257in}}%
\pgfpathlineto{\pgfqpoint{2.494241in}{1.289377in}}%
\pgfpathlineto{\pgfqpoint{2.499633in}{0.973776in}}%
\pgfpathlineto{\pgfqpoint{2.505025in}{1.100017in}}%
\pgfpathlineto{\pgfqpoint{2.510417in}{1.288390in}}%
\pgfpathlineto{\pgfqpoint{2.515809in}{1.099030in}}%
\pgfpathlineto{\pgfqpoint{2.521201in}{1.099030in}}%
\pgfpathlineto{\pgfqpoint{2.526593in}{1.051690in}}%
\pgfpathlineto{\pgfqpoint{2.531985in}{1.020130in}}%
\pgfpathlineto{\pgfqpoint{2.537377in}{1.241050in}}%
\pgfpathlineto{\pgfqpoint{2.542769in}{1.067470in}}%
\pgfpathlineto{\pgfqpoint{2.548161in}{1.225270in}}%
\pgfpathlineto{\pgfqpoint{2.553553in}{1.272610in}}%
\pgfpathlineto{\pgfqpoint{2.558945in}{1.146370in}}%
\pgfpathlineto{\pgfqpoint{2.564337in}{1.099030in}}%
\pgfpathlineto{\pgfqpoint{2.569729in}{1.146370in}}%
\pgfpathlineto{\pgfqpoint{2.575121in}{1.241050in}}%
\pgfpathlineto{\pgfqpoint{2.580513in}{1.067470in}}%
\pgfpathlineto{\pgfqpoint{2.585905in}{1.319950in}}%
\pgfpathlineto{\pgfqpoint{2.591297in}{1.067470in}}%
\pgfpathlineto{\pgfqpoint{2.596689in}{1.020130in}}%
\pgfpathlineto{\pgfqpoint{2.602081in}{0.878110in}}%
\pgfpathlineto{\pgfqpoint{2.607473in}{1.225270in}}%
\pgfpathlineto{\pgfqpoint{2.618257in}{1.083250in}}%
\pgfpathlineto{\pgfqpoint{2.623649in}{1.177930in}}%
\pgfpathlineto{\pgfqpoint{2.629041in}{1.099030in}}%
\pgfpathlineto{\pgfqpoint{2.634433in}{1.335730in}}%
\pgfpathlineto{\pgfqpoint{2.639825in}{1.035910in}}%
\pgfpathlineto{\pgfqpoint{2.645217in}{1.225270in}}%
\pgfpathlineto{\pgfqpoint{2.650609in}{1.256830in}}%
\pgfpathlineto{\pgfqpoint{2.656001in}{1.335730in}}%
\pgfpathlineto{\pgfqpoint{2.661393in}{1.083250in}}%
\pgfpathlineto{\pgfqpoint{2.666785in}{1.304170in}}%
\pgfpathlineto{\pgfqpoint{2.672177in}{1.177930in}}%
\pgfpathlineto{\pgfqpoint{2.677569in}{1.209490in}}%
\pgfpathlineto{\pgfqpoint{2.682961in}{1.162150in}}%
\pgfpathlineto{\pgfqpoint{2.688352in}{1.256830in}}%
\pgfpathlineto{\pgfqpoint{2.693744in}{0.814004in}}%
\pgfpathlineto{\pgfqpoint{2.699136in}{1.271624in}}%
\pgfpathlineto{\pgfqpoint{2.704528in}{1.129604in}}%
\pgfpathlineto{\pgfqpoint{2.709920in}{1.224284in}}%
\pgfpathlineto{\pgfqpoint{2.715312in}{0.971804in}}%
\pgfpathlineto{\pgfqpoint{2.720704in}{1.145384in}}%
\pgfpathlineto{\pgfqpoint{2.726096in}{1.050704in}}%
\pgfpathlineto{\pgfqpoint{2.731488in}{1.287404in}}%
\pgfpathlineto{\pgfqpoint{2.736880in}{1.334744in}}%
\pgfpathlineto{\pgfqpoint{2.742272in}{1.303184in}}%
\pgfpathlineto{\pgfqpoint{2.747664in}{1.287404in}}%
\pgfpathlineto{\pgfqpoint{2.753056in}{1.192724in}}%
\pgfpathlineto{\pgfqpoint{2.758448in}{1.176944in}}%
\pgfpathlineto{\pgfqpoint{2.763840in}{1.255844in}}%
\pgfpathlineto{\pgfqpoint{2.769232in}{1.366304in}}%
\pgfpathlineto{\pgfqpoint{2.774624in}{1.397864in}}%
\pgfpathlineto{\pgfqpoint{2.780016in}{1.255844in}}%
\pgfpathlineto{\pgfqpoint{2.785408in}{1.287404in}}%
\pgfpathlineto{\pgfqpoint{2.790800in}{1.176944in}}%
\pgfpathlineto{\pgfqpoint{2.796192in}{1.287404in}}%
\pgfpathlineto{\pgfqpoint{2.801584in}{1.271624in}}%
\pgfpathlineto{\pgfqpoint{2.806976in}{1.224284in}}%
\pgfpathlineto{\pgfqpoint{2.812368in}{1.334744in}}%
\pgfpathlineto{\pgfqpoint{2.817760in}{1.271624in}}%
\pgfpathlineto{\pgfqpoint{2.823152in}{1.287404in}}%
\pgfpathlineto{\pgfqpoint{2.828544in}{1.192724in}}%
\pgfpathlineto{\pgfqpoint{2.833936in}{1.287404in}}%
\pgfpathlineto{\pgfqpoint{2.839328in}{1.445204in}}%
\pgfpathlineto{\pgfqpoint{2.844720in}{1.224284in}}%
\pgfpathlineto{\pgfqpoint{2.850112in}{1.224284in}}%
\pgfpathlineto{\pgfqpoint{2.855504in}{0.971804in}}%
\pgfpathlineto{\pgfqpoint{2.860896in}{1.334744in}}%
\pgfpathlineto{\pgfqpoint{2.866288in}{1.303184in}}%
\pgfpathlineto{\pgfqpoint{2.871680in}{1.240064in}}%
\pgfpathlineto{\pgfqpoint{2.877072in}{1.161164in}}%
\pgfpathlineto{\pgfqpoint{2.882464in}{1.350524in}}%
\pgfpathlineto{\pgfqpoint{2.887856in}{1.286418in}}%
\pgfpathlineto{\pgfqpoint{2.893248in}{1.302198in}}%
\pgfpathlineto{\pgfqpoint{2.898639in}{1.175958in}}%
\pgfpathlineto{\pgfqpoint{2.904031in}{0.986598in}}%
\pgfpathlineto{\pgfqpoint{2.909423in}{1.207518in}}%
\pgfpathlineto{\pgfqpoint{2.914815in}{1.317978in}}%
\pgfpathlineto{\pgfqpoint{2.920207in}{1.317978in}}%
\pgfpathlineto{\pgfqpoint{2.930991in}{1.002378in}}%
\pgfpathlineto{\pgfqpoint{2.936383in}{1.396878in}}%
\pgfpathlineto{\pgfqpoint{2.941775in}{1.223298in}}%
\pgfpathlineto{\pgfqpoint{2.947167in}{1.239078in}}%
\pgfpathlineto{\pgfqpoint{2.952559in}{1.207518in}}%
\pgfpathlineto{\pgfqpoint{2.957951in}{1.239078in}}%
\pgfpathlineto{\pgfqpoint{2.963343in}{1.365318in}}%
\pgfpathlineto{\pgfqpoint{2.968735in}{1.081278in}}%
\pgfpathlineto{\pgfqpoint{2.974127in}{1.223298in}}%
\pgfpathlineto{\pgfqpoint{2.979519in}{1.175958in}}%
\pgfpathlineto{\pgfqpoint{2.984911in}{1.191738in}}%
\pgfpathlineto{\pgfqpoint{2.990303in}{1.223298in}}%
\pgfpathlineto{\pgfqpoint{2.995695in}{1.097058in}}%
\pgfpathlineto{\pgfqpoint{3.001087in}{1.270638in}}%
\pgfpathlineto{\pgfqpoint{3.006479in}{1.239078in}}%
\pgfpathlineto{\pgfqpoint{3.011871in}{1.286418in}}%
\pgfpathlineto{\pgfqpoint{3.017263in}{1.270638in}}%
\pgfpathlineto{\pgfqpoint{3.028047in}{1.302198in}}%
\pgfpathlineto{\pgfqpoint{3.033439in}{1.286418in}}%
\pgfpathlineto{\pgfqpoint{3.044223in}{1.286418in}}%
\pgfpathlineto{\pgfqpoint{3.049615in}{1.396878in}}%
\pgfpathlineto{\pgfqpoint{3.055007in}{1.270638in}}%
\pgfpathlineto{\pgfqpoint{3.060399in}{1.112838in}}%
\pgfpathlineto{\pgfqpoint{3.065791in}{1.239078in}}%
\pgfpathlineto{\pgfqpoint{3.071183in}{1.317978in}}%
\pgfpathlineto{\pgfqpoint{3.076575in}{1.364332in}}%
\pgfpathlineto{\pgfqpoint{3.081967in}{1.143412in}}%
\pgfpathlineto{\pgfqpoint{3.087359in}{1.206532in}}%
\pgfpathlineto{\pgfqpoint{3.092751in}{1.348552in}}%
\pgfpathlineto{\pgfqpoint{3.098143in}{1.332772in}}%
\pgfpathlineto{\pgfqpoint{3.103535in}{1.222312in}}%
\pgfpathlineto{\pgfqpoint{3.108926in}{1.364332in}}%
\pgfpathlineto{\pgfqpoint{3.114318in}{1.253872in}}%
\pgfpathlineto{\pgfqpoint{3.119710in}{1.048731in}}%
\pgfpathlineto{\pgfqpoint{3.125102in}{1.285432in}}%
\pgfpathlineto{\pgfqpoint{3.130494in}{1.159192in}}%
\pgfpathlineto{\pgfqpoint{3.135886in}{1.301212in}}%
\pgfpathlineto{\pgfqpoint{3.141278in}{1.190752in}}%
\pgfpathlineto{\pgfqpoint{3.146670in}{1.395892in}}%
\pgfpathlineto{\pgfqpoint{3.152062in}{1.174972in}}%
\pgfpathlineto{\pgfqpoint{3.157454in}{1.222312in}}%
\pgfpathlineto{\pgfqpoint{3.162846in}{1.190752in}}%
\pgfpathlineto{\pgfqpoint{3.168238in}{1.364332in}}%
\pgfpathlineto{\pgfqpoint{3.173630in}{1.348552in}}%
\pgfpathlineto{\pgfqpoint{3.179022in}{1.269652in}}%
\pgfpathlineto{\pgfqpoint{3.184414in}{1.253872in}}%
\pgfpathlineto{\pgfqpoint{3.189806in}{1.348552in}}%
\pgfpathlineto{\pgfqpoint{3.195198in}{1.206532in}}%
\pgfpathlineto{\pgfqpoint{3.200590in}{1.364332in}}%
\pgfpathlineto{\pgfqpoint{3.205982in}{1.301212in}}%
\pgfpathlineto{\pgfqpoint{3.211374in}{1.380112in}}%
\pgfpathlineto{\pgfqpoint{3.216766in}{1.427452in}}%
\pgfpathlineto{\pgfqpoint{3.222158in}{1.395892in}}%
\pgfpathlineto{\pgfqpoint{3.227550in}{1.348552in}}%
\pgfpathlineto{\pgfqpoint{3.232942in}{1.285432in}}%
\pgfpathlineto{\pgfqpoint{3.238334in}{1.411672in}}%
\pgfpathlineto{\pgfqpoint{3.243726in}{1.269652in}}%
\pgfpathlineto{\pgfqpoint{3.249118in}{1.364332in}}%
\pgfpathlineto{\pgfqpoint{3.254510in}{1.285432in}}%
\pgfpathlineto{\pgfqpoint{3.259902in}{1.316992in}}%
\pgfpathlineto{\pgfqpoint{3.265294in}{1.363345in}}%
\pgfpathlineto{\pgfqpoint{3.270686in}{1.300225in}}%
\pgfpathlineto{\pgfqpoint{3.281470in}{1.095085in}}%
\pgfpathlineto{\pgfqpoint{3.286862in}{1.284445in}}%
\pgfpathlineto{\pgfqpoint{3.292254in}{1.252885in}}%
\pgfpathlineto{\pgfqpoint{3.297646in}{1.347565in}}%
\pgfpathlineto{\pgfqpoint{3.308430in}{1.142425in}}%
\pgfpathlineto{\pgfqpoint{3.313822in}{1.394906in}}%
\pgfpathlineto{\pgfqpoint{3.319213in}{1.284445in}}%
\pgfpathlineto{\pgfqpoint{3.324605in}{1.473806in}}%
\pgfpathlineto{\pgfqpoint{3.329997in}{1.505366in}}%
\pgfpathlineto{\pgfqpoint{3.335389in}{1.221325in}}%
\pgfpathlineto{\pgfqpoint{3.346173in}{1.442246in}}%
\pgfpathlineto{\pgfqpoint{3.351565in}{1.347565in}}%
\pgfpathlineto{\pgfqpoint{3.356957in}{1.489586in}}%
\pgfpathlineto{\pgfqpoint{3.362349in}{1.221325in}}%
\pgfpathlineto{\pgfqpoint{3.367741in}{1.221325in}}%
\pgfpathlineto{\pgfqpoint{3.373133in}{1.268665in}}%
\pgfpathlineto{\pgfqpoint{3.378525in}{1.410686in}}%
\pgfpathlineto{\pgfqpoint{3.383917in}{1.237105in}}%
\pgfpathlineto{\pgfqpoint{3.389309in}{1.347565in}}%
\pgfpathlineto{\pgfqpoint{3.394701in}{1.347565in}}%
\pgfpathlineto{\pgfqpoint{3.400093in}{1.473806in}}%
\pgfpathlineto{\pgfqpoint{3.405485in}{1.363345in}}%
\pgfpathlineto{\pgfqpoint{3.410877in}{1.379126in}}%
\pgfpathlineto{\pgfqpoint{3.416269in}{1.284445in}}%
\pgfpathlineto{\pgfqpoint{3.421661in}{1.268665in}}%
\pgfpathlineto{\pgfqpoint{3.427053in}{1.489586in}}%
\pgfpathlineto{\pgfqpoint{3.432445in}{1.426466in}}%
\pgfpathlineto{\pgfqpoint{3.437837in}{1.379126in}}%
\pgfpathlineto{\pgfqpoint{3.443229in}{1.442246in}}%
\pgfpathlineto{\pgfqpoint{3.448621in}{1.347565in}}%
\pgfpathlineto{\pgfqpoint{3.454013in}{1.299239in}}%
\pgfpathlineto{\pgfqpoint{3.459405in}{1.188779in}}%
\pgfpathlineto{\pgfqpoint{3.464797in}{1.315019in}}%
\pgfpathlineto{\pgfqpoint{3.470189in}{1.472819in}}%
\pgfpathlineto{\pgfqpoint{3.475581in}{1.220339in}}%
\pgfpathlineto{\pgfqpoint{3.480973in}{1.362359in}}%
\pgfpathlineto{\pgfqpoint{3.486365in}{1.425479in}}%
\pgfpathlineto{\pgfqpoint{3.491757in}{1.362359in}}%
\pgfpathlineto{\pgfqpoint{3.497149in}{1.188779in}}%
\pgfpathlineto{\pgfqpoint{3.507933in}{1.283459in}}%
\pgfpathlineto{\pgfqpoint{3.513325in}{1.299239in}}%
\pgfpathlineto{\pgfqpoint{3.518717in}{1.378139in}}%
\pgfpathlineto{\pgfqpoint{3.524109in}{1.551719in}}%
\pgfpathlineto{\pgfqpoint{3.529500in}{1.441259in}}%
\pgfpathlineto{\pgfqpoint{3.534892in}{1.283459in}}%
\pgfpathlineto{\pgfqpoint{3.540284in}{1.315019in}}%
\pgfpathlineto{\pgfqpoint{3.545676in}{1.236119in}}%
\pgfpathlineto{\pgfqpoint{3.551068in}{1.393919in}}%
\pgfpathlineto{\pgfqpoint{3.556460in}{1.283459in}}%
\pgfpathlineto{\pgfqpoint{3.561852in}{1.315019in}}%
\pgfpathlineto{\pgfqpoint{3.567244in}{1.315019in}}%
\pgfpathlineto{\pgfqpoint{3.572636in}{1.472819in}}%
\pgfpathlineto{\pgfqpoint{3.578028in}{1.299239in}}%
\pgfpathlineto{\pgfqpoint{3.583420in}{1.457039in}}%
\pgfpathlineto{\pgfqpoint{3.588812in}{1.425479in}}%
\pgfpathlineto{\pgfqpoint{3.594204in}{1.378139in}}%
\pgfpathlineto{\pgfqpoint{3.599596in}{1.409699in}}%
\pgfpathlineto{\pgfqpoint{3.604988in}{1.330799in}}%
\pgfpathlineto{\pgfqpoint{3.610380in}{1.409699in}}%
\pgfpathlineto{\pgfqpoint{3.615772in}{1.236119in}}%
\pgfpathlineto{\pgfqpoint{3.621164in}{1.346579in}}%
\pgfpathlineto{\pgfqpoint{3.626556in}{1.172999in}}%
\pgfpathlineto{\pgfqpoint{3.631948in}{1.520159in}}%
\pgfpathlineto{\pgfqpoint{3.637340in}{1.346579in}}%
\pgfpathlineto{\pgfqpoint{3.642732in}{1.361373in}}%
\pgfpathlineto{\pgfqpoint{3.648124in}{1.361373in}}%
\pgfpathlineto{\pgfqpoint{3.653516in}{1.519173in}}%
\pgfpathlineto{\pgfqpoint{3.658908in}{1.424493in}}%
\pgfpathlineto{\pgfqpoint{3.664300in}{1.377153in}}%
\pgfpathlineto{\pgfqpoint{3.669692in}{1.408713in}}%
\pgfpathlineto{\pgfqpoint{3.675084in}{1.503393in}}%
\pgfpathlineto{\pgfqpoint{3.680476in}{1.345593in}}%
\pgfpathlineto{\pgfqpoint{3.685868in}{1.329813in}}%
\pgfpathlineto{\pgfqpoint{3.696652in}{1.093113in}}%
\pgfpathlineto{\pgfqpoint{3.702044in}{1.424493in}}%
\pgfpathlineto{\pgfqpoint{3.707436in}{1.250913in}}%
\pgfpathlineto{\pgfqpoint{3.712828in}{1.424493in}}%
\pgfpathlineto{\pgfqpoint{3.718220in}{1.377153in}}%
\pgfpathlineto{\pgfqpoint{3.723612in}{1.392933in}}%
\pgfpathlineto{\pgfqpoint{3.729004in}{1.361373in}}%
\pgfpathlineto{\pgfqpoint{3.734396in}{1.392933in}}%
\pgfpathlineto{\pgfqpoint{3.739787in}{1.203573in}}%
\pgfpathlineto{\pgfqpoint{3.745179in}{1.235133in}}%
\pgfpathlineto{\pgfqpoint{3.750571in}{1.440273in}}%
\pgfpathlineto{\pgfqpoint{3.755963in}{1.361373in}}%
\pgfpathlineto{\pgfqpoint{3.761355in}{1.314033in}}%
\pgfpathlineto{\pgfqpoint{3.766747in}{1.282473in}}%
\pgfpathlineto{\pgfqpoint{3.772139in}{1.487613in}}%
\pgfpathlineto{\pgfqpoint{3.777531in}{1.377153in}}%
\pgfpathlineto{\pgfqpoint{3.782923in}{1.408713in}}%
\pgfpathlineto{\pgfqpoint{3.788315in}{1.408713in}}%
\pgfpathlineto{\pgfqpoint{3.793707in}{1.329813in}}%
\pgfpathlineto{\pgfqpoint{3.799099in}{1.298253in}}%
\pgfpathlineto{\pgfqpoint{3.804491in}{1.424493in}}%
\pgfpathlineto{\pgfqpoint{3.809883in}{1.329813in}}%
\pgfpathlineto{\pgfqpoint{3.815275in}{1.440273in}}%
\pgfpathlineto{\pgfqpoint{3.820667in}{1.377153in}}%
\pgfpathlineto{\pgfqpoint{3.826059in}{1.440273in}}%
\pgfpathlineto{\pgfqpoint{3.831451in}{1.313047in}}%
\pgfpathlineto{\pgfqpoint{3.836843in}{1.486627in}}%
\pgfpathlineto{\pgfqpoint{3.842235in}{1.439287in}}%
\pgfpathlineto{\pgfqpoint{3.847627in}{1.518187in}}%
\pgfpathlineto{\pgfqpoint{3.853019in}{1.423507in}}%
\pgfpathlineto{\pgfqpoint{3.858411in}{1.439287in}}%
\pgfpathlineto{\pgfqpoint{3.863803in}{1.281487in}}%
\pgfpathlineto{\pgfqpoint{3.869195in}{1.281487in}}%
\pgfpathlineto{\pgfqpoint{3.874587in}{1.470847in}}%
\pgfpathlineto{\pgfqpoint{3.879979in}{1.360387in}}%
\pgfpathlineto{\pgfqpoint{3.885371in}{1.344607in}}%
\pgfpathlineto{\pgfqpoint{3.890763in}{1.360387in}}%
\pgfpathlineto{\pgfqpoint{3.896155in}{1.455067in}}%
\pgfpathlineto{\pgfqpoint{3.901547in}{1.486627in}}%
\pgfpathlineto{\pgfqpoint{3.906939in}{1.139467in}}%
\pgfpathlineto{\pgfqpoint{3.912331in}{1.455067in}}%
\pgfpathlineto{\pgfqpoint{3.923115in}{1.360387in}}%
\pgfpathlineto{\pgfqpoint{3.928507in}{1.407727in}}%
\pgfpathlineto{\pgfqpoint{3.933899in}{1.423507in}}%
\pgfpathlineto{\pgfqpoint{3.939291in}{1.423507in}}%
\pgfpathlineto{\pgfqpoint{3.944683in}{1.391947in}}%
\pgfpathlineto{\pgfqpoint{3.950074in}{1.265707in}}%
\pgfpathlineto{\pgfqpoint{3.955466in}{1.533967in}}%
\pgfpathlineto{\pgfqpoint{3.960858in}{1.376167in}}%
\pgfpathlineto{\pgfqpoint{3.966250in}{1.407727in}}%
\pgfpathlineto{\pgfqpoint{3.971642in}{1.407727in}}%
\pgfpathlineto{\pgfqpoint{3.977034in}{1.518187in}}%
\pgfpathlineto{\pgfqpoint{3.982426in}{1.376167in}}%
\pgfpathlineto{\pgfqpoint{3.987818in}{1.360387in}}%
\pgfpathlineto{\pgfqpoint{3.993210in}{1.565527in}}%
\pgfpathlineto{\pgfqpoint{3.998602in}{1.533967in}}%
\pgfpathlineto{\pgfqpoint{4.003994in}{1.439287in}}%
\pgfpathlineto{\pgfqpoint{4.009386in}{1.439287in}}%
\pgfpathlineto{\pgfqpoint{4.014778in}{1.391947in}}%
\pgfpathlineto{\pgfqpoint{4.020170in}{1.438301in}}%
\pgfpathlineto{\pgfqpoint{4.025562in}{1.280500in}}%
\pgfpathlineto{\pgfqpoint{4.030954in}{1.359400in}}%
\pgfpathlineto{\pgfqpoint{4.041738in}{1.454081in}}%
\pgfpathlineto{\pgfqpoint{4.047130in}{1.375181in}}%
\pgfpathlineto{\pgfqpoint{4.052522in}{1.454081in}}%
\pgfpathlineto{\pgfqpoint{4.057914in}{1.438301in}}%
\pgfpathlineto{\pgfqpoint{4.063306in}{1.501421in}}%
\pgfpathlineto{\pgfqpoint{4.068698in}{1.390961in}}%
\pgfpathlineto{\pgfqpoint{4.074090in}{1.438301in}}%
\pgfpathlineto{\pgfqpoint{4.079482in}{1.296280in}}%
\pgfpathlineto{\pgfqpoint{4.084874in}{1.359400in}}%
\pgfpathlineto{\pgfqpoint{4.090266in}{1.485641in}}%
\pgfpathlineto{\pgfqpoint{4.095658in}{1.343620in}}%
\pgfpathlineto{\pgfqpoint{4.101050in}{1.233160in}}%
\pgfpathlineto{\pgfqpoint{4.106442in}{1.296280in}}%
\pgfpathlineto{\pgfqpoint{4.111834in}{1.501421in}}%
\pgfpathlineto{\pgfqpoint{4.117226in}{1.501421in}}%
\pgfpathlineto{\pgfqpoint{4.122618in}{1.390961in}}%
\pgfpathlineto{\pgfqpoint{4.133402in}{1.233160in}}%
\pgfpathlineto{\pgfqpoint{4.138794in}{1.485641in}}%
\pgfpathlineto{\pgfqpoint{4.149578in}{1.327840in}}%
\pgfpathlineto{\pgfqpoint{4.154970in}{1.454081in}}%
\pgfpathlineto{\pgfqpoint{4.160361in}{1.454081in}}%
\pgfpathlineto{\pgfqpoint{4.165753in}{1.375181in}}%
\pgfpathlineto{\pgfqpoint{4.171145in}{1.406741in}}%
\pgfpathlineto{\pgfqpoint{4.176537in}{1.454081in}}%
\pgfpathlineto{\pgfqpoint{4.181929in}{1.422521in}}%
\pgfpathlineto{\pgfqpoint{4.187321in}{1.343620in}}%
\pgfpathlineto{\pgfqpoint{4.192713in}{1.327840in}}%
\pgfpathlineto{\pgfqpoint{4.198105in}{1.422521in}}%
\pgfpathlineto{\pgfqpoint{4.203497in}{1.406741in}}%
\pgfpathlineto{\pgfqpoint{4.208889in}{1.358414in}}%
\pgfpathlineto{\pgfqpoint{4.214281in}{1.279514in}}%
\pgfpathlineto{\pgfqpoint{4.219673in}{1.500434in}}%
\pgfpathlineto{\pgfqpoint{4.225065in}{1.358414in}}%
\pgfpathlineto{\pgfqpoint{4.230457in}{1.389974in}}%
\pgfpathlineto{\pgfqpoint{4.235849in}{1.311074in}}%
\pgfpathlineto{\pgfqpoint{4.241241in}{1.405754in}}%
\pgfpathlineto{\pgfqpoint{4.246633in}{1.263734in}}%
\pgfpathlineto{\pgfqpoint{4.252025in}{1.326854in}}%
\pgfpathlineto{\pgfqpoint{4.257417in}{1.468874in}}%
\pgfpathlineto{\pgfqpoint{4.262809in}{1.295294in}}%
\pgfpathlineto{\pgfqpoint{4.268201in}{1.516214in}}%
\pgfpathlineto{\pgfqpoint{4.273593in}{1.358414in}}%
\pgfpathlineto{\pgfqpoint{4.278985in}{1.389974in}}%
\pgfpathlineto{\pgfqpoint{4.284377in}{1.216394in}}%
\pgfpathlineto{\pgfqpoint{4.289769in}{1.263734in}}%
\pgfpathlineto{\pgfqpoint{4.295161in}{1.326854in}}%
\pgfpathlineto{\pgfqpoint{4.300553in}{1.358414in}}%
\pgfpathlineto{\pgfqpoint{4.305945in}{1.453094in}}%
\pgfpathlineto{\pgfqpoint{4.311337in}{1.437314in}}%
\pgfpathlineto{\pgfqpoint{4.316729in}{1.374194in}}%
\pgfpathlineto{\pgfqpoint{4.322121in}{1.389974in}}%
\pgfpathlineto{\pgfqpoint{4.327513in}{1.389974in}}%
\pgfpathlineto{\pgfqpoint{4.332905in}{1.216394in}}%
\pgfpathlineto{\pgfqpoint{4.338297in}{1.247954in}}%
\pgfpathlineto{\pgfqpoint{4.343689in}{1.453094in}}%
\pgfpathlineto{\pgfqpoint{4.349081in}{1.279514in}}%
\pgfpathlineto{\pgfqpoint{4.354473in}{1.263734in}}%
\pgfpathlineto{\pgfqpoint{4.359865in}{1.468874in}}%
\pgfpathlineto{\pgfqpoint{4.365257in}{1.405754in}}%
\pgfpathlineto{\pgfqpoint{4.370648in}{1.358414in}}%
\pgfpathlineto{\pgfqpoint{4.376040in}{1.374194in}}%
\pgfpathlineto{\pgfqpoint{4.381432in}{1.358414in}}%
\pgfpathlineto{\pgfqpoint{4.386824in}{1.453094in}}%
\pgfpathlineto{\pgfqpoint{4.392216in}{1.311074in}}%
\pgfpathlineto{\pgfqpoint{4.397608in}{1.325868in}}%
\pgfpathlineto{\pgfqpoint{4.403000in}{1.310088in}}%
\pgfpathlineto{\pgfqpoint{4.408392in}{1.531008in}}%
\pgfpathlineto{\pgfqpoint{4.413784in}{1.436328in}}%
\pgfpathlineto{\pgfqpoint{4.419176in}{1.168068in}}%
\pgfpathlineto{\pgfqpoint{4.424568in}{1.310088in}}%
\pgfpathlineto{\pgfqpoint{4.429960in}{1.310088in}}%
\pgfpathlineto{\pgfqpoint{4.435352in}{1.262748in}}%
\pgfpathlineto{\pgfqpoint{4.440744in}{1.341648in}}%
\pgfpathlineto{\pgfqpoint{4.446136in}{1.325868in}}%
\pgfpathlineto{\pgfqpoint{4.451528in}{1.357428in}}%
\pgfpathlineto{\pgfqpoint{4.456920in}{1.373208in}}%
\pgfpathlineto{\pgfqpoint{4.462312in}{1.436328in}}%
\pgfpathlineto{\pgfqpoint{4.467704in}{1.388988in}}%
\pgfpathlineto{\pgfqpoint{4.473096in}{1.373208in}}%
\pgfpathlineto{\pgfqpoint{4.478488in}{1.341648in}}%
\pgfpathlineto{\pgfqpoint{4.489272in}{1.310088in}}%
\pgfpathlineto{\pgfqpoint{4.494664in}{1.246968in}}%
\pgfpathlineto{\pgfqpoint{4.500056in}{1.357428in}}%
\pgfpathlineto{\pgfqpoint{4.505448in}{1.436328in}}%
\pgfpathlineto{\pgfqpoint{4.510840in}{1.357428in}}%
\pgfpathlineto{\pgfqpoint{4.516232in}{1.357428in}}%
\pgfpathlineto{\pgfqpoint{4.521624in}{1.373208in}}%
\pgfpathlineto{\pgfqpoint{4.527016in}{1.199628in}}%
\pgfpathlineto{\pgfqpoint{4.532408in}{1.215408in}}%
\pgfpathlineto{\pgfqpoint{4.537800in}{1.467888in}}%
\pgfpathlineto{\pgfqpoint{4.543192in}{1.420548in}}%
\pgfpathlineto{\pgfqpoint{4.548584in}{1.420548in}}%
\pgfpathlineto{\pgfqpoint{4.553976in}{1.373208in}}%
\pgfpathlineto{\pgfqpoint{4.559368in}{1.404768in}}%
\pgfpathlineto{\pgfqpoint{4.564760in}{1.467888in}}%
\pgfpathlineto{\pgfqpoint{4.570152in}{1.357428in}}%
\pgfpathlineto{\pgfqpoint{4.575544in}{1.420548in}}%
\pgfpathlineto{\pgfqpoint{4.580935in}{1.215408in}}%
\pgfpathlineto{\pgfqpoint{4.586327in}{1.245982in}}%
\pgfpathlineto{\pgfqpoint{4.591719in}{1.309102in}}%
\pgfpathlineto{\pgfqpoint{4.597111in}{1.419562in}}%
\pgfpathlineto{\pgfqpoint{4.602503in}{1.340662in}}%
\pgfpathlineto{\pgfqpoint{4.607895in}{1.451122in}}%
\pgfpathlineto{\pgfqpoint{4.613287in}{1.245982in}}%
\pgfpathlineto{\pgfqpoint{4.618679in}{1.324882in}}%
\pgfpathlineto{\pgfqpoint{4.624071in}{0.993501in}}%
\pgfpathlineto{\pgfqpoint{4.629463in}{1.261762in}}%
\pgfpathlineto{\pgfqpoint{4.634855in}{1.151302in}}%
\pgfpathlineto{\pgfqpoint{4.640247in}{1.182862in}}%
\pgfpathlineto{\pgfqpoint{4.645639in}{1.245982in}}%
\pgfpathlineto{\pgfqpoint{4.651031in}{1.403782in}}%
\pgfpathlineto{\pgfqpoint{4.656423in}{1.419562in}}%
\pgfpathlineto{\pgfqpoint{4.661815in}{1.309102in}}%
\pgfpathlineto{\pgfqpoint{4.667207in}{1.324882in}}%
\pgfpathlineto{\pgfqpoint{4.672599in}{1.388002in}}%
\pgfpathlineto{\pgfqpoint{4.677991in}{1.324882in}}%
\pgfpathlineto{\pgfqpoint{4.683383in}{1.182862in}}%
\pgfpathlineto{\pgfqpoint{4.688775in}{1.435342in}}%
\pgfpathlineto{\pgfqpoint{4.694167in}{1.388002in}}%
\pgfpathlineto{\pgfqpoint{4.699559in}{1.466902in}}%
\pgfpathlineto{\pgfqpoint{4.704951in}{1.340662in}}%
\pgfpathlineto{\pgfqpoint{4.715735in}{1.135522in}}%
\pgfpathlineto{\pgfqpoint{4.721127in}{1.356442in}}%
\pgfpathlineto{\pgfqpoint{4.726519in}{1.167082in}}%
\pgfpathlineto{\pgfqpoint{4.737303in}{1.466902in}}%
\pgfpathlineto{\pgfqpoint{4.742695in}{1.356442in}}%
\pgfpathlineto{\pgfqpoint{4.748087in}{1.356442in}}%
\pgfpathlineto{\pgfqpoint{4.753479in}{1.309102in}}%
\pgfpathlineto{\pgfqpoint{4.758871in}{1.293322in}}%
\pgfpathlineto{\pgfqpoint{4.764263in}{1.167082in}}%
\pgfpathlineto{\pgfqpoint{4.769655in}{1.103962in}}%
\pgfpathlineto{\pgfqpoint{4.775047in}{1.339675in}}%
\pgfpathlineto{\pgfqpoint{4.780439in}{1.371236in}}%
\pgfpathlineto{\pgfqpoint{4.785831in}{1.181875in}}%
\pgfpathlineto{\pgfqpoint{4.791222in}{1.339675in}}%
\pgfpathlineto{\pgfqpoint{4.796614in}{1.402796in}}%
\pgfpathlineto{\pgfqpoint{4.802006in}{1.402796in}}%
\pgfpathlineto{\pgfqpoint{4.807398in}{1.387016in}}%
\pgfpathlineto{\pgfqpoint{4.812790in}{1.166095in}}%
\pgfpathlineto{\pgfqpoint{4.818182in}{1.323895in}}%
\pgfpathlineto{\pgfqpoint{4.823574in}{1.339675in}}%
\pgfpathlineto{\pgfqpoint{4.828966in}{1.102975in}}%
\pgfpathlineto{\pgfqpoint{4.839750in}{1.465916in}}%
\pgfpathlineto{\pgfqpoint{4.855926in}{1.134535in}}%
\pgfpathlineto{\pgfqpoint{4.861318in}{1.418576in}}%
\pgfpathlineto{\pgfqpoint{4.866710in}{1.387016in}}%
\pgfpathlineto{\pgfqpoint{4.872102in}{1.276555in}}%
\pgfpathlineto{\pgfqpoint{4.877494in}{1.387016in}}%
\pgfpathlineto{\pgfqpoint{4.882886in}{1.292335in}}%
\pgfpathlineto{\pgfqpoint{4.893670in}{1.355455in}}%
\pgfpathlineto{\pgfqpoint{4.899062in}{1.260775in}}%
\pgfpathlineto{\pgfqpoint{4.904454in}{1.276555in}}%
\pgfpathlineto{\pgfqpoint{4.909846in}{1.355455in}}%
\pgfpathlineto{\pgfqpoint{4.915238in}{1.387016in}}%
\pgfpathlineto{\pgfqpoint{4.920630in}{1.166095in}}%
\pgfpathlineto{\pgfqpoint{4.926022in}{1.450136in}}%
\pgfpathlineto{\pgfqpoint{4.931414in}{1.276555in}}%
\pgfpathlineto{\pgfqpoint{4.936806in}{1.229215in}}%
\pgfpathlineto{\pgfqpoint{4.942198in}{1.387016in}}%
\pgfpathlineto{\pgfqpoint{4.947590in}{1.260775in}}%
\pgfpathlineto{\pgfqpoint{4.952982in}{1.308115in}}%
\pgfpathlineto{\pgfqpoint{4.958374in}{1.276555in}}%
\pgfpathlineto{\pgfqpoint{4.963766in}{1.180889in}}%
\pgfpathlineto{\pgfqpoint{4.974550in}{1.307129in}}%
\pgfpathlineto{\pgfqpoint{4.979942in}{1.275569in}}%
\pgfpathlineto{\pgfqpoint{4.985334in}{1.338689in}}%
\pgfpathlineto{\pgfqpoint{4.990726in}{1.370249in}}%
\pgfpathlineto{\pgfqpoint{4.996117in}{1.338689in}}%
\pgfpathlineto{\pgfqpoint{5.001509in}{1.338689in}}%
\pgfpathlineto{\pgfqpoint{5.006901in}{1.228229in}}%
\pgfpathlineto{\pgfqpoint{5.012293in}{1.354469in}}%
\pgfpathlineto{\pgfqpoint{5.017685in}{1.259789in}}%
\pgfpathlineto{\pgfqpoint{5.023077in}{1.370249in}}%
\pgfpathlineto{\pgfqpoint{5.028469in}{1.322909in}}%
\pgfpathlineto{\pgfqpoint{5.033861in}{1.433369in}}%
\pgfpathlineto{\pgfqpoint{5.039253in}{1.401809in}}%
\pgfpathlineto{\pgfqpoint{5.044645in}{1.338689in}}%
\pgfpathlineto{\pgfqpoint{5.050037in}{1.259789in}}%
\pgfpathlineto{\pgfqpoint{5.055429in}{1.244009in}}%
\pgfpathlineto{\pgfqpoint{5.060821in}{1.291349in}}%
\pgfpathlineto{\pgfqpoint{5.066213in}{1.496489in}}%
\pgfpathlineto{\pgfqpoint{5.071605in}{1.307129in}}%
\pgfpathlineto{\pgfqpoint{5.076997in}{1.417589in}}%
\pgfpathlineto{\pgfqpoint{5.082389in}{1.417589in}}%
\pgfpathlineto{\pgfqpoint{5.087781in}{1.464929in}}%
\pgfpathlineto{\pgfqpoint{5.093173in}{0.975749in}}%
\pgfpathlineto{\pgfqpoint{5.098565in}{1.259789in}}%
\pgfpathlineto{\pgfqpoint{5.103957in}{1.275569in}}%
\pgfpathlineto{\pgfqpoint{5.109349in}{1.196669in}}%
\pgfpathlineto{\pgfqpoint{5.114741in}{1.244009in}}%
\pgfpathlineto{\pgfqpoint{5.120133in}{1.354469in}}%
\pgfpathlineto{\pgfqpoint{5.125525in}{1.196669in}}%
\pgfpathlineto{\pgfqpoint{5.130917in}{1.259789in}}%
\pgfpathlineto{\pgfqpoint{5.136309in}{1.196669in}}%
\pgfpathlineto{\pgfqpoint{5.141701in}{1.275569in}}%
\pgfpathlineto{\pgfqpoint{5.147093in}{1.386029in}}%
\pgfpathlineto{\pgfqpoint{5.157877in}{1.037883in}}%
\pgfpathlineto{\pgfqpoint{5.163269in}{1.337703in}}%
\pgfpathlineto{\pgfqpoint{5.168661in}{1.432383in}}%
\pgfpathlineto{\pgfqpoint{5.174053in}{1.385043in}}%
\pgfpathlineto{\pgfqpoint{5.179445in}{1.195683in}}%
\pgfpathlineto{\pgfqpoint{5.184837in}{1.400823in}}%
\pgfpathlineto{\pgfqpoint{5.190229in}{1.227243in}}%
\pgfpathlineto{\pgfqpoint{5.195621in}{1.337703in}}%
\pgfpathlineto{\pgfqpoint{5.201013in}{1.211463in}}%
\pgfpathlineto{\pgfqpoint{5.206404in}{1.306143in}}%
\pgfpathlineto{\pgfqpoint{5.211796in}{1.101003in}}%
\pgfpathlineto{\pgfqpoint{5.217188in}{1.258803in}}%
\pgfpathlineto{\pgfqpoint{5.222580in}{1.353483in}}%
\pgfpathlineto{\pgfqpoint{5.227972in}{1.211463in}}%
\pgfpathlineto{\pgfqpoint{5.238756in}{1.274583in}}%
\pgfpathlineto{\pgfqpoint{5.244148in}{1.227243in}}%
\pgfpathlineto{\pgfqpoint{5.249540in}{1.053663in}}%
\pgfpathlineto{\pgfqpoint{5.254932in}{1.337703in}}%
\pgfpathlineto{\pgfqpoint{5.260324in}{1.195683in}}%
\pgfpathlineto{\pgfqpoint{5.265716in}{1.353483in}}%
\pgfpathlineto{\pgfqpoint{5.271108in}{1.116783in}}%
\pgfpathlineto{\pgfqpoint{5.276500in}{1.132563in}}%
\pgfpathlineto{\pgfqpoint{5.281892in}{1.274583in}}%
\pgfpathlineto{\pgfqpoint{5.287284in}{1.321923in}}%
\pgfpathlineto{\pgfqpoint{5.292676in}{1.400823in}}%
\pgfpathlineto{\pgfqpoint{5.298068in}{1.385043in}}%
\pgfpathlineto{\pgfqpoint{5.303460in}{1.258803in}}%
\pgfpathlineto{\pgfqpoint{5.308852in}{1.243023in}}%
\pgfpathlineto{\pgfqpoint{5.314244in}{1.211463in}}%
\pgfpathlineto{\pgfqpoint{5.319636in}{1.148343in}}%
\pgfpathlineto{\pgfqpoint{5.325028in}{1.337703in}}%
\pgfpathlineto{\pgfqpoint{5.330420in}{1.337703in}}%
\pgfpathlineto{\pgfqpoint{5.335812in}{1.416603in}}%
\pgfpathlineto{\pgfqpoint{5.341204in}{1.194697in}}%
\pgfpathlineto{\pgfqpoint{5.346596in}{1.226257in}}%
\pgfpathlineto{\pgfqpoint{5.351988in}{1.242037in}}%
\pgfpathlineto{\pgfqpoint{5.357380in}{1.289377in}}%
\pgfpathlineto{\pgfqpoint{5.368164in}{1.226257in}}%
\pgfpathlineto{\pgfqpoint{5.373556in}{1.384057in}}%
\pgfpathlineto{\pgfqpoint{5.378948in}{1.289377in}}%
\pgfpathlineto{\pgfqpoint{5.384340in}{1.368277in}}%
\pgfpathlineto{\pgfqpoint{5.389732in}{1.242037in}}%
\pgfpathlineto{\pgfqpoint{5.395124in}{1.273597in}}%
\pgfpathlineto{\pgfqpoint{5.400516in}{1.257817in}}%
\pgfpathlineto{\pgfqpoint{5.405908in}{1.210477in}}%
\pgfpathlineto{\pgfqpoint{5.411300in}{1.178917in}}%
\pgfpathlineto{\pgfqpoint{5.416691in}{1.178917in}}%
\pgfpathlineto{\pgfqpoint{5.422083in}{1.257817in}}%
\pgfpathlineto{\pgfqpoint{5.427475in}{1.257817in}}%
\pgfpathlineto{\pgfqpoint{5.432867in}{1.163137in}}%
\pgfpathlineto{\pgfqpoint{5.438259in}{1.305157in}}%
\pgfpathlineto{\pgfqpoint{5.449043in}{1.242037in}}%
\pgfpathlineto{\pgfqpoint{5.454435in}{1.226257in}}%
\pgfpathlineto{\pgfqpoint{5.459827in}{1.242037in}}%
\pgfpathlineto{\pgfqpoint{5.465219in}{1.242037in}}%
\pgfpathlineto{\pgfqpoint{5.470611in}{1.257817in}}%
\pgfpathlineto{\pgfqpoint{5.476003in}{1.336717in}}%
\pgfpathlineto{\pgfqpoint{5.481395in}{1.163137in}}%
\pgfpathlineto{\pgfqpoint{5.486787in}{1.289377in}}%
\pgfpathlineto{\pgfqpoint{5.492179in}{1.131577in}}%
\pgfpathlineto{\pgfqpoint{5.497571in}{1.257817in}}%
\pgfpathlineto{\pgfqpoint{5.502963in}{1.289377in}}%
\pgfpathlineto{\pgfqpoint{5.508355in}{1.305157in}}%
\pgfpathlineto{\pgfqpoint{5.513747in}{1.052676in}}%
\pgfpathlineto{\pgfqpoint{5.519139in}{1.305157in}}%
\pgfpathlineto{\pgfqpoint{5.524531in}{1.210477in}}%
\pgfpathlineto{\pgfqpoint{5.529923in}{1.162150in}}%
\pgfpathlineto{\pgfqpoint{5.535315in}{1.335730in}}%
\pgfpathlineto{\pgfqpoint{5.540707in}{1.430411in}}%
\pgfpathlineto{\pgfqpoint{5.546099in}{1.209490in}}%
\pgfpathlineto{\pgfqpoint{5.551491in}{1.146370in}}%
\pgfpathlineto{\pgfqpoint{5.556883in}{1.067470in}}%
\pgfpathlineto{\pgfqpoint{5.562275in}{1.383071in}}%
\pgfpathlineto{\pgfqpoint{5.567667in}{1.225270in}}%
\pgfpathlineto{\pgfqpoint{5.573059in}{1.256830in}}%
\pgfpathlineto{\pgfqpoint{5.578451in}{1.099030in}}%
\pgfpathlineto{\pgfqpoint{5.583843in}{1.335730in}}%
\pgfpathlineto{\pgfqpoint{5.589235in}{1.367291in}}%
\pgfpathlineto{\pgfqpoint{5.594627in}{1.367291in}}%
\pgfpathlineto{\pgfqpoint{5.600019in}{1.162150in}}%
\pgfpathlineto{\pgfqpoint{5.605411in}{1.130590in}}%
\pgfpathlineto{\pgfqpoint{5.610803in}{1.114810in}}%
\pgfpathlineto{\pgfqpoint{5.616195in}{1.004350in}}%
\pgfpathlineto{\pgfqpoint{5.621587in}{1.146370in}}%
\pgfpathlineto{\pgfqpoint{5.626978in}{1.083250in}}%
\pgfpathlineto{\pgfqpoint{5.632370in}{1.367291in}}%
\pgfpathlineto{\pgfqpoint{5.637762in}{1.335730in}}%
\pgfpathlineto{\pgfqpoint{5.643154in}{1.256830in}}%
\pgfpathlineto{\pgfqpoint{5.648546in}{1.256830in}}%
\pgfpathlineto{\pgfqpoint{5.653938in}{1.225270in}}%
\pgfpathlineto{\pgfqpoint{5.659330in}{1.272610in}}%
\pgfpathlineto{\pgfqpoint{5.664722in}{1.083250in}}%
\pgfpathlineto{\pgfqpoint{5.670114in}{1.177930in}}%
\pgfpathlineto{\pgfqpoint{5.675506in}{1.177930in}}%
\pgfpathlineto{\pgfqpoint{5.680898in}{1.351510in}}%
\pgfpathlineto{\pgfqpoint{5.686290in}{1.367291in}}%
\pgfpathlineto{\pgfqpoint{5.697074in}{1.225270in}}%
\pgfpathlineto{\pgfqpoint{5.702466in}{1.177930in}}%
\pgfpathlineto{\pgfqpoint{5.707858in}{1.319950in}}%
\pgfpathlineto{\pgfqpoint{5.713250in}{1.209490in}}%
\pgfpathlineto{\pgfqpoint{5.718642in}{1.240064in}}%
\pgfpathlineto{\pgfqpoint{5.724034in}{1.192724in}}%
\pgfpathlineto{\pgfqpoint{5.729426in}{1.366304in}}%
\pgfpathlineto{\pgfqpoint{5.734818in}{1.334744in}}%
\pgfpathlineto{\pgfqpoint{5.740210in}{1.318964in}}%
\pgfpathlineto{\pgfqpoint{5.745602in}{1.161164in}}%
\pgfpathlineto{\pgfqpoint{5.750994in}{1.334744in}}%
\pgfpathlineto{\pgfqpoint{5.756386in}{1.287404in}}%
\pgfpathlineto{\pgfqpoint{5.761778in}{1.303184in}}%
\pgfpathlineto{\pgfqpoint{5.767170in}{1.208504in}}%
\pgfpathlineto{\pgfqpoint{5.772562in}{1.192724in}}%
\pgfpathlineto{\pgfqpoint{5.783346in}{1.318964in}}%
\pgfpathlineto{\pgfqpoint{5.788738in}{1.287404in}}%
\pgfpathlineto{\pgfqpoint{5.794130in}{1.161164in}}%
\pgfpathlineto{\pgfqpoint{5.799522in}{1.098044in}}%
\pgfpathlineto{\pgfqpoint{5.804914in}{1.050704in}}%
\pgfpathlineto{\pgfqpoint{5.810306in}{1.255844in}}%
\pgfpathlineto{\pgfqpoint{5.815698in}{1.050704in}}%
\pgfpathlineto{\pgfqpoint{5.821090in}{1.271624in}}%
\pgfpathlineto{\pgfqpoint{5.826482in}{1.224284in}}%
\pgfpathlineto{\pgfqpoint{5.831874in}{1.287404in}}%
\pgfpathlineto{\pgfqpoint{5.837265in}{1.287404in}}%
\pgfpathlineto{\pgfqpoint{5.842657in}{1.176944in}}%
\pgfpathlineto{\pgfqpoint{5.848049in}{1.318964in}}%
\pgfpathlineto{\pgfqpoint{5.853441in}{1.318964in}}%
\pgfpathlineto{\pgfqpoint{5.858833in}{1.350524in}}%
\pgfpathlineto{\pgfqpoint{5.864225in}{0.987584in}}%
\pgfpathlineto{\pgfqpoint{5.869617in}{1.492544in}}%
\pgfpathlineto{\pgfqpoint{5.875009in}{1.287404in}}%
\pgfpathlineto{\pgfqpoint{5.880401in}{1.208504in}}%
\pgfpathlineto{\pgfqpoint{5.885793in}{1.082264in}}%
\pgfpathlineto{\pgfqpoint{5.891185in}{1.240064in}}%
\pgfpathlineto{\pgfqpoint{5.896577in}{1.255844in}}%
\pgfpathlineto{\pgfqpoint{5.901969in}{1.240064in}}%
\pgfpathlineto{\pgfqpoint{5.907361in}{1.333758in}}%
\pgfpathlineto{\pgfqpoint{5.912753in}{1.254858in}}%
\pgfpathlineto{\pgfqpoint{5.918145in}{1.286418in}}%
\pgfpathlineto{\pgfqpoint{5.923537in}{1.302198in}}%
\pgfpathlineto{\pgfqpoint{5.928929in}{1.254858in}}%
\pgfpathlineto{\pgfqpoint{5.934321in}{1.112838in}}%
\pgfpathlineto{\pgfqpoint{5.939713in}{1.286418in}}%
\pgfpathlineto{\pgfqpoint{5.945105in}{1.270638in}}%
\pgfpathlineto{\pgfqpoint{5.950497in}{1.223298in}}%
\pgfpathlineto{\pgfqpoint{5.955889in}{1.207518in}}%
\pgfpathlineto{\pgfqpoint{5.961281in}{1.317978in}}%
\pgfpathlineto{\pgfqpoint{5.966673in}{1.475778in}}%
\pgfpathlineto{\pgfqpoint{5.972065in}{1.396878in}}%
\pgfpathlineto{\pgfqpoint{5.977457in}{1.428438in}}%
\pgfpathlineto{\pgfqpoint{5.982849in}{1.302198in}}%
\pgfpathlineto{\pgfqpoint{5.988241in}{1.223298in}}%
\pgfpathlineto{\pgfqpoint{5.993633in}{1.333758in}}%
\pgfpathlineto{\pgfqpoint{5.999025in}{1.412658in}}%
\pgfpathlineto{\pgfqpoint{6.004417in}{1.317978in}}%
\pgfpathlineto{\pgfqpoint{6.009809in}{1.333758in}}%
\pgfpathlineto{\pgfqpoint{6.015201in}{1.270638in}}%
\pgfpathlineto{\pgfqpoint{6.020593in}{1.365318in}}%
\pgfpathlineto{\pgfqpoint{6.025985in}{1.491558in}}%
\pgfpathlineto{\pgfqpoint{6.031377in}{1.459998in}}%
\pgfpathlineto{\pgfqpoint{6.036769in}{1.365318in}}%
\pgfpathlineto{\pgfqpoint{6.047552in}{1.270638in}}%
\pgfpathlineto{\pgfqpoint{6.052944in}{1.175958in}}%
\pgfpathlineto{\pgfqpoint{6.058336in}{1.286418in}}%
\pgfpathlineto{\pgfqpoint{6.063728in}{1.175958in}}%
\pgfpathlineto{\pgfqpoint{6.074512in}{1.207518in}}%
\pgfpathlineto{\pgfqpoint{6.079904in}{1.112838in}}%
\pgfpathlineto{\pgfqpoint{6.085296in}{1.286418in}}%
\pgfpathlineto{\pgfqpoint{6.090688in}{1.254858in}}%
\pgfpathlineto{\pgfqpoint{6.096080in}{1.269652in}}%
\pgfpathlineto{\pgfqpoint{6.101472in}{1.348552in}}%
\pgfpathlineto{\pgfqpoint{6.106864in}{1.269652in}}%
\pgfpathlineto{\pgfqpoint{6.112256in}{1.301212in}}%
\pgfpathlineto{\pgfqpoint{6.117648in}{1.253872in}}%
\pgfpathlineto{\pgfqpoint{6.123040in}{1.459012in}}%
\pgfpathlineto{\pgfqpoint{6.128432in}{1.206532in}}%
\pgfpathlineto{\pgfqpoint{6.133824in}{1.080291in}}%
\pgfpathlineto{\pgfqpoint{6.139216in}{1.459012in}}%
\pgfpathlineto{\pgfqpoint{6.144608in}{1.364332in}}%
\pgfpathlineto{\pgfqpoint{6.150000in}{1.443232in}}%
\pgfpathlineto{\pgfqpoint{6.150000in}{1.443232in}}%
\pgfusepath{stroke}%
\end{pgfscope}%
\begin{pgfscope}%
\pgfsetrectcap%
\pgfsetmiterjoin%
\pgfsetlinewidth{0.803000pt}%
\definecolor{currentstroke}{rgb}{0.000000,0.000000,0.000000}%
\pgfsetstrokecolor{currentstroke}%
\pgfsetdash{}{0pt}%
\pgfpathmoveto{\pgfqpoint{0.758026in}{0.565123in}}%
\pgfpathlineto{\pgfqpoint{0.758026in}{3.350000in}}%
\pgfusepath{stroke}%
\end{pgfscope}%
\begin{pgfscope}%
\pgfsetrectcap%
\pgfsetmiterjoin%
\pgfsetlinewidth{0.803000pt}%
\definecolor{currentstroke}{rgb}{0.000000,0.000000,0.000000}%
\pgfsetstrokecolor{currentstroke}%
\pgfsetdash{}{0pt}%
\pgfpathmoveto{\pgfqpoint{6.150000in}{0.565123in}}%
\pgfpathlineto{\pgfqpoint{6.150000in}{3.350000in}}%
\pgfusepath{stroke}%
\end{pgfscope}%
\begin{pgfscope}%
\pgfsetrectcap%
\pgfsetmiterjoin%
\pgfsetlinewidth{0.803000pt}%
\definecolor{currentstroke}{rgb}{0.000000,0.000000,0.000000}%
\pgfsetstrokecolor{currentstroke}%
\pgfsetdash{}{0pt}%
\pgfpathmoveto{\pgfqpoint{0.758026in}{0.565123in}}%
\pgfpathlineto{\pgfqpoint{6.150000in}{0.565123in}}%
\pgfusepath{stroke}%
\end{pgfscope}%
\begin{pgfscope}%
\pgfsetrectcap%
\pgfsetmiterjoin%
\pgfsetlinewidth{0.803000pt}%
\definecolor{currentstroke}{rgb}{0.000000,0.000000,0.000000}%
\pgfsetstrokecolor{currentstroke}%
\pgfsetdash{}{0pt}%
\pgfpathmoveto{\pgfqpoint{0.758026in}{3.350000in}}%
\pgfpathlineto{\pgfqpoint{6.150000in}{3.350000in}}%
\pgfusepath{stroke}%
\end{pgfscope}%
\end{pgfpicture}%
\makeatother%
\endgroup%

%		\caption{Measured frequency content with a spectrum analyzer as a function of frequency of a sine signal.}
%		\label{fig:443-sine}
%	\end{figure}
%	
	\begin{figure}[H]
		\centering
		%% Creator: Matplotlib, PGF backend
%%
%% To include the figure in your LaTeX document, write
%%   \input{<filename>.pgf}
%%
%% Make sure the required packages are loaded in your preamble
%%   \usepackage{pgf}
%%
%% Also ensure that all the required font packages are loaded; for instance,
%% the lmodern package is sometimes necessary when using math font.
%%   \usepackage{lmodern}
%%
%% Figures using additional raster images can only be included by \input if
%% they are in the same directory as the main LaTeX file. For loading figures
%% from other directories you can use the `import` package
%%   \usepackage{import}
%%
%% and then include the figures with
%%   \import{<path to file>}{<filename>.pgf}
%%
%% Matplotlib used the following preamble
%%
\begingroup%
\makeatletter%
\begin{pgfpicture}%
\pgfpathrectangle{\pgfpointorigin}{\pgfqpoint{6.300000in}{3.500000in}}%
\pgfusepath{use as bounding box, clip}%
\begin{pgfscope}%
\pgfsetbuttcap%
\pgfsetmiterjoin%
\definecolor{currentfill}{rgb}{1.000000,1.000000,1.000000}%
\pgfsetfillcolor{currentfill}%
\pgfsetlinewidth{0.000000pt}%
\definecolor{currentstroke}{rgb}{1.000000,1.000000,1.000000}%
\pgfsetstrokecolor{currentstroke}%
\pgfsetdash{}{0pt}%
\pgfpathmoveto{\pgfqpoint{0.000000in}{0.000000in}}%
\pgfpathlineto{\pgfqpoint{6.300000in}{0.000000in}}%
\pgfpathlineto{\pgfqpoint{6.300000in}{3.500000in}}%
\pgfpathlineto{\pgfqpoint{0.000000in}{3.500000in}}%
\pgfpathlineto{\pgfqpoint{0.000000in}{0.000000in}}%
\pgfpathclose%
\pgfusepath{fill}%
\end{pgfscope}%
\begin{pgfscope}%
\pgfsetbuttcap%
\pgfsetmiterjoin%
\definecolor{currentfill}{rgb}{1.000000,1.000000,1.000000}%
\pgfsetfillcolor{currentfill}%
\pgfsetlinewidth{0.000000pt}%
\definecolor{currentstroke}{rgb}{0.000000,0.000000,0.000000}%
\pgfsetstrokecolor{currentstroke}%
\pgfsetstrokeopacity{0.000000}%
\pgfsetdash{}{0pt}%
\pgfpathmoveto{\pgfqpoint{0.758026in}{2.240123in}}%
\pgfpathlineto{\pgfqpoint{3.040278in}{2.240123in}}%
\pgfpathlineto{\pgfqpoint{3.040278in}{3.150926in}}%
\pgfpathlineto{\pgfqpoint{0.758026in}{3.150926in}}%
\pgfpathlineto{\pgfqpoint{0.758026in}{2.240123in}}%
\pgfpathclose%
\pgfusepath{fill}%
\end{pgfscope}%
\begin{pgfscope}%
\pgfpathrectangle{\pgfqpoint{0.758026in}{2.240123in}}{\pgfqpoint{2.282252in}{0.910803in}}%
\pgfusepath{clip}%
\pgfsetbuttcap%
\pgfsetroundjoin%
\pgfsetlinewidth{0.501875pt}%
\definecolor{currentstroke}{rgb}{0.690196,0.690196,0.690196}%
\pgfsetstrokecolor{currentstroke}%
\pgfsetstrokeopacity{0.700000}%
\pgfsetdash{{1.850000pt}{0.800000pt}}{0.000000pt}%
\pgfpathmoveto{\pgfqpoint{1.177215in}{2.240123in}}%
\pgfpathlineto{\pgfqpoint{1.177215in}{3.150926in}}%
\pgfusepath{stroke}%
\end{pgfscope}%
\begin{pgfscope}%
\pgfsetbuttcap%
\pgfsetroundjoin%
\definecolor{currentfill}{rgb}{0.000000,0.000000,0.000000}%
\pgfsetfillcolor{currentfill}%
\pgfsetlinewidth{0.803000pt}%
\definecolor{currentstroke}{rgb}{0.000000,0.000000,0.000000}%
\pgfsetstrokecolor{currentstroke}%
\pgfsetdash{}{0pt}%
\pgfsys@defobject{currentmarker}{\pgfqpoint{0.000000in}{-0.048611in}}{\pgfqpoint{0.000000in}{0.000000in}}{%
\pgfpathmoveto{\pgfqpoint{0.000000in}{0.000000in}}%
\pgfpathlineto{\pgfqpoint{0.000000in}{-0.048611in}}%
\pgfusepath{stroke,fill}%
}%
\begin{pgfscope}%
\pgfsys@transformshift{1.177215in}{2.240123in}%
\pgfsys@useobject{currentmarker}{}%
\end{pgfscope}%
\end{pgfscope}%
\begin{pgfscope}%
\definecolor{textcolor}{rgb}{0.000000,0.000000,0.000000}%
\pgfsetstrokecolor{textcolor}%
\pgfsetfillcolor{textcolor}%
\pgftext[x=1.177215in,y=2.142901in,,top]{\color{textcolor}\rmfamily\fontsize{10.000000}{12.000000}\selectfont \(\displaystyle {1}\)}%
\end{pgfscope}%
\begin{pgfscope}%
\pgfpathrectangle{\pgfqpoint{0.758026in}{2.240123in}}{\pgfqpoint{2.282252in}{0.910803in}}%
\pgfusepath{clip}%
\pgfsetbuttcap%
\pgfsetroundjoin%
\pgfsetlinewidth{0.501875pt}%
\definecolor{currentstroke}{rgb}{0.690196,0.690196,0.690196}%
\pgfsetstrokecolor{currentstroke}%
\pgfsetstrokeopacity{0.700000}%
\pgfsetdash{{1.850000pt}{0.800000pt}}{0.000000pt}%
\pgfpathmoveto{\pgfqpoint{1.642981in}{2.240123in}}%
\pgfpathlineto{\pgfqpoint{1.642981in}{3.150926in}}%
\pgfusepath{stroke}%
\end{pgfscope}%
\begin{pgfscope}%
\pgfsetbuttcap%
\pgfsetroundjoin%
\definecolor{currentfill}{rgb}{0.000000,0.000000,0.000000}%
\pgfsetfillcolor{currentfill}%
\pgfsetlinewidth{0.803000pt}%
\definecolor{currentstroke}{rgb}{0.000000,0.000000,0.000000}%
\pgfsetstrokecolor{currentstroke}%
\pgfsetdash{}{0pt}%
\pgfsys@defobject{currentmarker}{\pgfqpoint{0.000000in}{-0.048611in}}{\pgfqpoint{0.000000in}{0.000000in}}{%
\pgfpathmoveto{\pgfqpoint{0.000000in}{0.000000in}}%
\pgfpathlineto{\pgfqpoint{0.000000in}{-0.048611in}}%
\pgfusepath{stroke,fill}%
}%
\begin{pgfscope}%
\pgfsys@transformshift{1.642981in}{2.240123in}%
\pgfsys@useobject{currentmarker}{}%
\end{pgfscope}%
\end{pgfscope}%
\begin{pgfscope}%
\definecolor{textcolor}{rgb}{0.000000,0.000000,0.000000}%
\pgfsetstrokecolor{textcolor}%
\pgfsetfillcolor{textcolor}%
\pgftext[x=1.642981in,y=2.142901in,,top]{\color{textcolor}\rmfamily\fontsize{10.000000}{12.000000}\selectfont \(\displaystyle {2}\)}%
\end{pgfscope}%
\begin{pgfscope}%
\pgfpathrectangle{\pgfqpoint{0.758026in}{2.240123in}}{\pgfqpoint{2.282252in}{0.910803in}}%
\pgfusepath{clip}%
\pgfsetbuttcap%
\pgfsetroundjoin%
\pgfsetlinewidth{0.501875pt}%
\definecolor{currentstroke}{rgb}{0.690196,0.690196,0.690196}%
\pgfsetstrokecolor{currentstroke}%
\pgfsetstrokeopacity{0.700000}%
\pgfsetdash{{1.850000pt}{0.800000pt}}{0.000000pt}%
\pgfpathmoveto{\pgfqpoint{2.108746in}{2.240123in}}%
\pgfpathlineto{\pgfqpoint{2.108746in}{3.150926in}}%
\pgfusepath{stroke}%
\end{pgfscope}%
\begin{pgfscope}%
\pgfsetbuttcap%
\pgfsetroundjoin%
\definecolor{currentfill}{rgb}{0.000000,0.000000,0.000000}%
\pgfsetfillcolor{currentfill}%
\pgfsetlinewidth{0.803000pt}%
\definecolor{currentstroke}{rgb}{0.000000,0.000000,0.000000}%
\pgfsetstrokecolor{currentstroke}%
\pgfsetdash{}{0pt}%
\pgfsys@defobject{currentmarker}{\pgfqpoint{0.000000in}{-0.048611in}}{\pgfqpoint{0.000000in}{0.000000in}}{%
\pgfpathmoveto{\pgfqpoint{0.000000in}{0.000000in}}%
\pgfpathlineto{\pgfqpoint{0.000000in}{-0.048611in}}%
\pgfusepath{stroke,fill}%
}%
\begin{pgfscope}%
\pgfsys@transformshift{2.108746in}{2.240123in}%
\pgfsys@useobject{currentmarker}{}%
\end{pgfscope}%
\end{pgfscope}%
\begin{pgfscope}%
\definecolor{textcolor}{rgb}{0.000000,0.000000,0.000000}%
\pgfsetstrokecolor{textcolor}%
\pgfsetfillcolor{textcolor}%
\pgftext[x=2.108746in,y=2.142901in,,top]{\color{textcolor}\rmfamily\fontsize{10.000000}{12.000000}\selectfont \(\displaystyle {3}\)}%
\end{pgfscope}%
\begin{pgfscope}%
\pgfpathrectangle{\pgfqpoint{0.758026in}{2.240123in}}{\pgfqpoint{2.282252in}{0.910803in}}%
\pgfusepath{clip}%
\pgfsetbuttcap%
\pgfsetroundjoin%
\pgfsetlinewidth{0.501875pt}%
\definecolor{currentstroke}{rgb}{0.690196,0.690196,0.690196}%
\pgfsetstrokecolor{currentstroke}%
\pgfsetstrokeopacity{0.700000}%
\pgfsetdash{{1.850000pt}{0.800000pt}}{0.000000pt}%
\pgfpathmoveto{\pgfqpoint{2.574512in}{2.240123in}}%
\pgfpathlineto{\pgfqpoint{2.574512in}{3.150926in}}%
\pgfusepath{stroke}%
\end{pgfscope}%
\begin{pgfscope}%
\pgfsetbuttcap%
\pgfsetroundjoin%
\definecolor{currentfill}{rgb}{0.000000,0.000000,0.000000}%
\pgfsetfillcolor{currentfill}%
\pgfsetlinewidth{0.803000pt}%
\definecolor{currentstroke}{rgb}{0.000000,0.000000,0.000000}%
\pgfsetstrokecolor{currentstroke}%
\pgfsetdash{}{0pt}%
\pgfsys@defobject{currentmarker}{\pgfqpoint{0.000000in}{-0.048611in}}{\pgfqpoint{0.000000in}{0.000000in}}{%
\pgfpathmoveto{\pgfqpoint{0.000000in}{0.000000in}}%
\pgfpathlineto{\pgfqpoint{0.000000in}{-0.048611in}}%
\pgfusepath{stroke,fill}%
}%
\begin{pgfscope}%
\pgfsys@transformshift{2.574512in}{2.240123in}%
\pgfsys@useobject{currentmarker}{}%
\end{pgfscope}%
\end{pgfscope}%
\begin{pgfscope}%
\definecolor{textcolor}{rgb}{0.000000,0.000000,0.000000}%
\pgfsetstrokecolor{textcolor}%
\pgfsetfillcolor{textcolor}%
\pgftext[x=2.574512in,y=2.142901in,,top]{\color{textcolor}\rmfamily\fontsize{10.000000}{12.000000}\selectfont \(\displaystyle {4}\)}%
\end{pgfscope}%
\begin{pgfscope}%
\pgfpathrectangle{\pgfqpoint{0.758026in}{2.240123in}}{\pgfqpoint{2.282252in}{0.910803in}}%
\pgfusepath{clip}%
\pgfsetbuttcap%
\pgfsetroundjoin%
\pgfsetlinewidth{0.501875pt}%
\definecolor{currentstroke}{rgb}{0.690196,0.690196,0.690196}%
\pgfsetstrokecolor{currentstroke}%
\pgfsetstrokeopacity{0.700000}%
\pgfsetdash{{1.850000pt}{0.800000pt}}{0.000000pt}%
\pgfpathmoveto{\pgfqpoint{3.040278in}{2.240123in}}%
\pgfpathlineto{\pgfqpoint{3.040278in}{3.150926in}}%
\pgfusepath{stroke}%
\end{pgfscope}%
\begin{pgfscope}%
\pgfsetbuttcap%
\pgfsetroundjoin%
\definecolor{currentfill}{rgb}{0.000000,0.000000,0.000000}%
\pgfsetfillcolor{currentfill}%
\pgfsetlinewidth{0.803000pt}%
\definecolor{currentstroke}{rgb}{0.000000,0.000000,0.000000}%
\pgfsetstrokecolor{currentstroke}%
\pgfsetdash{}{0pt}%
\pgfsys@defobject{currentmarker}{\pgfqpoint{0.000000in}{-0.048611in}}{\pgfqpoint{0.000000in}{0.000000in}}{%
\pgfpathmoveto{\pgfqpoint{0.000000in}{0.000000in}}%
\pgfpathlineto{\pgfqpoint{0.000000in}{-0.048611in}}%
\pgfusepath{stroke,fill}%
}%
\begin{pgfscope}%
\pgfsys@transformshift{3.040278in}{2.240123in}%
\pgfsys@useobject{currentmarker}{}%
\end{pgfscope}%
\end{pgfscope}%
\begin{pgfscope}%
\definecolor{textcolor}{rgb}{0.000000,0.000000,0.000000}%
\pgfsetstrokecolor{textcolor}%
\pgfsetfillcolor{textcolor}%
\pgftext[x=3.040278in,y=2.142901in,,top]{\color{textcolor}\rmfamily\fontsize{10.000000}{12.000000}\selectfont \(\displaystyle {5}\)}%
\end{pgfscope}%
\begin{pgfscope}%
\definecolor{textcolor}{rgb}{0.000000,0.000000,0.000000}%
\pgfsetstrokecolor{textcolor}%
\pgfsetfillcolor{textcolor}%
\pgftext[x=1.899152in,y=1.963889in,,top]{\color{textcolor}\rmfamily\fontsize{10.000000}{12.000000}\selectfont Frequency (Hz)}%
\end{pgfscope}%
\begin{pgfscope}%
\definecolor{textcolor}{rgb}{0.000000,0.000000,0.000000}%
\pgfsetstrokecolor{textcolor}%
\pgfsetfillcolor{textcolor}%
\pgftext[x=3.040278in,y=1.977778in,right,top]{\color{textcolor}\rmfamily\fontsize{10.000000}{12.000000}\selectfont \(\displaystyle \times{10^{6}}{}\)}%
\end{pgfscope}%
\begin{pgfscope}%
\pgfpathrectangle{\pgfqpoint{0.758026in}{2.240123in}}{\pgfqpoint{2.282252in}{0.910803in}}%
\pgfusepath{clip}%
\pgfsetbuttcap%
\pgfsetroundjoin%
\pgfsetlinewidth{0.501875pt}%
\definecolor{currentstroke}{rgb}{0.690196,0.690196,0.690196}%
\pgfsetstrokecolor{currentstroke}%
\pgfsetstrokeopacity{0.700000}%
\pgfsetdash{{1.850000pt}{0.800000pt}}{0.000000pt}%
\pgfpathmoveto{\pgfqpoint{0.758026in}{2.340229in}}%
\pgfpathlineto{\pgfqpoint{3.040278in}{2.340229in}}%
\pgfusepath{stroke}%
\end{pgfscope}%
\begin{pgfscope}%
\pgfsetbuttcap%
\pgfsetroundjoin%
\definecolor{currentfill}{rgb}{0.000000,0.000000,0.000000}%
\pgfsetfillcolor{currentfill}%
\pgfsetlinewidth{0.803000pt}%
\definecolor{currentstroke}{rgb}{0.000000,0.000000,0.000000}%
\pgfsetstrokecolor{currentstroke}%
\pgfsetdash{}{0pt}%
\pgfsys@defobject{currentmarker}{\pgfqpoint{-0.048611in}{0.000000in}}{\pgfqpoint{-0.000000in}{0.000000in}}{%
\pgfpathmoveto{\pgfqpoint{-0.000000in}{0.000000in}}%
\pgfpathlineto{\pgfqpoint{-0.048611in}{0.000000in}}%
\pgfusepath{stroke,fill}%
}%
\begin{pgfscope}%
\pgfsys@transformshift{0.758026in}{2.340229in}%
\pgfsys@useobject{currentmarker}{}%
\end{pgfscope}%
\end{pgfscope}%
\begin{pgfscope}%
\definecolor{textcolor}{rgb}{0.000000,0.000000,0.000000}%
\pgfsetstrokecolor{textcolor}%
\pgfsetfillcolor{textcolor}%
\pgftext[x=0.344444in, y=2.292003in, left, base]{\color{textcolor}\rmfamily\fontsize{10.000000}{12.000000}\selectfont \(\displaystyle {\ensuremath{-}100}\)}%
\end{pgfscope}%
\begin{pgfscope}%
\pgfpathrectangle{\pgfqpoint{0.758026in}{2.240123in}}{\pgfqpoint{2.282252in}{0.910803in}}%
\pgfusepath{clip}%
\pgfsetbuttcap%
\pgfsetroundjoin%
\pgfsetlinewidth{0.501875pt}%
\definecolor{currentstroke}{rgb}{0.690196,0.690196,0.690196}%
\pgfsetstrokecolor{currentstroke}%
\pgfsetstrokeopacity{0.700000}%
\pgfsetdash{{1.850000pt}{0.800000pt}}{0.000000pt}%
\pgfpathmoveto{\pgfqpoint{0.758026in}{2.856319in}}%
\pgfpathlineto{\pgfqpoint{3.040278in}{2.856319in}}%
\pgfusepath{stroke}%
\end{pgfscope}%
\begin{pgfscope}%
\pgfsetbuttcap%
\pgfsetroundjoin%
\definecolor{currentfill}{rgb}{0.000000,0.000000,0.000000}%
\pgfsetfillcolor{currentfill}%
\pgfsetlinewidth{0.803000pt}%
\definecolor{currentstroke}{rgb}{0.000000,0.000000,0.000000}%
\pgfsetstrokecolor{currentstroke}%
\pgfsetdash{}{0pt}%
\pgfsys@defobject{currentmarker}{\pgfqpoint{-0.048611in}{0.000000in}}{\pgfqpoint{-0.000000in}{0.000000in}}{%
\pgfpathmoveto{\pgfqpoint{-0.000000in}{0.000000in}}%
\pgfpathlineto{\pgfqpoint{-0.048611in}{0.000000in}}%
\pgfusepath{stroke,fill}%
}%
\begin{pgfscope}%
\pgfsys@transformshift{0.758026in}{2.856319in}%
\pgfsys@useobject{currentmarker}{}%
\end{pgfscope}%
\end{pgfscope}%
\begin{pgfscope}%
\definecolor{textcolor}{rgb}{0.000000,0.000000,0.000000}%
\pgfsetstrokecolor{textcolor}%
\pgfsetfillcolor{textcolor}%
\pgftext[x=0.413889in, y=2.808094in, left, base]{\color{textcolor}\rmfamily\fontsize{10.000000}{12.000000}\selectfont \(\displaystyle {\ensuremath{-}50}\)}%
\end{pgfscope}%
\begin{pgfscope}%
\definecolor{textcolor}{rgb}{0.000000,0.000000,0.000000}%
\pgfsetstrokecolor{textcolor}%
\pgfsetfillcolor{textcolor}%
\pgftext[x=0.288889in,y=2.695525in,,bottom,rotate=90.000000]{\color{textcolor}\rmfamily\fontsize{10.000000}{12.000000}\selectfont Amplitude (dB)}%
\end{pgfscope}%
\begin{pgfscope}%
\pgfpathrectangle{\pgfqpoint{0.758026in}{2.240123in}}{\pgfqpoint{2.282252in}{0.910803in}}%
\pgfusepath{clip}%
\pgfsetrectcap%
\pgfsetroundjoin%
\pgfsetlinewidth{1.003750pt}%
\definecolor{currentstroke}{rgb}{0.000000,0.000000,0.000000}%
\pgfsetstrokecolor{currentstroke}%
\pgfsetdash{}{0pt}%
\pgfpathmoveto{\pgfqpoint{0.758026in}{2.590533in}}%
\pgfpathlineto{\pgfqpoint{0.760308in}{2.585372in}}%
\pgfpathlineto{\pgfqpoint{0.762590in}{2.538601in}}%
\pgfpathlineto{\pgfqpoint{0.764872in}{2.554084in}}%
\pgfpathlineto{\pgfqpoint{0.767155in}{2.522796in}}%
\pgfpathlineto{\pgfqpoint{0.769437in}{2.558600in}}%
\pgfpathlineto{\pgfqpoint{0.771719in}{2.537956in}}%
\pgfpathlineto{\pgfqpoint{0.774001in}{2.584082in}}%
\pgfpathlineto{\pgfqpoint{0.776284in}{2.542472in}}%
\pgfpathlineto{\pgfqpoint{0.778566in}{2.536988in}}%
\pgfpathlineto{\pgfqpoint{0.780848in}{2.536988in}}%
\pgfpathlineto{\pgfqpoint{0.783130in}{2.572792in}}%
\pgfpathlineto{\pgfqpoint{0.785413in}{2.520860in}}%
\pgfpathlineto{\pgfqpoint{0.787695in}{2.582791in}}%
\pgfpathlineto{\pgfqpoint{0.792259in}{2.535698in}}%
\pgfpathlineto{\pgfqpoint{0.794542in}{2.540859in}}%
\pgfpathlineto{\pgfqpoint{0.796824in}{2.566341in}}%
\pgfpathlineto{\pgfqpoint{0.799106in}{2.535053in}}%
\pgfpathlineto{\pgfqpoint{0.801388in}{2.581501in}}%
\pgfpathlineto{\pgfqpoint{0.803671in}{2.534730in}}%
\pgfpathlineto{\pgfqpoint{0.805953in}{2.616982in}}%
\pgfpathlineto{\pgfqpoint{0.808235in}{2.549891in}}%
\pgfpathlineto{\pgfqpoint{0.810517in}{2.554729in}}%
\pgfpathlineto{\pgfqpoint{0.812800in}{2.533763in}}%
\pgfpathlineto{\pgfqpoint{0.815082in}{2.580211in}}%
\pgfpathlineto{\pgfqpoint{0.817364in}{2.548923in}}%
\pgfpathlineto{\pgfqpoint{0.819646in}{2.533118in}}%
\pgfpathlineto{\pgfqpoint{0.821929in}{2.584727in}}%
\pgfpathlineto{\pgfqpoint{0.824211in}{2.517312in}}%
\pgfpathlineto{\pgfqpoint{0.828776in}{2.599242in}}%
\pgfpathlineto{\pgfqpoint{0.831058in}{2.552794in}}%
\pgfpathlineto{\pgfqpoint{0.833340in}{2.552471in}}%
\pgfpathlineto{\pgfqpoint{0.835622in}{2.536666in}}%
\pgfpathlineto{\pgfqpoint{0.837905in}{2.681171in}}%
\pgfpathlineto{\pgfqpoint{0.840187in}{2.546665in}}%
\pgfpathlineto{\pgfqpoint{0.842469in}{2.634078in}}%
\pgfpathlineto{\pgfqpoint{0.844751in}{2.551503in}}%
\pgfpathlineto{\pgfqpoint{0.847034in}{2.546020in}}%
\pgfpathlineto{\pgfqpoint{0.849316in}{2.545697in}}%
\pgfpathlineto{\pgfqpoint{0.851598in}{2.592145in}}%
\pgfpathlineto{\pgfqpoint{0.853880in}{2.550536in}}%
\pgfpathlineto{\pgfqpoint{0.856163in}{2.622466in}}%
\pgfpathlineto{\pgfqpoint{0.858445in}{2.576018in}}%
\pgfpathlineto{\pgfqpoint{0.860727in}{2.596339in}}%
\pgfpathlineto{\pgfqpoint{0.863009in}{2.528924in}}%
\pgfpathlineto{\pgfqpoint{0.865292in}{2.528924in}}%
\pgfpathlineto{\pgfqpoint{0.867574in}{2.513119in}}%
\pgfpathlineto{\pgfqpoint{0.869856in}{2.543762in}}%
\pgfpathlineto{\pgfqpoint{0.872138in}{2.600209in}}%
\pgfpathlineto{\pgfqpoint{0.874421in}{2.610531in}}%
\pgfpathlineto{\pgfqpoint{0.876703in}{2.543117in}}%
\pgfpathlineto{\pgfqpoint{0.878985in}{2.542794in}}%
\pgfpathlineto{\pgfqpoint{0.881267in}{2.635691in}}%
\pgfpathlineto{\pgfqpoint{0.883550in}{2.511506in}}%
\pgfpathlineto{\pgfqpoint{0.885832in}{2.578275in}}%
\pgfpathlineto{\pgfqpoint{0.888114in}{2.578275in}}%
\pgfpathlineto{\pgfqpoint{0.890396in}{2.732780in}}%
\pgfpathlineto{\pgfqpoint{0.892679in}{2.753101in}}%
\pgfpathlineto{\pgfqpoint{0.894961in}{2.649883in}}%
\pgfpathlineto{\pgfqpoint{0.899525in}{2.711169in}}%
\pgfpathlineto{\pgfqpoint{0.901808in}{2.597629in}}%
\pgfpathlineto{\pgfqpoint{0.904090in}{2.545697in}}%
\pgfpathlineto{\pgfqpoint{0.906372in}{2.571179in}}%
\pgfpathlineto{\pgfqpoint{0.908654in}{2.617627in}}%
\pgfpathlineto{\pgfqpoint{0.910937in}{2.555374in}}%
\pgfpathlineto{\pgfqpoint{0.913219in}{2.524086in}}%
\pgfpathlineto{\pgfqpoint{0.915501in}{2.575695in}}%
\pgfpathlineto{\pgfqpoint{0.917783in}{2.539246in}}%
\pgfpathlineto{\pgfqpoint{0.920066in}{2.559567in}}%
\pgfpathlineto{\pgfqpoint{0.922348in}{2.554406in}}%
\pgfpathlineto{\pgfqpoint{0.924630in}{2.554084in}}%
\pgfpathlineto{\pgfqpoint{0.926912in}{2.538278in}}%
\pgfpathlineto{\pgfqpoint{0.929195in}{2.600209in}}%
\pgfpathlineto{\pgfqpoint{0.931477in}{2.543117in}}%
\pgfpathlineto{\pgfqpoint{0.933759in}{2.558277in}}%
\pgfpathlineto{\pgfqpoint{0.936041in}{2.537633in}}%
\pgfpathlineto{\pgfqpoint{0.938324in}{2.867609in}}%
\pgfpathlineto{\pgfqpoint{0.940606in}{3.032435in}}%
\pgfpathlineto{\pgfqpoint{0.942888in}{3.099204in}}%
\pgfpathlineto{\pgfqpoint{0.945170in}{3.109526in}}%
\pgfpathlineto{\pgfqpoint{0.947453in}{3.073077in}}%
\pgfpathlineto{\pgfqpoint{0.949735in}{2.954054in}}%
\pgfpathlineto{\pgfqpoint{0.952017in}{2.654721in}}%
\pgfpathlineto{\pgfqpoint{0.954299in}{2.551181in}}%
\pgfpathlineto{\pgfqpoint{0.956582in}{2.617950in}}%
\pgfpathlineto{\pgfqpoint{0.958864in}{2.643754in}}%
\pgfpathlineto{\pgfqpoint{0.961146in}{2.519570in}}%
\pgfpathlineto{\pgfqpoint{0.963428in}{2.539891in}}%
\pgfpathlineto{\pgfqpoint{0.967993in}{2.503442in}}%
\pgfpathlineto{\pgfqpoint{0.970275in}{2.704395in}}%
\pgfpathlineto{\pgfqpoint{0.972557in}{2.647625in}}%
\pgfpathlineto{\pgfqpoint{0.974840in}{2.502797in}}%
\pgfpathlineto{\pgfqpoint{0.999944in}{2.500217in}}%
\pgfpathlineto{\pgfqpoint{1.011356in}{2.499249in}}%
\pgfpathlineto{\pgfqpoint{1.022767in}{2.498281in}}%
\pgfpathlineto{\pgfqpoint{1.025049in}{2.544407in}}%
\pgfpathlineto{\pgfqpoint{1.027331in}{2.497636in}}%
\pgfpathlineto{\pgfqpoint{1.029614in}{2.497636in}}%
\pgfpathlineto{\pgfqpoint{1.031896in}{2.641819in}}%
\pgfpathlineto{\pgfqpoint{1.034178in}{2.435060in}}%
\pgfpathlineto{\pgfqpoint{1.038743in}{2.434738in}}%
\pgfpathlineto{\pgfqpoint{1.041025in}{2.470542in}}%
\pgfpathlineto{\pgfqpoint{1.043307in}{2.434415in}}%
\pgfpathlineto{\pgfqpoint{1.045589in}{2.449575in}}%
\pgfpathlineto{\pgfqpoint{1.047872in}{2.433770in}}%
\pgfpathlineto{\pgfqpoint{1.050154in}{2.449253in}}%
\pgfpathlineto{\pgfqpoint{1.052436in}{2.454091in}}%
\pgfpathlineto{\pgfqpoint{1.054718in}{2.448608in}}%
\pgfpathlineto{\pgfqpoint{1.057001in}{2.432802in}}%
\pgfpathlineto{\pgfqpoint{1.068412in}{2.431835in}}%
\pgfpathlineto{\pgfqpoint{1.084388in}{2.430222in}}%
\pgfpathlineto{\pgfqpoint{1.086670in}{2.414739in}}%
\pgfpathlineto{\pgfqpoint{1.088952in}{2.450543in}}%
\pgfpathlineto{\pgfqpoint{1.091234in}{2.434738in}}%
\pgfpathlineto{\pgfqpoint{1.093517in}{2.398611in}}%
\pgfpathlineto{\pgfqpoint{1.095799in}{2.413772in}}%
\pgfpathlineto{\pgfqpoint{1.102646in}{2.413127in}}%
\pgfpathlineto{\pgfqpoint{1.104928in}{2.474735in}}%
\pgfpathlineto{\pgfqpoint{1.107210in}{2.427964in}}%
\pgfpathlineto{\pgfqpoint{1.109492in}{2.427964in}}%
\pgfpathlineto{\pgfqpoint{1.111775in}{2.396676in}}%
\pgfpathlineto{\pgfqpoint{1.114057in}{2.396354in}}%
\pgfpathlineto{\pgfqpoint{1.116339in}{2.411836in}}%
\pgfpathlineto{\pgfqpoint{1.118622in}{2.416675in}}%
\pgfpathlineto{\pgfqpoint{1.120904in}{2.416352in}}%
\pgfpathlineto{\pgfqpoint{1.123186in}{2.395708in}}%
\pgfpathlineto{\pgfqpoint{1.125468in}{2.441834in}}%
\pgfpathlineto{\pgfqpoint{1.130033in}{2.415707in}}%
\pgfpathlineto{\pgfqpoint{1.132315in}{2.430867in}}%
\pgfpathlineto{\pgfqpoint{1.136880in}{2.404740in}}%
\pgfpathlineto{\pgfqpoint{1.139162in}{2.466348in}}%
\pgfpathlineto{\pgfqpoint{1.141444in}{2.398934in}}%
\pgfpathlineto{\pgfqpoint{1.143726in}{2.398934in}}%
\pgfpathlineto{\pgfqpoint{1.148291in}{2.542794in}}%
\pgfpathlineto{\pgfqpoint{1.150573in}{2.418610in}}%
\pgfpathlineto{\pgfqpoint{1.152855in}{2.444415in}}%
\pgfpathlineto{\pgfqpoint{1.155138in}{2.449253in}}%
\pgfpathlineto{\pgfqpoint{1.159702in}{2.366356in}}%
\pgfpathlineto{\pgfqpoint{1.164267in}{2.505055in}}%
\pgfpathlineto{\pgfqpoint{1.166549in}{2.381193in}}%
\pgfpathlineto{\pgfqpoint{1.168831in}{2.396354in}}%
\pgfpathlineto{\pgfqpoint{1.171113in}{2.468284in}}%
\pgfpathlineto{\pgfqpoint{1.173396in}{2.442479in}}%
\pgfpathlineto{\pgfqpoint{1.177960in}{2.421190in}}%
\pgfpathlineto{\pgfqpoint{1.180242in}{2.421190in}}%
\pgfpathlineto{\pgfqpoint{1.182525in}{2.405708in}}%
\pgfpathlineto{\pgfqpoint{1.184807in}{2.446995in}}%
\pgfpathlineto{\pgfqpoint{1.189371in}{2.472799in}}%
\pgfpathlineto{\pgfqpoint{1.191654in}{2.426351in}}%
\pgfpathlineto{\pgfqpoint{1.193936in}{2.498604in}}%
\pgfpathlineto{\pgfqpoint{1.196218in}{2.467639in}}%
\pgfpathlineto{\pgfqpoint{1.198500in}{2.410869in}}%
\pgfpathlineto{\pgfqpoint{1.200783in}{2.539891in}}%
\pgfpathlineto{\pgfqpoint{1.203065in}{2.426351in}}%
\pgfpathlineto{\pgfqpoint{1.205347in}{2.472799in}}%
\pgfpathlineto{\pgfqpoint{1.207629in}{2.493443in}}%
\pgfpathlineto{\pgfqpoint{1.209912in}{2.498604in}}%
\pgfpathlineto{\pgfqpoint{1.212194in}{2.472799in}}%
\pgfpathlineto{\pgfqpoint{1.214476in}{2.498604in}}%
\pgfpathlineto{\pgfqpoint{1.216758in}{2.503765in}}%
\pgfpathlineto{\pgfqpoint{1.219041in}{2.524409in}}%
\pgfpathlineto{\pgfqpoint{1.221323in}{2.462478in}}%
\pgfpathlineto{\pgfqpoint{1.223605in}{2.529569in}}%
\pgfpathlineto{\pgfqpoint{1.225887in}{2.457317in}}%
\pgfpathlineto{\pgfqpoint{1.228170in}{2.477960in}}%
\pgfpathlineto{\pgfqpoint{1.230452in}{2.477960in}}%
\pgfpathlineto{\pgfqpoint{1.232734in}{2.483121in}}%
\pgfpathlineto{\pgfqpoint{1.235016in}{2.467639in}}%
\pgfpathlineto{\pgfqpoint{1.237299in}{2.477960in}}%
\pgfpathlineto{\pgfqpoint{1.239581in}{2.472799in}}%
\pgfpathlineto{\pgfqpoint{1.241863in}{2.462478in}}%
\pgfpathlineto{\pgfqpoint{1.244145in}{2.416030in}}%
\pgfpathlineto{\pgfqpoint{1.246428in}{2.446995in}}%
\pgfpathlineto{\pgfqpoint{1.248710in}{2.436673in}}%
\pgfpathlineto{\pgfqpoint{1.250992in}{2.431512in}}%
\pgfpathlineto{\pgfqpoint{1.253274in}{2.452156in}}%
\pgfpathlineto{\pgfqpoint{1.255557in}{2.421190in}}%
\pgfpathlineto{\pgfqpoint{1.257839in}{2.426351in}}%
\pgfpathlineto{\pgfqpoint{1.260121in}{2.436351in}}%
\pgfpathlineto{\pgfqpoint{1.262403in}{2.498281in}}%
\pgfpathlineto{\pgfqpoint{1.264686in}{2.420868in}}%
\pgfpathlineto{\pgfqpoint{1.266968in}{2.415707in}}%
\pgfpathlineto{\pgfqpoint{1.269250in}{2.431190in}}%
\pgfpathlineto{\pgfqpoint{1.271532in}{2.426029in}}%
\pgfpathlineto{\pgfqpoint{1.273815in}{2.317650in}}%
\pgfpathlineto{\pgfqpoint{1.276097in}{2.431190in}}%
\pgfpathlineto{\pgfqpoint{1.278379in}{2.374420in}}%
\pgfpathlineto{\pgfqpoint{1.280661in}{2.281523in}}%
\pgfpathlineto{\pgfqpoint{1.282944in}{2.415707in}}%
\pgfpathlineto{\pgfqpoint{1.285226in}{2.415707in}}%
\pgfpathlineto{\pgfqpoint{1.287508in}{2.405385in}}%
\pgfpathlineto{\pgfqpoint{1.289790in}{2.353776in}}%
\pgfpathlineto{\pgfqpoint{1.292073in}{2.446672in}}%
\pgfpathlineto{\pgfqpoint{1.294355in}{2.400224in}}%
\pgfpathlineto{\pgfqpoint{1.296637in}{2.560212in}}%
\pgfpathlineto{\pgfqpoint{1.298919in}{2.446672in}}%
\pgfpathlineto{\pgfqpoint{1.301202in}{2.462155in}}%
\pgfpathlineto{\pgfqpoint{1.303484in}{2.420868in}}%
\pgfpathlineto{\pgfqpoint{1.305766in}{2.456994in}}%
\pgfpathlineto{\pgfqpoint{1.308048in}{2.431190in}}%
\pgfpathlineto{\pgfqpoint{1.310331in}{2.426029in}}%
\pgfpathlineto{\pgfqpoint{1.312613in}{2.405385in}}%
\pgfpathlineto{\pgfqpoint{1.317177in}{2.446672in}}%
\pgfpathlineto{\pgfqpoint{1.319460in}{2.451833in}}%
\pgfpathlineto{\pgfqpoint{1.321742in}{2.436351in}}%
\pgfpathlineto{\pgfqpoint{1.324024in}{2.410546in}}%
\pgfpathlineto{\pgfqpoint{1.326306in}{2.415707in}}%
\pgfpathlineto{\pgfqpoint{1.330871in}{2.467316in}}%
\pgfpathlineto{\pgfqpoint{1.333153in}{2.436351in}}%
\pgfpathlineto{\pgfqpoint{1.335435in}{2.482799in}}%
\pgfpathlineto{\pgfqpoint{1.337718in}{2.420545in}}%
\pgfpathlineto{\pgfqpoint{1.340000in}{2.461833in}}%
\pgfpathlineto{\pgfqpoint{1.342282in}{2.482476in}}%
\pgfpathlineto{\pgfqpoint{1.344564in}{2.461833in}}%
\pgfpathlineto{\pgfqpoint{1.346847in}{2.461833in}}%
\pgfpathlineto{\pgfqpoint{1.349129in}{2.441189in}}%
\pgfpathlineto{\pgfqpoint{1.351411in}{2.410224in}}%
\pgfpathlineto{\pgfqpoint{1.353693in}{2.451511in}}%
\pgfpathlineto{\pgfqpoint{1.355976in}{2.436028in}}%
\pgfpathlineto{\pgfqpoint{1.358258in}{2.575372in}}%
\pgfpathlineto{\pgfqpoint{1.360540in}{2.399902in}}%
\pgfpathlineto{\pgfqpoint{1.362822in}{2.384419in}}%
\pgfpathlineto{\pgfqpoint{1.365105in}{2.456672in}}%
\pgfpathlineto{\pgfqpoint{1.367387in}{2.446350in}}%
\pgfpathlineto{\pgfqpoint{1.371951in}{2.379258in}}%
\pgfpathlineto{\pgfqpoint{1.374234in}{2.461833in}}%
\pgfpathlineto{\pgfqpoint{1.376516in}{2.410224in}}%
\pgfpathlineto{\pgfqpoint{1.378798in}{2.451511in}}%
\pgfpathlineto{\pgfqpoint{1.381080in}{2.466993in}}%
\pgfpathlineto{\pgfqpoint{1.383363in}{2.384419in}}%
\pgfpathlineto{\pgfqpoint{1.385645in}{2.456672in}}%
\pgfpathlineto{\pgfqpoint{1.387927in}{2.461833in}}%
\pgfpathlineto{\pgfqpoint{1.392492in}{2.405063in}}%
\pgfpathlineto{\pgfqpoint{1.394774in}{2.425706in}}%
\pgfpathlineto{\pgfqpoint{1.397056in}{2.436028in}}%
\pgfpathlineto{\pgfqpoint{1.399338in}{2.394741in}}%
\pgfpathlineto{\pgfqpoint{1.403903in}{2.482476in}}%
\pgfpathlineto{\pgfqpoint{1.408467in}{2.374097in}}%
\pgfpathlineto{\pgfqpoint{1.410750in}{2.456672in}}%
\pgfpathlineto{\pgfqpoint{1.413032in}{2.482476in}}%
\pgfpathlineto{\pgfqpoint{1.415314in}{2.466993in}}%
\pgfpathlineto{\pgfqpoint{1.417597in}{2.436028in}}%
\pgfpathlineto{\pgfqpoint{1.419879in}{2.471832in}}%
\pgfpathlineto{\pgfqpoint{1.422161in}{2.415062in}}%
\pgfpathlineto{\pgfqpoint{1.424443in}{2.415062in}}%
\pgfpathlineto{\pgfqpoint{1.429008in}{2.404740in}}%
\pgfpathlineto{\pgfqpoint{1.431290in}{2.451188in}}%
\pgfpathlineto{\pgfqpoint{1.435855in}{2.451188in}}%
\pgfpathlineto{\pgfqpoint{1.438137in}{2.420223in}}%
\pgfpathlineto{\pgfqpoint{1.440419in}{2.363453in}}%
\pgfpathlineto{\pgfqpoint{1.442701in}{2.373775in}}%
\pgfpathlineto{\pgfqpoint{1.449548in}{2.482154in}}%
\pgfpathlineto{\pgfqpoint{1.451830in}{2.363453in}}%
\pgfpathlineto{\pgfqpoint{1.454113in}{2.466671in}}%
\pgfpathlineto{\pgfqpoint{1.456395in}{2.415062in}}%
\pgfpathlineto{\pgfqpoint{1.458677in}{2.435705in}}%
\pgfpathlineto{\pgfqpoint{1.460959in}{2.466671in}}%
\pgfpathlineto{\pgfqpoint{1.463242in}{2.451188in}}%
\pgfpathlineto{\pgfqpoint{1.465524in}{2.425384in}}%
\pgfpathlineto{\pgfqpoint{1.467806in}{2.451188in}}%
\pgfpathlineto{\pgfqpoint{1.470088in}{2.394418in}}%
\pgfpathlineto{\pgfqpoint{1.472371in}{2.404740in}}%
\pgfpathlineto{\pgfqpoint{1.474653in}{2.399579in}}%
\pgfpathlineto{\pgfqpoint{1.476935in}{2.456349in}}%
\pgfpathlineto{\pgfqpoint{1.479217in}{2.399579in}}%
\pgfpathlineto{\pgfqpoint{1.481500in}{2.404740in}}%
\pgfpathlineto{\pgfqpoint{1.486064in}{2.487315in}}%
\pgfpathlineto{\pgfqpoint{1.488346in}{2.440866in}}%
\pgfpathlineto{\pgfqpoint{1.492911in}{2.476993in}}%
\pgfpathlineto{\pgfqpoint{1.495193in}{2.373775in}}%
\pgfpathlineto{\pgfqpoint{1.497475in}{2.415062in}}%
\pgfpathlineto{\pgfqpoint{1.499758in}{2.476670in}}%
\pgfpathlineto{\pgfqpoint{1.502040in}{2.414739in}}%
\pgfpathlineto{\pgfqpoint{1.504322in}{2.414739in}}%
\pgfpathlineto{\pgfqpoint{1.508887in}{2.388935in}}%
\pgfpathlineto{\pgfqpoint{1.511169in}{2.461187in}}%
\pgfpathlineto{\pgfqpoint{1.513451in}{2.404418in}}%
\pgfpathlineto{\pgfqpoint{1.515733in}{2.456027in}}%
\pgfpathlineto{\pgfqpoint{1.518016in}{2.471509in}}%
\pgfpathlineto{\pgfqpoint{1.520298in}{2.430222in}}%
\pgfpathlineto{\pgfqpoint{1.522580in}{2.414739in}}%
\pgfpathlineto{\pgfqpoint{1.524862in}{2.430222in}}%
\pgfpathlineto{\pgfqpoint{1.527145in}{2.461187in}}%
\pgfpathlineto{\pgfqpoint{1.529427in}{2.404418in}}%
\pgfpathlineto{\pgfqpoint{1.531709in}{2.486992in}}%
\pgfpathlineto{\pgfqpoint{1.533991in}{2.404418in}}%
\pgfpathlineto{\pgfqpoint{1.536274in}{2.388935in}}%
\pgfpathlineto{\pgfqpoint{1.538556in}{2.342487in}}%
\pgfpathlineto{\pgfqpoint{1.540838in}{2.456027in}}%
\pgfpathlineto{\pgfqpoint{1.545403in}{2.409578in}}%
\pgfpathlineto{\pgfqpoint{1.547685in}{2.440544in}}%
\pgfpathlineto{\pgfqpoint{1.549967in}{2.414739in}}%
\pgfpathlineto{\pgfqpoint{1.552249in}{2.492153in}}%
\pgfpathlineto{\pgfqpoint{1.554532in}{2.394096in}}%
\pgfpathlineto{\pgfqpoint{1.556814in}{2.456027in}}%
\pgfpathlineto{\pgfqpoint{1.559096in}{2.466348in}}%
\pgfpathlineto{\pgfqpoint{1.561378in}{2.492153in}}%
\pgfpathlineto{\pgfqpoint{1.563661in}{2.409578in}}%
\pgfpathlineto{\pgfqpoint{1.565943in}{2.481831in}}%
\pgfpathlineto{\pgfqpoint{1.568225in}{2.440544in}}%
\pgfpathlineto{\pgfqpoint{1.570507in}{2.450866in}}%
\pgfpathlineto{\pgfqpoint{1.572790in}{2.435383in}}%
\pgfpathlineto{\pgfqpoint{1.575072in}{2.466348in}}%
\pgfpathlineto{\pgfqpoint{1.577354in}{2.321520in}}%
\pgfpathlineto{\pgfqpoint{1.579636in}{2.471187in}}%
\pgfpathlineto{\pgfqpoint{1.581919in}{2.424739in}}%
\pgfpathlineto{\pgfqpoint{1.584201in}{2.455704in}}%
\pgfpathlineto{\pgfqpoint{1.586483in}{2.373130in}}%
\pgfpathlineto{\pgfqpoint{1.588765in}{2.429899in}}%
\pgfpathlineto{\pgfqpoint{1.591048in}{2.398934in}}%
\pgfpathlineto{\pgfqpoint{1.593330in}{2.476348in}}%
\pgfpathlineto{\pgfqpoint{1.595612in}{2.491830in}}%
\pgfpathlineto{\pgfqpoint{1.597894in}{2.481509in}}%
\pgfpathlineto{\pgfqpoint{1.600177in}{2.476348in}}%
\pgfpathlineto{\pgfqpoint{1.602459in}{2.445382in}}%
\pgfpathlineto{\pgfqpoint{1.604741in}{2.440221in}}%
\pgfpathlineto{\pgfqpoint{1.611588in}{2.512474in}}%
\pgfpathlineto{\pgfqpoint{1.613870in}{2.466026in}}%
\pgfpathlineto{\pgfqpoint{1.616152in}{2.476348in}}%
\pgfpathlineto{\pgfqpoint{1.618435in}{2.440221in}}%
\pgfpathlineto{\pgfqpoint{1.620717in}{2.476348in}}%
\pgfpathlineto{\pgfqpoint{1.622999in}{2.471187in}}%
\pgfpathlineto{\pgfqpoint{1.625281in}{2.455704in}}%
\pgfpathlineto{\pgfqpoint{1.627564in}{2.491830in}}%
\pgfpathlineto{\pgfqpoint{1.629846in}{2.471187in}}%
\pgfpathlineto{\pgfqpoint{1.632128in}{2.476348in}}%
\pgfpathlineto{\pgfqpoint{1.634410in}{2.445382in}}%
\pgfpathlineto{\pgfqpoint{1.636693in}{2.476348in}}%
\pgfpathlineto{\pgfqpoint{1.638975in}{2.527957in}}%
\pgfpathlineto{\pgfqpoint{1.641257in}{2.455704in}}%
\pgfpathlineto{\pgfqpoint{1.643539in}{2.455704in}}%
\pgfpathlineto{\pgfqpoint{1.645822in}{2.373130in}}%
\pgfpathlineto{\pgfqpoint{1.648104in}{2.491830in}}%
\pgfpathlineto{\pgfqpoint{1.650386in}{2.481509in}}%
\pgfpathlineto{\pgfqpoint{1.654951in}{2.435060in}}%
\pgfpathlineto{\pgfqpoint{1.657233in}{2.496991in}}%
\pgfpathlineto{\pgfqpoint{1.659515in}{2.476025in}}%
\pgfpathlineto{\pgfqpoint{1.661797in}{2.481186in}}%
\pgfpathlineto{\pgfqpoint{1.664080in}{2.439899in}}%
\pgfpathlineto{\pgfqpoint{1.666362in}{2.377968in}}%
\pgfpathlineto{\pgfqpoint{1.668644in}{2.450221in}}%
\pgfpathlineto{\pgfqpoint{1.670926in}{2.486347in}}%
\pgfpathlineto{\pgfqpoint{1.673209in}{2.486347in}}%
\pgfpathlineto{\pgfqpoint{1.677773in}{2.383129in}}%
\pgfpathlineto{\pgfqpoint{1.680055in}{2.512151in}}%
\pgfpathlineto{\pgfqpoint{1.682338in}{2.455381in}}%
\pgfpathlineto{\pgfqpoint{1.684620in}{2.460542in}}%
\pgfpathlineto{\pgfqpoint{1.686902in}{2.450221in}}%
\pgfpathlineto{\pgfqpoint{1.689184in}{2.460542in}}%
\pgfpathlineto{\pgfqpoint{1.691467in}{2.501830in}}%
\pgfpathlineto{\pgfqpoint{1.693749in}{2.408933in}}%
\pgfpathlineto{\pgfqpoint{1.696031in}{2.455381in}}%
\pgfpathlineto{\pgfqpoint{1.698313in}{2.439899in}}%
\pgfpathlineto{\pgfqpoint{1.700596in}{2.445060in}}%
\pgfpathlineto{\pgfqpoint{1.702878in}{2.455381in}}%
\pgfpathlineto{\pgfqpoint{1.705160in}{2.414094in}}%
\pgfpathlineto{\pgfqpoint{1.707443in}{2.470864in}}%
\pgfpathlineto{\pgfqpoint{1.709725in}{2.460542in}}%
\pgfpathlineto{\pgfqpoint{1.712007in}{2.476025in}}%
\pgfpathlineto{\pgfqpoint{1.714289in}{2.470864in}}%
\pgfpathlineto{\pgfqpoint{1.718854in}{2.481186in}}%
\pgfpathlineto{\pgfqpoint{1.721136in}{2.476025in}}%
\pgfpathlineto{\pgfqpoint{1.725701in}{2.476025in}}%
\pgfpathlineto{\pgfqpoint{1.727983in}{2.512151in}}%
\pgfpathlineto{\pgfqpoint{1.732547in}{2.419255in}}%
\pgfpathlineto{\pgfqpoint{1.737112in}{2.486347in}}%
\pgfpathlineto{\pgfqpoint{1.739394in}{2.501507in}}%
\pgfpathlineto{\pgfqpoint{1.741676in}{2.429254in}}%
\pgfpathlineto{\pgfqpoint{1.743959in}{2.449898in}}%
\pgfpathlineto{\pgfqpoint{1.746241in}{2.496346in}}%
\pgfpathlineto{\pgfqpoint{1.748523in}{2.491185in}}%
\pgfpathlineto{\pgfqpoint{1.750805in}{2.455059in}}%
\pgfpathlineto{\pgfqpoint{1.753088in}{2.501507in}}%
\pgfpathlineto{\pgfqpoint{1.755370in}{2.465381in}}%
\pgfpathlineto{\pgfqpoint{1.757652in}{2.398289in}}%
\pgfpathlineto{\pgfqpoint{1.759934in}{2.475703in}}%
\pgfpathlineto{\pgfqpoint{1.762217in}{2.434415in}}%
\pgfpathlineto{\pgfqpoint{1.764499in}{2.480863in}}%
\pgfpathlineto{\pgfqpoint{1.766781in}{2.444737in}}%
\pgfpathlineto{\pgfqpoint{1.769063in}{2.511829in}}%
\pgfpathlineto{\pgfqpoint{1.771346in}{2.439576in}}%
\pgfpathlineto{\pgfqpoint{1.773628in}{2.455059in}}%
\pgfpathlineto{\pgfqpoint{1.775910in}{2.444737in}}%
\pgfpathlineto{\pgfqpoint{1.778192in}{2.501507in}}%
\pgfpathlineto{\pgfqpoint{1.780475in}{2.496346in}}%
\pgfpathlineto{\pgfqpoint{1.782757in}{2.470542in}}%
\pgfpathlineto{\pgfqpoint{1.785039in}{2.465381in}}%
\pgfpathlineto{\pgfqpoint{1.787321in}{2.496346in}}%
\pgfpathlineto{\pgfqpoint{1.789604in}{2.449898in}}%
\pgfpathlineto{\pgfqpoint{1.791886in}{2.501507in}}%
\pgfpathlineto{\pgfqpoint{1.794168in}{2.480863in}}%
\pgfpathlineto{\pgfqpoint{1.798733in}{2.522151in}}%
\pgfpathlineto{\pgfqpoint{1.801015in}{2.511829in}}%
\pgfpathlineto{\pgfqpoint{1.805579in}{2.475703in}}%
\pgfpathlineto{\pgfqpoint{1.807862in}{2.516990in}}%
\pgfpathlineto{\pgfqpoint{1.810144in}{2.470542in}}%
\pgfpathlineto{\pgfqpoint{1.812426in}{2.501507in}}%
\pgfpathlineto{\pgfqpoint{1.814708in}{2.475703in}}%
\pgfpathlineto{\pgfqpoint{1.819273in}{2.501184in}}%
\pgfpathlineto{\pgfqpoint{1.821555in}{2.480541in}}%
\pgfpathlineto{\pgfqpoint{1.826120in}{2.413449in}}%
\pgfpathlineto{\pgfqpoint{1.828402in}{2.475380in}}%
\pgfpathlineto{\pgfqpoint{1.830684in}{2.465058in}}%
\pgfpathlineto{\pgfqpoint{1.832966in}{2.496024in}}%
\pgfpathlineto{\pgfqpoint{1.837531in}{2.428932in}}%
\pgfpathlineto{\pgfqpoint{1.839813in}{2.511506in}}%
\pgfpathlineto{\pgfqpoint{1.842095in}{2.475380in}}%
\pgfpathlineto{\pgfqpoint{1.844378in}{2.537311in}}%
\pgfpathlineto{\pgfqpoint{1.846660in}{2.547633in}}%
\pgfpathlineto{\pgfqpoint{1.848942in}{2.454736in}}%
\pgfpathlineto{\pgfqpoint{1.853507in}{2.526989in}}%
\pgfpathlineto{\pgfqpoint{1.855789in}{2.496024in}}%
\pgfpathlineto{\pgfqpoint{1.858071in}{2.542472in}}%
\pgfpathlineto{\pgfqpoint{1.860353in}{2.454736in}}%
\pgfpathlineto{\pgfqpoint{1.862636in}{2.454736in}}%
\pgfpathlineto{\pgfqpoint{1.864918in}{2.470219in}}%
\pgfpathlineto{\pgfqpoint{1.867200in}{2.516667in}}%
\pgfpathlineto{\pgfqpoint{1.869482in}{2.459897in}}%
\pgfpathlineto{\pgfqpoint{1.871765in}{2.496024in}}%
\pgfpathlineto{\pgfqpoint{1.874047in}{2.496024in}}%
\pgfpathlineto{\pgfqpoint{1.876329in}{2.537311in}}%
\pgfpathlineto{\pgfqpoint{1.878611in}{2.501184in}}%
\pgfpathlineto{\pgfqpoint{1.880894in}{2.506345in}}%
\pgfpathlineto{\pgfqpoint{1.883176in}{2.475380in}}%
\pgfpathlineto{\pgfqpoint{1.885458in}{2.470219in}}%
\pgfpathlineto{\pgfqpoint{1.887740in}{2.542472in}}%
\pgfpathlineto{\pgfqpoint{1.892305in}{2.506345in}}%
\pgfpathlineto{\pgfqpoint{1.894587in}{2.526989in}}%
\pgfpathlineto{\pgfqpoint{1.896869in}{2.496024in}}%
\pgfpathlineto{\pgfqpoint{1.899152in}{2.480218in}}%
\pgfpathlineto{\pgfqpoint{1.901434in}{2.444092in}}%
\pgfpathlineto{\pgfqpoint{1.905998in}{2.536988in}}%
\pgfpathlineto{\pgfqpoint{1.908281in}{2.454414in}}%
\pgfpathlineto{\pgfqpoint{1.910563in}{2.500862in}}%
\pgfpathlineto{\pgfqpoint{1.912845in}{2.521506in}}%
\pgfpathlineto{\pgfqpoint{1.915127in}{2.500862in}}%
\pgfpathlineto{\pgfqpoint{1.917410in}{2.444092in}}%
\pgfpathlineto{\pgfqpoint{1.921974in}{2.475057in}}%
\pgfpathlineto{\pgfqpoint{1.924256in}{2.480218in}}%
\pgfpathlineto{\pgfqpoint{1.926539in}{2.506023in}}%
\pgfpathlineto{\pgfqpoint{1.928821in}{2.562793in}}%
\pgfpathlineto{\pgfqpoint{1.933385in}{2.475057in}}%
\pgfpathlineto{\pgfqpoint{1.935668in}{2.485379in}}%
\pgfpathlineto{\pgfqpoint{1.937950in}{2.459575in}}%
\pgfpathlineto{\pgfqpoint{1.940232in}{2.511184in}}%
\pgfpathlineto{\pgfqpoint{1.942514in}{2.475057in}}%
\pgfpathlineto{\pgfqpoint{1.944797in}{2.485379in}}%
\pgfpathlineto{\pgfqpoint{1.947079in}{2.485379in}}%
\pgfpathlineto{\pgfqpoint{1.949361in}{2.536988in}}%
\pgfpathlineto{\pgfqpoint{1.951643in}{2.480218in}}%
\pgfpathlineto{\pgfqpoint{1.953926in}{2.531827in}}%
\pgfpathlineto{\pgfqpoint{1.956208in}{2.521506in}}%
\pgfpathlineto{\pgfqpoint{1.958490in}{2.506023in}}%
\pgfpathlineto{\pgfqpoint{1.960772in}{2.516345in}}%
\pgfpathlineto{\pgfqpoint{1.963055in}{2.490540in}}%
\pgfpathlineto{\pgfqpoint{1.965337in}{2.516345in}}%
\pgfpathlineto{\pgfqpoint{1.967619in}{2.459575in}}%
\pgfpathlineto{\pgfqpoint{1.969901in}{2.495701in}}%
\pgfpathlineto{\pgfqpoint{1.972184in}{2.438931in}}%
\pgfpathlineto{\pgfqpoint{1.974466in}{2.552471in}}%
\pgfpathlineto{\pgfqpoint{1.976748in}{2.495701in}}%
\pgfpathlineto{\pgfqpoint{1.979030in}{2.500539in}}%
\pgfpathlineto{\pgfqpoint{1.981313in}{2.500539in}}%
\pgfpathlineto{\pgfqpoint{1.983595in}{2.552148in}}%
\pgfpathlineto{\pgfqpoint{1.985877in}{2.521183in}}%
\pgfpathlineto{\pgfqpoint{1.988159in}{2.505700in}}%
\pgfpathlineto{\pgfqpoint{1.990442in}{2.516022in}}%
\pgfpathlineto{\pgfqpoint{1.992724in}{2.546988in}}%
\pgfpathlineto{\pgfqpoint{1.995006in}{2.495378in}}%
\pgfpathlineto{\pgfqpoint{1.997289in}{2.490218in}}%
\pgfpathlineto{\pgfqpoint{2.001853in}{2.412804in}}%
\pgfpathlineto{\pgfqpoint{2.004135in}{2.521183in}}%
\pgfpathlineto{\pgfqpoint{2.006418in}{2.464413in}}%
\pgfpathlineto{\pgfqpoint{2.008700in}{2.521183in}}%
\pgfpathlineto{\pgfqpoint{2.010982in}{2.505700in}}%
\pgfpathlineto{\pgfqpoint{2.013264in}{2.510861in}}%
\pgfpathlineto{\pgfqpoint{2.015547in}{2.500539in}}%
\pgfpathlineto{\pgfqpoint{2.017829in}{2.510861in}}%
\pgfpathlineto{\pgfqpoint{2.020111in}{2.448930in}}%
\pgfpathlineto{\pgfqpoint{2.022393in}{2.459252in}}%
\pgfpathlineto{\pgfqpoint{2.024676in}{2.526344in}}%
\pgfpathlineto{\pgfqpoint{2.029240in}{2.485057in}}%
\pgfpathlineto{\pgfqpoint{2.031522in}{2.474735in}}%
\pgfpathlineto{\pgfqpoint{2.033805in}{2.541827in}}%
\pgfpathlineto{\pgfqpoint{2.036087in}{2.505700in}}%
\pgfpathlineto{\pgfqpoint{2.038369in}{2.516022in}}%
\pgfpathlineto{\pgfqpoint{2.040651in}{2.516022in}}%
\pgfpathlineto{\pgfqpoint{2.042934in}{2.490218in}}%
\pgfpathlineto{\pgfqpoint{2.045216in}{2.479896in}}%
\pgfpathlineto{\pgfqpoint{2.047498in}{2.521183in}}%
\pgfpathlineto{\pgfqpoint{2.049780in}{2.490218in}}%
\pgfpathlineto{\pgfqpoint{2.052063in}{2.526344in}}%
\pgfpathlineto{\pgfqpoint{2.054345in}{2.505700in}}%
\pgfpathlineto{\pgfqpoint{2.056627in}{2.526344in}}%
\pgfpathlineto{\pgfqpoint{2.058909in}{2.484734in}}%
\pgfpathlineto{\pgfqpoint{2.061192in}{2.541504in}}%
\pgfpathlineto{\pgfqpoint{2.063474in}{2.526021in}}%
\pgfpathlineto{\pgfqpoint{2.065756in}{2.551826in}}%
\pgfpathlineto{\pgfqpoint{2.068038in}{2.520860in}}%
\pgfpathlineto{\pgfqpoint{2.070321in}{2.526021in}}%
\pgfpathlineto{\pgfqpoint{2.072603in}{2.474412in}}%
\pgfpathlineto{\pgfqpoint{2.074885in}{2.474412in}}%
\pgfpathlineto{\pgfqpoint{2.077167in}{2.536343in}}%
\pgfpathlineto{\pgfqpoint{2.079450in}{2.500217in}}%
\pgfpathlineto{\pgfqpoint{2.081732in}{2.495056in}}%
\pgfpathlineto{\pgfqpoint{2.084014in}{2.500217in}}%
\pgfpathlineto{\pgfqpoint{2.086296in}{2.531182in}}%
\pgfpathlineto{\pgfqpoint{2.088579in}{2.541504in}}%
\pgfpathlineto{\pgfqpoint{2.090861in}{2.427964in}}%
\pgfpathlineto{\pgfqpoint{2.093143in}{2.531182in}}%
\pgfpathlineto{\pgfqpoint{2.097708in}{2.500217in}}%
\pgfpathlineto{\pgfqpoint{2.099990in}{2.515700in}}%
\pgfpathlineto{\pgfqpoint{2.102272in}{2.520860in}}%
\pgfpathlineto{\pgfqpoint{2.104554in}{2.520860in}}%
\pgfpathlineto{\pgfqpoint{2.106837in}{2.510539in}}%
\pgfpathlineto{\pgfqpoint{2.109119in}{2.469251in}}%
\pgfpathlineto{\pgfqpoint{2.111401in}{2.556987in}}%
\pgfpathlineto{\pgfqpoint{2.113683in}{2.505378in}}%
\pgfpathlineto{\pgfqpoint{2.115966in}{2.515700in}}%
\pgfpathlineto{\pgfqpoint{2.118248in}{2.515700in}}%
\pgfpathlineto{\pgfqpoint{2.120530in}{2.551826in}}%
\pgfpathlineto{\pgfqpoint{2.122812in}{2.505378in}}%
\pgfpathlineto{\pgfqpoint{2.125095in}{2.500217in}}%
\pgfpathlineto{\pgfqpoint{2.127377in}{2.567309in}}%
\pgfpathlineto{\pgfqpoint{2.129659in}{2.556987in}}%
\pgfpathlineto{\pgfqpoint{2.131941in}{2.526021in}}%
\pgfpathlineto{\pgfqpoint{2.134224in}{2.526021in}}%
\pgfpathlineto{\pgfqpoint{2.136506in}{2.510539in}}%
\pgfpathlineto{\pgfqpoint{2.138788in}{2.525699in}}%
\pgfpathlineto{\pgfqpoint{2.141070in}{2.474090in}}%
\pgfpathlineto{\pgfqpoint{2.145635in}{2.515377in}}%
\pgfpathlineto{\pgfqpoint{2.147917in}{2.530860in}}%
\pgfpathlineto{\pgfqpoint{2.150199in}{2.505055in}}%
\pgfpathlineto{\pgfqpoint{2.152482in}{2.530860in}}%
\pgfpathlineto{\pgfqpoint{2.154764in}{2.525699in}}%
\pgfpathlineto{\pgfqpoint{2.157046in}{2.546342in}}%
\pgfpathlineto{\pgfqpoint{2.159328in}{2.510216in}}%
\pgfpathlineto{\pgfqpoint{2.161611in}{2.525699in}}%
\pgfpathlineto{\pgfqpoint{2.163893in}{2.479251in}}%
\pgfpathlineto{\pgfqpoint{2.166175in}{2.499894in}}%
\pgfpathlineto{\pgfqpoint{2.168457in}{2.541181in}}%
\pgfpathlineto{\pgfqpoint{2.173022in}{2.458607in}}%
\pgfpathlineto{\pgfqpoint{2.175304in}{2.479251in}}%
\pgfpathlineto{\pgfqpoint{2.177586in}{2.546342in}}%
\pgfpathlineto{\pgfqpoint{2.179869in}{2.546342in}}%
\pgfpathlineto{\pgfqpoint{2.184433in}{2.484412in}}%
\pgfpathlineto{\pgfqpoint{2.186715in}{2.458607in}}%
\pgfpathlineto{\pgfqpoint{2.188998in}{2.541181in}}%
\pgfpathlineto{\pgfqpoint{2.193562in}{2.489572in}}%
\pgfpathlineto{\pgfqpoint{2.195844in}{2.530860in}}%
\pgfpathlineto{\pgfqpoint{2.198127in}{2.530860in}}%
\pgfpathlineto{\pgfqpoint{2.200409in}{2.505055in}}%
\pgfpathlineto{\pgfqpoint{2.202691in}{2.515377in}}%
\pgfpathlineto{\pgfqpoint{2.204973in}{2.530860in}}%
\pgfpathlineto{\pgfqpoint{2.207256in}{2.520538in}}%
\pgfpathlineto{\pgfqpoint{2.209538in}{2.494733in}}%
\pgfpathlineto{\pgfqpoint{2.211820in}{2.489572in}}%
\pgfpathlineto{\pgfqpoint{2.214102in}{2.520538in}}%
\pgfpathlineto{\pgfqpoint{2.216385in}{2.515377in}}%
\pgfpathlineto{\pgfqpoint{2.218667in}{2.499572in}}%
\pgfpathlineto{\pgfqpoint{2.220949in}{2.473767in}}%
\pgfpathlineto{\pgfqpoint{2.223231in}{2.546020in}}%
\pgfpathlineto{\pgfqpoint{2.225514in}{2.499572in}}%
\pgfpathlineto{\pgfqpoint{2.227796in}{2.509893in}}%
\pgfpathlineto{\pgfqpoint{2.230078in}{2.484089in}}%
\pgfpathlineto{\pgfqpoint{2.232360in}{2.515054in}}%
\pgfpathlineto{\pgfqpoint{2.234643in}{2.468606in}}%
\pgfpathlineto{\pgfqpoint{2.236925in}{2.489250in}}%
\pgfpathlineto{\pgfqpoint{2.239207in}{2.535698in}}%
\pgfpathlineto{\pgfqpoint{2.241489in}{2.478928in}}%
\pgfpathlineto{\pgfqpoint{2.243772in}{2.551181in}}%
\pgfpathlineto{\pgfqpoint{2.246054in}{2.499572in}}%
\pgfpathlineto{\pgfqpoint{2.248336in}{2.509893in}}%
\pgfpathlineto{\pgfqpoint{2.250618in}{2.453124in}}%
\pgfpathlineto{\pgfqpoint{2.257465in}{2.499572in}}%
\pgfpathlineto{\pgfqpoint{2.259747in}{2.530537in}}%
\pgfpathlineto{\pgfqpoint{2.262030in}{2.525376in}}%
\pgfpathlineto{\pgfqpoint{2.264312in}{2.504733in}}%
\pgfpathlineto{\pgfqpoint{2.266594in}{2.509893in}}%
\pgfpathlineto{\pgfqpoint{2.268876in}{2.509893in}}%
\pgfpathlineto{\pgfqpoint{2.271159in}{2.453124in}}%
\pgfpathlineto{\pgfqpoint{2.273441in}{2.463445in}}%
\pgfpathlineto{\pgfqpoint{2.275723in}{2.530537in}}%
\pgfpathlineto{\pgfqpoint{2.278005in}{2.473767in}}%
\pgfpathlineto{\pgfqpoint{2.280288in}{2.468606in}}%
\pgfpathlineto{\pgfqpoint{2.282570in}{2.535698in}}%
\pgfpathlineto{\pgfqpoint{2.287134in}{2.499572in}}%
\pgfpathlineto{\pgfqpoint{2.289417in}{2.504733in}}%
\pgfpathlineto{\pgfqpoint{2.291699in}{2.499572in}}%
\pgfpathlineto{\pgfqpoint{2.293981in}{2.530537in}}%
\pgfpathlineto{\pgfqpoint{2.296264in}{2.484089in}}%
\pgfpathlineto{\pgfqpoint{2.298546in}{2.488927in}}%
\pgfpathlineto{\pgfqpoint{2.300828in}{2.483766in}}%
\pgfpathlineto{\pgfqpoint{2.303110in}{2.556019in}}%
\pgfpathlineto{\pgfqpoint{2.305393in}{2.525054in}}%
\pgfpathlineto{\pgfqpoint{2.307675in}{2.437318in}}%
\pgfpathlineto{\pgfqpoint{2.309957in}{2.483766in}}%
\pgfpathlineto{\pgfqpoint{2.312239in}{2.483766in}}%
\pgfpathlineto{\pgfqpoint{2.314522in}{2.468284in}}%
\pgfpathlineto{\pgfqpoint{2.316804in}{2.494088in}}%
\pgfpathlineto{\pgfqpoint{2.319086in}{2.488927in}}%
\pgfpathlineto{\pgfqpoint{2.321368in}{2.499249in}}%
\pgfpathlineto{\pgfqpoint{2.323651in}{2.504410in}}%
\pgfpathlineto{\pgfqpoint{2.325933in}{2.525054in}}%
\pgfpathlineto{\pgfqpoint{2.328215in}{2.509571in}}%
\pgfpathlineto{\pgfqpoint{2.330497in}{2.504410in}}%
\pgfpathlineto{\pgfqpoint{2.332780in}{2.494088in}}%
\pgfpathlineto{\pgfqpoint{2.337344in}{2.483766in}}%
\pgfpathlineto{\pgfqpoint{2.339626in}{2.463123in}}%
\pgfpathlineto{\pgfqpoint{2.344191in}{2.525054in}}%
\pgfpathlineto{\pgfqpoint{2.346473in}{2.499249in}}%
\pgfpathlineto{\pgfqpoint{2.348755in}{2.499249in}}%
\pgfpathlineto{\pgfqpoint{2.351038in}{2.504410in}}%
\pgfpathlineto{\pgfqpoint{2.353320in}{2.447640in}}%
\pgfpathlineto{\pgfqpoint{2.355602in}{2.452801in}}%
\pgfpathlineto{\pgfqpoint{2.357884in}{2.535375in}}%
\pgfpathlineto{\pgfqpoint{2.360167in}{2.519893in}}%
\pgfpathlineto{\pgfqpoint{2.362449in}{2.519893in}}%
\pgfpathlineto{\pgfqpoint{2.364731in}{2.504410in}}%
\pgfpathlineto{\pgfqpoint{2.367013in}{2.514732in}}%
\pgfpathlineto{\pgfqpoint{2.369296in}{2.535375in}}%
\pgfpathlineto{\pgfqpoint{2.371578in}{2.499249in}}%
\pgfpathlineto{\pgfqpoint{2.373860in}{2.519893in}}%
\pgfpathlineto{\pgfqpoint{2.376142in}{2.452801in}}%
\pgfpathlineto{\pgfqpoint{2.378425in}{2.462800in}}%
\pgfpathlineto{\pgfqpoint{2.380707in}{2.483444in}}%
\pgfpathlineto{\pgfqpoint{2.382989in}{2.519570in}}%
\pgfpathlineto{\pgfqpoint{2.385271in}{2.493766in}}%
\pgfpathlineto{\pgfqpoint{2.387554in}{2.529892in}}%
\pgfpathlineto{\pgfqpoint{2.389836in}{2.462800in}}%
\pgfpathlineto{\pgfqpoint{2.392118in}{2.488605in}}%
\pgfpathlineto{\pgfqpoint{2.394400in}{2.380226in}}%
\pgfpathlineto{\pgfqpoint{2.396683in}{2.467961in}}%
\pgfpathlineto{\pgfqpoint{2.398965in}{2.431835in}}%
\pgfpathlineto{\pgfqpoint{2.401247in}{2.442157in}}%
\pgfpathlineto{\pgfqpoint{2.403529in}{2.462800in}}%
\pgfpathlineto{\pgfqpoint{2.405812in}{2.514409in}}%
\pgfpathlineto{\pgfqpoint{2.408094in}{2.519570in}}%
\pgfpathlineto{\pgfqpoint{2.410376in}{2.483444in}}%
\pgfpathlineto{\pgfqpoint{2.412658in}{2.488605in}}%
\pgfpathlineto{\pgfqpoint{2.414941in}{2.509248in}}%
\pgfpathlineto{\pgfqpoint{2.417223in}{2.488605in}}%
\pgfpathlineto{\pgfqpoint{2.419505in}{2.442157in}}%
\pgfpathlineto{\pgfqpoint{2.421787in}{2.524731in}}%
\pgfpathlineto{\pgfqpoint{2.424070in}{2.509248in}}%
\pgfpathlineto{\pgfqpoint{2.426352in}{2.535053in}}%
\pgfpathlineto{\pgfqpoint{2.433199in}{2.426674in}}%
\pgfpathlineto{\pgfqpoint{2.435481in}{2.498927in}}%
\pgfpathlineto{\pgfqpoint{2.437763in}{2.436996in}}%
\pgfpathlineto{\pgfqpoint{2.442328in}{2.535053in}}%
\pgfpathlineto{\pgfqpoint{2.444610in}{2.498927in}}%
\pgfpathlineto{\pgfqpoint{2.446892in}{2.498927in}}%
\pgfpathlineto{\pgfqpoint{2.449174in}{2.483444in}}%
\pgfpathlineto{\pgfqpoint{2.451457in}{2.478283in}}%
\pgfpathlineto{\pgfqpoint{2.453739in}{2.436996in}}%
\pgfpathlineto{\pgfqpoint{2.456021in}{2.416352in}}%
\pgfpathlineto{\pgfqpoint{2.458303in}{2.493443in}}%
\pgfpathlineto{\pgfqpoint{2.460586in}{2.503765in}}%
\pgfpathlineto{\pgfqpoint{2.462868in}{2.441834in}}%
\pgfpathlineto{\pgfqpoint{2.465150in}{2.493443in}}%
\pgfpathlineto{\pgfqpoint{2.467432in}{2.514087in}}%
\pgfpathlineto{\pgfqpoint{2.469715in}{2.514087in}}%
\pgfpathlineto{\pgfqpoint{2.471997in}{2.508926in}}%
\pgfpathlineto{\pgfqpoint{2.474279in}{2.436673in}}%
\pgfpathlineto{\pgfqpoint{2.476561in}{2.488282in}}%
\pgfpathlineto{\pgfqpoint{2.478844in}{2.493443in}}%
\pgfpathlineto{\pgfqpoint{2.481126in}{2.416030in}}%
\pgfpathlineto{\pgfqpoint{2.485690in}{2.534730in}}%
\pgfpathlineto{\pgfqpoint{2.492537in}{2.426351in}}%
\pgfpathlineto{\pgfqpoint{2.494819in}{2.519248in}}%
\pgfpathlineto{\pgfqpoint{2.497102in}{2.508926in}}%
\pgfpathlineto{\pgfqpoint{2.499384in}{2.472799in}}%
\pgfpathlineto{\pgfqpoint{2.501666in}{2.508926in}}%
\pgfpathlineto{\pgfqpoint{2.503948in}{2.477960in}}%
\pgfpathlineto{\pgfqpoint{2.508513in}{2.498604in}}%
\pgfpathlineto{\pgfqpoint{2.510795in}{2.467639in}}%
\pgfpathlineto{\pgfqpoint{2.513077in}{2.472799in}}%
\pgfpathlineto{\pgfqpoint{2.515360in}{2.498604in}}%
\pgfpathlineto{\pgfqpoint{2.517642in}{2.508926in}}%
\pgfpathlineto{\pgfqpoint{2.519924in}{2.436673in}}%
\pgfpathlineto{\pgfqpoint{2.522206in}{2.529569in}}%
\pgfpathlineto{\pgfqpoint{2.524489in}{2.472799in}}%
\pgfpathlineto{\pgfqpoint{2.526771in}{2.457317in}}%
\pgfpathlineto{\pgfqpoint{2.529053in}{2.508926in}}%
\pgfpathlineto{\pgfqpoint{2.531335in}{2.467639in}}%
\pgfpathlineto{\pgfqpoint{2.533618in}{2.483121in}}%
\pgfpathlineto{\pgfqpoint{2.535900in}{2.472799in}}%
\pgfpathlineto{\pgfqpoint{2.538182in}{2.441512in}}%
\pgfpathlineto{\pgfqpoint{2.542747in}{2.482799in}}%
\pgfpathlineto{\pgfqpoint{2.545029in}{2.472477in}}%
\pgfpathlineto{\pgfqpoint{2.547311in}{2.493121in}}%
\pgfpathlineto{\pgfqpoint{2.549593in}{2.503442in}}%
\pgfpathlineto{\pgfqpoint{2.551876in}{2.493121in}}%
\pgfpathlineto{\pgfqpoint{2.554158in}{2.493121in}}%
\pgfpathlineto{\pgfqpoint{2.556440in}{2.456994in}}%
\pgfpathlineto{\pgfqpoint{2.558722in}{2.498281in}}%
\pgfpathlineto{\pgfqpoint{2.561005in}{2.467316in}}%
\pgfpathlineto{\pgfqpoint{2.563287in}{2.503442in}}%
\pgfpathlineto{\pgfqpoint{2.565569in}{2.487960in}}%
\pgfpathlineto{\pgfqpoint{2.567851in}{2.524086in}}%
\pgfpathlineto{\pgfqpoint{2.570134in}{2.513764in}}%
\pgfpathlineto{\pgfqpoint{2.574698in}{2.467316in}}%
\pgfpathlineto{\pgfqpoint{2.576980in}{2.462155in}}%
\pgfpathlineto{\pgfqpoint{2.579263in}{2.477638in}}%
\pgfpathlineto{\pgfqpoint{2.581545in}{2.544730in}}%
\pgfpathlineto{\pgfqpoint{2.583827in}{2.482799in}}%
\pgfpathlineto{\pgfqpoint{2.586110in}{2.518925in}}%
\pgfpathlineto{\pgfqpoint{2.588392in}{2.518925in}}%
\pgfpathlineto{\pgfqpoint{2.590674in}{2.534408in}}%
\pgfpathlineto{\pgfqpoint{2.592956in}{2.374420in}}%
\pgfpathlineto{\pgfqpoint{2.595239in}{2.467316in}}%
\pgfpathlineto{\pgfqpoint{2.597521in}{2.472477in}}%
\pgfpathlineto{\pgfqpoint{2.599803in}{2.446672in}}%
\pgfpathlineto{\pgfqpoint{2.602085in}{2.462155in}}%
\pgfpathlineto{\pgfqpoint{2.604368in}{2.498281in}}%
\pgfpathlineto{\pgfqpoint{2.606650in}{2.446672in}}%
\pgfpathlineto{\pgfqpoint{2.608932in}{2.467316in}}%
\pgfpathlineto{\pgfqpoint{2.611214in}{2.446672in}}%
\pgfpathlineto{\pgfqpoint{2.615779in}{2.508603in}}%
\pgfpathlineto{\pgfqpoint{2.620343in}{2.394741in}}%
\pgfpathlineto{\pgfqpoint{2.622626in}{2.492798in}}%
\pgfpathlineto{\pgfqpoint{2.624908in}{2.523763in}}%
\pgfpathlineto{\pgfqpoint{2.627190in}{2.508281in}}%
\pgfpathlineto{\pgfqpoint{2.629472in}{2.446350in}}%
\pgfpathlineto{\pgfqpoint{2.631755in}{2.513442in}}%
\pgfpathlineto{\pgfqpoint{2.634037in}{2.456672in}}%
\pgfpathlineto{\pgfqpoint{2.636319in}{2.492798in}}%
\pgfpathlineto{\pgfqpoint{2.638601in}{2.451511in}}%
\pgfpathlineto{\pgfqpoint{2.640884in}{2.482476in}}%
\pgfpathlineto{\pgfqpoint{2.643166in}{2.415384in}}%
\pgfpathlineto{\pgfqpoint{2.647730in}{2.497959in}}%
\pgfpathlineto{\pgfqpoint{2.650013in}{2.451511in}}%
\pgfpathlineto{\pgfqpoint{2.654577in}{2.472154in}}%
\pgfpathlineto{\pgfqpoint{2.656859in}{2.456672in}}%
\pgfpathlineto{\pgfqpoint{2.659142in}{2.399902in}}%
\pgfpathlineto{\pgfqpoint{2.661424in}{2.492798in}}%
\pgfpathlineto{\pgfqpoint{2.663706in}{2.446350in}}%
\pgfpathlineto{\pgfqpoint{2.665988in}{2.497959in}}%
\pgfpathlineto{\pgfqpoint{2.668271in}{2.420545in}}%
\pgfpathlineto{\pgfqpoint{2.670553in}{2.425706in}}%
\pgfpathlineto{\pgfqpoint{2.672835in}{2.472154in}}%
\pgfpathlineto{\pgfqpoint{2.675117in}{2.487637in}}%
\pgfpathlineto{\pgfqpoint{2.677400in}{2.513442in}}%
\pgfpathlineto{\pgfqpoint{2.679682in}{2.508281in}}%
\pgfpathlineto{\pgfqpoint{2.681964in}{2.466993in}}%
\pgfpathlineto{\pgfqpoint{2.684246in}{2.461833in}}%
\pgfpathlineto{\pgfqpoint{2.686529in}{2.451511in}}%
\pgfpathlineto{\pgfqpoint{2.688811in}{2.430867in}}%
\pgfpathlineto{\pgfqpoint{2.691093in}{2.492798in}}%
\pgfpathlineto{\pgfqpoint{2.693375in}{2.492798in}}%
\pgfpathlineto{\pgfqpoint{2.695658in}{2.518603in}}%
\pgfpathlineto{\pgfqpoint{2.697940in}{2.446027in}}%
\pgfpathlineto{\pgfqpoint{2.700222in}{2.456349in}}%
\pgfpathlineto{\pgfqpoint{2.702504in}{2.461510in}}%
\pgfpathlineto{\pgfqpoint{2.704787in}{2.476993in}}%
\pgfpathlineto{\pgfqpoint{2.709351in}{2.456349in}}%
\pgfpathlineto{\pgfqpoint{2.711633in}{2.507958in}}%
\pgfpathlineto{\pgfqpoint{2.713916in}{2.476993in}}%
\pgfpathlineto{\pgfqpoint{2.716198in}{2.502797in}}%
\pgfpathlineto{\pgfqpoint{2.718480in}{2.461510in}}%
\pgfpathlineto{\pgfqpoint{2.720762in}{2.471832in}}%
\pgfpathlineto{\pgfqpoint{2.723045in}{2.466671in}}%
\pgfpathlineto{\pgfqpoint{2.727609in}{2.440866in}}%
\pgfpathlineto{\pgfqpoint{2.729891in}{2.440866in}}%
\pgfpathlineto{\pgfqpoint{2.732174in}{2.466671in}}%
\pgfpathlineto{\pgfqpoint{2.734456in}{2.466671in}}%
\pgfpathlineto{\pgfqpoint{2.736738in}{2.435705in}}%
\pgfpathlineto{\pgfqpoint{2.739020in}{2.482154in}}%
\pgfpathlineto{\pgfqpoint{2.743585in}{2.461510in}}%
\pgfpathlineto{\pgfqpoint{2.745867in}{2.456349in}}%
\pgfpathlineto{\pgfqpoint{2.748149in}{2.461510in}}%
\pgfpathlineto{\pgfqpoint{2.750432in}{2.461510in}}%
\pgfpathlineto{\pgfqpoint{2.752714in}{2.466671in}}%
\pgfpathlineto{\pgfqpoint{2.754996in}{2.492475in}}%
\pgfpathlineto{\pgfqpoint{2.757278in}{2.435705in}}%
\pgfpathlineto{\pgfqpoint{2.759561in}{2.476993in}}%
\pgfpathlineto{\pgfqpoint{2.761843in}{2.425384in}}%
\pgfpathlineto{\pgfqpoint{2.764125in}{2.466671in}}%
\pgfpathlineto{\pgfqpoint{2.766407in}{2.476993in}}%
\pgfpathlineto{\pgfqpoint{2.768690in}{2.482154in}}%
\pgfpathlineto{\pgfqpoint{2.770972in}{2.399579in}}%
\pgfpathlineto{\pgfqpoint{2.773254in}{2.482154in}}%
\pgfpathlineto{\pgfqpoint{2.775536in}{2.451188in}}%
\pgfpathlineto{\pgfqpoint{2.777819in}{2.435383in}}%
\pgfpathlineto{\pgfqpoint{2.782383in}{2.523118in}}%
\pgfpathlineto{\pgfqpoint{2.784665in}{2.450866in}}%
\pgfpathlineto{\pgfqpoint{2.789230in}{2.404418in}}%
\pgfpathlineto{\pgfqpoint{2.791512in}{2.507636in}}%
\pgfpathlineto{\pgfqpoint{2.793794in}{2.456027in}}%
\pgfpathlineto{\pgfqpoint{2.796077in}{2.466348in}}%
\pgfpathlineto{\pgfqpoint{2.798359in}{2.414739in}}%
\pgfpathlineto{\pgfqpoint{2.800641in}{2.492153in}}%
\pgfpathlineto{\pgfqpoint{2.802923in}{2.502475in}}%
\pgfpathlineto{\pgfqpoint{2.805206in}{2.502475in}}%
\pgfpathlineto{\pgfqpoint{2.807488in}{2.435383in}}%
\pgfpathlineto{\pgfqpoint{2.809770in}{2.425061in}}%
\pgfpathlineto{\pgfqpoint{2.812052in}{2.419900in}}%
\pgfpathlineto{\pgfqpoint{2.814335in}{2.383774in}}%
\pgfpathlineto{\pgfqpoint{2.816617in}{2.430222in}}%
\pgfpathlineto{\pgfqpoint{2.818899in}{2.409578in}}%
\pgfpathlineto{\pgfqpoint{2.821181in}{2.502475in}}%
\pgfpathlineto{\pgfqpoint{2.823464in}{2.492153in}}%
\pgfpathlineto{\pgfqpoint{2.825746in}{2.466348in}}%
\pgfpathlineto{\pgfqpoint{2.828028in}{2.466348in}}%
\pgfpathlineto{\pgfqpoint{2.830310in}{2.456027in}}%
\pgfpathlineto{\pgfqpoint{2.832593in}{2.471509in}}%
\pgfpathlineto{\pgfqpoint{2.834875in}{2.409578in}}%
\pgfpathlineto{\pgfqpoint{2.837157in}{2.440544in}}%
\pgfpathlineto{\pgfqpoint{2.839439in}{2.440544in}}%
\pgfpathlineto{\pgfqpoint{2.841722in}{2.497314in}}%
\pgfpathlineto{\pgfqpoint{2.844004in}{2.502475in}}%
\pgfpathlineto{\pgfqpoint{2.848568in}{2.456027in}}%
\pgfpathlineto{\pgfqpoint{2.850851in}{2.440544in}}%
\pgfpathlineto{\pgfqpoint{2.853133in}{2.486992in}}%
\pgfpathlineto{\pgfqpoint{2.855415in}{2.450866in}}%
\pgfpathlineto{\pgfqpoint{2.857697in}{2.460865in}}%
\pgfpathlineto{\pgfqpoint{2.859980in}{2.445382in}}%
\pgfpathlineto{\pgfqpoint{2.862262in}{2.502152in}}%
\pgfpathlineto{\pgfqpoint{2.864544in}{2.491830in}}%
\pgfpathlineto{\pgfqpoint{2.866826in}{2.486669in}}%
\pgfpathlineto{\pgfqpoint{2.869109in}{2.435060in}}%
\pgfpathlineto{\pgfqpoint{2.871391in}{2.491830in}}%
\pgfpathlineto{\pgfqpoint{2.873673in}{2.476348in}}%
\pgfpathlineto{\pgfqpoint{2.875955in}{2.481509in}}%
\pgfpathlineto{\pgfqpoint{2.878238in}{2.450543in}}%
\pgfpathlineto{\pgfqpoint{2.880520in}{2.445382in}}%
\pgfpathlineto{\pgfqpoint{2.885085in}{2.486669in}}%
\pgfpathlineto{\pgfqpoint{2.887367in}{2.476348in}}%
\pgfpathlineto{\pgfqpoint{2.889649in}{2.435060in}}%
\pgfpathlineto{\pgfqpoint{2.894214in}{2.398934in}}%
\pgfpathlineto{\pgfqpoint{2.896496in}{2.466026in}}%
\pgfpathlineto{\pgfqpoint{2.898778in}{2.398934in}}%
\pgfpathlineto{\pgfqpoint{2.901060in}{2.471187in}}%
\pgfpathlineto{\pgfqpoint{2.903343in}{2.455704in}}%
\pgfpathlineto{\pgfqpoint{2.905625in}{2.476348in}}%
\pgfpathlineto{\pgfqpoint{2.907907in}{2.476348in}}%
\pgfpathlineto{\pgfqpoint{2.910189in}{2.440221in}}%
\pgfpathlineto{\pgfqpoint{2.912472in}{2.486669in}}%
\pgfpathlineto{\pgfqpoint{2.914754in}{2.486669in}}%
\pgfpathlineto{\pgfqpoint{2.917036in}{2.496991in}}%
\pgfpathlineto{\pgfqpoint{2.919318in}{2.378290in}}%
\pgfpathlineto{\pgfqpoint{2.921601in}{2.543439in}}%
\pgfpathlineto{\pgfqpoint{2.923883in}{2.476348in}}%
\pgfpathlineto{\pgfqpoint{2.926165in}{2.450543in}}%
\pgfpathlineto{\pgfqpoint{2.928447in}{2.409256in}}%
\pgfpathlineto{\pgfqpoint{2.930730in}{2.460865in}}%
\pgfpathlineto{\pgfqpoint{2.933012in}{2.466026in}}%
\pgfpathlineto{\pgfqpoint{2.935294in}{2.460865in}}%
\pgfpathlineto{\pgfqpoint{2.937576in}{2.491508in}}%
\pgfpathlineto{\pgfqpoint{2.939859in}{2.465703in}}%
\pgfpathlineto{\pgfqpoint{2.942141in}{2.476025in}}%
\pgfpathlineto{\pgfqpoint{2.944423in}{2.481186in}}%
\pgfpathlineto{\pgfqpoint{2.946705in}{2.465703in}}%
\pgfpathlineto{\pgfqpoint{2.948988in}{2.419255in}}%
\pgfpathlineto{\pgfqpoint{2.951270in}{2.476025in}}%
\pgfpathlineto{\pgfqpoint{2.953552in}{2.470864in}}%
\pgfpathlineto{\pgfqpoint{2.955834in}{2.455381in}}%
\pgfpathlineto{\pgfqpoint{2.958117in}{2.450221in}}%
\pgfpathlineto{\pgfqpoint{2.962681in}{2.537956in}}%
\pgfpathlineto{\pgfqpoint{2.964963in}{2.512151in}}%
\pgfpathlineto{\pgfqpoint{2.967246in}{2.522473in}}%
\pgfpathlineto{\pgfqpoint{2.971810in}{2.455381in}}%
\pgfpathlineto{\pgfqpoint{2.976375in}{2.517312in}}%
\pgfpathlineto{\pgfqpoint{2.978657in}{2.486347in}}%
\pgfpathlineto{\pgfqpoint{2.980939in}{2.491508in}}%
\pgfpathlineto{\pgfqpoint{2.983221in}{2.470864in}}%
\pgfpathlineto{\pgfqpoint{2.987786in}{2.543117in}}%
\pgfpathlineto{\pgfqpoint{2.990068in}{2.532795in}}%
\pgfpathlineto{\pgfqpoint{2.992350in}{2.501830in}}%
\pgfpathlineto{\pgfqpoint{2.996915in}{2.470864in}}%
\pgfpathlineto{\pgfqpoint{2.999197in}{2.439899in}}%
\pgfpathlineto{\pgfqpoint{3.001479in}{2.476025in}}%
\pgfpathlineto{\pgfqpoint{3.003762in}{2.439899in}}%
\pgfpathlineto{\pgfqpoint{3.008326in}{2.450221in}}%
\pgfpathlineto{\pgfqpoint{3.010608in}{2.419255in}}%
\pgfpathlineto{\pgfqpoint{3.012891in}{2.476025in}}%
\pgfpathlineto{\pgfqpoint{3.015173in}{2.465703in}}%
\pgfpathlineto{\pgfqpoint{3.017455in}{2.470542in}}%
\pgfpathlineto{\pgfqpoint{3.019737in}{2.496346in}}%
\pgfpathlineto{\pgfqpoint{3.022020in}{2.470542in}}%
\pgfpathlineto{\pgfqpoint{3.024302in}{2.480863in}}%
\pgfpathlineto{\pgfqpoint{3.026584in}{2.465381in}}%
\pgfpathlineto{\pgfqpoint{3.028866in}{2.532472in}}%
\pgfpathlineto{\pgfqpoint{3.031149in}{2.449898in}}%
\pgfpathlineto{\pgfqpoint{3.033431in}{2.408611in}}%
\pgfpathlineto{\pgfqpoint{3.035713in}{2.532472in}}%
\pgfpathlineto{\pgfqpoint{3.037995in}{2.501507in}}%
\pgfpathlineto{\pgfqpoint{3.040278in}{2.527312in}}%
\pgfpathlineto{\pgfqpoint{3.040278in}{2.527312in}}%
\pgfusepath{stroke}%
\end{pgfscope}%
\begin{pgfscope}%
\pgfsetrectcap%
\pgfsetmiterjoin%
\pgfsetlinewidth{0.803000pt}%
\definecolor{currentstroke}{rgb}{0.000000,0.000000,0.000000}%
\pgfsetstrokecolor{currentstroke}%
\pgfsetdash{}{0pt}%
\pgfpathmoveto{\pgfqpoint{0.758026in}{2.240123in}}%
\pgfpathlineto{\pgfqpoint{0.758026in}{3.150926in}}%
\pgfusepath{stroke}%
\end{pgfscope}%
\begin{pgfscope}%
\pgfsetrectcap%
\pgfsetmiterjoin%
\pgfsetlinewidth{0.803000pt}%
\definecolor{currentstroke}{rgb}{0.000000,0.000000,0.000000}%
\pgfsetstrokecolor{currentstroke}%
\pgfsetdash{}{0pt}%
\pgfpathmoveto{\pgfqpoint{3.040278in}{2.240123in}}%
\pgfpathlineto{\pgfqpoint{3.040278in}{3.150926in}}%
\pgfusepath{stroke}%
\end{pgfscope}%
\begin{pgfscope}%
\pgfsetrectcap%
\pgfsetmiterjoin%
\pgfsetlinewidth{0.803000pt}%
\definecolor{currentstroke}{rgb}{0.000000,0.000000,0.000000}%
\pgfsetstrokecolor{currentstroke}%
\pgfsetdash{}{0pt}%
\pgfpathmoveto{\pgfqpoint{0.758026in}{2.240123in}}%
\pgfpathlineto{\pgfqpoint{3.040278in}{2.240123in}}%
\pgfusepath{stroke}%
\end{pgfscope}%
\begin{pgfscope}%
\pgfsetrectcap%
\pgfsetmiterjoin%
\pgfsetlinewidth{0.803000pt}%
\definecolor{currentstroke}{rgb}{0.000000,0.000000,0.000000}%
\pgfsetstrokecolor{currentstroke}%
\pgfsetdash{}{0pt}%
\pgfpathmoveto{\pgfqpoint{0.758026in}{3.150926in}}%
\pgfpathlineto{\pgfqpoint{3.040278in}{3.150926in}}%
\pgfusepath{stroke}%
\end{pgfscope}%
\begin{pgfscope}%
\definecolor{textcolor}{rgb}{0.000000,0.000000,0.000000}%
\pgfsetstrokecolor{textcolor}%
\pgfsetfillcolor{textcolor}%
\pgftext[x=1.899152in,y=3.234260in,,base]{\color{textcolor}\rmfamily\fontsize{12.000000}{14.400000}\selectfont Sine Wave}%
\end{pgfscope}%
\begin{pgfscope}%
\pgfsetbuttcap%
\pgfsetmiterjoin%
\definecolor{currentfill}{rgb}{1.000000,1.000000,1.000000}%
\pgfsetfillcolor{currentfill}%
\pgfsetlinewidth{0.000000pt}%
\definecolor{currentstroke}{rgb}{0.000000,0.000000,0.000000}%
\pgfsetstrokecolor{currentstroke}%
\pgfsetstrokeopacity{0.000000}%
\pgfsetdash{}{0pt}%
\pgfpathmoveto{\pgfqpoint{3.833026in}{2.240123in}}%
\pgfpathlineto{\pgfqpoint{6.115278in}{2.240123in}}%
\pgfpathlineto{\pgfqpoint{6.115278in}{3.150926in}}%
\pgfpathlineto{\pgfqpoint{3.833026in}{3.150926in}}%
\pgfpathlineto{\pgfqpoint{3.833026in}{2.240123in}}%
\pgfpathclose%
\pgfusepath{fill}%
\end{pgfscope}%
\begin{pgfscope}%
\pgfpathrectangle{\pgfqpoint{3.833026in}{2.240123in}}{\pgfqpoint{2.282252in}{0.910803in}}%
\pgfusepath{clip}%
\pgfsetbuttcap%
\pgfsetroundjoin%
\pgfsetlinewidth{0.501875pt}%
\definecolor{currentstroke}{rgb}{0.690196,0.690196,0.690196}%
\pgfsetstrokecolor{currentstroke}%
\pgfsetstrokeopacity{0.700000}%
\pgfsetdash{{1.850000pt}{0.800000pt}}{0.000000pt}%
\pgfpathmoveto{\pgfqpoint{4.252215in}{2.240123in}}%
\pgfpathlineto{\pgfqpoint{4.252215in}{3.150926in}}%
\pgfusepath{stroke}%
\end{pgfscope}%
\begin{pgfscope}%
\pgfsetbuttcap%
\pgfsetroundjoin%
\definecolor{currentfill}{rgb}{0.000000,0.000000,0.000000}%
\pgfsetfillcolor{currentfill}%
\pgfsetlinewidth{0.803000pt}%
\definecolor{currentstroke}{rgb}{0.000000,0.000000,0.000000}%
\pgfsetstrokecolor{currentstroke}%
\pgfsetdash{}{0pt}%
\pgfsys@defobject{currentmarker}{\pgfqpoint{0.000000in}{-0.048611in}}{\pgfqpoint{0.000000in}{0.000000in}}{%
\pgfpathmoveto{\pgfqpoint{0.000000in}{0.000000in}}%
\pgfpathlineto{\pgfqpoint{0.000000in}{-0.048611in}}%
\pgfusepath{stroke,fill}%
}%
\begin{pgfscope}%
\pgfsys@transformshift{4.252215in}{2.240123in}%
\pgfsys@useobject{currentmarker}{}%
\end{pgfscope}%
\end{pgfscope}%
\begin{pgfscope}%
\definecolor{textcolor}{rgb}{0.000000,0.000000,0.000000}%
\pgfsetstrokecolor{textcolor}%
\pgfsetfillcolor{textcolor}%
\pgftext[x=4.252215in,y=2.142901in,,top]{\color{textcolor}\rmfamily\fontsize{10.000000}{12.000000}\selectfont \(\displaystyle {1}\)}%
\end{pgfscope}%
\begin{pgfscope}%
\pgfpathrectangle{\pgfqpoint{3.833026in}{2.240123in}}{\pgfqpoint{2.282252in}{0.910803in}}%
\pgfusepath{clip}%
\pgfsetbuttcap%
\pgfsetroundjoin%
\pgfsetlinewidth{0.501875pt}%
\definecolor{currentstroke}{rgb}{0.690196,0.690196,0.690196}%
\pgfsetstrokecolor{currentstroke}%
\pgfsetstrokeopacity{0.700000}%
\pgfsetdash{{1.850000pt}{0.800000pt}}{0.000000pt}%
\pgfpathmoveto{\pgfqpoint{4.717981in}{2.240123in}}%
\pgfpathlineto{\pgfqpoint{4.717981in}{3.150926in}}%
\pgfusepath{stroke}%
\end{pgfscope}%
\begin{pgfscope}%
\pgfsetbuttcap%
\pgfsetroundjoin%
\definecolor{currentfill}{rgb}{0.000000,0.000000,0.000000}%
\pgfsetfillcolor{currentfill}%
\pgfsetlinewidth{0.803000pt}%
\definecolor{currentstroke}{rgb}{0.000000,0.000000,0.000000}%
\pgfsetstrokecolor{currentstroke}%
\pgfsetdash{}{0pt}%
\pgfsys@defobject{currentmarker}{\pgfqpoint{0.000000in}{-0.048611in}}{\pgfqpoint{0.000000in}{0.000000in}}{%
\pgfpathmoveto{\pgfqpoint{0.000000in}{0.000000in}}%
\pgfpathlineto{\pgfqpoint{0.000000in}{-0.048611in}}%
\pgfusepath{stroke,fill}%
}%
\begin{pgfscope}%
\pgfsys@transformshift{4.717981in}{2.240123in}%
\pgfsys@useobject{currentmarker}{}%
\end{pgfscope}%
\end{pgfscope}%
\begin{pgfscope}%
\definecolor{textcolor}{rgb}{0.000000,0.000000,0.000000}%
\pgfsetstrokecolor{textcolor}%
\pgfsetfillcolor{textcolor}%
\pgftext[x=4.717981in,y=2.142901in,,top]{\color{textcolor}\rmfamily\fontsize{10.000000}{12.000000}\selectfont \(\displaystyle {2}\)}%
\end{pgfscope}%
\begin{pgfscope}%
\pgfpathrectangle{\pgfqpoint{3.833026in}{2.240123in}}{\pgfqpoint{2.282252in}{0.910803in}}%
\pgfusepath{clip}%
\pgfsetbuttcap%
\pgfsetroundjoin%
\pgfsetlinewidth{0.501875pt}%
\definecolor{currentstroke}{rgb}{0.690196,0.690196,0.690196}%
\pgfsetstrokecolor{currentstroke}%
\pgfsetstrokeopacity{0.700000}%
\pgfsetdash{{1.850000pt}{0.800000pt}}{0.000000pt}%
\pgfpathmoveto{\pgfqpoint{5.183746in}{2.240123in}}%
\pgfpathlineto{\pgfqpoint{5.183746in}{3.150926in}}%
\pgfusepath{stroke}%
\end{pgfscope}%
\begin{pgfscope}%
\pgfsetbuttcap%
\pgfsetroundjoin%
\definecolor{currentfill}{rgb}{0.000000,0.000000,0.000000}%
\pgfsetfillcolor{currentfill}%
\pgfsetlinewidth{0.803000pt}%
\definecolor{currentstroke}{rgb}{0.000000,0.000000,0.000000}%
\pgfsetstrokecolor{currentstroke}%
\pgfsetdash{}{0pt}%
\pgfsys@defobject{currentmarker}{\pgfqpoint{0.000000in}{-0.048611in}}{\pgfqpoint{0.000000in}{0.000000in}}{%
\pgfpathmoveto{\pgfqpoint{0.000000in}{0.000000in}}%
\pgfpathlineto{\pgfqpoint{0.000000in}{-0.048611in}}%
\pgfusepath{stroke,fill}%
}%
\begin{pgfscope}%
\pgfsys@transformshift{5.183746in}{2.240123in}%
\pgfsys@useobject{currentmarker}{}%
\end{pgfscope}%
\end{pgfscope}%
\begin{pgfscope}%
\definecolor{textcolor}{rgb}{0.000000,0.000000,0.000000}%
\pgfsetstrokecolor{textcolor}%
\pgfsetfillcolor{textcolor}%
\pgftext[x=5.183746in,y=2.142901in,,top]{\color{textcolor}\rmfamily\fontsize{10.000000}{12.000000}\selectfont \(\displaystyle {3}\)}%
\end{pgfscope}%
\begin{pgfscope}%
\pgfpathrectangle{\pgfqpoint{3.833026in}{2.240123in}}{\pgfqpoint{2.282252in}{0.910803in}}%
\pgfusepath{clip}%
\pgfsetbuttcap%
\pgfsetroundjoin%
\pgfsetlinewidth{0.501875pt}%
\definecolor{currentstroke}{rgb}{0.690196,0.690196,0.690196}%
\pgfsetstrokecolor{currentstroke}%
\pgfsetstrokeopacity{0.700000}%
\pgfsetdash{{1.850000pt}{0.800000pt}}{0.000000pt}%
\pgfpathmoveto{\pgfqpoint{5.649512in}{2.240123in}}%
\pgfpathlineto{\pgfqpoint{5.649512in}{3.150926in}}%
\pgfusepath{stroke}%
\end{pgfscope}%
\begin{pgfscope}%
\pgfsetbuttcap%
\pgfsetroundjoin%
\definecolor{currentfill}{rgb}{0.000000,0.000000,0.000000}%
\pgfsetfillcolor{currentfill}%
\pgfsetlinewidth{0.803000pt}%
\definecolor{currentstroke}{rgb}{0.000000,0.000000,0.000000}%
\pgfsetstrokecolor{currentstroke}%
\pgfsetdash{}{0pt}%
\pgfsys@defobject{currentmarker}{\pgfqpoint{0.000000in}{-0.048611in}}{\pgfqpoint{0.000000in}{0.000000in}}{%
\pgfpathmoveto{\pgfqpoint{0.000000in}{0.000000in}}%
\pgfpathlineto{\pgfqpoint{0.000000in}{-0.048611in}}%
\pgfusepath{stroke,fill}%
}%
\begin{pgfscope}%
\pgfsys@transformshift{5.649512in}{2.240123in}%
\pgfsys@useobject{currentmarker}{}%
\end{pgfscope}%
\end{pgfscope}%
\begin{pgfscope}%
\definecolor{textcolor}{rgb}{0.000000,0.000000,0.000000}%
\pgfsetstrokecolor{textcolor}%
\pgfsetfillcolor{textcolor}%
\pgftext[x=5.649512in,y=2.142901in,,top]{\color{textcolor}\rmfamily\fontsize{10.000000}{12.000000}\selectfont \(\displaystyle {4}\)}%
\end{pgfscope}%
\begin{pgfscope}%
\pgfpathrectangle{\pgfqpoint{3.833026in}{2.240123in}}{\pgfqpoint{2.282252in}{0.910803in}}%
\pgfusepath{clip}%
\pgfsetbuttcap%
\pgfsetroundjoin%
\pgfsetlinewidth{0.501875pt}%
\definecolor{currentstroke}{rgb}{0.690196,0.690196,0.690196}%
\pgfsetstrokecolor{currentstroke}%
\pgfsetstrokeopacity{0.700000}%
\pgfsetdash{{1.850000pt}{0.800000pt}}{0.000000pt}%
\pgfpathmoveto{\pgfqpoint{6.115278in}{2.240123in}}%
\pgfpathlineto{\pgfqpoint{6.115278in}{3.150926in}}%
\pgfusepath{stroke}%
\end{pgfscope}%
\begin{pgfscope}%
\pgfsetbuttcap%
\pgfsetroundjoin%
\definecolor{currentfill}{rgb}{0.000000,0.000000,0.000000}%
\pgfsetfillcolor{currentfill}%
\pgfsetlinewidth{0.803000pt}%
\definecolor{currentstroke}{rgb}{0.000000,0.000000,0.000000}%
\pgfsetstrokecolor{currentstroke}%
\pgfsetdash{}{0pt}%
\pgfsys@defobject{currentmarker}{\pgfqpoint{0.000000in}{-0.048611in}}{\pgfqpoint{0.000000in}{0.000000in}}{%
\pgfpathmoveto{\pgfqpoint{0.000000in}{0.000000in}}%
\pgfpathlineto{\pgfqpoint{0.000000in}{-0.048611in}}%
\pgfusepath{stroke,fill}%
}%
\begin{pgfscope}%
\pgfsys@transformshift{6.115278in}{2.240123in}%
\pgfsys@useobject{currentmarker}{}%
\end{pgfscope}%
\end{pgfscope}%
\begin{pgfscope}%
\definecolor{textcolor}{rgb}{0.000000,0.000000,0.000000}%
\pgfsetstrokecolor{textcolor}%
\pgfsetfillcolor{textcolor}%
\pgftext[x=6.115278in,y=2.142901in,,top]{\color{textcolor}\rmfamily\fontsize{10.000000}{12.000000}\selectfont \(\displaystyle {5}\)}%
\end{pgfscope}%
\begin{pgfscope}%
\definecolor{textcolor}{rgb}{0.000000,0.000000,0.000000}%
\pgfsetstrokecolor{textcolor}%
\pgfsetfillcolor{textcolor}%
\pgftext[x=4.974152in,y=1.963889in,,top]{\color{textcolor}\rmfamily\fontsize{10.000000}{12.000000}\selectfont Frequency (Hz)}%
\end{pgfscope}%
\begin{pgfscope}%
\definecolor{textcolor}{rgb}{0.000000,0.000000,0.000000}%
\pgfsetstrokecolor{textcolor}%
\pgfsetfillcolor{textcolor}%
\pgftext[x=6.115278in,y=1.977778in,right,top]{\color{textcolor}\rmfamily\fontsize{10.000000}{12.000000}\selectfont \(\displaystyle \times{10^{6}}{}\)}%
\end{pgfscope}%
\begin{pgfscope}%
\pgfpathrectangle{\pgfqpoint{3.833026in}{2.240123in}}{\pgfqpoint{2.282252in}{0.910803in}}%
\pgfusepath{clip}%
\pgfsetbuttcap%
\pgfsetroundjoin%
\pgfsetlinewidth{0.501875pt}%
\definecolor{currentstroke}{rgb}{0.690196,0.690196,0.690196}%
\pgfsetstrokecolor{currentstroke}%
\pgfsetstrokeopacity{0.700000}%
\pgfsetdash{{1.850000pt}{0.800000pt}}{0.000000pt}%
\pgfpathmoveto{\pgfqpoint{3.833026in}{2.359963in}}%
\pgfpathlineto{\pgfqpoint{6.115278in}{2.359963in}}%
\pgfusepath{stroke}%
\end{pgfscope}%
\begin{pgfscope}%
\pgfsetbuttcap%
\pgfsetroundjoin%
\definecolor{currentfill}{rgb}{0.000000,0.000000,0.000000}%
\pgfsetfillcolor{currentfill}%
\pgfsetlinewidth{0.803000pt}%
\definecolor{currentstroke}{rgb}{0.000000,0.000000,0.000000}%
\pgfsetstrokecolor{currentstroke}%
\pgfsetdash{}{0pt}%
\pgfsys@defobject{currentmarker}{\pgfqpoint{-0.048611in}{0.000000in}}{\pgfqpoint{-0.000000in}{0.000000in}}{%
\pgfpathmoveto{\pgfqpoint{-0.000000in}{0.000000in}}%
\pgfpathlineto{\pgfqpoint{-0.048611in}{0.000000in}}%
\pgfusepath{stroke,fill}%
}%
\begin{pgfscope}%
\pgfsys@transformshift{3.833026in}{2.359963in}%
\pgfsys@useobject{currentmarker}{}%
\end{pgfscope}%
\end{pgfscope}%
\begin{pgfscope}%
\definecolor{textcolor}{rgb}{0.000000,0.000000,0.000000}%
\pgfsetstrokecolor{textcolor}%
\pgfsetfillcolor{textcolor}%
\pgftext[x=3.419444in, y=2.311737in, left, base]{\color{textcolor}\rmfamily\fontsize{10.000000}{12.000000}\selectfont \(\displaystyle {\ensuremath{-}100}\)}%
\end{pgfscope}%
\begin{pgfscope}%
\pgfpathrectangle{\pgfqpoint{3.833026in}{2.240123in}}{\pgfqpoint{2.282252in}{0.910803in}}%
\pgfusepath{clip}%
\pgfsetbuttcap%
\pgfsetroundjoin%
\pgfsetlinewidth{0.501875pt}%
\definecolor{currentstroke}{rgb}{0.690196,0.690196,0.690196}%
\pgfsetstrokecolor{currentstroke}%
\pgfsetstrokeopacity{0.700000}%
\pgfsetdash{{1.850000pt}{0.800000pt}}{0.000000pt}%
\pgfpathmoveto{\pgfqpoint{3.833026in}{2.828257in}}%
\pgfpathlineto{\pgfqpoint{6.115278in}{2.828257in}}%
\pgfusepath{stroke}%
\end{pgfscope}%
\begin{pgfscope}%
\pgfsetbuttcap%
\pgfsetroundjoin%
\definecolor{currentfill}{rgb}{0.000000,0.000000,0.000000}%
\pgfsetfillcolor{currentfill}%
\pgfsetlinewidth{0.803000pt}%
\definecolor{currentstroke}{rgb}{0.000000,0.000000,0.000000}%
\pgfsetstrokecolor{currentstroke}%
\pgfsetdash{}{0pt}%
\pgfsys@defobject{currentmarker}{\pgfqpoint{-0.048611in}{0.000000in}}{\pgfqpoint{-0.000000in}{0.000000in}}{%
\pgfpathmoveto{\pgfqpoint{-0.000000in}{0.000000in}}%
\pgfpathlineto{\pgfqpoint{-0.048611in}{0.000000in}}%
\pgfusepath{stroke,fill}%
}%
\begin{pgfscope}%
\pgfsys@transformshift{3.833026in}{2.828257in}%
\pgfsys@useobject{currentmarker}{}%
\end{pgfscope}%
\end{pgfscope}%
\begin{pgfscope}%
\definecolor{textcolor}{rgb}{0.000000,0.000000,0.000000}%
\pgfsetstrokecolor{textcolor}%
\pgfsetfillcolor{textcolor}%
\pgftext[x=3.488889in, y=2.780032in, left, base]{\color{textcolor}\rmfamily\fontsize{10.000000}{12.000000}\selectfont \(\displaystyle {\ensuremath{-}50}\)}%
\end{pgfscope}%
\begin{pgfscope}%
\definecolor{textcolor}{rgb}{0.000000,0.000000,0.000000}%
\pgfsetstrokecolor{textcolor}%
\pgfsetfillcolor{textcolor}%
\pgftext[x=3.363889in,y=2.695525in,,bottom,rotate=90.000000]{\color{textcolor}\rmfamily\fontsize{10.000000}{12.000000}\selectfont Amplitude (dB)}%
\end{pgfscope}%
\begin{pgfscope}%
\pgfpathrectangle{\pgfqpoint{3.833026in}{2.240123in}}{\pgfqpoint{2.282252in}{0.910803in}}%
\pgfusepath{clip}%
\pgfsetrectcap%
\pgfsetroundjoin%
\pgfsetlinewidth{1.003750pt}%
\definecolor{currentstroke}{rgb}{0.000000,0.000000,0.000000}%
\pgfsetstrokecolor{currentstroke}%
\pgfsetdash{}{0pt}%
\pgfpathmoveto{\pgfqpoint{3.833026in}{2.558988in}}%
\pgfpathlineto{\pgfqpoint{3.835308in}{2.526207in}}%
\pgfpathlineto{\pgfqpoint{3.837590in}{2.582110in}}%
\pgfpathlineto{\pgfqpoint{3.839872in}{2.554012in}}%
\pgfpathlineto{\pgfqpoint{3.842155in}{2.581817in}}%
\pgfpathlineto{\pgfqpoint{3.846719in}{2.581524in}}%
\pgfpathlineto{\pgfqpoint{3.849001in}{2.539085in}}%
\pgfpathlineto{\pgfqpoint{3.851284in}{2.552841in}}%
\pgfpathlineto{\pgfqpoint{3.853566in}{2.538500in}}%
\pgfpathlineto{\pgfqpoint{3.855848in}{2.594695in}}%
\pgfpathlineto{\pgfqpoint{3.858130in}{2.538207in}}%
\pgfpathlineto{\pgfqpoint{3.860413in}{2.566012in}}%
\pgfpathlineto{\pgfqpoint{3.862695in}{2.523866in}}%
\pgfpathlineto{\pgfqpoint{3.864977in}{2.537622in}}%
\pgfpathlineto{\pgfqpoint{3.867259in}{2.626305in}}%
\pgfpathlineto{\pgfqpoint{3.869542in}{2.570110in}}%
\pgfpathlineto{\pgfqpoint{3.871824in}{2.607281in}}%
\pgfpathlineto{\pgfqpoint{3.874106in}{2.541427in}}%
\pgfpathlineto{\pgfqpoint{3.878671in}{2.592646in}}%
\pgfpathlineto{\pgfqpoint{3.880953in}{2.536158in}}%
\pgfpathlineto{\pgfqpoint{3.885517in}{2.568646in}}%
\pgfpathlineto{\pgfqpoint{3.887800in}{2.563671in}}%
\pgfpathlineto{\pgfqpoint{3.890082in}{2.563671in}}%
\pgfpathlineto{\pgfqpoint{3.894646in}{2.549037in}}%
\pgfpathlineto{\pgfqpoint{3.899211in}{2.623671in}}%
\pgfpathlineto{\pgfqpoint{3.901493in}{2.567183in}}%
\pgfpathlineto{\pgfqpoint{3.903776in}{2.604354in}}%
\pgfpathlineto{\pgfqpoint{3.908340in}{2.585329in}}%
\pgfpathlineto{\pgfqpoint{3.910622in}{2.608451in}}%
\pgfpathlineto{\pgfqpoint{3.912905in}{2.585037in}}%
\pgfpathlineto{\pgfqpoint{3.915187in}{2.537915in}}%
\pgfpathlineto{\pgfqpoint{3.917469in}{2.537622in}}%
\pgfpathlineto{\pgfqpoint{3.922034in}{2.579476in}}%
\pgfpathlineto{\pgfqpoint{3.924316in}{2.583866in}}%
\pgfpathlineto{\pgfqpoint{3.926598in}{2.616646in}}%
\pgfpathlineto{\pgfqpoint{3.931163in}{2.517719in}}%
\pgfpathlineto{\pgfqpoint{3.933445in}{2.602012in}}%
\pgfpathlineto{\pgfqpoint{3.938009in}{2.568646in}}%
\pgfpathlineto{\pgfqpoint{3.940292in}{2.587378in}}%
\pgfpathlineto{\pgfqpoint{3.942574in}{2.577720in}}%
\pgfpathlineto{\pgfqpoint{3.944856in}{2.577427in}}%
\pgfpathlineto{\pgfqpoint{3.949421in}{2.614598in}}%
\pgfpathlineto{\pgfqpoint{3.951703in}{2.581524in}}%
\pgfpathlineto{\pgfqpoint{3.953985in}{2.585915in}}%
\pgfpathlineto{\pgfqpoint{3.956267in}{2.576549in}}%
\pgfpathlineto{\pgfqpoint{3.958550in}{2.557524in}}%
\pgfpathlineto{\pgfqpoint{3.960832in}{2.613427in}}%
\pgfpathlineto{\pgfqpoint{3.963114in}{2.599378in}}%
\pgfpathlineto{\pgfqpoint{3.965396in}{2.613134in}}%
\pgfpathlineto{\pgfqpoint{3.967679in}{2.561329in}}%
\pgfpathlineto{\pgfqpoint{3.969961in}{2.570695in}}%
\pgfpathlineto{\pgfqpoint{3.972243in}{2.664061in}}%
\pgfpathlineto{\pgfqpoint{3.974525in}{2.560744in}}%
\pgfpathlineto{\pgfqpoint{3.976808in}{2.574793in}}%
\pgfpathlineto{\pgfqpoint{3.979090in}{2.574500in}}%
\pgfpathlineto{\pgfqpoint{3.983654in}{2.527378in}}%
\pgfpathlineto{\pgfqpoint{3.985937in}{2.527085in}}%
\pgfpathlineto{\pgfqpoint{3.988219in}{2.573622in}}%
\pgfpathlineto{\pgfqpoint{3.990501in}{2.526793in}}%
\pgfpathlineto{\pgfqpoint{3.992783in}{2.512451in}}%
\pgfpathlineto{\pgfqpoint{3.995066in}{2.530890in}}%
\pgfpathlineto{\pgfqpoint{3.997348in}{2.568354in}}%
\pgfpathlineto{\pgfqpoint{4.001912in}{2.553719in}}%
\pgfpathlineto{\pgfqpoint{4.004195in}{2.609915in}}%
\pgfpathlineto{\pgfqpoint{4.006477in}{2.633037in}}%
\pgfpathlineto{\pgfqpoint{4.008759in}{2.749818in}}%
\pgfpathlineto{\pgfqpoint{4.011041in}{2.585915in}}%
\pgfpathlineto{\pgfqpoint{4.013324in}{2.894696in}}%
\pgfpathlineto{\pgfqpoint{4.015606in}{3.044258in}}%
\pgfpathlineto{\pgfqpoint{4.017888in}{3.100160in}}%
\pgfpathlineto{\pgfqpoint{4.020170in}{3.109526in}}%
\pgfpathlineto{\pgfqpoint{4.022453in}{3.071770in}}%
\pgfpathlineto{\pgfqpoint{4.024735in}{2.973135in}}%
\pgfpathlineto{\pgfqpoint{4.027017in}{2.701525in}}%
\pgfpathlineto{\pgfqpoint{4.029299in}{2.630988in}}%
\pgfpathlineto{\pgfqpoint{4.031582in}{2.654110in}}%
\pgfpathlineto{\pgfqpoint{4.033864in}{2.729037in}}%
\pgfpathlineto{\pgfqpoint{4.038428in}{2.522402in}}%
\pgfpathlineto{\pgfqpoint{4.040711in}{2.597329in}}%
\pgfpathlineto{\pgfqpoint{4.042993in}{2.526793in}}%
\pgfpathlineto{\pgfqpoint{4.045275in}{2.582695in}}%
\pgfpathlineto{\pgfqpoint{4.047557in}{2.568646in}}%
\pgfpathlineto{\pgfqpoint{4.049840in}{2.629232in}}%
\pgfpathlineto{\pgfqpoint{4.054404in}{2.507183in}}%
\pgfpathlineto{\pgfqpoint{4.081791in}{2.504549in}}%
\pgfpathlineto{\pgfqpoint{4.086356in}{2.504256in}}%
\pgfpathlineto{\pgfqpoint{4.088638in}{2.518012in}}%
\pgfpathlineto{\pgfqpoint{4.090920in}{2.503963in}}%
\pgfpathlineto{\pgfqpoint{4.095485in}{2.503378in}}%
\pgfpathlineto{\pgfqpoint{4.097767in}{2.517427in}}%
\pgfpathlineto{\pgfqpoint{4.100049in}{2.503085in}}%
\pgfpathlineto{\pgfqpoint{4.102331in}{2.516841in}}%
\pgfpathlineto{\pgfqpoint{4.104614in}{2.516841in}}%
\pgfpathlineto{\pgfqpoint{4.106896in}{2.502500in}}%
\pgfpathlineto{\pgfqpoint{4.132001in}{2.500158in}}%
\pgfpathlineto{\pgfqpoint{4.143412in}{2.499280in}}%
\pgfpathlineto{\pgfqpoint{4.157105in}{2.498110in}}%
\pgfpathlineto{\pgfqpoint{4.159388in}{2.483768in}}%
\pgfpathlineto{\pgfqpoint{4.161670in}{2.497817in}}%
\pgfpathlineto{\pgfqpoint{4.163952in}{2.497524in}}%
\pgfpathlineto{\pgfqpoint{4.166234in}{2.469134in}}%
\pgfpathlineto{\pgfqpoint{4.168517in}{2.469134in}}%
\pgfpathlineto{\pgfqpoint{4.170799in}{2.496939in}}%
\pgfpathlineto{\pgfqpoint{4.173081in}{2.496646in}}%
\pgfpathlineto{\pgfqpoint{4.175363in}{2.529134in}}%
\pgfpathlineto{\pgfqpoint{4.177646in}{2.468256in}}%
\pgfpathlineto{\pgfqpoint{4.179928in}{2.467963in}}%
\pgfpathlineto{\pgfqpoint{4.182210in}{2.495768in}}%
\pgfpathlineto{\pgfqpoint{4.186775in}{2.453329in}}%
\pgfpathlineto{\pgfqpoint{4.189057in}{2.438987in}}%
\pgfpathlineto{\pgfqpoint{4.191339in}{2.438987in}}%
\pgfpathlineto{\pgfqpoint{4.193622in}{2.480841in}}%
\pgfpathlineto{\pgfqpoint{4.195904in}{2.452451in}}%
\pgfpathlineto{\pgfqpoint{4.198186in}{2.438402in}}%
\pgfpathlineto{\pgfqpoint{4.200468in}{2.470890in}}%
\pgfpathlineto{\pgfqpoint{4.202751in}{2.437817in}}%
\pgfpathlineto{\pgfqpoint{4.205033in}{2.437817in}}%
\pgfpathlineto{\pgfqpoint{4.207315in}{2.451573in}}%
\pgfpathlineto{\pgfqpoint{4.209597in}{2.437231in}}%
\pgfpathlineto{\pgfqpoint{4.211880in}{2.479378in}}%
\pgfpathlineto{\pgfqpoint{4.214162in}{2.450987in}}%
\pgfpathlineto{\pgfqpoint{4.216444in}{2.492841in}}%
\pgfpathlineto{\pgfqpoint{4.218726in}{2.450695in}}%
\pgfpathlineto{\pgfqpoint{4.221009in}{2.478500in}}%
\pgfpathlineto{\pgfqpoint{4.223291in}{2.436061in}}%
\pgfpathlineto{\pgfqpoint{4.227855in}{2.435768in}}%
\pgfpathlineto{\pgfqpoint{4.230138in}{2.468256in}}%
\pgfpathlineto{\pgfqpoint{4.232420in}{2.435182in}}%
\pgfpathlineto{\pgfqpoint{4.236984in}{2.434890in}}%
\pgfpathlineto{\pgfqpoint{4.239267in}{2.462695in}}%
\pgfpathlineto{\pgfqpoint{4.241549in}{2.434597in}}%
\pgfpathlineto{\pgfqpoint{4.246113in}{2.434012in}}%
\pgfpathlineto{\pgfqpoint{4.250678in}{2.578890in}}%
\pgfpathlineto{\pgfqpoint{4.252960in}{2.587963in}}%
\pgfpathlineto{\pgfqpoint{4.255242in}{2.564549in}}%
\pgfpathlineto{\pgfqpoint{4.259807in}{2.405329in}}%
\pgfpathlineto{\pgfqpoint{4.262089in}{2.447475in}}%
\pgfpathlineto{\pgfqpoint{4.264371in}{2.438109in}}%
\pgfpathlineto{\pgfqpoint{4.266654in}{2.447475in}}%
\pgfpathlineto{\pgfqpoint{4.268936in}{2.405329in}}%
\pgfpathlineto{\pgfqpoint{4.271218in}{2.405329in}}%
\pgfpathlineto{\pgfqpoint{4.275783in}{2.433426in}}%
\pgfpathlineto{\pgfqpoint{4.278065in}{2.419378in}}%
\pgfpathlineto{\pgfqpoint{4.280347in}{2.447475in}}%
\pgfpathlineto{\pgfqpoint{4.282629in}{2.419378in}}%
\pgfpathlineto{\pgfqpoint{4.284912in}{2.433426in}}%
\pgfpathlineto{\pgfqpoint{4.287194in}{2.433426in}}%
\pgfpathlineto{\pgfqpoint{4.289476in}{2.452158in}}%
\pgfpathlineto{\pgfqpoint{4.291758in}{2.433426in}}%
\pgfpathlineto{\pgfqpoint{4.294041in}{2.447475in}}%
\pgfpathlineto{\pgfqpoint{4.298605in}{2.433426in}}%
\pgfpathlineto{\pgfqpoint{4.300887in}{2.405329in}}%
\pgfpathlineto{\pgfqpoint{4.303170in}{2.419378in}}%
\pgfpathlineto{\pgfqpoint{4.305452in}{2.419378in}}%
\pgfpathlineto{\pgfqpoint{4.307734in}{2.405329in}}%
\pgfpathlineto{\pgfqpoint{4.310016in}{2.438109in}}%
\pgfpathlineto{\pgfqpoint{4.312299in}{2.405329in}}%
\pgfpathlineto{\pgfqpoint{4.314581in}{2.419378in}}%
\pgfpathlineto{\pgfqpoint{4.316863in}{2.405329in}}%
\pgfpathlineto{\pgfqpoint{4.319145in}{2.419378in}}%
\pgfpathlineto{\pgfqpoint{4.323710in}{2.419378in}}%
\pgfpathlineto{\pgfqpoint{4.325992in}{2.405329in}}%
\pgfpathlineto{\pgfqpoint{4.328274in}{2.419378in}}%
\pgfpathlineto{\pgfqpoint{4.330557in}{2.405329in}}%
\pgfpathlineto{\pgfqpoint{4.332839in}{2.419378in}}%
\pgfpathlineto{\pgfqpoint{4.335121in}{2.405036in}}%
\pgfpathlineto{\pgfqpoint{4.348815in}{2.405036in}}%
\pgfpathlineto{\pgfqpoint{4.351097in}{2.419085in}}%
\pgfpathlineto{\pgfqpoint{4.353379in}{2.405036in}}%
\pgfpathlineto{\pgfqpoint{4.385331in}{2.405036in}}%
\pgfpathlineto{\pgfqpoint{4.387613in}{2.376938in}}%
\pgfpathlineto{\pgfqpoint{4.392177in}{2.376938in}}%
\pgfpathlineto{\pgfqpoint{4.394460in}{2.419085in}}%
\pgfpathlineto{\pgfqpoint{4.396742in}{2.405036in}}%
\pgfpathlineto{\pgfqpoint{4.399024in}{2.405036in}}%
\pgfpathlineto{\pgfqpoint{4.401306in}{2.376938in}}%
\pgfpathlineto{\pgfqpoint{4.403589in}{2.376938in}}%
\pgfpathlineto{\pgfqpoint{4.405871in}{2.381621in}}%
\pgfpathlineto{\pgfqpoint{4.408153in}{2.395670in}}%
\pgfpathlineto{\pgfqpoint{4.410435in}{2.395670in}}%
\pgfpathlineto{\pgfqpoint{4.412718in}{2.381329in}}%
\pgfpathlineto{\pgfqpoint{4.415000in}{2.428158in}}%
\pgfpathlineto{\pgfqpoint{4.417282in}{2.517134in}}%
\pgfpathlineto{\pgfqpoint{4.419564in}{2.404743in}}%
\pgfpathlineto{\pgfqpoint{4.421847in}{2.451573in}}%
\pgfpathlineto{\pgfqpoint{4.424129in}{2.390695in}}%
\pgfpathlineto{\pgfqpoint{4.426411in}{2.451573in}}%
\pgfpathlineto{\pgfqpoint{4.430976in}{2.409426in}}%
\pgfpathlineto{\pgfqpoint{4.433258in}{2.418792in}}%
\pgfpathlineto{\pgfqpoint{4.435540in}{2.367280in}}%
\pgfpathlineto{\pgfqpoint{4.437822in}{2.432841in}}%
\pgfpathlineto{\pgfqpoint{4.440105in}{2.409426in}}%
\pgfpathlineto{\pgfqpoint{4.442387in}{2.418792in}}%
\pgfpathlineto{\pgfqpoint{4.444669in}{2.376646in}}%
\pgfpathlineto{\pgfqpoint{4.446951in}{2.437524in}}%
\pgfpathlineto{\pgfqpoint{4.449234in}{2.395377in}}%
\pgfpathlineto{\pgfqpoint{4.456080in}{2.428158in}}%
\pgfpathlineto{\pgfqpoint{4.458363in}{2.442207in}}%
\pgfpathlineto{\pgfqpoint{4.460645in}{2.479670in}}%
\pgfpathlineto{\pgfqpoint{4.462927in}{2.432841in}}%
\pgfpathlineto{\pgfqpoint{4.465209in}{2.437524in}}%
\pgfpathlineto{\pgfqpoint{4.467492in}{2.521817in}}%
\pgfpathlineto{\pgfqpoint{4.469774in}{2.554598in}}%
\pgfpathlineto{\pgfqpoint{4.472056in}{2.507768in}}%
\pgfpathlineto{\pgfqpoint{4.474338in}{2.521817in}}%
\pgfpathlineto{\pgfqpoint{4.476621in}{2.456256in}}%
\pgfpathlineto{\pgfqpoint{4.481185in}{2.905818in}}%
\pgfpathlineto{\pgfqpoint{4.483467in}{2.985428in}}%
\pgfpathlineto{\pgfqpoint{4.485750in}{3.013526in}}%
\pgfpathlineto{\pgfqpoint{4.488032in}{2.990111in}}%
\pgfpathlineto{\pgfqpoint{4.490314in}{2.905818in}}%
\pgfpathlineto{\pgfqpoint{4.494879in}{2.507475in}}%
\pgfpathlineto{\pgfqpoint{4.497161in}{2.535573in}}%
\pgfpathlineto{\pgfqpoint{4.499443in}{2.512158in}}%
\pgfpathlineto{\pgfqpoint{4.501726in}{2.502792in}}%
\pgfpathlineto{\pgfqpoint{4.504008in}{2.507475in}}%
\pgfpathlineto{\pgfqpoint{4.506290in}{2.465329in}}%
\pgfpathlineto{\pgfqpoint{4.508572in}{2.366987in}}%
\pgfpathlineto{\pgfqpoint{4.510855in}{2.540256in}}%
\pgfpathlineto{\pgfqpoint{4.515419in}{2.577720in}}%
\pgfpathlineto{\pgfqpoint{4.517701in}{2.544939in}}%
\pgfpathlineto{\pgfqpoint{4.522266in}{2.381036in}}%
\pgfpathlineto{\pgfqpoint{4.524548in}{2.376353in}}%
\pgfpathlineto{\pgfqpoint{4.526830in}{2.432548in}}%
\pgfpathlineto{\pgfqpoint{4.529113in}{2.395085in}}%
\pgfpathlineto{\pgfqpoint{4.531395in}{2.376353in}}%
\pgfpathlineto{\pgfqpoint{4.533677in}{2.404451in}}%
\pgfpathlineto{\pgfqpoint{4.535959in}{2.362304in}}%
\pgfpathlineto{\pgfqpoint{4.538242in}{2.376353in}}%
\pgfpathlineto{\pgfqpoint{4.540524in}{2.376353in}}%
\pgfpathlineto{\pgfqpoint{4.542806in}{2.362304in}}%
\pgfpathlineto{\pgfqpoint{4.545088in}{2.390402in}}%
\pgfpathlineto{\pgfqpoint{4.547371in}{2.404451in}}%
\pgfpathlineto{\pgfqpoint{4.549653in}{2.432548in}}%
\pgfpathlineto{\pgfqpoint{4.551935in}{2.376353in}}%
\pgfpathlineto{\pgfqpoint{4.554217in}{2.441914in}}%
\pgfpathlineto{\pgfqpoint{4.556500in}{2.409134in}}%
\pgfpathlineto{\pgfqpoint{4.558782in}{2.418500in}}%
\pgfpathlineto{\pgfqpoint{4.561064in}{2.376353in}}%
\pgfpathlineto{\pgfqpoint{4.563346in}{2.390402in}}%
\pgfpathlineto{\pgfqpoint{4.565629in}{2.362304in}}%
\pgfpathlineto{\pgfqpoint{4.567911in}{2.376353in}}%
\pgfpathlineto{\pgfqpoint{4.570193in}{2.366987in}}%
\pgfpathlineto{\pgfqpoint{4.572475in}{2.390402in}}%
\pgfpathlineto{\pgfqpoint{4.574758in}{2.441622in}}%
\pgfpathlineto{\pgfqpoint{4.577040in}{2.427573in}}%
\pgfpathlineto{\pgfqpoint{4.579322in}{2.404158in}}%
\pgfpathlineto{\pgfqpoint{4.581604in}{2.394792in}}%
\pgfpathlineto{\pgfqpoint{4.583887in}{2.394792in}}%
\pgfpathlineto{\pgfqpoint{4.586169in}{2.441622in}}%
\pgfpathlineto{\pgfqpoint{4.588451in}{2.408841in}}%
\pgfpathlineto{\pgfqpoint{4.590733in}{2.441622in}}%
\pgfpathlineto{\pgfqpoint{4.593016in}{2.408841in}}%
\pgfpathlineto{\pgfqpoint{4.595298in}{2.441622in}}%
\pgfpathlineto{\pgfqpoint{4.597580in}{2.427573in}}%
\pgfpathlineto{\pgfqpoint{4.599862in}{2.390109in}}%
\pgfpathlineto{\pgfqpoint{4.602145in}{2.404158in}}%
\pgfpathlineto{\pgfqpoint{4.604427in}{2.446304in}}%
\pgfpathlineto{\pgfqpoint{4.608991in}{2.380743in}}%
\pgfpathlineto{\pgfqpoint{4.611274in}{2.394792in}}%
\pgfpathlineto{\pgfqpoint{4.613556in}{2.455670in}}%
\pgfpathlineto{\pgfqpoint{4.618120in}{2.362012in}}%
\pgfpathlineto{\pgfqpoint{4.620403in}{2.418207in}}%
\pgfpathlineto{\pgfqpoint{4.622685in}{2.390109in}}%
\pgfpathlineto{\pgfqpoint{4.624967in}{2.441622in}}%
\pgfpathlineto{\pgfqpoint{4.627249in}{2.441622in}}%
\pgfpathlineto{\pgfqpoint{4.629532in}{2.376060in}}%
\pgfpathlineto{\pgfqpoint{4.631814in}{2.427573in}}%
\pgfpathlineto{\pgfqpoint{4.634096in}{2.394792in}}%
\pgfpathlineto{\pgfqpoint{4.636378in}{2.394792in}}%
\pgfpathlineto{\pgfqpoint{4.640943in}{2.450987in}}%
\pgfpathlineto{\pgfqpoint{4.643225in}{2.380743in}}%
\pgfpathlineto{\pgfqpoint{4.645507in}{2.390109in}}%
\pgfpathlineto{\pgfqpoint{4.647790in}{2.427573in}}%
\pgfpathlineto{\pgfqpoint{4.650072in}{2.394792in}}%
\pgfpathlineto{\pgfqpoint{4.652354in}{2.464744in}}%
\pgfpathlineto{\pgfqpoint{4.656919in}{2.361719in}}%
\pgfpathlineto{\pgfqpoint{4.659201in}{2.422597in}}%
\pgfpathlineto{\pgfqpoint{4.661483in}{2.422597in}}%
\pgfpathlineto{\pgfqpoint{4.663765in}{2.474109in}}%
\pgfpathlineto{\pgfqpoint{4.666048in}{2.464744in}}%
\pgfpathlineto{\pgfqpoint{4.670612in}{2.417914in}}%
\pgfpathlineto{\pgfqpoint{4.672894in}{2.450695in}}%
\pgfpathlineto{\pgfqpoint{4.675177in}{2.375768in}}%
\pgfpathlineto{\pgfqpoint{4.677459in}{2.441329in}}%
\pgfpathlineto{\pgfqpoint{4.679741in}{2.417914in}}%
\pgfpathlineto{\pgfqpoint{4.684306in}{2.431963in}}%
\pgfpathlineto{\pgfqpoint{4.686588in}{2.422597in}}%
\pgfpathlineto{\pgfqpoint{4.688870in}{2.394499in}}%
\pgfpathlineto{\pgfqpoint{4.691152in}{2.417914in}}%
\pgfpathlineto{\pgfqpoint{4.693435in}{2.389816in}}%
\pgfpathlineto{\pgfqpoint{4.695717in}{2.417914in}}%
\pgfpathlineto{\pgfqpoint{4.697999in}{2.389816in}}%
\pgfpathlineto{\pgfqpoint{4.700281in}{2.389816in}}%
\pgfpathlineto{\pgfqpoint{4.702564in}{2.403865in}}%
\pgfpathlineto{\pgfqpoint{4.704846in}{2.450695in}}%
\pgfpathlineto{\pgfqpoint{4.711693in}{2.361719in}}%
\pgfpathlineto{\pgfqpoint{4.713975in}{2.394499in}}%
\pgfpathlineto{\pgfqpoint{4.716257in}{2.497524in}}%
\pgfpathlineto{\pgfqpoint{4.720822in}{2.553719in}}%
\pgfpathlineto{\pgfqpoint{4.723104in}{2.431963in}}%
\pgfpathlineto{\pgfqpoint{4.725386in}{2.413231in}}%
\pgfpathlineto{\pgfqpoint{4.729951in}{2.399182in}}%
\pgfpathlineto{\pgfqpoint{4.732233in}{2.371085in}}%
\pgfpathlineto{\pgfqpoint{4.734515in}{2.431670in}}%
\pgfpathlineto{\pgfqpoint{4.736797in}{2.361426in}}%
\pgfpathlineto{\pgfqpoint{4.739080in}{2.412939in}}%
\pgfpathlineto{\pgfqpoint{4.741362in}{2.384841in}}%
\pgfpathlineto{\pgfqpoint{4.743644in}{2.384841in}}%
\pgfpathlineto{\pgfqpoint{4.748209in}{2.450402in}}%
\pgfpathlineto{\pgfqpoint{4.752773in}{2.384841in}}%
\pgfpathlineto{\pgfqpoint{4.755055in}{2.412939in}}%
\pgfpathlineto{\pgfqpoint{4.757338in}{2.394207in}}%
\pgfpathlineto{\pgfqpoint{4.759620in}{2.356743in}}%
\pgfpathlineto{\pgfqpoint{4.764184in}{2.398890in}}%
\pgfpathlineto{\pgfqpoint{4.766467in}{2.445719in}}%
\pgfpathlineto{\pgfqpoint{4.768749in}{2.426987in}}%
\pgfpathlineto{\pgfqpoint{4.771031in}{2.422304in}}%
\pgfpathlineto{\pgfqpoint{4.773313in}{2.389524in}}%
\pgfpathlineto{\pgfqpoint{4.775596in}{2.403573in}}%
\pgfpathlineto{\pgfqpoint{4.777878in}{2.389524in}}%
\pgfpathlineto{\pgfqpoint{4.780160in}{2.384841in}}%
\pgfpathlineto{\pgfqpoint{4.782443in}{2.445719in}}%
\pgfpathlineto{\pgfqpoint{4.784725in}{2.469134in}}%
\pgfpathlineto{\pgfqpoint{4.787007in}{2.455085in}}%
\pgfpathlineto{\pgfqpoint{4.791572in}{2.403573in}}%
\pgfpathlineto{\pgfqpoint{4.793854in}{2.478500in}}%
\pgfpathlineto{\pgfqpoint{4.796136in}{2.441036in}}%
\pgfpathlineto{\pgfqpoint{4.798418in}{2.436353in}}%
\pgfpathlineto{\pgfqpoint{4.800701in}{2.398890in}}%
\pgfpathlineto{\pgfqpoint{4.802983in}{2.441036in}}%
\pgfpathlineto{\pgfqpoint{4.805265in}{2.389524in}}%
\pgfpathlineto{\pgfqpoint{4.807547in}{2.445719in}}%
\pgfpathlineto{\pgfqpoint{4.812112in}{2.445719in}}%
\pgfpathlineto{\pgfqpoint{4.814394in}{2.482890in}}%
\pgfpathlineto{\pgfqpoint{4.816676in}{2.417329in}}%
\pgfpathlineto{\pgfqpoint{4.818959in}{2.431378in}}%
\pgfpathlineto{\pgfqpoint{4.821241in}{2.407963in}}%
\pgfpathlineto{\pgfqpoint{4.823523in}{2.478207in}}%
\pgfpathlineto{\pgfqpoint{4.825805in}{2.492256in}}%
\pgfpathlineto{\pgfqpoint{4.828088in}{2.496939in}}%
\pgfpathlineto{\pgfqpoint{4.830370in}{2.370499in}}%
\pgfpathlineto{\pgfqpoint{4.834934in}{2.436061in}}%
\pgfpathlineto{\pgfqpoint{4.837217in}{2.431378in}}%
\pgfpathlineto{\pgfqpoint{4.841781in}{2.370499in}}%
\pgfpathlineto{\pgfqpoint{4.844063in}{2.459475in}}%
\pgfpathlineto{\pgfqpoint{4.846346in}{2.431378in}}%
\pgfpathlineto{\pgfqpoint{4.848628in}{2.450109in}}%
\pgfpathlineto{\pgfqpoint{4.850910in}{2.431378in}}%
\pgfpathlineto{\pgfqpoint{4.853192in}{2.370499in}}%
\pgfpathlineto{\pgfqpoint{4.855475in}{2.426695in}}%
\pgfpathlineto{\pgfqpoint{4.857757in}{2.412646in}}%
\pgfpathlineto{\pgfqpoint{4.860039in}{2.445426in}}%
\pgfpathlineto{\pgfqpoint{4.862321in}{2.426695in}}%
\pgfpathlineto{\pgfqpoint{4.864604in}{2.501622in}}%
\pgfpathlineto{\pgfqpoint{4.866886in}{2.431378in}}%
\pgfpathlineto{\pgfqpoint{4.869168in}{2.417329in}}%
\pgfpathlineto{\pgfqpoint{4.871450in}{2.431378in}}%
\pgfpathlineto{\pgfqpoint{4.873733in}{2.412646in}}%
\pgfpathlineto{\pgfqpoint{4.876015in}{2.492256in}}%
\pgfpathlineto{\pgfqpoint{4.878297in}{2.417329in}}%
\pgfpathlineto{\pgfqpoint{4.880579in}{2.468841in}}%
\pgfpathlineto{\pgfqpoint{4.882862in}{2.464158in}}%
\pgfpathlineto{\pgfqpoint{4.885144in}{2.384548in}}%
\pgfpathlineto{\pgfqpoint{4.887426in}{2.375182in}}%
\pgfpathlineto{\pgfqpoint{4.889708in}{2.389231in}}%
\pgfpathlineto{\pgfqpoint{4.891991in}{2.384548in}}%
\pgfpathlineto{\pgfqpoint{4.894273in}{2.426402in}}%
\pgfpathlineto{\pgfqpoint{4.896555in}{2.417036in}}%
\pgfpathlineto{\pgfqpoint{4.898837in}{2.435768in}}%
\pgfpathlineto{\pgfqpoint{4.903402in}{2.384256in}}%
\pgfpathlineto{\pgfqpoint{4.905684in}{2.449817in}}%
\pgfpathlineto{\pgfqpoint{4.907966in}{2.370207in}}%
\pgfpathlineto{\pgfqpoint{4.910249in}{2.426402in}}%
\pgfpathlineto{\pgfqpoint{4.912531in}{2.426402in}}%
\pgfpathlineto{\pgfqpoint{4.914813in}{2.454500in}}%
\pgfpathlineto{\pgfqpoint{4.917095in}{2.374890in}}%
\pgfpathlineto{\pgfqpoint{4.919378in}{2.435768in}}%
\pgfpathlineto{\pgfqpoint{4.921660in}{2.440451in}}%
\pgfpathlineto{\pgfqpoint{4.923942in}{2.417036in}}%
\pgfpathlineto{\pgfqpoint{4.926224in}{2.435768in}}%
\pgfpathlineto{\pgfqpoint{4.928507in}{2.421719in}}%
\pgfpathlineto{\pgfqpoint{4.930789in}{2.463866in}}%
\pgfpathlineto{\pgfqpoint{4.933071in}{2.402987in}}%
\pgfpathlineto{\pgfqpoint{4.937636in}{2.529427in}}%
\pgfpathlineto{\pgfqpoint{4.942200in}{2.482597in}}%
\pgfpathlineto{\pgfqpoint{4.944482in}{2.566890in}}%
\pgfpathlineto{\pgfqpoint{4.946765in}{2.824452in}}%
\pgfpathlineto{\pgfqpoint{4.949047in}{2.932160in}}%
\pgfpathlineto{\pgfqpoint{4.951329in}{2.969623in}}%
\pgfpathlineto{\pgfqpoint{4.953611in}{2.960257in}}%
\pgfpathlineto{\pgfqpoint{4.955894in}{2.899379in}}%
\pgfpathlineto{\pgfqpoint{4.958176in}{2.749525in}}%
\pgfpathlineto{\pgfqpoint{4.960458in}{2.445134in}}%
\pgfpathlineto{\pgfqpoint{4.962740in}{2.501329in}}%
\pgfpathlineto{\pgfqpoint{4.965023in}{2.501329in}}%
\pgfpathlineto{\pgfqpoint{4.967305in}{2.529427in}}%
\pgfpathlineto{\pgfqpoint{4.969587in}{2.435768in}}%
\pgfpathlineto{\pgfqpoint{4.971869in}{2.515378in}}%
\pgfpathlineto{\pgfqpoint{4.974152in}{2.440158in}}%
\pgfpathlineto{\pgfqpoint{4.976434in}{2.416743in}}%
\pgfpathlineto{\pgfqpoint{4.980998in}{2.519768in}}%
\pgfpathlineto{\pgfqpoint{4.983281in}{2.533817in}}%
\pgfpathlineto{\pgfqpoint{4.987845in}{2.444841in}}%
\pgfpathlineto{\pgfqpoint{4.990127in}{2.454207in}}%
\pgfpathlineto{\pgfqpoint{4.992410in}{2.435475in}}%
\pgfpathlineto{\pgfqpoint{4.994692in}{2.402695in}}%
\pgfpathlineto{\pgfqpoint{4.996974in}{2.388646in}}%
\pgfpathlineto{\pgfqpoint{4.999256in}{2.444841in}}%
\pgfpathlineto{\pgfqpoint{5.001539in}{2.374597in}}%
\pgfpathlineto{\pgfqpoint{5.003821in}{2.463573in}}%
\pgfpathlineto{\pgfqpoint{5.006103in}{2.463573in}}%
\pgfpathlineto{\pgfqpoint{5.008385in}{2.449524in}}%
\pgfpathlineto{\pgfqpoint{5.010668in}{2.402695in}}%
\pgfpathlineto{\pgfqpoint{5.012950in}{2.477622in}}%
\pgfpathlineto{\pgfqpoint{5.015232in}{2.496353in}}%
\pgfpathlineto{\pgfqpoint{5.017514in}{2.524451in}}%
\pgfpathlineto{\pgfqpoint{5.019797in}{2.486988in}}%
\pgfpathlineto{\pgfqpoint{5.022079in}{2.501036in}}%
\pgfpathlineto{\pgfqpoint{5.024361in}{2.529134in}}%
\pgfpathlineto{\pgfqpoint{5.026643in}{2.505719in}}%
\pgfpathlineto{\pgfqpoint{5.028926in}{2.468256in}}%
\pgfpathlineto{\pgfqpoint{5.031208in}{2.575963in}}%
\pgfpathlineto{\pgfqpoint{5.033490in}{2.575963in}}%
\pgfpathlineto{\pgfqpoint{5.035772in}{2.519768in}}%
\pgfpathlineto{\pgfqpoint{5.038055in}{2.557232in}}%
\pgfpathlineto{\pgfqpoint{5.042619in}{2.599378in}}%
\pgfpathlineto{\pgfqpoint{5.047184in}{2.547866in}}%
\pgfpathlineto{\pgfqpoint{5.049466in}{2.529134in}}%
\pgfpathlineto{\pgfqpoint{5.051748in}{2.533817in}}%
\pgfpathlineto{\pgfqpoint{5.054030in}{2.561622in}}%
\pgfpathlineto{\pgfqpoint{5.056313in}{2.524158in}}%
\pgfpathlineto{\pgfqpoint{5.058595in}{2.575671in}}%
\pgfpathlineto{\pgfqpoint{5.060877in}{2.542890in}}%
\pgfpathlineto{\pgfqpoint{5.063159in}{2.538207in}}%
\pgfpathlineto{\pgfqpoint{5.065442in}{2.599085in}}%
\pgfpathlineto{\pgfqpoint{5.067724in}{2.561622in}}%
\pgfpathlineto{\pgfqpoint{5.070006in}{2.561622in}}%
\pgfpathlineto{\pgfqpoint{5.072289in}{2.655281in}}%
\pgfpathlineto{\pgfqpoint{5.076853in}{2.533524in}}%
\pgfpathlineto{\pgfqpoint{5.079135in}{2.524158in}}%
\pgfpathlineto{\pgfqpoint{5.081418in}{2.449231in}}%
\pgfpathlineto{\pgfqpoint{5.083700in}{2.467963in}}%
\pgfpathlineto{\pgfqpoint{5.088264in}{2.561622in}}%
\pgfpathlineto{\pgfqpoint{5.090547in}{2.561622in}}%
\pgfpathlineto{\pgfqpoint{5.092829in}{2.627183in}}%
\pgfpathlineto{\pgfqpoint{5.095111in}{2.613134in}}%
\pgfpathlineto{\pgfqpoint{5.097393in}{2.575671in}}%
\pgfpathlineto{\pgfqpoint{5.101958in}{2.622500in}}%
\pgfpathlineto{\pgfqpoint{5.106522in}{2.538207in}}%
\pgfpathlineto{\pgfqpoint{5.108805in}{2.538207in}}%
\pgfpathlineto{\pgfqpoint{5.111087in}{2.585037in}}%
\pgfpathlineto{\pgfqpoint{5.113369in}{2.575671in}}%
\pgfpathlineto{\pgfqpoint{5.115651in}{2.538207in}}%
\pgfpathlineto{\pgfqpoint{5.117934in}{2.556939in}}%
\pgfpathlineto{\pgfqpoint{5.120216in}{2.608451in}}%
\pgfpathlineto{\pgfqpoint{5.122498in}{2.510110in}}%
\pgfpathlineto{\pgfqpoint{5.124780in}{2.575671in}}%
\pgfpathlineto{\pgfqpoint{5.127063in}{2.585037in}}%
\pgfpathlineto{\pgfqpoint{5.129345in}{2.603768in}}%
\pgfpathlineto{\pgfqpoint{5.131627in}{2.542890in}}%
\pgfpathlineto{\pgfqpoint{5.133909in}{2.594110in}}%
\pgfpathlineto{\pgfqpoint{5.138474in}{2.626890in}}%
\pgfpathlineto{\pgfqpoint{5.140756in}{2.547280in}}%
\pgfpathlineto{\pgfqpoint{5.143038in}{2.612842in}}%
\pgfpathlineto{\pgfqpoint{5.145321in}{2.589427in}}%
\pgfpathlineto{\pgfqpoint{5.149885in}{2.650305in}}%
\pgfpathlineto{\pgfqpoint{5.154450in}{2.448939in}}%
\pgfpathlineto{\pgfqpoint{5.156732in}{2.594110in}}%
\pgfpathlineto{\pgfqpoint{5.159014in}{2.608159in}}%
\pgfpathlineto{\pgfqpoint{5.163579in}{2.575378in}}%
\pgfpathlineto{\pgfqpoint{5.165861in}{2.636256in}}%
\pgfpathlineto{\pgfqpoint{5.168143in}{2.561329in}}%
\pgfpathlineto{\pgfqpoint{5.170425in}{2.617525in}}%
\pgfpathlineto{\pgfqpoint{5.174990in}{2.589427in}}%
\pgfpathlineto{\pgfqpoint{5.177272in}{2.617525in}}%
\pgfpathlineto{\pgfqpoint{5.179554in}{2.533232in}}%
\pgfpathlineto{\pgfqpoint{5.181837in}{2.673720in}}%
\pgfpathlineto{\pgfqpoint{5.184119in}{2.608159in}}%
\pgfpathlineto{\pgfqpoint{5.186401in}{2.636256in}}%
\pgfpathlineto{\pgfqpoint{5.188683in}{2.537915in}}%
\pgfpathlineto{\pgfqpoint{5.190966in}{2.594110in}}%
\pgfpathlineto{\pgfqpoint{5.193248in}{2.575378in}}%
\pgfpathlineto{\pgfqpoint{5.195530in}{2.631573in}}%
\pgfpathlineto{\pgfqpoint{5.197812in}{2.584744in}}%
\pgfpathlineto{\pgfqpoint{5.200095in}{2.612842in}}%
\pgfpathlineto{\pgfqpoint{5.202377in}{2.608159in}}%
\pgfpathlineto{\pgfqpoint{5.204659in}{2.594110in}}%
\pgfpathlineto{\pgfqpoint{5.206941in}{2.533232in}}%
\pgfpathlineto{\pgfqpoint{5.209224in}{2.580061in}}%
\pgfpathlineto{\pgfqpoint{5.211506in}{2.561329in}}%
\pgfpathlineto{\pgfqpoint{5.213788in}{2.556354in}}%
\pgfpathlineto{\pgfqpoint{5.218353in}{2.607866in}}%
\pgfpathlineto{\pgfqpoint{5.220635in}{2.579768in}}%
\pgfpathlineto{\pgfqpoint{5.222917in}{2.593817in}}%
\pgfpathlineto{\pgfqpoint{5.227482in}{2.472061in}}%
\pgfpathlineto{\pgfqpoint{5.229764in}{2.462695in}}%
\pgfpathlineto{\pgfqpoint{5.232046in}{2.509524in}}%
\pgfpathlineto{\pgfqpoint{5.234328in}{2.500158in}}%
\pgfpathlineto{\pgfqpoint{5.236611in}{2.467378in}}%
\pgfpathlineto{\pgfqpoint{5.238893in}{2.369036in}}%
\pgfpathlineto{\pgfqpoint{5.243457in}{2.458012in}}%
\pgfpathlineto{\pgfqpoint{5.245740in}{2.490792in}}%
\pgfpathlineto{\pgfqpoint{5.248022in}{2.458012in}}%
\pgfpathlineto{\pgfqpoint{5.250304in}{2.481427in}}%
\pgfpathlineto{\pgfqpoint{5.252586in}{2.556354in}}%
\pgfpathlineto{\pgfqpoint{5.254869in}{2.453329in}}%
\pgfpathlineto{\pgfqpoint{5.259433in}{2.509524in}}%
\pgfpathlineto{\pgfqpoint{5.261715in}{2.462695in}}%
\pgfpathlineto{\pgfqpoint{5.263998in}{2.495475in}}%
\pgfpathlineto{\pgfqpoint{5.266280in}{2.490792in}}%
\pgfpathlineto{\pgfqpoint{5.268562in}{2.425231in}}%
\pgfpathlineto{\pgfqpoint{5.270844in}{2.429914in}}%
\pgfpathlineto{\pgfqpoint{5.273127in}{2.425231in}}%
\pgfpathlineto{\pgfqpoint{5.275409in}{2.415865in}}%
\pgfpathlineto{\pgfqpoint{5.277691in}{2.495475in}}%
\pgfpathlineto{\pgfqpoint{5.279973in}{2.425231in}}%
\pgfpathlineto{\pgfqpoint{5.282256in}{2.429914in}}%
\pgfpathlineto{\pgfqpoint{5.284538in}{2.411182in}}%
\pgfpathlineto{\pgfqpoint{5.286820in}{2.448646in}}%
\pgfpathlineto{\pgfqpoint{5.289102in}{2.420548in}}%
\pgfpathlineto{\pgfqpoint{5.291385in}{2.476744in}}%
\pgfpathlineto{\pgfqpoint{5.293667in}{2.467085in}}%
\pgfpathlineto{\pgfqpoint{5.295949in}{2.523280in}}%
\pgfpathlineto{\pgfqpoint{5.298231in}{2.429622in}}%
\pgfpathlineto{\pgfqpoint{5.300514in}{2.401524in}}%
\pgfpathlineto{\pgfqpoint{5.305078in}{2.490500in}}%
\pgfpathlineto{\pgfqpoint{5.307360in}{2.448353in}}%
\pgfpathlineto{\pgfqpoint{5.309643in}{2.443670in}}%
\pgfpathlineto{\pgfqpoint{5.311925in}{2.415573in}}%
\pgfpathlineto{\pgfqpoint{5.314207in}{2.424939in}}%
\pgfpathlineto{\pgfqpoint{5.316489in}{2.462402in}}%
\pgfpathlineto{\pgfqpoint{5.318772in}{2.434304in}}%
\pgfpathlineto{\pgfqpoint{5.321054in}{2.504549in}}%
\pgfpathlineto{\pgfqpoint{5.323336in}{2.438987in}}%
\pgfpathlineto{\pgfqpoint{5.325618in}{2.485817in}}%
\pgfpathlineto{\pgfqpoint{5.327901in}{2.481134in}}%
\pgfpathlineto{\pgfqpoint{5.330183in}{2.434304in}}%
\pgfpathlineto{\pgfqpoint{5.334747in}{2.396841in}}%
\pgfpathlineto{\pgfqpoint{5.337030in}{2.467085in}}%
\pgfpathlineto{\pgfqpoint{5.339312in}{2.462402in}}%
\pgfpathlineto{\pgfqpoint{5.341594in}{2.401524in}}%
\pgfpathlineto{\pgfqpoint{5.346159in}{2.532646in}}%
\pgfpathlineto{\pgfqpoint{5.348441in}{2.420256in}}%
\pgfpathlineto{\pgfqpoint{5.350723in}{2.438987in}}%
\pgfpathlineto{\pgfqpoint{5.353005in}{2.373426in}}%
\pgfpathlineto{\pgfqpoint{5.357570in}{2.537329in}}%
\pgfpathlineto{\pgfqpoint{5.359852in}{2.527963in}}%
\pgfpathlineto{\pgfqpoint{5.362134in}{2.448353in}}%
\pgfpathlineto{\pgfqpoint{5.366699in}{2.504549in}}%
\pgfpathlineto{\pgfqpoint{5.368981in}{2.513914in}}%
\pgfpathlineto{\pgfqpoint{5.371264in}{2.490500in}}%
\pgfpathlineto{\pgfqpoint{5.373546in}{2.518305in}}%
\pgfpathlineto{\pgfqpoint{5.380393in}{2.396548in}}%
\pgfpathlineto{\pgfqpoint{5.382675in}{2.410597in}}%
\pgfpathlineto{\pgfqpoint{5.384957in}{2.410597in}}%
\pgfpathlineto{\pgfqpoint{5.387239in}{2.354402in}}%
\pgfpathlineto{\pgfqpoint{5.389522in}{2.438695in}}%
\pgfpathlineto{\pgfqpoint{5.391804in}{2.424646in}}%
\pgfpathlineto{\pgfqpoint{5.394086in}{2.401231in}}%
\pgfpathlineto{\pgfqpoint{5.396368in}{2.448061in}}%
\pgfpathlineto{\pgfqpoint{5.398651in}{2.527671in}}%
\pgfpathlineto{\pgfqpoint{5.400933in}{2.396548in}}%
\pgfpathlineto{\pgfqpoint{5.403215in}{2.476158in}}%
\pgfpathlineto{\pgfqpoint{5.405497in}{2.471475in}}%
\pgfpathlineto{\pgfqpoint{5.410062in}{2.471475in}}%
\pgfpathlineto{\pgfqpoint{5.412344in}{2.752452in}}%
\pgfpathlineto{\pgfqpoint{5.414626in}{2.888257in}}%
\pgfpathlineto{\pgfqpoint{5.416909in}{2.944452in}}%
\pgfpathlineto{\pgfqpoint{5.419191in}{2.944452in}}%
\pgfpathlineto{\pgfqpoint{5.421473in}{2.902306in}}%
\pgfpathlineto{\pgfqpoint{5.423755in}{2.785232in}}%
\pgfpathlineto{\pgfqpoint{5.426038in}{2.466792in}}%
\pgfpathlineto{\pgfqpoint{5.428320in}{2.438695in}}%
\pgfpathlineto{\pgfqpoint{5.430602in}{2.480841in}}%
\pgfpathlineto{\pgfqpoint{5.432884in}{2.401231in}}%
\pgfpathlineto{\pgfqpoint{5.435167in}{2.434012in}}%
\pgfpathlineto{\pgfqpoint{5.437449in}{2.405914in}}%
\pgfpathlineto{\pgfqpoint{5.439731in}{2.443378in}}%
\pgfpathlineto{\pgfqpoint{5.442013in}{2.434012in}}%
\pgfpathlineto{\pgfqpoint{5.446578in}{2.499573in}}%
\pgfpathlineto{\pgfqpoint{5.448860in}{2.448061in}}%
\pgfpathlineto{\pgfqpoint{5.451142in}{2.429329in}}%
\pgfpathlineto{\pgfqpoint{5.453425in}{2.424353in}}%
\pgfpathlineto{\pgfqpoint{5.455707in}{2.400938in}}%
\pgfpathlineto{\pgfqpoint{5.457989in}{2.410304in}}%
\pgfpathlineto{\pgfqpoint{5.460271in}{2.372841in}}%
\pgfpathlineto{\pgfqpoint{5.462554in}{2.457134in}}%
\pgfpathlineto{\pgfqpoint{5.464836in}{2.429036in}}%
\pgfpathlineto{\pgfqpoint{5.467118in}{2.386890in}}%
\pgfpathlineto{\pgfqpoint{5.469400in}{2.410304in}}%
\pgfpathlineto{\pgfqpoint{5.473965in}{2.340060in}}%
\pgfpathlineto{\pgfqpoint{5.476247in}{2.382207in}}%
\pgfpathlineto{\pgfqpoint{5.478529in}{2.354109in}}%
\pgfpathlineto{\pgfqpoint{5.480812in}{2.372841in}}%
\pgfpathlineto{\pgfqpoint{5.483094in}{2.457134in}}%
\pgfpathlineto{\pgfqpoint{5.485376in}{2.438402in}}%
\pgfpathlineto{\pgfqpoint{5.487658in}{2.485231in}}%
\pgfpathlineto{\pgfqpoint{5.492223in}{2.433719in}}%
\pgfpathlineto{\pgfqpoint{5.494505in}{2.386890in}}%
\pgfpathlineto{\pgfqpoint{5.501352in}{2.489914in}}%
\pgfpathlineto{\pgfqpoint{5.503634in}{2.424353in}}%
\pgfpathlineto{\pgfqpoint{5.505916in}{2.433719in}}%
\pgfpathlineto{\pgfqpoint{5.508199in}{2.382207in}}%
\pgfpathlineto{\pgfqpoint{5.510481in}{2.405621in}}%
\pgfpathlineto{\pgfqpoint{5.512763in}{2.414987in}}%
\pgfpathlineto{\pgfqpoint{5.515045in}{2.396256in}}%
\pgfpathlineto{\pgfqpoint{5.517328in}{2.419670in}}%
\pgfpathlineto{\pgfqpoint{5.519610in}{2.424353in}}%
\pgfpathlineto{\pgfqpoint{5.521892in}{2.396256in}}%
\pgfpathlineto{\pgfqpoint{5.524174in}{2.386890in}}%
\pgfpathlineto{\pgfqpoint{5.526457in}{2.358792in}}%
\pgfpathlineto{\pgfqpoint{5.528739in}{2.354109in}}%
\pgfpathlineto{\pgfqpoint{5.531021in}{2.354109in}}%
\pgfpathlineto{\pgfqpoint{5.535586in}{2.452158in}}%
\pgfpathlineto{\pgfqpoint{5.537868in}{2.386597in}}%
\pgfpathlineto{\pgfqpoint{5.540150in}{2.419378in}}%
\pgfpathlineto{\pgfqpoint{5.542432in}{2.433426in}}%
\pgfpathlineto{\pgfqpoint{5.546997in}{2.400646in}}%
\pgfpathlineto{\pgfqpoint{5.549279in}{2.358499in}}%
\pgfpathlineto{\pgfqpoint{5.551561in}{2.358499in}}%
\pgfpathlineto{\pgfqpoint{5.553844in}{2.410012in}}%
\pgfpathlineto{\pgfqpoint{5.556126in}{2.339768in}}%
\pgfpathlineto{\pgfqpoint{5.560690in}{2.428743in}}%
\pgfpathlineto{\pgfqpoint{5.562973in}{2.410012in}}%
\pgfpathlineto{\pgfqpoint{5.565255in}{2.400646in}}%
\pgfpathlineto{\pgfqpoint{5.567537in}{2.442792in}}%
\pgfpathlineto{\pgfqpoint{5.569819in}{2.367865in}}%
\pgfpathlineto{\pgfqpoint{5.572102in}{2.339768in}}%
\pgfpathlineto{\pgfqpoint{5.574384in}{2.410012in}}%
\pgfpathlineto{\pgfqpoint{5.576666in}{2.395963in}}%
\pgfpathlineto{\pgfqpoint{5.578948in}{2.353816in}}%
\pgfpathlineto{\pgfqpoint{5.581231in}{2.410012in}}%
\pgfpathlineto{\pgfqpoint{5.585795in}{2.353816in}}%
\pgfpathlineto{\pgfqpoint{5.588077in}{2.414695in}}%
\pgfpathlineto{\pgfqpoint{5.590360in}{2.353816in}}%
\pgfpathlineto{\pgfqpoint{5.592642in}{2.419378in}}%
\pgfpathlineto{\pgfqpoint{5.597206in}{2.311670in}}%
\pgfpathlineto{\pgfqpoint{5.599489in}{2.410012in}}%
\pgfpathlineto{\pgfqpoint{5.601771in}{2.400646in}}%
\pgfpathlineto{\pgfqpoint{5.604053in}{2.433426in}}%
\pgfpathlineto{\pgfqpoint{5.606335in}{2.414695in}}%
\pgfpathlineto{\pgfqpoint{5.608618in}{2.353816in}}%
\pgfpathlineto{\pgfqpoint{5.610900in}{2.419378in}}%
\pgfpathlineto{\pgfqpoint{5.613182in}{2.419085in}}%
\pgfpathlineto{\pgfqpoint{5.615464in}{2.461231in}}%
\pgfpathlineto{\pgfqpoint{5.617747in}{2.358207in}}%
\pgfpathlineto{\pgfqpoint{5.620029in}{2.372255in}}%
\pgfpathlineto{\pgfqpoint{5.622311in}{2.400353in}}%
\pgfpathlineto{\pgfqpoint{5.624593in}{2.362890in}}%
\pgfpathlineto{\pgfqpoint{5.626876in}{2.419085in}}%
\pgfpathlineto{\pgfqpoint{5.629158in}{2.419085in}}%
\pgfpathlineto{\pgfqpoint{5.631440in}{2.395670in}}%
\pgfpathlineto{\pgfqpoint{5.633722in}{2.428451in}}%
\pgfpathlineto{\pgfqpoint{5.636005in}{2.334792in}}%
\pgfpathlineto{\pgfqpoint{5.638287in}{2.386304in}}%
\pgfpathlineto{\pgfqpoint{5.640569in}{2.367573in}}%
\pgfpathlineto{\pgfqpoint{5.642851in}{2.442500in}}%
\pgfpathlineto{\pgfqpoint{5.645134in}{2.442500in}}%
\pgfpathlineto{\pgfqpoint{5.651980in}{2.559573in}}%
\pgfpathlineto{\pgfqpoint{5.654263in}{2.512744in}}%
\pgfpathlineto{\pgfqpoint{5.656545in}{2.395670in}}%
\pgfpathlineto{\pgfqpoint{5.658827in}{2.400353in}}%
\pgfpathlineto{\pgfqpoint{5.661110in}{2.344158in}}%
\pgfpathlineto{\pgfqpoint{5.663392in}{2.353524in}}%
\pgfpathlineto{\pgfqpoint{5.667956in}{2.419085in}}%
\pgfpathlineto{\pgfqpoint{5.670239in}{2.390987in}}%
\pgfpathlineto{\pgfqpoint{5.672521in}{2.386304in}}%
\pgfpathlineto{\pgfqpoint{5.674803in}{2.325426in}}%
\pgfpathlineto{\pgfqpoint{5.681650in}{2.419085in}}%
\pgfpathlineto{\pgfqpoint{5.683932in}{2.419085in}}%
\pgfpathlineto{\pgfqpoint{5.686214in}{2.362890in}}%
\pgfpathlineto{\pgfqpoint{5.688497in}{2.405036in}}%
\pgfpathlineto{\pgfqpoint{5.690779in}{2.348841in}}%
\pgfpathlineto{\pgfqpoint{5.693061in}{2.357914in}}%
\pgfpathlineto{\pgfqpoint{5.695343in}{2.339182in}}%
\pgfpathlineto{\pgfqpoint{5.697626in}{2.395377in}}%
\pgfpathlineto{\pgfqpoint{5.699908in}{2.418792in}}%
\pgfpathlineto{\pgfqpoint{5.702190in}{2.418792in}}%
\pgfpathlineto{\pgfqpoint{5.704472in}{2.404743in}}%
\pgfpathlineto{\pgfqpoint{5.706755in}{2.414109in}}%
\pgfpathlineto{\pgfqpoint{5.709037in}{2.339182in}}%
\pgfpathlineto{\pgfqpoint{5.713601in}{2.400060in}}%
\pgfpathlineto{\pgfqpoint{5.715884in}{2.386012in}}%
\pgfpathlineto{\pgfqpoint{5.718166in}{2.400060in}}%
\pgfpathlineto{\pgfqpoint{5.720448in}{2.390695in}}%
\pgfpathlineto{\pgfqpoint{5.722730in}{2.395377in}}%
\pgfpathlineto{\pgfqpoint{5.725013in}{2.357914in}}%
\pgfpathlineto{\pgfqpoint{5.727295in}{2.432841in}}%
\pgfpathlineto{\pgfqpoint{5.729577in}{2.395377in}}%
\pgfpathlineto{\pgfqpoint{5.731859in}{2.465622in}}%
\pgfpathlineto{\pgfqpoint{5.734142in}{2.339182in}}%
\pgfpathlineto{\pgfqpoint{5.736424in}{2.353231in}}%
\pgfpathlineto{\pgfqpoint{5.738706in}{2.329816in}}%
\pgfpathlineto{\pgfqpoint{5.740988in}{2.381329in}}%
\pgfpathlineto{\pgfqpoint{5.743271in}{2.329816in}}%
\pgfpathlineto{\pgfqpoint{5.745553in}{2.423475in}}%
\pgfpathlineto{\pgfqpoint{5.750117in}{2.362597in}}%
\pgfpathlineto{\pgfqpoint{5.752400in}{2.362597in}}%
\pgfpathlineto{\pgfqpoint{5.754682in}{2.325133in}}%
\pgfpathlineto{\pgfqpoint{5.756964in}{2.386012in}}%
\pgfpathlineto{\pgfqpoint{5.759246in}{2.353231in}}%
\pgfpathlineto{\pgfqpoint{5.761529in}{2.367280in}}%
\pgfpathlineto{\pgfqpoint{5.763811in}{2.390695in}}%
\pgfpathlineto{\pgfqpoint{5.766093in}{2.390695in}}%
\pgfpathlineto{\pgfqpoint{5.768375in}{2.395377in}}%
\pgfpathlineto{\pgfqpoint{5.770658in}{2.339182in}}%
\pgfpathlineto{\pgfqpoint{5.772940in}{2.357621in}}%
\pgfpathlineto{\pgfqpoint{5.775222in}{2.338890in}}%
\pgfpathlineto{\pgfqpoint{5.777504in}{2.371670in}}%
\pgfpathlineto{\pgfqpoint{5.779787in}{2.385719in}}%
\pgfpathlineto{\pgfqpoint{5.782069in}{2.385719in}}%
\pgfpathlineto{\pgfqpoint{5.784351in}{2.338890in}}%
\pgfpathlineto{\pgfqpoint{5.786633in}{2.381036in}}%
\pgfpathlineto{\pgfqpoint{5.788916in}{2.366987in}}%
\pgfpathlineto{\pgfqpoint{5.791198in}{2.399768in}}%
\pgfpathlineto{\pgfqpoint{5.793480in}{2.357621in}}%
\pgfpathlineto{\pgfqpoint{5.795762in}{2.352938in}}%
\pgfpathlineto{\pgfqpoint{5.798045in}{2.376353in}}%
\pgfpathlineto{\pgfqpoint{5.800327in}{2.413817in}}%
\pgfpathlineto{\pgfqpoint{5.802609in}{2.423182in}}%
\pgfpathlineto{\pgfqpoint{5.807174in}{2.338890in}}%
\pgfpathlineto{\pgfqpoint{5.809456in}{2.352938in}}%
\pgfpathlineto{\pgfqpoint{5.811738in}{2.385719in}}%
\pgfpathlineto{\pgfqpoint{5.814020in}{2.366987in}}%
\pgfpathlineto{\pgfqpoint{5.816303in}{2.413817in}}%
\pgfpathlineto{\pgfqpoint{5.818585in}{2.390402in}}%
\pgfpathlineto{\pgfqpoint{5.820867in}{2.409134in}}%
\pgfpathlineto{\pgfqpoint{5.823149in}{2.399768in}}%
\pgfpathlineto{\pgfqpoint{5.825432in}{2.366987in}}%
\pgfpathlineto{\pgfqpoint{5.827714in}{2.371670in}}%
\pgfpathlineto{\pgfqpoint{5.829996in}{2.362304in}}%
\pgfpathlineto{\pgfqpoint{5.832278in}{2.334207in}}%
\pgfpathlineto{\pgfqpoint{5.834561in}{2.343572in}}%
\pgfpathlineto{\pgfqpoint{5.839125in}{2.381036in}}%
\pgfpathlineto{\pgfqpoint{5.841407in}{2.409134in}}%
\pgfpathlineto{\pgfqpoint{5.843690in}{2.357621in}}%
\pgfpathlineto{\pgfqpoint{5.845972in}{2.343572in}}%
\pgfpathlineto{\pgfqpoint{5.848254in}{2.366987in}}%
\pgfpathlineto{\pgfqpoint{5.850536in}{2.451280in}}%
\pgfpathlineto{\pgfqpoint{5.852819in}{2.329231in}}%
\pgfpathlineto{\pgfqpoint{5.855101in}{2.390109in}}%
\pgfpathlineto{\pgfqpoint{5.857383in}{2.399475in}}%
\pgfpathlineto{\pgfqpoint{5.859665in}{2.497817in}}%
\pgfpathlineto{\pgfqpoint{5.861948in}{2.329231in}}%
\pgfpathlineto{\pgfqpoint{5.866512in}{2.465036in}}%
\pgfpathlineto{\pgfqpoint{5.868794in}{2.446304in}}%
\pgfpathlineto{\pgfqpoint{5.871077in}{2.483768in}}%
\pgfpathlineto{\pgfqpoint{5.873359in}{2.432256in}}%
\pgfpathlineto{\pgfqpoint{5.882488in}{2.914599in}}%
\pgfpathlineto{\pgfqpoint{5.884770in}{2.933330in}}%
\pgfpathlineto{\pgfqpoint{5.887052in}{2.900550in}}%
\pgfpathlineto{\pgfqpoint{5.889335in}{2.811574in}}%
\pgfpathlineto{\pgfqpoint{5.893899in}{2.455670in}}%
\pgfpathlineto{\pgfqpoint{5.896181in}{2.446304in}}%
\pgfpathlineto{\pgfqpoint{5.900746in}{2.577427in}}%
\pgfpathlineto{\pgfqpoint{5.905310in}{2.399475in}}%
\pgfpathlineto{\pgfqpoint{5.907593in}{2.371377in}}%
\pgfpathlineto{\pgfqpoint{5.909875in}{2.427573in}}%
\pgfpathlineto{\pgfqpoint{5.912157in}{2.446304in}}%
\pgfpathlineto{\pgfqpoint{5.914439in}{2.535280in}}%
\pgfpathlineto{\pgfqpoint{5.919004in}{2.366694in}}%
\pgfpathlineto{\pgfqpoint{5.921286in}{2.390109in}}%
\pgfpathlineto{\pgfqpoint{5.923568in}{2.338597in}}%
\pgfpathlineto{\pgfqpoint{5.925851in}{2.371377in}}%
\pgfpathlineto{\pgfqpoint{5.928133in}{2.352646in}}%
\pgfpathlineto{\pgfqpoint{5.930415in}{2.366694in}}%
\pgfpathlineto{\pgfqpoint{5.932697in}{2.352353in}}%
\pgfpathlineto{\pgfqpoint{5.934980in}{2.352353in}}%
\pgfpathlineto{\pgfqpoint{5.937262in}{2.366402in}}%
\pgfpathlineto{\pgfqpoint{5.939544in}{2.357036in}}%
\pgfpathlineto{\pgfqpoint{5.941826in}{2.371085in}}%
\pgfpathlineto{\pgfqpoint{5.944109in}{2.352353in}}%
\pgfpathlineto{\pgfqpoint{5.946391in}{2.342987in}}%
\pgfpathlineto{\pgfqpoint{5.948673in}{2.380451in}}%
\pgfpathlineto{\pgfqpoint{5.950955in}{2.371085in}}%
\pgfpathlineto{\pgfqpoint{5.953238in}{2.352353in}}%
\pgfpathlineto{\pgfqpoint{5.955520in}{2.399182in}}%
\pgfpathlineto{\pgfqpoint{5.957802in}{2.324255in}}%
\pgfpathlineto{\pgfqpoint{5.960085in}{2.413231in}}%
\pgfpathlineto{\pgfqpoint{5.962367in}{2.352353in}}%
\pgfpathlineto{\pgfqpoint{5.964649in}{2.352353in}}%
\pgfpathlineto{\pgfqpoint{5.966931in}{2.366402in}}%
\pgfpathlineto{\pgfqpoint{5.969214in}{2.408548in}}%
\pgfpathlineto{\pgfqpoint{5.971496in}{2.389816in}}%
\pgfpathlineto{\pgfqpoint{5.973778in}{2.380451in}}%
\pgfpathlineto{\pgfqpoint{5.976060in}{2.361719in}}%
\pgfpathlineto{\pgfqpoint{5.978343in}{2.371085in}}%
\pgfpathlineto{\pgfqpoint{5.980625in}{2.394499in}}%
\pgfpathlineto{\pgfqpoint{5.982907in}{2.338304in}}%
\pgfpathlineto{\pgfqpoint{5.987472in}{2.375768in}}%
\pgfpathlineto{\pgfqpoint{5.989754in}{2.347670in}}%
\pgfpathlineto{\pgfqpoint{5.994318in}{2.431963in}}%
\pgfpathlineto{\pgfqpoint{5.996601in}{2.357036in}}%
\pgfpathlineto{\pgfqpoint{5.998883in}{2.375768in}}%
\pgfpathlineto{\pgfqpoint{6.001165in}{2.371085in}}%
\pgfpathlineto{\pgfqpoint{6.003447in}{2.389816in}}%
\pgfpathlineto{\pgfqpoint{6.005730in}{2.324255in}}%
\pgfpathlineto{\pgfqpoint{6.010294in}{2.389816in}}%
\pgfpathlineto{\pgfqpoint{6.012576in}{2.295865in}}%
\pgfpathlineto{\pgfqpoint{6.017141in}{2.375475in}}%
\pgfpathlineto{\pgfqpoint{6.019423in}{2.328646in}}%
\pgfpathlineto{\pgfqpoint{6.021705in}{2.389524in}}%
\pgfpathlineto{\pgfqpoint{6.023988in}{2.370792in}}%
\pgfpathlineto{\pgfqpoint{6.026270in}{2.370792in}}%
\pgfpathlineto{\pgfqpoint{6.028552in}{2.347377in}}%
\pgfpathlineto{\pgfqpoint{6.030834in}{2.295865in}}%
\pgfpathlineto{\pgfqpoint{6.033117in}{2.352060in}}%
\pgfpathlineto{\pgfqpoint{6.035399in}{2.338011in}}%
\pgfpathlineto{\pgfqpoint{6.037681in}{2.366109in}}%
\pgfpathlineto{\pgfqpoint{6.039963in}{2.356743in}}%
\pgfpathlineto{\pgfqpoint{6.042246in}{2.295865in}}%
\pgfpathlineto{\pgfqpoint{6.044528in}{2.384841in}}%
\pgfpathlineto{\pgfqpoint{6.046810in}{2.323963in}}%
\pgfpathlineto{\pgfqpoint{6.051375in}{2.370792in}}%
\pgfpathlineto{\pgfqpoint{6.053657in}{2.338011in}}%
\pgfpathlineto{\pgfqpoint{6.055939in}{2.328646in}}%
\pgfpathlineto{\pgfqpoint{6.062786in}{2.370792in}}%
\pgfpathlineto{\pgfqpoint{6.065068in}{2.356743in}}%
\pgfpathlineto{\pgfqpoint{6.067350in}{2.361426in}}%
\pgfpathlineto{\pgfqpoint{6.069633in}{2.338011in}}%
\pgfpathlineto{\pgfqpoint{6.071915in}{2.384841in}}%
\pgfpathlineto{\pgfqpoint{6.074197in}{2.394207in}}%
\pgfpathlineto{\pgfqpoint{6.076479in}{2.412939in}}%
\pgfpathlineto{\pgfqpoint{6.078762in}{2.338011in}}%
\pgfpathlineto{\pgfqpoint{6.081044in}{2.366109in}}%
\pgfpathlineto{\pgfqpoint{6.083326in}{2.342694in}}%
\pgfpathlineto{\pgfqpoint{6.085608in}{2.352060in}}%
\pgfpathlineto{\pgfqpoint{6.087891in}{2.352060in}}%
\pgfpathlineto{\pgfqpoint{6.090173in}{2.366109in}}%
\pgfpathlineto{\pgfqpoint{6.092455in}{2.370499in}}%
\pgfpathlineto{\pgfqpoint{6.094737in}{2.398597in}}%
\pgfpathlineto{\pgfqpoint{6.097020in}{2.342402in}}%
\pgfpathlineto{\pgfqpoint{6.099302in}{2.389231in}}%
\pgfpathlineto{\pgfqpoint{6.101584in}{2.281523in}}%
\pgfpathlineto{\pgfqpoint{6.106149in}{2.370499in}}%
\pgfpathlineto{\pgfqpoint{6.108431in}{2.365816in}}%
\pgfpathlineto{\pgfqpoint{6.110713in}{2.328353in}}%
\pgfpathlineto{\pgfqpoint{6.115278in}{2.534402in}}%
\pgfpathlineto{\pgfqpoint{6.115278in}{2.534402in}}%
\pgfusepath{stroke}%
\end{pgfscope}%
\begin{pgfscope}%
\pgfsetrectcap%
\pgfsetmiterjoin%
\pgfsetlinewidth{0.803000pt}%
\definecolor{currentstroke}{rgb}{0.000000,0.000000,0.000000}%
\pgfsetstrokecolor{currentstroke}%
\pgfsetdash{}{0pt}%
\pgfpathmoveto{\pgfqpoint{3.833026in}{2.240123in}}%
\pgfpathlineto{\pgfqpoint{3.833026in}{3.150926in}}%
\pgfusepath{stroke}%
\end{pgfscope}%
\begin{pgfscope}%
\pgfsetrectcap%
\pgfsetmiterjoin%
\pgfsetlinewidth{0.803000pt}%
\definecolor{currentstroke}{rgb}{0.000000,0.000000,0.000000}%
\pgfsetstrokecolor{currentstroke}%
\pgfsetdash{}{0pt}%
\pgfpathmoveto{\pgfqpoint{6.115278in}{2.240123in}}%
\pgfpathlineto{\pgfqpoint{6.115278in}{3.150926in}}%
\pgfusepath{stroke}%
\end{pgfscope}%
\begin{pgfscope}%
\pgfsetrectcap%
\pgfsetmiterjoin%
\pgfsetlinewidth{0.803000pt}%
\definecolor{currentstroke}{rgb}{0.000000,0.000000,0.000000}%
\pgfsetstrokecolor{currentstroke}%
\pgfsetdash{}{0pt}%
\pgfpathmoveto{\pgfqpoint{3.833026in}{2.240123in}}%
\pgfpathlineto{\pgfqpoint{6.115278in}{2.240123in}}%
\pgfusepath{stroke}%
\end{pgfscope}%
\begin{pgfscope}%
\pgfsetrectcap%
\pgfsetmiterjoin%
\pgfsetlinewidth{0.803000pt}%
\definecolor{currentstroke}{rgb}{0.000000,0.000000,0.000000}%
\pgfsetstrokecolor{currentstroke}%
\pgfsetdash{}{0pt}%
\pgfpathmoveto{\pgfqpoint{3.833026in}{3.150926in}}%
\pgfpathlineto{\pgfqpoint{6.115278in}{3.150926in}}%
\pgfusepath{stroke}%
\end{pgfscope}%
\begin{pgfscope}%
\definecolor{textcolor}{rgb}{0.000000,0.000000,0.000000}%
\pgfsetstrokecolor{textcolor}%
\pgfsetfillcolor{textcolor}%
\pgftext[x=4.974152in,y=3.234260in,,base]{\color{textcolor}\rmfamily\fontsize{12.000000}{14.400000}\selectfont Rectangular Wave}%
\end{pgfscope}%
\begin{pgfscope}%
\pgfsetbuttcap%
\pgfsetmiterjoin%
\definecolor{currentfill}{rgb}{1.000000,1.000000,1.000000}%
\pgfsetfillcolor{currentfill}%
\pgfsetlinewidth{0.000000pt}%
\definecolor{currentstroke}{rgb}{0.000000,0.000000,0.000000}%
\pgfsetstrokecolor{currentstroke}%
\pgfsetstrokeopacity{0.000000}%
\pgfsetdash{}{0pt}%
\pgfpathmoveto{\pgfqpoint{0.758026in}{0.565123in}}%
\pgfpathlineto{\pgfqpoint{3.040278in}{0.565123in}}%
\pgfpathlineto{\pgfqpoint{3.040278in}{1.475926in}}%
\pgfpathlineto{\pgfqpoint{0.758026in}{1.475926in}}%
\pgfpathlineto{\pgfqpoint{0.758026in}{0.565123in}}%
\pgfpathclose%
\pgfusepath{fill}%
\end{pgfscope}%
\begin{pgfscope}%
\pgfpathrectangle{\pgfqpoint{0.758026in}{0.565123in}}{\pgfqpoint{2.282252in}{0.910803in}}%
\pgfusepath{clip}%
\pgfsetbuttcap%
\pgfsetroundjoin%
\pgfsetlinewidth{0.501875pt}%
\definecolor{currentstroke}{rgb}{0.690196,0.690196,0.690196}%
\pgfsetstrokecolor{currentstroke}%
\pgfsetstrokeopacity{0.700000}%
\pgfsetdash{{1.850000pt}{0.800000pt}}{0.000000pt}%
\pgfpathmoveto{\pgfqpoint{1.177215in}{0.565123in}}%
\pgfpathlineto{\pgfqpoint{1.177215in}{1.475926in}}%
\pgfusepath{stroke}%
\end{pgfscope}%
\begin{pgfscope}%
\pgfsetbuttcap%
\pgfsetroundjoin%
\definecolor{currentfill}{rgb}{0.000000,0.000000,0.000000}%
\pgfsetfillcolor{currentfill}%
\pgfsetlinewidth{0.803000pt}%
\definecolor{currentstroke}{rgb}{0.000000,0.000000,0.000000}%
\pgfsetstrokecolor{currentstroke}%
\pgfsetdash{}{0pt}%
\pgfsys@defobject{currentmarker}{\pgfqpoint{0.000000in}{-0.048611in}}{\pgfqpoint{0.000000in}{0.000000in}}{%
\pgfpathmoveto{\pgfqpoint{0.000000in}{0.000000in}}%
\pgfpathlineto{\pgfqpoint{0.000000in}{-0.048611in}}%
\pgfusepath{stroke,fill}%
}%
\begin{pgfscope}%
\pgfsys@transformshift{1.177215in}{0.565123in}%
\pgfsys@useobject{currentmarker}{}%
\end{pgfscope}%
\end{pgfscope}%
\begin{pgfscope}%
\definecolor{textcolor}{rgb}{0.000000,0.000000,0.000000}%
\pgfsetstrokecolor{textcolor}%
\pgfsetfillcolor{textcolor}%
\pgftext[x=1.177215in,y=0.467901in,,top]{\color{textcolor}\rmfamily\fontsize{10.000000}{12.000000}\selectfont \(\displaystyle {1}\)}%
\end{pgfscope}%
\begin{pgfscope}%
\pgfpathrectangle{\pgfqpoint{0.758026in}{0.565123in}}{\pgfqpoint{2.282252in}{0.910803in}}%
\pgfusepath{clip}%
\pgfsetbuttcap%
\pgfsetroundjoin%
\pgfsetlinewidth{0.501875pt}%
\definecolor{currentstroke}{rgb}{0.690196,0.690196,0.690196}%
\pgfsetstrokecolor{currentstroke}%
\pgfsetstrokeopacity{0.700000}%
\pgfsetdash{{1.850000pt}{0.800000pt}}{0.000000pt}%
\pgfpathmoveto{\pgfqpoint{1.642981in}{0.565123in}}%
\pgfpathlineto{\pgfqpoint{1.642981in}{1.475926in}}%
\pgfusepath{stroke}%
\end{pgfscope}%
\begin{pgfscope}%
\pgfsetbuttcap%
\pgfsetroundjoin%
\definecolor{currentfill}{rgb}{0.000000,0.000000,0.000000}%
\pgfsetfillcolor{currentfill}%
\pgfsetlinewidth{0.803000pt}%
\definecolor{currentstroke}{rgb}{0.000000,0.000000,0.000000}%
\pgfsetstrokecolor{currentstroke}%
\pgfsetdash{}{0pt}%
\pgfsys@defobject{currentmarker}{\pgfqpoint{0.000000in}{-0.048611in}}{\pgfqpoint{0.000000in}{0.000000in}}{%
\pgfpathmoveto{\pgfqpoint{0.000000in}{0.000000in}}%
\pgfpathlineto{\pgfqpoint{0.000000in}{-0.048611in}}%
\pgfusepath{stroke,fill}%
}%
\begin{pgfscope}%
\pgfsys@transformshift{1.642981in}{0.565123in}%
\pgfsys@useobject{currentmarker}{}%
\end{pgfscope}%
\end{pgfscope}%
\begin{pgfscope}%
\definecolor{textcolor}{rgb}{0.000000,0.000000,0.000000}%
\pgfsetstrokecolor{textcolor}%
\pgfsetfillcolor{textcolor}%
\pgftext[x=1.642981in,y=0.467901in,,top]{\color{textcolor}\rmfamily\fontsize{10.000000}{12.000000}\selectfont \(\displaystyle {2}\)}%
\end{pgfscope}%
\begin{pgfscope}%
\pgfpathrectangle{\pgfqpoint{0.758026in}{0.565123in}}{\pgfqpoint{2.282252in}{0.910803in}}%
\pgfusepath{clip}%
\pgfsetbuttcap%
\pgfsetroundjoin%
\pgfsetlinewidth{0.501875pt}%
\definecolor{currentstroke}{rgb}{0.690196,0.690196,0.690196}%
\pgfsetstrokecolor{currentstroke}%
\pgfsetstrokeopacity{0.700000}%
\pgfsetdash{{1.850000pt}{0.800000pt}}{0.000000pt}%
\pgfpathmoveto{\pgfqpoint{2.108746in}{0.565123in}}%
\pgfpathlineto{\pgfqpoint{2.108746in}{1.475926in}}%
\pgfusepath{stroke}%
\end{pgfscope}%
\begin{pgfscope}%
\pgfsetbuttcap%
\pgfsetroundjoin%
\definecolor{currentfill}{rgb}{0.000000,0.000000,0.000000}%
\pgfsetfillcolor{currentfill}%
\pgfsetlinewidth{0.803000pt}%
\definecolor{currentstroke}{rgb}{0.000000,0.000000,0.000000}%
\pgfsetstrokecolor{currentstroke}%
\pgfsetdash{}{0pt}%
\pgfsys@defobject{currentmarker}{\pgfqpoint{0.000000in}{-0.048611in}}{\pgfqpoint{0.000000in}{0.000000in}}{%
\pgfpathmoveto{\pgfqpoint{0.000000in}{0.000000in}}%
\pgfpathlineto{\pgfqpoint{0.000000in}{-0.048611in}}%
\pgfusepath{stroke,fill}%
}%
\begin{pgfscope}%
\pgfsys@transformshift{2.108746in}{0.565123in}%
\pgfsys@useobject{currentmarker}{}%
\end{pgfscope}%
\end{pgfscope}%
\begin{pgfscope}%
\definecolor{textcolor}{rgb}{0.000000,0.000000,0.000000}%
\pgfsetstrokecolor{textcolor}%
\pgfsetfillcolor{textcolor}%
\pgftext[x=2.108746in,y=0.467901in,,top]{\color{textcolor}\rmfamily\fontsize{10.000000}{12.000000}\selectfont \(\displaystyle {3}\)}%
\end{pgfscope}%
\begin{pgfscope}%
\pgfpathrectangle{\pgfqpoint{0.758026in}{0.565123in}}{\pgfqpoint{2.282252in}{0.910803in}}%
\pgfusepath{clip}%
\pgfsetbuttcap%
\pgfsetroundjoin%
\pgfsetlinewidth{0.501875pt}%
\definecolor{currentstroke}{rgb}{0.690196,0.690196,0.690196}%
\pgfsetstrokecolor{currentstroke}%
\pgfsetstrokeopacity{0.700000}%
\pgfsetdash{{1.850000pt}{0.800000pt}}{0.000000pt}%
\pgfpathmoveto{\pgfqpoint{2.574512in}{0.565123in}}%
\pgfpathlineto{\pgfqpoint{2.574512in}{1.475926in}}%
\pgfusepath{stroke}%
\end{pgfscope}%
\begin{pgfscope}%
\pgfsetbuttcap%
\pgfsetroundjoin%
\definecolor{currentfill}{rgb}{0.000000,0.000000,0.000000}%
\pgfsetfillcolor{currentfill}%
\pgfsetlinewidth{0.803000pt}%
\definecolor{currentstroke}{rgb}{0.000000,0.000000,0.000000}%
\pgfsetstrokecolor{currentstroke}%
\pgfsetdash{}{0pt}%
\pgfsys@defobject{currentmarker}{\pgfqpoint{0.000000in}{-0.048611in}}{\pgfqpoint{0.000000in}{0.000000in}}{%
\pgfpathmoveto{\pgfqpoint{0.000000in}{0.000000in}}%
\pgfpathlineto{\pgfqpoint{0.000000in}{-0.048611in}}%
\pgfusepath{stroke,fill}%
}%
\begin{pgfscope}%
\pgfsys@transformshift{2.574512in}{0.565123in}%
\pgfsys@useobject{currentmarker}{}%
\end{pgfscope}%
\end{pgfscope}%
\begin{pgfscope}%
\definecolor{textcolor}{rgb}{0.000000,0.000000,0.000000}%
\pgfsetstrokecolor{textcolor}%
\pgfsetfillcolor{textcolor}%
\pgftext[x=2.574512in,y=0.467901in,,top]{\color{textcolor}\rmfamily\fontsize{10.000000}{12.000000}\selectfont \(\displaystyle {4}\)}%
\end{pgfscope}%
\begin{pgfscope}%
\pgfpathrectangle{\pgfqpoint{0.758026in}{0.565123in}}{\pgfqpoint{2.282252in}{0.910803in}}%
\pgfusepath{clip}%
\pgfsetbuttcap%
\pgfsetroundjoin%
\pgfsetlinewidth{0.501875pt}%
\definecolor{currentstroke}{rgb}{0.690196,0.690196,0.690196}%
\pgfsetstrokecolor{currentstroke}%
\pgfsetstrokeopacity{0.700000}%
\pgfsetdash{{1.850000pt}{0.800000pt}}{0.000000pt}%
\pgfpathmoveto{\pgfqpoint{3.040278in}{0.565123in}}%
\pgfpathlineto{\pgfqpoint{3.040278in}{1.475926in}}%
\pgfusepath{stroke}%
\end{pgfscope}%
\begin{pgfscope}%
\pgfsetbuttcap%
\pgfsetroundjoin%
\definecolor{currentfill}{rgb}{0.000000,0.000000,0.000000}%
\pgfsetfillcolor{currentfill}%
\pgfsetlinewidth{0.803000pt}%
\definecolor{currentstroke}{rgb}{0.000000,0.000000,0.000000}%
\pgfsetstrokecolor{currentstroke}%
\pgfsetdash{}{0pt}%
\pgfsys@defobject{currentmarker}{\pgfqpoint{0.000000in}{-0.048611in}}{\pgfqpoint{0.000000in}{0.000000in}}{%
\pgfpathmoveto{\pgfqpoint{0.000000in}{0.000000in}}%
\pgfpathlineto{\pgfqpoint{0.000000in}{-0.048611in}}%
\pgfusepath{stroke,fill}%
}%
\begin{pgfscope}%
\pgfsys@transformshift{3.040278in}{0.565123in}%
\pgfsys@useobject{currentmarker}{}%
\end{pgfscope}%
\end{pgfscope}%
\begin{pgfscope}%
\definecolor{textcolor}{rgb}{0.000000,0.000000,0.000000}%
\pgfsetstrokecolor{textcolor}%
\pgfsetfillcolor{textcolor}%
\pgftext[x=3.040278in,y=0.467901in,,top]{\color{textcolor}\rmfamily\fontsize{10.000000}{12.000000}\selectfont \(\displaystyle {5}\)}%
\end{pgfscope}%
\begin{pgfscope}%
\definecolor{textcolor}{rgb}{0.000000,0.000000,0.000000}%
\pgfsetstrokecolor{textcolor}%
\pgfsetfillcolor{textcolor}%
\pgftext[x=1.899152in,y=0.288889in,,top]{\color{textcolor}\rmfamily\fontsize{10.000000}{12.000000}\selectfont Frequency (Hz)}%
\end{pgfscope}%
\begin{pgfscope}%
\definecolor{textcolor}{rgb}{0.000000,0.000000,0.000000}%
\pgfsetstrokecolor{textcolor}%
\pgfsetfillcolor{textcolor}%
\pgftext[x=3.040278in,y=0.302778in,right,top]{\color{textcolor}\rmfamily\fontsize{10.000000}{12.000000}\selectfont \(\displaystyle \times{10^{6}}{}\)}%
\end{pgfscope}%
\begin{pgfscope}%
\pgfpathrectangle{\pgfqpoint{0.758026in}{0.565123in}}{\pgfqpoint{2.282252in}{0.910803in}}%
\pgfusepath{clip}%
\pgfsetbuttcap%
\pgfsetroundjoin%
\pgfsetlinewidth{0.501875pt}%
\definecolor{currentstroke}{rgb}{0.690196,0.690196,0.690196}%
\pgfsetstrokecolor{currentstroke}%
\pgfsetstrokeopacity{0.700000}%
\pgfsetdash{{1.850000pt}{0.800000pt}}{0.000000pt}%
\pgfpathmoveto{\pgfqpoint{0.758026in}{0.736496in}}%
\pgfpathlineto{\pgfqpoint{3.040278in}{0.736496in}}%
\pgfusepath{stroke}%
\end{pgfscope}%
\begin{pgfscope}%
\pgfsetbuttcap%
\pgfsetroundjoin%
\definecolor{currentfill}{rgb}{0.000000,0.000000,0.000000}%
\pgfsetfillcolor{currentfill}%
\pgfsetlinewidth{0.803000pt}%
\definecolor{currentstroke}{rgb}{0.000000,0.000000,0.000000}%
\pgfsetstrokecolor{currentstroke}%
\pgfsetdash{}{0pt}%
\pgfsys@defobject{currentmarker}{\pgfqpoint{-0.048611in}{0.000000in}}{\pgfqpoint{-0.000000in}{0.000000in}}{%
\pgfpathmoveto{\pgfqpoint{-0.000000in}{0.000000in}}%
\pgfpathlineto{\pgfqpoint{-0.048611in}{0.000000in}}%
\pgfusepath{stroke,fill}%
}%
\begin{pgfscope}%
\pgfsys@transformshift{0.758026in}{0.736496in}%
\pgfsys@useobject{currentmarker}{}%
\end{pgfscope}%
\end{pgfscope}%
\begin{pgfscope}%
\definecolor{textcolor}{rgb}{0.000000,0.000000,0.000000}%
\pgfsetstrokecolor{textcolor}%
\pgfsetfillcolor{textcolor}%
\pgftext[x=0.344444in, y=0.688270in, left, base]{\color{textcolor}\rmfamily\fontsize{10.000000}{12.000000}\selectfont \(\displaystyle {\ensuremath{-}100}\)}%
\end{pgfscope}%
\begin{pgfscope}%
\pgfpathrectangle{\pgfqpoint{0.758026in}{0.565123in}}{\pgfqpoint{2.282252in}{0.910803in}}%
\pgfusepath{clip}%
\pgfsetbuttcap%
\pgfsetroundjoin%
\pgfsetlinewidth{0.501875pt}%
\definecolor{currentstroke}{rgb}{0.690196,0.690196,0.690196}%
\pgfsetstrokecolor{currentstroke}%
\pgfsetstrokeopacity{0.700000}%
\pgfsetdash{{1.850000pt}{0.800000pt}}{0.000000pt}%
\pgfpathmoveto{\pgfqpoint{0.758026in}{1.192538in}}%
\pgfpathlineto{\pgfqpoint{3.040278in}{1.192538in}}%
\pgfusepath{stroke}%
\end{pgfscope}%
\begin{pgfscope}%
\pgfsetbuttcap%
\pgfsetroundjoin%
\definecolor{currentfill}{rgb}{0.000000,0.000000,0.000000}%
\pgfsetfillcolor{currentfill}%
\pgfsetlinewidth{0.803000pt}%
\definecolor{currentstroke}{rgb}{0.000000,0.000000,0.000000}%
\pgfsetstrokecolor{currentstroke}%
\pgfsetdash{}{0pt}%
\pgfsys@defobject{currentmarker}{\pgfqpoint{-0.048611in}{0.000000in}}{\pgfqpoint{-0.000000in}{0.000000in}}{%
\pgfpathmoveto{\pgfqpoint{-0.000000in}{0.000000in}}%
\pgfpathlineto{\pgfqpoint{-0.048611in}{0.000000in}}%
\pgfusepath{stroke,fill}%
}%
\begin{pgfscope}%
\pgfsys@transformshift{0.758026in}{1.192538in}%
\pgfsys@useobject{currentmarker}{}%
\end{pgfscope}%
\end{pgfscope}%
\begin{pgfscope}%
\definecolor{textcolor}{rgb}{0.000000,0.000000,0.000000}%
\pgfsetstrokecolor{textcolor}%
\pgfsetfillcolor{textcolor}%
\pgftext[x=0.413889in, y=1.144313in, left, base]{\color{textcolor}\rmfamily\fontsize{10.000000}{12.000000}\selectfont \(\displaystyle {\ensuremath{-}50}\)}%
\end{pgfscope}%
\begin{pgfscope}%
\definecolor{textcolor}{rgb}{0.000000,0.000000,0.000000}%
\pgfsetstrokecolor{textcolor}%
\pgfsetfillcolor{textcolor}%
\pgftext[x=0.288889in,y=1.020525in,,bottom,rotate=90.000000]{\color{textcolor}\rmfamily\fontsize{10.000000}{12.000000}\selectfont Amplitude (dB)}%
\end{pgfscope}%
\begin{pgfscope}%
\pgfpathrectangle{\pgfqpoint{0.758026in}{0.565123in}}{\pgfqpoint{2.282252in}{0.910803in}}%
\pgfusepath{clip}%
\pgfsetrectcap%
\pgfsetroundjoin%
\pgfsetlinewidth{1.003750pt}%
\definecolor{currentstroke}{rgb}{0.000000,0.000000,0.000000}%
\pgfsetstrokecolor{currentstroke}%
\pgfsetdash{}{0pt}%
\pgfpathmoveto{\pgfqpoint{0.758026in}{0.957676in}}%
\pgfpathlineto{\pgfqpoint{0.760308in}{0.912072in}}%
\pgfpathlineto{\pgfqpoint{0.762590in}{0.911787in}}%
\pgfpathlineto{\pgfqpoint{0.764872in}{0.925468in}}%
\pgfpathlineto{\pgfqpoint{0.767155in}{0.957106in}}%
\pgfpathlineto{\pgfqpoint{0.771719in}{0.911217in}}%
\pgfpathlineto{\pgfqpoint{0.776284in}{0.969933in}}%
\pgfpathlineto{\pgfqpoint{0.778566in}{0.914922in}}%
\pgfpathlineto{\pgfqpoint{0.783130in}{0.942000in}}%
\pgfpathlineto{\pgfqpoint{0.787695in}{0.896111in}}%
\pgfpathlineto{\pgfqpoint{0.789977in}{0.909507in}}%
\pgfpathlineto{\pgfqpoint{0.792259in}{0.936584in}}%
\pgfpathlineto{\pgfqpoint{0.794542in}{0.913782in}}%
\pgfpathlineto{\pgfqpoint{0.796824in}{0.922618in}}%
\pgfpathlineto{\pgfqpoint{0.799106in}{0.958817in}}%
\pgfpathlineto{\pgfqpoint{0.803671in}{0.894685in}}%
\pgfpathlineto{\pgfqpoint{0.810517in}{0.939720in}}%
\pgfpathlineto{\pgfqpoint{0.812800in}{0.912072in}}%
\pgfpathlineto{\pgfqpoint{0.815082in}{0.962237in}}%
\pgfpathlineto{\pgfqpoint{0.817364in}{0.952831in}}%
\pgfpathlineto{\pgfqpoint{0.819646in}{0.893260in}}%
\pgfpathlineto{\pgfqpoint{0.821929in}{0.934304in}}%
\pgfpathlineto{\pgfqpoint{0.824211in}{0.947700in}}%
\pgfpathlineto{\pgfqpoint{0.826493in}{0.947415in}}%
\pgfpathlineto{\pgfqpoint{0.828776in}{0.924328in}}%
\pgfpathlineto{\pgfqpoint{0.831058in}{0.951691in}}%
\pgfpathlineto{\pgfqpoint{0.835622in}{0.919198in}}%
\pgfpathlineto{\pgfqpoint{0.837905in}{0.951121in}}%
\pgfpathlineto{\pgfqpoint{0.840187in}{0.959957in}}%
\pgfpathlineto{\pgfqpoint{0.842469in}{0.932309in}}%
\pgfpathlineto{\pgfqpoint{0.844751in}{0.987034in}}%
\pgfpathlineto{\pgfqpoint{0.847034in}{0.932024in}}%
\pgfpathlineto{\pgfqpoint{0.851598in}{0.995585in}}%
\pgfpathlineto{\pgfqpoint{0.853880in}{0.949696in}}%
\pgfpathlineto{\pgfqpoint{0.856163in}{0.935729in}}%
\pgfpathlineto{\pgfqpoint{0.858445in}{0.990454in}}%
\pgfpathlineto{\pgfqpoint{0.860727in}{0.903521in}}%
\pgfpathlineto{\pgfqpoint{0.863009in}{0.962522in}}%
\pgfpathlineto{\pgfqpoint{0.865292in}{0.948841in}}%
\pgfpathlineto{\pgfqpoint{0.867574in}{1.007841in}}%
\pgfpathlineto{\pgfqpoint{0.869856in}{0.902666in}}%
\pgfpathlineto{\pgfqpoint{0.872138in}{0.934304in}}%
\pgfpathlineto{\pgfqpoint{0.874421in}{0.947985in}}%
\pgfpathlineto{\pgfqpoint{0.876703in}{0.943140in}}%
\pgfpathlineto{\pgfqpoint{0.878985in}{0.974778in}}%
\pgfpathlineto{\pgfqpoint{0.881267in}{0.970218in}}%
\pgfpathlineto{\pgfqpoint{0.883550in}{0.901526in}}%
\pgfpathlineto{\pgfqpoint{0.888114in}{0.960527in}}%
\pgfpathlineto{\pgfqpoint{0.890396in}{0.964802in}}%
\pgfpathlineto{\pgfqpoint{0.892679in}{0.900671in}}%
\pgfpathlineto{\pgfqpoint{0.899525in}{0.941145in}}%
\pgfpathlineto{\pgfqpoint{0.901808in}{0.963947in}}%
\pgfpathlineto{\pgfqpoint{0.904090in}{0.899816in}}%
\pgfpathlineto{\pgfqpoint{0.906372in}{0.967937in}}%
\pgfpathlineto{\pgfqpoint{0.908654in}{0.917773in}}%
\pgfpathlineto{\pgfqpoint{0.910937in}{0.899246in}}%
\pgfpathlineto{\pgfqpoint{0.913219in}{0.971928in}}%
\pgfpathlineto{\pgfqpoint{0.915501in}{0.903521in}}%
\pgfpathlineto{\pgfqpoint{0.917783in}{0.884995in}}%
\pgfpathlineto{\pgfqpoint{0.922348in}{0.943995in}}%
\pgfpathlineto{\pgfqpoint{0.924630in}{0.911787in}}%
\pgfpathlineto{\pgfqpoint{0.926912in}{0.938865in}}%
\pgfpathlineto{\pgfqpoint{0.929195in}{0.925183in}}%
\pgfpathlineto{\pgfqpoint{0.931477in}{0.938580in}}%
\pgfpathlineto{\pgfqpoint{0.933759in}{0.965657in}}%
\pgfpathlineto{\pgfqpoint{0.936041in}{0.929174in}}%
\pgfpathlineto{\pgfqpoint{0.938324in}{1.225316in}}%
\pgfpathlineto{\pgfqpoint{0.940606in}{1.370965in}}%
\pgfpathlineto{\pgfqpoint{0.942888in}{1.425405in}}%
\pgfpathlineto{\pgfqpoint{0.945170in}{1.434526in}}%
\pgfpathlineto{\pgfqpoint{0.947453in}{1.402318in}}%
\pgfpathlineto{\pgfqpoint{0.949735in}{1.301704in}}%
\pgfpathlineto{\pgfqpoint{0.952017in}{1.041759in}}%
\pgfpathlineto{\pgfqpoint{0.954299in}{0.909222in}}%
\pgfpathlineto{\pgfqpoint{0.958864in}{1.063991in}}%
\pgfpathlineto{\pgfqpoint{0.961146in}{0.940575in}}%
\pgfpathlineto{\pgfqpoint{0.963428in}{0.894685in}}%
\pgfpathlineto{\pgfqpoint{0.965711in}{0.894685in}}%
\pgfpathlineto{\pgfqpoint{0.967993in}{0.880719in}}%
\pgfpathlineto{\pgfqpoint{0.970275in}{0.880434in}}%
\pgfpathlineto{\pgfqpoint{0.974840in}{0.907512in}}%
\pgfpathlineto{\pgfqpoint{0.977122in}{0.879864in}}%
\pgfpathlineto{\pgfqpoint{0.979404in}{0.879864in}}%
\pgfpathlineto{\pgfqpoint{0.981686in}{0.906942in}}%
\pgfpathlineto{\pgfqpoint{0.983969in}{0.879294in}}%
\pgfpathlineto{\pgfqpoint{1.006791in}{0.877299in}}%
\pgfpathlineto{\pgfqpoint{1.018202in}{0.876444in}}%
\pgfpathlineto{\pgfqpoint{1.031896in}{0.875304in}}%
\pgfpathlineto{\pgfqpoint{1.034178in}{0.820294in}}%
\pgfpathlineto{\pgfqpoint{1.036460in}{0.861052in}}%
\pgfpathlineto{\pgfqpoint{1.038743in}{0.833690in}}%
\pgfpathlineto{\pgfqpoint{1.041025in}{0.819723in}}%
\pgfpathlineto{\pgfqpoint{1.045589in}{0.819438in}}%
\pgfpathlineto{\pgfqpoint{1.047872in}{0.846516in}}%
\pgfpathlineto{\pgfqpoint{1.050154in}{0.819153in}}%
\pgfpathlineto{\pgfqpoint{1.059283in}{0.818298in}}%
\pgfpathlineto{\pgfqpoint{1.061565in}{0.877299in}}%
\pgfpathlineto{\pgfqpoint{1.063847in}{0.817728in}}%
\pgfpathlineto{\pgfqpoint{1.072976in}{0.817158in}}%
\pgfpathlineto{\pgfqpoint{1.075259in}{0.830554in}}%
\pgfpathlineto{\pgfqpoint{1.077541in}{0.816588in}}%
\pgfpathlineto{\pgfqpoint{1.079823in}{0.816588in}}%
\pgfpathlineto{\pgfqpoint{1.082105in}{0.829984in}}%
\pgfpathlineto{\pgfqpoint{1.084388in}{0.829699in}}%
\pgfpathlineto{\pgfqpoint{1.086670in}{0.843381in}}%
\pgfpathlineto{\pgfqpoint{1.088952in}{0.815733in}}%
\pgfpathlineto{\pgfqpoint{1.093517in}{0.815448in}}%
\pgfpathlineto{\pgfqpoint{1.095799in}{0.828844in}}%
\pgfpathlineto{\pgfqpoint{1.098081in}{0.828559in}}%
\pgfpathlineto{\pgfqpoint{1.100363in}{0.814593in}}%
\pgfpathlineto{\pgfqpoint{1.102646in}{0.814593in}}%
\pgfpathlineto{\pgfqpoint{1.104928in}{0.827989in}}%
\pgfpathlineto{\pgfqpoint{1.107210in}{0.827704in}}%
\pgfpathlineto{\pgfqpoint{1.114057in}{0.786090in}}%
\pgfpathlineto{\pgfqpoint{1.120904in}{0.831125in}}%
\pgfpathlineto{\pgfqpoint{1.123186in}{0.785520in}}%
\pgfpathlineto{\pgfqpoint{1.125468in}{0.812598in}}%
\pgfpathlineto{\pgfqpoint{1.130033in}{0.784950in}}%
\pgfpathlineto{\pgfqpoint{1.132315in}{0.834830in}}%
\pgfpathlineto{\pgfqpoint{1.134597in}{0.784380in}}%
\pgfpathlineto{\pgfqpoint{1.136880in}{0.811743in}}%
\pgfpathlineto{\pgfqpoint{1.141444in}{0.783810in}}%
\pgfpathlineto{\pgfqpoint{1.148291in}{0.783240in}}%
\pgfpathlineto{\pgfqpoint{1.150573in}{0.810318in}}%
\pgfpathlineto{\pgfqpoint{1.152855in}{0.782955in}}%
\pgfpathlineto{\pgfqpoint{1.155138in}{0.782670in}}%
\pgfpathlineto{\pgfqpoint{1.157420in}{0.796066in}}%
\pgfpathlineto{\pgfqpoint{1.159702in}{0.782385in}}%
\pgfpathlineto{\pgfqpoint{1.164267in}{0.781815in}}%
\pgfpathlineto{\pgfqpoint{1.166549in}{0.754452in}}%
\pgfpathlineto{\pgfqpoint{1.168831in}{0.795211in}}%
\pgfpathlineto{\pgfqpoint{1.171113in}{0.767564in}}%
\pgfpathlineto{\pgfqpoint{1.173396in}{0.794926in}}%
\pgfpathlineto{\pgfqpoint{1.175678in}{0.780960in}}%
\pgfpathlineto{\pgfqpoint{1.177960in}{0.808037in}}%
\pgfpathlineto{\pgfqpoint{1.180242in}{0.808037in}}%
\pgfpathlineto{\pgfqpoint{1.182525in}{0.753312in}}%
\pgfpathlineto{\pgfqpoint{1.184807in}{0.766994in}}%
\pgfpathlineto{\pgfqpoint{1.187089in}{0.739631in}}%
\pgfpathlineto{\pgfqpoint{1.189371in}{0.757873in}}%
\pgfpathlineto{\pgfqpoint{1.191654in}{0.725950in}}%
\pgfpathlineto{\pgfqpoint{1.193936in}{0.725950in}}%
\pgfpathlineto{\pgfqpoint{1.196218in}{0.739631in}}%
\pgfpathlineto{\pgfqpoint{1.198500in}{0.771554in}}%
\pgfpathlineto{\pgfqpoint{1.200783in}{0.739631in}}%
\pgfpathlineto{\pgfqpoint{1.203065in}{0.739631in}}%
\pgfpathlineto{\pgfqpoint{1.205347in}{0.780675in}}%
\pgfpathlineto{\pgfqpoint{1.207629in}{0.739631in}}%
\pgfpathlineto{\pgfqpoint{1.209912in}{0.725950in}}%
\pgfpathlineto{\pgfqpoint{1.212194in}{0.753312in}}%
\pgfpathlineto{\pgfqpoint{1.214476in}{0.744191in}}%
\pgfpathlineto{\pgfqpoint{1.216758in}{0.744191in}}%
\pgfpathlineto{\pgfqpoint{1.219041in}{0.798917in}}%
\pgfpathlineto{\pgfqpoint{1.221323in}{0.725950in}}%
\pgfpathlineto{\pgfqpoint{1.223605in}{0.753312in}}%
\pgfpathlineto{\pgfqpoint{1.230452in}{0.753312in}}%
\pgfpathlineto{\pgfqpoint{1.232734in}{0.725950in}}%
\pgfpathlineto{\pgfqpoint{1.250992in}{0.725950in}}%
\pgfpathlineto{\pgfqpoint{1.253274in}{0.739631in}}%
\pgfpathlineto{\pgfqpoint{1.257839in}{0.739631in}}%
\pgfpathlineto{\pgfqpoint{1.260121in}{0.725665in}}%
\pgfpathlineto{\pgfqpoint{1.266968in}{0.725665in}}%
\pgfpathlineto{\pgfqpoint{1.269250in}{0.739346in}}%
\pgfpathlineto{\pgfqpoint{1.271532in}{0.725665in}}%
\pgfpathlineto{\pgfqpoint{1.273815in}{0.739346in}}%
\pgfpathlineto{\pgfqpoint{1.276097in}{0.725665in}}%
\pgfpathlineto{\pgfqpoint{1.282944in}{0.725665in}}%
\pgfpathlineto{\pgfqpoint{1.287508in}{0.670940in}}%
\pgfpathlineto{\pgfqpoint{1.289790in}{0.670940in}}%
\pgfpathlineto{\pgfqpoint{1.292073in}{0.684621in}}%
\pgfpathlineto{\pgfqpoint{1.294355in}{0.670940in}}%
\pgfpathlineto{\pgfqpoint{1.296637in}{0.689181in}}%
\pgfpathlineto{\pgfqpoint{1.298919in}{0.670940in}}%
\pgfpathlineto{\pgfqpoint{1.301202in}{0.698302in}}%
\pgfpathlineto{\pgfqpoint{1.305766in}{0.670940in}}%
\pgfpathlineto{\pgfqpoint{1.308048in}{0.684621in}}%
\pgfpathlineto{\pgfqpoint{1.310331in}{0.670940in}}%
\pgfpathlineto{\pgfqpoint{1.317177in}{0.670940in}}%
\pgfpathlineto{\pgfqpoint{1.319460in}{0.684621in}}%
\pgfpathlineto{\pgfqpoint{1.321742in}{0.670940in}}%
\pgfpathlineto{\pgfqpoint{1.326306in}{0.670940in}}%
\pgfpathlineto{\pgfqpoint{1.328589in}{0.684621in}}%
\pgfpathlineto{\pgfqpoint{1.330871in}{0.684621in}}%
\pgfpathlineto{\pgfqpoint{1.333153in}{0.670940in}}%
\pgfpathlineto{\pgfqpoint{1.335435in}{0.721104in}}%
\pgfpathlineto{\pgfqpoint{1.337718in}{0.643292in}}%
\pgfpathlineto{\pgfqpoint{1.340000in}{0.684336in}}%
\pgfpathlineto{\pgfqpoint{1.342282in}{0.643292in}}%
\pgfpathlineto{\pgfqpoint{1.344564in}{0.656973in}}%
\pgfpathlineto{\pgfqpoint{1.346847in}{0.656973in}}%
\pgfpathlineto{\pgfqpoint{1.349129in}{0.670654in}}%
\pgfpathlineto{\pgfqpoint{1.353693in}{0.757303in}}%
\pgfpathlineto{\pgfqpoint{1.355976in}{0.656973in}}%
\pgfpathlineto{\pgfqpoint{1.358258in}{0.670654in}}%
\pgfpathlineto{\pgfqpoint{1.360540in}{0.661534in}}%
\pgfpathlineto{\pgfqpoint{1.362822in}{0.670654in}}%
\pgfpathlineto{\pgfqpoint{1.365105in}{0.670654in}}%
\pgfpathlineto{\pgfqpoint{1.369669in}{0.643292in}}%
\pgfpathlineto{\pgfqpoint{1.371951in}{0.643292in}}%
\pgfpathlineto{\pgfqpoint{1.374234in}{0.684336in}}%
\pgfpathlineto{\pgfqpoint{1.376516in}{0.656973in}}%
\pgfpathlineto{\pgfqpoint{1.378798in}{0.684336in}}%
\pgfpathlineto{\pgfqpoint{1.381080in}{0.643292in}}%
\pgfpathlineto{\pgfqpoint{1.385645in}{0.743621in}}%
\pgfpathlineto{\pgfqpoint{1.387927in}{0.688896in}}%
\pgfpathlineto{\pgfqpoint{1.392492in}{0.798346in}}%
\pgfpathlineto{\pgfqpoint{1.394774in}{0.793786in}}%
\pgfpathlineto{\pgfqpoint{1.397056in}{0.784665in}}%
\pgfpathlineto{\pgfqpoint{1.399338in}{0.757303in}}%
\pgfpathlineto{\pgfqpoint{1.401621in}{0.707138in}}%
\pgfpathlineto{\pgfqpoint{1.406185in}{1.163181in}}%
\pgfpathlineto{\pgfqpoint{1.408467in}{1.245268in}}%
\pgfpathlineto{\pgfqpoint{1.410750in}{1.268071in}}%
\pgfpathlineto{\pgfqpoint{1.413032in}{1.245268in}}%
\pgfpathlineto{\pgfqpoint{1.415314in}{1.167741in}}%
\pgfpathlineto{\pgfqpoint{1.419879in}{0.788941in}}%
\pgfpathlineto{\pgfqpoint{1.422161in}{0.784380in}}%
\pgfpathlineto{\pgfqpoint{1.424443in}{0.802622in}}%
\pgfpathlineto{\pgfqpoint{1.426726in}{0.747897in}}%
\pgfpathlineto{\pgfqpoint{1.429008in}{0.747897in}}%
\pgfpathlineto{\pgfqpoint{1.431290in}{0.702292in}}%
\pgfpathlineto{\pgfqpoint{1.433572in}{0.715974in}}%
\pgfpathlineto{\pgfqpoint{1.438137in}{0.779820in}}%
\pgfpathlineto{\pgfqpoint{1.440419in}{0.811743in}}%
\pgfpathlineto{\pgfqpoint{1.444984in}{0.670369in}}%
\pgfpathlineto{\pgfqpoint{1.447266in}{0.670369in}}%
\pgfpathlineto{\pgfqpoint{1.449548in}{0.684051in}}%
\pgfpathlineto{\pgfqpoint{1.451830in}{0.656688in}}%
\pgfpathlineto{\pgfqpoint{1.454113in}{0.656688in}}%
\pgfpathlineto{\pgfqpoint{1.456395in}{0.643007in}}%
\pgfpathlineto{\pgfqpoint{1.458677in}{0.670369in}}%
\pgfpathlineto{\pgfqpoint{1.460959in}{0.656688in}}%
\pgfpathlineto{\pgfqpoint{1.463242in}{0.670369in}}%
\pgfpathlineto{\pgfqpoint{1.465524in}{0.643007in}}%
\pgfpathlineto{\pgfqpoint{1.467806in}{0.697732in}}%
\pgfpathlineto{\pgfqpoint{1.470088in}{0.656688in}}%
\pgfpathlineto{\pgfqpoint{1.472371in}{0.702292in}}%
\pgfpathlineto{\pgfqpoint{1.474653in}{0.715974in}}%
\pgfpathlineto{\pgfqpoint{1.479217in}{0.643007in}}%
\pgfpathlineto{\pgfqpoint{1.481500in}{0.670369in}}%
\pgfpathlineto{\pgfqpoint{1.486064in}{0.670369in}}%
\pgfpathlineto{\pgfqpoint{1.488346in}{0.674930in}}%
\pgfpathlineto{\pgfqpoint{1.490629in}{0.702292in}}%
\pgfpathlineto{\pgfqpoint{1.497475in}{0.656688in}}%
\pgfpathlineto{\pgfqpoint{1.499758in}{0.656403in}}%
\pgfpathlineto{\pgfqpoint{1.502040in}{0.702007in}}%
\pgfpathlineto{\pgfqpoint{1.504322in}{0.702007in}}%
\pgfpathlineto{\pgfqpoint{1.511169in}{0.656403in}}%
\pgfpathlineto{\pgfqpoint{1.513451in}{0.656403in}}%
\pgfpathlineto{\pgfqpoint{1.515733in}{0.660964in}}%
\pgfpathlineto{\pgfqpoint{1.518016in}{0.683766in}}%
\pgfpathlineto{\pgfqpoint{1.520298in}{0.674645in}}%
\pgfpathlineto{\pgfqpoint{1.522580in}{0.674645in}}%
\pgfpathlineto{\pgfqpoint{1.524862in}{0.670084in}}%
\pgfpathlineto{\pgfqpoint{1.527145in}{0.670084in}}%
\pgfpathlineto{\pgfqpoint{1.529427in}{0.642722in}}%
\pgfpathlineto{\pgfqpoint{1.531709in}{0.674645in}}%
\pgfpathlineto{\pgfqpoint{1.533991in}{0.674645in}}%
\pgfpathlineto{\pgfqpoint{1.536274in}{0.683766in}}%
\pgfpathlineto{\pgfqpoint{1.538556in}{0.642722in}}%
\pgfpathlineto{\pgfqpoint{1.540838in}{0.670084in}}%
\pgfpathlineto{\pgfqpoint{1.543120in}{0.656403in}}%
\pgfpathlineto{\pgfqpoint{1.545403in}{0.702007in}}%
\pgfpathlineto{\pgfqpoint{1.549967in}{0.670084in}}%
\pgfpathlineto{\pgfqpoint{1.552249in}{0.683766in}}%
\pgfpathlineto{\pgfqpoint{1.554532in}{0.660964in}}%
\pgfpathlineto{\pgfqpoint{1.556814in}{0.683766in}}%
\pgfpathlineto{\pgfqpoint{1.559096in}{0.670084in}}%
\pgfpathlineto{\pgfqpoint{1.561378in}{0.670084in}}%
\pgfpathlineto{\pgfqpoint{1.563661in}{0.683766in}}%
\pgfpathlineto{\pgfqpoint{1.565943in}{0.688326in}}%
\pgfpathlineto{\pgfqpoint{1.568225in}{0.656403in}}%
\pgfpathlineto{\pgfqpoint{1.570507in}{0.702007in}}%
\pgfpathlineto{\pgfqpoint{1.572790in}{0.702007in}}%
\pgfpathlineto{\pgfqpoint{1.575072in}{0.656403in}}%
\pgfpathlineto{\pgfqpoint{1.577354in}{0.683481in}}%
\pgfpathlineto{\pgfqpoint{1.579636in}{0.688041in}}%
\pgfpathlineto{\pgfqpoint{1.581919in}{0.683481in}}%
\pgfpathlineto{\pgfqpoint{1.584201in}{0.692602in}}%
\pgfpathlineto{\pgfqpoint{1.586483in}{0.688041in}}%
\pgfpathlineto{\pgfqpoint{1.588765in}{0.742766in}}%
\pgfpathlineto{\pgfqpoint{1.591048in}{0.710843in}}%
\pgfpathlineto{\pgfqpoint{1.593330in}{0.656118in}}%
\pgfpathlineto{\pgfqpoint{1.595612in}{0.688041in}}%
\pgfpathlineto{\pgfqpoint{1.597894in}{0.669799in}}%
\pgfpathlineto{\pgfqpoint{1.600177in}{0.619635in}}%
\pgfpathlineto{\pgfqpoint{1.602459in}{0.683481in}}%
\pgfpathlineto{\pgfqpoint{1.604741in}{0.656118in}}%
\pgfpathlineto{\pgfqpoint{1.607023in}{0.683481in}}%
\pgfpathlineto{\pgfqpoint{1.609306in}{0.660679in}}%
\pgfpathlineto{\pgfqpoint{1.611588in}{0.697162in}}%
\pgfpathlineto{\pgfqpoint{1.613870in}{0.715404in}}%
\pgfpathlineto{\pgfqpoint{1.616152in}{0.701722in}}%
\pgfpathlineto{\pgfqpoint{1.618435in}{0.697162in}}%
\pgfpathlineto{\pgfqpoint{1.620717in}{0.642437in}}%
\pgfpathlineto{\pgfqpoint{1.622999in}{0.719964in}}%
\pgfpathlineto{\pgfqpoint{1.627564in}{0.633316in}}%
\pgfpathlineto{\pgfqpoint{1.629846in}{0.683481in}}%
\pgfpathlineto{\pgfqpoint{1.632128in}{0.651558in}}%
\pgfpathlineto{\pgfqpoint{1.634410in}{0.710843in}}%
\pgfpathlineto{\pgfqpoint{1.638975in}{0.678920in}}%
\pgfpathlineto{\pgfqpoint{1.643539in}{0.715404in}}%
\pgfpathlineto{\pgfqpoint{1.645822in}{0.724525in}}%
\pgfpathlineto{\pgfqpoint{1.650386in}{0.678920in}}%
\pgfpathlineto{\pgfqpoint{1.652668in}{0.665239in}}%
\pgfpathlineto{\pgfqpoint{1.659515in}{0.664954in}}%
\pgfpathlineto{\pgfqpoint{1.661797in}{0.678635in}}%
\pgfpathlineto{\pgfqpoint{1.664080in}{0.664954in}}%
\pgfpathlineto{\pgfqpoint{1.666362in}{0.664954in}}%
\pgfpathlineto{\pgfqpoint{1.668644in}{0.696877in}}%
\pgfpathlineto{\pgfqpoint{1.670926in}{0.651273in}}%
\pgfpathlineto{\pgfqpoint{1.673209in}{0.678635in}}%
\pgfpathlineto{\pgfqpoint{1.675491in}{0.664954in}}%
\pgfpathlineto{\pgfqpoint{1.677773in}{0.705998in}}%
\pgfpathlineto{\pgfqpoint{1.680055in}{0.696877in}}%
\pgfpathlineto{\pgfqpoint{1.682338in}{0.651273in}}%
\pgfpathlineto{\pgfqpoint{1.684620in}{0.637591in}}%
\pgfpathlineto{\pgfqpoint{1.686902in}{0.719679in}}%
\pgfpathlineto{\pgfqpoint{1.689184in}{0.715119in}}%
\pgfpathlineto{\pgfqpoint{1.691467in}{0.664954in}}%
\pgfpathlineto{\pgfqpoint{1.693749in}{0.669514in}}%
\pgfpathlineto{\pgfqpoint{1.698313in}{0.692317in}}%
\pgfpathlineto{\pgfqpoint{1.700596in}{0.724240in}}%
\pgfpathlineto{\pgfqpoint{1.702878in}{0.683196in}}%
\pgfpathlineto{\pgfqpoint{1.705160in}{0.715119in}}%
\pgfpathlineto{\pgfqpoint{1.707443in}{0.669514in}}%
\pgfpathlineto{\pgfqpoint{1.709725in}{0.701437in}}%
\pgfpathlineto{\pgfqpoint{1.712007in}{0.655833in}}%
\pgfpathlineto{\pgfqpoint{1.714289in}{0.696877in}}%
\pgfpathlineto{\pgfqpoint{1.716572in}{0.705998in}}%
\pgfpathlineto{\pgfqpoint{1.718854in}{0.651273in}}%
\pgfpathlineto{\pgfqpoint{1.721136in}{0.651273in}}%
\pgfpathlineto{\pgfqpoint{1.723418in}{0.696877in}}%
\pgfpathlineto{\pgfqpoint{1.725701in}{0.664954in}}%
\pgfpathlineto{\pgfqpoint{1.727983in}{0.678635in}}%
\pgfpathlineto{\pgfqpoint{1.730265in}{0.678635in}}%
\pgfpathlineto{\pgfqpoint{1.732547in}{0.669514in}}%
\pgfpathlineto{\pgfqpoint{1.734830in}{0.669514in}}%
\pgfpathlineto{\pgfqpoint{1.737112in}{0.728800in}}%
\pgfpathlineto{\pgfqpoint{1.739394in}{0.637306in}}%
\pgfpathlineto{\pgfqpoint{1.743959in}{0.719394in}}%
\pgfpathlineto{\pgfqpoint{1.746241in}{0.733075in}}%
\pgfpathlineto{\pgfqpoint{1.750805in}{0.650988in}}%
\pgfpathlineto{\pgfqpoint{1.753088in}{0.655548in}}%
\pgfpathlineto{\pgfqpoint{1.755370in}{0.678350in}}%
\pgfpathlineto{\pgfqpoint{1.757652in}{0.664669in}}%
\pgfpathlineto{\pgfqpoint{1.762217in}{0.664669in}}%
\pgfpathlineto{\pgfqpoint{1.764499in}{0.637306in}}%
\pgfpathlineto{\pgfqpoint{1.766781in}{0.723954in}}%
\pgfpathlineto{\pgfqpoint{1.769063in}{0.637306in}}%
\pgfpathlineto{\pgfqpoint{1.771346in}{0.682911in}}%
\pgfpathlineto{\pgfqpoint{1.773628in}{0.650988in}}%
\pgfpathlineto{\pgfqpoint{1.775910in}{0.664669in}}%
\pgfpathlineto{\pgfqpoint{1.778192in}{0.705713in}}%
\pgfpathlineto{\pgfqpoint{1.780475in}{0.664669in}}%
\pgfpathlineto{\pgfqpoint{1.782757in}{0.755877in}}%
\pgfpathlineto{\pgfqpoint{1.787321in}{0.655548in}}%
\pgfpathlineto{\pgfqpoint{1.789604in}{0.669229in}}%
\pgfpathlineto{\pgfqpoint{1.791886in}{0.650988in}}%
\pgfpathlineto{\pgfqpoint{1.794168in}{0.650988in}}%
\pgfpathlineto{\pgfqpoint{1.796450in}{0.682911in}}%
\pgfpathlineto{\pgfqpoint{1.798733in}{0.737636in}}%
\pgfpathlineto{\pgfqpoint{1.801015in}{0.650988in}}%
\pgfpathlineto{\pgfqpoint{1.805579in}{0.678350in}}%
\pgfpathlineto{\pgfqpoint{1.807862in}{0.705713in}}%
\pgfpathlineto{\pgfqpoint{1.810144in}{0.696592in}}%
\pgfpathlineto{\pgfqpoint{1.812426in}{0.769559in}}%
\pgfpathlineto{\pgfqpoint{1.814708in}{0.664669in}}%
\pgfpathlineto{\pgfqpoint{1.816991in}{0.710273in}}%
\pgfpathlineto{\pgfqpoint{1.819273in}{0.655263in}}%
\pgfpathlineto{\pgfqpoint{1.821555in}{0.650703in}}%
\pgfpathlineto{\pgfqpoint{1.823837in}{0.678065in}}%
\pgfpathlineto{\pgfqpoint{1.826120in}{0.637021in}}%
\pgfpathlineto{\pgfqpoint{1.828402in}{0.650703in}}%
\pgfpathlineto{\pgfqpoint{1.830684in}{0.637021in}}%
\pgfpathlineto{\pgfqpoint{1.832966in}{0.696307in}}%
\pgfpathlineto{\pgfqpoint{1.835249in}{0.691746in}}%
\pgfpathlineto{\pgfqpoint{1.837531in}{0.691746in}}%
\pgfpathlineto{\pgfqpoint{1.839813in}{0.723669in}}%
\pgfpathlineto{\pgfqpoint{1.842095in}{0.637021in}}%
\pgfpathlineto{\pgfqpoint{1.844378in}{0.723669in}}%
\pgfpathlineto{\pgfqpoint{1.846660in}{0.655263in}}%
\pgfpathlineto{\pgfqpoint{1.848942in}{0.696307in}}%
\pgfpathlineto{\pgfqpoint{1.851224in}{0.668944in}}%
\pgfpathlineto{\pgfqpoint{1.853507in}{0.696307in}}%
\pgfpathlineto{\pgfqpoint{1.855789in}{0.678065in}}%
\pgfpathlineto{\pgfqpoint{1.858071in}{0.668944in}}%
\pgfpathlineto{\pgfqpoint{1.860353in}{0.719109in}}%
\pgfpathlineto{\pgfqpoint{1.862636in}{0.664384in}}%
\pgfpathlineto{\pgfqpoint{1.864918in}{0.705428in}}%
\pgfpathlineto{\pgfqpoint{1.867200in}{0.709988in}}%
\pgfpathlineto{\pgfqpoint{1.869482in}{0.778395in}}%
\pgfpathlineto{\pgfqpoint{1.871765in}{1.042899in}}%
\pgfpathlineto{\pgfqpoint{1.874047in}{1.143229in}}%
\pgfpathlineto{\pgfqpoint{1.876329in}{1.179712in}}%
\pgfpathlineto{\pgfqpoint{1.878611in}{1.175152in}}%
\pgfpathlineto{\pgfqpoint{1.880894in}{1.111306in}}%
\pgfpathlineto{\pgfqpoint{1.883176in}{0.969933in}}%
\pgfpathlineto{\pgfqpoint{1.885458in}{0.705428in}}%
\pgfpathlineto{\pgfqpoint{1.887740in}{0.751032in}}%
\pgfpathlineto{\pgfqpoint{1.890023in}{0.664384in}}%
\pgfpathlineto{\pgfqpoint{1.892305in}{0.751032in}}%
\pgfpathlineto{\pgfqpoint{1.894587in}{0.696307in}}%
\pgfpathlineto{\pgfqpoint{1.896869in}{0.769274in}}%
\pgfpathlineto{\pgfqpoint{1.899152in}{0.668659in}}%
\pgfpathlineto{\pgfqpoint{1.901434in}{0.691461in}}%
\pgfpathlineto{\pgfqpoint{1.903716in}{0.746187in}}%
\pgfpathlineto{\pgfqpoint{1.905998in}{0.686901in}}%
\pgfpathlineto{\pgfqpoint{1.908281in}{0.755307in}}%
\pgfpathlineto{\pgfqpoint{1.910563in}{0.718824in}}%
\pgfpathlineto{\pgfqpoint{1.912845in}{0.709703in}}%
\pgfpathlineto{\pgfqpoint{1.915127in}{0.732505in}}%
\pgfpathlineto{\pgfqpoint{1.917410in}{0.682341in}}%
\pgfpathlineto{\pgfqpoint{1.919692in}{0.718824in}}%
\pgfpathlineto{\pgfqpoint{1.924256in}{0.673220in}}%
\pgfpathlineto{\pgfqpoint{1.926539in}{0.696022in}}%
\pgfpathlineto{\pgfqpoint{1.928821in}{0.682341in}}%
\pgfpathlineto{\pgfqpoint{1.931103in}{0.732505in}}%
\pgfpathlineto{\pgfqpoint{1.933385in}{0.686901in}}%
\pgfpathlineto{\pgfqpoint{1.935668in}{0.691461in}}%
\pgfpathlineto{\pgfqpoint{1.937950in}{0.741626in}}%
\pgfpathlineto{\pgfqpoint{1.940232in}{0.718824in}}%
\pgfpathlineto{\pgfqpoint{1.942514in}{0.732505in}}%
\pgfpathlineto{\pgfqpoint{1.944797in}{0.727945in}}%
\pgfpathlineto{\pgfqpoint{1.947079in}{0.705143in}}%
\pgfpathlineto{\pgfqpoint{1.951643in}{0.727945in}}%
\pgfpathlineto{\pgfqpoint{1.953926in}{0.700582in}}%
\pgfpathlineto{\pgfqpoint{1.956208in}{0.750747in}}%
\pgfpathlineto{\pgfqpoint{1.958490in}{0.737066in}}%
\pgfpathlineto{\pgfqpoint{1.960772in}{0.700582in}}%
\pgfpathlineto{\pgfqpoint{1.963055in}{0.746187in}}%
\pgfpathlineto{\pgfqpoint{1.965337in}{0.659538in}}%
\pgfpathlineto{\pgfqpoint{1.967619in}{0.732505in}}%
\pgfpathlineto{\pgfqpoint{1.969901in}{0.732505in}}%
\pgfpathlineto{\pgfqpoint{1.972184in}{0.718824in}}%
\pgfpathlineto{\pgfqpoint{1.974466in}{0.718824in}}%
\pgfpathlineto{\pgfqpoint{1.976748in}{0.664099in}}%
\pgfpathlineto{\pgfqpoint{1.979030in}{0.741341in}}%
\pgfpathlineto{\pgfqpoint{1.981313in}{0.718539in}}%
\pgfpathlineto{\pgfqpoint{1.983595in}{0.713979in}}%
\pgfpathlineto{\pgfqpoint{1.985877in}{0.691176in}}%
\pgfpathlineto{\pgfqpoint{1.988159in}{0.704858in}}%
\pgfpathlineto{\pgfqpoint{1.990442in}{0.755022in}}%
\pgfpathlineto{\pgfqpoint{1.992724in}{0.723099in}}%
\pgfpathlineto{\pgfqpoint{1.995006in}{0.736781in}}%
\pgfpathlineto{\pgfqpoint{1.997289in}{0.663814in}}%
\pgfpathlineto{\pgfqpoint{1.999571in}{0.713979in}}%
\pgfpathlineto{\pgfqpoint{2.001853in}{0.700297in}}%
\pgfpathlineto{\pgfqpoint{2.004135in}{0.727660in}}%
\pgfpathlineto{\pgfqpoint{2.006418in}{0.677495in}}%
\pgfpathlineto{\pgfqpoint{2.008700in}{0.686616in}}%
\pgfpathlineto{\pgfqpoint{2.010982in}{0.736781in}}%
\pgfpathlineto{\pgfqpoint{2.013264in}{0.741341in}}%
\pgfpathlineto{\pgfqpoint{2.015547in}{0.713979in}}%
\pgfpathlineto{\pgfqpoint{2.017829in}{0.713979in}}%
\pgfpathlineto{\pgfqpoint{2.022393in}{0.736781in}}%
\pgfpathlineto{\pgfqpoint{2.024676in}{0.700297in}}%
\pgfpathlineto{\pgfqpoint{2.026958in}{0.750462in}}%
\pgfpathlineto{\pgfqpoint{2.031522in}{0.709418in}}%
\pgfpathlineto{\pgfqpoint{2.033805in}{0.654693in}}%
\pgfpathlineto{\pgfqpoint{2.036087in}{0.704858in}}%
\pgfpathlineto{\pgfqpoint{2.038369in}{0.700297in}}%
\pgfpathlineto{\pgfqpoint{2.040651in}{0.741341in}}%
\pgfpathlineto{\pgfqpoint{2.045216in}{0.668374in}}%
\pgfpathlineto{\pgfqpoint{2.047498in}{0.736781in}}%
\pgfpathlineto{\pgfqpoint{2.049780in}{0.686616in}}%
\pgfpathlineto{\pgfqpoint{2.052063in}{0.759583in}}%
\pgfpathlineto{\pgfqpoint{2.054345in}{0.700297in}}%
\pgfpathlineto{\pgfqpoint{2.058909in}{0.741056in}}%
\pgfpathlineto{\pgfqpoint{2.065756in}{0.695452in}}%
\pgfpathlineto{\pgfqpoint{2.068038in}{0.709133in}}%
\pgfpathlineto{\pgfqpoint{2.070321in}{0.713694in}}%
\pgfpathlineto{\pgfqpoint{2.072603in}{0.727375in}}%
\pgfpathlineto{\pgfqpoint{2.074885in}{0.731935in}}%
\pgfpathlineto{\pgfqpoint{2.077167in}{0.731935in}}%
\pgfpathlineto{\pgfqpoint{2.079450in}{0.782100in}}%
\pgfpathlineto{\pgfqpoint{2.081732in}{0.745617in}}%
\pgfpathlineto{\pgfqpoint{2.084014in}{0.736496in}}%
\pgfpathlineto{\pgfqpoint{2.086296in}{0.695452in}}%
\pgfpathlineto{\pgfqpoint{2.088579in}{0.704573in}}%
\pgfpathlineto{\pgfqpoint{2.090861in}{0.741056in}}%
\pgfpathlineto{\pgfqpoint{2.093143in}{0.736496in}}%
\pgfpathlineto{\pgfqpoint{2.095425in}{0.718254in}}%
\pgfpathlineto{\pgfqpoint{2.099990in}{0.700012in}}%
\pgfpathlineto{\pgfqpoint{2.102272in}{0.718254in}}%
\pgfpathlineto{\pgfqpoint{2.104554in}{0.722814in}}%
\pgfpathlineto{\pgfqpoint{2.106837in}{0.677210in}}%
\pgfpathlineto{\pgfqpoint{2.109119in}{0.759298in}}%
\pgfpathlineto{\pgfqpoint{2.111401in}{0.658968in}}%
\pgfpathlineto{\pgfqpoint{2.113683in}{0.731935in}}%
\pgfpathlineto{\pgfqpoint{2.115966in}{0.759298in}}%
\pgfpathlineto{\pgfqpoint{2.118248in}{0.736496in}}%
\pgfpathlineto{\pgfqpoint{2.120530in}{0.750177in}}%
\pgfpathlineto{\pgfqpoint{2.122812in}{0.772979in}}%
\pgfpathlineto{\pgfqpoint{2.125095in}{0.722814in}}%
\pgfpathlineto{\pgfqpoint{2.131941in}{0.700012in}}%
\pgfpathlineto{\pgfqpoint{2.134224in}{0.658968in}}%
\pgfpathlineto{\pgfqpoint{2.136506in}{0.731935in}}%
\pgfpathlineto{\pgfqpoint{2.138788in}{0.667804in}}%
\pgfpathlineto{\pgfqpoint{2.141070in}{0.713408in}}%
\pgfpathlineto{\pgfqpoint{2.143353in}{0.717969in}}%
\pgfpathlineto{\pgfqpoint{2.145635in}{0.704288in}}%
\pgfpathlineto{\pgfqpoint{2.147917in}{0.676925in}}%
\pgfpathlineto{\pgfqpoint{2.150199in}{0.740771in}}%
\pgfpathlineto{\pgfqpoint{2.152482in}{0.708848in}}%
\pgfpathlineto{\pgfqpoint{2.154764in}{0.740771in}}%
\pgfpathlineto{\pgfqpoint{2.157046in}{0.695167in}}%
\pgfpathlineto{\pgfqpoint{2.159328in}{0.768134in}}%
\pgfpathlineto{\pgfqpoint{2.161611in}{0.727090in}}%
\pgfpathlineto{\pgfqpoint{2.166175in}{0.740771in}}%
\pgfpathlineto{\pgfqpoint{2.168457in}{0.722529in}}%
\pgfpathlineto{\pgfqpoint{2.170740in}{0.722529in}}%
\pgfpathlineto{\pgfqpoint{2.173022in}{0.681485in}}%
\pgfpathlineto{\pgfqpoint{2.175304in}{0.740771in}}%
\pgfpathlineto{\pgfqpoint{2.177586in}{0.717969in}}%
\pgfpathlineto{\pgfqpoint{2.179869in}{0.736211in}}%
\pgfpathlineto{\pgfqpoint{2.182151in}{0.695167in}}%
\pgfpathlineto{\pgfqpoint{2.184433in}{0.749892in}}%
\pgfpathlineto{\pgfqpoint{2.186715in}{0.713408in}}%
\pgfpathlineto{\pgfqpoint{2.188998in}{0.722529in}}%
\pgfpathlineto{\pgfqpoint{2.191280in}{0.754452in}}%
\pgfpathlineto{\pgfqpoint{2.193562in}{0.722529in}}%
\pgfpathlineto{\pgfqpoint{2.195844in}{0.736211in}}%
\pgfpathlineto{\pgfqpoint{2.198127in}{0.695167in}}%
\pgfpathlineto{\pgfqpoint{2.200409in}{0.704288in}}%
\pgfpathlineto{\pgfqpoint{2.202691in}{0.690606in}}%
\pgfpathlineto{\pgfqpoint{2.204973in}{0.713408in}}%
\pgfpathlineto{\pgfqpoint{2.207256in}{0.713408in}}%
\pgfpathlineto{\pgfqpoint{2.211820in}{0.699727in}}%
\pgfpathlineto{\pgfqpoint{2.214102in}{0.663244in}}%
\pgfpathlineto{\pgfqpoint{2.216385in}{0.740771in}}%
\pgfpathlineto{\pgfqpoint{2.220949in}{0.694882in}}%
\pgfpathlineto{\pgfqpoint{2.223231in}{0.722244in}}%
\pgfpathlineto{\pgfqpoint{2.225514in}{0.717684in}}%
\pgfpathlineto{\pgfqpoint{2.230078in}{0.667519in}}%
\pgfpathlineto{\pgfqpoint{2.232360in}{0.694882in}}%
\pgfpathlineto{\pgfqpoint{2.234643in}{0.694882in}}%
\pgfpathlineto{\pgfqpoint{2.236925in}{0.726805in}}%
\pgfpathlineto{\pgfqpoint{2.239207in}{0.708563in}}%
\pgfpathlineto{\pgfqpoint{2.241489in}{0.667519in}}%
\pgfpathlineto{\pgfqpoint{2.246054in}{0.767849in}}%
\pgfpathlineto{\pgfqpoint{2.250618in}{0.704003in}}%
\pgfpathlineto{\pgfqpoint{2.255183in}{0.749607in}}%
\pgfpathlineto{\pgfqpoint{2.257465in}{0.685761in}}%
\pgfpathlineto{\pgfqpoint{2.262030in}{0.653838in}}%
\pgfpathlineto{\pgfqpoint{2.266594in}{0.735926in}}%
\pgfpathlineto{\pgfqpoint{2.268876in}{0.713123in}}%
\pgfpathlineto{\pgfqpoint{2.271159in}{0.735926in}}%
\pgfpathlineto{\pgfqpoint{2.273441in}{0.713123in}}%
\pgfpathlineto{\pgfqpoint{2.275723in}{0.722244in}}%
\pgfpathlineto{\pgfqpoint{2.280288in}{0.667519in}}%
\pgfpathlineto{\pgfqpoint{2.282570in}{0.722244in}}%
\pgfpathlineto{\pgfqpoint{2.287134in}{0.754167in}}%
\pgfpathlineto{\pgfqpoint{2.289417in}{0.681200in}}%
\pgfpathlineto{\pgfqpoint{2.291699in}{0.713123in}}%
\pgfpathlineto{\pgfqpoint{2.293981in}{0.662959in}}%
\pgfpathlineto{\pgfqpoint{2.296264in}{0.740486in}}%
\pgfpathlineto{\pgfqpoint{2.298546in}{0.726520in}}%
\pgfpathlineto{\pgfqpoint{2.300828in}{0.767564in}}%
\pgfpathlineto{\pgfqpoint{2.303110in}{0.731080in}}%
\pgfpathlineto{\pgfqpoint{2.305393in}{0.735641in}}%
\pgfpathlineto{\pgfqpoint{2.307675in}{0.744761in}}%
\pgfpathlineto{\pgfqpoint{2.309957in}{0.703718in}}%
\pgfpathlineto{\pgfqpoint{2.312239in}{0.717399in}}%
\pgfpathlineto{\pgfqpoint{2.316804in}{0.767564in}}%
\pgfpathlineto{\pgfqpoint{2.319086in}{0.744761in}}%
\pgfpathlineto{\pgfqpoint{2.321368in}{0.758443in}}%
\pgfpathlineto{\pgfqpoint{2.323651in}{0.749322in}}%
\pgfpathlineto{\pgfqpoint{2.325933in}{0.703718in}}%
\pgfpathlineto{\pgfqpoint{2.328215in}{0.753882in}}%
\pgfpathlineto{\pgfqpoint{2.330497in}{0.690036in}}%
\pgfpathlineto{\pgfqpoint{2.332780in}{0.699157in}}%
\pgfpathlineto{\pgfqpoint{2.335062in}{0.717399in}}%
\pgfpathlineto{\pgfqpoint{2.339626in}{1.068552in}}%
\pgfpathlineto{\pgfqpoint{2.341909in}{1.123277in}}%
\pgfpathlineto{\pgfqpoint{2.344191in}{1.123277in}}%
\pgfpathlineto{\pgfqpoint{2.346473in}{1.082233in}}%
\pgfpathlineto{\pgfqpoint{2.348755in}{0.963662in}}%
\pgfpathlineto{\pgfqpoint{2.351038in}{0.753882in}}%
\pgfpathlineto{\pgfqpoint{2.353320in}{0.721959in}}%
\pgfpathlineto{\pgfqpoint{2.355602in}{0.749322in}}%
\pgfpathlineto{\pgfqpoint{2.357884in}{0.731080in}}%
\pgfpathlineto{\pgfqpoint{2.360167in}{0.671795in}}%
\pgfpathlineto{\pgfqpoint{2.362449in}{0.753882in}}%
\pgfpathlineto{\pgfqpoint{2.364731in}{0.721959in}}%
\pgfpathlineto{\pgfqpoint{2.367013in}{0.721959in}}%
\pgfpathlineto{\pgfqpoint{2.369296in}{0.753882in}}%
\pgfpathlineto{\pgfqpoint{2.371578in}{0.735641in}}%
\pgfpathlineto{\pgfqpoint{2.373860in}{0.763003in}}%
\pgfpathlineto{\pgfqpoint{2.376142in}{0.708278in}}%
\pgfpathlineto{\pgfqpoint{2.378425in}{0.698872in}}%
\pgfpathlineto{\pgfqpoint{2.382989in}{0.721674in}}%
\pgfpathlineto{\pgfqpoint{2.385271in}{0.717114in}}%
\pgfpathlineto{\pgfqpoint{2.387554in}{0.739916in}}%
\pgfpathlineto{\pgfqpoint{2.389836in}{0.698872in}}%
\pgfpathlineto{\pgfqpoint{2.394400in}{0.744476in}}%
\pgfpathlineto{\pgfqpoint{2.396683in}{0.753597in}}%
\pgfpathlineto{\pgfqpoint{2.398965in}{0.685191in}}%
\pgfpathlineto{\pgfqpoint{2.401247in}{0.758158in}}%
\pgfpathlineto{\pgfqpoint{2.403529in}{0.758158in}}%
\pgfpathlineto{\pgfqpoint{2.405812in}{0.762718in}}%
\pgfpathlineto{\pgfqpoint{2.412658in}{0.694312in}}%
\pgfpathlineto{\pgfqpoint{2.414941in}{0.744476in}}%
\pgfpathlineto{\pgfqpoint{2.417223in}{0.698872in}}%
\pgfpathlineto{\pgfqpoint{2.419505in}{0.730795in}}%
\pgfpathlineto{\pgfqpoint{2.421787in}{0.739916in}}%
\pgfpathlineto{\pgfqpoint{2.424070in}{0.712553in}}%
\pgfpathlineto{\pgfqpoint{2.426352in}{0.721674in}}%
\pgfpathlineto{\pgfqpoint{2.428634in}{0.689751in}}%
\pgfpathlineto{\pgfqpoint{2.430916in}{0.698872in}}%
\pgfpathlineto{\pgfqpoint{2.433199in}{0.739916in}}%
\pgfpathlineto{\pgfqpoint{2.435481in}{0.712553in}}%
\pgfpathlineto{\pgfqpoint{2.437763in}{0.762718in}}%
\pgfpathlineto{\pgfqpoint{2.440045in}{0.758158in}}%
\pgfpathlineto{\pgfqpoint{2.442328in}{0.744476in}}%
\pgfpathlineto{\pgfqpoint{2.444610in}{0.762718in}}%
\pgfpathlineto{\pgfqpoint{2.446892in}{0.726235in}}%
\pgfpathlineto{\pgfqpoint{2.449174in}{0.730795in}}%
\pgfpathlineto{\pgfqpoint{2.451457in}{0.717114in}}%
\pgfpathlineto{\pgfqpoint{2.453739in}{0.730795in}}%
\pgfpathlineto{\pgfqpoint{2.456021in}{0.717114in}}%
\pgfpathlineto{\pgfqpoint{2.458303in}{0.744191in}}%
\pgfpathlineto{\pgfqpoint{2.460586in}{0.689466in}}%
\pgfpathlineto{\pgfqpoint{2.462868in}{0.707708in}}%
\pgfpathlineto{\pgfqpoint{2.465150in}{0.757873in}}%
\pgfpathlineto{\pgfqpoint{2.467432in}{0.730510in}}%
\pgfpathlineto{\pgfqpoint{2.469715in}{0.730510in}}%
\pgfpathlineto{\pgfqpoint{2.474279in}{0.689466in}}%
\pgfpathlineto{\pgfqpoint{2.476561in}{0.680345in}}%
\pgfpathlineto{\pgfqpoint{2.478844in}{0.744191in}}%
\pgfpathlineto{\pgfqpoint{2.481126in}{0.725950in}}%
\pgfpathlineto{\pgfqpoint{2.483408in}{0.725950in}}%
\pgfpathlineto{\pgfqpoint{2.485690in}{0.721389in}}%
\pgfpathlineto{\pgfqpoint{2.487973in}{0.748752in}}%
\pgfpathlineto{\pgfqpoint{2.490255in}{0.730510in}}%
\pgfpathlineto{\pgfqpoint{2.492537in}{0.680345in}}%
\pgfpathlineto{\pgfqpoint{2.494819in}{0.721389in}}%
\pgfpathlineto{\pgfqpoint{2.497102in}{0.639302in}}%
\pgfpathlineto{\pgfqpoint{2.499384in}{0.725950in}}%
\pgfpathlineto{\pgfqpoint{2.501666in}{0.680345in}}%
\pgfpathlineto{\pgfqpoint{2.503948in}{0.707708in}}%
\pgfpathlineto{\pgfqpoint{2.506231in}{0.753312in}}%
\pgfpathlineto{\pgfqpoint{2.508513in}{0.753312in}}%
\pgfpathlineto{\pgfqpoint{2.513077in}{0.657543in}}%
\pgfpathlineto{\pgfqpoint{2.515360in}{0.675785in}}%
\pgfpathlineto{\pgfqpoint{2.517642in}{0.730510in}}%
\pgfpathlineto{\pgfqpoint{2.519924in}{0.703148in}}%
\pgfpathlineto{\pgfqpoint{2.522206in}{0.712268in}}%
\pgfpathlineto{\pgfqpoint{2.524489in}{0.680345in}}%
\pgfpathlineto{\pgfqpoint{2.526771in}{0.698587in}}%
\pgfpathlineto{\pgfqpoint{2.529053in}{0.730510in}}%
\pgfpathlineto{\pgfqpoint{2.531335in}{0.735071in}}%
\pgfpathlineto{\pgfqpoint{2.535900in}{0.766994in}}%
\pgfpathlineto{\pgfqpoint{2.538182in}{0.766708in}}%
\pgfpathlineto{\pgfqpoint{2.540464in}{0.721104in}}%
\pgfpathlineto{\pgfqpoint{2.542747in}{0.702863in}}%
\pgfpathlineto{\pgfqpoint{2.545029in}{0.743906in}}%
\pgfpathlineto{\pgfqpoint{2.547311in}{0.711983in}}%
\pgfpathlineto{\pgfqpoint{2.549593in}{0.639017in}}%
\pgfpathlineto{\pgfqpoint{2.551876in}{0.693742in}}%
\pgfpathlineto{\pgfqpoint{2.554158in}{0.698302in}}%
\pgfpathlineto{\pgfqpoint{2.556440in}{0.639017in}}%
\pgfpathlineto{\pgfqpoint{2.558722in}{0.771269in}}%
\pgfpathlineto{\pgfqpoint{2.561005in}{0.693742in}}%
\pgfpathlineto{\pgfqpoint{2.563287in}{0.762148in}}%
\pgfpathlineto{\pgfqpoint{2.565569in}{0.753027in}}%
\pgfpathlineto{\pgfqpoint{2.567851in}{0.702863in}}%
\pgfpathlineto{\pgfqpoint{2.570134in}{0.707423in}}%
\pgfpathlineto{\pgfqpoint{2.572416in}{0.707423in}}%
\pgfpathlineto{\pgfqpoint{2.574698in}{0.725665in}}%
\pgfpathlineto{\pgfqpoint{2.576980in}{0.680060in}}%
\pgfpathlineto{\pgfqpoint{2.579263in}{0.702863in}}%
\pgfpathlineto{\pgfqpoint{2.581545in}{0.743906in}}%
\pgfpathlineto{\pgfqpoint{2.583827in}{0.707423in}}%
\pgfpathlineto{\pgfqpoint{2.586110in}{0.734785in}}%
\pgfpathlineto{\pgfqpoint{2.588392in}{0.730225in}}%
\pgfpathlineto{\pgfqpoint{2.590674in}{0.711983in}}%
\pgfpathlineto{\pgfqpoint{2.592956in}{0.707423in}}%
\pgfpathlineto{\pgfqpoint{2.595239in}{0.716544in}}%
\pgfpathlineto{\pgfqpoint{2.597521in}{0.702863in}}%
\pgfpathlineto{\pgfqpoint{2.599803in}{0.743906in}}%
\pgfpathlineto{\pgfqpoint{2.602085in}{0.725665in}}%
\pgfpathlineto{\pgfqpoint{2.606650in}{0.753027in}}%
\pgfpathlineto{\pgfqpoint{2.608932in}{0.725665in}}%
\pgfpathlineto{\pgfqpoint{2.611214in}{0.730225in}}%
\pgfpathlineto{\pgfqpoint{2.613497in}{0.661819in}}%
\pgfpathlineto{\pgfqpoint{2.618061in}{0.757303in}}%
\pgfpathlineto{\pgfqpoint{2.620343in}{0.711698in}}%
\pgfpathlineto{\pgfqpoint{2.622626in}{0.752742in}}%
\pgfpathlineto{\pgfqpoint{2.624908in}{0.729940in}}%
\pgfpathlineto{\pgfqpoint{2.627190in}{0.675215in}}%
\pgfpathlineto{\pgfqpoint{2.629472in}{0.711698in}}%
\pgfpathlineto{\pgfqpoint{2.631755in}{0.693457in}}%
\pgfpathlineto{\pgfqpoint{2.634037in}{0.684336in}}%
\pgfpathlineto{\pgfqpoint{2.636319in}{0.702577in}}%
\pgfpathlineto{\pgfqpoint{2.638601in}{0.739061in}}%
\pgfpathlineto{\pgfqpoint{2.640884in}{0.739061in}}%
\pgfpathlineto{\pgfqpoint{2.643166in}{0.720819in}}%
\pgfpathlineto{\pgfqpoint{2.645448in}{0.711698in}}%
\pgfpathlineto{\pgfqpoint{2.647730in}{0.693457in}}%
\pgfpathlineto{\pgfqpoint{2.650013in}{0.716259in}}%
\pgfpathlineto{\pgfqpoint{2.652295in}{0.725380in}}%
\pgfpathlineto{\pgfqpoint{2.654577in}{0.711698in}}%
\pgfpathlineto{\pgfqpoint{2.656859in}{0.725380in}}%
\pgfpathlineto{\pgfqpoint{2.659142in}{0.693457in}}%
\pgfpathlineto{\pgfqpoint{2.661424in}{0.757303in}}%
\pgfpathlineto{\pgfqpoint{2.663706in}{0.720819in}}%
\pgfpathlineto{\pgfqpoint{2.665988in}{0.716259in}}%
\pgfpathlineto{\pgfqpoint{2.668271in}{0.643292in}}%
\pgfpathlineto{\pgfqpoint{2.670553in}{0.702577in}}%
\pgfpathlineto{\pgfqpoint{2.672835in}{0.711698in}}%
\pgfpathlineto{\pgfqpoint{2.675117in}{0.698017in}}%
\pgfpathlineto{\pgfqpoint{2.677400in}{0.698017in}}%
\pgfpathlineto{\pgfqpoint{2.679682in}{0.725380in}}%
\pgfpathlineto{\pgfqpoint{2.681964in}{0.647852in}}%
\pgfpathlineto{\pgfqpoint{2.684246in}{0.656973in}}%
\pgfpathlineto{\pgfqpoint{2.686529in}{0.707138in}}%
\pgfpathlineto{\pgfqpoint{2.688811in}{0.684336in}}%
\pgfpathlineto{\pgfqpoint{2.691093in}{0.729940in}}%
\pgfpathlineto{\pgfqpoint{2.695658in}{0.606808in}}%
\pgfpathlineto{\pgfqpoint{2.697940in}{0.693172in}}%
\pgfpathlineto{\pgfqpoint{2.700222in}{0.688611in}}%
\pgfpathlineto{\pgfqpoint{2.702504in}{0.711413in}}%
\pgfpathlineto{\pgfqpoint{2.704787in}{0.606523in}}%
\pgfpathlineto{\pgfqpoint{2.707069in}{0.620205in}}%
\pgfpathlineto{\pgfqpoint{2.709351in}{0.679490in}}%
\pgfpathlineto{\pgfqpoint{2.716198in}{0.679490in}}%
\pgfpathlineto{\pgfqpoint{2.718480in}{0.720534in}}%
\pgfpathlineto{\pgfqpoint{2.720762in}{0.679490in}}%
\pgfpathlineto{\pgfqpoint{2.723045in}{0.702292in}}%
\pgfpathlineto{\pgfqpoint{2.725327in}{0.693172in}}%
\pgfpathlineto{\pgfqpoint{2.727609in}{0.706853in}}%
\pgfpathlineto{\pgfqpoint{2.729891in}{0.656688in}}%
\pgfpathlineto{\pgfqpoint{2.732174in}{0.697732in}}%
\pgfpathlineto{\pgfqpoint{2.734456in}{0.711413in}}%
\pgfpathlineto{\pgfqpoint{2.736738in}{0.702292in}}%
\pgfpathlineto{\pgfqpoint{2.739020in}{0.702292in}}%
\pgfpathlineto{\pgfqpoint{2.741303in}{0.665809in}}%
\pgfpathlineto{\pgfqpoint{2.743585in}{0.688611in}}%
\pgfpathlineto{\pgfqpoint{2.745867in}{0.624765in}}%
\pgfpathlineto{\pgfqpoint{2.750432in}{0.647567in}}%
\pgfpathlineto{\pgfqpoint{2.752714in}{0.679490in}}%
\pgfpathlineto{\pgfqpoint{2.754996in}{0.670369in}}%
\pgfpathlineto{\pgfqpoint{2.757278in}{0.697732in}}%
\pgfpathlineto{\pgfqpoint{2.759561in}{0.665809in}}%
\pgfpathlineto{\pgfqpoint{2.761843in}{0.661249in}}%
\pgfpathlineto{\pgfqpoint{2.764125in}{0.706853in}}%
\pgfpathlineto{\pgfqpoint{2.766407in}{0.624765in}}%
\pgfpathlineto{\pgfqpoint{2.768690in}{0.661249in}}%
\pgfpathlineto{\pgfqpoint{2.770972in}{0.715974in}}%
\pgfpathlineto{\pgfqpoint{2.773254in}{0.620205in}}%
\pgfpathlineto{\pgfqpoint{2.775536in}{0.629326in}}%
\pgfpathlineto{\pgfqpoint{2.780101in}{0.688326in}}%
\pgfpathlineto{\pgfqpoint{2.782383in}{0.651843in}}%
\pgfpathlineto{\pgfqpoint{2.784665in}{0.743051in}}%
\pgfpathlineto{\pgfqpoint{2.786948in}{0.674645in}}%
\pgfpathlineto{\pgfqpoint{2.789230in}{0.683766in}}%
\pgfpathlineto{\pgfqpoint{2.791512in}{0.660964in}}%
\pgfpathlineto{\pgfqpoint{2.793794in}{0.670084in}}%
\pgfpathlineto{\pgfqpoint{2.796077in}{0.670084in}}%
\pgfpathlineto{\pgfqpoint{2.798359in}{0.702007in}}%
\pgfpathlineto{\pgfqpoint{2.800641in}{0.752172in}}%
\pgfpathlineto{\pgfqpoint{2.802923in}{0.834260in}}%
\pgfpathlineto{\pgfqpoint{2.805206in}{0.998435in}}%
\pgfpathlineto{\pgfqpoint{2.807488in}{1.071402in}}%
\pgfpathlineto{\pgfqpoint{2.809770in}{1.089644in}}%
\pgfpathlineto{\pgfqpoint{2.812052in}{1.057721in}}%
\pgfpathlineto{\pgfqpoint{2.814335in}{0.966512in}}%
\pgfpathlineto{\pgfqpoint{2.816617in}{0.752172in}}%
\pgfpathlineto{\pgfqpoint{2.818899in}{0.665524in}}%
\pgfpathlineto{\pgfqpoint{2.821181in}{0.724810in}}%
\pgfpathlineto{\pgfqpoint{2.823464in}{0.692887in}}%
\pgfpathlineto{\pgfqpoint{2.825746in}{0.711128in}}%
\pgfpathlineto{\pgfqpoint{2.828028in}{0.697447in}}%
\pgfpathlineto{\pgfqpoint{2.830310in}{0.633601in}}%
\pgfpathlineto{\pgfqpoint{2.832593in}{0.706568in}}%
\pgfpathlineto{\pgfqpoint{2.834875in}{0.647282in}}%
\pgfpathlineto{\pgfqpoint{2.837157in}{0.738491in}}%
\pgfpathlineto{\pgfqpoint{2.841722in}{0.642722in}}%
\pgfpathlineto{\pgfqpoint{2.844004in}{0.702007in}}%
\pgfpathlineto{\pgfqpoint{2.846286in}{0.642722in}}%
\pgfpathlineto{\pgfqpoint{2.848568in}{0.683766in}}%
\pgfpathlineto{\pgfqpoint{2.850851in}{0.679205in}}%
\pgfpathlineto{\pgfqpoint{2.853133in}{0.729370in}}%
\pgfpathlineto{\pgfqpoint{2.855415in}{0.638161in}}%
\pgfpathlineto{\pgfqpoint{2.857697in}{0.651558in}}%
\pgfpathlineto{\pgfqpoint{2.859980in}{0.646997in}}%
\pgfpathlineto{\pgfqpoint{2.862262in}{0.710843in}}%
\pgfpathlineto{\pgfqpoint{2.864544in}{0.642437in}}%
\pgfpathlineto{\pgfqpoint{2.869109in}{0.710843in}}%
\pgfpathlineto{\pgfqpoint{2.871391in}{0.697162in}}%
\pgfpathlineto{\pgfqpoint{2.873673in}{0.742766in}}%
\pgfpathlineto{\pgfqpoint{2.875955in}{0.729085in}}%
\pgfpathlineto{\pgfqpoint{2.878238in}{0.701722in}}%
\pgfpathlineto{\pgfqpoint{2.880520in}{0.706283in}}%
\pgfpathlineto{\pgfqpoint{2.882802in}{0.624195in}}%
\pgfpathlineto{\pgfqpoint{2.885085in}{0.710843in}}%
\pgfpathlineto{\pgfqpoint{2.887367in}{0.660679in}}%
\pgfpathlineto{\pgfqpoint{2.889649in}{0.692602in}}%
\pgfpathlineto{\pgfqpoint{2.891931in}{0.669799in}}%
\pgfpathlineto{\pgfqpoint{2.894214in}{0.697162in}}%
\pgfpathlineto{\pgfqpoint{2.896496in}{0.619635in}}%
\pgfpathlineto{\pgfqpoint{2.898778in}{0.733645in}}%
\pgfpathlineto{\pgfqpoint{2.901060in}{0.710843in}}%
\pgfpathlineto{\pgfqpoint{2.903343in}{0.701722in}}%
\pgfpathlineto{\pgfqpoint{2.905625in}{0.665239in}}%
\pgfpathlineto{\pgfqpoint{2.907907in}{0.660679in}}%
\pgfpathlineto{\pgfqpoint{2.910189in}{0.674360in}}%
\pgfpathlineto{\pgfqpoint{2.912472in}{0.710843in}}%
\pgfpathlineto{\pgfqpoint{2.914754in}{0.674360in}}%
\pgfpathlineto{\pgfqpoint{2.917036in}{0.674360in}}%
\pgfpathlineto{\pgfqpoint{2.919318in}{0.715404in}}%
\pgfpathlineto{\pgfqpoint{2.921601in}{0.697162in}}%
\pgfpathlineto{\pgfqpoint{2.923883in}{0.697162in}}%
\pgfpathlineto{\pgfqpoint{2.926165in}{0.660679in}}%
\pgfpathlineto{\pgfqpoint{2.928447in}{0.660679in}}%
\pgfpathlineto{\pgfqpoint{2.930730in}{0.628756in}}%
\pgfpathlineto{\pgfqpoint{2.933012in}{0.628756in}}%
\pgfpathlineto{\pgfqpoint{2.937576in}{0.701437in}}%
\pgfpathlineto{\pgfqpoint{2.939859in}{0.705998in}}%
\pgfpathlineto{\pgfqpoint{2.942141in}{0.701437in}}%
\pgfpathlineto{\pgfqpoint{2.946705in}{0.651273in}}%
\pgfpathlineto{\pgfqpoint{2.948988in}{0.655833in}}%
\pgfpathlineto{\pgfqpoint{2.951270in}{0.687756in}}%
\pgfpathlineto{\pgfqpoint{2.953552in}{0.642152in}}%
\pgfpathlineto{\pgfqpoint{2.955834in}{0.655833in}}%
\pgfpathlineto{\pgfqpoint{2.958117in}{0.637591in}}%
\pgfpathlineto{\pgfqpoint{2.960399in}{0.646712in}}%
\pgfpathlineto{\pgfqpoint{2.964963in}{0.737921in}}%
\pgfpathlineto{\pgfqpoint{2.967246in}{0.674075in}}%
\pgfpathlineto{\pgfqpoint{2.969528in}{0.728800in}}%
\pgfpathlineto{\pgfqpoint{2.971810in}{0.633031in}}%
\pgfpathlineto{\pgfqpoint{2.974092in}{0.642152in}}%
\pgfpathlineto{\pgfqpoint{2.976375in}{0.669514in}}%
\pgfpathlineto{\pgfqpoint{2.978657in}{0.669514in}}%
\pgfpathlineto{\pgfqpoint{2.980939in}{0.710558in}}%
\pgfpathlineto{\pgfqpoint{2.983221in}{0.687756in}}%
\pgfpathlineto{\pgfqpoint{2.985504in}{0.642152in}}%
\pgfpathlineto{\pgfqpoint{2.987786in}{0.651273in}}%
\pgfpathlineto{\pgfqpoint{2.990068in}{0.687756in}}%
\pgfpathlineto{\pgfqpoint{2.992350in}{0.669514in}}%
\pgfpathlineto{\pgfqpoint{2.994633in}{0.728800in}}%
\pgfpathlineto{\pgfqpoint{2.996915in}{0.664954in}}%
\pgfpathlineto{\pgfqpoint{2.999197in}{0.692317in}}%
\pgfpathlineto{\pgfqpoint{3.001479in}{0.669514in}}%
\pgfpathlineto{\pgfqpoint{3.003762in}{0.710558in}}%
\pgfpathlineto{\pgfqpoint{3.006044in}{0.724240in}}%
\pgfpathlineto{\pgfqpoint{3.010608in}{0.678635in}}%
\pgfpathlineto{\pgfqpoint{3.012891in}{0.751602in}}%
\pgfpathlineto{\pgfqpoint{3.015173in}{0.701437in}}%
\pgfpathlineto{\pgfqpoint{3.017455in}{0.682911in}}%
\pgfpathlineto{\pgfqpoint{3.022020in}{0.705713in}}%
\pgfpathlineto{\pgfqpoint{3.024302in}{0.628185in}}%
\pgfpathlineto{\pgfqpoint{3.026584in}{0.678350in}}%
\pgfpathlineto{\pgfqpoint{3.028866in}{0.682911in}}%
\pgfpathlineto{\pgfqpoint{3.031149in}{0.692031in}}%
\pgfpathlineto{\pgfqpoint{3.037995in}{0.746757in}}%
\pgfpathlineto{\pgfqpoint{3.040278in}{0.719394in}}%
\pgfpathlineto{\pgfqpoint{3.040278in}{0.719394in}}%
\pgfusepath{stroke}%
\end{pgfscope}%
\begin{pgfscope}%
\pgfsetrectcap%
\pgfsetmiterjoin%
\pgfsetlinewidth{0.803000pt}%
\definecolor{currentstroke}{rgb}{0.000000,0.000000,0.000000}%
\pgfsetstrokecolor{currentstroke}%
\pgfsetdash{}{0pt}%
\pgfpathmoveto{\pgfqpoint{0.758026in}{0.565123in}}%
\pgfpathlineto{\pgfqpoint{0.758026in}{1.475926in}}%
\pgfusepath{stroke}%
\end{pgfscope}%
\begin{pgfscope}%
\pgfsetrectcap%
\pgfsetmiterjoin%
\pgfsetlinewidth{0.803000pt}%
\definecolor{currentstroke}{rgb}{0.000000,0.000000,0.000000}%
\pgfsetstrokecolor{currentstroke}%
\pgfsetdash{}{0pt}%
\pgfpathmoveto{\pgfqpoint{3.040278in}{0.565123in}}%
\pgfpathlineto{\pgfqpoint{3.040278in}{1.475926in}}%
\pgfusepath{stroke}%
\end{pgfscope}%
\begin{pgfscope}%
\pgfsetrectcap%
\pgfsetmiterjoin%
\pgfsetlinewidth{0.803000pt}%
\definecolor{currentstroke}{rgb}{0.000000,0.000000,0.000000}%
\pgfsetstrokecolor{currentstroke}%
\pgfsetdash{}{0pt}%
\pgfpathmoveto{\pgfqpoint{0.758026in}{0.565123in}}%
\pgfpathlineto{\pgfqpoint{3.040278in}{0.565123in}}%
\pgfusepath{stroke}%
\end{pgfscope}%
\begin{pgfscope}%
\pgfsetrectcap%
\pgfsetmiterjoin%
\pgfsetlinewidth{0.803000pt}%
\definecolor{currentstroke}{rgb}{0.000000,0.000000,0.000000}%
\pgfsetstrokecolor{currentstroke}%
\pgfsetdash{}{0pt}%
\pgfpathmoveto{\pgfqpoint{0.758026in}{1.475926in}}%
\pgfpathlineto{\pgfqpoint{3.040278in}{1.475926in}}%
\pgfusepath{stroke}%
\end{pgfscope}%
\begin{pgfscope}%
\definecolor{textcolor}{rgb}{0.000000,0.000000,0.000000}%
\pgfsetstrokecolor{textcolor}%
\pgfsetfillcolor{textcolor}%
\pgftext[x=1.899152in,y=1.559260in,,base]{\color{textcolor}\rmfamily\fontsize{12.000000}{14.400000}\selectfont Triangle Wave}%
\end{pgfscope}%
\begin{pgfscope}%
\pgfsetbuttcap%
\pgfsetmiterjoin%
\definecolor{currentfill}{rgb}{1.000000,1.000000,1.000000}%
\pgfsetfillcolor{currentfill}%
\pgfsetlinewidth{0.000000pt}%
\definecolor{currentstroke}{rgb}{0.000000,0.000000,0.000000}%
\pgfsetstrokecolor{currentstroke}%
\pgfsetstrokeopacity{0.000000}%
\pgfsetdash{}{0pt}%
\pgfpathmoveto{\pgfqpoint{3.833026in}{0.565123in}}%
\pgfpathlineto{\pgfqpoint{6.115278in}{0.565123in}}%
\pgfpathlineto{\pgfqpoint{6.115278in}{1.475926in}}%
\pgfpathlineto{\pgfqpoint{3.833026in}{1.475926in}}%
\pgfpathlineto{\pgfqpoint{3.833026in}{0.565123in}}%
\pgfpathclose%
\pgfusepath{fill}%
\end{pgfscope}%
\begin{pgfscope}%
\pgfpathrectangle{\pgfqpoint{3.833026in}{0.565123in}}{\pgfqpoint{2.282252in}{0.910803in}}%
\pgfusepath{clip}%
\pgfsetbuttcap%
\pgfsetroundjoin%
\pgfsetlinewidth{0.501875pt}%
\definecolor{currentstroke}{rgb}{0.690196,0.690196,0.690196}%
\pgfsetstrokecolor{currentstroke}%
\pgfsetstrokeopacity{0.700000}%
\pgfsetdash{{1.850000pt}{0.800000pt}}{0.000000pt}%
\pgfpathmoveto{\pgfqpoint{4.252215in}{0.565123in}}%
\pgfpathlineto{\pgfqpoint{4.252215in}{1.475926in}}%
\pgfusepath{stroke}%
\end{pgfscope}%
\begin{pgfscope}%
\pgfsetbuttcap%
\pgfsetroundjoin%
\definecolor{currentfill}{rgb}{0.000000,0.000000,0.000000}%
\pgfsetfillcolor{currentfill}%
\pgfsetlinewidth{0.803000pt}%
\definecolor{currentstroke}{rgb}{0.000000,0.000000,0.000000}%
\pgfsetstrokecolor{currentstroke}%
\pgfsetdash{}{0pt}%
\pgfsys@defobject{currentmarker}{\pgfqpoint{0.000000in}{-0.048611in}}{\pgfqpoint{0.000000in}{0.000000in}}{%
\pgfpathmoveto{\pgfqpoint{0.000000in}{0.000000in}}%
\pgfpathlineto{\pgfqpoint{0.000000in}{-0.048611in}}%
\pgfusepath{stroke,fill}%
}%
\begin{pgfscope}%
\pgfsys@transformshift{4.252215in}{0.565123in}%
\pgfsys@useobject{currentmarker}{}%
\end{pgfscope}%
\end{pgfscope}%
\begin{pgfscope}%
\definecolor{textcolor}{rgb}{0.000000,0.000000,0.000000}%
\pgfsetstrokecolor{textcolor}%
\pgfsetfillcolor{textcolor}%
\pgftext[x=4.252215in,y=0.467901in,,top]{\color{textcolor}\rmfamily\fontsize{10.000000}{12.000000}\selectfont \(\displaystyle {1}\)}%
\end{pgfscope}%
\begin{pgfscope}%
\pgfpathrectangle{\pgfqpoint{3.833026in}{0.565123in}}{\pgfqpoint{2.282252in}{0.910803in}}%
\pgfusepath{clip}%
\pgfsetbuttcap%
\pgfsetroundjoin%
\pgfsetlinewidth{0.501875pt}%
\definecolor{currentstroke}{rgb}{0.690196,0.690196,0.690196}%
\pgfsetstrokecolor{currentstroke}%
\pgfsetstrokeopacity{0.700000}%
\pgfsetdash{{1.850000pt}{0.800000pt}}{0.000000pt}%
\pgfpathmoveto{\pgfqpoint{4.717981in}{0.565123in}}%
\pgfpathlineto{\pgfqpoint{4.717981in}{1.475926in}}%
\pgfusepath{stroke}%
\end{pgfscope}%
\begin{pgfscope}%
\pgfsetbuttcap%
\pgfsetroundjoin%
\definecolor{currentfill}{rgb}{0.000000,0.000000,0.000000}%
\pgfsetfillcolor{currentfill}%
\pgfsetlinewidth{0.803000pt}%
\definecolor{currentstroke}{rgb}{0.000000,0.000000,0.000000}%
\pgfsetstrokecolor{currentstroke}%
\pgfsetdash{}{0pt}%
\pgfsys@defobject{currentmarker}{\pgfqpoint{0.000000in}{-0.048611in}}{\pgfqpoint{0.000000in}{0.000000in}}{%
\pgfpathmoveto{\pgfqpoint{0.000000in}{0.000000in}}%
\pgfpathlineto{\pgfqpoint{0.000000in}{-0.048611in}}%
\pgfusepath{stroke,fill}%
}%
\begin{pgfscope}%
\pgfsys@transformshift{4.717981in}{0.565123in}%
\pgfsys@useobject{currentmarker}{}%
\end{pgfscope}%
\end{pgfscope}%
\begin{pgfscope}%
\definecolor{textcolor}{rgb}{0.000000,0.000000,0.000000}%
\pgfsetstrokecolor{textcolor}%
\pgfsetfillcolor{textcolor}%
\pgftext[x=4.717981in,y=0.467901in,,top]{\color{textcolor}\rmfamily\fontsize{10.000000}{12.000000}\selectfont \(\displaystyle {2}\)}%
\end{pgfscope}%
\begin{pgfscope}%
\pgfpathrectangle{\pgfqpoint{3.833026in}{0.565123in}}{\pgfqpoint{2.282252in}{0.910803in}}%
\pgfusepath{clip}%
\pgfsetbuttcap%
\pgfsetroundjoin%
\pgfsetlinewidth{0.501875pt}%
\definecolor{currentstroke}{rgb}{0.690196,0.690196,0.690196}%
\pgfsetstrokecolor{currentstroke}%
\pgfsetstrokeopacity{0.700000}%
\pgfsetdash{{1.850000pt}{0.800000pt}}{0.000000pt}%
\pgfpathmoveto{\pgfqpoint{5.183746in}{0.565123in}}%
\pgfpathlineto{\pgfqpoint{5.183746in}{1.475926in}}%
\pgfusepath{stroke}%
\end{pgfscope}%
\begin{pgfscope}%
\pgfsetbuttcap%
\pgfsetroundjoin%
\definecolor{currentfill}{rgb}{0.000000,0.000000,0.000000}%
\pgfsetfillcolor{currentfill}%
\pgfsetlinewidth{0.803000pt}%
\definecolor{currentstroke}{rgb}{0.000000,0.000000,0.000000}%
\pgfsetstrokecolor{currentstroke}%
\pgfsetdash{}{0pt}%
\pgfsys@defobject{currentmarker}{\pgfqpoint{0.000000in}{-0.048611in}}{\pgfqpoint{0.000000in}{0.000000in}}{%
\pgfpathmoveto{\pgfqpoint{0.000000in}{0.000000in}}%
\pgfpathlineto{\pgfqpoint{0.000000in}{-0.048611in}}%
\pgfusepath{stroke,fill}%
}%
\begin{pgfscope}%
\pgfsys@transformshift{5.183746in}{0.565123in}%
\pgfsys@useobject{currentmarker}{}%
\end{pgfscope}%
\end{pgfscope}%
\begin{pgfscope}%
\definecolor{textcolor}{rgb}{0.000000,0.000000,0.000000}%
\pgfsetstrokecolor{textcolor}%
\pgfsetfillcolor{textcolor}%
\pgftext[x=5.183746in,y=0.467901in,,top]{\color{textcolor}\rmfamily\fontsize{10.000000}{12.000000}\selectfont \(\displaystyle {3}\)}%
\end{pgfscope}%
\begin{pgfscope}%
\pgfpathrectangle{\pgfqpoint{3.833026in}{0.565123in}}{\pgfqpoint{2.282252in}{0.910803in}}%
\pgfusepath{clip}%
\pgfsetbuttcap%
\pgfsetroundjoin%
\pgfsetlinewidth{0.501875pt}%
\definecolor{currentstroke}{rgb}{0.690196,0.690196,0.690196}%
\pgfsetstrokecolor{currentstroke}%
\pgfsetstrokeopacity{0.700000}%
\pgfsetdash{{1.850000pt}{0.800000pt}}{0.000000pt}%
\pgfpathmoveto{\pgfqpoint{5.649512in}{0.565123in}}%
\pgfpathlineto{\pgfqpoint{5.649512in}{1.475926in}}%
\pgfusepath{stroke}%
\end{pgfscope}%
\begin{pgfscope}%
\pgfsetbuttcap%
\pgfsetroundjoin%
\definecolor{currentfill}{rgb}{0.000000,0.000000,0.000000}%
\pgfsetfillcolor{currentfill}%
\pgfsetlinewidth{0.803000pt}%
\definecolor{currentstroke}{rgb}{0.000000,0.000000,0.000000}%
\pgfsetstrokecolor{currentstroke}%
\pgfsetdash{}{0pt}%
\pgfsys@defobject{currentmarker}{\pgfqpoint{0.000000in}{-0.048611in}}{\pgfqpoint{0.000000in}{0.000000in}}{%
\pgfpathmoveto{\pgfqpoint{0.000000in}{0.000000in}}%
\pgfpathlineto{\pgfqpoint{0.000000in}{-0.048611in}}%
\pgfusepath{stroke,fill}%
}%
\begin{pgfscope}%
\pgfsys@transformshift{5.649512in}{0.565123in}%
\pgfsys@useobject{currentmarker}{}%
\end{pgfscope}%
\end{pgfscope}%
\begin{pgfscope}%
\definecolor{textcolor}{rgb}{0.000000,0.000000,0.000000}%
\pgfsetstrokecolor{textcolor}%
\pgfsetfillcolor{textcolor}%
\pgftext[x=5.649512in,y=0.467901in,,top]{\color{textcolor}\rmfamily\fontsize{10.000000}{12.000000}\selectfont \(\displaystyle {4}\)}%
\end{pgfscope}%
\begin{pgfscope}%
\pgfpathrectangle{\pgfqpoint{3.833026in}{0.565123in}}{\pgfqpoint{2.282252in}{0.910803in}}%
\pgfusepath{clip}%
\pgfsetbuttcap%
\pgfsetroundjoin%
\pgfsetlinewidth{0.501875pt}%
\definecolor{currentstroke}{rgb}{0.690196,0.690196,0.690196}%
\pgfsetstrokecolor{currentstroke}%
\pgfsetstrokeopacity{0.700000}%
\pgfsetdash{{1.850000pt}{0.800000pt}}{0.000000pt}%
\pgfpathmoveto{\pgfqpoint{6.115278in}{0.565123in}}%
\pgfpathlineto{\pgfqpoint{6.115278in}{1.475926in}}%
\pgfusepath{stroke}%
\end{pgfscope}%
\begin{pgfscope}%
\pgfsetbuttcap%
\pgfsetroundjoin%
\definecolor{currentfill}{rgb}{0.000000,0.000000,0.000000}%
\pgfsetfillcolor{currentfill}%
\pgfsetlinewidth{0.803000pt}%
\definecolor{currentstroke}{rgb}{0.000000,0.000000,0.000000}%
\pgfsetstrokecolor{currentstroke}%
\pgfsetdash{}{0pt}%
\pgfsys@defobject{currentmarker}{\pgfqpoint{0.000000in}{-0.048611in}}{\pgfqpoint{0.000000in}{0.000000in}}{%
\pgfpathmoveto{\pgfqpoint{0.000000in}{0.000000in}}%
\pgfpathlineto{\pgfqpoint{0.000000in}{-0.048611in}}%
\pgfusepath{stroke,fill}%
}%
\begin{pgfscope}%
\pgfsys@transformshift{6.115278in}{0.565123in}%
\pgfsys@useobject{currentmarker}{}%
\end{pgfscope}%
\end{pgfscope}%
\begin{pgfscope}%
\definecolor{textcolor}{rgb}{0.000000,0.000000,0.000000}%
\pgfsetstrokecolor{textcolor}%
\pgfsetfillcolor{textcolor}%
\pgftext[x=6.115278in,y=0.467901in,,top]{\color{textcolor}\rmfamily\fontsize{10.000000}{12.000000}\selectfont \(\displaystyle {5}\)}%
\end{pgfscope}%
\begin{pgfscope}%
\definecolor{textcolor}{rgb}{0.000000,0.000000,0.000000}%
\pgfsetstrokecolor{textcolor}%
\pgfsetfillcolor{textcolor}%
\pgftext[x=4.974152in,y=0.288889in,,top]{\color{textcolor}\rmfamily\fontsize{10.000000}{12.000000}\selectfont Frequency (Hz)}%
\end{pgfscope}%
\begin{pgfscope}%
\definecolor{textcolor}{rgb}{0.000000,0.000000,0.000000}%
\pgfsetstrokecolor{textcolor}%
\pgfsetfillcolor{textcolor}%
\pgftext[x=6.115278in,y=0.302778in,right,top]{\color{textcolor}\rmfamily\fontsize{10.000000}{12.000000}\selectfont \(\displaystyle \times{10^{6}}{}\)}%
\end{pgfscope}%
\begin{pgfscope}%
\pgfpathrectangle{\pgfqpoint{3.833026in}{0.565123in}}{\pgfqpoint{2.282252in}{0.910803in}}%
\pgfusepath{clip}%
\pgfsetbuttcap%
\pgfsetroundjoin%
\pgfsetlinewidth{0.501875pt}%
\definecolor{currentstroke}{rgb}{0.690196,0.690196,0.690196}%
\pgfsetstrokecolor{currentstroke}%
\pgfsetstrokeopacity{0.700000}%
\pgfsetdash{{1.850000pt}{0.800000pt}}{0.000000pt}%
\pgfpathmoveto{\pgfqpoint{3.833026in}{0.689386in}}%
\pgfpathlineto{\pgfqpoint{6.115278in}{0.689386in}}%
\pgfusepath{stroke}%
\end{pgfscope}%
\begin{pgfscope}%
\pgfsetbuttcap%
\pgfsetroundjoin%
\definecolor{currentfill}{rgb}{0.000000,0.000000,0.000000}%
\pgfsetfillcolor{currentfill}%
\pgfsetlinewidth{0.803000pt}%
\definecolor{currentstroke}{rgb}{0.000000,0.000000,0.000000}%
\pgfsetstrokecolor{currentstroke}%
\pgfsetdash{}{0pt}%
\pgfsys@defobject{currentmarker}{\pgfqpoint{-0.048611in}{0.000000in}}{\pgfqpoint{-0.000000in}{0.000000in}}{%
\pgfpathmoveto{\pgfqpoint{-0.000000in}{0.000000in}}%
\pgfpathlineto{\pgfqpoint{-0.048611in}{0.000000in}}%
\pgfusepath{stroke,fill}%
}%
\begin{pgfscope}%
\pgfsys@transformshift{3.833026in}{0.689386in}%
\pgfsys@useobject{currentmarker}{}%
\end{pgfscope}%
\end{pgfscope}%
\begin{pgfscope}%
\definecolor{textcolor}{rgb}{0.000000,0.000000,0.000000}%
\pgfsetstrokecolor{textcolor}%
\pgfsetfillcolor{textcolor}%
\pgftext[x=3.419444in, y=0.641161in, left, base]{\color{textcolor}\rmfamily\fontsize{10.000000}{12.000000}\selectfont \(\displaystyle {\ensuremath{-}100}\)}%
\end{pgfscope}%
\begin{pgfscope}%
\pgfpathrectangle{\pgfqpoint{3.833026in}{0.565123in}}{\pgfqpoint{2.282252in}{0.910803in}}%
\pgfusepath{clip}%
\pgfsetbuttcap%
\pgfsetroundjoin%
\pgfsetlinewidth{0.501875pt}%
\definecolor{currentstroke}{rgb}{0.690196,0.690196,0.690196}%
\pgfsetstrokecolor{currentstroke}%
\pgfsetstrokeopacity{0.700000}%
\pgfsetdash{{1.850000pt}{0.800000pt}}{0.000000pt}%
\pgfpathmoveto{\pgfqpoint{3.833026in}{1.185939in}}%
\pgfpathlineto{\pgfqpoint{6.115278in}{1.185939in}}%
\pgfusepath{stroke}%
\end{pgfscope}%
\begin{pgfscope}%
\pgfsetbuttcap%
\pgfsetroundjoin%
\definecolor{currentfill}{rgb}{0.000000,0.000000,0.000000}%
\pgfsetfillcolor{currentfill}%
\pgfsetlinewidth{0.803000pt}%
\definecolor{currentstroke}{rgb}{0.000000,0.000000,0.000000}%
\pgfsetstrokecolor{currentstroke}%
\pgfsetdash{}{0pt}%
\pgfsys@defobject{currentmarker}{\pgfqpoint{-0.048611in}{0.000000in}}{\pgfqpoint{-0.000000in}{0.000000in}}{%
\pgfpathmoveto{\pgfqpoint{-0.000000in}{0.000000in}}%
\pgfpathlineto{\pgfqpoint{-0.048611in}{0.000000in}}%
\pgfusepath{stroke,fill}%
}%
\begin{pgfscope}%
\pgfsys@transformshift{3.833026in}{1.185939in}%
\pgfsys@useobject{currentmarker}{}%
\end{pgfscope}%
\end{pgfscope}%
\begin{pgfscope}%
\definecolor{textcolor}{rgb}{0.000000,0.000000,0.000000}%
\pgfsetstrokecolor{textcolor}%
\pgfsetfillcolor{textcolor}%
\pgftext[x=3.488889in, y=1.137714in, left, base]{\color{textcolor}\rmfamily\fontsize{10.000000}{12.000000}\selectfont \(\displaystyle {\ensuremath{-}50}\)}%
\end{pgfscope}%
\begin{pgfscope}%
\definecolor{textcolor}{rgb}{0.000000,0.000000,0.000000}%
\pgfsetstrokecolor{textcolor}%
\pgfsetfillcolor{textcolor}%
\pgftext[x=3.363889in,y=1.020525in,,bottom,rotate=90.000000]{\color{textcolor}\rmfamily\fontsize{10.000000}{12.000000}\selectfont Amplitude (dB)}%
\end{pgfscope}%
\begin{pgfscope}%
\pgfpathrectangle{\pgfqpoint{3.833026in}{0.565123in}}{\pgfqpoint{2.282252in}{0.910803in}}%
\pgfusepath{clip}%
\pgfsetrectcap%
\pgfsetroundjoin%
\pgfsetlinewidth{1.003750pt}%
\definecolor{currentstroke}{rgb}{0.000000,0.000000,0.000000}%
\pgfsetstrokecolor{currentstroke}%
\pgfsetdash{}{0pt}%
\pgfpathmoveto{\pgfqpoint{3.833026in}{0.865662in}}%
\pgfpathlineto{\pgfqpoint{3.835308in}{0.880559in}}%
\pgfpathlineto{\pgfqpoint{3.837590in}{0.974594in}}%
\pgfpathlineto{\pgfqpoint{3.839872in}{0.880248in}}%
\pgfpathlineto{\pgfqpoint{3.842155in}{0.865042in}}%
\pgfpathlineto{\pgfqpoint{3.846719in}{0.864731in}}%
\pgfpathlineto{\pgfqpoint{3.849001in}{0.924007in}}%
\pgfpathlineto{\pgfqpoint{3.851284in}{0.923697in}}%
\pgfpathlineto{\pgfqpoint{3.853566in}{0.908490in}}%
\pgfpathlineto{\pgfqpoint{3.855848in}{0.923387in}}%
\pgfpathlineto{\pgfqpoint{3.860413in}{0.863179in}}%
\pgfpathlineto{\pgfqpoint{3.864977in}{0.897628in}}%
\pgfpathlineto{\pgfqpoint{3.867259in}{0.892352in}}%
\pgfpathlineto{\pgfqpoint{3.869542in}{0.922145in}}%
\pgfpathlineto{\pgfqpoint{3.871824in}{0.877145in}}%
\pgfpathlineto{\pgfqpoint{3.874106in}{0.861938in}}%
\pgfpathlineto{\pgfqpoint{3.878671in}{0.891421in}}%
\pgfpathlineto{\pgfqpoint{3.880953in}{0.891111in}}%
\pgfpathlineto{\pgfqpoint{3.883235in}{0.945731in}}%
\pgfpathlineto{\pgfqpoint{3.887800in}{0.890490in}}%
\pgfpathlineto{\pgfqpoint{3.890082in}{0.875593in}}%
\pgfpathlineto{\pgfqpoint{3.892364in}{0.924938in}}%
\pgfpathlineto{\pgfqpoint{3.894646in}{0.934559in}}%
\pgfpathlineto{\pgfqpoint{3.896929in}{0.934559in}}%
\pgfpathlineto{\pgfqpoint{3.899211in}{0.944180in}}%
\pgfpathlineto{\pgfqpoint{3.901493in}{0.904145in}}%
\pgfpathlineto{\pgfqpoint{3.903776in}{0.893904in}}%
\pgfpathlineto{\pgfqpoint{3.906058in}{0.859145in}}%
\pgfpathlineto{\pgfqpoint{3.908340in}{0.858835in}}%
\pgfpathlineto{\pgfqpoint{3.910622in}{0.972732in}}%
\pgfpathlineto{\pgfqpoint{3.912905in}{0.873421in}}%
\pgfpathlineto{\pgfqpoint{3.915187in}{0.888007in}}%
\pgfpathlineto{\pgfqpoint{3.917469in}{0.912524in}}%
\pgfpathlineto{\pgfqpoint{3.919751in}{0.907559in}}%
\pgfpathlineto{\pgfqpoint{3.922034in}{0.857593in}}%
\pgfpathlineto{\pgfqpoint{3.924316in}{0.956594in}}%
\pgfpathlineto{\pgfqpoint{3.926598in}{0.936731in}}%
\pgfpathlineto{\pgfqpoint{3.928880in}{0.886766in}}%
\pgfpathlineto{\pgfqpoint{3.931163in}{0.916249in}}%
\pgfpathlineto{\pgfqpoint{3.933445in}{0.921214in}}%
\pgfpathlineto{\pgfqpoint{3.935727in}{0.935800in}}%
\pgfpathlineto{\pgfqpoint{3.938009in}{0.915628in}}%
\pgfpathlineto{\pgfqpoint{3.940292in}{0.940456in}}%
\pgfpathlineto{\pgfqpoint{3.942574in}{0.870628in}}%
\pgfpathlineto{\pgfqpoint{3.944856in}{0.905076in}}%
\pgfpathlineto{\pgfqpoint{3.947138in}{0.904766in}}%
\pgfpathlineto{\pgfqpoint{3.951703in}{0.929283in}}%
\pgfpathlineto{\pgfqpoint{3.953985in}{0.948835in}}%
\pgfpathlineto{\pgfqpoint{3.956267in}{0.938904in}}%
\pgfpathlineto{\pgfqpoint{3.958550in}{0.978318in}}%
\pgfpathlineto{\pgfqpoint{3.960832in}{0.888628in}}%
\pgfpathlineto{\pgfqpoint{3.963114in}{0.888628in}}%
\pgfpathlineto{\pgfqpoint{3.965396in}{0.883352in}}%
\pgfpathlineto{\pgfqpoint{3.967679in}{0.853248in}}%
\pgfpathlineto{\pgfqpoint{3.969961in}{0.853248in}}%
\pgfpathlineto{\pgfqpoint{3.972243in}{0.882731in}}%
\pgfpathlineto{\pgfqpoint{3.976808in}{0.986697in}}%
\pgfpathlineto{\pgfqpoint{3.979090in}{0.897007in}}%
\pgfpathlineto{\pgfqpoint{3.983654in}{0.926490in}}%
\pgfpathlineto{\pgfqpoint{3.985937in}{0.881490in}}%
\pgfpathlineto{\pgfqpoint{3.988219in}{0.866283in}}%
\pgfpathlineto{\pgfqpoint{3.992783in}{0.865973in}}%
\pgfpathlineto{\pgfqpoint{3.995066in}{0.880559in}}%
\pgfpathlineto{\pgfqpoint{3.997348in}{0.850766in}}%
\pgfpathlineto{\pgfqpoint{4.001912in}{0.899800in}}%
\pgfpathlineto{\pgfqpoint{4.004195in}{0.924628in}}%
\pgfpathlineto{\pgfqpoint{4.006477in}{0.914387in}}%
\pgfpathlineto{\pgfqpoint{4.008759in}{1.023318in}}%
\pgfpathlineto{\pgfqpoint{4.011041in}{0.894214in}}%
\pgfpathlineto{\pgfqpoint{4.013324in}{1.206732in}}%
\pgfpathlineto{\pgfqpoint{4.015606in}{1.360353in}}%
\pgfpathlineto{\pgfqpoint{4.017888in}{1.424595in}}%
\pgfpathlineto{\pgfqpoint{4.020170in}{1.434526in}}%
\pgfpathlineto{\pgfqpoint{4.022453in}{1.399457in}}%
\pgfpathlineto{\pgfqpoint{4.024735in}{1.289905in}}%
\pgfpathlineto{\pgfqpoint{4.027017in}{1.001904in}}%
\pgfpathlineto{\pgfqpoint{4.031582in}{0.897007in}}%
\pgfpathlineto{\pgfqpoint{4.033864in}{1.100594in}}%
\pgfpathlineto{\pgfqpoint{4.036146in}{0.866904in}}%
\pgfpathlineto{\pgfqpoint{4.038428in}{0.896386in}}%
\pgfpathlineto{\pgfqpoint{4.040711in}{0.846731in}}%
\pgfpathlineto{\pgfqpoint{4.042993in}{0.846421in}}%
\pgfpathlineto{\pgfqpoint{4.045275in}{0.875904in}}%
\pgfpathlineto{\pgfqpoint{4.047557in}{0.935490in}}%
\pgfpathlineto{\pgfqpoint{4.049840in}{0.860697in}}%
\pgfpathlineto{\pgfqpoint{4.052122in}{0.860386in}}%
\pgfpathlineto{\pgfqpoint{4.054404in}{0.875283in}}%
\pgfpathlineto{\pgfqpoint{4.056686in}{0.845179in}}%
\pgfpathlineto{\pgfqpoint{4.081791in}{0.842697in}}%
\pgfpathlineto{\pgfqpoint{4.093202in}{0.841766in}}%
\pgfpathlineto{\pgfqpoint{4.116025in}{0.839593in}}%
\pgfpathlineto{\pgfqpoint{4.127436in}{0.838662in}}%
\pgfpathlineto{\pgfqpoint{4.145694in}{0.836800in}}%
\pgfpathlineto{\pgfqpoint{4.147976in}{0.807007in}}%
\pgfpathlineto{\pgfqpoint{4.150259in}{0.836490in}}%
\pgfpathlineto{\pgfqpoint{4.154823in}{0.836179in}}%
\pgfpathlineto{\pgfqpoint{4.159388in}{0.775972in}}%
\pgfpathlineto{\pgfqpoint{4.163952in}{0.775662in}}%
\pgfpathlineto{\pgfqpoint{4.166234in}{0.790248in}}%
\pgfpathlineto{\pgfqpoint{4.168517in}{0.775352in}}%
\pgfpathlineto{\pgfqpoint{4.170799in}{0.789938in}}%
\pgfpathlineto{\pgfqpoint{4.173081in}{0.774731in}}%
\pgfpathlineto{\pgfqpoint{4.175363in}{0.774421in}}%
\pgfpathlineto{\pgfqpoint{4.177646in}{0.789317in}}%
\pgfpathlineto{\pgfqpoint{4.179928in}{0.774110in}}%
\pgfpathlineto{\pgfqpoint{4.182210in}{0.788696in}}%
\pgfpathlineto{\pgfqpoint{4.184492in}{0.773800in}}%
\pgfpathlineto{\pgfqpoint{4.205033in}{0.771938in}}%
\pgfpathlineto{\pgfqpoint{4.207315in}{0.786524in}}%
\pgfpathlineto{\pgfqpoint{4.209597in}{0.771317in}}%
\pgfpathlineto{\pgfqpoint{4.218726in}{0.770696in}}%
\pgfpathlineto{\pgfqpoint{4.221009in}{0.785283in}}%
\pgfpathlineto{\pgfqpoint{4.223291in}{0.770076in}}%
\pgfpathlineto{\pgfqpoint{4.232420in}{0.769145in}}%
\pgfpathlineto{\pgfqpoint{4.234702in}{0.992594in}}%
\pgfpathlineto{\pgfqpoint{4.239267in}{0.773490in}}%
\pgfpathlineto{\pgfqpoint{4.241549in}{0.857904in}}%
\pgfpathlineto{\pgfqpoint{4.243831in}{0.802972in}}%
\pgfpathlineto{\pgfqpoint{4.248396in}{1.279353in}}%
\pgfpathlineto{\pgfqpoint{4.250678in}{1.353526in}}%
\pgfpathlineto{\pgfqpoint{4.252960in}{1.373078in}}%
\pgfpathlineto{\pgfqpoint{4.255242in}{1.338319in}}%
\pgfpathlineto{\pgfqpoint{4.257525in}{1.239008in}}%
\pgfpathlineto{\pgfqpoint{4.262089in}{0.841766in}}%
\pgfpathlineto{\pgfqpoint{4.264371in}{0.836800in}}%
\pgfpathlineto{\pgfqpoint{4.266654in}{0.866593in}}%
\pgfpathlineto{\pgfqpoint{4.268936in}{0.851697in}}%
\pgfpathlineto{\pgfqpoint{4.273500in}{0.752386in}}%
\pgfpathlineto{\pgfqpoint{4.275783in}{0.737489in}}%
\pgfpathlineto{\pgfqpoint{4.280347in}{0.881490in}}%
\pgfpathlineto{\pgfqpoint{4.282629in}{0.896386in}}%
\pgfpathlineto{\pgfqpoint{4.287194in}{0.737489in}}%
\pgfpathlineto{\pgfqpoint{4.300887in}{0.737489in}}%
\pgfpathlineto{\pgfqpoint{4.303170in}{0.752386in}}%
\pgfpathlineto{\pgfqpoint{4.305452in}{0.737489in}}%
\pgfpathlineto{\pgfqpoint{4.387613in}{0.737179in}}%
\pgfpathlineto{\pgfqpoint{4.389895in}{0.752076in}}%
\pgfpathlineto{\pgfqpoint{4.392177in}{0.737179in}}%
\pgfpathlineto{\pgfqpoint{4.394460in}{0.752076in}}%
\pgfpathlineto{\pgfqpoint{4.396742in}{0.737179in}}%
\pgfpathlineto{\pgfqpoint{4.405871in}{0.737179in}}%
\pgfpathlineto{\pgfqpoint{4.408153in}{0.757041in}}%
\pgfpathlineto{\pgfqpoint{4.410435in}{0.766972in}}%
\pgfpathlineto{\pgfqpoint{4.412718in}{0.736869in}}%
\pgfpathlineto{\pgfqpoint{4.428693in}{0.736869in}}%
\pgfpathlineto{\pgfqpoint{4.430976in}{0.766662in}}%
\pgfpathlineto{\pgfqpoint{4.433258in}{0.736869in}}%
\pgfpathlineto{\pgfqpoint{4.442387in}{0.736869in}}%
\pgfpathlineto{\pgfqpoint{4.444669in}{0.726938in}}%
\pgfpathlineto{\pgfqpoint{4.446951in}{0.736869in}}%
\pgfpathlineto{\pgfqpoint{4.451516in}{0.736869in}}%
\pgfpathlineto{\pgfqpoint{4.453798in}{0.751765in}}%
\pgfpathlineto{\pgfqpoint{4.456080in}{0.736869in}}%
\pgfpathlineto{\pgfqpoint{4.458363in}{0.736869in}}%
\pgfpathlineto{\pgfqpoint{4.460645in}{0.756731in}}%
\pgfpathlineto{\pgfqpoint{4.462927in}{0.736869in}}%
\pgfpathlineto{\pgfqpoint{4.465209in}{0.771627in}}%
\pgfpathlineto{\pgfqpoint{4.467492in}{0.736869in}}%
\pgfpathlineto{\pgfqpoint{4.469774in}{0.811352in}}%
\pgfpathlineto{\pgfqpoint{4.472056in}{0.781559in}}%
\pgfpathlineto{\pgfqpoint{4.474338in}{0.786524in}}%
\pgfpathlineto{\pgfqpoint{4.476621in}{0.751765in}}%
\pgfpathlineto{\pgfqpoint{4.481185in}{1.228456in}}%
\pgfpathlineto{\pgfqpoint{4.483467in}{1.317836in}}%
\pgfpathlineto{\pgfqpoint{4.485750in}{1.347629in}}%
\pgfpathlineto{\pgfqpoint{4.488032in}{1.317836in}}%
\pgfpathlineto{\pgfqpoint{4.490314in}{1.233422in}}%
\pgfpathlineto{\pgfqpoint{4.494879in}{0.801110in}}%
\pgfpathlineto{\pgfqpoint{4.497161in}{0.825938in}}%
\pgfpathlineto{\pgfqpoint{4.499443in}{0.801110in}}%
\pgfpathlineto{\pgfqpoint{4.501726in}{0.721662in}}%
\pgfpathlineto{\pgfqpoint{4.506290in}{0.751455in}}%
\pgfpathlineto{\pgfqpoint{4.508572in}{0.736558in}}%
\pgfpathlineto{\pgfqpoint{4.513137in}{0.845800in}}%
\pgfpathlineto{\pgfqpoint{4.515419in}{0.860697in}}%
\pgfpathlineto{\pgfqpoint{4.519984in}{0.751455in}}%
\pgfpathlineto{\pgfqpoint{4.522266in}{0.736558in}}%
\pgfpathlineto{\pgfqpoint{4.551935in}{0.736558in}}%
\pgfpathlineto{\pgfqpoint{4.554217in}{0.706765in}}%
\pgfpathlineto{\pgfqpoint{4.556500in}{0.736558in}}%
\pgfpathlineto{\pgfqpoint{4.558782in}{0.736558in}}%
\pgfpathlineto{\pgfqpoint{4.561064in}{0.706765in}}%
\pgfpathlineto{\pgfqpoint{4.563346in}{0.736558in}}%
\pgfpathlineto{\pgfqpoint{4.570193in}{0.736558in}}%
\pgfpathlineto{\pgfqpoint{4.572475in}{0.706765in}}%
\pgfpathlineto{\pgfqpoint{4.574758in}{0.736248in}}%
\pgfpathlineto{\pgfqpoint{4.588451in}{0.736248in}}%
\pgfpathlineto{\pgfqpoint{4.590733in}{0.706455in}}%
\pgfpathlineto{\pgfqpoint{4.593016in}{0.706455in}}%
\pgfpathlineto{\pgfqpoint{4.595298in}{0.736248in}}%
\pgfpathlineto{\pgfqpoint{4.597580in}{0.706455in}}%
\pgfpathlineto{\pgfqpoint{4.602145in}{0.706455in}}%
\pgfpathlineto{\pgfqpoint{4.604427in}{0.721351in}}%
\pgfpathlineto{\pgfqpoint{4.606709in}{0.706455in}}%
\pgfpathlineto{\pgfqpoint{4.608991in}{0.736248in}}%
\pgfpathlineto{\pgfqpoint{4.611274in}{0.706455in}}%
\pgfpathlineto{\pgfqpoint{4.613556in}{0.736248in}}%
\pgfpathlineto{\pgfqpoint{4.615838in}{0.706455in}}%
\pgfpathlineto{\pgfqpoint{4.618120in}{0.706455in}}%
\pgfpathlineto{\pgfqpoint{4.620403in}{0.736248in}}%
\pgfpathlineto{\pgfqpoint{4.622685in}{0.706455in}}%
\pgfpathlineto{\pgfqpoint{4.627249in}{0.706455in}}%
\pgfpathlineto{\pgfqpoint{4.629532in}{0.736248in}}%
\pgfpathlineto{\pgfqpoint{4.631814in}{0.706455in}}%
\pgfpathlineto{\pgfqpoint{4.640943in}{0.706455in}}%
\pgfpathlineto{\pgfqpoint{4.643225in}{0.736248in}}%
\pgfpathlineto{\pgfqpoint{4.645507in}{0.706455in}}%
\pgfpathlineto{\pgfqpoint{4.647790in}{0.736248in}}%
\pgfpathlineto{\pgfqpoint{4.650072in}{0.706455in}}%
\pgfpathlineto{\pgfqpoint{4.652354in}{0.706144in}}%
\pgfpathlineto{\pgfqpoint{4.654636in}{0.735938in}}%
\pgfpathlineto{\pgfqpoint{4.656919in}{0.735938in}}%
\pgfpathlineto{\pgfqpoint{4.659201in}{0.706144in}}%
\pgfpathlineto{\pgfqpoint{4.661483in}{0.735938in}}%
\pgfpathlineto{\pgfqpoint{4.663765in}{0.706144in}}%
\pgfpathlineto{\pgfqpoint{4.670612in}{0.706144in}}%
\pgfpathlineto{\pgfqpoint{4.672894in}{0.721041in}}%
\pgfpathlineto{\pgfqpoint{4.675177in}{0.750834in}}%
\pgfpathlineto{\pgfqpoint{4.679741in}{0.750834in}}%
\pgfpathlineto{\pgfqpoint{4.682023in}{0.706144in}}%
\pgfpathlineto{\pgfqpoint{4.686588in}{0.706144in}}%
\pgfpathlineto{\pgfqpoint{4.688870in}{0.676351in}}%
\pgfpathlineto{\pgfqpoint{4.691152in}{0.706144in}}%
\pgfpathlineto{\pgfqpoint{4.695717in}{0.706144in}}%
\pgfpathlineto{\pgfqpoint{4.697999in}{0.721041in}}%
\pgfpathlineto{\pgfqpoint{4.700281in}{0.765731in}}%
\pgfpathlineto{\pgfqpoint{4.702564in}{0.735938in}}%
\pgfpathlineto{\pgfqpoint{4.704846in}{0.785593in}}%
\pgfpathlineto{\pgfqpoint{4.707128in}{0.889869in}}%
\pgfpathlineto{\pgfqpoint{4.709410in}{0.765731in}}%
\pgfpathlineto{\pgfqpoint{4.713975in}{1.187801in}}%
\pgfpathlineto{\pgfqpoint{4.716257in}{1.287112in}}%
\pgfpathlineto{\pgfqpoint{4.718539in}{1.321871in}}%
\pgfpathlineto{\pgfqpoint{4.720822in}{1.292077in}}%
\pgfpathlineto{\pgfqpoint{4.723104in}{1.222560in}}%
\pgfpathlineto{\pgfqpoint{4.725386in}{1.043801in}}%
\pgfpathlineto{\pgfqpoint{4.727668in}{0.775662in}}%
\pgfpathlineto{\pgfqpoint{4.729951in}{0.730972in}}%
\pgfpathlineto{\pgfqpoint{4.732233in}{0.810421in}}%
\pgfpathlineto{\pgfqpoint{4.734515in}{0.745558in}}%
\pgfpathlineto{\pgfqpoint{4.736797in}{0.745558in}}%
\pgfpathlineto{\pgfqpoint{4.739080in}{0.730662in}}%
\pgfpathlineto{\pgfqpoint{4.743644in}{0.730662in}}%
\pgfpathlineto{\pgfqpoint{4.745926in}{0.805145in}}%
\pgfpathlineto{\pgfqpoint{4.748209in}{0.745558in}}%
\pgfpathlineto{\pgfqpoint{4.750491in}{0.730662in}}%
\pgfpathlineto{\pgfqpoint{4.848628in}{0.730351in}}%
\pgfpathlineto{\pgfqpoint{4.850910in}{0.715455in}}%
\pgfpathlineto{\pgfqpoint{4.853192in}{0.730351in}}%
\pgfpathlineto{\pgfqpoint{4.857757in}{0.730351in}}%
\pgfpathlineto{\pgfqpoint{4.860039in}{0.700558in}}%
\pgfpathlineto{\pgfqpoint{4.862321in}{0.730351in}}%
\pgfpathlineto{\pgfqpoint{4.866886in}{0.730351in}}%
\pgfpathlineto{\pgfqpoint{4.869168in}{0.700558in}}%
\pgfpathlineto{\pgfqpoint{4.873733in}{0.700558in}}%
\pgfpathlineto{\pgfqpoint{4.876015in}{0.730351in}}%
\pgfpathlineto{\pgfqpoint{4.903402in}{0.730041in}}%
\pgfpathlineto{\pgfqpoint{4.907966in}{0.700248in}}%
\pgfpathlineto{\pgfqpoint{4.912531in}{0.700248in}}%
\pgfpathlineto{\pgfqpoint{4.914813in}{0.730041in}}%
\pgfpathlineto{\pgfqpoint{4.917095in}{0.730041in}}%
\pgfpathlineto{\pgfqpoint{4.919378in}{0.700248in}}%
\pgfpathlineto{\pgfqpoint{4.921660in}{0.700248in}}%
\pgfpathlineto{\pgfqpoint{4.923942in}{0.730041in}}%
\pgfpathlineto{\pgfqpoint{4.926224in}{0.730041in}}%
\pgfpathlineto{\pgfqpoint{4.928507in}{0.700248in}}%
\pgfpathlineto{\pgfqpoint{4.933071in}{0.700248in}}%
\pgfpathlineto{\pgfqpoint{4.935353in}{0.744938in}}%
\pgfpathlineto{\pgfqpoint{4.937636in}{0.749903in}}%
\pgfpathlineto{\pgfqpoint{4.939918in}{0.774731in}}%
\pgfpathlineto{\pgfqpoint{4.942200in}{0.779696in}}%
\pgfpathlineto{\pgfqpoint{4.944482in}{0.844248in}}%
\pgfpathlineto{\pgfqpoint{4.946765in}{1.142180in}}%
\pgfpathlineto{\pgfqpoint{4.949047in}{1.261353in}}%
\pgfpathlineto{\pgfqpoint{4.951329in}{1.296112in}}%
\pgfpathlineto{\pgfqpoint{4.953611in}{1.286181in}}%
\pgfpathlineto{\pgfqpoint{4.955894in}{1.216663in}}%
\pgfpathlineto{\pgfqpoint{4.958176in}{1.062732in}}%
\pgfpathlineto{\pgfqpoint{4.960458in}{0.730041in}}%
\pgfpathlineto{\pgfqpoint{4.962740in}{0.759834in}}%
\pgfpathlineto{\pgfqpoint{4.965023in}{0.749903in}}%
\pgfpathlineto{\pgfqpoint{4.967305in}{0.764800in}}%
\pgfpathlineto{\pgfqpoint{4.969587in}{0.690317in}}%
\pgfpathlineto{\pgfqpoint{4.971869in}{0.804524in}}%
\pgfpathlineto{\pgfqpoint{4.974152in}{0.670144in}}%
\pgfpathlineto{\pgfqpoint{4.978716in}{0.774421in}}%
\pgfpathlineto{\pgfqpoint{4.980998in}{0.819110in}}%
\pgfpathlineto{\pgfqpoint{4.983281in}{0.734696in}}%
\pgfpathlineto{\pgfqpoint{4.987845in}{0.670144in}}%
\pgfpathlineto{\pgfqpoint{5.017514in}{0.670144in}}%
\pgfpathlineto{\pgfqpoint{5.019797in}{0.699938in}}%
\pgfpathlineto{\pgfqpoint{5.022079in}{0.699938in}}%
\pgfpathlineto{\pgfqpoint{5.024361in}{0.670144in}}%
\pgfpathlineto{\pgfqpoint{5.026643in}{0.685041in}}%
\pgfpathlineto{\pgfqpoint{5.028926in}{0.670144in}}%
\pgfpathlineto{\pgfqpoint{5.031208in}{0.685041in}}%
\pgfpathlineto{\pgfqpoint{5.033490in}{0.670144in}}%
\pgfpathlineto{\pgfqpoint{5.056313in}{0.669834in}}%
\pgfpathlineto{\pgfqpoint{5.058595in}{0.684731in}}%
\pgfpathlineto{\pgfqpoint{5.060877in}{0.669834in}}%
\pgfpathlineto{\pgfqpoint{5.063159in}{0.684731in}}%
\pgfpathlineto{\pgfqpoint{5.065442in}{0.669834in}}%
\pgfpathlineto{\pgfqpoint{5.067724in}{0.669834in}}%
\pgfpathlineto{\pgfqpoint{5.070006in}{0.699627in}}%
\pgfpathlineto{\pgfqpoint{5.072289in}{0.669834in}}%
\pgfpathlineto{\pgfqpoint{5.074571in}{0.684731in}}%
\pgfpathlineto{\pgfqpoint{5.076853in}{0.669834in}}%
\pgfpathlineto{\pgfqpoint{5.083700in}{0.669834in}}%
\pgfpathlineto{\pgfqpoint{5.085982in}{0.684731in}}%
\pgfpathlineto{\pgfqpoint{5.088264in}{0.669834in}}%
\pgfpathlineto{\pgfqpoint{5.090547in}{0.669834in}}%
\pgfpathlineto{\pgfqpoint{5.092829in}{0.699627in}}%
\pgfpathlineto{\pgfqpoint{5.095111in}{0.669834in}}%
\pgfpathlineto{\pgfqpoint{5.097393in}{0.669834in}}%
\pgfpathlineto{\pgfqpoint{5.099676in}{0.684731in}}%
\pgfpathlineto{\pgfqpoint{5.101958in}{0.669834in}}%
\pgfpathlineto{\pgfqpoint{5.122498in}{0.669834in}}%
\pgfpathlineto{\pgfqpoint{5.124780in}{0.714524in}}%
\pgfpathlineto{\pgfqpoint{5.127063in}{0.669834in}}%
\pgfpathlineto{\pgfqpoint{5.147603in}{0.669524in}}%
\pgfpathlineto{\pgfqpoint{5.149885in}{0.714213in}}%
\pgfpathlineto{\pgfqpoint{5.152167in}{0.684420in}}%
\pgfpathlineto{\pgfqpoint{5.154450in}{0.684420in}}%
\pgfpathlineto{\pgfqpoint{5.156732in}{0.669524in}}%
\pgfpathlineto{\pgfqpoint{5.159014in}{0.719179in}}%
\pgfpathlineto{\pgfqpoint{5.161296in}{0.669524in}}%
\pgfpathlineto{\pgfqpoint{5.163579in}{0.714213in}}%
\pgfpathlineto{\pgfqpoint{5.168143in}{0.684420in}}%
\pgfpathlineto{\pgfqpoint{5.170425in}{0.714213in}}%
\pgfpathlineto{\pgfqpoint{5.172708in}{0.888007in}}%
\pgfpathlineto{\pgfqpoint{5.174990in}{0.763869in}}%
\pgfpathlineto{\pgfqpoint{5.177272in}{0.734076in}}%
\pgfpathlineto{\pgfqpoint{5.179554in}{1.106491in}}%
\pgfpathlineto{\pgfqpoint{5.181837in}{1.225663in}}%
\pgfpathlineto{\pgfqpoint{5.184119in}{1.280284in}}%
\pgfpathlineto{\pgfqpoint{5.186401in}{1.275319in}}%
\pgfpathlineto{\pgfqpoint{5.188683in}{1.215732in}}%
\pgfpathlineto{\pgfqpoint{5.190966in}{1.081663in}}%
\pgfpathlineto{\pgfqpoint{5.193248in}{0.714213in}}%
\pgfpathlineto{\pgfqpoint{5.195530in}{0.714213in}}%
\pgfpathlineto{\pgfqpoint{5.197812in}{0.863179in}}%
\pgfpathlineto{\pgfqpoint{5.200095in}{0.753938in}}%
\pgfpathlineto{\pgfqpoint{5.202377in}{0.714213in}}%
\pgfpathlineto{\pgfqpoint{5.204659in}{0.704282in}}%
\pgfpathlineto{\pgfqpoint{5.206941in}{0.669524in}}%
\pgfpathlineto{\pgfqpoint{5.209224in}{0.669524in}}%
\pgfpathlineto{\pgfqpoint{5.211506in}{0.704282in}}%
\pgfpathlineto{\pgfqpoint{5.213788in}{0.788386in}}%
\pgfpathlineto{\pgfqpoint{5.218353in}{0.669213in}}%
\pgfpathlineto{\pgfqpoint{5.225199in}{0.669213in}}%
\pgfpathlineto{\pgfqpoint{5.227482in}{0.684110in}}%
\pgfpathlineto{\pgfqpoint{5.229764in}{0.669213in}}%
\pgfpathlineto{\pgfqpoint{5.236611in}{0.669213in}}%
\pgfpathlineto{\pgfqpoint{5.238893in}{0.699007in}}%
\pgfpathlineto{\pgfqpoint{5.241175in}{0.669213in}}%
\pgfpathlineto{\pgfqpoint{5.243457in}{0.669213in}}%
\pgfpathlineto{\pgfqpoint{5.245740in}{0.703972in}}%
\pgfpathlineto{\pgfqpoint{5.250304in}{0.669213in}}%
\pgfpathlineto{\pgfqpoint{5.257151in}{0.669213in}}%
\pgfpathlineto{\pgfqpoint{5.259433in}{0.684110in}}%
\pgfpathlineto{\pgfqpoint{5.261715in}{0.684110in}}%
\pgfpathlineto{\pgfqpoint{5.263998in}{0.669213in}}%
\pgfpathlineto{\pgfqpoint{5.266280in}{0.684110in}}%
\pgfpathlineto{\pgfqpoint{5.268562in}{0.669213in}}%
\pgfpathlineto{\pgfqpoint{5.273127in}{0.669213in}}%
\pgfpathlineto{\pgfqpoint{5.275409in}{0.684110in}}%
\pgfpathlineto{\pgfqpoint{5.277691in}{0.669213in}}%
\pgfpathlineto{\pgfqpoint{5.334747in}{0.668903in}}%
\pgfpathlineto{\pgfqpoint{5.337030in}{0.698696in}}%
\pgfpathlineto{\pgfqpoint{5.339312in}{0.668903in}}%
\pgfpathlineto{\pgfqpoint{5.341594in}{0.683800in}}%
\pgfpathlineto{\pgfqpoint{5.343876in}{0.683800in}}%
\pgfpathlineto{\pgfqpoint{5.346159in}{0.668903in}}%
\pgfpathlineto{\pgfqpoint{5.368981in}{0.668903in}}%
\pgfpathlineto{\pgfqpoint{5.371264in}{0.698696in}}%
\pgfpathlineto{\pgfqpoint{5.373546in}{0.668593in}}%
\pgfpathlineto{\pgfqpoint{5.378110in}{0.668593in}}%
\pgfpathlineto{\pgfqpoint{5.380393in}{0.683489in}}%
\pgfpathlineto{\pgfqpoint{5.382675in}{0.728179in}}%
\pgfpathlineto{\pgfqpoint{5.384957in}{0.668593in}}%
\pgfpathlineto{\pgfqpoint{5.387239in}{0.683489in}}%
\pgfpathlineto{\pgfqpoint{5.389522in}{0.668593in}}%
\pgfpathlineto{\pgfqpoint{5.394086in}{0.668593in}}%
\pgfpathlineto{\pgfqpoint{5.396368in}{0.703351in}}%
\pgfpathlineto{\pgfqpoint{5.398651in}{0.787765in}}%
\pgfpathlineto{\pgfqpoint{5.400933in}{0.698386in}}%
\pgfpathlineto{\pgfqpoint{5.403215in}{0.748041in}}%
\pgfpathlineto{\pgfqpoint{5.407780in}{0.698386in}}%
\pgfpathlineto{\pgfqpoint{5.410062in}{0.718248in}}%
\pgfpathlineto{\pgfqpoint{5.412344in}{1.060870in}}%
\pgfpathlineto{\pgfqpoint{5.414626in}{1.204870in}}%
\pgfpathlineto{\pgfqpoint{5.416909in}{1.264457in}}%
\pgfpathlineto{\pgfqpoint{5.419191in}{1.264457in}}%
\pgfpathlineto{\pgfqpoint{5.421473in}{1.219767in}}%
\pgfpathlineto{\pgfqpoint{5.423755in}{1.095628in}}%
\pgfpathlineto{\pgfqpoint{5.426038in}{0.728179in}}%
\pgfpathlineto{\pgfqpoint{5.428320in}{0.713282in}}%
\pgfpathlineto{\pgfqpoint{5.430602in}{0.718248in}}%
\pgfpathlineto{\pgfqpoint{5.432884in}{0.728179in}}%
\pgfpathlineto{\pgfqpoint{5.435167in}{0.728179in}}%
\pgfpathlineto{\pgfqpoint{5.437449in}{0.668593in}}%
\pgfpathlineto{\pgfqpoint{5.442013in}{0.668593in}}%
\pgfpathlineto{\pgfqpoint{5.446578in}{0.733145in}}%
\pgfpathlineto{\pgfqpoint{5.448860in}{0.777834in}}%
\pgfpathlineto{\pgfqpoint{5.451142in}{0.668593in}}%
\pgfpathlineto{\pgfqpoint{5.460271in}{0.668282in}}%
\pgfpathlineto{\pgfqpoint{5.462554in}{0.668282in}}%
\pgfpathlineto{\pgfqpoint{5.464836in}{0.683179in}}%
\pgfpathlineto{\pgfqpoint{5.467118in}{0.668282in}}%
\pgfpathlineto{\pgfqpoint{5.487658in}{0.668282in}}%
\pgfpathlineto{\pgfqpoint{5.489941in}{0.683179in}}%
\pgfpathlineto{\pgfqpoint{5.492223in}{0.668282in}}%
\pgfpathlineto{\pgfqpoint{5.510481in}{0.668282in}}%
\pgfpathlineto{\pgfqpoint{5.512763in}{0.683179in}}%
\pgfpathlineto{\pgfqpoint{5.515045in}{0.668282in}}%
\pgfpathlineto{\pgfqpoint{5.517328in}{0.698075in}}%
\pgfpathlineto{\pgfqpoint{5.519610in}{0.668282in}}%
\pgfpathlineto{\pgfqpoint{5.521892in}{0.683179in}}%
\pgfpathlineto{\pgfqpoint{5.524174in}{0.668282in}}%
\pgfpathlineto{\pgfqpoint{5.556126in}{0.667972in}}%
\pgfpathlineto{\pgfqpoint{5.558408in}{0.682869in}}%
\pgfpathlineto{\pgfqpoint{5.560690in}{0.667972in}}%
\pgfpathlineto{\pgfqpoint{5.601771in}{0.667972in}}%
\pgfpathlineto{\pgfqpoint{5.604053in}{0.697765in}}%
\pgfpathlineto{\pgfqpoint{5.608618in}{0.667972in}}%
\pgfpathlineto{\pgfqpoint{5.615464in}{0.667662in}}%
\pgfpathlineto{\pgfqpoint{5.617747in}{0.697455in}}%
\pgfpathlineto{\pgfqpoint{5.620029in}{0.667662in}}%
\pgfpathlineto{\pgfqpoint{5.629158in}{0.667662in}}%
\pgfpathlineto{\pgfqpoint{5.631440in}{0.682558in}}%
\pgfpathlineto{\pgfqpoint{5.633722in}{0.712351in}}%
\pgfpathlineto{\pgfqpoint{5.636005in}{0.712351in}}%
\pgfpathlineto{\pgfqpoint{5.638287in}{0.667662in}}%
\pgfpathlineto{\pgfqpoint{5.640569in}{0.697455in}}%
\pgfpathlineto{\pgfqpoint{5.642851in}{0.682558in}}%
\pgfpathlineto{\pgfqpoint{5.645134in}{1.020214in}}%
\pgfpathlineto{\pgfqpoint{5.647416in}{1.179112in}}%
\pgfpathlineto{\pgfqpoint{5.649698in}{1.238698in}}%
\pgfpathlineto{\pgfqpoint{5.651980in}{1.253595in}}%
\pgfpathlineto{\pgfqpoint{5.654263in}{1.213870in}}%
\pgfpathlineto{\pgfqpoint{5.656545in}{1.104629in}}%
\pgfpathlineto{\pgfqpoint{5.658827in}{0.816628in}}%
\pgfpathlineto{\pgfqpoint{5.661110in}{0.702420in}}%
\pgfpathlineto{\pgfqpoint{5.663392in}{0.712351in}}%
\pgfpathlineto{\pgfqpoint{5.665674in}{0.702420in}}%
\pgfpathlineto{\pgfqpoint{5.667956in}{0.727248in}}%
\pgfpathlineto{\pgfqpoint{5.672521in}{0.697455in}}%
\pgfpathlineto{\pgfqpoint{5.674803in}{0.707386in}}%
\pgfpathlineto{\pgfqpoint{5.677085in}{0.702420in}}%
\pgfpathlineto{\pgfqpoint{5.681650in}{0.771938in}}%
\pgfpathlineto{\pgfqpoint{5.683932in}{0.682558in}}%
\pgfpathlineto{\pgfqpoint{5.686214in}{0.667662in}}%
\pgfpathlineto{\pgfqpoint{5.690779in}{0.667662in}}%
\pgfpathlineto{\pgfqpoint{5.693061in}{0.652455in}}%
\pgfpathlineto{\pgfqpoint{5.695343in}{0.667351in}}%
\pgfpathlineto{\pgfqpoint{5.715884in}{0.667351in}}%
\pgfpathlineto{\pgfqpoint{5.718166in}{0.682248in}}%
\pgfpathlineto{\pgfqpoint{5.720448in}{0.667351in}}%
\pgfpathlineto{\pgfqpoint{5.752400in}{0.667351in}}%
\pgfpathlineto{\pgfqpoint{5.754682in}{0.637558in}}%
\pgfpathlineto{\pgfqpoint{5.756964in}{0.667351in}}%
\pgfpathlineto{\pgfqpoint{5.759246in}{0.667351in}}%
\pgfpathlineto{\pgfqpoint{5.761529in}{0.637558in}}%
\pgfpathlineto{\pgfqpoint{5.766093in}{0.637558in}}%
\pgfpathlineto{\pgfqpoint{5.768375in}{0.667351in}}%
\pgfpathlineto{\pgfqpoint{5.770658in}{0.667351in}}%
\pgfpathlineto{\pgfqpoint{5.772940in}{0.637248in}}%
\pgfpathlineto{\pgfqpoint{5.777504in}{0.637248in}}%
\pgfpathlineto{\pgfqpoint{5.779787in}{0.652144in}}%
\pgfpathlineto{\pgfqpoint{5.782069in}{0.637248in}}%
\pgfpathlineto{\pgfqpoint{5.791198in}{0.637248in}}%
\pgfpathlineto{\pgfqpoint{5.793480in}{0.667041in}}%
\pgfpathlineto{\pgfqpoint{5.795762in}{0.637248in}}%
\pgfpathlineto{\pgfqpoint{5.798045in}{0.637248in}}%
\pgfpathlineto{\pgfqpoint{5.800327in}{0.667041in}}%
\pgfpathlineto{\pgfqpoint{5.802609in}{0.637248in}}%
\pgfpathlineto{\pgfqpoint{5.804891in}{0.667041in}}%
\pgfpathlineto{\pgfqpoint{5.809456in}{0.667041in}}%
\pgfpathlineto{\pgfqpoint{5.811738in}{0.637248in}}%
\pgfpathlineto{\pgfqpoint{5.814020in}{0.667041in}}%
\pgfpathlineto{\pgfqpoint{5.820867in}{0.667041in}}%
\pgfpathlineto{\pgfqpoint{5.823149in}{0.652144in}}%
\pgfpathlineto{\pgfqpoint{5.825432in}{0.652144in}}%
\pgfpathlineto{\pgfqpoint{5.827714in}{0.637248in}}%
\pgfpathlineto{\pgfqpoint{5.829996in}{0.667041in}}%
\pgfpathlineto{\pgfqpoint{5.832278in}{0.667041in}}%
\pgfpathlineto{\pgfqpoint{5.834561in}{0.637248in}}%
\pgfpathlineto{\pgfqpoint{5.836843in}{0.667041in}}%
\pgfpathlineto{\pgfqpoint{5.839125in}{0.637248in}}%
\pgfpathlineto{\pgfqpoint{5.845972in}{0.637248in}}%
\pgfpathlineto{\pgfqpoint{5.848254in}{0.667041in}}%
\pgfpathlineto{\pgfqpoint{5.852819in}{0.636937in}}%
\pgfpathlineto{\pgfqpoint{5.859665in}{0.636937in}}%
\pgfpathlineto{\pgfqpoint{5.868794in}{0.701489in}}%
\pgfpathlineto{\pgfqpoint{5.871077in}{0.666731in}}%
\pgfpathlineto{\pgfqpoint{5.873359in}{0.681627in}}%
\pgfpathlineto{\pgfqpoint{5.880206in}{1.148387in}}%
\pgfpathlineto{\pgfqpoint{5.882488in}{1.242732in}}%
\pgfpathlineto{\pgfqpoint{5.884770in}{1.247698in}}%
\pgfpathlineto{\pgfqpoint{5.887052in}{1.212939in}}%
\pgfpathlineto{\pgfqpoint{5.889335in}{1.113629in}}%
\pgfpathlineto{\pgfqpoint{5.893899in}{0.691558in}}%
\pgfpathlineto{\pgfqpoint{5.896181in}{0.845490in}}%
\pgfpathlineto{\pgfqpoint{5.898464in}{0.681627in}}%
\pgfpathlineto{\pgfqpoint{5.900746in}{0.850455in}}%
\pgfpathlineto{\pgfqpoint{5.903028in}{0.671696in}}%
\pgfpathlineto{\pgfqpoint{5.905310in}{0.656799in}}%
\pgfpathlineto{\pgfqpoint{5.909875in}{0.701489in}}%
\pgfpathlineto{\pgfqpoint{5.912157in}{0.761076in}}%
\pgfpathlineto{\pgfqpoint{5.914439in}{0.785903in}}%
\pgfpathlineto{\pgfqpoint{5.921286in}{0.607144in}}%
\pgfpathlineto{\pgfqpoint{5.923568in}{0.641903in}}%
\pgfpathlineto{\pgfqpoint{5.925851in}{0.607144in}}%
\pgfpathlineto{\pgfqpoint{5.928133in}{0.607144in}}%
\pgfpathlineto{\pgfqpoint{5.930415in}{0.622041in}}%
\pgfpathlineto{\pgfqpoint{5.932697in}{0.606834in}}%
\pgfpathlineto{\pgfqpoint{5.934980in}{0.636627in}}%
\pgfpathlineto{\pgfqpoint{5.937262in}{0.606834in}}%
\pgfpathlineto{\pgfqpoint{5.939544in}{0.606834in}}%
\pgfpathlineto{\pgfqpoint{5.941826in}{0.621730in}}%
\pgfpathlineto{\pgfqpoint{5.944109in}{0.621730in}}%
\pgfpathlineto{\pgfqpoint{5.946391in}{0.606834in}}%
\pgfpathlineto{\pgfqpoint{5.948673in}{0.621730in}}%
\pgfpathlineto{\pgfqpoint{5.950955in}{0.606834in}}%
\pgfpathlineto{\pgfqpoint{5.955520in}{0.606834in}}%
\pgfpathlineto{\pgfqpoint{5.957802in}{0.621730in}}%
\pgfpathlineto{\pgfqpoint{5.960085in}{0.606834in}}%
\pgfpathlineto{\pgfqpoint{5.962367in}{0.621730in}}%
\pgfpathlineto{\pgfqpoint{5.964649in}{0.606834in}}%
\pgfpathlineto{\pgfqpoint{5.966931in}{0.606834in}}%
\pgfpathlineto{\pgfqpoint{5.969214in}{0.636627in}}%
\pgfpathlineto{\pgfqpoint{5.971496in}{0.621730in}}%
\pgfpathlineto{\pgfqpoint{5.973778in}{0.621730in}}%
\pgfpathlineto{\pgfqpoint{5.976060in}{0.636627in}}%
\pgfpathlineto{\pgfqpoint{5.978343in}{0.606834in}}%
\pgfpathlineto{\pgfqpoint{5.980625in}{0.666420in}}%
\pgfpathlineto{\pgfqpoint{5.985189in}{0.606834in}}%
\pgfpathlineto{\pgfqpoint{5.994318in}{0.606834in}}%
\pgfpathlineto{\pgfqpoint{5.996601in}{0.666420in}}%
\pgfpathlineto{\pgfqpoint{5.998883in}{0.606834in}}%
\pgfpathlineto{\pgfqpoint{6.001165in}{0.606834in}}%
\pgfpathlineto{\pgfqpoint{6.003447in}{0.621730in}}%
\pgfpathlineto{\pgfqpoint{6.005730in}{0.606834in}}%
\pgfpathlineto{\pgfqpoint{6.012576in}{0.606523in}}%
\pgfpathlineto{\pgfqpoint{6.014859in}{0.621420in}}%
\pgfpathlineto{\pgfqpoint{6.017141in}{0.606523in}}%
\pgfpathlineto{\pgfqpoint{6.019423in}{0.671075in}}%
\pgfpathlineto{\pgfqpoint{6.021705in}{0.606523in}}%
\pgfpathlineto{\pgfqpoint{6.026270in}{0.606523in}}%
\pgfpathlineto{\pgfqpoint{6.028552in}{0.621420in}}%
\pgfpathlineto{\pgfqpoint{6.030834in}{0.606523in}}%
\pgfpathlineto{\pgfqpoint{6.035399in}{0.606523in}}%
\pgfpathlineto{\pgfqpoint{6.037681in}{0.636317in}}%
\pgfpathlineto{\pgfqpoint{6.039963in}{0.651213in}}%
\pgfpathlineto{\pgfqpoint{6.042246in}{0.606523in}}%
\pgfpathlineto{\pgfqpoint{6.046810in}{0.606523in}}%
\pgfpathlineto{\pgfqpoint{6.049092in}{0.651213in}}%
\pgfpathlineto{\pgfqpoint{6.051375in}{0.621420in}}%
\pgfpathlineto{\pgfqpoint{6.053657in}{0.606523in}}%
\pgfpathlineto{\pgfqpoint{6.055939in}{0.621420in}}%
\pgfpathlineto{\pgfqpoint{6.058221in}{0.606523in}}%
\pgfpathlineto{\pgfqpoint{6.060504in}{0.606523in}}%
\pgfpathlineto{\pgfqpoint{6.062786in}{0.651213in}}%
\pgfpathlineto{\pgfqpoint{6.065068in}{0.666110in}}%
\pgfpathlineto{\pgfqpoint{6.067350in}{0.606523in}}%
\pgfpathlineto{\pgfqpoint{6.069633in}{0.651213in}}%
\pgfpathlineto{\pgfqpoint{6.071915in}{0.621420in}}%
\pgfpathlineto{\pgfqpoint{6.074197in}{0.606523in}}%
\pgfpathlineto{\pgfqpoint{6.078762in}{0.606523in}}%
\pgfpathlineto{\pgfqpoint{6.081044in}{0.636317in}}%
\pgfpathlineto{\pgfqpoint{6.083326in}{0.606523in}}%
\pgfpathlineto{\pgfqpoint{6.085608in}{0.606523in}}%
\pgfpathlineto{\pgfqpoint{6.087891in}{0.626386in}}%
\pgfpathlineto{\pgfqpoint{6.090173in}{0.621420in}}%
\pgfpathlineto{\pgfqpoint{6.092455in}{0.636006in}}%
\pgfpathlineto{\pgfqpoint{6.094737in}{0.621110in}}%
\pgfpathlineto{\pgfqpoint{6.097020in}{0.655868in}}%
\pgfpathlineto{\pgfqpoint{6.099302in}{0.670765in}}%
\pgfpathlineto{\pgfqpoint{6.101584in}{0.849524in}}%
\pgfpathlineto{\pgfqpoint{6.103866in}{0.670765in}}%
\pgfpathlineto{\pgfqpoint{6.106149in}{0.725386in}}%
\pgfpathlineto{\pgfqpoint{6.108431in}{0.665800in}}%
\pgfpathlineto{\pgfqpoint{6.112995in}{1.122629in}}%
\pgfpathlineto{\pgfqpoint{6.115278in}{1.212008in}}%
\pgfpathlineto{\pgfqpoint{6.115278in}{1.212008in}}%
\pgfusepath{stroke}%
\end{pgfscope}%
\begin{pgfscope}%
\pgfsetrectcap%
\pgfsetmiterjoin%
\pgfsetlinewidth{0.803000pt}%
\definecolor{currentstroke}{rgb}{0.000000,0.000000,0.000000}%
\pgfsetstrokecolor{currentstroke}%
\pgfsetdash{}{0pt}%
\pgfpathmoveto{\pgfqpoint{3.833026in}{0.565123in}}%
\pgfpathlineto{\pgfqpoint{3.833026in}{1.475926in}}%
\pgfusepath{stroke}%
\end{pgfscope}%
\begin{pgfscope}%
\pgfsetrectcap%
\pgfsetmiterjoin%
\pgfsetlinewidth{0.803000pt}%
\definecolor{currentstroke}{rgb}{0.000000,0.000000,0.000000}%
\pgfsetstrokecolor{currentstroke}%
\pgfsetdash{}{0pt}%
\pgfpathmoveto{\pgfqpoint{6.115278in}{0.565123in}}%
\pgfpathlineto{\pgfqpoint{6.115278in}{1.475926in}}%
\pgfusepath{stroke}%
\end{pgfscope}%
\begin{pgfscope}%
\pgfsetrectcap%
\pgfsetmiterjoin%
\pgfsetlinewidth{0.803000pt}%
\definecolor{currentstroke}{rgb}{0.000000,0.000000,0.000000}%
\pgfsetstrokecolor{currentstroke}%
\pgfsetdash{}{0pt}%
\pgfpathmoveto{\pgfqpoint{3.833026in}{0.565123in}}%
\pgfpathlineto{\pgfqpoint{6.115278in}{0.565123in}}%
\pgfusepath{stroke}%
\end{pgfscope}%
\begin{pgfscope}%
\pgfsetrectcap%
\pgfsetmiterjoin%
\pgfsetlinewidth{0.803000pt}%
\definecolor{currentstroke}{rgb}{0.000000,0.000000,0.000000}%
\pgfsetstrokecolor{currentstroke}%
\pgfsetdash{}{0pt}%
\pgfpathmoveto{\pgfqpoint{3.833026in}{1.475926in}}%
\pgfpathlineto{\pgfqpoint{6.115278in}{1.475926in}}%
\pgfusepath{stroke}%
\end{pgfscope}%
\begin{pgfscope}%
\definecolor{textcolor}{rgb}{0.000000,0.000000,0.000000}%
\pgfsetstrokecolor{textcolor}%
\pgfsetfillcolor{textcolor}%
\pgftext[x=4.974152in,y=1.559260in,,base]{\color{textcolor}\rmfamily\fontsize{12.000000}{14.400000}\selectfont Sawtooth Wave}%
\end{pgfscope}%
\end{pgfpicture}%
\makeatother%
\endgroup%

		\caption{Measured frequency content with a spectrum analyzer as a function of frequency of different signal shapes.}
		\label{fig:four_signals}
	\end{figure}
	
	\subsubsection{}
	
	Figure \ref{fig:444} shows different amplitude modulations of a 500 kHz carrier sine wave. The top two graphs display modulation frequencies of 50 kHz at modulation indexes of 70 and 20 \% and the bottom graphs show similar modulations but at a modulation frequency of 25 kHz.
	The frequency contents of the 50 kHz modulated signal are not always higher than the contents of the 25 kHz modulated signal. For the same carrier frequency the side bands occur at $f_\mathrm{carrier} \pm f_\mathrm{modulation}$. Therefore for 50 kHz the bands are at 450 kHz and 550 kHz and for 25 kHz the side bands are at 475 kHz and 525 kHz. Thus the lower side band of the 50 kHz signal is lower than the lower side band of the 25 kHz signal.
	
	The modulation index does not have an affect on the frequency of the side bands, but on the amplitudes of the bands.
	It is also observable that the modulation frequency has little affect on the relative amplitudes of the harmonics and mainly determines the frequency distance to the carrier wave.
		
	\begin{figure}[H]
		\centering
		%% Creator: Matplotlib, PGF backend
%%
%% To include the figure in your LaTeX document, write
%%   \input{<filename>.pgf}
%%
%% Make sure the required packages are loaded in your preamble
%%   \usepackage{pgf}
%%
%% Also ensure that all the required font packages are loaded; for instance,
%% the lmodern package is sometimes necessary when using math font.
%%   \usepackage{lmodern}
%%
%% Figures using additional raster images can only be included by \input if
%% they are in the same directory as the main LaTeX file. For loading figures
%% from other directories you can use the `import` package
%%   \usepackage{import}
%%
%% and then include the figures with
%%   \import{<path to file>}{<filename>.pgf}
%%
%% Matplotlib used the following preamble
%%
\begingroup%
\makeatletter%
\begin{pgfpicture}%
\pgfpathrectangle{\pgfpointorigin}{\pgfqpoint{6.300000in}{3.500000in}}%
\pgfusepath{use as bounding box, clip}%
\begin{pgfscope}%
\pgfsetbuttcap%
\pgfsetmiterjoin%
\definecolor{currentfill}{rgb}{1.000000,1.000000,1.000000}%
\pgfsetfillcolor{currentfill}%
\pgfsetlinewidth{0.000000pt}%
\definecolor{currentstroke}{rgb}{1.000000,1.000000,1.000000}%
\pgfsetstrokecolor{currentstroke}%
\pgfsetdash{}{0pt}%
\pgfpathmoveto{\pgfqpoint{0.000000in}{0.000000in}}%
\pgfpathlineto{\pgfqpoint{6.300000in}{0.000000in}}%
\pgfpathlineto{\pgfqpoint{6.300000in}{3.500000in}}%
\pgfpathlineto{\pgfqpoint{0.000000in}{3.500000in}}%
\pgfpathlineto{\pgfqpoint{0.000000in}{0.000000in}}%
\pgfpathclose%
\pgfusepath{fill}%
\end{pgfscope}%
\begin{pgfscope}%
\pgfsetbuttcap%
\pgfsetmiterjoin%
\definecolor{currentfill}{rgb}{1.000000,1.000000,1.000000}%
\pgfsetfillcolor{currentfill}%
\pgfsetlinewidth{0.000000pt}%
\definecolor{currentstroke}{rgb}{0.000000,0.000000,0.000000}%
\pgfsetstrokecolor{currentstroke}%
\pgfsetstrokeopacity{0.000000}%
\pgfsetdash{}{0pt}%
\pgfpathmoveto{\pgfqpoint{0.758026in}{2.240123in}}%
\pgfpathlineto{\pgfqpoint{2.866666in}{2.240123in}}%
\pgfpathlineto{\pgfqpoint{2.866666in}{3.141667in}}%
\pgfpathlineto{\pgfqpoint{0.758026in}{3.141667in}}%
\pgfpathlineto{\pgfqpoint{0.758026in}{2.240123in}}%
\pgfpathclose%
\pgfusepath{fill}%
\end{pgfscope}%
\begin{pgfscope}%
\pgfpathrectangle{\pgfqpoint{0.758026in}{2.240123in}}{\pgfqpoint{2.108640in}{0.901543in}}%
\pgfusepath{clip}%
\pgfsetbuttcap%
\pgfsetroundjoin%
\pgfsetlinewidth{0.501875pt}%
\definecolor{currentstroke}{rgb}{0.690196,0.690196,0.690196}%
\pgfsetstrokecolor{currentstroke}%
\pgfsetstrokeopacity{0.700000}%
\pgfsetdash{{1.850000pt}{0.800000pt}}{0.000000pt}%
\pgfpathmoveto{\pgfqpoint{0.758026in}{2.240123in}}%
\pgfpathlineto{\pgfqpoint{0.758026in}{3.141667in}}%
\pgfusepath{stroke}%
\end{pgfscope}%
\begin{pgfscope}%
\pgfsetbuttcap%
\pgfsetroundjoin%
\definecolor{currentfill}{rgb}{0.000000,0.000000,0.000000}%
\pgfsetfillcolor{currentfill}%
\pgfsetlinewidth{0.803000pt}%
\definecolor{currentstroke}{rgb}{0.000000,0.000000,0.000000}%
\pgfsetstrokecolor{currentstroke}%
\pgfsetdash{}{0pt}%
\pgfsys@defobject{currentmarker}{\pgfqpoint{0.000000in}{-0.048611in}}{\pgfqpoint{0.000000in}{0.000000in}}{%
\pgfpathmoveto{\pgfqpoint{0.000000in}{0.000000in}}%
\pgfpathlineto{\pgfqpoint{0.000000in}{-0.048611in}}%
\pgfusepath{stroke,fill}%
}%
\begin{pgfscope}%
\pgfsys@transformshift{0.758026in}{2.240123in}%
\pgfsys@useobject{currentmarker}{}%
\end{pgfscope}%
\end{pgfscope}%
\begin{pgfscope}%
\definecolor{textcolor}{rgb}{0.000000,0.000000,0.000000}%
\pgfsetstrokecolor{textcolor}%
\pgfsetfillcolor{textcolor}%
\pgftext[x=0.758026in,y=2.142901in,,top]{\color{textcolor}\rmfamily\fontsize{10.000000}{12.000000}\selectfont \(\displaystyle {400000}\)}%
\end{pgfscope}%
\begin{pgfscope}%
\pgfpathrectangle{\pgfqpoint{0.758026in}{2.240123in}}{\pgfqpoint{2.108640in}{0.901543in}}%
\pgfusepath{clip}%
\pgfsetbuttcap%
\pgfsetroundjoin%
\pgfsetlinewidth{0.501875pt}%
\definecolor{currentstroke}{rgb}{0.690196,0.690196,0.690196}%
\pgfsetstrokecolor{currentstroke}%
\pgfsetstrokeopacity{0.700000}%
\pgfsetdash{{1.850000pt}{0.800000pt}}{0.000000pt}%
\pgfpathmoveto{\pgfqpoint{1.285186in}{2.240123in}}%
\pgfpathlineto{\pgfqpoint{1.285186in}{3.141667in}}%
\pgfusepath{stroke}%
\end{pgfscope}%
\begin{pgfscope}%
\pgfsetbuttcap%
\pgfsetroundjoin%
\definecolor{currentfill}{rgb}{0.000000,0.000000,0.000000}%
\pgfsetfillcolor{currentfill}%
\pgfsetlinewidth{0.803000pt}%
\definecolor{currentstroke}{rgb}{0.000000,0.000000,0.000000}%
\pgfsetstrokecolor{currentstroke}%
\pgfsetdash{}{0pt}%
\pgfsys@defobject{currentmarker}{\pgfqpoint{0.000000in}{-0.048611in}}{\pgfqpoint{0.000000in}{0.000000in}}{%
\pgfpathmoveto{\pgfqpoint{0.000000in}{0.000000in}}%
\pgfpathlineto{\pgfqpoint{0.000000in}{-0.048611in}}%
\pgfusepath{stroke,fill}%
}%
\begin{pgfscope}%
\pgfsys@transformshift{1.285186in}{2.240123in}%
\pgfsys@useobject{currentmarker}{}%
\end{pgfscope}%
\end{pgfscope}%
\begin{pgfscope}%
\definecolor{textcolor}{rgb}{0.000000,0.000000,0.000000}%
\pgfsetstrokecolor{textcolor}%
\pgfsetfillcolor{textcolor}%
\pgftext[x=1.285186in,y=2.142901in,,top]{\color{textcolor}\rmfamily\fontsize{10.000000}{12.000000}\selectfont \(\displaystyle {450000}\)}%
\end{pgfscope}%
\begin{pgfscope}%
\pgfpathrectangle{\pgfqpoint{0.758026in}{2.240123in}}{\pgfqpoint{2.108640in}{0.901543in}}%
\pgfusepath{clip}%
\pgfsetbuttcap%
\pgfsetroundjoin%
\pgfsetlinewidth{0.501875pt}%
\definecolor{currentstroke}{rgb}{0.690196,0.690196,0.690196}%
\pgfsetstrokecolor{currentstroke}%
\pgfsetstrokeopacity{0.700000}%
\pgfsetdash{{1.850000pt}{0.800000pt}}{0.000000pt}%
\pgfpathmoveto{\pgfqpoint{1.812346in}{2.240123in}}%
\pgfpathlineto{\pgfqpoint{1.812346in}{3.141667in}}%
\pgfusepath{stroke}%
\end{pgfscope}%
\begin{pgfscope}%
\pgfsetbuttcap%
\pgfsetroundjoin%
\definecolor{currentfill}{rgb}{0.000000,0.000000,0.000000}%
\pgfsetfillcolor{currentfill}%
\pgfsetlinewidth{0.803000pt}%
\definecolor{currentstroke}{rgb}{0.000000,0.000000,0.000000}%
\pgfsetstrokecolor{currentstroke}%
\pgfsetdash{}{0pt}%
\pgfsys@defobject{currentmarker}{\pgfqpoint{0.000000in}{-0.048611in}}{\pgfqpoint{0.000000in}{0.000000in}}{%
\pgfpathmoveto{\pgfqpoint{0.000000in}{0.000000in}}%
\pgfpathlineto{\pgfqpoint{0.000000in}{-0.048611in}}%
\pgfusepath{stroke,fill}%
}%
\begin{pgfscope}%
\pgfsys@transformshift{1.812346in}{2.240123in}%
\pgfsys@useobject{currentmarker}{}%
\end{pgfscope}%
\end{pgfscope}%
\begin{pgfscope}%
\definecolor{textcolor}{rgb}{0.000000,0.000000,0.000000}%
\pgfsetstrokecolor{textcolor}%
\pgfsetfillcolor{textcolor}%
\pgftext[x=1.812346in,y=2.142901in,,top]{\color{textcolor}\rmfamily\fontsize{10.000000}{12.000000}\selectfont \(\displaystyle {500000}\)}%
\end{pgfscope}%
\begin{pgfscope}%
\pgfpathrectangle{\pgfqpoint{0.758026in}{2.240123in}}{\pgfqpoint{2.108640in}{0.901543in}}%
\pgfusepath{clip}%
\pgfsetbuttcap%
\pgfsetroundjoin%
\pgfsetlinewidth{0.501875pt}%
\definecolor{currentstroke}{rgb}{0.690196,0.690196,0.690196}%
\pgfsetstrokecolor{currentstroke}%
\pgfsetstrokeopacity{0.700000}%
\pgfsetdash{{1.850000pt}{0.800000pt}}{0.000000pt}%
\pgfpathmoveto{\pgfqpoint{2.339506in}{2.240123in}}%
\pgfpathlineto{\pgfqpoint{2.339506in}{3.141667in}}%
\pgfusepath{stroke}%
\end{pgfscope}%
\begin{pgfscope}%
\pgfsetbuttcap%
\pgfsetroundjoin%
\definecolor{currentfill}{rgb}{0.000000,0.000000,0.000000}%
\pgfsetfillcolor{currentfill}%
\pgfsetlinewidth{0.803000pt}%
\definecolor{currentstroke}{rgb}{0.000000,0.000000,0.000000}%
\pgfsetstrokecolor{currentstroke}%
\pgfsetdash{}{0pt}%
\pgfsys@defobject{currentmarker}{\pgfqpoint{0.000000in}{-0.048611in}}{\pgfqpoint{0.000000in}{0.000000in}}{%
\pgfpathmoveto{\pgfqpoint{0.000000in}{0.000000in}}%
\pgfpathlineto{\pgfqpoint{0.000000in}{-0.048611in}}%
\pgfusepath{stroke,fill}%
}%
\begin{pgfscope}%
\pgfsys@transformshift{2.339506in}{2.240123in}%
\pgfsys@useobject{currentmarker}{}%
\end{pgfscope}%
\end{pgfscope}%
\begin{pgfscope}%
\definecolor{textcolor}{rgb}{0.000000,0.000000,0.000000}%
\pgfsetstrokecolor{textcolor}%
\pgfsetfillcolor{textcolor}%
\pgftext[x=2.339506in,y=2.142901in,,top]{\color{textcolor}\rmfamily\fontsize{10.000000}{12.000000}\selectfont \(\displaystyle {550000}\)}%
\end{pgfscope}%
\begin{pgfscope}%
\pgfpathrectangle{\pgfqpoint{0.758026in}{2.240123in}}{\pgfqpoint{2.108640in}{0.901543in}}%
\pgfusepath{clip}%
\pgfsetbuttcap%
\pgfsetroundjoin%
\pgfsetlinewidth{0.501875pt}%
\definecolor{currentstroke}{rgb}{0.690196,0.690196,0.690196}%
\pgfsetstrokecolor{currentstroke}%
\pgfsetstrokeopacity{0.700000}%
\pgfsetdash{{1.850000pt}{0.800000pt}}{0.000000pt}%
\pgfpathmoveto{\pgfqpoint{2.866666in}{2.240123in}}%
\pgfpathlineto{\pgfqpoint{2.866666in}{3.141667in}}%
\pgfusepath{stroke}%
\end{pgfscope}%
\begin{pgfscope}%
\pgfsetbuttcap%
\pgfsetroundjoin%
\definecolor{currentfill}{rgb}{0.000000,0.000000,0.000000}%
\pgfsetfillcolor{currentfill}%
\pgfsetlinewidth{0.803000pt}%
\definecolor{currentstroke}{rgb}{0.000000,0.000000,0.000000}%
\pgfsetstrokecolor{currentstroke}%
\pgfsetdash{}{0pt}%
\pgfsys@defobject{currentmarker}{\pgfqpoint{0.000000in}{-0.048611in}}{\pgfqpoint{0.000000in}{0.000000in}}{%
\pgfpathmoveto{\pgfqpoint{0.000000in}{0.000000in}}%
\pgfpathlineto{\pgfqpoint{0.000000in}{-0.048611in}}%
\pgfusepath{stroke,fill}%
}%
\begin{pgfscope}%
\pgfsys@transformshift{2.866666in}{2.240123in}%
\pgfsys@useobject{currentmarker}{}%
\end{pgfscope}%
\end{pgfscope}%
\begin{pgfscope}%
\definecolor{textcolor}{rgb}{0.000000,0.000000,0.000000}%
\pgfsetstrokecolor{textcolor}%
\pgfsetfillcolor{textcolor}%
\pgftext[x=2.866666in,y=2.142901in,,top]{\color{textcolor}\rmfamily\fontsize{10.000000}{12.000000}\selectfont \(\displaystyle {600000}\)}%
\end{pgfscope}%
\begin{pgfscope}%
\definecolor{textcolor}{rgb}{0.000000,0.000000,0.000000}%
\pgfsetstrokecolor{textcolor}%
\pgfsetfillcolor{textcolor}%
\pgftext[x=1.812346in,y=1.963889in,,top]{\color{textcolor}\rmfamily\fontsize{10.000000}{12.000000}\selectfont Frequency (Hz)}%
\end{pgfscope}%
\begin{pgfscope}%
\pgfpathrectangle{\pgfqpoint{0.758026in}{2.240123in}}{\pgfqpoint{2.108640in}{0.901543in}}%
\pgfusepath{clip}%
\pgfsetbuttcap%
\pgfsetroundjoin%
\pgfsetlinewidth{0.501875pt}%
\definecolor{currentstroke}{rgb}{0.690196,0.690196,0.690196}%
\pgfsetstrokecolor{currentstroke}%
\pgfsetstrokeopacity{0.700000}%
\pgfsetdash{{1.850000pt}{0.800000pt}}{0.000000pt}%
\pgfpathmoveto{\pgfqpoint{0.758026in}{2.371142in}}%
\pgfpathlineto{\pgfqpoint{2.866666in}{2.371142in}}%
\pgfusepath{stroke}%
\end{pgfscope}%
\begin{pgfscope}%
\pgfsetbuttcap%
\pgfsetroundjoin%
\definecolor{currentfill}{rgb}{0.000000,0.000000,0.000000}%
\pgfsetfillcolor{currentfill}%
\pgfsetlinewidth{0.803000pt}%
\definecolor{currentstroke}{rgb}{0.000000,0.000000,0.000000}%
\pgfsetstrokecolor{currentstroke}%
\pgfsetdash{}{0pt}%
\pgfsys@defobject{currentmarker}{\pgfqpoint{-0.048611in}{0.000000in}}{\pgfqpoint{-0.000000in}{0.000000in}}{%
\pgfpathmoveto{\pgfqpoint{-0.000000in}{0.000000in}}%
\pgfpathlineto{\pgfqpoint{-0.048611in}{0.000000in}}%
\pgfusepath{stroke,fill}%
}%
\begin{pgfscope}%
\pgfsys@transformshift{0.758026in}{2.371142in}%
\pgfsys@useobject{currentmarker}{}%
\end{pgfscope}%
\end{pgfscope}%
\begin{pgfscope}%
\definecolor{textcolor}{rgb}{0.000000,0.000000,0.000000}%
\pgfsetstrokecolor{textcolor}%
\pgfsetfillcolor{textcolor}%
\pgftext[x=0.344444in, y=2.322917in, left, base]{\color{textcolor}\rmfamily\fontsize{10.000000}{12.000000}\selectfont \(\displaystyle {\ensuremath{-}100}\)}%
\end{pgfscope}%
\begin{pgfscope}%
\pgfpathrectangle{\pgfqpoint{0.758026in}{2.240123in}}{\pgfqpoint{2.108640in}{0.901543in}}%
\pgfusepath{clip}%
\pgfsetbuttcap%
\pgfsetroundjoin%
\pgfsetlinewidth{0.501875pt}%
\definecolor{currentstroke}{rgb}{0.690196,0.690196,0.690196}%
\pgfsetstrokecolor{currentstroke}%
\pgfsetstrokeopacity{0.700000}%
\pgfsetdash{{1.850000pt}{0.800000pt}}{0.000000pt}%
\pgfpathmoveto{\pgfqpoint{0.758026in}{2.880201in}}%
\pgfpathlineto{\pgfqpoint{2.866666in}{2.880201in}}%
\pgfusepath{stroke}%
\end{pgfscope}%
\begin{pgfscope}%
\pgfsetbuttcap%
\pgfsetroundjoin%
\definecolor{currentfill}{rgb}{0.000000,0.000000,0.000000}%
\pgfsetfillcolor{currentfill}%
\pgfsetlinewidth{0.803000pt}%
\definecolor{currentstroke}{rgb}{0.000000,0.000000,0.000000}%
\pgfsetstrokecolor{currentstroke}%
\pgfsetdash{}{0pt}%
\pgfsys@defobject{currentmarker}{\pgfqpoint{-0.048611in}{0.000000in}}{\pgfqpoint{-0.000000in}{0.000000in}}{%
\pgfpathmoveto{\pgfqpoint{-0.000000in}{0.000000in}}%
\pgfpathlineto{\pgfqpoint{-0.048611in}{0.000000in}}%
\pgfusepath{stroke,fill}%
}%
\begin{pgfscope}%
\pgfsys@transformshift{0.758026in}{2.880201in}%
\pgfsys@useobject{currentmarker}{}%
\end{pgfscope}%
\end{pgfscope}%
\begin{pgfscope}%
\definecolor{textcolor}{rgb}{0.000000,0.000000,0.000000}%
\pgfsetstrokecolor{textcolor}%
\pgfsetfillcolor{textcolor}%
\pgftext[x=0.413889in, y=2.831976in, left, base]{\color{textcolor}\rmfamily\fontsize{10.000000}{12.000000}\selectfont \(\displaystyle {\ensuremath{-}50}\)}%
\end{pgfscope}%
\begin{pgfscope}%
\definecolor{textcolor}{rgb}{0.000000,0.000000,0.000000}%
\pgfsetstrokecolor{textcolor}%
\pgfsetfillcolor{textcolor}%
\pgftext[x=0.288889in,y=2.690895in,,bottom,rotate=90.000000]{\color{textcolor}\rmfamily\fontsize{10.000000}{12.000000}\selectfont Amplitude (dB)}%
\end{pgfscope}%
\begin{pgfscope}%
\pgfpathrectangle{\pgfqpoint{0.758026in}{2.240123in}}{\pgfqpoint{2.108640in}{0.901543in}}%
\pgfusepath{clip}%
\pgfsetrectcap%
\pgfsetroundjoin%
\pgfsetlinewidth{1.003750pt}%
\definecolor{currentstroke}{rgb}{0.000000,0.000000,0.000000}%
\pgfsetstrokecolor{currentstroke}%
\pgfsetdash{}{0pt}%
\pgfpathmoveto{\pgfqpoint{0.754863in}{2.718257in}}%
\pgfpathlineto{\pgfqpoint{0.764352in}{2.570630in}}%
\pgfpathlineto{\pgfqpoint{0.773840in}{2.596083in}}%
\pgfpathlineto{\pgfqpoint{0.783329in}{2.718257in}}%
\pgfpathlineto{\pgfqpoint{0.792818in}{2.550267in}}%
\pgfpathlineto{\pgfqpoint{0.802307in}{2.621218in}}%
\pgfpathlineto{\pgfqpoint{0.811796in}{2.733211in}}%
\pgfpathlineto{\pgfqpoint{0.821285in}{2.514315in}}%
\pgfpathlineto{\pgfqpoint{0.830774in}{2.539768in}}%
\pgfpathlineto{\pgfqpoint{0.840263in}{2.672123in}}%
\pgfpathlineto{\pgfqpoint{0.849752in}{2.631399in}}%
\pgfpathlineto{\pgfqpoint{0.859240in}{2.575402in}}%
\pgfpathlineto{\pgfqpoint{0.868729in}{2.565221in}}%
\pgfpathlineto{\pgfqpoint{0.878218in}{2.534359in}}%
\pgfpathlineto{\pgfqpoint{0.887707in}{2.483453in}}%
\pgfpathlineto{\pgfqpoint{0.897196in}{2.580175in}}%
\pgfpathlineto{\pgfqpoint{0.906685in}{2.554722in}}%
\pgfpathlineto{\pgfqpoint{0.916174in}{2.569994in}}%
\pgfpathlineto{\pgfqpoint{0.925663in}{2.559812in}}%
\pgfpathlineto{\pgfqpoint{0.935151in}{2.569994in}}%
\pgfpathlineto{\pgfqpoint{0.944640in}{2.534359in}}%
\pgfpathlineto{\pgfqpoint{0.954129in}{2.518769in}}%
\pgfpathlineto{\pgfqpoint{0.963618in}{2.554404in}}%
\pgfpathlineto{\pgfqpoint{0.973107in}{2.539132in}}%
\pgfpathlineto{\pgfqpoint{0.982596in}{2.584947in}}%
\pgfpathlineto{\pgfqpoint{0.992085in}{2.564585in}}%
\pgfpathlineto{\pgfqpoint{1.001574in}{2.564585in}}%
\pgfpathlineto{\pgfqpoint{1.011063in}{2.478045in}}%
\pgfpathlineto{\pgfqpoint{1.020551in}{2.569675in}}%
\pgfpathlineto{\pgfqpoint{1.030040in}{2.589720in}}%
\pgfpathlineto{\pgfqpoint{1.039529in}{2.543904in}}%
\pgfpathlineto{\pgfqpoint{1.049018in}{2.528632in}}%
\pgfpathlineto{\pgfqpoint{1.058507in}{2.554085in}}%
\pgfpathlineto{\pgfqpoint{1.067996in}{2.564267in}}%
\pgfpathlineto{\pgfqpoint{1.077485in}{2.543904in}}%
\pgfpathlineto{\pgfqpoint{1.086974in}{2.528632in}}%
\pgfpathlineto{\pgfqpoint{1.096462in}{2.548995in}}%
\pgfpathlineto{\pgfqpoint{1.105951in}{2.599583in}}%
\pgfpathlineto{\pgfqpoint{1.115440in}{2.538495in}}%
\pgfpathlineto{\pgfqpoint{1.124929in}{2.487590in}}%
\pgfpathlineto{\pgfqpoint{1.134418in}{2.589401in}}%
\pgfpathlineto{\pgfqpoint{1.143907in}{2.558858in}}%
\pgfpathlineto{\pgfqpoint{1.153396in}{2.538495in}}%
\pgfpathlineto{\pgfqpoint{1.162885in}{2.502861in}}%
\pgfpathlineto{\pgfqpoint{1.172374in}{2.594492in}}%
\pgfpathlineto{\pgfqpoint{1.181862in}{2.507634in}}%
\pgfpathlineto{\pgfqpoint{1.191351in}{2.441456in}}%
\pgfpathlineto{\pgfqpoint{1.200840in}{2.507634in}}%
\pgfpathlineto{\pgfqpoint{1.210329in}{2.497453in}}%
\pgfpathlineto{\pgfqpoint{1.219818in}{2.497453in}}%
\pgfpathlineto{\pgfqpoint{1.229307in}{2.522906in}}%
\pgfpathlineto{\pgfqpoint{1.238796in}{2.645080in}}%
\pgfpathlineto{\pgfqpoint{1.248285in}{2.807660in}}%
\pgfpathlineto{\pgfqpoint{1.257773in}{2.894200in}}%
\pgfpathlineto{\pgfqpoint{1.267262in}{2.955287in}}%
\pgfpathlineto{\pgfqpoint{1.276751in}{2.990922in}}%
\pgfpathlineto{\pgfqpoint{1.286240in}{3.011284in}}%
\pgfpathlineto{\pgfqpoint{1.295729in}{3.006193in}}%
\pgfpathlineto{\pgfqpoint{1.305218in}{2.980740in}}%
\pgfpathlineto{\pgfqpoint{1.314707in}{2.934925in}}%
\pgfpathlineto{\pgfqpoint{1.324196in}{2.868429in}}%
\pgfpathlineto{\pgfqpoint{1.333685in}{2.761527in}}%
\pgfpathlineto{\pgfqpoint{1.343173in}{2.578266in}}%
\pgfpathlineto{\pgfqpoint{1.352662in}{2.496816in}}%
\pgfpathlineto{\pgfqpoint{1.362151in}{2.588447in}}%
\pgfpathlineto{\pgfqpoint{1.371640in}{2.512088in}}%
\pgfpathlineto{\pgfqpoint{1.381129in}{2.562994in}}%
\pgfpathlineto{\pgfqpoint{1.390618in}{2.527360in}}%
\pgfpathlineto{\pgfqpoint{1.400107in}{2.511770in}}%
\pgfpathlineto{\pgfqpoint{1.409596in}{2.511770in}}%
\pgfpathlineto{\pgfqpoint{1.419084in}{2.527042in}}%
\pgfpathlineto{\pgfqpoint{1.428573in}{2.552495in}}%
\pgfpathlineto{\pgfqpoint{1.438062in}{2.471045in}}%
\pgfpathlineto{\pgfqpoint{1.447551in}{2.547404in}}%
\pgfpathlineto{\pgfqpoint{1.457040in}{2.496498in}}%
\pgfpathlineto{\pgfqpoint{1.466529in}{2.537223in}}%
\pgfpathlineto{\pgfqpoint{1.476018in}{2.597992in}}%
\pgfpathlineto{\pgfqpoint{1.485507in}{2.475818in}}%
\pgfpathlineto{\pgfqpoint{1.494995in}{2.557267in}}%
\pgfpathlineto{\pgfqpoint{1.504484in}{2.562358in}}%
\pgfpathlineto{\pgfqpoint{1.513973in}{2.552176in}}%
\pgfpathlineto{\pgfqpoint{1.523462in}{2.521633in}}%
\pgfpathlineto{\pgfqpoint{1.532951in}{2.562358in}}%
\pgfpathlineto{\pgfqpoint{1.542440in}{2.526724in}}%
\pgfpathlineto{\pgfqpoint{1.551929in}{2.475499in}}%
\pgfpathlineto{\pgfqpoint{1.561418in}{2.516224in}}%
\pgfpathlineto{\pgfqpoint{1.570907in}{2.521315in}}%
\pgfpathlineto{\pgfqpoint{1.580395in}{2.582402in}}%
\pgfpathlineto{\pgfqpoint{1.589884in}{2.541677in}}%
\pgfpathlineto{\pgfqpoint{1.599373in}{2.485681in}}%
\pgfpathlineto{\pgfqpoint{1.608862in}{2.521315in}}%
\pgfpathlineto{\pgfqpoint{1.618351in}{2.526405in}}%
\pgfpathlineto{\pgfqpoint{1.627840in}{2.536268in}}%
\pgfpathlineto{\pgfqpoint{1.637329in}{2.587174in}}%
\pgfpathlineto{\pgfqpoint{1.646818in}{2.424275in}}%
\pgfpathlineto{\pgfqpoint{1.665795in}{2.551540in}}%
\pgfpathlineto{\pgfqpoint{1.675284in}{2.505725in}}%
\pgfpathlineto{\pgfqpoint{1.684773in}{2.515906in}}%
\pgfpathlineto{\pgfqpoint{1.694262in}{2.485362in}}%
\pgfpathlineto{\pgfqpoint{1.703751in}{2.485044in}}%
\pgfpathlineto{\pgfqpoint{1.713240in}{2.479954in}}%
\pgfpathlineto{\pgfqpoint{1.722729in}{2.551222in}}%
\pgfpathlineto{\pgfqpoint{1.732218in}{2.551222in}}%
\pgfpathlineto{\pgfqpoint{1.741706in}{2.500316in}}%
\pgfpathlineto{\pgfqpoint{1.751195in}{2.525769in}}%
\pgfpathlineto{\pgfqpoint{1.760684in}{2.541041in}}%
\pgfpathlineto{\pgfqpoint{1.770173in}{2.810842in}}%
\pgfpathlineto{\pgfqpoint{1.779662in}{2.937789in}}%
\pgfpathlineto{\pgfqpoint{1.789151in}{3.019238in}}%
\pgfpathlineto{\pgfqpoint{1.798640in}{3.065053in}}%
\pgfpathlineto{\pgfqpoint{1.808129in}{3.095597in}}%
\pgfpathlineto{\pgfqpoint{1.817617in}{3.100687in}}%
\pgfpathlineto{\pgfqpoint{1.827106in}{3.090506in}}%
\pgfpathlineto{\pgfqpoint{1.836595in}{3.059963in}}%
\pgfpathlineto{\pgfqpoint{1.846084in}{3.009057in}}%
\pgfpathlineto{\pgfqpoint{1.855573in}{2.922199in}}%
\pgfpathlineto{\pgfqpoint{1.865062in}{2.784753in}}%
\pgfpathlineto{\pgfqpoint{1.874551in}{2.509861in}}%
\pgfpathlineto{\pgfqpoint{1.884040in}{2.545495in}}%
\pgfpathlineto{\pgfqpoint{1.893529in}{2.504770in}}%
\pgfpathlineto{\pgfqpoint{1.903017in}{2.494589in}}%
\pgfpathlineto{\pgfqpoint{1.912506in}{2.514952in}}%
\pgfpathlineto{\pgfqpoint{1.921995in}{2.520042in}}%
\pgfpathlineto{\pgfqpoint{1.931484in}{2.570630in}}%
\pgfpathlineto{\pgfqpoint{1.940973in}{2.478999in}}%
\pgfpathlineto{\pgfqpoint{1.950462in}{2.540086in}}%
\pgfpathlineto{\pgfqpoint{1.959951in}{2.524815in}}%
\pgfpathlineto{\pgfqpoint{1.969440in}{2.555358in}}%
\pgfpathlineto{\pgfqpoint{1.978928in}{2.504452in}}%
\pgfpathlineto{\pgfqpoint{1.988417in}{2.540086in}}%
\pgfpathlineto{\pgfqpoint{1.997906in}{2.555358in}}%
\pgfpathlineto{\pgfqpoint{2.007395in}{2.509225in}}%
\pgfpathlineto{\pgfqpoint{2.016884in}{2.443047in}}%
\pgfpathlineto{\pgfqpoint{2.026373in}{2.560130in}}%
\pgfpathlineto{\pgfqpoint{2.035862in}{2.519406in}}%
\pgfpathlineto{\pgfqpoint{2.045351in}{2.509225in}}%
\pgfpathlineto{\pgfqpoint{2.054840in}{2.585583in}}%
\pgfpathlineto{\pgfqpoint{2.064328in}{2.580493in}}%
\pgfpathlineto{\pgfqpoint{2.073817in}{2.493953in}}%
\pgfpathlineto{\pgfqpoint{2.083306in}{2.549631in}}%
\pgfpathlineto{\pgfqpoint{2.092795in}{2.554722in}}%
\pgfpathlineto{\pgfqpoint{2.102284in}{2.498725in}}%
\pgfpathlineto{\pgfqpoint{2.111773in}{2.569994in}}%
\pgfpathlineto{\pgfqpoint{2.121262in}{2.539450in}}%
\pgfpathlineto{\pgfqpoint{2.130751in}{2.717621in}}%
\pgfpathlineto{\pgfqpoint{2.140239in}{2.442729in}}%
\pgfpathlineto{\pgfqpoint{2.149728in}{2.529269in}}%
\pgfpathlineto{\pgfqpoint{2.159217in}{2.498407in}}%
\pgfpathlineto{\pgfqpoint{2.168706in}{2.534041in}}%
\pgfpathlineto{\pgfqpoint{2.178195in}{2.488226in}}%
\pgfpathlineto{\pgfqpoint{2.187684in}{2.686759in}}%
\pgfpathlineto{\pgfqpoint{2.197173in}{2.528951in}}%
\pgfpathlineto{\pgfqpoint{2.206662in}{2.508588in}}%
\pgfpathlineto{\pgfqpoint{2.216150in}{2.518769in}}%
\pgfpathlineto{\pgfqpoint{2.225639in}{2.564585in}}%
\pgfpathlineto{\pgfqpoint{2.235128in}{2.482817in}}%
\pgfpathlineto{\pgfqpoint{2.244617in}{2.467545in}}%
\pgfpathlineto{\pgfqpoint{2.254106in}{2.472636in}}%
\pgfpathlineto{\pgfqpoint{2.263595in}{2.498089in}}%
\pgfpathlineto{\pgfqpoint{2.273084in}{2.457364in}}%
\pgfpathlineto{\pgfqpoint{2.282573in}{2.599901in}}%
\pgfpathlineto{\pgfqpoint{2.292062in}{2.594810in}}%
\pgfpathlineto{\pgfqpoint{2.301550in}{2.772981in}}%
\pgfpathlineto{\pgfqpoint{2.311039in}{2.879565in}}%
\pgfpathlineto{\pgfqpoint{2.320528in}{2.945743in}}%
\pgfpathlineto{\pgfqpoint{2.330017in}{2.986467in}}%
\pgfpathlineto{\pgfqpoint{2.339506in}{3.006830in}}%
\pgfpathlineto{\pgfqpoint{2.348995in}{3.006830in}}%
\pgfpathlineto{\pgfqpoint{2.358484in}{2.986467in}}%
\pgfpathlineto{\pgfqpoint{2.367973in}{2.950833in}}%
\pgfpathlineto{\pgfqpoint{2.377461in}{2.889746in}}%
\pgfpathlineto{\pgfqpoint{2.386950in}{2.787616in}}%
\pgfpathlineto{\pgfqpoint{2.396439in}{2.619627in}}%
\pgfpathlineto{\pgfqpoint{2.405928in}{2.502543in}}%
\pgfpathlineto{\pgfqpoint{2.415417in}{2.436366in}}%
\pgfpathlineto{\pgfqpoint{2.424906in}{2.466909in}}%
\pgfpathlineto{\pgfqpoint{2.434395in}{2.441456in}}%
\pgfpathlineto{\pgfqpoint{2.443884in}{2.517815in}}%
\pgfpathlineto{\pgfqpoint{2.453373in}{2.405822in}}%
\pgfpathlineto{\pgfqpoint{2.472350in}{2.497134in}}%
\pgfpathlineto{\pgfqpoint{2.481839in}{2.471682in}}%
\pgfpathlineto{\pgfqpoint{2.491328in}{2.512406in}}%
\pgfpathlineto{\pgfqpoint{2.500817in}{2.481863in}}%
\pgfpathlineto{\pgfqpoint{2.510306in}{2.405504in}}%
\pgfpathlineto{\pgfqpoint{2.519795in}{2.420776in}}%
\pgfpathlineto{\pgfqpoint{2.529284in}{2.471682in}}%
\pgfpathlineto{\pgfqpoint{2.538772in}{2.420457in}}%
\pgfpathlineto{\pgfqpoint{2.548261in}{2.481545in}}%
\pgfpathlineto{\pgfqpoint{2.557750in}{2.486635in}}%
\pgfpathlineto{\pgfqpoint{2.567239in}{2.456092in}}%
\pgfpathlineto{\pgfqpoint{2.576728in}{2.456092in}}%
\pgfpathlineto{\pgfqpoint{2.586217in}{2.466273in}}%
\pgfpathlineto{\pgfqpoint{2.595706in}{2.435729in}}%
\pgfpathlineto{\pgfqpoint{2.605195in}{2.420457in}}%
\pgfpathlineto{\pgfqpoint{2.614684in}{2.465955in}}%
\pgfpathlineto{\pgfqpoint{2.624172in}{2.455773in}}%
\pgfpathlineto{\pgfqpoint{2.633661in}{2.425230in}}%
\pgfpathlineto{\pgfqpoint{2.643150in}{2.435411in}}%
\pgfpathlineto{\pgfqpoint{2.652639in}{2.404868in}}%
\pgfpathlineto{\pgfqpoint{2.662128in}{2.486317in}}%
\pgfpathlineto{\pgfqpoint{2.671617in}{2.481226in}}%
\pgfpathlineto{\pgfqpoint{2.681106in}{2.450683in}}%
\pgfpathlineto{\pgfqpoint{2.690595in}{2.404549in}}%
\pgfpathlineto{\pgfqpoint{2.700083in}{2.435093in}}%
\pgfpathlineto{\pgfqpoint{2.709572in}{2.491089in}}%
\pgfpathlineto{\pgfqpoint{2.719061in}{2.424912in}}%
\pgfpathlineto{\pgfqpoint{2.728550in}{2.455455in}}%
\pgfpathlineto{\pgfqpoint{2.738039in}{2.450365in}}%
\pgfpathlineto{\pgfqpoint{2.747528in}{2.419821in}}%
\pgfpathlineto{\pgfqpoint{2.757017in}{2.485999in}}%
\pgfpathlineto{\pgfqpoint{2.766506in}{2.465318in}}%
\pgfpathlineto{\pgfqpoint{2.775995in}{2.465318in}}%
\pgfpathlineto{\pgfqpoint{2.785483in}{2.404231in}}%
\pgfpathlineto{\pgfqpoint{2.794972in}{2.470409in}}%
\pgfpathlineto{\pgfqpoint{2.804461in}{2.419503in}}%
\pgfpathlineto{\pgfqpoint{2.813950in}{2.439865in}}%
\pgfpathlineto{\pgfqpoint{2.823439in}{2.470409in}}%
\pgfpathlineto{\pgfqpoint{2.832928in}{2.419503in}}%
\pgfpathlineto{\pgfqpoint{2.842417in}{2.419185in}}%
\pgfpathlineto{\pgfqpoint{2.851906in}{2.439547in}}%
\pgfpathlineto{\pgfqpoint{2.861394in}{2.449728in}}%
\pgfpathlineto{\pgfqpoint{2.870883in}{2.480272in}}%
\pgfpathlineto{\pgfqpoint{2.870883in}{2.480272in}}%
\pgfusepath{stroke}%
\end{pgfscope}%
\begin{pgfscope}%
\pgfsetrectcap%
\pgfsetmiterjoin%
\pgfsetlinewidth{0.803000pt}%
\definecolor{currentstroke}{rgb}{0.000000,0.000000,0.000000}%
\pgfsetstrokecolor{currentstroke}%
\pgfsetdash{}{0pt}%
\pgfpathmoveto{\pgfqpoint{0.758026in}{2.240123in}}%
\pgfpathlineto{\pgfqpoint{0.758026in}{3.141667in}}%
\pgfusepath{stroke}%
\end{pgfscope}%
\begin{pgfscope}%
\pgfsetrectcap%
\pgfsetmiterjoin%
\pgfsetlinewidth{0.803000pt}%
\definecolor{currentstroke}{rgb}{0.000000,0.000000,0.000000}%
\pgfsetstrokecolor{currentstroke}%
\pgfsetdash{}{0pt}%
\pgfpathmoveto{\pgfqpoint{2.866666in}{2.240123in}}%
\pgfpathlineto{\pgfqpoint{2.866666in}{3.141667in}}%
\pgfusepath{stroke}%
\end{pgfscope}%
\begin{pgfscope}%
\pgfsetrectcap%
\pgfsetmiterjoin%
\pgfsetlinewidth{0.803000pt}%
\definecolor{currentstroke}{rgb}{0.000000,0.000000,0.000000}%
\pgfsetstrokecolor{currentstroke}%
\pgfsetdash{}{0pt}%
\pgfpathmoveto{\pgfqpoint{0.758026in}{2.240123in}}%
\pgfpathlineto{\pgfqpoint{2.866666in}{2.240123in}}%
\pgfusepath{stroke}%
\end{pgfscope}%
\begin{pgfscope}%
\pgfsetrectcap%
\pgfsetmiterjoin%
\pgfsetlinewidth{0.803000pt}%
\definecolor{currentstroke}{rgb}{0.000000,0.000000,0.000000}%
\pgfsetstrokecolor{currentstroke}%
\pgfsetdash{}{0pt}%
\pgfpathmoveto{\pgfqpoint{0.758026in}{3.141667in}}%
\pgfpathlineto{\pgfqpoint{2.866666in}{3.141667in}}%
\pgfusepath{stroke}%
\end{pgfscope}%
\begin{pgfscope}%
\definecolor{textcolor}{rgb}{0.000000,0.000000,0.000000}%
\pgfsetstrokecolor{textcolor}%
\pgfsetfillcolor{textcolor}%
\pgftext[x=1.812346in,y=3.225000in,,base]{\color{textcolor}\rmfamily\fontsize{12.000000}{14.400000}\selectfont 50 kHz, 70\%}%
\end{pgfscope}%
\begin{pgfscope}%
\pgfsetbuttcap%
\pgfsetmiterjoin%
\definecolor{currentfill}{rgb}{1.000000,1.000000,1.000000}%
\pgfsetfillcolor{currentfill}%
\pgfsetlinewidth{0.000000pt}%
\definecolor{currentstroke}{rgb}{0.000000,0.000000,0.000000}%
\pgfsetstrokecolor{currentstroke}%
\pgfsetstrokeopacity{0.000000}%
\pgfsetdash{}{0pt}%
\pgfpathmoveto{\pgfqpoint{3.833026in}{2.240123in}}%
\pgfpathlineto{\pgfqpoint{5.941666in}{2.240123in}}%
\pgfpathlineto{\pgfqpoint{5.941666in}{3.141667in}}%
\pgfpathlineto{\pgfqpoint{3.833026in}{3.141667in}}%
\pgfpathlineto{\pgfqpoint{3.833026in}{2.240123in}}%
\pgfpathclose%
\pgfusepath{fill}%
\end{pgfscope}%
\begin{pgfscope}%
\pgfpathrectangle{\pgfqpoint{3.833026in}{2.240123in}}{\pgfqpoint{2.108640in}{0.901543in}}%
\pgfusepath{clip}%
\pgfsetbuttcap%
\pgfsetroundjoin%
\pgfsetlinewidth{0.501875pt}%
\definecolor{currentstroke}{rgb}{0.690196,0.690196,0.690196}%
\pgfsetstrokecolor{currentstroke}%
\pgfsetstrokeopacity{0.700000}%
\pgfsetdash{{1.850000pt}{0.800000pt}}{0.000000pt}%
\pgfpathmoveto{\pgfqpoint{3.833026in}{2.240123in}}%
\pgfpathlineto{\pgfqpoint{3.833026in}{3.141667in}}%
\pgfusepath{stroke}%
\end{pgfscope}%
\begin{pgfscope}%
\pgfsetbuttcap%
\pgfsetroundjoin%
\definecolor{currentfill}{rgb}{0.000000,0.000000,0.000000}%
\pgfsetfillcolor{currentfill}%
\pgfsetlinewidth{0.803000pt}%
\definecolor{currentstroke}{rgb}{0.000000,0.000000,0.000000}%
\pgfsetstrokecolor{currentstroke}%
\pgfsetdash{}{0pt}%
\pgfsys@defobject{currentmarker}{\pgfqpoint{0.000000in}{-0.048611in}}{\pgfqpoint{0.000000in}{0.000000in}}{%
\pgfpathmoveto{\pgfqpoint{0.000000in}{0.000000in}}%
\pgfpathlineto{\pgfqpoint{0.000000in}{-0.048611in}}%
\pgfusepath{stroke,fill}%
}%
\begin{pgfscope}%
\pgfsys@transformshift{3.833026in}{2.240123in}%
\pgfsys@useobject{currentmarker}{}%
\end{pgfscope}%
\end{pgfscope}%
\begin{pgfscope}%
\definecolor{textcolor}{rgb}{0.000000,0.000000,0.000000}%
\pgfsetstrokecolor{textcolor}%
\pgfsetfillcolor{textcolor}%
\pgftext[x=3.833026in,y=2.142901in,,top]{\color{textcolor}\rmfamily\fontsize{10.000000}{12.000000}\selectfont \(\displaystyle {400000}\)}%
\end{pgfscope}%
\begin{pgfscope}%
\pgfpathrectangle{\pgfqpoint{3.833026in}{2.240123in}}{\pgfqpoint{2.108640in}{0.901543in}}%
\pgfusepath{clip}%
\pgfsetbuttcap%
\pgfsetroundjoin%
\pgfsetlinewidth{0.501875pt}%
\definecolor{currentstroke}{rgb}{0.690196,0.690196,0.690196}%
\pgfsetstrokecolor{currentstroke}%
\pgfsetstrokeopacity{0.700000}%
\pgfsetdash{{1.850000pt}{0.800000pt}}{0.000000pt}%
\pgfpathmoveto{\pgfqpoint{4.360186in}{2.240123in}}%
\pgfpathlineto{\pgfqpoint{4.360186in}{3.141667in}}%
\pgfusepath{stroke}%
\end{pgfscope}%
\begin{pgfscope}%
\pgfsetbuttcap%
\pgfsetroundjoin%
\definecolor{currentfill}{rgb}{0.000000,0.000000,0.000000}%
\pgfsetfillcolor{currentfill}%
\pgfsetlinewidth{0.803000pt}%
\definecolor{currentstroke}{rgb}{0.000000,0.000000,0.000000}%
\pgfsetstrokecolor{currentstroke}%
\pgfsetdash{}{0pt}%
\pgfsys@defobject{currentmarker}{\pgfqpoint{0.000000in}{-0.048611in}}{\pgfqpoint{0.000000in}{0.000000in}}{%
\pgfpathmoveto{\pgfqpoint{0.000000in}{0.000000in}}%
\pgfpathlineto{\pgfqpoint{0.000000in}{-0.048611in}}%
\pgfusepath{stroke,fill}%
}%
\begin{pgfscope}%
\pgfsys@transformshift{4.360186in}{2.240123in}%
\pgfsys@useobject{currentmarker}{}%
\end{pgfscope}%
\end{pgfscope}%
\begin{pgfscope}%
\definecolor{textcolor}{rgb}{0.000000,0.000000,0.000000}%
\pgfsetstrokecolor{textcolor}%
\pgfsetfillcolor{textcolor}%
\pgftext[x=4.360186in,y=2.142901in,,top]{\color{textcolor}\rmfamily\fontsize{10.000000}{12.000000}\selectfont \(\displaystyle {450000}\)}%
\end{pgfscope}%
\begin{pgfscope}%
\pgfpathrectangle{\pgfqpoint{3.833026in}{2.240123in}}{\pgfqpoint{2.108640in}{0.901543in}}%
\pgfusepath{clip}%
\pgfsetbuttcap%
\pgfsetroundjoin%
\pgfsetlinewidth{0.501875pt}%
\definecolor{currentstroke}{rgb}{0.690196,0.690196,0.690196}%
\pgfsetstrokecolor{currentstroke}%
\pgfsetstrokeopacity{0.700000}%
\pgfsetdash{{1.850000pt}{0.800000pt}}{0.000000pt}%
\pgfpathmoveto{\pgfqpoint{4.887346in}{2.240123in}}%
\pgfpathlineto{\pgfqpoint{4.887346in}{3.141667in}}%
\pgfusepath{stroke}%
\end{pgfscope}%
\begin{pgfscope}%
\pgfsetbuttcap%
\pgfsetroundjoin%
\definecolor{currentfill}{rgb}{0.000000,0.000000,0.000000}%
\pgfsetfillcolor{currentfill}%
\pgfsetlinewidth{0.803000pt}%
\definecolor{currentstroke}{rgb}{0.000000,0.000000,0.000000}%
\pgfsetstrokecolor{currentstroke}%
\pgfsetdash{}{0pt}%
\pgfsys@defobject{currentmarker}{\pgfqpoint{0.000000in}{-0.048611in}}{\pgfqpoint{0.000000in}{0.000000in}}{%
\pgfpathmoveto{\pgfqpoint{0.000000in}{0.000000in}}%
\pgfpathlineto{\pgfqpoint{0.000000in}{-0.048611in}}%
\pgfusepath{stroke,fill}%
}%
\begin{pgfscope}%
\pgfsys@transformshift{4.887346in}{2.240123in}%
\pgfsys@useobject{currentmarker}{}%
\end{pgfscope}%
\end{pgfscope}%
\begin{pgfscope}%
\definecolor{textcolor}{rgb}{0.000000,0.000000,0.000000}%
\pgfsetstrokecolor{textcolor}%
\pgfsetfillcolor{textcolor}%
\pgftext[x=4.887346in,y=2.142901in,,top]{\color{textcolor}\rmfamily\fontsize{10.000000}{12.000000}\selectfont \(\displaystyle {500000}\)}%
\end{pgfscope}%
\begin{pgfscope}%
\pgfpathrectangle{\pgfqpoint{3.833026in}{2.240123in}}{\pgfqpoint{2.108640in}{0.901543in}}%
\pgfusepath{clip}%
\pgfsetbuttcap%
\pgfsetroundjoin%
\pgfsetlinewidth{0.501875pt}%
\definecolor{currentstroke}{rgb}{0.690196,0.690196,0.690196}%
\pgfsetstrokecolor{currentstroke}%
\pgfsetstrokeopacity{0.700000}%
\pgfsetdash{{1.850000pt}{0.800000pt}}{0.000000pt}%
\pgfpathmoveto{\pgfqpoint{5.414506in}{2.240123in}}%
\pgfpathlineto{\pgfqpoint{5.414506in}{3.141667in}}%
\pgfusepath{stroke}%
\end{pgfscope}%
\begin{pgfscope}%
\pgfsetbuttcap%
\pgfsetroundjoin%
\definecolor{currentfill}{rgb}{0.000000,0.000000,0.000000}%
\pgfsetfillcolor{currentfill}%
\pgfsetlinewidth{0.803000pt}%
\definecolor{currentstroke}{rgb}{0.000000,0.000000,0.000000}%
\pgfsetstrokecolor{currentstroke}%
\pgfsetdash{}{0pt}%
\pgfsys@defobject{currentmarker}{\pgfqpoint{0.000000in}{-0.048611in}}{\pgfqpoint{0.000000in}{0.000000in}}{%
\pgfpathmoveto{\pgfqpoint{0.000000in}{0.000000in}}%
\pgfpathlineto{\pgfqpoint{0.000000in}{-0.048611in}}%
\pgfusepath{stroke,fill}%
}%
\begin{pgfscope}%
\pgfsys@transformshift{5.414506in}{2.240123in}%
\pgfsys@useobject{currentmarker}{}%
\end{pgfscope}%
\end{pgfscope}%
\begin{pgfscope}%
\definecolor{textcolor}{rgb}{0.000000,0.000000,0.000000}%
\pgfsetstrokecolor{textcolor}%
\pgfsetfillcolor{textcolor}%
\pgftext[x=5.414506in,y=2.142901in,,top]{\color{textcolor}\rmfamily\fontsize{10.000000}{12.000000}\selectfont \(\displaystyle {550000}\)}%
\end{pgfscope}%
\begin{pgfscope}%
\pgfpathrectangle{\pgfqpoint{3.833026in}{2.240123in}}{\pgfqpoint{2.108640in}{0.901543in}}%
\pgfusepath{clip}%
\pgfsetbuttcap%
\pgfsetroundjoin%
\pgfsetlinewidth{0.501875pt}%
\definecolor{currentstroke}{rgb}{0.690196,0.690196,0.690196}%
\pgfsetstrokecolor{currentstroke}%
\pgfsetstrokeopacity{0.700000}%
\pgfsetdash{{1.850000pt}{0.800000pt}}{0.000000pt}%
\pgfpathmoveto{\pgfqpoint{5.941666in}{2.240123in}}%
\pgfpathlineto{\pgfqpoint{5.941666in}{3.141667in}}%
\pgfusepath{stroke}%
\end{pgfscope}%
\begin{pgfscope}%
\pgfsetbuttcap%
\pgfsetroundjoin%
\definecolor{currentfill}{rgb}{0.000000,0.000000,0.000000}%
\pgfsetfillcolor{currentfill}%
\pgfsetlinewidth{0.803000pt}%
\definecolor{currentstroke}{rgb}{0.000000,0.000000,0.000000}%
\pgfsetstrokecolor{currentstroke}%
\pgfsetdash{}{0pt}%
\pgfsys@defobject{currentmarker}{\pgfqpoint{0.000000in}{-0.048611in}}{\pgfqpoint{0.000000in}{0.000000in}}{%
\pgfpathmoveto{\pgfqpoint{0.000000in}{0.000000in}}%
\pgfpathlineto{\pgfqpoint{0.000000in}{-0.048611in}}%
\pgfusepath{stroke,fill}%
}%
\begin{pgfscope}%
\pgfsys@transformshift{5.941666in}{2.240123in}%
\pgfsys@useobject{currentmarker}{}%
\end{pgfscope}%
\end{pgfscope}%
\begin{pgfscope}%
\definecolor{textcolor}{rgb}{0.000000,0.000000,0.000000}%
\pgfsetstrokecolor{textcolor}%
\pgfsetfillcolor{textcolor}%
\pgftext[x=5.941666in,y=2.142901in,,top]{\color{textcolor}\rmfamily\fontsize{10.000000}{12.000000}\selectfont \(\displaystyle {600000}\)}%
\end{pgfscope}%
\begin{pgfscope}%
\definecolor{textcolor}{rgb}{0.000000,0.000000,0.000000}%
\pgfsetstrokecolor{textcolor}%
\pgfsetfillcolor{textcolor}%
\pgftext[x=4.887346in,y=1.963889in,,top]{\color{textcolor}\rmfamily\fontsize{10.000000}{12.000000}\selectfont Frequency (Hz)}%
\end{pgfscope}%
\begin{pgfscope}%
\pgfpathrectangle{\pgfqpoint{3.833026in}{2.240123in}}{\pgfqpoint{2.108640in}{0.901543in}}%
\pgfusepath{clip}%
\pgfsetbuttcap%
\pgfsetroundjoin%
\pgfsetlinewidth{0.501875pt}%
\definecolor{currentstroke}{rgb}{0.690196,0.690196,0.690196}%
\pgfsetstrokecolor{currentstroke}%
\pgfsetstrokeopacity{0.700000}%
\pgfsetdash{{1.850000pt}{0.800000pt}}{0.000000pt}%
\pgfpathmoveto{\pgfqpoint{3.833026in}{2.391128in}}%
\pgfpathlineto{\pgfqpoint{5.941666in}{2.391128in}}%
\pgfusepath{stroke}%
\end{pgfscope}%
\begin{pgfscope}%
\pgfsetbuttcap%
\pgfsetroundjoin%
\definecolor{currentfill}{rgb}{0.000000,0.000000,0.000000}%
\pgfsetfillcolor{currentfill}%
\pgfsetlinewidth{0.803000pt}%
\definecolor{currentstroke}{rgb}{0.000000,0.000000,0.000000}%
\pgfsetstrokecolor{currentstroke}%
\pgfsetdash{}{0pt}%
\pgfsys@defobject{currentmarker}{\pgfqpoint{-0.048611in}{0.000000in}}{\pgfqpoint{-0.000000in}{0.000000in}}{%
\pgfpathmoveto{\pgfqpoint{-0.000000in}{0.000000in}}%
\pgfpathlineto{\pgfqpoint{-0.048611in}{0.000000in}}%
\pgfusepath{stroke,fill}%
}%
\begin{pgfscope}%
\pgfsys@transformshift{3.833026in}{2.391128in}%
\pgfsys@useobject{currentmarker}{}%
\end{pgfscope}%
\end{pgfscope}%
\begin{pgfscope}%
\definecolor{textcolor}{rgb}{0.000000,0.000000,0.000000}%
\pgfsetstrokecolor{textcolor}%
\pgfsetfillcolor{textcolor}%
\pgftext[x=3.419444in, y=2.342903in, left, base]{\color{textcolor}\rmfamily\fontsize{10.000000}{12.000000}\selectfont \(\displaystyle {\ensuremath{-}100}\)}%
\end{pgfscope}%
\begin{pgfscope}%
\pgfpathrectangle{\pgfqpoint{3.833026in}{2.240123in}}{\pgfqpoint{2.108640in}{0.901543in}}%
\pgfusepath{clip}%
\pgfsetbuttcap%
\pgfsetroundjoin%
\pgfsetlinewidth{0.501875pt}%
\definecolor{currentstroke}{rgb}{0.690196,0.690196,0.690196}%
\pgfsetstrokecolor{currentstroke}%
\pgfsetstrokeopacity{0.700000}%
\pgfsetdash{{1.850000pt}{0.800000pt}}{0.000000pt}%
\pgfpathmoveto{\pgfqpoint{3.833026in}{2.876090in}}%
\pgfpathlineto{\pgfqpoint{5.941666in}{2.876090in}}%
\pgfusepath{stroke}%
\end{pgfscope}%
\begin{pgfscope}%
\pgfsetbuttcap%
\pgfsetroundjoin%
\definecolor{currentfill}{rgb}{0.000000,0.000000,0.000000}%
\pgfsetfillcolor{currentfill}%
\pgfsetlinewidth{0.803000pt}%
\definecolor{currentstroke}{rgb}{0.000000,0.000000,0.000000}%
\pgfsetstrokecolor{currentstroke}%
\pgfsetdash{}{0pt}%
\pgfsys@defobject{currentmarker}{\pgfqpoint{-0.048611in}{0.000000in}}{\pgfqpoint{-0.000000in}{0.000000in}}{%
\pgfpathmoveto{\pgfqpoint{-0.000000in}{0.000000in}}%
\pgfpathlineto{\pgfqpoint{-0.048611in}{0.000000in}}%
\pgfusepath{stroke,fill}%
}%
\begin{pgfscope}%
\pgfsys@transformshift{3.833026in}{2.876090in}%
\pgfsys@useobject{currentmarker}{}%
\end{pgfscope}%
\end{pgfscope}%
\begin{pgfscope}%
\definecolor{textcolor}{rgb}{0.000000,0.000000,0.000000}%
\pgfsetstrokecolor{textcolor}%
\pgfsetfillcolor{textcolor}%
\pgftext[x=3.488889in, y=2.827864in, left, base]{\color{textcolor}\rmfamily\fontsize{10.000000}{12.000000}\selectfont \(\displaystyle {\ensuremath{-}50}\)}%
\end{pgfscope}%
\begin{pgfscope}%
\definecolor{textcolor}{rgb}{0.000000,0.000000,0.000000}%
\pgfsetstrokecolor{textcolor}%
\pgfsetfillcolor{textcolor}%
\pgftext[x=3.363889in,y=2.690895in,,bottom,rotate=90.000000]{\color{textcolor}\rmfamily\fontsize{10.000000}{12.000000}\selectfont Amplitude (dB)}%
\end{pgfscope}%
\begin{pgfscope}%
\pgfpathrectangle{\pgfqpoint{3.833026in}{2.240123in}}{\pgfqpoint{2.108640in}{0.901543in}}%
\pgfusepath{clip}%
\pgfsetrectcap%
\pgfsetroundjoin%
\pgfsetlinewidth{1.003750pt}%
\definecolor{currentstroke}{rgb}{0.000000,0.000000,0.000000}%
\pgfsetstrokecolor{currentstroke}%
\pgfsetdash{}{0pt}%
\pgfpathmoveto{\pgfqpoint{3.829863in}{2.571473in}}%
\pgfpathlineto{\pgfqpoint{3.839352in}{2.566624in}}%
\pgfpathlineto{\pgfqpoint{3.848840in}{2.532676in}}%
\pgfpathlineto{\pgfqpoint{3.858329in}{2.629669in}}%
\pgfpathlineto{\pgfqpoint{3.867818in}{2.508428in}}%
\pgfpathlineto{\pgfqpoint{3.877307in}{2.546922in}}%
\pgfpathlineto{\pgfqpoint{3.886796in}{2.571170in}}%
\pgfpathlineto{\pgfqpoint{3.896285in}{2.561471in}}%
\pgfpathlineto{\pgfqpoint{3.905774in}{2.517824in}}%
\pgfpathlineto{\pgfqpoint{3.915263in}{2.595418in}}%
\pgfpathlineto{\pgfqpoint{3.924752in}{2.474178in}}%
\pgfpathlineto{\pgfqpoint{3.934240in}{2.566321in}}%
\pgfpathlineto{\pgfqpoint{3.943729in}{2.629366in}}%
\pgfpathlineto{\pgfqpoint{3.953218in}{2.590265in}}%
\pgfpathlineto{\pgfqpoint{3.962707in}{2.517521in}}%
\pgfpathlineto{\pgfqpoint{3.972196in}{2.473875in}}%
\pgfpathlineto{\pgfqpoint{3.981685in}{2.585416in}}%
\pgfpathlineto{\pgfqpoint{3.991174in}{2.566017in}}%
\pgfpathlineto{\pgfqpoint{4.000663in}{2.580566in}}%
\pgfpathlineto{\pgfqpoint{4.010151in}{2.551469in}}%
\pgfpathlineto{\pgfqpoint{4.019640in}{2.575717in}}%
\pgfpathlineto{\pgfqpoint{4.029129in}{2.580263in}}%
\pgfpathlineto{\pgfqpoint{4.038618in}{2.531767in}}%
\pgfpathlineto{\pgfqpoint{4.048107in}{2.546316in}}%
\pgfpathlineto{\pgfqpoint{4.057596in}{2.599662in}}%
\pgfpathlineto{\pgfqpoint{4.067085in}{2.556015in}}%
\pgfpathlineto{\pgfqpoint{4.076574in}{2.546316in}}%
\pgfpathlineto{\pgfqpoint{4.086063in}{2.556015in}}%
\pgfpathlineto{\pgfqpoint{4.095551in}{2.570564in}}%
\pgfpathlineto{\pgfqpoint{4.105040in}{2.492667in}}%
\pgfpathlineto{\pgfqpoint{4.114529in}{2.599359in}}%
\pgfpathlineto{\pgfqpoint{4.124018in}{2.560562in}}%
\pgfpathlineto{\pgfqpoint{4.133507in}{2.599359in}}%
\pgfpathlineto{\pgfqpoint{4.142996in}{2.570261in}}%
\pgfpathlineto{\pgfqpoint{4.152485in}{2.560562in}}%
\pgfpathlineto{\pgfqpoint{4.161974in}{2.570261in}}%
\pgfpathlineto{\pgfqpoint{4.171462in}{2.531464in}}%
\pgfpathlineto{\pgfqpoint{4.180951in}{2.555409in}}%
\pgfpathlineto{\pgfqpoint{4.190440in}{2.540860in}}%
\pgfpathlineto{\pgfqpoint{4.199929in}{2.589356in}}%
\pgfpathlineto{\pgfqpoint{4.209418in}{2.565108in}}%
\pgfpathlineto{\pgfqpoint{4.218907in}{2.584507in}}%
\pgfpathlineto{\pgfqpoint{4.228396in}{2.579657in}}%
\pgfpathlineto{\pgfqpoint{4.237885in}{2.560259in}}%
\pgfpathlineto{\pgfqpoint{4.247374in}{2.560259in}}%
\pgfpathlineto{\pgfqpoint{4.256862in}{2.579354in}}%
\pgfpathlineto{\pgfqpoint{4.266351in}{2.526008in}}%
\pgfpathlineto{\pgfqpoint{4.275840in}{2.516309in}}%
\pgfpathlineto{\pgfqpoint{4.285329in}{2.574504in}}%
\pgfpathlineto{\pgfqpoint{4.294818in}{2.496910in}}%
\pgfpathlineto{\pgfqpoint{4.304307in}{2.540557in}}%
\pgfpathlineto{\pgfqpoint{4.313796in}{2.598752in}}%
\pgfpathlineto{\pgfqpoint{4.323285in}{2.719690in}}%
\pgfpathlineto{\pgfqpoint{4.332773in}{2.802133in}}%
\pgfpathlineto{\pgfqpoint{4.342262in}{2.860328in}}%
\pgfpathlineto{\pgfqpoint{4.351751in}{2.894276in}}%
\pgfpathlineto{\pgfqpoint{4.361240in}{2.908825in}}%
\pgfpathlineto{\pgfqpoint{4.370729in}{2.908825in}}%
\pgfpathlineto{\pgfqpoint{4.380218in}{2.884576in}}%
\pgfpathlineto{\pgfqpoint{4.389707in}{2.840930in}}%
\pgfpathlineto{\pgfqpoint{4.399196in}{2.782431in}}%
\pgfpathlineto{\pgfqpoint{4.408685in}{2.675740in}}%
\pgfpathlineto{\pgfqpoint{4.418173in}{2.535101in}}%
\pgfpathlineto{\pgfqpoint{4.427662in}{2.520552in}}%
\pgfpathlineto{\pgfqpoint{4.437151in}{2.569048in}}%
\pgfpathlineto{\pgfqpoint{4.446640in}{2.593297in}}%
\pgfpathlineto{\pgfqpoint{4.456129in}{2.602996in}}%
\pgfpathlineto{\pgfqpoint{4.465618in}{2.525402in}}%
\pgfpathlineto{\pgfqpoint{4.475107in}{2.505700in}}%
\pgfpathlineto{\pgfqpoint{4.484596in}{2.520249in}}%
\pgfpathlineto{\pgfqpoint{4.494084in}{2.481452in}}%
\pgfpathlineto{\pgfqpoint{4.503573in}{2.529948in}}%
\pgfpathlineto{\pgfqpoint{4.513062in}{2.549347in}}%
\pgfpathlineto{\pgfqpoint{4.522551in}{2.549347in}}%
\pgfpathlineto{\pgfqpoint{4.532040in}{2.481452in}}%
\pgfpathlineto{\pgfqpoint{4.541529in}{2.525099in}}%
\pgfpathlineto{\pgfqpoint{4.551018in}{2.592690in}}%
\pgfpathlineto{\pgfqpoint{4.560507in}{2.500548in}}%
\pgfpathlineto{\pgfqpoint{4.569995in}{2.592690in}}%
\pgfpathlineto{\pgfqpoint{4.579484in}{2.558743in}}%
\pgfpathlineto{\pgfqpoint{4.588973in}{2.573292in}}%
\pgfpathlineto{\pgfqpoint{4.598462in}{2.549044in}}%
\pgfpathlineto{\pgfqpoint{4.607951in}{2.471450in}}%
\pgfpathlineto{\pgfqpoint{4.617440in}{2.563593in}}%
\pgfpathlineto{\pgfqpoint{4.626929in}{2.548741in}}%
\pgfpathlineto{\pgfqpoint{4.636418in}{2.626335in}}%
\pgfpathlineto{\pgfqpoint{4.645907in}{2.548741in}}%
\pgfpathlineto{\pgfqpoint{4.655395in}{2.592387in}}%
\pgfpathlineto{\pgfqpoint{4.664884in}{2.592387in}}%
\pgfpathlineto{\pgfqpoint{4.674373in}{2.543891in}}%
\pgfpathlineto{\pgfqpoint{4.683862in}{2.563290in}}%
\pgfpathlineto{\pgfqpoint{4.693351in}{2.543891in}}%
\pgfpathlineto{\pgfqpoint{4.702840in}{2.567836in}}%
\pgfpathlineto{\pgfqpoint{4.712329in}{2.562986in}}%
\pgfpathlineto{\pgfqpoint{4.721818in}{2.499941in}}%
\pgfpathlineto{\pgfqpoint{4.731306in}{2.543588in}}%
\pgfpathlineto{\pgfqpoint{4.740795in}{2.524190in}}%
\pgfpathlineto{\pgfqpoint{4.750284in}{2.587234in}}%
\pgfpathlineto{\pgfqpoint{4.759773in}{2.553287in}}%
\pgfpathlineto{\pgfqpoint{4.769262in}{2.587234in}}%
\pgfpathlineto{\pgfqpoint{4.778751in}{2.552984in}}%
\pgfpathlineto{\pgfqpoint{4.788240in}{2.455992in}}%
\pgfpathlineto{\pgfqpoint{4.797729in}{2.552984in}}%
\pgfpathlineto{\pgfqpoint{4.807218in}{2.523886in}}%
\pgfpathlineto{\pgfqpoint{4.816706in}{2.543285in}}%
\pgfpathlineto{\pgfqpoint{4.826195in}{2.557834in}}%
\pgfpathlineto{\pgfqpoint{4.835684in}{2.538435in}}%
\pgfpathlineto{\pgfqpoint{4.845173in}{2.829412in}}%
\pgfpathlineto{\pgfqpoint{4.854662in}{2.950349in}}%
\pgfpathlineto{\pgfqpoint{4.864151in}{3.027943in}}%
\pgfpathlineto{\pgfqpoint{4.873640in}{3.071590in}}%
\pgfpathlineto{\pgfqpoint{4.883129in}{3.095838in}}%
\pgfpathlineto{\pgfqpoint{4.892617in}{3.100687in}}%
\pgfpathlineto{\pgfqpoint{4.902106in}{3.095838in}}%
\pgfpathlineto{\pgfqpoint{4.911595in}{3.066740in}}%
\pgfpathlineto{\pgfqpoint{4.921084in}{3.018244in}}%
\pgfpathlineto{\pgfqpoint{4.930573in}{2.935497in}}%
\pgfpathlineto{\pgfqpoint{4.940062in}{2.804558in}}%
\pgfpathlineto{\pgfqpoint{4.949551in}{2.489333in}}%
\pgfpathlineto{\pgfqpoint{4.959040in}{2.547528in}}%
\pgfpathlineto{\pgfqpoint{4.968529in}{2.528130in}}%
\pgfpathlineto{\pgfqpoint{4.978017in}{2.542679in}}%
\pgfpathlineto{\pgfqpoint{4.987506in}{2.532979in}}%
\pgfpathlineto{\pgfqpoint{4.996995in}{2.557228in}}%
\pgfpathlineto{\pgfqpoint{5.006484in}{2.489030in}}%
\pgfpathlineto{\pgfqpoint{5.015973in}{2.537526in}}%
\pgfpathlineto{\pgfqpoint{5.025462in}{2.498729in}}%
\pgfpathlineto{\pgfqpoint{5.034951in}{2.498729in}}%
\pgfpathlineto{\pgfqpoint{5.044440in}{2.547225in}}%
\pgfpathlineto{\pgfqpoint{5.053928in}{2.566624in}}%
\pgfpathlineto{\pgfqpoint{5.063417in}{2.474481in}}%
\pgfpathlineto{\pgfqpoint{5.072906in}{2.532676in}}%
\pgfpathlineto{\pgfqpoint{5.082395in}{2.440231in}}%
\pgfpathlineto{\pgfqpoint{5.091884in}{2.556621in}}%
\pgfpathlineto{\pgfqpoint{5.101373in}{2.498426in}}%
\pgfpathlineto{\pgfqpoint{5.110862in}{2.551772in}}%
\pgfpathlineto{\pgfqpoint{5.120351in}{2.590569in}}%
\pgfpathlineto{\pgfqpoint{5.129840in}{2.566321in}}%
\pgfpathlineto{\pgfqpoint{5.139328in}{2.556621in}}%
\pgfpathlineto{\pgfqpoint{5.148817in}{2.585719in}}%
\pgfpathlineto{\pgfqpoint{5.167795in}{2.575717in}}%
\pgfpathlineto{\pgfqpoint{5.177284in}{2.532070in}}%
\pgfpathlineto{\pgfqpoint{5.186773in}{2.561168in}}%
\pgfpathlineto{\pgfqpoint{5.196262in}{2.536920in}}%
\pgfpathlineto{\pgfqpoint{5.205751in}{2.551469in}}%
\pgfpathlineto{\pgfqpoint{5.215239in}{2.502972in}}%
\pgfpathlineto{\pgfqpoint{5.234217in}{2.570564in}}%
\pgfpathlineto{\pgfqpoint{5.243706in}{2.497820in}}%
\pgfpathlineto{\pgfqpoint{5.253195in}{2.531767in}}%
\pgfpathlineto{\pgfqpoint{5.262684in}{2.497820in}}%
\pgfpathlineto{\pgfqpoint{5.272173in}{2.541466in}}%
\pgfpathlineto{\pgfqpoint{5.281662in}{2.546316in}}%
\pgfpathlineto{\pgfqpoint{5.291150in}{2.541466in}}%
\pgfpathlineto{\pgfqpoint{5.300639in}{2.522068in}}%
\pgfpathlineto{\pgfqpoint{5.310128in}{2.512065in}}%
\pgfpathlineto{\pgfqpoint{5.319617in}{2.516915in}}%
\pgfpathlineto{\pgfqpoint{5.329106in}{2.512065in}}%
\pgfpathlineto{\pgfqpoint{5.338595in}{2.473269in}}%
\pgfpathlineto{\pgfqpoint{5.348084in}{2.473269in}}%
\pgfpathlineto{\pgfqpoint{5.357573in}{2.550862in}}%
\pgfpathlineto{\pgfqpoint{5.367062in}{2.497517in}}%
\pgfpathlineto{\pgfqpoint{5.376550in}{2.686652in}}%
\pgfpathlineto{\pgfqpoint{5.386039in}{2.788190in}}%
\pgfpathlineto{\pgfqpoint{5.395528in}{2.851235in}}%
\pgfpathlineto{\pgfqpoint{5.405017in}{2.885183in}}%
\pgfpathlineto{\pgfqpoint{5.414506in}{2.909431in}}%
\pgfpathlineto{\pgfqpoint{5.423995in}{2.909431in}}%
\pgfpathlineto{\pgfqpoint{5.433484in}{2.890032in}}%
\pgfpathlineto{\pgfqpoint{5.442973in}{2.856085in}}%
\pgfpathlineto{\pgfqpoint{5.452461in}{2.793040in}}%
\pgfpathlineto{\pgfqpoint{5.461950in}{2.695745in}}%
\pgfpathlineto{\pgfqpoint{5.471439in}{2.535707in}}%
\pgfpathlineto{\pgfqpoint{5.480928in}{2.496910in}}%
\pgfpathlineto{\pgfqpoint{5.490417in}{2.550256in}}%
\pgfpathlineto{\pgfqpoint{5.499906in}{2.487211in}}%
\pgfpathlineto{\pgfqpoint{5.509395in}{2.453264in}}%
\pgfpathlineto{\pgfqpoint{5.518884in}{2.487211in}}%
\pgfpathlineto{\pgfqpoint{5.528373in}{2.482362in}}%
\pgfpathlineto{\pgfqpoint{5.537861in}{2.452961in}}%
\pgfpathlineto{\pgfqpoint{5.547350in}{2.467510in}}%
\pgfpathlineto{\pgfqpoint{5.556839in}{2.467510in}}%
\pgfpathlineto{\pgfqpoint{5.566328in}{2.472359in}}%
\pgfpathlineto{\pgfqpoint{5.575817in}{2.486908in}}%
\pgfpathlineto{\pgfqpoint{5.585306in}{2.520855in}}%
\pgfpathlineto{\pgfqpoint{5.594795in}{2.452961in}}%
\pgfpathlineto{\pgfqpoint{5.604284in}{2.501457in}}%
\pgfpathlineto{\pgfqpoint{5.613772in}{2.501154in}}%
\pgfpathlineto{\pgfqpoint{5.623261in}{2.530252in}}%
\pgfpathlineto{\pgfqpoint{5.632750in}{2.423560in}}%
\pgfpathlineto{\pgfqpoint{5.642239in}{2.452658in}}%
\pgfpathlineto{\pgfqpoint{5.651728in}{2.457507in}}%
\pgfpathlineto{\pgfqpoint{5.661217in}{2.481755in}}%
\pgfpathlineto{\pgfqpoint{5.670706in}{2.472056in}}%
\pgfpathlineto{\pgfqpoint{5.680195in}{2.544800in}}%
\pgfpathlineto{\pgfqpoint{5.689684in}{2.486302in}}%
\pgfpathlineto{\pgfqpoint{5.699172in}{2.457204in}}%
\pgfpathlineto{\pgfqpoint{5.708661in}{2.452355in}}%
\pgfpathlineto{\pgfqpoint{5.718150in}{2.452355in}}%
\pgfpathlineto{\pgfqpoint{5.727639in}{2.496001in}}%
\pgfpathlineto{\pgfqpoint{5.737128in}{2.496001in}}%
\pgfpathlineto{\pgfqpoint{5.746617in}{2.423257in}}%
\pgfpathlineto{\pgfqpoint{5.756106in}{2.437806in}}%
\pgfpathlineto{\pgfqpoint{5.765595in}{2.437503in}}%
\pgfpathlineto{\pgfqpoint{5.775083in}{2.466600in}}%
\pgfpathlineto{\pgfqpoint{5.784572in}{2.437503in}}%
\pgfpathlineto{\pgfqpoint{5.794061in}{2.422954in}}%
\pgfpathlineto{\pgfqpoint{5.803550in}{2.422954in}}%
\pgfpathlineto{\pgfqpoint{5.813039in}{2.481149in}}%
\pgfpathlineto{\pgfqpoint{5.822528in}{2.471450in}}%
\pgfpathlineto{\pgfqpoint{5.832017in}{2.466600in}}%
\pgfpathlineto{\pgfqpoint{5.841506in}{2.456598in}}%
\pgfpathlineto{\pgfqpoint{5.850995in}{2.466297in}}%
\pgfpathlineto{\pgfqpoint{5.860483in}{2.422651in}}%
\pgfpathlineto{\pgfqpoint{5.869972in}{2.456598in}}%
\pgfpathlineto{\pgfqpoint{5.879461in}{2.422651in}}%
\pgfpathlineto{\pgfqpoint{5.888950in}{2.475996in}}%
\pgfpathlineto{\pgfqpoint{5.898439in}{2.514793in}}%
\pgfpathlineto{\pgfqpoint{5.907928in}{2.437200in}}%
\pgfpathlineto{\pgfqpoint{5.917417in}{2.422348in}}%
\pgfpathlineto{\pgfqpoint{5.926906in}{2.451445in}}%
\pgfpathlineto{\pgfqpoint{5.936394in}{2.436896in}}%
\pgfpathlineto{\pgfqpoint{5.945883in}{2.470844in}}%
\pgfpathlineto{\pgfqpoint{5.945883in}{2.470844in}}%
\pgfusepath{stroke}%
\end{pgfscope}%
\begin{pgfscope}%
\pgfsetrectcap%
\pgfsetmiterjoin%
\pgfsetlinewidth{0.803000pt}%
\definecolor{currentstroke}{rgb}{0.000000,0.000000,0.000000}%
\pgfsetstrokecolor{currentstroke}%
\pgfsetdash{}{0pt}%
\pgfpathmoveto{\pgfqpoint{3.833026in}{2.240123in}}%
\pgfpathlineto{\pgfqpoint{3.833026in}{3.141667in}}%
\pgfusepath{stroke}%
\end{pgfscope}%
\begin{pgfscope}%
\pgfsetrectcap%
\pgfsetmiterjoin%
\pgfsetlinewidth{0.803000pt}%
\definecolor{currentstroke}{rgb}{0.000000,0.000000,0.000000}%
\pgfsetstrokecolor{currentstroke}%
\pgfsetdash{}{0pt}%
\pgfpathmoveto{\pgfqpoint{5.941666in}{2.240123in}}%
\pgfpathlineto{\pgfqpoint{5.941666in}{3.141667in}}%
\pgfusepath{stroke}%
\end{pgfscope}%
\begin{pgfscope}%
\pgfsetrectcap%
\pgfsetmiterjoin%
\pgfsetlinewidth{0.803000pt}%
\definecolor{currentstroke}{rgb}{0.000000,0.000000,0.000000}%
\pgfsetstrokecolor{currentstroke}%
\pgfsetdash{}{0pt}%
\pgfpathmoveto{\pgfqpoint{3.833026in}{2.240123in}}%
\pgfpathlineto{\pgfqpoint{5.941666in}{2.240123in}}%
\pgfusepath{stroke}%
\end{pgfscope}%
\begin{pgfscope}%
\pgfsetrectcap%
\pgfsetmiterjoin%
\pgfsetlinewidth{0.803000pt}%
\definecolor{currentstroke}{rgb}{0.000000,0.000000,0.000000}%
\pgfsetstrokecolor{currentstroke}%
\pgfsetdash{}{0pt}%
\pgfpathmoveto{\pgfqpoint{3.833026in}{3.141667in}}%
\pgfpathlineto{\pgfqpoint{5.941666in}{3.141667in}}%
\pgfusepath{stroke}%
\end{pgfscope}%
\begin{pgfscope}%
\definecolor{textcolor}{rgb}{0.000000,0.000000,0.000000}%
\pgfsetstrokecolor{textcolor}%
\pgfsetfillcolor{textcolor}%
\pgftext[x=4.887346in,y=3.225000in,,base]{\color{textcolor}\rmfamily\fontsize{12.000000}{14.400000}\selectfont 50 kHz, 20\%}%
\end{pgfscope}%
\begin{pgfscope}%
\pgfsetbuttcap%
\pgfsetmiterjoin%
\definecolor{currentfill}{rgb}{1.000000,1.000000,1.000000}%
\pgfsetfillcolor{currentfill}%
\pgfsetlinewidth{0.000000pt}%
\definecolor{currentstroke}{rgb}{0.000000,0.000000,0.000000}%
\pgfsetstrokecolor{currentstroke}%
\pgfsetstrokeopacity{0.000000}%
\pgfsetdash{}{0pt}%
\pgfpathmoveto{\pgfqpoint{0.758026in}{0.565123in}}%
\pgfpathlineto{\pgfqpoint{2.866666in}{0.565123in}}%
\pgfpathlineto{\pgfqpoint{2.866666in}{1.466667in}}%
\pgfpathlineto{\pgfqpoint{0.758026in}{1.466667in}}%
\pgfpathlineto{\pgfqpoint{0.758026in}{0.565123in}}%
\pgfpathclose%
\pgfusepath{fill}%
\end{pgfscope}%
\begin{pgfscope}%
\pgfpathrectangle{\pgfqpoint{0.758026in}{0.565123in}}{\pgfqpoint{2.108640in}{0.901543in}}%
\pgfusepath{clip}%
\pgfsetbuttcap%
\pgfsetroundjoin%
\pgfsetlinewidth{0.501875pt}%
\definecolor{currentstroke}{rgb}{0.690196,0.690196,0.690196}%
\pgfsetstrokecolor{currentstroke}%
\pgfsetstrokeopacity{0.700000}%
\pgfsetdash{{1.850000pt}{0.800000pt}}{0.000000pt}%
\pgfpathmoveto{\pgfqpoint{0.758026in}{0.565123in}}%
\pgfpathlineto{\pgfqpoint{0.758026in}{1.466667in}}%
\pgfusepath{stroke}%
\end{pgfscope}%
\begin{pgfscope}%
\pgfsetbuttcap%
\pgfsetroundjoin%
\definecolor{currentfill}{rgb}{0.000000,0.000000,0.000000}%
\pgfsetfillcolor{currentfill}%
\pgfsetlinewidth{0.803000pt}%
\definecolor{currentstroke}{rgb}{0.000000,0.000000,0.000000}%
\pgfsetstrokecolor{currentstroke}%
\pgfsetdash{}{0pt}%
\pgfsys@defobject{currentmarker}{\pgfqpoint{0.000000in}{-0.048611in}}{\pgfqpoint{0.000000in}{0.000000in}}{%
\pgfpathmoveto{\pgfqpoint{0.000000in}{0.000000in}}%
\pgfpathlineto{\pgfqpoint{0.000000in}{-0.048611in}}%
\pgfusepath{stroke,fill}%
}%
\begin{pgfscope}%
\pgfsys@transformshift{0.758026in}{0.565123in}%
\pgfsys@useobject{currentmarker}{}%
\end{pgfscope}%
\end{pgfscope}%
\begin{pgfscope}%
\definecolor{textcolor}{rgb}{0.000000,0.000000,0.000000}%
\pgfsetstrokecolor{textcolor}%
\pgfsetfillcolor{textcolor}%
\pgftext[x=0.758026in,y=0.467901in,,top]{\color{textcolor}\rmfamily\fontsize{10.000000}{12.000000}\selectfont \(\displaystyle {400000}\)}%
\end{pgfscope}%
\begin{pgfscope}%
\pgfpathrectangle{\pgfqpoint{0.758026in}{0.565123in}}{\pgfqpoint{2.108640in}{0.901543in}}%
\pgfusepath{clip}%
\pgfsetbuttcap%
\pgfsetroundjoin%
\pgfsetlinewidth{0.501875pt}%
\definecolor{currentstroke}{rgb}{0.690196,0.690196,0.690196}%
\pgfsetstrokecolor{currentstroke}%
\pgfsetstrokeopacity{0.700000}%
\pgfsetdash{{1.850000pt}{0.800000pt}}{0.000000pt}%
\pgfpathmoveto{\pgfqpoint{1.285186in}{0.565123in}}%
\pgfpathlineto{\pgfqpoint{1.285186in}{1.466667in}}%
\pgfusepath{stroke}%
\end{pgfscope}%
\begin{pgfscope}%
\pgfsetbuttcap%
\pgfsetroundjoin%
\definecolor{currentfill}{rgb}{0.000000,0.000000,0.000000}%
\pgfsetfillcolor{currentfill}%
\pgfsetlinewidth{0.803000pt}%
\definecolor{currentstroke}{rgb}{0.000000,0.000000,0.000000}%
\pgfsetstrokecolor{currentstroke}%
\pgfsetdash{}{0pt}%
\pgfsys@defobject{currentmarker}{\pgfqpoint{0.000000in}{-0.048611in}}{\pgfqpoint{0.000000in}{0.000000in}}{%
\pgfpathmoveto{\pgfqpoint{0.000000in}{0.000000in}}%
\pgfpathlineto{\pgfqpoint{0.000000in}{-0.048611in}}%
\pgfusepath{stroke,fill}%
}%
\begin{pgfscope}%
\pgfsys@transformshift{1.285186in}{0.565123in}%
\pgfsys@useobject{currentmarker}{}%
\end{pgfscope}%
\end{pgfscope}%
\begin{pgfscope}%
\definecolor{textcolor}{rgb}{0.000000,0.000000,0.000000}%
\pgfsetstrokecolor{textcolor}%
\pgfsetfillcolor{textcolor}%
\pgftext[x=1.285186in,y=0.467901in,,top]{\color{textcolor}\rmfamily\fontsize{10.000000}{12.000000}\selectfont \(\displaystyle {450000}\)}%
\end{pgfscope}%
\begin{pgfscope}%
\pgfpathrectangle{\pgfqpoint{0.758026in}{0.565123in}}{\pgfqpoint{2.108640in}{0.901543in}}%
\pgfusepath{clip}%
\pgfsetbuttcap%
\pgfsetroundjoin%
\pgfsetlinewidth{0.501875pt}%
\definecolor{currentstroke}{rgb}{0.690196,0.690196,0.690196}%
\pgfsetstrokecolor{currentstroke}%
\pgfsetstrokeopacity{0.700000}%
\pgfsetdash{{1.850000pt}{0.800000pt}}{0.000000pt}%
\pgfpathmoveto{\pgfqpoint{1.812346in}{0.565123in}}%
\pgfpathlineto{\pgfqpoint{1.812346in}{1.466667in}}%
\pgfusepath{stroke}%
\end{pgfscope}%
\begin{pgfscope}%
\pgfsetbuttcap%
\pgfsetroundjoin%
\definecolor{currentfill}{rgb}{0.000000,0.000000,0.000000}%
\pgfsetfillcolor{currentfill}%
\pgfsetlinewidth{0.803000pt}%
\definecolor{currentstroke}{rgb}{0.000000,0.000000,0.000000}%
\pgfsetstrokecolor{currentstroke}%
\pgfsetdash{}{0pt}%
\pgfsys@defobject{currentmarker}{\pgfqpoint{0.000000in}{-0.048611in}}{\pgfqpoint{0.000000in}{0.000000in}}{%
\pgfpathmoveto{\pgfqpoint{0.000000in}{0.000000in}}%
\pgfpathlineto{\pgfqpoint{0.000000in}{-0.048611in}}%
\pgfusepath{stroke,fill}%
}%
\begin{pgfscope}%
\pgfsys@transformshift{1.812346in}{0.565123in}%
\pgfsys@useobject{currentmarker}{}%
\end{pgfscope}%
\end{pgfscope}%
\begin{pgfscope}%
\definecolor{textcolor}{rgb}{0.000000,0.000000,0.000000}%
\pgfsetstrokecolor{textcolor}%
\pgfsetfillcolor{textcolor}%
\pgftext[x=1.812346in,y=0.467901in,,top]{\color{textcolor}\rmfamily\fontsize{10.000000}{12.000000}\selectfont \(\displaystyle {500000}\)}%
\end{pgfscope}%
\begin{pgfscope}%
\pgfpathrectangle{\pgfqpoint{0.758026in}{0.565123in}}{\pgfqpoint{2.108640in}{0.901543in}}%
\pgfusepath{clip}%
\pgfsetbuttcap%
\pgfsetroundjoin%
\pgfsetlinewidth{0.501875pt}%
\definecolor{currentstroke}{rgb}{0.690196,0.690196,0.690196}%
\pgfsetstrokecolor{currentstroke}%
\pgfsetstrokeopacity{0.700000}%
\pgfsetdash{{1.850000pt}{0.800000pt}}{0.000000pt}%
\pgfpathmoveto{\pgfqpoint{2.339506in}{0.565123in}}%
\pgfpathlineto{\pgfqpoint{2.339506in}{1.466667in}}%
\pgfusepath{stroke}%
\end{pgfscope}%
\begin{pgfscope}%
\pgfsetbuttcap%
\pgfsetroundjoin%
\definecolor{currentfill}{rgb}{0.000000,0.000000,0.000000}%
\pgfsetfillcolor{currentfill}%
\pgfsetlinewidth{0.803000pt}%
\definecolor{currentstroke}{rgb}{0.000000,0.000000,0.000000}%
\pgfsetstrokecolor{currentstroke}%
\pgfsetdash{}{0pt}%
\pgfsys@defobject{currentmarker}{\pgfqpoint{0.000000in}{-0.048611in}}{\pgfqpoint{0.000000in}{0.000000in}}{%
\pgfpathmoveto{\pgfqpoint{0.000000in}{0.000000in}}%
\pgfpathlineto{\pgfqpoint{0.000000in}{-0.048611in}}%
\pgfusepath{stroke,fill}%
}%
\begin{pgfscope}%
\pgfsys@transformshift{2.339506in}{0.565123in}%
\pgfsys@useobject{currentmarker}{}%
\end{pgfscope}%
\end{pgfscope}%
\begin{pgfscope}%
\definecolor{textcolor}{rgb}{0.000000,0.000000,0.000000}%
\pgfsetstrokecolor{textcolor}%
\pgfsetfillcolor{textcolor}%
\pgftext[x=2.339506in,y=0.467901in,,top]{\color{textcolor}\rmfamily\fontsize{10.000000}{12.000000}\selectfont \(\displaystyle {550000}\)}%
\end{pgfscope}%
\begin{pgfscope}%
\pgfpathrectangle{\pgfqpoint{0.758026in}{0.565123in}}{\pgfqpoint{2.108640in}{0.901543in}}%
\pgfusepath{clip}%
\pgfsetbuttcap%
\pgfsetroundjoin%
\pgfsetlinewidth{0.501875pt}%
\definecolor{currentstroke}{rgb}{0.690196,0.690196,0.690196}%
\pgfsetstrokecolor{currentstroke}%
\pgfsetstrokeopacity{0.700000}%
\pgfsetdash{{1.850000pt}{0.800000pt}}{0.000000pt}%
\pgfpathmoveto{\pgfqpoint{2.866666in}{0.565123in}}%
\pgfpathlineto{\pgfqpoint{2.866666in}{1.466667in}}%
\pgfusepath{stroke}%
\end{pgfscope}%
\begin{pgfscope}%
\pgfsetbuttcap%
\pgfsetroundjoin%
\definecolor{currentfill}{rgb}{0.000000,0.000000,0.000000}%
\pgfsetfillcolor{currentfill}%
\pgfsetlinewidth{0.803000pt}%
\definecolor{currentstroke}{rgb}{0.000000,0.000000,0.000000}%
\pgfsetstrokecolor{currentstroke}%
\pgfsetdash{}{0pt}%
\pgfsys@defobject{currentmarker}{\pgfqpoint{0.000000in}{-0.048611in}}{\pgfqpoint{0.000000in}{0.000000in}}{%
\pgfpathmoveto{\pgfqpoint{0.000000in}{0.000000in}}%
\pgfpathlineto{\pgfqpoint{0.000000in}{-0.048611in}}%
\pgfusepath{stroke,fill}%
}%
\begin{pgfscope}%
\pgfsys@transformshift{2.866666in}{0.565123in}%
\pgfsys@useobject{currentmarker}{}%
\end{pgfscope}%
\end{pgfscope}%
\begin{pgfscope}%
\definecolor{textcolor}{rgb}{0.000000,0.000000,0.000000}%
\pgfsetstrokecolor{textcolor}%
\pgfsetfillcolor{textcolor}%
\pgftext[x=2.866666in,y=0.467901in,,top]{\color{textcolor}\rmfamily\fontsize{10.000000}{12.000000}\selectfont \(\displaystyle {600000}\)}%
\end{pgfscope}%
\begin{pgfscope}%
\definecolor{textcolor}{rgb}{0.000000,0.000000,0.000000}%
\pgfsetstrokecolor{textcolor}%
\pgfsetfillcolor{textcolor}%
\pgftext[x=1.812346in,y=0.288889in,,top]{\color{textcolor}\rmfamily\fontsize{10.000000}{12.000000}\selectfont Frequency (Hz)}%
\end{pgfscope}%
\begin{pgfscope}%
\pgfpathrectangle{\pgfqpoint{0.758026in}{0.565123in}}{\pgfqpoint{2.108640in}{0.901543in}}%
\pgfusepath{clip}%
\pgfsetbuttcap%
\pgfsetroundjoin%
\pgfsetlinewidth{0.501875pt}%
\definecolor{currentstroke}{rgb}{0.690196,0.690196,0.690196}%
\pgfsetstrokecolor{currentstroke}%
\pgfsetstrokeopacity{0.700000}%
\pgfsetdash{{1.850000pt}{0.800000pt}}{0.000000pt}%
\pgfpathmoveto{\pgfqpoint{0.758026in}{0.695038in}}%
\pgfpathlineto{\pgfqpoint{2.866666in}{0.695038in}}%
\pgfusepath{stroke}%
\end{pgfscope}%
\begin{pgfscope}%
\pgfsetbuttcap%
\pgfsetroundjoin%
\definecolor{currentfill}{rgb}{0.000000,0.000000,0.000000}%
\pgfsetfillcolor{currentfill}%
\pgfsetlinewidth{0.803000pt}%
\definecolor{currentstroke}{rgb}{0.000000,0.000000,0.000000}%
\pgfsetstrokecolor{currentstroke}%
\pgfsetdash{}{0pt}%
\pgfsys@defobject{currentmarker}{\pgfqpoint{-0.048611in}{0.000000in}}{\pgfqpoint{-0.000000in}{0.000000in}}{%
\pgfpathmoveto{\pgfqpoint{-0.000000in}{0.000000in}}%
\pgfpathlineto{\pgfqpoint{-0.048611in}{0.000000in}}%
\pgfusepath{stroke,fill}%
}%
\begin{pgfscope}%
\pgfsys@transformshift{0.758026in}{0.695038in}%
\pgfsys@useobject{currentmarker}{}%
\end{pgfscope}%
\end{pgfscope}%
\begin{pgfscope}%
\definecolor{textcolor}{rgb}{0.000000,0.000000,0.000000}%
\pgfsetstrokecolor{textcolor}%
\pgfsetfillcolor{textcolor}%
\pgftext[x=0.344444in, y=0.646812in, left, base]{\color{textcolor}\rmfamily\fontsize{10.000000}{12.000000}\selectfont \(\displaystyle {\ensuremath{-}100}\)}%
\end{pgfscope}%
\begin{pgfscope}%
\pgfpathrectangle{\pgfqpoint{0.758026in}{0.565123in}}{\pgfqpoint{2.108640in}{0.901543in}}%
\pgfusepath{clip}%
\pgfsetbuttcap%
\pgfsetroundjoin%
\pgfsetlinewidth{0.501875pt}%
\definecolor{currentstroke}{rgb}{0.690196,0.690196,0.690196}%
\pgfsetstrokecolor{currentstroke}%
\pgfsetstrokeopacity{0.700000}%
\pgfsetdash{{1.850000pt}{0.800000pt}}{0.000000pt}%
\pgfpathmoveto{\pgfqpoint{0.758026in}{1.197850in}}%
\pgfpathlineto{\pgfqpoint{2.866666in}{1.197850in}}%
\pgfusepath{stroke}%
\end{pgfscope}%
\begin{pgfscope}%
\pgfsetbuttcap%
\pgfsetroundjoin%
\definecolor{currentfill}{rgb}{0.000000,0.000000,0.000000}%
\pgfsetfillcolor{currentfill}%
\pgfsetlinewidth{0.803000pt}%
\definecolor{currentstroke}{rgb}{0.000000,0.000000,0.000000}%
\pgfsetstrokecolor{currentstroke}%
\pgfsetdash{}{0pt}%
\pgfsys@defobject{currentmarker}{\pgfqpoint{-0.048611in}{0.000000in}}{\pgfqpoint{-0.000000in}{0.000000in}}{%
\pgfpathmoveto{\pgfqpoint{-0.000000in}{0.000000in}}%
\pgfpathlineto{\pgfqpoint{-0.048611in}{0.000000in}}%
\pgfusepath{stroke,fill}%
}%
\begin{pgfscope}%
\pgfsys@transformshift{0.758026in}{1.197850in}%
\pgfsys@useobject{currentmarker}{}%
\end{pgfscope}%
\end{pgfscope}%
\begin{pgfscope}%
\definecolor{textcolor}{rgb}{0.000000,0.000000,0.000000}%
\pgfsetstrokecolor{textcolor}%
\pgfsetfillcolor{textcolor}%
\pgftext[x=0.413889in, y=1.149625in, left, base]{\color{textcolor}\rmfamily\fontsize{10.000000}{12.000000}\selectfont \(\displaystyle {\ensuremath{-}50}\)}%
\end{pgfscope}%
\begin{pgfscope}%
\definecolor{textcolor}{rgb}{0.000000,0.000000,0.000000}%
\pgfsetstrokecolor{textcolor}%
\pgfsetfillcolor{textcolor}%
\pgftext[x=0.288889in,y=1.015895in,,bottom,rotate=90.000000]{\color{textcolor}\rmfamily\fontsize{10.000000}{12.000000}\selectfont Amplitude (dB)}%
\end{pgfscope}%
\begin{pgfscope}%
\pgfpathrectangle{\pgfqpoint{0.758026in}{0.565123in}}{\pgfqpoint{2.108640in}{0.901543in}}%
\pgfusepath{clip}%
\pgfsetrectcap%
\pgfsetroundjoin%
\pgfsetlinewidth{1.003750pt}%
\definecolor{currentstroke}{rgb}{0.000000,0.000000,0.000000}%
\pgfsetstrokecolor{currentstroke}%
\pgfsetdash{}{0pt}%
\pgfpathmoveto{\pgfqpoint{0.754863in}{0.902134in}}%
\pgfpathlineto{\pgfqpoint{0.764352in}{0.866937in}}%
\pgfpathlineto{\pgfqpoint{0.773840in}{0.866937in}}%
\pgfpathlineto{\pgfqpoint{0.783329in}{0.967499in}}%
\pgfpathlineto{\pgfqpoint{0.792818in}{0.927274in}}%
\pgfpathlineto{\pgfqpoint{0.802307in}{0.866622in}}%
\pgfpathlineto{\pgfqpoint{0.811796in}{0.967185in}}%
\pgfpathlineto{\pgfqpoint{0.821285in}{0.811313in}}%
\pgfpathlineto{\pgfqpoint{0.830774in}{0.841482in}}%
\pgfpathlineto{\pgfqpoint{0.840263in}{0.876679in}}%
\pgfpathlineto{\pgfqpoint{0.849752in}{0.831426in}}%
\pgfpathlineto{\pgfqpoint{0.859240in}{0.911876in}}%
\pgfpathlineto{\pgfqpoint{0.868729in}{0.891763in}}%
\pgfpathlineto{\pgfqpoint{0.887707in}{0.871336in}}%
\pgfpathlineto{\pgfqpoint{0.897196in}{0.881393in}}%
\pgfpathlineto{\pgfqpoint{0.906685in}{0.876364in}}%
\pgfpathlineto{\pgfqpoint{0.916174in}{0.881393in}}%
\pgfpathlineto{\pgfqpoint{0.925663in}{0.846196in}}%
\pgfpathlineto{\pgfqpoint{0.935151in}{0.861280in}}%
\pgfpathlineto{\pgfqpoint{0.944640in}{0.896477in}}%
\pgfpathlineto{\pgfqpoint{0.954129in}{0.840853in}}%
\pgfpathlineto{\pgfqpoint{0.963618in}{0.901191in}}%
\pgfpathlineto{\pgfqpoint{0.973107in}{0.765431in}}%
\pgfpathlineto{\pgfqpoint{0.982596in}{0.886106in}}%
\pgfpathlineto{\pgfqpoint{0.992085in}{0.891135in}}%
\pgfpathlineto{\pgfqpoint{1.001574in}{0.820741in}}%
\pgfpathlineto{\pgfqpoint{1.011063in}{0.865994in}}%
\pgfpathlineto{\pgfqpoint{1.020551in}{0.871022in}}%
\pgfpathlineto{\pgfqpoint{1.030040in}{0.845567in}}%
\pgfpathlineto{\pgfqpoint{1.039529in}{0.870708in}}%
\pgfpathlineto{\pgfqpoint{1.049018in}{0.850595in}}%
\pgfpathlineto{\pgfqpoint{1.058507in}{0.900877in}}%
\pgfpathlineto{\pgfqpoint{1.067996in}{0.900877in}}%
\pgfpathlineto{\pgfqpoint{1.077485in}{0.910933in}}%
\pgfpathlineto{\pgfqpoint{1.086974in}{0.830483in}}%
\pgfpathlineto{\pgfqpoint{1.096462in}{0.865680in}}%
\pgfpathlineto{\pgfqpoint{1.105951in}{0.810056in}}%
\pgfpathlineto{\pgfqpoint{1.115440in}{0.905590in}}%
\pgfpathlineto{\pgfqpoint{1.124929in}{0.850281in}}%
\pgfpathlineto{\pgfqpoint{1.134418in}{0.875422in}}%
\pgfpathlineto{\pgfqpoint{1.143907in}{0.850281in}}%
\pgfpathlineto{\pgfqpoint{1.153396in}{0.875422in}}%
\pgfpathlineto{\pgfqpoint{1.162885in}{0.865365in}}%
\pgfpathlineto{\pgfqpoint{1.172374in}{0.905590in}}%
\pgfpathlineto{\pgfqpoint{1.181862in}{0.885164in}}%
\pgfpathlineto{\pgfqpoint{1.191351in}{0.849967in}}%
\pgfpathlineto{\pgfqpoint{1.200840in}{0.784601in}}%
\pgfpathlineto{\pgfqpoint{1.210329in}{0.895220in}}%
\pgfpathlineto{\pgfqpoint{1.219818in}{0.870079in}}%
\pgfpathlineto{\pgfqpoint{1.229307in}{0.865051in}}%
\pgfpathlineto{\pgfqpoint{1.248285in}{0.884849in}}%
\pgfpathlineto{\pgfqpoint{1.257773in}{0.839596in}}%
\pgfpathlineto{\pgfqpoint{1.267262in}{0.904962in}}%
\pgfpathlineto{\pgfqpoint{1.276751in}{0.894906in}}%
\pgfpathlineto{\pgfqpoint{1.286240in}{0.879821in}}%
\pgfpathlineto{\pgfqpoint{1.295729in}{0.799371in}}%
\pgfpathlineto{\pgfqpoint{1.305218in}{0.879821in}}%
\pgfpathlineto{\pgfqpoint{1.314707in}{0.859709in}}%
\pgfpathlineto{\pgfqpoint{1.324196in}{0.844310in}}%
\pgfpathlineto{\pgfqpoint{1.333685in}{0.879507in}}%
\pgfpathlineto{\pgfqpoint{1.343173in}{0.829226in}}%
\pgfpathlineto{\pgfqpoint{1.352662in}{0.919732in}}%
\pgfpathlineto{\pgfqpoint{1.362151in}{0.859395in}}%
\pgfpathlineto{\pgfqpoint{1.371640in}{0.864423in}}%
\pgfpathlineto{\pgfqpoint{1.381129in}{0.829226in}}%
\pgfpathlineto{\pgfqpoint{1.390618in}{0.879507in}}%
\pgfpathlineto{\pgfqpoint{1.400107in}{0.843996in}}%
\pgfpathlineto{\pgfqpoint{1.409596in}{0.854052in}}%
\pgfpathlineto{\pgfqpoint{1.419084in}{0.849024in}}%
\pgfpathlineto{\pgfqpoint{1.428573in}{0.859080in}}%
\pgfpathlineto{\pgfqpoint{1.438062in}{0.899305in}}%
\pgfpathlineto{\pgfqpoint{1.447551in}{0.833940in}}%
\pgfpathlineto{\pgfqpoint{1.457040in}{0.919418in}}%
\pgfpathlineto{\pgfqpoint{1.466529in}{0.808799in}}%
\pgfpathlineto{\pgfqpoint{1.476018in}{0.843682in}}%
\pgfpathlineto{\pgfqpoint{1.485507in}{0.883907in}}%
\pgfpathlineto{\pgfqpoint{1.494995in}{0.813513in}}%
\pgfpathlineto{\pgfqpoint{1.504484in}{1.004582in}}%
\pgfpathlineto{\pgfqpoint{1.513973in}{1.150397in}}%
\pgfpathlineto{\pgfqpoint{1.523462in}{1.235876in}}%
\pgfpathlineto{\pgfqpoint{1.532951in}{1.286157in}}%
\pgfpathlineto{\pgfqpoint{1.542440in}{1.321354in}}%
\pgfpathlineto{\pgfqpoint{1.551929in}{1.336124in}}%
\pgfpathlineto{\pgfqpoint{1.561418in}{1.326068in}}%
\pgfpathlineto{\pgfqpoint{1.570907in}{1.300927in}}%
\pgfpathlineto{\pgfqpoint{1.580395in}{1.250646in}}%
\pgfpathlineto{\pgfqpoint{1.589884in}{1.175224in}}%
\pgfpathlineto{\pgfqpoint{1.599373in}{1.059577in}}%
\pgfpathlineto{\pgfqpoint{1.608862in}{0.893649in}}%
\pgfpathlineto{\pgfqpoint{1.618351in}{0.893649in}}%
\pgfpathlineto{\pgfqpoint{1.627840in}{0.832997in}}%
\pgfpathlineto{\pgfqpoint{1.637329in}{0.838025in}}%
\pgfpathlineto{\pgfqpoint{1.646818in}{0.858137in}}%
\pgfpathlineto{\pgfqpoint{1.656306in}{0.782716in}}%
\pgfpathlineto{\pgfqpoint{1.665795in}{0.792772in}}%
\pgfpathlineto{\pgfqpoint{1.675284in}{0.817912in}}%
\pgfpathlineto{\pgfqpoint{1.684773in}{0.832997in}}%
\pgfpathlineto{\pgfqpoint{1.694262in}{0.762603in}}%
\pgfpathlineto{\pgfqpoint{1.703751in}{0.847767in}}%
\pgfpathlineto{\pgfqpoint{1.713240in}{0.872908in}}%
\pgfpathlineto{\pgfqpoint{1.722729in}{0.857823in}}%
\pgfpathlineto{\pgfqpoint{1.732218in}{0.857823in}}%
\pgfpathlineto{\pgfqpoint{1.741706in}{0.797486in}}%
\pgfpathlineto{\pgfqpoint{1.751195in}{0.887992in}}%
\pgfpathlineto{\pgfqpoint{1.760684in}{0.872908in}}%
\pgfpathlineto{\pgfqpoint{1.770173in}{1.134370in}}%
\pgfpathlineto{\pgfqpoint{1.779662in}{1.264787in}}%
\pgfpathlineto{\pgfqpoint{1.789151in}{1.340209in}}%
\pgfpathlineto{\pgfqpoint{1.798640in}{1.385462in}}%
\pgfpathlineto{\pgfqpoint{1.808129in}{1.415631in}}%
\pgfpathlineto{\pgfqpoint{1.817617in}{1.425687in}}%
\pgfpathlineto{\pgfqpoint{1.827106in}{1.415631in}}%
\pgfpathlineto{\pgfqpoint{1.836595in}{1.385462in}}%
\pgfpathlineto{\pgfqpoint{1.846084in}{1.330153in}}%
\pgfpathlineto{\pgfqpoint{1.855573in}{1.244361in}}%
\pgfpathlineto{\pgfqpoint{1.865062in}{1.113629in}}%
\pgfpathlineto{\pgfqpoint{1.874551in}{0.882335in}}%
\pgfpathlineto{\pgfqpoint{1.884040in}{0.867251in}}%
\pgfpathlineto{\pgfqpoint{1.893529in}{0.867251in}}%
\pgfpathlineto{\pgfqpoint{1.903017in}{0.837082in}}%
\pgfpathlineto{\pgfqpoint{1.912506in}{0.907476in}}%
\pgfpathlineto{\pgfqpoint{1.921995in}{0.796857in}}%
\pgfpathlineto{\pgfqpoint{1.931484in}{0.811627in}}%
\pgfpathlineto{\pgfqpoint{1.940973in}{0.821684in}}%
\pgfpathlineto{\pgfqpoint{1.950462in}{0.826712in}}%
\pgfpathlineto{\pgfqpoint{1.959951in}{0.856880in}}%
\pgfpathlineto{\pgfqpoint{1.969440in}{0.816655in}}%
\pgfpathlineto{\pgfqpoint{1.978928in}{0.892077in}}%
\pgfpathlineto{\pgfqpoint{1.988417in}{0.846824in}}%
\pgfpathlineto{\pgfqpoint{1.997906in}{0.791515in}}%
\pgfpathlineto{\pgfqpoint{2.007395in}{0.876679in}}%
\pgfpathlineto{\pgfqpoint{2.016884in}{0.796229in}}%
\pgfpathlineto{\pgfqpoint{2.026373in}{0.856566in}}%
\pgfpathlineto{\pgfqpoint{2.035862in}{1.077804in}}%
\pgfpathlineto{\pgfqpoint{2.045351in}{1.183395in}}%
\pgfpathlineto{\pgfqpoint{2.054840in}{1.258816in}}%
\pgfpathlineto{\pgfqpoint{2.064328in}{1.304070in}}%
\pgfpathlineto{\pgfqpoint{2.073817in}{1.329210in}}%
\pgfpathlineto{\pgfqpoint{2.083306in}{1.333924in}}%
\pgfpathlineto{\pgfqpoint{2.092795in}{1.313812in}}%
\pgfpathlineto{\pgfqpoint{2.102284in}{1.283643in}}%
\pgfpathlineto{\pgfqpoint{2.111773in}{1.228333in}}%
\pgfpathlineto{\pgfqpoint{2.121262in}{1.132799in}}%
\pgfpathlineto{\pgfqpoint{2.130751in}{0.992011in}}%
\pgfpathlineto{\pgfqpoint{2.140239in}{0.876364in}}%
\pgfpathlineto{\pgfqpoint{2.149728in}{0.891449in}}%
\pgfpathlineto{\pgfqpoint{2.159217in}{0.780516in}}%
\pgfpathlineto{\pgfqpoint{2.168706in}{0.845881in}}%
\pgfpathlineto{\pgfqpoint{2.178195in}{0.886106in}}%
\pgfpathlineto{\pgfqpoint{2.187684in}{0.825769in}}%
\pgfpathlineto{\pgfqpoint{2.197173in}{0.855938in}}%
\pgfpathlineto{\pgfqpoint{2.206662in}{0.820741in}}%
\pgfpathlineto{\pgfqpoint{2.216150in}{0.815713in}}%
\pgfpathlineto{\pgfqpoint{2.225639in}{0.881078in}}%
\pgfpathlineto{\pgfqpoint{2.235128in}{0.875736in}}%
\pgfpathlineto{\pgfqpoint{2.254106in}{0.775173in}}%
\pgfpathlineto{\pgfqpoint{2.263595in}{0.895848in}}%
\pgfpathlineto{\pgfqpoint{2.273084in}{0.775173in}}%
\pgfpathlineto{\pgfqpoint{2.282573in}{0.865680in}}%
\pgfpathlineto{\pgfqpoint{2.292062in}{0.795286in}}%
\pgfpathlineto{\pgfqpoint{2.311039in}{0.764803in}}%
\pgfpathlineto{\pgfqpoint{2.320528in}{0.815084in}}%
\pgfpathlineto{\pgfqpoint{2.330017in}{0.825140in}}%
\pgfpathlineto{\pgfqpoint{2.339506in}{0.830169in}}%
\pgfpathlineto{\pgfqpoint{2.348995in}{0.774859in}}%
\pgfpathlineto{\pgfqpoint{2.358484in}{0.920675in}}%
\pgfpathlineto{\pgfqpoint{2.367973in}{0.800000in}}%
\pgfpathlineto{\pgfqpoint{2.377461in}{0.779887in}}%
\pgfpathlineto{\pgfqpoint{2.386950in}{0.789629in}}%
\pgfpathlineto{\pgfqpoint{2.396439in}{0.744376in}}%
\pgfpathlineto{\pgfqpoint{2.405928in}{0.880136in}}%
\pgfpathlineto{\pgfqpoint{2.415417in}{0.824826in}}%
\pgfpathlineto{\pgfqpoint{2.424906in}{0.744376in}}%
\pgfpathlineto{\pgfqpoint{2.434395in}{0.774545in}}%
\pgfpathlineto{\pgfqpoint{2.443884in}{0.744376in}}%
\pgfpathlineto{\pgfqpoint{2.453373in}{0.774545in}}%
\pgfpathlineto{\pgfqpoint{2.462861in}{0.784287in}}%
\pgfpathlineto{\pgfqpoint{2.472350in}{0.779259in}}%
\pgfpathlineto{\pgfqpoint{2.481839in}{0.728977in}}%
\pgfpathlineto{\pgfqpoint{2.491328in}{0.744062in}}%
\pgfpathlineto{\pgfqpoint{2.500817in}{0.844624in}}%
\pgfpathlineto{\pgfqpoint{2.510306in}{0.744062in}}%
\pgfpathlineto{\pgfqpoint{2.519795in}{0.774231in}}%
\pgfpathlineto{\pgfqpoint{2.529284in}{0.774231in}}%
\pgfpathlineto{\pgfqpoint{2.538772in}{0.743748in}}%
\pgfpathlineto{\pgfqpoint{2.548261in}{0.743748in}}%
\pgfpathlineto{\pgfqpoint{2.557750in}{0.773916in}}%
\pgfpathlineto{\pgfqpoint{2.567239in}{0.778944in}}%
\pgfpathlineto{\pgfqpoint{2.576728in}{0.763860in}}%
\pgfpathlineto{\pgfqpoint{2.586217in}{0.829226in}}%
\pgfpathlineto{\pgfqpoint{2.595706in}{0.773916in}}%
\pgfpathlineto{\pgfqpoint{2.605195in}{0.743748in}}%
\pgfpathlineto{\pgfqpoint{2.614684in}{0.773602in}}%
\pgfpathlineto{\pgfqpoint{2.624172in}{0.763546in}}%
\pgfpathlineto{\pgfqpoint{2.633661in}{0.818855in}}%
\pgfpathlineto{\pgfqpoint{2.643150in}{0.838968in}}%
\pgfpathlineto{\pgfqpoint{2.652639in}{0.748461in}}%
\pgfpathlineto{\pgfqpoint{2.662128in}{0.808799in}}%
\pgfpathlineto{\pgfqpoint{2.671617in}{0.728349in}}%
\pgfpathlineto{\pgfqpoint{2.681106in}{0.728349in}}%
\pgfpathlineto{\pgfqpoint{2.690595in}{0.803457in}}%
\pgfpathlineto{\pgfqpoint{2.700083in}{0.788372in}}%
\pgfpathlineto{\pgfqpoint{2.709572in}{0.743119in}}%
\pgfpathlineto{\pgfqpoint{2.719061in}{0.793400in}}%
\pgfpathlineto{\pgfqpoint{2.728550in}{0.743119in}}%
\pgfpathlineto{\pgfqpoint{2.738039in}{0.773288in}}%
\pgfpathlineto{\pgfqpoint{2.747528in}{0.763232in}}%
\pgfpathlineto{\pgfqpoint{2.757017in}{0.743119in}}%
\pgfpathlineto{\pgfqpoint{2.766506in}{0.757889in}}%
\pgfpathlineto{\pgfqpoint{2.775995in}{0.727720in}}%
\pgfpathlineto{\pgfqpoint{2.785483in}{0.788058in}}%
\pgfpathlineto{\pgfqpoint{2.794972in}{0.793086in}}%
\pgfpathlineto{\pgfqpoint{2.804461in}{0.742805in}}%
\pgfpathlineto{\pgfqpoint{2.813950in}{0.818227in}}%
\pgfpathlineto{\pgfqpoint{2.823439in}{0.742805in}}%
\pgfpathlineto{\pgfqpoint{2.832928in}{0.788058in}}%
\pgfpathlineto{\pgfqpoint{2.842417in}{0.757575in}}%
\pgfpathlineto{\pgfqpoint{2.851906in}{0.792772in}}%
\pgfpathlineto{\pgfqpoint{2.861394in}{0.777687in}}%
\pgfpathlineto{\pgfqpoint{2.870883in}{0.747519in}}%
\pgfpathlineto{\pgfqpoint{2.870883in}{0.747519in}}%
\pgfusepath{stroke}%
\end{pgfscope}%
\begin{pgfscope}%
\pgfsetrectcap%
\pgfsetmiterjoin%
\pgfsetlinewidth{0.803000pt}%
\definecolor{currentstroke}{rgb}{0.000000,0.000000,0.000000}%
\pgfsetstrokecolor{currentstroke}%
\pgfsetdash{}{0pt}%
\pgfpathmoveto{\pgfqpoint{0.758026in}{0.565123in}}%
\pgfpathlineto{\pgfqpoint{0.758026in}{1.466667in}}%
\pgfusepath{stroke}%
\end{pgfscope}%
\begin{pgfscope}%
\pgfsetrectcap%
\pgfsetmiterjoin%
\pgfsetlinewidth{0.803000pt}%
\definecolor{currentstroke}{rgb}{0.000000,0.000000,0.000000}%
\pgfsetstrokecolor{currentstroke}%
\pgfsetdash{}{0pt}%
\pgfpathmoveto{\pgfqpoint{2.866666in}{0.565123in}}%
\pgfpathlineto{\pgfqpoint{2.866666in}{1.466667in}}%
\pgfusepath{stroke}%
\end{pgfscope}%
\begin{pgfscope}%
\pgfsetrectcap%
\pgfsetmiterjoin%
\pgfsetlinewidth{0.803000pt}%
\definecolor{currentstroke}{rgb}{0.000000,0.000000,0.000000}%
\pgfsetstrokecolor{currentstroke}%
\pgfsetdash{}{0pt}%
\pgfpathmoveto{\pgfqpoint{0.758026in}{0.565123in}}%
\pgfpathlineto{\pgfqpoint{2.866666in}{0.565123in}}%
\pgfusepath{stroke}%
\end{pgfscope}%
\begin{pgfscope}%
\pgfsetrectcap%
\pgfsetmiterjoin%
\pgfsetlinewidth{0.803000pt}%
\definecolor{currentstroke}{rgb}{0.000000,0.000000,0.000000}%
\pgfsetstrokecolor{currentstroke}%
\pgfsetdash{}{0pt}%
\pgfpathmoveto{\pgfqpoint{0.758026in}{1.466667in}}%
\pgfpathlineto{\pgfqpoint{2.866666in}{1.466667in}}%
\pgfusepath{stroke}%
\end{pgfscope}%
\begin{pgfscope}%
\definecolor{textcolor}{rgb}{0.000000,0.000000,0.000000}%
\pgfsetstrokecolor{textcolor}%
\pgfsetfillcolor{textcolor}%
\pgftext[x=1.812346in,y=1.550000in,,base]{\color{textcolor}\rmfamily\fontsize{12.000000}{14.400000}\selectfont 25 kHz, 70\%}%
\end{pgfscope}%
\begin{pgfscope}%
\pgfsetbuttcap%
\pgfsetmiterjoin%
\definecolor{currentfill}{rgb}{1.000000,1.000000,1.000000}%
\pgfsetfillcolor{currentfill}%
\pgfsetlinewidth{0.000000pt}%
\definecolor{currentstroke}{rgb}{0.000000,0.000000,0.000000}%
\pgfsetstrokecolor{currentstroke}%
\pgfsetstrokeopacity{0.000000}%
\pgfsetdash{}{0pt}%
\pgfpathmoveto{\pgfqpoint{3.833026in}{0.565123in}}%
\pgfpathlineto{\pgfqpoint{5.941666in}{0.565123in}}%
\pgfpathlineto{\pgfqpoint{5.941666in}{1.466667in}}%
\pgfpathlineto{\pgfqpoint{3.833026in}{1.466667in}}%
\pgfpathlineto{\pgfqpoint{3.833026in}{0.565123in}}%
\pgfpathclose%
\pgfusepath{fill}%
\end{pgfscope}%
\begin{pgfscope}%
\pgfpathrectangle{\pgfqpoint{3.833026in}{0.565123in}}{\pgfqpoint{2.108640in}{0.901543in}}%
\pgfusepath{clip}%
\pgfsetbuttcap%
\pgfsetroundjoin%
\pgfsetlinewidth{0.501875pt}%
\definecolor{currentstroke}{rgb}{0.690196,0.690196,0.690196}%
\pgfsetstrokecolor{currentstroke}%
\pgfsetstrokeopacity{0.700000}%
\pgfsetdash{{1.850000pt}{0.800000pt}}{0.000000pt}%
\pgfpathmoveto{\pgfqpoint{3.833026in}{0.565123in}}%
\pgfpathlineto{\pgfqpoint{3.833026in}{1.466667in}}%
\pgfusepath{stroke}%
\end{pgfscope}%
\begin{pgfscope}%
\pgfsetbuttcap%
\pgfsetroundjoin%
\definecolor{currentfill}{rgb}{0.000000,0.000000,0.000000}%
\pgfsetfillcolor{currentfill}%
\pgfsetlinewidth{0.803000pt}%
\definecolor{currentstroke}{rgb}{0.000000,0.000000,0.000000}%
\pgfsetstrokecolor{currentstroke}%
\pgfsetdash{}{0pt}%
\pgfsys@defobject{currentmarker}{\pgfqpoint{0.000000in}{-0.048611in}}{\pgfqpoint{0.000000in}{0.000000in}}{%
\pgfpathmoveto{\pgfqpoint{0.000000in}{0.000000in}}%
\pgfpathlineto{\pgfqpoint{0.000000in}{-0.048611in}}%
\pgfusepath{stroke,fill}%
}%
\begin{pgfscope}%
\pgfsys@transformshift{3.833026in}{0.565123in}%
\pgfsys@useobject{currentmarker}{}%
\end{pgfscope}%
\end{pgfscope}%
\begin{pgfscope}%
\definecolor{textcolor}{rgb}{0.000000,0.000000,0.000000}%
\pgfsetstrokecolor{textcolor}%
\pgfsetfillcolor{textcolor}%
\pgftext[x=3.833026in,y=0.467901in,,top]{\color{textcolor}\rmfamily\fontsize{10.000000}{12.000000}\selectfont \(\displaystyle {400000}\)}%
\end{pgfscope}%
\begin{pgfscope}%
\pgfpathrectangle{\pgfqpoint{3.833026in}{0.565123in}}{\pgfqpoint{2.108640in}{0.901543in}}%
\pgfusepath{clip}%
\pgfsetbuttcap%
\pgfsetroundjoin%
\pgfsetlinewidth{0.501875pt}%
\definecolor{currentstroke}{rgb}{0.690196,0.690196,0.690196}%
\pgfsetstrokecolor{currentstroke}%
\pgfsetstrokeopacity{0.700000}%
\pgfsetdash{{1.850000pt}{0.800000pt}}{0.000000pt}%
\pgfpathmoveto{\pgfqpoint{4.360186in}{0.565123in}}%
\pgfpathlineto{\pgfqpoint{4.360186in}{1.466667in}}%
\pgfusepath{stroke}%
\end{pgfscope}%
\begin{pgfscope}%
\pgfsetbuttcap%
\pgfsetroundjoin%
\definecolor{currentfill}{rgb}{0.000000,0.000000,0.000000}%
\pgfsetfillcolor{currentfill}%
\pgfsetlinewidth{0.803000pt}%
\definecolor{currentstroke}{rgb}{0.000000,0.000000,0.000000}%
\pgfsetstrokecolor{currentstroke}%
\pgfsetdash{}{0pt}%
\pgfsys@defobject{currentmarker}{\pgfqpoint{0.000000in}{-0.048611in}}{\pgfqpoint{0.000000in}{0.000000in}}{%
\pgfpathmoveto{\pgfqpoint{0.000000in}{0.000000in}}%
\pgfpathlineto{\pgfqpoint{0.000000in}{-0.048611in}}%
\pgfusepath{stroke,fill}%
}%
\begin{pgfscope}%
\pgfsys@transformshift{4.360186in}{0.565123in}%
\pgfsys@useobject{currentmarker}{}%
\end{pgfscope}%
\end{pgfscope}%
\begin{pgfscope}%
\definecolor{textcolor}{rgb}{0.000000,0.000000,0.000000}%
\pgfsetstrokecolor{textcolor}%
\pgfsetfillcolor{textcolor}%
\pgftext[x=4.360186in,y=0.467901in,,top]{\color{textcolor}\rmfamily\fontsize{10.000000}{12.000000}\selectfont \(\displaystyle {450000}\)}%
\end{pgfscope}%
\begin{pgfscope}%
\pgfpathrectangle{\pgfqpoint{3.833026in}{0.565123in}}{\pgfqpoint{2.108640in}{0.901543in}}%
\pgfusepath{clip}%
\pgfsetbuttcap%
\pgfsetroundjoin%
\pgfsetlinewidth{0.501875pt}%
\definecolor{currentstroke}{rgb}{0.690196,0.690196,0.690196}%
\pgfsetstrokecolor{currentstroke}%
\pgfsetstrokeopacity{0.700000}%
\pgfsetdash{{1.850000pt}{0.800000pt}}{0.000000pt}%
\pgfpathmoveto{\pgfqpoint{4.887346in}{0.565123in}}%
\pgfpathlineto{\pgfqpoint{4.887346in}{1.466667in}}%
\pgfusepath{stroke}%
\end{pgfscope}%
\begin{pgfscope}%
\pgfsetbuttcap%
\pgfsetroundjoin%
\definecolor{currentfill}{rgb}{0.000000,0.000000,0.000000}%
\pgfsetfillcolor{currentfill}%
\pgfsetlinewidth{0.803000pt}%
\definecolor{currentstroke}{rgb}{0.000000,0.000000,0.000000}%
\pgfsetstrokecolor{currentstroke}%
\pgfsetdash{}{0pt}%
\pgfsys@defobject{currentmarker}{\pgfqpoint{0.000000in}{-0.048611in}}{\pgfqpoint{0.000000in}{0.000000in}}{%
\pgfpathmoveto{\pgfqpoint{0.000000in}{0.000000in}}%
\pgfpathlineto{\pgfqpoint{0.000000in}{-0.048611in}}%
\pgfusepath{stroke,fill}%
}%
\begin{pgfscope}%
\pgfsys@transformshift{4.887346in}{0.565123in}%
\pgfsys@useobject{currentmarker}{}%
\end{pgfscope}%
\end{pgfscope}%
\begin{pgfscope}%
\definecolor{textcolor}{rgb}{0.000000,0.000000,0.000000}%
\pgfsetstrokecolor{textcolor}%
\pgfsetfillcolor{textcolor}%
\pgftext[x=4.887346in,y=0.467901in,,top]{\color{textcolor}\rmfamily\fontsize{10.000000}{12.000000}\selectfont \(\displaystyle {500000}\)}%
\end{pgfscope}%
\begin{pgfscope}%
\pgfpathrectangle{\pgfqpoint{3.833026in}{0.565123in}}{\pgfqpoint{2.108640in}{0.901543in}}%
\pgfusepath{clip}%
\pgfsetbuttcap%
\pgfsetroundjoin%
\pgfsetlinewidth{0.501875pt}%
\definecolor{currentstroke}{rgb}{0.690196,0.690196,0.690196}%
\pgfsetstrokecolor{currentstroke}%
\pgfsetstrokeopacity{0.700000}%
\pgfsetdash{{1.850000pt}{0.800000pt}}{0.000000pt}%
\pgfpathmoveto{\pgfqpoint{5.414506in}{0.565123in}}%
\pgfpathlineto{\pgfqpoint{5.414506in}{1.466667in}}%
\pgfusepath{stroke}%
\end{pgfscope}%
\begin{pgfscope}%
\pgfsetbuttcap%
\pgfsetroundjoin%
\definecolor{currentfill}{rgb}{0.000000,0.000000,0.000000}%
\pgfsetfillcolor{currentfill}%
\pgfsetlinewidth{0.803000pt}%
\definecolor{currentstroke}{rgb}{0.000000,0.000000,0.000000}%
\pgfsetstrokecolor{currentstroke}%
\pgfsetdash{}{0pt}%
\pgfsys@defobject{currentmarker}{\pgfqpoint{0.000000in}{-0.048611in}}{\pgfqpoint{0.000000in}{0.000000in}}{%
\pgfpathmoveto{\pgfqpoint{0.000000in}{0.000000in}}%
\pgfpathlineto{\pgfqpoint{0.000000in}{-0.048611in}}%
\pgfusepath{stroke,fill}%
}%
\begin{pgfscope}%
\pgfsys@transformshift{5.414506in}{0.565123in}%
\pgfsys@useobject{currentmarker}{}%
\end{pgfscope}%
\end{pgfscope}%
\begin{pgfscope}%
\definecolor{textcolor}{rgb}{0.000000,0.000000,0.000000}%
\pgfsetstrokecolor{textcolor}%
\pgfsetfillcolor{textcolor}%
\pgftext[x=5.414506in,y=0.467901in,,top]{\color{textcolor}\rmfamily\fontsize{10.000000}{12.000000}\selectfont \(\displaystyle {550000}\)}%
\end{pgfscope}%
\begin{pgfscope}%
\pgfpathrectangle{\pgfqpoint{3.833026in}{0.565123in}}{\pgfqpoint{2.108640in}{0.901543in}}%
\pgfusepath{clip}%
\pgfsetbuttcap%
\pgfsetroundjoin%
\pgfsetlinewidth{0.501875pt}%
\definecolor{currentstroke}{rgb}{0.690196,0.690196,0.690196}%
\pgfsetstrokecolor{currentstroke}%
\pgfsetstrokeopacity{0.700000}%
\pgfsetdash{{1.850000pt}{0.800000pt}}{0.000000pt}%
\pgfpathmoveto{\pgfqpoint{5.941666in}{0.565123in}}%
\pgfpathlineto{\pgfqpoint{5.941666in}{1.466667in}}%
\pgfusepath{stroke}%
\end{pgfscope}%
\begin{pgfscope}%
\pgfsetbuttcap%
\pgfsetroundjoin%
\definecolor{currentfill}{rgb}{0.000000,0.000000,0.000000}%
\pgfsetfillcolor{currentfill}%
\pgfsetlinewidth{0.803000pt}%
\definecolor{currentstroke}{rgb}{0.000000,0.000000,0.000000}%
\pgfsetstrokecolor{currentstroke}%
\pgfsetdash{}{0pt}%
\pgfsys@defobject{currentmarker}{\pgfqpoint{0.000000in}{-0.048611in}}{\pgfqpoint{0.000000in}{0.000000in}}{%
\pgfpathmoveto{\pgfqpoint{0.000000in}{0.000000in}}%
\pgfpathlineto{\pgfqpoint{0.000000in}{-0.048611in}}%
\pgfusepath{stroke,fill}%
}%
\begin{pgfscope}%
\pgfsys@transformshift{5.941666in}{0.565123in}%
\pgfsys@useobject{currentmarker}{}%
\end{pgfscope}%
\end{pgfscope}%
\begin{pgfscope}%
\definecolor{textcolor}{rgb}{0.000000,0.000000,0.000000}%
\pgfsetstrokecolor{textcolor}%
\pgfsetfillcolor{textcolor}%
\pgftext[x=5.941666in,y=0.467901in,,top]{\color{textcolor}\rmfamily\fontsize{10.000000}{12.000000}\selectfont \(\displaystyle {600000}\)}%
\end{pgfscope}%
\begin{pgfscope}%
\definecolor{textcolor}{rgb}{0.000000,0.000000,0.000000}%
\pgfsetstrokecolor{textcolor}%
\pgfsetfillcolor{textcolor}%
\pgftext[x=4.887346in,y=0.288889in,,top]{\color{textcolor}\rmfamily\fontsize{10.000000}{12.000000}\selectfont Frequency (Hz)}%
\end{pgfscope}%
\begin{pgfscope}%
\pgfpathrectangle{\pgfqpoint{3.833026in}{0.565123in}}{\pgfqpoint{2.108640in}{0.901543in}}%
\pgfusepath{clip}%
\pgfsetbuttcap%
\pgfsetroundjoin%
\pgfsetlinewidth{0.501875pt}%
\definecolor{currentstroke}{rgb}{0.690196,0.690196,0.690196}%
\pgfsetstrokecolor{currentstroke}%
\pgfsetstrokeopacity{0.700000}%
\pgfsetdash{{1.850000pt}{0.800000pt}}{0.000000pt}%
\pgfpathmoveto{\pgfqpoint{3.833026in}{0.721338in}}%
\pgfpathlineto{\pgfqpoint{5.941666in}{0.721338in}}%
\pgfusepath{stroke}%
\end{pgfscope}%
\begin{pgfscope}%
\pgfsetbuttcap%
\pgfsetroundjoin%
\definecolor{currentfill}{rgb}{0.000000,0.000000,0.000000}%
\pgfsetfillcolor{currentfill}%
\pgfsetlinewidth{0.803000pt}%
\definecolor{currentstroke}{rgb}{0.000000,0.000000,0.000000}%
\pgfsetstrokecolor{currentstroke}%
\pgfsetdash{}{0pt}%
\pgfsys@defobject{currentmarker}{\pgfqpoint{-0.048611in}{0.000000in}}{\pgfqpoint{-0.000000in}{0.000000in}}{%
\pgfpathmoveto{\pgfqpoint{-0.000000in}{0.000000in}}%
\pgfpathlineto{\pgfqpoint{-0.048611in}{0.000000in}}%
\pgfusepath{stroke,fill}%
}%
\begin{pgfscope}%
\pgfsys@transformshift{3.833026in}{0.721338in}%
\pgfsys@useobject{currentmarker}{}%
\end{pgfscope}%
\end{pgfscope}%
\begin{pgfscope}%
\definecolor{textcolor}{rgb}{0.000000,0.000000,0.000000}%
\pgfsetstrokecolor{textcolor}%
\pgfsetfillcolor{textcolor}%
\pgftext[x=3.419444in, y=0.673113in, left, base]{\color{textcolor}\rmfamily\fontsize{10.000000}{12.000000}\selectfont \(\displaystyle {\ensuremath{-}100}\)}%
\end{pgfscope}%
\begin{pgfscope}%
\pgfpathrectangle{\pgfqpoint{3.833026in}{0.565123in}}{\pgfqpoint{2.108640in}{0.901543in}}%
\pgfusepath{clip}%
\pgfsetbuttcap%
\pgfsetroundjoin%
\pgfsetlinewidth{0.501875pt}%
\definecolor{currentstroke}{rgb}{0.690196,0.690196,0.690196}%
\pgfsetstrokecolor{currentstroke}%
\pgfsetstrokeopacity{0.700000}%
\pgfsetdash{{1.850000pt}{0.800000pt}}{0.000000pt}%
\pgfpathmoveto{\pgfqpoint{3.833026in}{1.202739in}}%
\pgfpathlineto{\pgfqpoint{5.941666in}{1.202739in}}%
\pgfusepath{stroke}%
\end{pgfscope}%
\begin{pgfscope}%
\pgfsetbuttcap%
\pgfsetroundjoin%
\definecolor{currentfill}{rgb}{0.000000,0.000000,0.000000}%
\pgfsetfillcolor{currentfill}%
\pgfsetlinewidth{0.803000pt}%
\definecolor{currentstroke}{rgb}{0.000000,0.000000,0.000000}%
\pgfsetstrokecolor{currentstroke}%
\pgfsetdash{}{0pt}%
\pgfsys@defobject{currentmarker}{\pgfqpoint{-0.048611in}{0.000000in}}{\pgfqpoint{-0.000000in}{0.000000in}}{%
\pgfpathmoveto{\pgfqpoint{-0.000000in}{0.000000in}}%
\pgfpathlineto{\pgfqpoint{-0.048611in}{0.000000in}}%
\pgfusepath{stroke,fill}%
}%
\begin{pgfscope}%
\pgfsys@transformshift{3.833026in}{1.202739in}%
\pgfsys@useobject{currentmarker}{}%
\end{pgfscope}%
\end{pgfscope}%
\begin{pgfscope}%
\definecolor{textcolor}{rgb}{0.000000,0.000000,0.000000}%
\pgfsetstrokecolor{textcolor}%
\pgfsetfillcolor{textcolor}%
\pgftext[x=3.488889in, y=1.154513in, left, base]{\color{textcolor}\rmfamily\fontsize{10.000000}{12.000000}\selectfont \(\displaystyle {\ensuremath{-}50}\)}%
\end{pgfscope}%
\begin{pgfscope}%
\definecolor{textcolor}{rgb}{0.000000,0.000000,0.000000}%
\pgfsetstrokecolor{textcolor}%
\pgfsetfillcolor{textcolor}%
\pgftext[x=3.363889in,y=1.015895in,,bottom,rotate=90.000000]{\color{textcolor}\rmfamily\fontsize{10.000000}{12.000000}\selectfont Amplitude (dB)}%
\end{pgfscope}%
\begin{pgfscope}%
\pgfpathrectangle{\pgfqpoint{3.833026in}{0.565123in}}{\pgfqpoint{2.108640in}{0.901543in}}%
\pgfusepath{clip}%
\pgfsetrectcap%
\pgfsetroundjoin%
\pgfsetlinewidth{1.003750pt}%
\definecolor{currentstroke}{rgb}{0.000000,0.000000,0.000000}%
\pgfsetstrokecolor{currentstroke}%
\pgfsetdash{}{0pt}%
\pgfpathmoveto{\pgfqpoint{3.829863in}{0.861847in}}%
\pgfpathlineto{\pgfqpoint{3.839352in}{0.866661in}}%
\pgfpathlineto{\pgfqpoint{3.848840in}{0.962941in}}%
\pgfpathlineto{\pgfqpoint{3.858329in}{0.895545in}}%
\pgfpathlineto{\pgfqpoint{3.867818in}{0.914801in}}%
\pgfpathlineto{\pgfqpoint{3.877307in}{0.914500in}}%
\pgfpathlineto{\pgfqpoint{3.886796in}{0.818220in}}%
\pgfpathlineto{\pgfqpoint{3.896285in}{0.851918in}}%
\pgfpathlineto{\pgfqpoint{3.905774in}{0.895244in}}%
\pgfpathlineto{\pgfqpoint{3.915263in}{0.880802in}}%
\pgfpathlineto{\pgfqpoint{3.924752in}{0.837476in}}%
\pgfpathlineto{\pgfqpoint{3.934240in}{0.904872in}}%
\pgfpathlineto{\pgfqpoint{3.943729in}{0.837476in}}%
\pgfpathlineto{\pgfqpoint{3.953218in}{0.914199in}}%
\pgfpathlineto{\pgfqpoint{3.962707in}{0.846803in}}%
\pgfpathlineto{\pgfqpoint{3.972196in}{0.846803in}}%
\pgfpathlineto{\pgfqpoint{3.981685in}{0.856431in}}%
\pgfpathlineto{\pgfqpoint{3.991174in}{0.914199in}}%
\pgfpathlineto{\pgfqpoint{4.000663in}{0.938269in}}%
\pgfpathlineto{\pgfqpoint{4.010151in}{0.861245in}}%
\pgfpathlineto{\pgfqpoint{4.019640in}{0.866059in}}%
\pgfpathlineto{\pgfqpoint{4.029129in}{0.865758in}}%
\pgfpathlineto{\pgfqpoint{4.038618in}{0.889828in}}%
\pgfpathlineto{\pgfqpoint{4.048107in}{0.937968in}}%
\pgfpathlineto{\pgfqpoint{4.057596in}{0.870572in}}%
\pgfpathlineto{\pgfqpoint{4.067085in}{0.860944in}}%
\pgfpathlineto{\pgfqpoint{4.076574in}{0.894642in}}%
\pgfpathlineto{\pgfqpoint{4.086063in}{0.918712in}}%
\pgfpathlineto{\pgfqpoint{4.095551in}{0.918712in}}%
\pgfpathlineto{\pgfqpoint{4.105040in}{0.884713in}}%
\pgfpathlineto{\pgfqpoint{4.114529in}{0.932853in}}%
\pgfpathlineto{\pgfqpoint{4.124018in}{0.865457in}}%
\pgfpathlineto{\pgfqpoint{4.133507in}{0.894341in}}%
\pgfpathlineto{\pgfqpoint{4.142996in}{0.908783in}}%
\pgfpathlineto{\pgfqpoint{4.152485in}{0.918411in}}%
\pgfpathlineto{\pgfqpoint{4.161974in}{0.879899in}}%
\pgfpathlineto{\pgfqpoint{4.171462in}{0.860643in}}%
\pgfpathlineto{\pgfqpoint{4.180951in}{0.855528in}}%
\pgfpathlineto{\pgfqpoint{4.190440in}{0.865156in}}%
\pgfpathlineto{\pgfqpoint{4.199929in}{0.894040in}}%
\pgfpathlineto{\pgfqpoint{4.209418in}{0.898854in}}%
\pgfpathlineto{\pgfqpoint{4.218907in}{0.913296in}}%
\pgfpathlineto{\pgfqpoint{4.228396in}{0.860342in}}%
\pgfpathlineto{\pgfqpoint{4.237885in}{0.879598in}}%
\pgfpathlineto{\pgfqpoint{4.247374in}{0.860342in}}%
\pgfpathlineto{\pgfqpoint{4.256862in}{0.835971in}}%
\pgfpathlineto{\pgfqpoint{4.266351in}{0.884112in}}%
\pgfpathlineto{\pgfqpoint{4.275840in}{0.903368in}}%
\pgfpathlineto{\pgfqpoint{4.285329in}{0.941880in}}%
\pgfpathlineto{\pgfqpoint{4.294818in}{0.937066in}}%
\pgfpathlineto{\pgfqpoint{4.304307in}{0.874484in}}%
\pgfpathlineto{\pgfqpoint{4.313796in}{0.864856in}}%
\pgfpathlineto{\pgfqpoint{4.323285in}{0.830857in}}%
\pgfpathlineto{\pgfqpoint{4.332773in}{0.850113in}}%
\pgfpathlineto{\pgfqpoint{4.342262in}{0.888625in}}%
\pgfpathlineto{\pgfqpoint{4.351751in}{0.864555in}}%
\pgfpathlineto{\pgfqpoint{4.361240in}{0.854927in}}%
\pgfpathlineto{\pgfqpoint{4.370729in}{0.864555in}}%
\pgfpathlineto{\pgfqpoint{4.380218in}{0.878997in}}%
\pgfpathlineto{\pgfqpoint{4.399196in}{0.830556in}}%
\pgfpathlineto{\pgfqpoint{4.408685in}{0.878696in}}%
\pgfpathlineto{\pgfqpoint{4.418173in}{0.864254in}}%
\pgfpathlineto{\pgfqpoint{4.427662in}{0.893138in}}%
\pgfpathlineto{\pgfqpoint{4.437151in}{0.869068in}}%
\pgfpathlineto{\pgfqpoint{4.446640in}{0.835370in}}%
\pgfpathlineto{\pgfqpoint{4.456129in}{0.830556in}}%
\pgfpathlineto{\pgfqpoint{4.465618in}{0.787230in}}%
\pgfpathlineto{\pgfqpoint{4.475107in}{0.873581in}}%
\pgfpathlineto{\pgfqpoint{4.494084in}{0.835069in}}%
\pgfpathlineto{\pgfqpoint{4.503573in}{0.873581in}}%
\pgfpathlineto{\pgfqpoint{4.513062in}{0.849511in}}%
\pgfpathlineto{\pgfqpoint{4.522551in}{0.863953in}}%
\pgfpathlineto{\pgfqpoint{4.532040in}{0.902465in}}%
\pgfpathlineto{\pgfqpoint{4.541529in}{0.868767in}}%
\pgfpathlineto{\pgfqpoint{4.551018in}{0.916606in}}%
\pgfpathlineto{\pgfqpoint{4.560507in}{0.921420in}}%
\pgfpathlineto{\pgfqpoint{4.569995in}{0.849210in}}%
\pgfpathlineto{\pgfqpoint{4.579484in}{0.935862in}}%
\pgfpathlineto{\pgfqpoint{4.588973in}{1.065840in}}%
\pgfpathlineto{\pgfqpoint{4.598462in}{1.147678in}}%
\pgfpathlineto{\pgfqpoint{4.607951in}{1.195819in}}%
\pgfpathlineto{\pgfqpoint{4.617440in}{1.224703in}}%
\pgfpathlineto{\pgfqpoint{4.626929in}{1.238844in}}%
\pgfpathlineto{\pgfqpoint{4.636418in}{1.234030in}}%
\pgfpathlineto{\pgfqpoint{4.645907in}{1.205146in}}%
\pgfpathlineto{\pgfqpoint{4.655395in}{1.157006in}}%
\pgfpathlineto{\pgfqpoint{4.664884in}{1.089609in}}%
\pgfpathlineto{\pgfqpoint{4.674373in}{0.974073in}}%
\pgfpathlineto{\pgfqpoint{4.683862in}{0.844095in}}%
\pgfpathlineto{\pgfqpoint{4.693351in}{0.892235in}}%
\pgfpathlineto{\pgfqpoint{4.702840in}{0.863050in}}%
\pgfpathlineto{\pgfqpoint{4.712329in}{0.877492in}}%
\pgfpathlineto{\pgfqpoint{4.721818in}{0.887120in}}%
\pgfpathlineto{\pgfqpoint{4.731306in}{0.848608in}}%
\pgfpathlineto{\pgfqpoint{4.740795in}{0.834166in}}%
\pgfpathlineto{\pgfqpoint{4.750284in}{0.838980in}}%
\pgfpathlineto{\pgfqpoint{4.759773in}{0.911190in}}%
\pgfpathlineto{\pgfqpoint{4.769262in}{0.872678in}}%
\pgfpathlineto{\pgfqpoint{4.778751in}{0.862749in}}%
\pgfpathlineto{\pgfqpoint{4.788240in}{0.848307in}}%
\pgfpathlineto{\pgfqpoint{4.797729in}{0.886819in}}%
\pgfpathlineto{\pgfqpoint{4.807218in}{0.853121in}}%
\pgfpathlineto{\pgfqpoint{4.816706in}{0.867563in}}%
\pgfpathlineto{\pgfqpoint{4.826195in}{0.853121in}}%
\pgfpathlineto{\pgfqpoint{4.835684in}{0.882005in}}%
\pgfpathlineto{\pgfqpoint{4.845173in}{1.156404in}}%
\pgfpathlineto{\pgfqpoint{4.854662in}{1.276453in}}%
\pgfpathlineto{\pgfqpoint{4.864151in}{1.353477in}}%
\pgfpathlineto{\pgfqpoint{4.873640in}{1.396803in}}%
\pgfpathlineto{\pgfqpoint{4.883129in}{1.420873in}}%
\pgfpathlineto{\pgfqpoint{4.892617in}{1.425687in}}%
\pgfpathlineto{\pgfqpoint{4.902106in}{1.420873in}}%
\pgfpathlineto{\pgfqpoint{4.911595in}{1.391989in}}%
\pgfpathlineto{\pgfqpoint{4.921084in}{1.339035in}}%
\pgfpathlineto{\pgfqpoint{4.930573in}{1.261710in}}%
\pgfpathlineto{\pgfqpoint{4.940062in}{1.131732in}}%
\pgfpathlineto{\pgfqpoint{4.949551in}{0.852520in}}%
\pgfpathlineto{\pgfqpoint{4.959040in}{0.891032in}}%
\pgfpathlineto{\pgfqpoint{4.968529in}{0.857334in}}%
\pgfpathlineto{\pgfqpoint{4.978017in}{0.818822in}}%
\pgfpathlineto{\pgfqpoint{4.987506in}{0.891032in}}%
\pgfpathlineto{\pgfqpoint{4.996995in}{0.770681in}}%
\pgfpathlineto{\pgfqpoint{5.006484in}{0.857033in}}%
\pgfpathlineto{\pgfqpoint{5.015973in}{0.900359in}}%
\pgfpathlineto{\pgfqpoint{5.025462in}{0.832963in}}%
\pgfpathlineto{\pgfqpoint{5.034951in}{0.890731in}}%
\pgfpathlineto{\pgfqpoint{5.044440in}{0.890731in}}%
\pgfpathlineto{\pgfqpoint{5.053928in}{0.832963in}}%
\pgfpathlineto{\pgfqpoint{5.063417in}{0.818521in}}%
\pgfpathlineto{\pgfqpoint{5.072906in}{0.857033in}}%
\pgfpathlineto{\pgfqpoint{5.082395in}{0.847104in}}%
\pgfpathlineto{\pgfqpoint{5.091884in}{0.866360in}}%
\pgfpathlineto{\pgfqpoint{5.101373in}{0.856732in}}%
\pgfpathlineto{\pgfqpoint{5.110862in}{0.996338in}}%
\pgfpathlineto{\pgfqpoint{5.120351in}{1.097432in}}%
\pgfpathlineto{\pgfqpoint{5.129840in}{1.169642in}}%
\pgfpathlineto{\pgfqpoint{5.139328in}{1.208154in}}%
\pgfpathlineto{\pgfqpoint{5.148817in}{1.232224in}}%
\pgfpathlineto{\pgfqpoint{5.158306in}{1.236738in}}%
\pgfpathlineto{\pgfqpoint{5.167795in}{1.222296in}}%
\pgfpathlineto{\pgfqpoint{5.177284in}{1.188598in}}%
\pgfpathlineto{\pgfqpoint{5.186773in}{1.135643in}}%
\pgfpathlineto{\pgfqpoint{5.196262in}{1.044177in}}%
\pgfpathlineto{\pgfqpoint{5.205751in}{0.919013in}}%
\pgfpathlineto{\pgfqpoint{5.215239in}{0.943083in}}%
\pgfpathlineto{\pgfqpoint{5.224728in}{0.769779in}}%
\pgfpathlineto{\pgfqpoint{5.234217in}{0.832060in}}%
\pgfpathlineto{\pgfqpoint{5.243706in}{0.841688in}}%
\pgfpathlineto{\pgfqpoint{5.253195in}{0.827246in}}%
\pgfpathlineto{\pgfqpoint{5.262684in}{0.817618in}}%
\pgfpathlineto{\pgfqpoint{5.272173in}{0.865758in}}%
\pgfpathlineto{\pgfqpoint{5.281662in}{0.846502in}}%
\pgfpathlineto{\pgfqpoint{5.291150in}{0.783920in}}%
\pgfpathlineto{\pgfqpoint{5.300639in}{0.885014in}}%
\pgfpathlineto{\pgfqpoint{5.310128in}{0.802875in}}%
\pgfpathlineto{\pgfqpoint{5.319617in}{0.788433in}}%
\pgfpathlineto{\pgfqpoint{5.329106in}{0.812503in}}%
\pgfpathlineto{\pgfqpoint{5.338595in}{0.802875in}}%
\pgfpathlineto{\pgfqpoint{5.348084in}{0.812503in}}%
\pgfpathlineto{\pgfqpoint{5.357573in}{0.812503in}}%
\pgfpathlineto{\pgfqpoint{5.367062in}{0.831759in}}%
\pgfpathlineto{\pgfqpoint{5.376550in}{0.822131in}}%
\pgfpathlineto{\pgfqpoint{5.386039in}{0.802574in}}%
\pgfpathlineto{\pgfqpoint{5.405017in}{0.831458in}}%
\pgfpathlineto{\pgfqpoint{5.414506in}{0.850714in}}%
\pgfpathlineto{\pgfqpoint{5.423995in}{0.754434in}}%
\pgfpathlineto{\pgfqpoint{5.433484in}{0.797760in}}%
\pgfpathlineto{\pgfqpoint{5.442973in}{0.831458in}}%
\pgfpathlineto{\pgfqpoint{5.452461in}{0.817016in}}%
\pgfpathlineto{\pgfqpoint{5.461950in}{0.811901in}}%
\pgfpathlineto{\pgfqpoint{5.471439in}{0.811901in}}%
\pgfpathlineto{\pgfqpoint{5.480928in}{0.860041in}}%
\pgfpathlineto{\pgfqpoint{5.490417in}{0.754133in}}%
\pgfpathlineto{\pgfqpoint{5.499906in}{0.850413in}}%
\pgfpathlineto{\pgfqpoint{5.509395in}{0.855227in}}%
\pgfpathlineto{\pgfqpoint{5.518884in}{0.811901in}}%
\pgfpathlineto{\pgfqpoint{5.528373in}{0.845599in}}%
\pgfpathlineto{\pgfqpoint{5.537861in}{0.826043in}}%
\pgfpathlineto{\pgfqpoint{5.547350in}{0.864555in}}%
\pgfpathlineto{\pgfqpoint{5.556839in}{0.840485in}}%
\pgfpathlineto{\pgfqpoint{5.566328in}{0.835671in}}%
\pgfpathlineto{\pgfqpoint{5.575817in}{0.816415in}}%
\pgfpathlineto{\pgfqpoint{5.585306in}{0.753832in}}%
\pgfpathlineto{\pgfqpoint{5.594795in}{0.835671in}}%
\pgfpathlineto{\pgfqpoint{5.604284in}{0.768274in}}%
\pgfpathlineto{\pgfqpoint{5.613772in}{0.820928in}}%
\pgfpathlineto{\pgfqpoint{5.623261in}{0.767974in}}%
\pgfpathlineto{\pgfqpoint{5.632750in}{0.782416in}}%
\pgfpathlineto{\pgfqpoint{5.642239in}{0.753532in}}%
\pgfpathlineto{\pgfqpoint{5.651728in}{0.801672in}}%
\pgfpathlineto{\pgfqpoint{5.661217in}{0.825742in}}%
\pgfpathlineto{\pgfqpoint{5.670706in}{0.782416in}}%
\pgfpathlineto{\pgfqpoint{5.680195in}{0.825742in}}%
\pgfpathlineto{\pgfqpoint{5.689684in}{0.796557in}}%
\pgfpathlineto{\pgfqpoint{5.699172in}{0.782115in}}%
\pgfpathlineto{\pgfqpoint{5.708661in}{0.782115in}}%
\pgfpathlineto{\pgfqpoint{5.718150in}{0.772487in}}%
\pgfpathlineto{\pgfqpoint{5.727639in}{0.786929in}}%
\pgfpathlineto{\pgfqpoint{5.737128in}{0.810999in}}%
\pgfpathlineto{\pgfqpoint{5.746617in}{0.796557in}}%
\pgfpathlineto{\pgfqpoint{5.756106in}{0.767673in}}%
\pgfpathlineto{\pgfqpoint{5.765595in}{0.786628in}}%
\pgfpathlineto{\pgfqpoint{5.775083in}{0.752930in}}%
\pgfpathlineto{\pgfqpoint{5.784572in}{0.781814in}}%
\pgfpathlineto{\pgfqpoint{5.794061in}{0.752930in}}%
\pgfpathlineto{\pgfqpoint{5.803550in}{0.815512in}}%
\pgfpathlineto{\pgfqpoint{5.813039in}{0.810698in}}%
\pgfpathlineto{\pgfqpoint{5.822528in}{0.801070in}}%
\pgfpathlineto{\pgfqpoint{5.832017in}{0.801070in}}%
\pgfpathlineto{\pgfqpoint{5.841506in}{0.863351in}}%
\pgfpathlineto{\pgfqpoint{5.850995in}{0.795955in}}%
\pgfpathlineto{\pgfqpoint{5.860483in}{0.815211in}}%
\pgfpathlineto{\pgfqpoint{5.869972in}{0.767071in}}%
\pgfpathlineto{\pgfqpoint{5.879461in}{0.767071in}}%
\pgfpathlineto{\pgfqpoint{5.888950in}{0.752629in}}%
\pgfpathlineto{\pgfqpoint{5.898439in}{0.781513in}}%
\pgfpathlineto{\pgfqpoint{5.907928in}{0.815211in}}%
\pgfpathlineto{\pgfqpoint{5.917417in}{0.781212in}}%
\pgfpathlineto{\pgfqpoint{5.926906in}{0.810096in}}%
\pgfpathlineto{\pgfqpoint{5.936394in}{0.752328in}}%
\pgfpathlineto{\pgfqpoint{5.945883in}{0.766770in}}%
\pgfpathlineto{\pgfqpoint{5.945883in}{0.766770in}}%
\pgfusepath{stroke}%
\end{pgfscope}%
\begin{pgfscope}%
\pgfsetrectcap%
\pgfsetmiterjoin%
\pgfsetlinewidth{0.803000pt}%
\definecolor{currentstroke}{rgb}{0.000000,0.000000,0.000000}%
\pgfsetstrokecolor{currentstroke}%
\pgfsetdash{}{0pt}%
\pgfpathmoveto{\pgfqpoint{3.833026in}{0.565123in}}%
\pgfpathlineto{\pgfqpoint{3.833026in}{1.466667in}}%
\pgfusepath{stroke}%
\end{pgfscope}%
\begin{pgfscope}%
\pgfsetrectcap%
\pgfsetmiterjoin%
\pgfsetlinewidth{0.803000pt}%
\definecolor{currentstroke}{rgb}{0.000000,0.000000,0.000000}%
\pgfsetstrokecolor{currentstroke}%
\pgfsetdash{}{0pt}%
\pgfpathmoveto{\pgfqpoint{5.941666in}{0.565123in}}%
\pgfpathlineto{\pgfqpoint{5.941666in}{1.466667in}}%
\pgfusepath{stroke}%
\end{pgfscope}%
\begin{pgfscope}%
\pgfsetrectcap%
\pgfsetmiterjoin%
\pgfsetlinewidth{0.803000pt}%
\definecolor{currentstroke}{rgb}{0.000000,0.000000,0.000000}%
\pgfsetstrokecolor{currentstroke}%
\pgfsetdash{}{0pt}%
\pgfpathmoveto{\pgfqpoint{3.833026in}{0.565123in}}%
\pgfpathlineto{\pgfqpoint{5.941666in}{0.565123in}}%
\pgfusepath{stroke}%
\end{pgfscope}%
\begin{pgfscope}%
\pgfsetrectcap%
\pgfsetmiterjoin%
\pgfsetlinewidth{0.803000pt}%
\definecolor{currentstroke}{rgb}{0.000000,0.000000,0.000000}%
\pgfsetstrokecolor{currentstroke}%
\pgfsetdash{}{0pt}%
\pgfpathmoveto{\pgfqpoint{3.833026in}{1.466667in}}%
\pgfpathlineto{\pgfqpoint{5.941666in}{1.466667in}}%
\pgfusepath{stroke}%
\end{pgfscope}%
\begin{pgfscope}%
\definecolor{textcolor}{rgb}{0.000000,0.000000,0.000000}%
\pgfsetstrokecolor{textcolor}%
\pgfsetfillcolor{textcolor}%
\pgftext[x=4.887346in,y=1.550000in,,base]{\color{textcolor}\rmfamily\fontsize{12.000000}{14.400000}\selectfont 25 kHz, 20\%}%
\end{pgfscope}%
\end{pgfpicture}%
\makeatother%
\endgroup%

		\caption{Measured frequency content with a spectrum analyzer as a function of frequency of sine waves with different modulation frequencies and different modulation indexes.}
		\label{fig:444}
	\end{figure}
	
	\subsubsection{}
	
	\begin{figure}[H]
		\centering
		%% Creator: Matplotlib, PGF backend
%%
%% To include the figure in your LaTeX document, write
%%   \input{<filename>.pgf}
%%
%% Make sure the required packages are loaded in your preamble
%%   \usepackage{pgf}
%%
%% Also ensure that all the required font packages are loaded; for instance,
%% the lmodern package is sometimes necessary when using math font.
%%   \usepackage{lmodern}
%%
%% Figures using additional raster images can only be included by \input if
%% they are in the same directory as the main LaTeX file. For loading figures
%% from other directories you can use the `import` package
%%   \usepackage{import}
%%
%% and then include the figures with
%%   \import{<path to file>}{<filename>.pgf}
%%
%% Matplotlib used the following preamble
%%
\begingroup%
\makeatletter%
\begin{pgfpicture}%
\pgfpathrectangle{\pgfpointorigin}{\pgfqpoint{6.300000in}{3.000000in}}%
\pgfusepath{use as bounding box, clip}%
\begin{pgfscope}%
\pgfsetbuttcap%
\pgfsetmiterjoin%
\definecolor{currentfill}{rgb}{1.000000,1.000000,1.000000}%
\pgfsetfillcolor{currentfill}%
\pgfsetlinewidth{0.000000pt}%
\definecolor{currentstroke}{rgb}{1.000000,1.000000,1.000000}%
\pgfsetstrokecolor{currentstroke}%
\pgfsetdash{}{0pt}%
\pgfpathmoveto{\pgfqpoint{0.000000in}{0.000000in}}%
\pgfpathlineto{\pgfqpoint{6.300000in}{0.000000in}}%
\pgfpathlineto{\pgfqpoint{6.300000in}{3.000000in}}%
\pgfpathlineto{\pgfqpoint{0.000000in}{3.000000in}}%
\pgfpathlineto{\pgfqpoint{0.000000in}{0.000000in}}%
\pgfpathclose%
\pgfusepath{fill}%
\end{pgfscope}%
\begin{pgfscope}%
\pgfsetbuttcap%
\pgfsetmiterjoin%
\definecolor{currentfill}{rgb}{1.000000,1.000000,1.000000}%
\pgfsetfillcolor{currentfill}%
\pgfsetlinewidth{0.000000pt}%
\definecolor{currentstroke}{rgb}{0.000000,0.000000,0.000000}%
\pgfsetstrokecolor{currentstroke}%
\pgfsetstrokeopacity{0.000000}%
\pgfsetdash{}{0pt}%
\pgfpathmoveto{\pgfqpoint{0.758026in}{0.565123in}}%
\pgfpathlineto{\pgfqpoint{6.061265in}{0.565123in}}%
\pgfpathlineto{\pgfqpoint{6.061265in}{2.850000in}}%
\pgfpathlineto{\pgfqpoint{0.758026in}{2.850000in}}%
\pgfpathlineto{\pgfqpoint{0.758026in}{0.565123in}}%
\pgfpathclose%
\pgfusepath{fill}%
\end{pgfscope}%
\begin{pgfscope}%
\pgfpathrectangle{\pgfqpoint{0.758026in}{0.565123in}}{\pgfqpoint{5.303239in}{2.284877in}}%
\pgfusepath{clip}%
\pgfsetbuttcap%
\pgfsetroundjoin%
\pgfsetlinewidth{0.501875pt}%
\definecolor{currentstroke}{rgb}{0.690196,0.690196,0.690196}%
\pgfsetstrokecolor{currentstroke}%
\pgfsetstrokeopacity{0.700000}%
\pgfsetdash{{1.850000pt}{0.800000pt}}{0.000000pt}%
\pgfpathmoveto{\pgfqpoint{1.489507in}{0.565123in}}%
\pgfpathlineto{\pgfqpoint{1.489507in}{2.850000in}}%
\pgfusepath{stroke}%
\end{pgfscope}%
\begin{pgfscope}%
\pgfsetbuttcap%
\pgfsetroundjoin%
\definecolor{currentfill}{rgb}{0.000000,0.000000,0.000000}%
\pgfsetfillcolor{currentfill}%
\pgfsetlinewidth{0.803000pt}%
\definecolor{currentstroke}{rgb}{0.000000,0.000000,0.000000}%
\pgfsetstrokecolor{currentstroke}%
\pgfsetdash{}{0pt}%
\pgfsys@defobject{currentmarker}{\pgfqpoint{0.000000in}{-0.048611in}}{\pgfqpoint{0.000000in}{0.000000in}}{%
\pgfpathmoveto{\pgfqpoint{0.000000in}{0.000000in}}%
\pgfpathlineto{\pgfqpoint{0.000000in}{-0.048611in}}%
\pgfusepath{stroke,fill}%
}%
\begin{pgfscope}%
\pgfsys@transformshift{1.489507in}{0.565123in}%
\pgfsys@useobject{currentmarker}{}%
\end{pgfscope}%
\end{pgfscope}%
\begin{pgfscope}%
\definecolor{textcolor}{rgb}{0.000000,0.000000,0.000000}%
\pgfsetstrokecolor{textcolor}%
\pgfsetfillcolor{textcolor}%
\pgftext[x=1.489507in,y=0.467901in,,top]{\color{textcolor}\rmfamily\fontsize{10.000000}{12.000000}\selectfont \(\displaystyle {0.5}\)}%
\end{pgfscope}%
\begin{pgfscope}%
\pgfpathrectangle{\pgfqpoint{0.758026in}{0.565123in}}{\pgfqpoint{5.303239in}{2.284877in}}%
\pgfusepath{clip}%
\pgfsetbuttcap%
\pgfsetroundjoin%
\pgfsetlinewidth{0.501875pt}%
\definecolor{currentstroke}{rgb}{0.690196,0.690196,0.690196}%
\pgfsetstrokecolor{currentstroke}%
\pgfsetstrokeopacity{0.700000}%
\pgfsetdash{{1.850000pt}{0.800000pt}}{0.000000pt}%
\pgfpathmoveto{\pgfqpoint{2.403859in}{0.565123in}}%
\pgfpathlineto{\pgfqpoint{2.403859in}{2.850000in}}%
\pgfusepath{stroke}%
\end{pgfscope}%
\begin{pgfscope}%
\pgfsetbuttcap%
\pgfsetroundjoin%
\definecolor{currentfill}{rgb}{0.000000,0.000000,0.000000}%
\pgfsetfillcolor{currentfill}%
\pgfsetlinewidth{0.803000pt}%
\definecolor{currentstroke}{rgb}{0.000000,0.000000,0.000000}%
\pgfsetstrokecolor{currentstroke}%
\pgfsetdash{}{0pt}%
\pgfsys@defobject{currentmarker}{\pgfqpoint{0.000000in}{-0.048611in}}{\pgfqpoint{0.000000in}{0.000000in}}{%
\pgfpathmoveto{\pgfqpoint{0.000000in}{0.000000in}}%
\pgfpathlineto{\pgfqpoint{0.000000in}{-0.048611in}}%
\pgfusepath{stroke,fill}%
}%
\begin{pgfscope}%
\pgfsys@transformshift{2.403859in}{0.565123in}%
\pgfsys@useobject{currentmarker}{}%
\end{pgfscope}%
\end{pgfscope}%
\begin{pgfscope}%
\definecolor{textcolor}{rgb}{0.000000,0.000000,0.000000}%
\pgfsetstrokecolor{textcolor}%
\pgfsetfillcolor{textcolor}%
\pgftext[x=2.403859in,y=0.467901in,,top]{\color{textcolor}\rmfamily\fontsize{10.000000}{12.000000}\selectfont \(\displaystyle {1.0}\)}%
\end{pgfscope}%
\begin{pgfscope}%
\pgfpathrectangle{\pgfqpoint{0.758026in}{0.565123in}}{\pgfqpoint{5.303239in}{2.284877in}}%
\pgfusepath{clip}%
\pgfsetbuttcap%
\pgfsetroundjoin%
\pgfsetlinewidth{0.501875pt}%
\definecolor{currentstroke}{rgb}{0.690196,0.690196,0.690196}%
\pgfsetstrokecolor{currentstroke}%
\pgfsetstrokeopacity{0.700000}%
\pgfsetdash{{1.850000pt}{0.800000pt}}{0.000000pt}%
\pgfpathmoveto{\pgfqpoint{3.318210in}{0.565123in}}%
\pgfpathlineto{\pgfqpoint{3.318210in}{2.850000in}}%
\pgfusepath{stroke}%
\end{pgfscope}%
\begin{pgfscope}%
\pgfsetbuttcap%
\pgfsetroundjoin%
\definecolor{currentfill}{rgb}{0.000000,0.000000,0.000000}%
\pgfsetfillcolor{currentfill}%
\pgfsetlinewidth{0.803000pt}%
\definecolor{currentstroke}{rgb}{0.000000,0.000000,0.000000}%
\pgfsetstrokecolor{currentstroke}%
\pgfsetdash{}{0pt}%
\pgfsys@defobject{currentmarker}{\pgfqpoint{0.000000in}{-0.048611in}}{\pgfqpoint{0.000000in}{0.000000in}}{%
\pgfpathmoveto{\pgfqpoint{0.000000in}{0.000000in}}%
\pgfpathlineto{\pgfqpoint{0.000000in}{-0.048611in}}%
\pgfusepath{stroke,fill}%
}%
\begin{pgfscope}%
\pgfsys@transformshift{3.318210in}{0.565123in}%
\pgfsys@useobject{currentmarker}{}%
\end{pgfscope}%
\end{pgfscope}%
\begin{pgfscope}%
\definecolor{textcolor}{rgb}{0.000000,0.000000,0.000000}%
\pgfsetstrokecolor{textcolor}%
\pgfsetfillcolor{textcolor}%
\pgftext[x=3.318210in,y=0.467901in,,top]{\color{textcolor}\rmfamily\fontsize{10.000000}{12.000000}\selectfont \(\displaystyle {1.5}\)}%
\end{pgfscope}%
\begin{pgfscope}%
\pgfpathrectangle{\pgfqpoint{0.758026in}{0.565123in}}{\pgfqpoint{5.303239in}{2.284877in}}%
\pgfusepath{clip}%
\pgfsetbuttcap%
\pgfsetroundjoin%
\pgfsetlinewidth{0.501875pt}%
\definecolor{currentstroke}{rgb}{0.690196,0.690196,0.690196}%
\pgfsetstrokecolor{currentstroke}%
\pgfsetstrokeopacity{0.700000}%
\pgfsetdash{{1.850000pt}{0.800000pt}}{0.000000pt}%
\pgfpathmoveto{\pgfqpoint{4.232562in}{0.565123in}}%
\pgfpathlineto{\pgfqpoint{4.232562in}{2.850000in}}%
\pgfusepath{stroke}%
\end{pgfscope}%
\begin{pgfscope}%
\pgfsetbuttcap%
\pgfsetroundjoin%
\definecolor{currentfill}{rgb}{0.000000,0.000000,0.000000}%
\pgfsetfillcolor{currentfill}%
\pgfsetlinewidth{0.803000pt}%
\definecolor{currentstroke}{rgb}{0.000000,0.000000,0.000000}%
\pgfsetstrokecolor{currentstroke}%
\pgfsetdash{}{0pt}%
\pgfsys@defobject{currentmarker}{\pgfqpoint{0.000000in}{-0.048611in}}{\pgfqpoint{0.000000in}{0.000000in}}{%
\pgfpathmoveto{\pgfqpoint{0.000000in}{0.000000in}}%
\pgfpathlineto{\pgfqpoint{0.000000in}{-0.048611in}}%
\pgfusepath{stroke,fill}%
}%
\begin{pgfscope}%
\pgfsys@transformshift{4.232562in}{0.565123in}%
\pgfsys@useobject{currentmarker}{}%
\end{pgfscope}%
\end{pgfscope}%
\begin{pgfscope}%
\definecolor{textcolor}{rgb}{0.000000,0.000000,0.000000}%
\pgfsetstrokecolor{textcolor}%
\pgfsetfillcolor{textcolor}%
\pgftext[x=4.232562in,y=0.467901in,,top]{\color{textcolor}\rmfamily\fontsize{10.000000}{12.000000}\selectfont \(\displaystyle {2.0}\)}%
\end{pgfscope}%
\begin{pgfscope}%
\pgfpathrectangle{\pgfqpoint{0.758026in}{0.565123in}}{\pgfqpoint{5.303239in}{2.284877in}}%
\pgfusepath{clip}%
\pgfsetbuttcap%
\pgfsetroundjoin%
\pgfsetlinewidth{0.501875pt}%
\definecolor{currentstroke}{rgb}{0.690196,0.690196,0.690196}%
\pgfsetstrokecolor{currentstroke}%
\pgfsetstrokeopacity{0.700000}%
\pgfsetdash{{1.850000pt}{0.800000pt}}{0.000000pt}%
\pgfpathmoveto{\pgfqpoint{5.146914in}{0.565123in}}%
\pgfpathlineto{\pgfqpoint{5.146914in}{2.850000in}}%
\pgfusepath{stroke}%
\end{pgfscope}%
\begin{pgfscope}%
\pgfsetbuttcap%
\pgfsetroundjoin%
\definecolor{currentfill}{rgb}{0.000000,0.000000,0.000000}%
\pgfsetfillcolor{currentfill}%
\pgfsetlinewidth{0.803000pt}%
\definecolor{currentstroke}{rgb}{0.000000,0.000000,0.000000}%
\pgfsetstrokecolor{currentstroke}%
\pgfsetdash{}{0pt}%
\pgfsys@defobject{currentmarker}{\pgfqpoint{0.000000in}{-0.048611in}}{\pgfqpoint{0.000000in}{0.000000in}}{%
\pgfpathmoveto{\pgfqpoint{0.000000in}{0.000000in}}%
\pgfpathlineto{\pgfqpoint{0.000000in}{-0.048611in}}%
\pgfusepath{stroke,fill}%
}%
\begin{pgfscope}%
\pgfsys@transformshift{5.146914in}{0.565123in}%
\pgfsys@useobject{currentmarker}{}%
\end{pgfscope}%
\end{pgfscope}%
\begin{pgfscope}%
\definecolor{textcolor}{rgb}{0.000000,0.000000,0.000000}%
\pgfsetstrokecolor{textcolor}%
\pgfsetfillcolor{textcolor}%
\pgftext[x=5.146914in,y=0.467901in,,top]{\color{textcolor}\rmfamily\fontsize{10.000000}{12.000000}\selectfont \(\displaystyle {2.5}\)}%
\end{pgfscope}%
\begin{pgfscope}%
\pgfpathrectangle{\pgfqpoint{0.758026in}{0.565123in}}{\pgfqpoint{5.303239in}{2.284877in}}%
\pgfusepath{clip}%
\pgfsetbuttcap%
\pgfsetroundjoin%
\pgfsetlinewidth{0.501875pt}%
\definecolor{currentstroke}{rgb}{0.690196,0.690196,0.690196}%
\pgfsetstrokecolor{currentstroke}%
\pgfsetstrokeopacity{0.700000}%
\pgfsetdash{{1.850000pt}{0.800000pt}}{0.000000pt}%
\pgfpathmoveto{\pgfqpoint{6.061265in}{0.565123in}}%
\pgfpathlineto{\pgfqpoint{6.061265in}{2.850000in}}%
\pgfusepath{stroke}%
\end{pgfscope}%
\begin{pgfscope}%
\pgfsetbuttcap%
\pgfsetroundjoin%
\definecolor{currentfill}{rgb}{0.000000,0.000000,0.000000}%
\pgfsetfillcolor{currentfill}%
\pgfsetlinewidth{0.803000pt}%
\definecolor{currentstroke}{rgb}{0.000000,0.000000,0.000000}%
\pgfsetstrokecolor{currentstroke}%
\pgfsetdash{}{0pt}%
\pgfsys@defobject{currentmarker}{\pgfqpoint{0.000000in}{-0.048611in}}{\pgfqpoint{0.000000in}{0.000000in}}{%
\pgfpathmoveto{\pgfqpoint{0.000000in}{0.000000in}}%
\pgfpathlineto{\pgfqpoint{0.000000in}{-0.048611in}}%
\pgfusepath{stroke,fill}%
}%
\begin{pgfscope}%
\pgfsys@transformshift{6.061265in}{0.565123in}%
\pgfsys@useobject{currentmarker}{}%
\end{pgfscope}%
\end{pgfscope}%
\begin{pgfscope}%
\definecolor{textcolor}{rgb}{0.000000,0.000000,0.000000}%
\pgfsetstrokecolor{textcolor}%
\pgfsetfillcolor{textcolor}%
\pgftext[x=6.061265in,y=0.467901in,,top]{\color{textcolor}\rmfamily\fontsize{10.000000}{12.000000}\selectfont \(\displaystyle {3.0}\)}%
\end{pgfscope}%
\begin{pgfscope}%
\definecolor{textcolor}{rgb}{0.000000,0.000000,0.000000}%
\pgfsetstrokecolor{textcolor}%
\pgfsetfillcolor{textcolor}%
\pgftext[x=3.409645in,y=0.288889in,,top]{\color{textcolor}\rmfamily\fontsize{10.000000}{12.000000}\selectfont Frequency (Hz)}%
\end{pgfscope}%
\begin{pgfscope}%
\definecolor{textcolor}{rgb}{0.000000,0.000000,0.000000}%
\pgfsetstrokecolor{textcolor}%
\pgfsetfillcolor{textcolor}%
\pgftext[x=6.061265in,y=0.302778in,right,top]{\color{textcolor}\rmfamily\fontsize{10.000000}{12.000000}\selectfont \(\displaystyle \times{10^{6}}{}\)}%
\end{pgfscope}%
\begin{pgfscope}%
\pgfpathrectangle{\pgfqpoint{0.758026in}{0.565123in}}{\pgfqpoint{5.303239in}{2.284877in}}%
\pgfusepath{clip}%
\pgfsetbuttcap%
\pgfsetroundjoin%
\pgfsetlinewidth{0.501875pt}%
\definecolor{currentstroke}{rgb}{0.690196,0.690196,0.690196}%
\pgfsetstrokecolor{currentstroke}%
\pgfsetstrokeopacity{0.700000}%
\pgfsetdash{{1.850000pt}{0.800000pt}}{0.000000pt}%
\pgfpathmoveto{\pgfqpoint{0.758026in}{0.988756in}}%
\pgfpathlineto{\pgfqpoint{6.061265in}{0.988756in}}%
\pgfusepath{stroke}%
\end{pgfscope}%
\begin{pgfscope}%
\pgfsetbuttcap%
\pgfsetroundjoin%
\definecolor{currentfill}{rgb}{0.000000,0.000000,0.000000}%
\pgfsetfillcolor{currentfill}%
\pgfsetlinewidth{0.803000pt}%
\definecolor{currentstroke}{rgb}{0.000000,0.000000,0.000000}%
\pgfsetstrokecolor{currentstroke}%
\pgfsetdash{}{0pt}%
\pgfsys@defobject{currentmarker}{\pgfqpoint{-0.048611in}{0.000000in}}{\pgfqpoint{-0.000000in}{0.000000in}}{%
\pgfpathmoveto{\pgfqpoint{-0.000000in}{0.000000in}}%
\pgfpathlineto{\pgfqpoint{-0.048611in}{0.000000in}}%
\pgfusepath{stroke,fill}%
}%
\begin{pgfscope}%
\pgfsys@transformshift{0.758026in}{0.988756in}%
\pgfsys@useobject{currentmarker}{}%
\end{pgfscope}%
\end{pgfscope}%
\begin{pgfscope}%
\definecolor{textcolor}{rgb}{0.000000,0.000000,0.000000}%
\pgfsetstrokecolor{textcolor}%
\pgfsetfillcolor{textcolor}%
\pgftext[x=0.344444in, y=0.940530in, left, base]{\color{textcolor}\rmfamily\fontsize{10.000000}{12.000000}\selectfont \(\displaystyle {\ensuremath{-}100}\)}%
\end{pgfscope}%
\begin{pgfscope}%
\pgfpathrectangle{\pgfqpoint{0.758026in}{0.565123in}}{\pgfqpoint{5.303239in}{2.284877in}}%
\pgfusepath{clip}%
\pgfsetbuttcap%
\pgfsetroundjoin%
\pgfsetlinewidth{0.501875pt}%
\definecolor{currentstroke}{rgb}{0.690196,0.690196,0.690196}%
\pgfsetstrokecolor{currentstroke}%
\pgfsetstrokeopacity{0.700000}%
\pgfsetdash{{1.850000pt}{0.800000pt}}{0.000000pt}%
\pgfpathmoveto{\pgfqpoint{0.758026in}{1.427931in}}%
\pgfpathlineto{\pgfqpoint{6.061265in}{1.427931in}}%
\pgfusepath{stroke}%
\end{pgfscope}%
\begin{pgfscope}%
\pgfsetbuttcap%
\pgfsetroundjoin%
\definecolor{currentfill}{rgb}{0.000000,0.000000,0.000000}%
\pgfsetfillcolor{currentfill}%
\pgfsetlinewidth{0.803000pt}%
\definecolor{currentstroke}{rgb}{0.000000,0.000000,0.000000}%
\pgfsetstrokecolor{currentstroke}%
\pgfsetdash{}{0pt}%
\pgfsys@defobject{currentmarker}{\pgfqpoint{-0.048611in}{0.000000in}}{\pgfqpoint{-0.000000in}{0.000000in}}{%
\pgfpathmoveto{\pgfqpoint{-0.000000in}{0.000000in}}%
\pgfpathlineto{\pgfqpoint{-0.048611in}{0.000000in}}%
\pgfusepath{stroke,fill}%
}%
\begin{pgfscope}%
\pgfsys@transformshift{0.758026in}{1.427931in}%
\pgfsys@useobject{currentmarker}{}%
\end{pgfscope}%
\end{pgfscope}%
\begin{pgfscope}%
\definecolor{textcolor}{rgb}{0.000000,0.000000,0.000000}%
\pgfsetstrokecolor{textcolor}%
\pgfsetfillcolor{textcolor}%
\pgftext[x=0.413889in, y=1.379705in, left, base]{\color{textcolor}\rmfamily\fontsize{10.000000}{12.000000}\selectfont \(\displaystyle {\ensuremath{-}80}\)}%
\end{pgfscope}%
\begin{pgfscope}%
\pgfpathrectangle{\pgfqpoint{0.758026in}{0.565123in}}{\pgfqpoint{5.303239in}{2.284877in}}%
\pgfusepath{clip}%
\pgfsetbuttcap%
\pgfsetroundjoin%
\pgfsetlinewidth{0.501875pt}%
\definecolor{currentstroke}{rgb}{0.690196,0.690196,0.690196}%
\pgfsetstrokecolor{currentstroke}%
\pgfsetstrokeopacity{0.700000}%
\pgfsetdash{{1.850000pt}{0.800000pt}}{0.000000pt}%
\pgfpathmoveto{\pgfqpoint{0.758026in}{1.867106in}}%
\pgfpathlineto{\pgfqpoint{6.061265in}{1.867106in}}%
\pgfusepath{stroke}%
\end{pgfscope}%
\begin{pgfscope}%
\pgfsetbuttcap%
\pgfsetroundjoin%
\definecolor{currentfill}{rgb}{0.000000,0.000000,0.000000}%
\pgfsetfillcolor{currentfill}%
\pgfsetlinewidth{0.803000pt}%
\definecolor{currentstroke}{rgb}{0.000000,0.000000,0.000000}%
\pgfsetstrokecolor{currentstroke}%
\pgfsetdash{}{0pt}%
\pgfsys@defobject{currentmarker}{\pgfqpoint{-0.048611in}{0.000000in}}{\pgfqpoint{-0.000000in}{0.000000in}}{%
\pgfpathmoveto{\pgfqpoint{-0.000000in}{0.000000in}}%
\pgfpathlineto{\pgfqpoint{-0.048611in}{0.000000in}}%
\pgfusepath{stroke,fill}%
}%
\begin{pgfscope}%
\pgfsys@transformshift{0.758026in}{1.867106in}%
\pgfsys@useobject{currentmarker}{}%
\end{pgfscope}%
\end{pgfscope}%
\begin{pgfscope}%
\definecolor{textcolor}{rgb}{0.000000,0.000000,0.000000}%
\pgfsetstrokecolor{textcolor}%
\pgfsetfillcolor{textcolor}%
\pgftext[x=0.413889in, y=1.818880in, left, base]{\color{textcolor}\rmfamily\fontsize{10.000000}{12.000000}\selectfont \(\displaystyle {\ensuremath{-}60}\)}%
\end{pgfscope}%
\begin{pgfscope}%
\pgfpathrectangle{\pgfqpoint{0.758026in}{0.565123in}}{\pgfqpoint{5.303239in}{2.284877in}}%
\pgfusepath{clip}%
\pgfsetbuttcap%
\pgfsetroundjoin%
\pgfsetlinewidth{0.501875pt}%
\definecolor{currentstroke}{rgb}{0.690196,0.690196,0.690196}%
\pgfsetstrokecolor{currentstroke}%
\pgfsetstrokeopacity{0.700000}%
\pgfsetdash{{1.850000pt}{0.800000pt}}{0.000000pt}%
\pgfpathmoveto{\pgfqpoint{0.758026in}{2.306281in}}%
\pgfpathlineto{\pgfqpoint{6.061265in}{2.306281in}}%
\pgfusepath{stroke}%
\end{pgfscope}%
\begin{pgfscope}%
\pgfsetbuttcap%
\pgfsetroundjoin%
\definecolor{currentfill}{rgb}{0.000000,0.000000,0.000000}%
\pgfsetfillcolor{currentfill}%
\pgfsetlinewidth{0.803000pt}%
\definecolor{currentstroke}{rgb}{0.000000,0.000000,0.000000}%
\pgfsetstrokecolor{currentstroke}%
\pgfsetdash{}{0pt}%
\pgfsys@defobject{currentmarker}{\pgfqpoint{-0.048611in}{0.000000in}}{\pgfqpoint{-0.000000in}{0.000000in}}{%
\pgfpathmoveto{\pgfqpoint{-0.000000in}{0.000000in}}%
\pgfpathlineto{\pgfqpoint{-0.048611in}{0.000000in}}%
\pgfusepath{stroke,fill}%
}%
\begin{pgfscope}%
\pgfsys@transformshift{0.758026in}{2.306281in}%
\pgfsys@useobject{currentmarker}{}%
\end{pgfscope}%
\end{pgfscope}%
\begin{pgfscope}%
\definecolor{textcolor}{rgb}{0.000000,0.000000,0.000000}%
\pgfsetstrokecolor{textcolor}%
\pgfsetfillcolor{textcolor}%
\pgftext[x=0.413889in, y=2.258055in, left, base]{\color{textcolor}\rmfamily\fontsize{10.000000}{12.000000}\selectfont \(\displaystyle {\ensuremath{-}40}\)}%
\end{pgfscope}%
\begin{pgfscope}%
\pgfpathrectangle{\pgfqpoint{0.758026in}{0.565123in}}{\pgfqpoint{5.303239in}{2.284877in}}%
\pgfusepath{clip}%
\pgfsetbuttcap%
\pgfsetroundjoin%
\pgfsetlinewidth{0.501875pt}%
\definecolor{currentstroke}{rgb}{0.690196,0.690196,0.690196}%
\pgfsetstrokecolor{currentstroke}%
\pgfsetstrokeopacity{0.700000}%
\pgfsetdash{{1.850000pt}{0.800000pt}}{0.000000pt}%
\pgfpathmoveto{\pgfqpoint{0.758026in}{2.745456in}}%
\pgfpathlineto{\pgfqpoint{6.061265in}{2.745456in}}%
\pgfusepath{stroke}%
\end{pgfscope}%
\begin{pgfscope}%
\pgfsetbuttcap%
\pgfsetroundjoin%
\definecolor{currentfill}{rgb}{0.000000,0.000000,0.000000}%
\pgfsetfillcolor{currentfill}%
\pgfsetlinewidth{0.803000pt}%
\definecolor{currentstroke}{rgb}{0.000000,0.000000,0.000000}%
\pgfsetstrokecolor{currentstroke}%
\pgfsetdash{}{0pt}%
\pgfsys@defobject{currentmarker}{\pgfqpoint{-0.048611in}{0.000000in}}{\pgfqpoint{-0.000000in}{0.000000in}}{%
\pgfpathmoveto{\pgfqpoint{-0.000000in}{0.000000in}}%
\pgfpathlineto{\pgfqpoint{-0.048611in}{0.000000in}}%
\pgfusepath{stroke,fill}%
}%
\begin{pgfscope}%
\pgfsys@transformshift{0.758026in}{2.745456in}%
\pgfsys@useobject{currentmarker}{}%
\end{pgfscope}%
\end{pgfscope}%
\begin{pgfscope}%
\definecolor{textcolor}{rgb}{0.000000,0.000000,0.000000}%
\pgfsetstrokecolor{textcolor}%
\pgfsetfillcolor{textcolor}%
\pgftext[x=0.413889in, y=2.697230in, left, base]{\color{textcolor}\rmfamily\fontsize{10.000000}{12.000000}\selectfont \(\displaystyle {\ensuremath{-}20}\)}%
\end{pgfscope}%
\begin{pgfscope}%
\definecolor{textcolor}{rgb}{0.000000,0.000000,0.000000}%
\pgfsetstrokecolor{textcolor}%
\pgfsetfillcolor{textcolor}%
\pgftext[x=0.288889in,y=1.707562in,,bottom,rotate=90.000000]{\color{textcolor}\rmfamily\fontsize{10.000000}{12.000000}\selectfont Amplitude (dB)}%
\end{pgfscope}%
\begin{pgfscope}%
\pgfpathrectangle{\pgfqpoint{0.758026in}{0.565123in}}{\pgfqpoint{5.303239in}{2.284877in}}%
\pgfusepath{clip}%
\pgfsetrectcap%
\pgfsetroundjoin%
\pgfsetlinewidth{1.003750pt}%
\definecolor{currentstroke}{rgb}{0.000000,0.000000,0.000000}%
\pgfsetstrokecolor{currentstroke}%
\pgfsetdash{}{0pt}%
\pgfpathmoveto{\pgfqpoint{0.758026in}{1.455379in}}%
\pgfpathlineto{\pgfqpoint{0.766986in}{1.378524in}}%
\pgfpathlineto{\pgfqpoint{0.775947in}{1.509590in}}%
\pgfpathlineto{\pgfqpoint{0.784908in}{1.443714in}}%
\pgfpathlineto{\pgfqpoint{0.793868in}{1.508904in}}%
\pgfpathlineto{\pgfqpoint{0.811790in}{1.508217in}}%
\pgfpathlineto{\pgfqpoint{0.820750in}{1.408717in}}%
\pgfpathlineto{\pgfqpoint{0.829711in}{1.440969in}}%
\pgfpathlineto{\pgfqpoint{0.838672in}{1.407344in}}%
\pgfpathlineto{\pgfqpoint{0.847632in}{1.539097in}}%
\pgfpathlineto{\pgfqpoint{0.856593in}{1.406658in}}%
\pgfpathlineto{\pgfqpoint{0.865553in}{1.471848in}}%
\pgfpathlineto{\pgfqpoint{0.874514in}{1.373034in}}%
\pgfpathlineto{\pgfqpoint{0.883475in}{1.405286in}}%
\pgfpathlineto{\pgfqpoint{0.892435in}{1.613208in}}%
\pgfpathlineto{\pgfqpoint{0.901396in}{1.481455in}}%
\pgfpathlineto{\pgfqpoint{0.910357in}{1.568604in}}%
\pgfpathlineto{\pgfqpoint{0.919317in}{1.414206in}}%
\pgfpathlineto{\pgfqpoint{0.928278in}{1.480083in}}%
\pgfpathlineto{\pgfqpoint{0.937239in}{1.534293in}}%
\pgfpathlineto{\pgfqpoint{0.946199in}{1.401855in}}%
\pgfpathlineto{\pgfqpoint{0.955160in}{1.434793in}}%
\pgfpathlineto{\pgfqpoint{0.964121in}{1.478024in}}%
\pgfpathlineto{\pgfqpoint{0.973081in}{1.466359in}}%
\pgfpathlineto{\pgfqpoint{0.982042in}{1.466359in}}%
\pgfpathlineto{\pgfqpoint{0.991002in}{1.443714in}}%
\pgfpathlineto{\pgfqpoint{0.999963in}{1.432048in}}%
\pgfpathlineto{\pgfqpoint{1.008924in}{1.541842in}}%
\pgfpathlineto{\pgfqpoint{1.017884in}{1.607032in}}%
\pgfpathlineto{\pgfqpoint{1.026845in}{1.474593in}}%
\pgfpathlineto{\pgfqpoint{1.035806in}{1.561742in}}%
\pgfpathlineto{\pgfqpoint{1.053727in}{1.517138in}}%
\pgfpathlineto{\pgfqpoint{1.062688in}{1.571349in}}%
\pgfpathlineto{\pgfqpoint{1.071648in}{1.516452in}}%
\pgfpathlineto{\pgfqpoint{1.080609in}{1.405972in}}%
\pgfpathlineto{\pgfqpoint{1.089570in}{1.405286in}}%
\pgfpathlineto{\pgfqpoint{1.098530in}{1.460183in}}%
\pgfpathlineto{\pgfqpoint{1.107491in}{1.503414in}}%
\pgfpathlineto{\pgfqpoint{1.116452in}{1.513707in}}%
\pgfpathlineto{\pgfqpoint{1.125412in}{1.590563in}}%
\pgfpathlineto{\pgfqpoint{1.143333in}{1.358623in}}%
\pgfpathlineto{\pgfqpoint{1.152294in}{1.556252in}}%
\pgfpathlineto{\pgfqpoint{1.161255in}{1.522628in}}%
\pgfpathlineto{\pgfqpoint{1.170215in}{1.478024in}}%
\pgfpathlineto{\pgfqpoint{1.179176in}{1.521942in}}%
\pgfpathlineto{\pgfqpoint{1.188137in}{1.499297in}}%
\pgfpathlineto{\pgfqpoint{1.197097in}{1.498610in}}%
\pgfpathlineto{\pgfqpoint{1.215019in}{1.585759in}}%
\pgfpathlineto{\pgfqpoint{1.223979in}{1.508217in}}%
\pgfpathlineto{\pgfqpoint{1.232940in}{1.518511in}}%
\pgfpathlineto{\pgfqpoint{1.241901in}{1.496552in}}%
\pgfpathlineto{\pgfqpoint{1.250861in}{1.451948in}}%
\pgfpathlineto{\pgfqpoint{1.259822in}{1.583014in}}%
\pgfpathlineto{\pgfqpoint{1.268783in}{1.550076in}}%
\pgfpathlineto{\pgfqpoint{1.277743in}{1.582328in}}%
\pgfpathlineto{\pgfqpoint{1.286704in}{1.460869in}}%
\pgfpathlineto{\pgfqpoint{1.295664in}{1.482828in}}%
\pgfpathlineto{\pgfqpoint{1.304625in}{1.701729in}}%
\pgfpathlineto{\pgfqpoint{1.313586in}{1.459496in}}%
\pgfpathlineto{\pgfqpoint{1.322546in}{1.492435in}}%
\pgfpathlineto{\pgfqpoint{1.331507in}{1.491748in}}%
\pgfpathlineto{\pgfqpoint{1.340468in}{1.447145in}}%
\pgfpathlineto{\pgfqpoint{1.349428in}{1.381268in}}%
\pgfpathlineto{\pgfqpoint{1.358389in}{1.380582in}}%
\pgfpathlineto{\pgfqpoint{1.367350in}{1.489690in}}%
\pgfpathlineto{\pgfqpoint{1.376310in}{1.379896in}}%
\pgfpathlineto{\pgfqpoint{1.385271in}{1.346272in}}%
\pgfpathlineto{\pgfqpoint{1.394232in}{1.389503in}}%
\pgfpathlineto{\pgfqpoint{1.403192in}{1.477338in}}%
\pgfpathlineto{\pgfqpoint{1.412153in}{1.454693in}}%
\pgfpathlineto{\pgfqpoint{1.421113in}{1.443027in}}%
\pgfpathlineto{\pgfqpoint{1.430074in}{1.574780in}}%
\pgfpathlineto{\pgfqpoint{1.439035in}{1.628991in}}%
\pgfpathlineto{\pgfqpoint{1.447995in}{1.902789in}}%
\pgfpathlineto{\pgfqpoint{1.456956in}{1.518511in}}%
\pgfpathlineto{\pgfqpoint{1.465917in}{2.242463in}}%
\pgfpathlineto{\pgfqpoint{1.474877in}{2.593117in}}%
\pgfpathlineto{\pgfqpoint{1.483838in}{2.724183in}}%
\pgfpathlineto{\pgfqpoint{1.492799in}{2.746142in}}%
\pgfpathlineto{\pgfqpoint{1.501759in}{2.657621in}}%
\pgfpathlineto{\pgfqpoint{1.510720in}{2.426368in}}%
\pgfpathlineto{\pgfqpoint{1.519681in}{1.789564in}}%
\pgfpathlineto{\pgfqpoint{1.528641in}{1.624187in}}%
\pgfpathlineto{\pgfqpoint{1.537602in}{1.678398in}}%
\pgfpathlineto{\pgfqpoint{1.546563in}{1.854068in}}%
\pgfpathlineto{\pgfqpoint{1.555523in}{1.578897in}}%
\pgfpathlineto{\pgfqpoint{1.564484in}{1.369603in}}%
\pgfpathlineto{\pgfqpoint{1.573444in}{1.545273in}}%
\pgfpathlineto{\pgfqpoint{1.582405in}{1.379896in}}%
\pgfpathlineto{\pgfqpoint{1.591366in}{1.510962in}}%
\pgfpathlineto{\pgfqpoint{1.600326in}{1.478024in}}%
\pgfpathlineto{\pgfqpoint{1.609287in}{1.620070in}}%
\pgfpathlineto{\pgfqpoint{1.618248in}{1.443714in}}%
\pgfpathlineto{\pgfqpoint{1.627208in}{1.333920in}}%
\pgfpathlineto{\pgfqpoint{1.680972in}{1.330489in}}%
\pgfpathlineto{\pgfqpoint{1.698894in}{1.329803in}}%
\pgfpathlineto{\pgfqpoint{1.734736in}{1.327744in}}%
\pgfpathlineto{\pgfqpoint{1.752657in}{1.327058in}}%
\pgfpathlineto{\pgfqpoint{1.761618in}{1.359310in}}%
\pgfpathlineto{\pgfqpoint{1.770579in}{1.326371in}}%
\pgfpathlineto{\pgfqpoint{1.788500in}{1.324999in}}%
\pgfpathlineto{\pgfqpoint{1.797461in}{1.357937in}}%
\pgfpathlineto{\pgfqpoint{1.806421in}{1.324313in}}%
\pgfpathlineto{\pgfqpoint{1.815382in}{1.356565in}}%
\pgfpathlineto{\pgfqpoint{1.824343in}{1.356565in}}%
\pgfpathlineto{\pgfqpoint{1.833303in}{1.322940in}}%
\pgfpathlineto{\pgfqpoint{1.869146in}{1.320882in}}%
\pgfpathlineto{\pgfqpoint{1.887067in}{1.320196in}}%
\pgfpathlineto{\pgfqpoint{1.931870in}{1.317451in}}%
\pgfpathlineto{\pgfqpoint{1.949792in}{1.316765in}}%
\pgfpathlineto{\pgfqpoint{1.985634in}{1.314706in}}%
\pgfpathlineto{\pgfqpoint{2.003555in}{1.314020in}}%
\pgfpathlineto{\pgfqpoint{2.030437in}{1.312647in}}%
\pgfpathlineto{\pgfqpoint{2.039398in}{1.279023in}}%
\pgfpathlineto{\pgfqpoint{2.048359in}{1.311961in}}%
\pgfpathlineto{\pgfqpoint{2.057319in}{1.311275in}}%
\pgfpathlineto{\pgfqpoint{2.066280in}{1.244712in}}%
\pgfpathlineto{\pgfqpoint{2.075241in}{1.244712in}}%
\pgfpathlineto{\pgfqpoint{2.084201in}{1.309902in}}%
\pgfpathlineto{\pgfqpoint{2.093162in}{1.309216in}}%
\pgfpathlineto{\pgfqpoint{2.102123in}{1.385386in}}%
\pgfpathlineto{\pgfqpoint{2.111083in}{1.242654in}}%
\pgfpathlineto{\pgfqpoint{2.120044in}{1.241968in}}%
\pgfpathlineto{\pgfqpoint{2.129005in}{1.307158in}}%
\pgfpathlineto{\pgfqpoint{2.137965in}{1.252261in}}%
\pgfpathlineto{\pgfqpoint{2.146926in}{1.207657in}}%
\pgfpathlineto{\pgfqpoint{2.155886in}{1.174033in}}%
\pgfpathlineto{\pgfqpoint{2.164847in}{1.174033in}}%
\pgfpathlineto{\pgfqpoint{2.173808in}{1.272161in}}%
\pgfpathlineto{\pgfqpoint{2.182768in}{1.205598in}}%
\pgfpathlineto{\pgfqpoint{2.191729in}{1.172660in}}%
\pgfpathlineto{\pgfqpoint{2.200690in}{1.248830in}}%
\pgfpathlineto{\pgfqpoint{2.209650in}{1.171288in}}%
\pgfpathlineto{\pgfqpoint{2.218611in}{1.171288in}}%
\pgfpathlineto{\pgfqpoint{2.227572in}{1.203540in}}%
\pgfpathlineto{\pgfqpoint{2.236532in}{1.169915in}}%
\pgfpathlineto{\pgfqpoint{2.245493in}{1.268730in}}%
\pgfpathlineto{\pgfqpoint{2.254454in}{1.202167in}}%
\pgfpathlineto{\pgfqpoint{2.263414in}{1.300295in}}%
\pgfpathlineto{\pgfqpoint{2.272375in}{1.201481in}}%
\pgfpathlineto{\pgfqpoint{2.281336in}{1.266671in}}%
\pgfpathlineto{\pgfqpoint{2.290296in}{1.167171in}}%
\pgfpathlineto{\pgfqpoint{2.308217in}{1.166484in}}%
\pgfpathlineto{\pgfqpoint{2.317178in}{1.242654in}}%
\pgfpathlineto{\pgfqpoint{2.326139in}{1.165112in}}%
\pgfpathlineto{\pgfqpoint{2.344060in}{1.164426in}}%
\pgfpathlineto{\pgfqpoint{2.353021in}{1.229616in}}%
\pgfpathlineto{\pgfqpoint{2.361981in}{1.163739in}}%
\pgfpathlineto{\pgfqpoint{2.379903in}{1.162367in}}%
\pgfpathlineto{\pgfqpoint{2.388863in}{1.349016in}}%
\pgfpathlineto{\pgfqpoint{2.397824in}{1.502041in}}%
\pgfpathlineto{\pgfqpoint{2.406785in}{1.523314in}}%
\pgfpathlineto{\pgfqpoint{2.415745in}{1.468417in}}%
\pgfpathlineto{\pgfqpoint{2.424706in}{1.237850in}}%
\pgfpathlineto{\pgfqpoint{2.433666in}{1.095118in}}%
\pgfpathlineto{\pgfqpoint{2.442627in}{1.193933in}}%
\pgfpathlineto{\pgfqpoint{2.451588in}{1.171974in}}%
\pgfpathlineto{\pgfqpoint{2.460548in}{1.193933in}}%
\pgfpathlineto{\pgfqpoint{2.469509in}{1.095118in}}%
\pgfpathlineto{\pgfqpoint{2.478470in}{1.095118in}}%
\pgfpathlineto{\pgfqpoint{2.496391in}{1.160995in}}%
\pgfpathlineto{\pgfqpoint{2.505352in}{1.128057in}}%
\pgfpathlineto{\pgfqpoint{2.514312in}{1.193933in}}%
\pgfpathlineto{\pgfqpoint{2.523273in}{1.128057in}}%
\pgfpathlineto{\pgfqpoint{2.532234in}{1.160995in}}%
\pgfpathlineto{\pgfqpoint{2.541194in}{1.160995in}}%
\pgfpathlineto{\pgfqpoint{2.550155in}{1.204912in}}%
\pgfpathlineto{\pgfqpoint{2.559116in}{1.160995in}}%
\pgfpathlineto{\pgfqpoint{2.568076in}{1.193933in}}%
\pgfpathlineto{\pgfqpoint{2.577037in}{1.171974in}}%
\pgfpathlineto{\pgfqpoint{2.585997in}{1.160995in}}%
\pgfpathlineto{\pgfqpoint{2.594958in}{1.095118in}}%
\pgfpathlineto{\pgfqpoint{2.603919in}{1.128057in}}%
\pgfpathlineto{\pgfqpoint{2.612879in}{1.128057in}}%
\pgfpathlineto{\pgfqpoint{2.621840in}{1.095118in}}%
\pgfpathlineto{\pgfqpoint{2.630801in}{1.171974in}}%
\pgfpathlineto{\pgfqpoint{2.639761in}{1.095118in}}%
\pgfpathlineto{\pgfqpoint{2.648722in}{1.128057in}}%
\pgfpathlineto{\pgfqpoint{2.657683in}{1.095118in}}%
\pgfpathlineto{\pgfqpoint{2.666643in}{1.128057in}}%
\pgfpathlineto{\pgfqpoint{2.684565in}{1.128057in}}%
\pgfpathlineto{\pgfqpoint{2.693525in}{1.095118in}}%
\pgfpathlineto{\pgfqpoint{2.702486in}{1.128057in}}%
\pgfpathlineto{\pgfqpoint{2.711447in}{1.095118in}}%
\pgfpathlineto{\pgfqpoint{2.720407in}{1.128057in}}%
\pgfpathlineto{\pgfqpoint{2.729368in}{1.094432in}}%
\pgfpathlineto{\pgfqpoint{2.783132in}{1.094432in}}%
\pgfpathlineto{\pgfqpoint{2.792092in}{1.127370in}}%
\pgfpathlineto{\pgfqpoint{2.801053in}{1.094432in}}%
\pgfpathlineto{\pgfqpoint{2.926502in}{1.094432in}}%
\pgfpathlineto{\pgfqpoint{2.935463in}{1.028556in}}%
\pgfpathlineto{\pgfqpoint{2.953384in}{1.028556in}}%
\pgfpathlineto{\pgfqpoint{2.962345in}{1.127370in}}%
\pgfpathlineto{\pgfqpoint{2.971305in}{1.094432in}}%
\pgfpathlineto{\pgfqpoint{2.980266in}{1.094432in}}%
\pgfpathlineto{\pgfqpoint{2.989227in}{1.028556in}}%
\pgfpathlineto{\pgfqpoint{2.998187in}{1.028556in}}%
\pgfpathlineto{\pgfqpoint{3.007148in}{1.039535in}}%
\pgfpathlineto{\pgfqpoint{3.016108in}{1.072473in}}%
\pgfpathlineto{\pgfqpoint{3.025069in}{1.072473in}}%
\pgfpathlineto{\pgfqpoint{3.034030in}{1.038849in}}%
\pgfpathlineto{\pgfqpoint{3.042990in}{1.148643in}}%
\pgfpathlineto{\pgfqpoint{3.051951in}{1.357251in}}%
\pgfpathlineto{\pgfqpoint{3.060912in}{1.093746in}}%
\pgfpathlineto{\pgfqpoint{3.069872in}{1.203540in}}%
\pgfpathlineto{\pgfqpoint{3.078833in}{1.060808in}}%
\pgfpathlineto{\pgfqpoint{3.087794in}{1.203540in}}%
\pgfpathlineto{\pgfqpoint{3.096754in}{1.148643in}}%
\pgfpathlineto{\pgfqpoint{3.105715in}{1.104725in}}%
\pgfpathlineto{\pgfqpoint{3.114676in}{1.126684in}}%
\pgfpathlineto{\pgfqpoint{3.123636in}{1.005911in}}%
\pgfpathlineto{\pgfqpoint{3.132597in}{1.159622in}}%
\pgfpathlineto{\pgfqpoint{3.141558in}{1.104725in}}%
\pgfpathlineto{\pgfqpoint{3.150518in}{1.126684in}}%
\pgfpathlineto{\pgfqpoint{3.159479in}{1.027870in}}%
\pgfpathlineto{\pgfqpoint{3.168439in}{1.170602in}}%
\pgfpathlineto{\pgfqpoint{3.177400in}{1.071787in}}%
\pgfpathlineto{\pgfqpoint{3.186361in}{1.093746in}}%
\pgfpathlineto{\pgfqpoint{3.195321in}{1.104725in}}%
\pgfpathlineto{\pgfqpoint{3.204282in}{1.148643in}}%
\pgfpathlineto{\pgfqpoint{3.213243in}{1.181581in}}%
\pgfpathlineto{\pgfqpoint{3.222203in}{1.269416in}}%
\pgfpathlineto{\pgfqpoint{3.231164in}{1.159622in}}%
\pgfpathlineto{\pgfqpoint{3.240125in}{1.170602in}}%
\pgfpathlineto{\pgfqpoint{3.249085in}{1.368230in}}%
\pgfpathlineto{\pgfqpoint{3.258046in}{1.445086in}}%
\pgfpathlineto{\pgfqpoint{3.267007in}{1.335292in}}%
\pgfpathlineto{\pgfqpoint{3.275967in}{1.368230in}}%
\pgfpathlineto{\pgfqpoint{3.284928in}{1.214519in}}%
\pgfpathlineto{\pgfqpoint{3.293889in}{1.818385in}}%
\pgfpathlineto{\pgfqpoint{3.302849in}{2.268539in}}%
\pgfpathlineto{\pgfqpoint{3.311810in}{2.455189in}}%
\pgfpathlineto{\pgfqpoint{3.320770in}{2.521065in}}%
\pgfpathlineto{\pgfqpoint{3.329731in}{2.466168in}}%
\pgfpathlineto{\pgfqpoint{3.338692in}{2.268539in}}%
\pgfpathlineto{\pgfqpoint{3.356613in}{1.334606in}}%
\pgfpathlineto{\pgfqpoint{3.365574in}{1.400482in}}%
\pgfpathlineto{\pgfqpoint{3.374534in}{1.345585in}}%
\pgfpathlineto{\pgfqpoint{3.383495in}{1.323627in}}%
\pgfpathlineto{\pgfqpoint{3.392456in}{1.334606in}}%
\pgfpathlineto{\pgfqpoint{3.401416in}{1.235792in}}%
\pgfpathlineto{\pgfqpoint{3.410377in}{1.005225in}}%
\pgfpathlineto{\pgfqpoint{3.419338in}{1.411462in}}%
\pgfpathlineto{\pgfqpoint{3.428298in}{1.466359in}}%
\pgfpathlineto{\pgfqpoint{3.437259in}{1.499297in}}%
\pgfpathlineto{\pgfqpoint{3.446219in}{1.422441in}}%
\pgfpathlineto{\pgfqpoint{3.455180in}{1.246771in}}%
\pgfpathlineto{\pgfqpoint{3.464141in}{1.038163in}}%
\pgfpathlineto{\pgfqpoint{3.473101in}{1.027183in}}%
\pgfpathlineto{\pgfqpoint{3.482062in}{1.158936in}}%
\pgfpathlineto{\pgfqpoint{3.491023in}{1.071101in}}%
\pgfpathlineto{\pgfqpoint{3.499983in}{1.027183in}}%
\pgfpathlineto{\pgfqpoint{3.508944in}{1.093060in}}%
\pgfpathlineto{\pgfqpoint{3.517905in}{0.994245in}}%
\pgfpathlineto{\pgfqpoint{3.526865in}{1.027183in}}%
\pgfpathlineto{\pgfqpoint{3.535826in}{1.027183in}}%
\pgfpathlineto{\pgfqpoint{3.544787in}{0.994245in}}%
\pgfpathlineto{\pgfqpoint{3.553747in}{1.060122in}}%
\pgfpathlineto{\pgfqpoint{3.562708in}{1.093060in}}%
\pgfpathlineto{\pgfqpoint{3.571669in}{1.158936in}}%
\pgfpathlineto{\pgfqpoint{3.580629in}{1.027183in}}%
\pgfpathlineto{\pgfqpoint{3.589590in}{1.180895in}}%
\pgfpathlineto{\pgfqpoint{3.598550in}{1.104039in}}%
\pgfpathlineto{\pgfqpoint{3.607511in}{1.125998in}}%
\pgfpathlineto{\pgfqpoint{3.616472in}{1.027183in}}%
\pgfpathlineto{\pgfqpoint{3.625432in}{1.060122in}}%
\pgfpathlineto{\pgfqpoint{3.634393in}{0.994245in}}%
\pgfpathlineto{\pgfqpoint{3.643354in}{1.027183in}}%
\pgfpathlineto{\pgfqpoint{3.652314in}{1.005225in}}%
\pgfpathlineto{\pgfqpoint{3.661275in}{1.060122in}}%
\pgfpathlineto{\pgfqpoint{3.670236in}{1.180209in}}%
\pgfpathlineto{\pgfqpoint{3.679196in}{1.147270in}}%
\pgfpathlineto{\pgfqpoint{3.688157in}{1.092374in}}%
\pgfpathlineto{\pgfqpoint{3.697118in}{1.070415in}}%
\pgfpathlineto{\pgfqpoint{3.706078in}{1.070415in}}%
\pgfpathlineto{\pgfqpoint{3.715039in}{1.180209in}}%
\pgfpathlineto{\pgfqpoint{3.724000in}{1.103353in}}%
\pgfpathlineto{\pgfqpoint{3.732960in}{1.180209in}}%
\pgfpathlineto{\pgfqpoint{3.741921in}{1.103353in}}%
\pgfpathlineto{\pgfqpoint{3.750881in}{1.180209in}}%
\pgfpathlineto{\pgfqpoint{3.759842in}{1.147270in}}%
\pgfpathlineto{\pgfqpoint{3.768803in}{1.059435in}}%
\pgfpathlineto{\pgfqpoint{3.777763in}{1.092374in}}%
\pgfpathlineto{\pgfqpoint{3.786724in}{1.191188in}}%
\pgfpathlineto{\pgfqpoint{3.795685in}{1.092374in}}%
\pgfpathlineto{\pgfqpoint{3.804645in}{1.037477in}}%
\pgfpathlineto{\pgfqpoint{3.813606in}{1.070415in}}%
\pgfpathlineto{\pgfqpoint{3.822567in}{1.213147in}}%
\pgfpathlineto{\pgfqpoint{3.831527in}{1.114332in}}%
\pgfpathlineto{\pgfqpoint{3.840488in}{0.993559in}}%
\pgfpathlineto{\pgfqpoint{3.849449in}{1.125312in}}%
\pgfpathlineto{\pgfqpoint{3.858409in}{1.059435in}}%
\pgfpathlineto{\pgfqpoint{3.867370in}{1.180209in}}%
\pgfpathlineto{\pgfqpoint{3.876330in}{1.180209in}}%
\pgfpathlineto{\pgfqpoint{3.885291in}{1.026497in}}%
\pgfpathlineto{\pgfqpoint{3.894252in}{1.147270in}}%
\pgfpathlineto{\pgfqpoint{3.903212in}{1.070415in}}%
\pgfpathlineto{\pgfqpoint{3.912173in}{1.070415in}}%
\pgfpathlineto{\pgfqpoint{3.930094in}{1.202167in}}%
\pgfpathlineto{\pgfqpoint{3.939055in}{1.037477in}}%
\pgfpathlineto{\pgfqpoint{3.948016in}{1.059435in}}%
\pgfpathlineto{\pgfqpoint{3.956976in}{1.147270in}}%
\pgfpathlineto{\pgfqpoint{3.965937in}{1.070415in}}%
\pgfpathlineto{\pgfqpoint{3.974898in}{1.234419in}}%
\pgfpathlineto{\pgfqpoint{3.983858in}{1.124625in}}%
\pgfpathlineto{\pgfqpoint{3.992819in}{0.992873in}}%
\pgfpathlineto{\pgfqpoint{4.001780in}{1.135605in}}%
\pgfpathlineto{\pgfqpoint{4.010740in}{1.135605in}}%
\pgfpathlineto{\pgfqpoint{4.019701in}{1.256378in}}%
\pgfpathlineto{\pgfqpoint{4.028661in}{1.234419in}}%
\pgfpathlineto{\pgfqpoint{4.046583in}{1.124625in}}%
\pgfpathlineto{\pgfqpoint{4.055543in}{1.201481in}}%
\pgfpathlineto{\pgfqpoint{4.064504in}{1.025811in}}%
\pgfpathlineto{\pgfqpoint{4.073465in}{1.179522in}}%
\pgfpathlineto{\pgfqpoint{4.082425in}{1.124625in}}%
\pgfpathlineto{\pgfqpoint{4.091386in}{1.146584in}}%
\pgfpathlineto{\pgfqpoint{4.100347in}{1.157564in}}%
\pgfpathlineto{\pgfqpoint{4.109307in}{1.135605in}}%
\pgfpathlineto{\pgfqpoint{4.118268in}{1.069729in}}%
\pgfpathlineto{\pgfqpoint{4.127229in}{1.124625in}}%
\pgfpathlineto{\pgfqpoint{4.136189in}{1.058749in}}%
\pgfpathlineto{\pgfqpoint{4.145150in}{1.124625in}}%
\pgfpathlineto{\pgfqpoint{4.154111in}{1.058749in}}%
\pgfpathlineto{\pgfqpoint{4.163071in}{1.058749in}}%
\pgfpathlineto{\pgfqpoint{4.172032in}{1.091687in}}%
\pgfpathlineto{\pgfqpoint{4.180992in}{1.201481in}}%
\pgfpathlineto{\pgfqpoint{4.189953in}{1.124625in}}%
\pgfpathlineto{\pgfqpoint{4.198914in}{1.036790in}}%
\pgfpathlineto{\pgfqpoint{4.207874in}{0.992873in}}%
\pgfpathlineto{\pgfqpoint{4.216835in}{1.069729in}}%
\pgfpathlineto{\pgfqpoint{4.225796in}{1.311275in}}%
\pgfpathlineto{\pgfqpoint{4.234756in}{1.388130in}}%
\pgfpathlineto{\pgfqpoint{4.243717in}{1.443027in}}%
\pgfpathlineto{\pgfqpoint{4.252678in}{1.157564in}}%
\pgfpathlineto{\pgfqpoint{4.261638in}{1.113646in}}%
\pgfpathlineto{\pgfqpoint{4.270599in}{1.091687in}}%
\pgfpathlineto{\pgfqpoint{4.279560in}{1.080708in}}%
\pgfpathlineto{\pgfqpoint{4.288520in}{1.014832in}}%
\pgfpathlineto{\pgfqpoint{4.297481in}{1.156877in}}%
\pgfpathlineto{\pgfqpoint{4.306441in}{0.992187in}}%
\pgfpathlineto{\pgfqpoint{4.315402in}{1.112960in}}%
\pgfpathlineto{\pgfqpoint{4.324363in}{1.047084in}}%
\pgfpathlineto{\pgfqpoint{4.333323in}{1.047084in}}%
\pgfpathlineto{\pgfqpoint{4.351245in}{1.200795in}}%
\pgfpathlineto{\pgfqpoint{4.360205in}{1.134919in}}%
\pgfpathlineto{\pgfqpoint{4.369166in}{1.047084in}}%
\pgfpathlineto{\pgfqpoint{4.378127in}{1.112960in}}%
\pgfpathlineto{\pgfqpoint{4.387087in}{1.069042in}}%
\pgfpathlineto{\pgfqpoint{4.396048in}{0.981207in}}%
\pgfpathlineto{\pgfqpoint{4.405009in}{1.025125in}}%
\pgfpathlineto{\pgfqpoint{4.413969in}{1.080022in}}%
\pgfpathlineto{\pgfqpoint{4.422930in}{1.189815in}}%
\pgfpathlineto{\pgfqpoint{4.431891in}{1.145898in}}%
\pgfpathlineto{\pgfqpoint{4.440851in}{1.134919in}}%
\pgfpathlineto{\pgfqpoint{4.449812in}{1.058063in}}%
\pgfpathlineto{\pgfqpoint{4.458772in}{1.091001in}}%
\pgfpathlineto{\pgfqpoint{4.467733in}{1.058063in}}%
\pgfpathlineto{\pgfqpoint{4.476694in}{1.047084in}}%
\pgfpathlineto{\pgfqpoint{4.485654in}{1.189815in}}%
\pgfpathlineto{\pgfqpoint{4.494615in}{1.244712in}}%
\pgfpathlineto{\pgfqpoint{4.503576in}{1.211774in}}%
\pgfpathlineto{\pgfqpoint{4.512536in}{1.145898in}}%
\pgfpathlineto{\pgfqpoint{4.521497in}{1.091001in}}%
\pgfpathlineto{\pgfqpoint{4.530458in}{1.266671in}}%
\pgfpathlineto{\pgfqpoint{4.539418in}{1.178836in}}%
\pgfpathlineto{\pgfqpoint{4.548379in}{1.167857in}}%
\pgfpathlineto{\pgfqpoint{4.557340in}{1.080022in}}%
\pgfpathlineto{\pgfqpoint{4.566300in}{1.178836in}}%
\pgfpathlineto{\pgfqpoint{4.575261in}{1.058063in}}%
\pgfpathlineto{\pgfqpoint{4.584222in}{1.189815in}}%
\pgfpathlineto{\pgfqpoint{4.602143in}{1.189815in}}%
\pgfpathlineto{\pgfqpoint{4.611103in}{1.276964in}}%
\pgfpathlineto{\pgfqpoint{4.620064in}{1.123253in}}%
\pgfpathlineto{\pgfqpoint{4.629025in}{1.156191in}}%
\pgfpathlineto{\pgfqpoint{4.637985in}{1.101294in}}%
\pgfpathlineto{\pgfqpoint{4.646946in}{1.265985in}}%
\pgfpathlineto{\pgfqpoint{4.655907in}{1.298923in}}%
\pgfpathlineto{\pgfqpoint{4.664867in}{1.309902in}}%
\pgfpathlineto{\pgfqpoint{4.673828in}{1.013459in}}%
\pgfpathlineto{\pgfqpoint{4.691749in}{1.167171in}}%
\pgfpathlineto{\pgfqpoint{4.700710in}{1.156191in}}%
\pgfpathlineto{\pgfqpoint{4.709671in}{1.079336in}}%
\pgfpathlineto{\pgfqpoint{4.718631in}{1.013459in}}%
\pgfpathlineto{\pgfqpoint{4.727592in}{1.222067in}}%
\pgfpathlineto{\pgfqpoint{4.736552in}{1.156191in}}%
\pgfpathlineto{\pgfqpoint{4.745513in}{1.200109in}}%
\pgfpathlineto{\pgfqpoint{4.754474in}{1.156191in}}%
\pgfpathlineto{\pgfqpoint{4.763434in}{1.013459in}}%
\pgfpathlineto{\pgfqpoint{4.772395in}{1.145212in}}%
\pgfpathlineto{\pgfqpoint{4.781356in}{1.112274in}}%
\pgfpathlineto{\pgfqpoint{4.790316in}{1.189129in}}%
\pgfpathlineto{\pgfqpoint{4.799277in}{1.145212in}}%
\pgfpathlineto{\pgfqpoint{4.808238in}{1.320882in}}%
\pgfpathlineto{\pgfqpoint{4.817198in}{1.156191in}}%
\pgfpathlineto{\pgfqpoint{4.826159in}{1.123253in}}%
\pgfpathlineto{\pgfqpoint{4.835120in}{1.156191in}}%
\pgfpathlineto{\pgfqpoint{4.844080in}{1.112274in}}%
\pgfpathlineto{\pgfqpoint{4.853041in}{1.298923in}}%
\pgfpathlineto{\pgfqpoint{4.862002in}{1.123253in}}%
\pgfpathlineto{\pgfqpoint{4.870962in}{1.244026in}}%
\pgfpathlineto{\pgfqpoint{4.879923in}{1.233047in}}%
\pgfpathlineto{\pgfqpoint{4.888883in}{1.046397in}}%
\pgfpathlineto{\pgfqpoint{4.897844in}{1.024439in}}%
\pgfpathlineto{\pgfqpoint{4.906805in}{1.057377in}}%
\pgfpathlineto{\pgfqpoint{4.915765in}{1.046397in}}%
\pgfpathlineto{\pgfqpoint{4.924726in}{1.144526in}}%
\pgfpathlineto{\pgfqpoint{4.933687in}{1.122567in}}%
\pgfpathlineto{\pgfqpoint{4.942647in}{1.166484in}}%
\pgfpathlineto{\pgfqpoint{4.951608in}{1.111587in}}%
\pgfpathlineto{\pgfqpoint{4.960569in}{1.045711in}}%
\pgfpathlineto{\pgfqpoint{4.969529in}{1.199422in}}%
\pgfpathlineto{\pgfqpoint{4.978490in}{1.012773in}}%
\pgfpathlineto{\pgfqpoint{4.987451in}{1.144526in}}%
\pgfpathlineto{\pgfqpoint{4.996411in}{1.144526in}}%
\pgfpathlineto{\pgfqpoint{5.005372in}{1.210402in}}%
\pgfpathlineto{\pgfqpoint{5.014333in}{1.023752in}}%
\pgfpathlineto{\pgfqpoint{5.023293in}{1.166484in}}%
\pgfpathlineto{\pgfqpoint{5.032254in}{1.177464in}}%
\pgfpathlineto{\pgfqpoint{5.041214in}{1.122567in}}%
\pgfpathlineto{\pgfqpoint{5.050175in}{1.166484in}}%
\pgfpathlineto{\pgfqpoint{5.059136in}{1.133546in}}%
\pgfpathlineto{\pgfqpoint{5.068096in}{1.232361in}}%
\pgfpathlineto{\pgfqpoint{5.077057in}{1.089629in}}%
\pgfpathlineto{\pgfqpoint{5.086018in}{1.221381in}}%
\pgfpathlineto{\pgfqpoint{5.094978in}{1.386072in}}%
\pgfpathlineto{\pgfqpoint{5.112900in}{1.276278in}}%
\pgfpathlineto{\pgfqpoint{5.121860in}{1.473907in}}%
\pgfpathlineto{\pgfqpoint{5.130821in}{2.077772in}}%
\pgfpathlineto{\pgfqpoint{5.139782in}{2.330298in}}%
\pgfpathlineto{\pgfqpoint{5.148742in}{2.418133in}}%
\pgfpathlineto{\pgfqpoint{5.157703in}{2.396174in}}%
\pgfpathlineto{\pgfqpoint{5.166664in}{2.253442in}}%
\pgfpathlineto{\pgfqpoint{5.175624in}{1.902102in}}%
\pgfpathlineto{\pgfqpoint{5.184585in}{1.188443in}}%
\pgfpathlineto{\pgfqpoint{5.193545in}{1.320196in}}%
\pgfpathlineto{\pgfqpoint{5.202506in}{1.320196in}}%
\pgfpathlineto{\pgfqpoint{5.211467in}{1.386072in}}%
\pgfpathlineto{\pgfqpoint{5.220427in}{1.166484in}}%
\pgfpathlineto{\pgfqpoint{5.229388in}{1.353134in}}%
\pgfpathlineto{\pgfqpoint{5.238349in}{1.176777in}}%
\pgfpathlineto{\pgfqpoint{5.247309in}{1.121881in}}%
\pgfpathlineto{\pgfqpoint{5.256270in}{1.231674in}}%
\pgfpathlineto{\pgfqpoint{5.265231in}{1.363427in}}%
\pgfpathlineto{\pgfqpoint{5.274191in}{1.396365in}}%
\pgfpathlineto{\pgfqpoint{5.283152in}{1.308530in}}%
\pgfpathlineto{\pgfqpoint{5.292113in}{1.187757in}}%
\pgfpathlineto{\pgfqpoint{5.301073in}{1.209716in}}%
\pgfpathlineto{\pgfqpoint{5.310034in}{1.165798in}}%
\pgfpathlineto{\pgfqpoint{5.318994in}{1.088942in}}%
\pgfpathlineto{\pgfqpoint{5.327955in}{1.056004in}}%
\pgfpathlineto{\pgfqpoint{5.336916in}{1.187757in}}%
\pgfpathlineto{\pgfqpoint{5.345876in}{1.023066in}}%
\pgfpathlineto{\pgfqpoint{5.354837in}{1.231674in}}%
\pgfpathlineto{\pgfqpoint{5.363798in}{1.231674in}}%
\pgfpathlineto{\pgfqpoint{5.372758in}{1.198736in}}%
\pgfpathlineto{\pgfqpoint{5.381719in}{1.088942in}}%
\pgfpathlineto{\pgfqpoint{5.390680in}{1.264612in}}%
\pgfpathlineto{\pgfqpoint{5.399640in}{1.308530in}}%
\pgfpathlineto{\pgfqpoint{5.408601in}{1.374406in}}%
\pgfpathlineto{\pgfqpoint{5.417562in}{1.286571in}}%
\pgfpathlineto{\pgfqpoint{5.426522in}{1.319509in}}%
\pgfpathlineto{\pgfqpoint{5.435483in}{1.385386in}}%
\pgfpathlineto{\pgfqpoint{5.444444in}{1.330489in}}%
\pgfpathlineto{\pgfqpoint{5.453404in}{1.242654in}}%
\pgfpathlineto{\pgfqpoint{5.462365in}{1.495179in}}%
\pgfpathlineto{\pgfqpoint{5.471325in}{1.495179in}}%
\pgfpathlineto{\pgfqpoint{5.480286in}{1.363427in}}%
\pgfpathlineto{\pgfqpoint{5.489247in}{1.451262in}}%
\pgfpathlineto{\pgfqpoint{5.498207in}{1.495179in}}%
\pgfpathlineto{\pgfqpoint{5.507168in}{1.550076in}}%
\pgfpathlineto{\pgfqpoint{5.516129in}{1.473221in}}%
\pgfpathlineto{\pgfqpoint{5.534050in}{1.385386in}}%
\pgfpathlineto{\pgfqpoint{5.543011in}{1.396365in}}%
\pgfpathlineto{\pgfqpoint{5.551971in}{1.461555in}}%
\pgfpathlineto{\pgfqpoint{5.560932in}{1.373720in}}%
\pgfpathlineto{\pgfqpoint{5.569893in}{1.494493in}}%
\pgfpathlineto{\pgfqpoint{5.578853in}{1.417638in}}%
\pgfpathlineto{\pgfqpoint{5.587814in}{1.406658in}}%
\pgfpathlineto{\pgfqpoint{5.596775in}{1.549390in}}%
\pgfpathlineto{\pgfqpoint{5.605735in}{1.461555in}}%
\pgfpathlineto{\pgfqpoint{5.614696in}{1.461555in}}%
\pgfpathlineto{\pgfqpoint{5.623656in}{1.681143in}}%
\pgfpathlineto{\pgfqpoint{5.632617in}{1.494493in}}%
\pgfpathlineto{\pgfqpoint{5.641578in}{1.395679in}}%
\pgfpathlineto{\pgfqpoint{5.650538in}{1.373720in}}%
\pgfpathlineto{\pgfqpoint{5.659499in}{1.198050in}}%
\pgfpathlineto{\pgfqpoint{5.668460in}{1.241968in}}%
\pgfpathlineto{\pgfqpoint{5.677420in}{1.384699in}}%
\pgfpathlineto{\pgfqpoint{5.686381in}{1.461555in}}%
\pgfpathlineto{\pgfqpoint{5.695342in}{1.461555in}}%
\pgfpathlineto{\pgfqpoint{5.704302in}{1.615266in}}%
\pgfpathlineto{\pgfqpoint{5.713263in}{1.582328in}}%
\pgfpathlineto{\pgfqpoint{5.722224in}{1.494493in}}%
\pgfpathlineto{\pgfqpoint{5.731184in}{1.560369in}}%
\pgfpathlineto{\pgfqpoint{5.740145in}{1.604287in}}%
\pgfpathlineto{\pgfqpoint{5.749105in}{1.516452in}}%
\pgfpathlineto{\pgfqpoint{5.758066in}{1.406658in}}%
\pgfpathlineto{\pgfqpoint{5.767027in}{1.406658in}}%
\pgfpathlineto{\pgfqpoint{5.775987in}{1.516452in}}%
\pgfpathlineto{\pgfqpoint{5.784948in}{1.494493in}}%
\pgfpathlineto{\pgfqpoint{5.793909in}{1.406658in}}%
\pgfpathlineto{\pgfqpoint{5.802869in}{1.450576in}}%
\pgfpathlineto{\pgfqpoint{5.811830in}{1.571349in}}%
\pgfpathlineto{\pgfqpoint{5.820791in}{1.340782in}}%
\pgfpathlineto{\pgfqpoint{5.829751in}{1.494493in}}%
\pgfpathlineto{\pgfqpoint{5.838712in}{1.516452in}}%
\pgfpathlineto{\pgfqpoint{5.847673in}{1.560369in}}%
\pgfpathlineto{\pgfqpoint{5.856633in}{1.417638in}}%
\pgfpathlineto{\pgfqpoint{5.865594in}{1.537724in}}%
\pgfpathlineto{\pgfqpoint{5.874555in}{1.581642in}}%
\pgfpathlineto{\pgfqpoint{5.883515in}{1.614580in}}%
\pgfpathlineto{\pgfqpoint{5.892476in}{1.427931in}}%
\pgfpathlineto{\pgfqpoint{5.901436in}{1.581642in}}%
\pgfpathlineto{\pgfqpoint{5.910397in}{1.526745in}}%
\pgfpathlineto{\pgfqpoint{5.919358in}{1.592621in}}%
\pgfpathlineto{\pgfqpoint{5.928318in}{1.669477in}}%
\pgfpathlineto{\pgfqpoint{5.946240in}{1.197364in}}%
\pgfpathlineto{\pgfqpoint{5.955200in}{1.537724in}}%
\pgfpathlineto{\pgfqpoint{5.964161in}{1.570663in}}%
\pgfpathlineto{\pgfqpoint{5.973122in}{1.537724in}}%
\pgfpathlineto{\pgfqpoint{5.982082in}{1.493807in}}%
\pgfpathlineto{\pgfqpoint{5.991043in}{1.636539in}}%
\pgfpathlineto{\pgfqpoint{6.000004in}{1.460869in}}%
\pgfpathlineto{\pgfqpoint{6.008964in}{1.592621in}}%
\pgfpathlineto{\pgfqpoint{6.026886in}{1.526745in}}%
\pgfpathlineto{\pgfqpoint{6.035846in}{1.592621in}}%
\pgfpathlineto{\pgfqpoint{6.044807in}{1.394993in}}%
\pgfpathlineto{\pgfqpoint{6.053767in}{1.724374in}}%
\pgfpathlineto{\pgfqpoint{6.062728in}{1.570663in}}%
\pgfpathlineto{\pgfqpoint{6.062728in}{1.570663in}}%
\pgfusepath{stroke}%
\end{pgfscope}%
\begin{pgfscope}%
\pgfpathrectangle{\pgfqpoint{0.758026in}{0.565123in}}{\pgfqpoint{5.303239in}{2.284877in}}%
\pgfusepath{clip}%
\pgfsetrectcap%
\pgfsetroundjoin%
\pgfsetlinewidth{1.003750pt}%
\definecolor{currentstroke}{rgb}{1.000000,0.000000,0.000000}%
\pgfsetstrokecolor{currentstroke}%
\pgfsetdash{}{0pt}%
\pgfpathmoveto{\pgfqpoint{0.758026in}{1.589876in}}%
\pgfpathlineto{\pgfqpoint{0.766986in}{1.534980in}}%
\pgfpathlineto{\pgfqpoint{0.775947in}{1.348330in}}%
\pgfpathlineto{\pgfqpoint{0.784908in}{1.435479in}}%
\pgfpathlineto{\pgfqpoint{0.793868in}{1.423813in}}%
\pgfpathlineto{\pgfqpoint{0.802829in}{1.456065in}}%
\pgfpathlineto{\pgfqpoint{0.811790in}{1.357251in}}%
\pgfpathlineto{\pgfqpoint{0.820750in}{1.466359in}}%
\pgfpathlineto{\pgfqpoint{0.829711in}{1.542528in}}%
\pgfpathlineto{\pgfqpoint{0.838672in}{1.487631in}}%
\pgfpathlineto{\pgfqpoint{0.847632in}{1.388130in}}%
\pgfpathlineto{\pgfqpoint{0.856593in}{1.387444in}}%
\pgfpathlineto{\pgfqpoint{0.865553in}{1.508217in}}%
\pgfpathlineto{\pgfqpoint{0.874514in}{1.463614in}}%
\pgfpathlineto{\pgfqpoint{0.883475in}{1.375092in}}%
\pgfpathlineto{\pgfqpoint{0.892435in}{1.561742in}}%
\pgfpathlineto{\pgfqpoint{0.901396in}{1.472534in}}%
\pgfpathlineto{\pgfqpoint{0.910357in}{1.351761in}}%
\pgfpathlineto{\pgfqpoint{0.919317in}{1.439596in}}%
\pgfpathlineto{\pgfqpoint{0.928278in}{1.438910in}}%
\pgfpathlineto{\pgfqpoint{0.937239in}{1.405286in}}%
\pgfpathlineto{\pgfqpoint{0.946199in}{1.427244in}}%
\pgfpathlineto{\pgfqpoint{0.955160in}{1.470476in}}%
\pgfpathlineto{\pgfqpoint{0.964121in}{1.392934in}}%
\pgfpathlineto{\pgfqpoint{0.973081in}{1.425186in}}%
\pgfpathlineto{\pgfqpoint{0.982042in}{1.425186in}}%
\pgfpathlineto{\pgfqpoint{0.991002in}{1.446458in}}%
\pgfpathlineto{\pgfqpoint{0.999963in}{1.434793in}}%
\pgfpathlineto{\pgfqpoint{1.008924in}{1.467045in}}%
\pgfpathlineto{\pgfqpoint{1.017884in}{1.467045in}}%
\pgfpathlineto{\pgfqpoint{1.026845in}{1.444400in}}%
\pgfpathlineto{\pgfqpoint{1.035806in}{1.344899in}}%
\pgfpathlineto{\pgfqpoint{1.044766in}{1.498610in}}%
\pgfpathlineto{\pgfqpoint{1.053727in}{1.454007in}}%
\pgfpathlineto{\pgfqpoint{1.062688in}{1.508904in}}%
\pgfpathlineto{\pgfqpoint{1.071648in}{1.364799in}}%
\pgfpathlineto{\pgfqpoint{1.080609in}{1.507531in}}%
\pgfpathlineto{\pgfqpoint{1.098530in}{1.462927in}}%
\pgfpathlineto{\pgfqpoint{1.107491in}{1.451262in}}%
\pgfpathlineto{\pgfqpoint{1.116452in}{1.539097in}}%
\pgfpathlineto{\pgfqpoint{1.125412in}{1.427931in}}%
\pgfpathlineto{\pgfqpoint{1.134373in}{1.438910in}}%
\pgfpathlineto{\pgfqpoint{1.143333in}{1.405286in}}%
\pgfpathlineto{\pgfqpoint{1.152294in}{1.493121in}}%
\pgfpathlineto{\pgfqpoint{1.161255in}{1.404600in}}%
\pgfpathlineto{\pgfqpoint{1.170215in}{1.524686in}}%
\pgfpathlineto{\pgfqpoint{1.179176in}{1.480083in}}%
\pgfpathlineto{\pgfqpoint{1.188137in}{1.480083in}}%
\pgfpathlineto{\pgfqpoint{1.197097in}{1.435479in}}%
\pgfpathlineto{\pgfqpoint{1.206058in}{1.402541in}}%
\pgfpathlineto{\pgfqpoint{1.215019in}{1.489690in}}%
\pgfpathlineto{\pgfqpoint{1.223979in}{1.445086in}}%
\pgfpathlineto{\pgfqpoint{1.232940in}{1.368230in}}%
\pgfpathlineto{\pgfqpoint{1.241901in}{1.400482in}}%
\pgfpathlineto{\pgfqpoint{1.250861in}{1.465672in}}%
\pgfpathlineto{\pgfqpoint{1.259822in}{1.486945in}}%
\pgfpathlineto{\pgfqpoint{1.268783in}{1.453321in}}%
\pgfpathlineto{\pgfqpoint{1.277743in}{1.486259in}}%
\pgfpathlineto{\pgfqpoint{1.286704in}{1.496552in}}%
\pgfpathlineto{\pgfqpoint{1.295664in}{1.462927in}}%
\pgfpathlineto{\pgfqpoint{1.304625in}{1.506845in}}%
\pgfpathlineto{\pgfqpoint{1.313586in}{1.484200in}}%
\pgfpathlineto{\pgfqpoint{1.322546in}{1.517138in}}%
\pgfpathlineto{\pgfqpoint{1.331507in}{1.461555in}}%
\pgfpathlineto{\pgfqpoint{1.349428in}{1.438224in}}%
\pgfpathlineto{\pgfqpoint{1.358389in}{1.350389in}}%
\pgfpathlineto{\pgfqpoint{1.367350in}{1.503414in}}%
\pgfpathlineto{\pgfqpoint{1.376310in}{1.459496in}}%
\pgfpathlineto{\pgfqpoint{1.385271in}{1.447831in}}%
\pgfpathlineto{\pgfqpoint{1.394232in}{2.501165in}}%
\pgfpathlineto{\pgfqpoint{1.403192in}{2.501165in}}%
\pgfpathlineto{\pgfqpoint{1.412153in}{1.380582in}}%
\pgfpathlineto{\pgfqpoint{1.421113in}{1.445772in}}%
\pgfpathlineto{\pgfqpoint{1.430074in}{1.467731in}}%
\pgfpathlineto{\pgfqpoint{1.439035in}{1.499983in}}%
\pgfpathlineto{\pgfqpoint{1.447995in}{1.444400in}}%
\pgfpathlineto{\pgfqpoint{1.456956in}{1.520569in}}%
\pgfpathlineto{\pgfqpoint{1.465917in}{1.344899in}}%
\pgfpathlineto{\pgfqpoint{1.474877in}{1.333234in}}%
\pgfpathlineto{\pgfqpoint{1.483838in}{2.573217in}}%
\pgfpathlineto{\pgfqpoint{1.492799in}{2.628114in}}%
\pgfpathlineto{\pgfqpoint{1.510720in}{1.451948in}}%
\pgfpathlineto{\pgfqpoint{1.519681in}{1.451948in}}%
\pgfpathlineto{\pgfqpoint{1.528641in}{1.495179in}}%
\pgfpathlineto{\pgfqpoint{1.537602in}{1.494493in}}%
\pgfpathlineto{\pgfqpoint{1.546563in}{1.471848in}}%
\pgfpathlineto{\pgfqpoint{1.555523in}{1.471848in}}%
\pgfpathlineto{\pgfqpoint{1.564484in}{1.350389in}}%
\pgfpathlineto{\pgfqpoint{1.573444in}{2.304908in}}%
\pgfpathlineto{\pgfqpoint{1.582405in}{2.513516in}}%
\pgfpathlineto{\pgfqpoint{1.591366in}{2.184135in}}%
\pgfpathlineto{\pgfqpoint{1.600326in}{1.534980in}}%
\pgfpathlineto{\pgfqpoint{1.609287in}{1.502041in}}%
\pgfpathlineto{\pgfqpoint{1.618248in}{1.336665in}}%
\pgfpathlineto{\pgfqpoint{1.627208in}{1.259123in}}%
\pgfpathlineto{\pgfqpoint{1.636169in}{1.412834in}}%
\pgfpathlineto{\pgfqpoint{1.645130in}{1.269416in}}%
\pgfpathlineto{\pgfqpoint{1.654090in}{1.291375in}}%
\pgfpathlineto{\pgfqpoint{1.663051in}{1.191874in}}%
\pgfpathlineto{\pgfqpoint{1.672012in}{1.388817in}}%
\pgfpathlineto{\pgfqpoint{1.680972in}{1.333920in}}%
\pgfpathlineto{\pgfqpoint{1.689933in}{1.344213in}}%
\pgfpathlineto{\pgfqpoint{1.698894in}{1.442341in}}%
\pgfpathlineto{\pgfqpoint{1.707854in}{1.200795in}}%
\pgfpathlineto{\pgfqpoint{1.716815in}{1.134232in}}%
\pgfpathlineto{\pgfqpoint{1.725775in}{1.133546in}}%
\pgfpathlineto{\pgfqpoint{1.734736in}{1.154819in}}%
\pgfpathlineto{\pgfqpoint{1.743697in}{1.154819in}}%
\pgfpathlineto{\pgfqpoint{1.752657in}{1.132174in}}%
\pgfpathlineto{\pgfqpoint{1.770579in}{1.130801in}}%
\pgfpathlineto{\pgfqpoint{1.779539in}{1.086884in}}%
\pgfpathlineto{\pgfqpoint{1.788500in}{1.195991in}}%
\pgfpathlineto{\pgfqpoint{1.797461in}{1.195305in}}%
\pgfpathlineto{\pgfqpoint{1.806421in}{1.261181in}}%
\pgfpathlineto{\pgfqpoint{1.815382in}{1.183640in}}%
\pgfpathlineto{\pgfqpoint{1.824343in}{1.150015in}}%
\pgfpathlineto{\pgfqpoint{1.833303in}{1.128057in}}%
\pgfpathlineto{\pgfqpoint{1.842264in}{1.160308in}}%
\pgfpathlineto{\pgfqpoint{1.851225in}{1.127370in}}%
\pgfpathlineto{\pgfqpoint{1.860185in}{1.082767in}}%
\pgfpathlineto{\pgfqpoint{1.869146in}{1.148643in}}%
\pgfpathlineto{\pgfqpoint{1.878106in}{1.180209in}}%
\pgfpathlineto{\pgfqpoint{1.887067in}{1.180209in}}%
\pgfpathlineto{\pgfqpoint{1.896028in}{1.124625in}}%
\pgfpathlineto{\pgfqpoint{1.904988in}{1.080022in}}%
\pgfpathlineto{\pgfqpoint{1.913949in}{1.123939in}}%
\pgfpathlineto{\pgfqpoint{1.922910in}{1.090315in}}%
\pgfpathlineto{\pgfqpoint{1.931870in}{1.111587in}}%
\pgfpathlineto{\pgfqpoint{1.940831in}{1.122567in}}%
\pgfpathlineto{\pgfqpoint{1.958752in}{1.121194in}}%
\pgfpathlineto{\pgfqpoint{1.967713in}{1.077277in}}%
\pgfpathlineto{\pgfqpoint{1.976674in}{1.120508in}}%
\pgfpathlineto{\pgfqpoint{1.985634in}{1.229616in}}%
\pgfpathlineto{\pgfqpoint{1.994595in}{1.152074in}}%
\pgfpathlineto{\pgfqpoint{2.003555in}{1.185012in}}%
\pgfpathlineto{\pgfqpoint{2.012516in}{1.085511in}}%
\pgfpathlineto{\pgfqpoint{2.021477in}{1.107470in}}%
\pgfpathlineto{\pgfqpoint{2.030437in}{1.073846in}}%
\pgfpathlineto{\pgfqpoint{2.039398in}{1.117763in}}%
\pgfpathlineto{\pgfqpoint{2.048359in}{1.073160in}}%
\pgfpathlineto{\pgfqpoint{2.057319in}{1.138350in}}%
\pgfpathlineto{\pgfqpoint{2.066280in}{1.148643in}}%
\pgfpathlineto{\pgfqpoint{2.075241in}{1.136977in}}%
\pgfpathlineto{\pgfqpoint{2.084201in}{1.071101in}}%
\pgfpathlineto{\pgfqpoint{2.093162in}{1.081394in}}%
\pgfpathlineto{\pgfqpoint{2.102123in}{1.080708in}}%
\pgfpathlineto{\pgfqpoint{2.111083in}{1.135605in}}%
\pgfpathlineto{\pgfqpoint{2.120044in}{1.101980in}}%
\pgfpathlineto{\pgfqpoint{2.129005in}{1.145898in}}%
\pgfpathlineto{\pgfqpoint{2.137965in}{1.112274in}}%
\pgfpathlineto{\pgfqpoint{2.146926in}{1.111587in}}%
\pgfpathlineto{\pgfqpoint{2.155886in}{0.979149in}}%
\pgfpathlineto{\pgfqpoint{2.164847in}{1.066984in}}%
\pgfpathlineto{\pgfqpoint{2.173808in}{1.209029in}}%
\pgfpathlineto{\pgfqpoint{2.182768in}{1.142467in}}%
\pgfpathlineto{\pgfqpoint{2.191729in}{1.086884in}}%
\pgfpathlineto{\pgfqpoint{2.200690in}{1.108843in}}%
\pgfpathlineto{\pgfqpoint{2.209650in}{1.218636in}}%
\pgfpathlineto{\pgfqpoint{2.218611in}{1.206971in}}%
\pgfpathlineto{\pgfqpoint{2.227572in}{1.283140in}}%
\pgfpathlineto{\pgfqpoint{2.236532in}{1.117763in}}%
\pgfpathlineto{\pgfqpoint{2.245493in}{0.996990in}}%
\pgfpathlineto{\pgfqpoint{2.254454in}{1.040221in}}%
\pgfpathlineto{\pgfqpoint{2.263414in}{1.204226in}}%
\pgfpathlineto{\pgfqpoint{2.272375in}{1.028556in}}%
\pgfpathlineto{\pgfqpoint{2.281336in}{1.060808in}}%
\pgfpathlineto{\pgfqpoint{2.290296in}{1.049142in}}%
\pgfpathlineto{\pgfqpoint{2.299257in}{0.993559in}}%
\pgfpathlineto{\pgfqpoint{2.308217in}{1.213147in}}%
\pgfpathlineto{\pgfqpoint{2.326139in}{1.124625in}}%
\pgfpathlineto{\pgfqpoint{2.335099in}{1.003166in}}%
\pgfpathlineto{\pgfqpoint{2.344060in}{1.112274in}}%
\pgfpathlineto{\pgfqpoint{2.353021in}{0.968856in}}%
\pgfpathlineto{\pgfqpoint{2.361981in}{1.067670in}}%
\pgfpathlineto{\pgfqpoint{2.370942in}{0.968169in}}%
\pgfpathlineto{\pgfqpoint{2.379903in}{0.989442in}}%
\pgfpathlineto{\pgfqpoint{2.388863in}{1.099236in}}%
\pgfpathlineto{\pgfqpoint{2.397824in}{0.955818in}}%
\pgfpathlineto{\pgfqpoint{2.406785in}{1.174719in}}%
\pgfpathlineto{\pgfqpoint{2.415745in}{0.999049in}}%
\pgfpathlineto{\pgfqpoint{2.424706in}{1.119822in}}%
\pgfpathlineto{\pgfqpoint{2.433666in}{1.097863in}}%
\pgfpathlineto{\pgfqpoint{2.442627in}{0.988069in}}%
\pgfpathlineto{\pgfqpoint{2.451588in}{0.933173in}}%
\pgfpathlineto{\pgfqpoint{2.460548in}{0.988069in}}%
\pgfpathlineto{\pgfqpoint{2.469509in}{1.021008in}}%
\pgfpathlineto{\pgfqpoint{2.478470in}{0.955131in}}%
\pgfpathlineto{\pgfqpoint{2.487430in}{0.999049in}}%
\pgfpathlineto{\pgfqpoint{2.496391in}{1.108843in}}%
\pgfpathlineto{\pgfqpoint{2.505352in}{1.042966in}}%
\pgfpathlineto{\pgfqpoint{2.514312in}{0.999049in}}%
\pgfpathlineto{\pgfqpoint{2.523273in}{0.999049in}}%
\pgfpathlineto{\pgfqpoint{2.532234in}{1.021008in}}%
\pgfpathlineto{\pgfqpoint{2.541194in}{1.075904in}}%
\pgfpathlineto{\pgfqpoint{2.550155in}{0.977090in}}%
\pgfpathlineto{\pgfqpoint{2.559116in}{0.999049in}}%
\pgfpathlineto{\pgfqpoint{2.568076in}{1.064925in}}%
\pgfpathlineto{\pgfqpoint{2.577037in}{1.042966in}}%
\pgfpathlineto{\pgfqpoint{2.585997in}{1.053946in}}%
\pgfpathlineto{\pgfqpoint{2.594958in}{1.021008in}}%
\pgfpathlineto{\pgfqpoint{2.603919in}{1.152760in}}%
\pgfpathlineto{\pgfqpoint{2.612879in}{0.988069in}}%
\pgfpathlineto{\pgfqpoint{2.621840in}{0.988069in}}%
\pgfpathlineto{\pgfqpoint{2.630801in}{1.031987in}}%
\pgfpathlineto{\pgfqpoint{2.639761in}{0.933173in}}%
\pgfpathlineto{\pgfqpoint{2.648722in}{0.933173in}}%
\pgfpathlineto{\pgfqpoint{2.657683in}{0.966111in}}%
\pgfpathlineto{\pgfqpoint{2.666643in}{0.988069in}}%
\pgfpathlineto{\pgfqpoint{2.675604in}{0.922193in}}%
\pgfpathlineto{\pgfqpoint{2.684565in}{0.988069in}}%
\pgfpathlineto{\pgfqpoint{2.693525in}{0.922193in}}%
\pgfpathlineto{\pgfqpoint{2.702486in}{0.988069in}}%
\pgfpathlineto{\pgfqpoint{2.711447in}{0.933173in}}%
\pgfpathlineto{\pgfqpoint{2.720407in}{0.933173in}}%
\pgfpathlineto{\pgfqpoint{2.729368in}{0.866610in}}%
\pgfpathlineto{\pgfqpoint{2.738328in}{0.888569in}}%
\pgfpathlineto{\pgfqpoint{2.747289in}{0.855631in}}%
\pgfpathlineto{\pgfqpoint{2.756250in}{0.899548in}}%
\pgfpathlineto{\pgfqpoint{2.765210in}{0.888569in}}%
\pgfpathlineto{\pgfqpoint{2.774171in}{0.987383in}}%
\pgfpathlineto{\pgfqpoint{2.783132in}{0.965425in}}%
\pgfpathlineto{\pgfqpoint{2.792092in}{1.031301in}}%
\pgfpathlineto{\pgfqpoint{2.801053in}{0.899548in}}%
\pgfpathlineto{\pgfqpoint{2.810014in}{0.932486in}}%
\pgfpathlineto{\pgfqpoint{2.818974in}{1.009342in}}%
\pgfpathlineto{\pgfqpoint{2.827935in}{0.822693in}}%
\pgfpathlineto{\pgfqpoint{2.836896in}{1.009342in}}%
\pgfpathlineto{\pgfqpoint{2.845856in}{0.932486in}}%
\pgfpathlineto{\pgfqpoint{2.854817in}{0.987383in}}%
\pgfpathlineto{\pgfqpoint{2.863777in}{0.921507in}}%
\pgfpathlineto{\pgfqpoint{2.872738in}{1.086198in}}%
\pgfpathlineto{\pgfqpoint{2.881699in}{0.888569in}}%
\pgfpathlineto{\pgfqpoint{2.890659in}{0.932486in}}%
\pgfpathlineto{\pgfqpoint{2.899620in}{0.866610in}}%
\pgfpathlineto{\pgfqpoint{2.908581in}{0.998363in}}%
\pgfpathlineto{\pgfqpoint{2.917541in}{0.888569in}}%
\pgfpathlineto{\pgfqpoint{2.926502in}{0.899548in}}%
\pgfpathlineto{\pgfqpoint{2.935463in}{0.855631in}}%
\pgfpathlineto{\pgfqpoint{2.944423in}{0.888569in}}%
\pgfpathlineto{\pgfqpoint{2.953384in}{0.888569in}}%
\pgfpathlineto{\pgfqpoint{2.962345in}{0.899548in}}%
\pgfpathlineto{\pgfqpoint{2.971305in}{0.855631in}}%
\pgfpathlineto{\pgfqpoint{2.980266in}{0.855631in}}%
\pgfpathlineto{\pgfqpoint{2.989227in}{0.899548in}}%
\pgfpathlineto{\pgfqpoint{2.998187in}{0.866610in}}%
\pgfpathlineto{\pgfqpoint{3.007148in}{0.987383in}}%
\pgfpathlineto{\pgfqpoint{3.025069in}{0.921507in}}%
\pgfpathlineto{\pgfqpoint{3.034030in}{0.986697in}}%
\pgfpathlineto{\pgfqpoint{3.042990in}{0.931800in}}%
\pgfpathlineto{\pgfqpoint{3.051951in}{0.898862in}}%
\pgfpathlineto{\pgfqpoint{3.060912in}{0.986697in}}%
\pgfpathlineto{\pgfqpoint{3.069872in}{0.942780in}}%
\pgfpathlineto{\pgfqpoint{3.078833in}{0.920821in}}%
\pgfpathlineto{\pgfqpoint{3.087794in}{0.931800in}}%
\pgfpathlineto{\pgfqpoint{3.096754in}{0.854945in}}%
\pgfpathlineto{\pgfqpoint{3.105715in}{0.854945in}}%
\pgfpathlineto{\pgfqpoint{3.114676in}{0.931800in}}%
\pgfpathlineto{\pgfqpoint{3.123636in}{0.898862in}}%
\pgfpathlineto{\pgfqpoint{3.132597in}{1.030615in}}%
\pgfpathlineto{\pgfqpoint{3.141558in}{1.063553in}}%
\pgfpathlineto{\pgfqpoint{3.150518in}{0.920821in}}%
\pgfpathlineto{\pgfqpoint{3.159479in}{1.019635in}}%
\pgfpathlineto{\pgfqpoint{3.168439in}{1.162367in}}%
\pgfpathlineto{\pgfqpoint{3.177400in}{1.107470in}}%
\pgfpathlineto{\pgfqpoint{3.186361in}{1.019635in}}%
\pgfpathlineto{\pgfqpoint{3.195321in}{1.063553in}}%
\pgfpathlineto{\pgfqpoint{3.204282in}{0.997676in}}%
\pgfpathlineto{\pgfqpoint{3.213243in}{1.030615in}}%
\pgfpathlineto{\pgfqpoint{3.222203in}{2.194428in}}%
\pgfpathlineto{\pgfqpoint{3.231164in}{2.293243in}}%
\pgfpathlineto{\pgfqpoint{3.240125in}{1.788191in}}%
\pgfpathlineto{\pgfqpoint{3.249085in}{1.107470in}}%
\pgfpathlineto{\pgfqpoint{3.258046in}{1.239223in}}%
\pgfpathlineto{\pgfqpoint{3.267007in}{1.305099in}}%
\pgfpathlineto{\pgfqpoint{3.275967in}{1.217264in}}%
\pgfpathlineto{\pgfqpoint{3.284928in}{1.195305in}}%
\pgfpathlineto{\pgfqpoint{3.293889in}{1.184326in}}%
\pgfpathlineto{\pgfqpoint{3.302849in}{1.392934in}}%
\pgfpathlineto{\pgfqpoint{3.311810in}{2.161490in}}%
\pgfpathlineto{\pgfqpoint{3.320770in}{2.424995in}}%
\pgfpathlineto{\pgfqpoint{3.329731in}{2.172470in}}%
\pgfpathlineto{\pgfqpoint{3.338692in}{1.316078in}}%
\pgfpathlineto{\pgfqpoint{3.347652in}{1.305099in}}%
\pgfpathlineto{\pgfqpoint{3.356613in}{1.227557in}}%
\pgfpathlineto{\pgfqpoint{3.365574in}{1.183640in}}%
\pgfpathlineto{\pgfqpoint{3.374534in}{1.293433in}}%
\pgfpathlineto{\pgfqpoint{3.383495in}{1.238536in}}%
\pgfpathlineto{\pgfqpoint{3.392456in}{1.084825in}}%
\pgfpathlineto{\pgfqpoint{3.401416in}{1.743588in}}%
\pgfpathlineto{\pgfqpoint{3.410377in}{2.292557in}}%
\pgfpathlineto{\pgfqpoint{3.419338in}{2.193742in}}%
\pgfpathlineto{\pgfqpoint{3.428298in}{1.315392in}}%
\pgfpathlineto{\pgfqpoint{3.437259in}{1.425186in}}%
\pgfpathlineto{\pgfqpoint{3.446219in}{1.139722in}}%
\pgfpathlineto{\pgfqpoint{3.455180in}{1.051887in}}%
\pgfpathlineto{\pgfqpoint{3.464141in}{1.139722in}}%
\pgfpathlineto{\pgfqpoint{3.473101in}{1.150701in}}%
\pgfpathlineto{\pgfqpoint{3.482062in}{1.029928in}}%
\pgfpathlineto{\pgfqpoint{3.491023in}{0.964052in}}%
\pgfpathlineto{\pgfqpoint{3.499983in}{1.051887in}}%
\pgfpathlineto{\pgfqpoint{3.508944in}{0.986011in}}%
\pgfpathlineto{\pgfqpoint{3.517905in}{1.029928in}}%
\pgfpathlineto{\pgfqpoint{3.526865in}{1.238536in}}%
\pgfpathlineto{\pgfqpoint{3.535826in}{0.953073in}}%
\pgfpathlineto{\pgfqpoint{3.544787in}{0.920135in}}%
\pgfpathlineto{\pgfqpoint{3.553747in}{0.931114in}}%
\pgfpathlineto{\pgfqpoint{3.562708in}{0.854258in}}%
\pgfpathlineto{\pgfqpoint{3.580629in}{0.920135in}}%
\pgfpathlineto{\pgfqpoint{3.589590in}{0.887196in}}%
\pgfpathlineto{\pgfqpoint{3.598550in}{0.898176in}}%
\pgfpathlineto{\pgfqpoint{3.607511in}{0.854258in}}%
\pgfpathlineto{\pgfqpoint{3.616472in}{0.887196in}}%
\pgfpathlineto{\pgfqpoint{3.625432in}{0.887196in}}%
\pgfpathlineto{\pgfqpoint{3.634393in}{0.931114in}}%
\pgfpathlineto{\pgfqpoint{3.643354in}{0.854258in}}%
\pgfpathlineto{\pgfqpoint{3.652314in}{0.854258in}}%
\pgfpathlineto{\pgfqpoint{3.661275in}{0.941407in}}%
\pgfpathlineto{\pgfqpoint{3.670236in}{0.897490in}}%
\pgfpathlineto{\pgfqpoint{3.679196in}{0.864551in}}%
\pgfpathlineto{\pgfqpoint{3.688157in}{0.820634in}}%
\pgfpathlineto{\pgfqpoint{3.697118in}{0.864551in}}%
\pgfpathlineto{\pgfqpoint{3.706078in}{0.864551in}}%
\pgfpathlineto{\pgfqpoint{3.715039in}{0.886510in}}%
\pgfpathlineto{\pgfqpoint{3.724000in}{0.952386in}}%
\pgfpathlineto{\pgfqpoint{3.732960in}{1.029242in}}%
\pgfpathlineto{\pgfqpoint{3.741921in}{0.897490in}}%
\pgfpathlineto{\pgfqpoint{3.750881in}{0.996304in}}%
\pgfpathlineto{\pgfqpoint{3.759842in}{0.952386in}}%
\pgfpathlineto{\pgfqpoint{3.768803in}{0.897490in}}%
\pgfpathlineto{\pgfqpoint{3.795685in}{0.897490in}}%
\pgfpathlineto{\pgfqpoint{3.804645in}{0.919448in}}%
\pgfpathlineto{\pgfqpoint{3.813606in}{0.864551in}}%
\pgfpathlineto{\pgfqpoint{3.822567in}{0.919448in}}%
\pgfpathlineto{\pgfqpoint{3.831527in}{0.930428in}}%
\pgfpathlineto{\pgfqpoint{3.840488in}{0.930428in}}%
\pgfpathlineto{\pgfqpoint{3.849449in}{0.897490in}}%
\pgfpathlineto{\pgfqpoint{3.858409in}{1.018263in}}%
\pgfpathlineto{\pgfqpoint{3.867370in}{0.952386in}}%
\pgfpathlineto{\pgfqpoint{3.876330in}{0.963366in}}%
\pgfpathlineto{\pgfqpoint{3.885291in}{0.919448in}}%
\pgfpathlineto{\pgfqpoint{3.894252in}{0.963366in}}%
\pgfpathlineto{\pgfqpoint{3.903212in}{0.919448in}}%
\pgfpathlineto{\pgfqpoint{3.912173in}{0.952386in}}%
\pgfpathlineto{\pgfqpoint{3.921134in}{1.018263in}}%
\pgfpathlineto{\pgfqpoint{3.930094in}{0.963366in}}%
\pgfpathlineto{\pgfqpoint{3.939055in}{0.930428in}}%
\pgfpathlineto{\pgfqpoint{3.948016in}{0.853572in}}%
\pgfpathlineto{\pgfqpoint{3.956976in}{0.897490in}}%
\pgfpathlineto{\pgfqpoint{3.965937in}{0.963366in}}%
\pgfpathlineto{\pgfqpoint{3.974898in}{0.984638in}}%
\pgfpathlineto{\pgfqpoint{3.983858in}{0.929742in}}%
\pgfpathlineto{\pgfqpoint{3.992819in}{0.984638in}}%
\pgfpathlineto{\pgfqpoint{4.001780in}{0.918762in}}%
\pgfpathlineto{\pgfqpoint{4.010740in}{0.929742in}}%
\pgfpathlineto{\pgfqpoint{4.019701in}{0.896803in}}%
\pgfpathlineto{\pgfqpoint{4.028661in}{0.951700in}}%
\pgfpathlineto{\pgfqpoint{4.037622in}{0.995618in}}%
\pgfpathlineto{\pgfqpoint{4.046583in}{0.885824in}}%
\pgfpathlineto{\pgfqpoint{4.055543in}{0.819948in}}%
\pgfpathlineto{\pgfqpoint{4.064504in}{0.918762in}}%
\pgfpathlineto{\pgfqpoint{4.073465in}{1.072473in}}%
\pgfpathlineto{\pgfqpoint{4.082425in}{0.852886in}}%
\pgfpathlineto{\pgfqpoint{4.091386in}{0.951700in}}%
\pgfpathlineto{\pgfqpoint{4.109307in}{0.929742in}}%
\pgfpathlineto{\pgfqpoint{4.118268in}{0.852886in}}%
\pgfpathlineto{\pgfqpoint{4.127229in}{0.951700in}}%
\pgfpathlineto{\pgfqpoint{4.136189in}{0.940721in}}%
\pgfpathlineto{\pgfqpoint{4.145150in}{0.940721in}}%
\pgfpathlineto{\pgfqpoint{4.154111in}{0.984638in}}%
\pgfpathlineto{\pgfqpoint{4.163071in}{0.962680in}}%
\pgfpathlineto{\pgfqpoint{4.172032in}{0.984638in}}%
\pgfpathlineto{\pgfqpoint{4.180992in}{0.940721in}}%
\pgfpathlineto{\pgfqpoint{4.189953in}{0.885824in}}%
\pgfpathlineto{\pgfqpoint{4.198914in}{0.918762in}}%
\pgfpathlineto{\pgfqpoint{4.207874in}{0.907783in}}%
\pgfpathlineto{\pgfqpoint{4.216835in}{0.984638in}}%
\pgfpathlineto{\pgfqpoint{4.225796in}{0.962680in}}%
\pgfpathlineto{\pgfqpoint{4.234756in}{1.017577in}}%
\pgfpathlineto{\pgfqpoint{4.243717in}{0.962680in}}%
\pgfpathlineto{\pgfqpoint{4.252678in}{0.808968in}}%
\pgfpathlineto{\pgfqpoint{4.261638in}{0.907783in}}%
\pgfpathlineto{\pgfqpoint{4.270599in}{0.896803in}}%
\pgfpathlineto{\pgfqpoint{4.279560in}{0.995618in}}%
\pgfpathlineto{\pgfqpoint{4.288520in}{0.907783in}}%
\pgfpathlineto{\pgfqpoint{4.297481in}{0.874158in}}%
\pgfpathlineto{\pgfqpoint{4.306441in}{0.808282in}}%
\pgfpathlineto{\pgfqpoint{4.315402in}{0.874158in}}%
\pgfpathlineto{\pgfqpoint{4.324363in}{0.994932in}}%
\pgfpathlineto{\pgfqpoint{4.333323in}{0.961993in}}%
\pgfpathlineto{\pgfqpoint{4.342284in}{0.907097in}}%
\pgfpathlineto{\pgfqpoint{4.351245in}{0.841220in}}%
\pgfpathlineto{\pgfqpoint{4.360205in}{0.983952in}}%
\pgfpathlineto{\pgfqpoint{4.369166in}{0.929055in}}%
\pgfpathlineto{\pgfqpoint{4.378127in}{0.841220in}}%
\pgfpathlineto{\pgfqpoint{4.387087in}{0.983952in}}%
\pgfpathlineto{\pgfqpoint{4.396048in}{0.940035in}}%
\pgfpathlineto{\pgfqpoint{4.405009in}{0.994932in}}%
\pgfpathlineto{\pgfqpoint{4.413969in}{0.961993in}}%
\pgfpathlineto{\pgfqpoint{4.422930in}{0.951014in}}%
\pgfpathlineto{\pgfqpoint{4.440851in}{0.907097in}}%
\pgfpathlineto{\pgfqpoint{4.449812in}{0.951014in}}%
\pgfpathlineto{\pgfqpoint{4.458772in}{0.841220in}}%
\pgfpathlineto{\pgfqpoint{4.467733in}{0.918076in}}%
\pgfpathlineto{\pgfqpoint{4.476694in}{0.863179in}}%
\pgfpathlineto{\pgfqpoint{4.485654in}{0.852200in}}%
\pgfpathlineto{\pgfqpoint{4.494615in}{0.885138in}}%
\pgfpathlineto{\pgfqpoint{4.503576in}{0.852200in}}%
\pgfpathlineto{\pgfqpoint{4.512536in}{0.874158in}}%
\pgfpathlineto{\pgfqpoint{4.521497in}{0.961993in}}%
\pgfpathlineto{\pgfqpoint{4.530458in}{0.874158in}}%
\pgfpathlineto{\pgfqpoint{4.539418in}{0.951014in}}%
\pgfpathlineto{\pgfqpoint{4.548379in}{0.907097in}}%
\pgfpathlineto{\pgfqpoint{4.557340in}{0.885138in}}%
\pgfpathlineto{\pgfqpoint{4.566300in}{0.907097in}}%
\pgfpathlineto{\pgfqpoint{4.584222in}{0.972973in}}%
\pgfpathlineto{\pgfqpoint{4.593182in}{0.885138in}}%
\pgfpathlineto{\pgfqpoint{4.602143in}{0.874158in}}%
\pgfpathlineto{\pgfqpoint{4.611103in}{0.950328in}}%
\pgfpathlineto{\pgfqpoint{4.620064in}{0.994245in}}%
\pgfpathlineto{\pgfqpoint{4.629025in}{1.049142in}}%
\pgfpathlineto{\pgfqpoint{4.637985in}{0.917390in}}%
\pgfpathlineto{\pgfqpoint{4.646946in}{0.994245in}}%
\pgfpathlineto{\pgfqpoint{4.655907in}{0.972287in}}%
\pgfpathlineto{\pgfqpoint{4.664867in}{0.873472in}}%
\pgfpathlineto{\pgfqpoint{4.673828in}{1.038163in}}%
\pgfpathlineto{\pgfqpoint{4.682789in}{0.972287in}}%
\pgfpathlineto{\pgfqpoint{4.691749in}{0.873472in}}%
\pgfpathlineto{\pgfqpoint{4.700710in}{0.818575in}}%
\pgfpathlineto{\pgfqpoint{4.709671in}{0.873472in}}%
\pgfpathlineto{\pgfqpoint{4.718631in}{0.939348in}}%
\pgfpathlineto{\pgfqpoint{4.727592in}{0.840534in}}%
\pgfpathlineto{\pgfqpoint{4.736552in}{0.818575in}}%
\pgfpathlineto{\pgfqpoint{4.745513in}{0.884452in}}%
\pgfpathlineto{\pgfqpoint{4.754474in}{0.851513in}}%
\pgfpathlineto{\pgfqpoint{4.763434in}{0.851513in}}%
\pgfpathlineto{\pgfqpoint{4.772395in}{0.983266in}}%
\pgfpathlineto{\pgfqpoint{4.781356in}{0.840534in}}%
\pgfpathlineto{\pgfqpoint{4.790316in}{0.950328in}}%
\pgfpathlineto{\pgfqpoint{4.799277in}{0.818575in}}%
\pgfpathlineto{\pgfqpoint{4.817198in}{0.928369in}}%
\pgfpathlineto{\pgfqpoint{4.826159in}{0.807596in}}%
\pgfpathlineto{\pgfqpoint{4.835120in}{0.884452in}}%
\pgfpathlineto{\pgfqpoint{4.844080in}{0.884452in}}%
\pgfpathlineto{\pgfqpoint{4.853041in}{0.873472in}}%
\pgfpathlineto{\pgfqpoint{4.862002in}{0.807596in}}%
\pgfpathlineto{\pgfqpoint{4.870962in}{0.884452in}}%
\pgfpathlineto{\pgfqpoint{4.879923in}{0.785637in}}%
\pgfpathlineto{\pgfqpoint{4.888883in}{0.906410in}}%
\pgfpathlineto{\pgfqpoint{4.897844in}{0.884452in}}%
\pgfpathlineto{\pgfqpoint{4.906805in}{0.884452in}}%
\pgfpathlineto{\pgfqpoint{4.915765in}{0.840534in}}%
\pgfpathlineto{\pgfqpoint{4.924726in}{0.872786in}}%
\pgfpathlineto{\pgfqpoint{4.933687in}{0.916704in}}%
\pgfpathlineto{\pgfqpoint{4.942647in}{0.839848in}}%
\pgfpathlineto{\pgfqpoint{4.951608in}{0.927683in}}%
\pgfpathlineto{\pgfqpoint{4.960569in}{0.916704in}}%
\pgfpathlineto{\pgfqpoint{4.969529in}{1.125312in}}%
\pgfpathlineto{\pgfqpoint{4.978490in}{0.905724in}}%
\pgfpathlineto{\pgfqpoint{4.987451in}{0.905724in}}%
\pgfpathlineto{\pgfqpoint{4.996411in}{1.015518in}}%
\pgfpathlineto{\pgfqpoint{5.005372in}{0.883765in}}%
\pgfpathlineto{\pgfqpoint{5.014333in}{0.938662in}}%
\pgfpathlineto{\pgfqpoint{5.023293in}{1.081394in}}%
\pgfpathlineto{\pgfqpoint{5.041214in}{0.971600in}}%
\pgfpathlineto{\pgfqpoint{5.050175in}{1.860930in}}%
\pgfpathlineto{\pgfqpoint{5.059136in}{2.201290in}}%
\pgfpathlineto{\pgfqpoint{5.068096in}{2.014641in}}%
\pgfpathlineto{\pgfqpoint{5.077057in}{0.938662in}}%
\pgfpathlineto{\pgfqpoint{5.086018in}{1.136291in}}%
\pgfpathlineto{\pgfqpoint{5.094978in}{1.125312in}}%
\pgfpathlineto{\pgfqpoint{5.103939in}{1.136291in}}%
\pgfpathlineto{\pgfqpoint{5.112900in}{1.015518in}}%
\pgfpathlineto{\pgfqpoint{5.121860in}{1.081394in}}%
\pgfpathlineto{\pgfqpoint{5.130821in}{1.048456in}}%
\pgfpathlineto{\pgfqpoint{5.148742in}{2.333043in}}%
\pgfpathlineto{\pgfqpoint{5.157703in}{2.267167in}}%
\pgfpathlineto{\pgfqpoint{5.166664in}{1.311961in}}%
\pgfpathlineto{\pgfqpoint{5.175624in}{1.125312in}}%
\pgfpathlineto{\pgfqpoint{5.184585in}{1.180209in}}%
\pgfpathlineto{\pgfqpoint{5.193545in}{1.114332in}}%
\pgfpathlineto{\pgfqpoint{5.202506in}{1.235105in}}%
\pgfpathlineto{\pgfqpoint{5.211467in}{1.070415in}}%
\pgfpathlineto{\pgfqpoint{5.220427in}{1.059435in}}%
\pgfpathlineto{\pgfqpoint{5.229388in}{0.959935in}}%
\pgfpathlineto{\pgfqpoint{5.238349in}{2.189625in}}%
\pgfpathlineto{\pgfqpoint{5.247309in}{2.189625in}}%
\pgfpathlineto{\pgfqpoint{5.256270in}{1.168543in}}%
\pgfpathlineto{\pgfqpoint{5.265231in}{1.311275in}}%
\pgfpathlineto{\pgfqpoint{5.274191in}{0.981894in}}%
\pgfpathlineto{\pgfqpoint{5.283152in}{0.916017in}}%
\pgfpathlineto{\pgfqpoint{5.292113in}{1.014832in}}%
\pgfpathlineto{\pgfqpoint{5.301073in}{1.069729in}}%
\pgfpathlineto{\pgfqpoint{5.310034in}{1.036790in}}%
\pgfpathlineto{\pgfqpoint{5.318994in}{0.883079in}}%
\pgfpathlineto{\pgfqpoint{5.327955in}{0.905038in}}%
\pgfpathlineto{\pgfqpoint{5.336916in}{0.948955in}}%
\pgfpathlineto{\pgfqpoint{5.345876in}{0.883079in}}%
\pgfpathlineto{\pgfqpoint{5.354837in}{1.113646in}}%
\pgfpathlineto{\pgfqpoint{5.363798in}{1.014832in}}%
\pgfpathlineto{\pgfqpoint{5.372758in}{0.981894in}}%
\pgfpathlineto{\pgfqpoint{5.381719in}{0.850141in}}%
\pgfpathlineto{\pgfqpoint{5.390680in}{0.894059in}}%
\pgfpathlineto{\pgfqpoint{5.399640in}{0.839162in}}%
\pgfpathlineto{\pgfqpoint{5.408601in}{0.872100in}}%
\pgfpathlineto{\pgfqpoint{5.417562in}{0.806224in}}%
\pgfpathlineto{\pgfqpoint{5.426522in}{0.806224in}}%
\pgfpathlineto{\pgfqpoint{5.435483in}{0.883079in}}%
\pgfpathlineto{\pgfqpoint{5.444444in}{0.839162in}}%
\pgfpathlineto{\pgfqpoint{5.453404in}{0.905038in}}%
\pgfpathlineto{\pgfqpoint{5.462365in}{0.850141in}}%
\pgfpathlineto{\pgfqpoint{5.471325in}{0.992873in}}%
\pgfpathlineto{\pgfqpoint{5.480286in}{0.839162in}}%
\pgfpathlineto{\pgfqpoint{5.489247in}{0.883079in}}%
\pgfpathlineto{\pgfqpoint{5.516129in}{0.817203in}}%
\pgfpathlineto{\pgfqpoint{5.525089in}{0.948955in}}%
\pgfpathlineto{\pgfqpoint{5.534050in}{0.970914in}}%
\pgfpathlineto{\pgfqpoint{5.543011in}{0.872100in}}%
\pgfpathlineto{\pgfqpoint{5.551971in}{0.805537in}}%
\pgfpathlineto{\pgfqpoint{5.560932in}{0.926310in}}%
\pgfpathlineto{\pgfqpoint{5.569893in}{0.981207in}}%
\pgfpathlineto{\pgfqpoint{5.587814in}{0.915331in}}%
\pgfpathlineto{\pgfqpoint{5.596775in}{0.838475in}}%
\pgfpathlineto{\pgfqpoint{5.605735in}{1.014145in}}%
\pgfpathlineto{\pgfqpoint{5.614696in}{0.871414in}}%
\pgfpathlineto{\pgfqpoint{5.623656in}{0.871414in}}%
\pgfpathlineto{\pgfqpoint{5.632617in}{0.783579in}}%
\pgfpathlineto{\pgfqpoint{5.641578in}{0.849455in}}%
\pgfpathlineto{\pgfqpoint{5.650538in}{0.882393in}}%
\pgfpathlineto{\pgfqpoint{5.659499in}{0.750640in}}%
\pgfpathlineto{\pgfqpoint{5.668460in}{0.772599in}}%
\pgfpathlineto{\pgfqpoint{5.677420in}{0.871414in}}%
\pgfpathlineto{\pgfqpoint{5.686381in}{0.783579in}}%
\pgfpathlineto{\pgfqpoint{5.695342in}{0.860434in}}%
\pgfpathlineto{\pgfqpoint{5.704302in}{0.838475in}}%
\pgfpathlineto{\pgfqpoint{5.713263in}{0.992187in}}%
\pgfpathlineto{\pgfqpoint{5.722224in}{0.937290in}}%
\pgfpathlineto{\pgfqpoint{5.731184in}{0.827496in}}%
\pgfpathlineto{\pgfqpoint{5.740145in}{0.970228in}}%
\pgfpathlineto{\pgfqpoint{5.749105in}{0.860434in}}%
\pgfpathlineto{\pgfqpoint{5.758066in}{0.893372in}}%
\pgfpathlineto{\pgfqpoint{5.767027in}{0.970228in}}%
\pgfpathlineto{\pgfqpoint{5.775987in}{0.860434in}}%
\pgfpathlineto{\pgfqpoint{5.784948in}{0.904352in}}%
\pgfpathlineto{\pgfqpoint{5.793909in}{0.816517in}}%
\pgfpathlineto{\pgfqpoint{5.802869in}{0.893372in}}%
\pgfpathlineto{\pgfqpoint{5.811830in}{0.827496in}}%
\pgfpathlineto{\pgfqpoint{5.820791in}{0.882393in}}%
\pgfpathlineto{\pgfqpoint{5.829751in}{0.871414in}}%
\pgfpathlineto{\pgfqpoint{5.838712in}{0.783579in}}%
\pgfpathlineto{\pgfqpoint{5.847673in}{0.937290in}}%
\pgfpathlineto{\pgfqpoint{5.856633in}{0.816517in}}%
\pgfpathlineto{\pgfqpoint{5.865594in}{0.804851in}}%
\pgfpathlineto{\pgfqpoint{5.874555in}{0.936604in}}%
\pgfpathlineto{\pgfqpoint{5.883515in}{0.848769in}}%
\pgfpathlineto{\pgfqpoint{5.901436in}{0.782892in}}%
\pgfpathlineto{\pgfqpoint{5.910397in}{0.936604in}}%
\pgfpathlineto{\pgfqpoint{5.919358in}{0.925624in}}%
\pgfpathlineto{\pgfqpoint{5.928318in}{0.815831in}}%
\pgfpathlineto{\pgfqpoint{5.937279in}{0.870727in}}%
\pgfpathlineto{\pgfqpoint{5.946240in}{0.881707in}}%
\pgfpathlineto{\pgfqpoint{5.955200in}{0.837789in}}%
\pgfpathlineto{\pgfqpoint{5.964161in}{0.881707in}}%
\pgfpathlineto{\pgfqpoint{5.973122in}{0.859748in}}%
\pgfpathlineto{\pgfqpoint{5.982082in}{0.804851in}}%
\pgfpathlineto{\pgfqpoint{5.991043in}{0.815831in}}%
\pgfpathlineto{\pgfqpoint{6.008964in}{0.793872in}}%
\pgfpathlineto{\pgfqpoint{6.017925in}{0.914645in}}%
\pgfpathlineto{\pgfqpoint{6.026886in}{0.914645in}}%
\pgfpathlineto{\pgfqpoint{6.035846in}{0.815831in}}%
\pgfpathlineto{\pgfqpoint{6.044807in}{0.870727in}}%
\pgfpathlineto{\pgfqpoint{6.053767in}{0.837789in}}%
\pgfpathlineto{\pgfqpoint{6.062728in}{1.024439in}}%
\pgfpathlineto{\pgfqpoint{6.062728in}{1.024439in}}%
\pgfusepath{stroke}%
\end{pgfscope}%
\begin{pgfscope}%
\pgfsetrectcap%
\pgfsetmiterjoin%
\pgfsetlinewidth{0.803000pt}%
\definecolor{currentstroke}{rgb}{0.000000,0.000000,0.000000}%
\pgfsetstrokecolor{currentstroke}%
\pgfsetdash{}{0pt}%
\pgfpathmoveto{\pgfqpoint{0.758026in}{0.565123in}}%
\pgfpathlineto{\pgfqpoint{0.758026in}{2.850000in}}%
\pgfusepath{stroke}%
\end{pgfscope}%
\begin{pgfscope}%
\pgfsetrectcap%
\pgfsetmiterjoin%
\pgfsetlinewidth{0.803000pt}%
\definecolor{currentstroke}{rgb}{0.000000,0.000000,0.000000}%
\pgfsetstrokecolor{currentstroke}%
\pgfsetdash{}{0pt}%
\pgfpathmoveto{\pgfqpoint{6.061265in}{0.565123in}}%
\pgfpathlineto{\pgfqpoint{6.061265in}{2.850000in}}%
\pgfusepath{stroke}%
\end{pgfscope}%
\begin{pgfscope}%
\pgfsetrectcap%
\pgfsetmiterjoin%
\pgfsetlinewidth{0.803000pt}%
\definecolor{currentstroke}{rgb}{0.000000,0.000000,0.000000}%
\pgfsetstrokecolor{currentstroke}%
\pgfsetdash{}{0pt}%
\pgfpathmoveto{\pgfqpoint{0.758026in}{0.565123in}}%
\pgfpathlineto{\pgfqpoint{6.061265in}{0.565123in}}%
\pgfusepath{stroke}%
\end{pgfscope}%
\begin{pgfscope}%
\pgfsetrectcap%
\pgfsetmiterjoin%
\pgfsetlinewidth{0.803000pt}%
\definecolor{currentstroke}{rgb}{0.000000,0.000000,0.000000}%
\pgfsetstrokecolor{currentstroke}%
\pgfsetdash{}{0pt}%
\pgfpathmoveto{\pgfqpoint{0.758026in}{2.850000in}}%
\pgfpathlineto{\pgfqpoint{6.061265in}{2.850000in}}%
\pgfusepath{stroke}%
\end{pgfscope}%
\begin{pgfscope}%
\pgfsetbuttcap%
\pgfsetmiterjoin%
\definecolor{currentfill}{rgb}{1.000000,1.000000,1.000000}%
\pgfsetfillcolor{currentfill}%
\pgfsetfillopacity{0.800000}%
\pgfsetlinewidth{1.003750pt}%
\definecolor{currentstroke}{rgb}{0.800000,0.800000,0.800000}%
\pgfsetstrokecolor{currentstroke}%
\pgfsetstrokeopacity{0.800000}%
\pgfsetdash{}{0pt}%
\pgfpathmoveto{\pgfqpoint{0.855248in}{0.634568in}}%
\pgfpathlineto{\pgfqpoint{2.035033in}{0.634568in}}%
\pgfpathquadraticcurveto{\pgfqpoint{2.062811in}{0.634568in}}{\pgfqpoint{2.062811in}{0.662346in}}%
\pgfpathlineto{\pgfqpoint{2.062811in}{1.035802in}}%
\pgfpathquadraticcurveto{\pgfqpoint{2.062811in}{1.063580in}}{\pgfqpoint{2.035033in}{1.063580in}}%
\pgfpathlineto{\pgfqpoint{0.855248in}{1.063580in}}%
\pgfpathquadraticcurveto{\pgfqpoint{0.827470in}{1.063580in}}{\pgfqpoint{0.827470in}{1.035802in}}%
\pgfpathlineto{\pgfqpoint{0.827470in}{0.662346in}}%
\pgfpathquadraticcurveto{\pgfqpoint{0.827470in}{0.634568in}}{\pgfqpoint{0.855248in}{0.634568in}}%
\pgfpathlineto{\pgfqpoint{0.855248in}{0.634568in}}%
\pgfpathclose%
\pgfusepath{stroke,fill}%
\end{pgfscope}%
\begin{pgfscope}%
\pgfsetrectcap%
\pgfsetroundjoin%
\pgfsetlinewidth{1.003750pt}%
\definecolor{currentstroke}{rgb}{0.000000,0.000000,0.000000}%
\pgfsetstrokecolor{currentstroke}%
\pgfsetdash{}{0pt}%
\pgfpathmoveto{\pgfqpoint{0.883026in}{0.959413in}}%
\pgfpathlineto{\pgfqpoint{1.021915in}{0.959413in}}%
\pgfpathlineto{\pgfqpoint{1.160803in}{0.959413in}}%
\pgfusepath{stroke}%
\end{pgfscope}%
\begin{pgfscope}%
\definecolor{textcolor}{rgb}{0.000000,0.000000,0.000000}%
\pgfsetstrokecolor{textcolor}%
\pgfsetfillcolor{textcolor}%
\pgftext[x=1.271915in,y=0.910802in,left,base]{\color{textcolor}\rmfamily\fontsize{10.000000}{12.000000}\selectfont Rectangular}%
\end{pgfscope}%
\begin{pgfscope}%
\pgfsetrectcap%
\pgfsetroundjoin%
\pgfsetlinewidth{1.003750pt}%
\definecolor{currentstroke}{rgb}{1.000000,0.000000,0.000000}%
\pgfsetstrokecolor{currentstroke}%
\pgfsetdash{}{0pt}%
\pgfpathmoveto{\pgfqpoint{0.883026in}{0.765741in}}%
\pgfpathlineto{\pgfqpoint{1.021915in}{0.765741in}}%
\pgfpathlineto{\pgfqpoint{1.160803in}{0.765741in}}%
\pgfusepath{stroke}%
\end{pgfscope}%
\begin{pgfscope}%
\definecolor{textcolor}{rgb}{0.000000,0.000000,0.000000}%
\pgfsetstrokecolor{textcolor}%
\pgfsetfillcolor{textcolor}%
\pgftext[x=1.271915in,y=0.717129in,left,base]{\color{textcolor}\rmfamily\fontsize{10.000000}{12.000000}\selectfont Modulated}%
\end{pgfscope}%
\end{pgfpicture}%
\makeatother%
\endgroup%

		\caption{Measured frequency content of rectangular waves and a rectangular modulated sine wave.}
		\label{fig:445}
	\end{figure}
	
	The rectangle modulated sine signal shows some similarities to pure rectangle signal. The modulated signal shows side bands (fig. \ref{fig:444} red line) like discussed in the previous section. Moreover these bands including the center band repeat at uneven harmonics like known from a regular rectangle spectrum as can be seen in given figure.
	
	\subsubsection{}
	
	\begin{figure}[H]
		\centering
		%% Creator: Matplotlib, PGF backend
%%
%% To include the figure in your LaTeX document, write
%%   \input{<filename>.pgf}
%%
%% Make sure the required packages are loaded in your preamble
%%   \usepackage{pgf}
%%
%% Also ensure that all the required font packages are loaded; for instance,
%% the lmodern package is sometimes necessary when using math font.
%%   \usepackage{lmodern}
%%
%% Figures using additional raster images can only be included by \input if
%% they are in the same directory as the main LaTeX file. For loading figures
%% from other directories you can use the `import` package
%%   \usepackage{import}
%%
%% and then include the figures with
%%   \import{<path to file>}{<filename>.pgf}
%%
%% Matplotlib used the following preamble
%%
\begingroup%
\makeatletter%
\begin{pgfpicture}%
\pgfpathrectangle{\pgfpointorigin}{\pgfqpoint{6.300000in}{3.500000in}}%
\pgfusepath{use as bounding box, clip}%
\begin{pgfscope}%
\pgfsetbuttcap%
\pgfsetmiterjoin%
\definecolor{currentfill}{rgb}{1.000000,1.000000,1.000000}%
\pgfsetfillcolor{currentfill}%
\pgfsetlinewidth{0.000000pt}%
\definecolor{currentstroke}{rgb}{1.000000,1.000000,1.000000}%
\pgfsetstrokecolor{currentstroke}%
\pgfsetdash{}{0pt}%
\pgfpathmoveto{\pgfqpoint{0.000000in}{0.000000in}}%
\pgfpathlineto{\pgfqpoint{6.300000in}{0.000000in}}%
\pgfpathlineto{\pgfqpoint{6.300000in}{3.500000in}}%
\pgfpathlineto{\pgfqpoint{0.000000in}{3.500000in}}%
\pgfpathlineto{\pgfqpoint{0.000000in}{0.000000in}}%
\pgfpathclose%
\pgfusepath{fill}%
\end{pgfscope}%
\begin{pgfscope}%
\pgfsetbuttcap%
\pgfsetmiterjoin%
\definecolor{currentfill}{rgb}{1.000000,1.000000,1.000000}%
\pgfsetfillcolor{currentfill}%
\pgfsetlinewidth{0.000000pt}%
\definecolor{currentstroke}{rgb}{0.000000,0.000000,0.000000}%
\pgfsetstrokecolor{currentstroke}%
\pgfsetstrokeopacity{0.000000}%
\pgfsetdash{}{0pt}%
\pgfpathmoveto{\pgfqpoint{0.758026in}{2.240123in}}%
\pgfpathlineto{\pgfqpoint{2.866666in}{2.240123in}}%
\pgfpathlineto{\pgfqpoint{2.866666in}{3.150926in}}%
\pgfpathlineto{\pgfqpoint{0.758026in}{3.150926in}}%
\pgfpathlineto{\pgfqpoint{0.758026in}{2.240123in}}%
\pgfpathclose%
\pgfusepath{fill}%
\end{pgfscope}%
\begin{pgfscope}%
\pgfpathrectangle{\pgfqpoint{0.758026in}{2.240123in}}{\pgfqpoint{2.108640in}{0.910803in}}%
\pgfusepath{clip}%
\pgfsetbuttcap%
\pgfsetroundjoin%
\pgfsetlinewidth{0.501875pt}%
\definecolor{currentstroke}{rgb}{0.690196,0.690196,0.690196}%
\pgfsetstrokecolor{currentstroke}%
\pgfsetstrokeopacity{0.700000}%
\pgfsetdash{{1.850000pt}{0.800000pt}}{0.000000pt}%
\pgfpathmoveto{\pgfqpoint{0.758026in}{2.240123in}}%
\pgfpathlineto{\pgfqpoint{0.758026in}{3.150926in}}%
\pgfusepath{stroke}%
\end{pgfscope}%
\begin{pgfscope}%
\pgfsetbuttcap%
\pgfsetroundjoin%
\definecolor{currentfill}{rgb}{0.000000,0.000000,0.000000}%
\pgfsetfillcolor{currentfill}%
\pgfsetlinewidth{0.803000pt}%
\definecolor{currentstroke}{rgb}{0.000000,0.000000,0.000000}%
\pgfsetstrokecolor{currentstroke}%
\pgfsetdash{}{0pt}%
\pgfsys@defobject{currentmarker}{\pgfqpoint{0.000000in}{-0.048611in}}{\pgfqpoint{0.000000in}{0.000000in}}{%
\pgfpathmoveto{\pgfqpoint{0.000000in}{0.000000in}}%
\pgfpathlineto{\pgfqpoint{0.000000in}{-0.048611in}}%
\pgfusepath{stroke,fill}%
}%
\begin{pgfscope}%
\pgfsys@transformshift{0.758026in}{2.240123in}%
\pgfsys@useobject{currentmarker}{}%
\end{pgfscope}%
\end{pgfscope}%
\begin{pgfscope}%
\definecolor{textcolor}{rgb}{0.000000,0.000000,0.000000}%
\pgfsetstrokecolor{textcolor}%
\pgfsetfillcolor{textcolor}%
\pgftext[x=0.758026in,y=2.142901in,,top]{\color{textcolor}\rmfamily\fontsize{10.000000}{12.000000}\selectfont \(\displaystyle {200000}\)}%
\end{pgfscope}%
\begin{pgfscope}%
\pgfpathrectangle{\pgfqpoint{0.758026in}{2.240123in}}{\pgfqpoint{2.108640in}{0.910803in}}%
\pgfusepath{clip}%
\pgfsetbuttcap%
\pgfsetroundjoin%
\pgfsetlinewidth{0.501875pt}%
\definecolor{currentstroke}{rgb}{0.690196,0.690196,0.690196}%
\pgfsetstrokecolor{currentstroke}%
\pgfsetstrokeopacity{0.700000}%
\pgfsetdash{{1.850000pt}{0.800000pt}}{0.000000pt}%
\pgfpathmoveto{\pgfqpoint{1.460906in}{2.240123in}}%
\pgfpathlineto{\pgfqpoint{1.460906in}{3.150926in}}%
\pgfusepath{stroke}%
\end{pgfscope}%
\begin{pgfscope}%
\pgfsetbuttcap%
\pgfsetroundjoin%
\definecolor{currentfill}{rgb}{0.000000,0.000000,0.000000}%
\pgfsetfillcolor{currentfill}%
\pgfsetlinewidth{0.803000pt}%
\definecolor{currentstroke}{rgb}{0.000000,0.000000,0.000000}%
\pgfsetstrokecolor{currentstroke}%
\pgfsetdash{}{0pt}%
\pgfsys@defobject{currentmarker}{\pgfqpoint{0.000000in}{-0.048611in}}{\pgfqpoint{0.000000in}{0.000000in}}{%
\pgfpathmoveto{\pgfqpoint{0.000000in}{0.000000in}}%
\pgfpathlineto{\pgfqpoint{0.000000in}{-0.048611in}}%
\pgfusepath{stroke,fill}%
}%
\begin{pgfscope}%
\pgfsys@transformshift{1.460906in}{2.240123in}%
\pgfsys@useobject{currentmarker}{}%
\end{pgfscope}%
\end{pgfscope}%
\begin{pgfscope}%
\definecolor{textcolor}{rgb}{0.000000,0.000000,0.000000}%
\pgfsetstrokecolor{textcolor}%
\pgfsetfillcolor{textcolor}%
\pgftext[x=1.460906in,y=2.142901in,,top]{\color{textcolor}\rmfamily\fontsize{10.000000}{12.000000}\selectfont \(\displaystyle {400000}\)}%
\end{pgfscope}%
\begin{pgfscope}%
\pgfpathrectangle{\pgfqpoint{0.758026in}{2.240123in}}{\pgfqpoint{2.108640in}{0.910803in}}%
\pgfusepath{clip}%
\pgfsetbuttcap%
\pgfsetroundjoin%
\pgfsetlinewidth{0.501875pt}%
\definecolor{currentstroke}{rgb}{0.690196,0.690196,0.690196}%
\pgfsetstrokecolor{currentstroke}%
\pgfsetstrokeopacity{0.700000}%
\pgfsetdash{{1.850000pt}{0.800000pt}}{0.000000pt}%
\pgfpathmoveto{\pgfqpoint{2.163786in}{2.240123in}}%
\pgfpathlineto{\pgfqpoint{2.163786in}{3.150926in}}%
\pgfusepath{stroke}%
\end{pgfscope}%
\begin{pgfscope}%
\pgfsetbuttcap%
\pgfsetroundjoin%
\definecolor{currentfill}{rgb}{0.000000,0.000000,0.000000}%
\pgfsetfillcolor{currentfill}%
\pgfsetlinewidth{0.803000pt}%
\definecolor{currentstroke}{rgb}{0.000000,0.000000,0.000000}%
\pgfsetstrokecolor{currentstroke}%
\pgfsetdash{}{0pt}%
\pgfsys@defobject{currentmarker}{\pgfqpoint{0.000000in}{-0.048611in}}{\pgfqpoint{0.000000in}{0.000000in}}{%
\pgfpathmoveto{\pgfqpoint{0.000000in}{0.000000in}}%
\pgfpathlineto{\pgfqpoint{0.000000in}{-0.048611in}}%
\pgfusepath{stroke,fill}%
}%
\begin{pgfscope}%
\pgfsys@transformshift{2.163786in}{2.240123in}%
\pgfsys@useobject{currentmarker}{}%
\end{pgfscope}%
\end{pgfscope}%
\begin{pgfscope}%
\definecolor{textcolor}{rgb}{0.000000,0.000000,0.000000}%
\pgfsetstrokecolor{textcolor}%
\pgfsetfillcolor{textcolor}%
\pgftext[x=2.163786in,y=2.142901in,,top]{\color{textcolor}\rmfamily\fontsize{10.000000}{12.000000}\selectfont \(\displaystyle {600000}\)}%
\end{pgfscope}%
\begin{pgfscope}%
\pgfpathrectangle{\pgfqpoint{0.758026in}{2.240123in}}{\pgfqpoint{2.108640in}{0.910803in}}%
\pgfusepath{clip}%
\pgfsetbuttcap%
\pgfsetroundjoin%
\pgfsetlinewidth{0.501875pt}%
\definecolor{currentstroke}{rgb}{0.690196,0.690196,0.690196}%
\pgfsetstrokecolor{currentstroke}%
\pgfsetstrokeopacity{0.700000}%
\pgfsetdash{{1.850000pt}{0.800000pt}}{0.000000pt}%
\pgfpathmoveto{\pgfqpoint{2.866666in}{2.240123in}}%
\pgfpathlineto{\pgfqpoint{2.866666in}{3.150926in}}%
\pgfusepath{stroke}%
\end{pgfscope}%
\begin{pgfscope}%
\pgfsetbuttcap%
\pgfsetroundjoin%
\definecolor{currentfill}{rgb}{0.000000,0.000000,0.000000}%
\pgfsetfillcolor{currentfill}%
\pgfsetlinewidth{0.803000pt}%
\definecolor{currentstroke}{rgb}{0.000000,0.000000,0.000000}%
\pgfsetstrokecolor{currentstroke}%
\pgfsetdash{}{0pt}%
\pgfsys@defobject{currentmarker}{\pgfqpoint{0.000000in}{-0.048611in}}{\pgfqpoint{0.000000in}{0.000000in}}{%
\pgfpathmoveto{\pgfqpoint{0.000000in}{0.000000in}}%
\pgfpathlineto{\pgfqpoint{0.000000in}{-0.048611in}}%
\pgfusepath{stroke,fill}%
}%
\begin{pgfscope}%
\pgfsys@transformshift{2.866666in}{2.240123in}%
\pgfsys@useobject{currentmarker}{}%
\end{pgfscope}%
\end{pgfscope}%
\begin{pgfscope}%
\definecolor{textcolor}{rgb}{0.000000,0.000000,0.000000}%
\pgfsetstrokecolor{textcolor}%
\pgfsetfillcolor{textcolor}%
\pgftext[x=2.866666in,y=2.142901in,,top]{\color{textcolor}\rmfamily\fontsize{10.000000}{12.000000}\selectfont \(\displaystyle {800000}\)}%
\end{pgfscope}%
\begin{pgfscope}%
\definecolor{textcolor}{rgb}{0.000000,0.000000,0.000000}%
\pgfsetstrokecolor{textcolor}%
\pgfsetfillcolor{textcolor}%
\pgftext[x=1.812346in,y=1.963889in,,top]{\color{textcolor}\rmfamily\fontsize{10.000000}{12.000000}\selectfont Frequency (Hz)}%
\end{pgfscope}%
\begin{pgfscope}%
\pgfpathrectangle{\pgfqpoint{0.758026in}{2.240123in}}{\pgfqpoint{2.108640in}{0.910803in}}%
\pgfusepath{clip}%
\pgfsetbuttcap%
\pgfsetroundjoin%
\pgfsetlinewidth{0.501875pt}%
\definecolor{currentstroke}{rgb}{0.690196,0.690196,0.690196}%
\pgfsetstrokecolor{currentstroke}%
\pgfsetstrokeopacity{0.700000}%
\pgfsetdash{{1.850000pt}{0.800000pt}}{0.000000pt}%
\pgfpathmoveto{\pgfqpoint{0.758026in}{2.358113in}}%
\pgfpathlineto{\pgfqpoint{2.866666in}{2.358113in}}%
\pgfusepath{stroke}%
\end{pgfscope}%
\begin{pgfscope}%
\pgfsetbuttcap%
\pgfsetroundjoin%
\definecolor{currentfill}{rgb}{0.000000,0.000000,0.000000}%
\pgfsetfillcolor{currentfill}%
\pgfsetlinewidth{0.803000pt}%
\definecolor{currentstroke}{rgb}{0.000000,0.000000,0.000000}%
\pgfsetstrokecolor{currentstroke}%
\pgfsetdash{}{0pt}%
\pgfsys@defobject{currentmarker}{\pgfqpoint{-0.048611in}{0.000000in}}{\pgfqpoint{-0.000000in}{0.000000in}}{%
\pgfpathmoveto{\pgfqpoint{-0.000000in}{0.000000in}}%
\pgfpathlineto{\pgfqpoint{-0.048611in}{0.000000in}}%
\pgfusepath{stroke,fill}%
}%
\begin{pgfscope}%
\pgfsys@transformshift{0.758026in}{2.358113in}%
\pgfsys@useobject{currentmarker}{}%
\end{pgfscope}%
\end{pgfscope}%
\begin{pgfscope}%
\definecolor{textcolor}{rgb}{0.000000,0.000000,0.000000}%
\pgfsetstrokecolor{textcolor}%
\pgfsetfillcolor{textcolor}%
\pgftext[x=0.344444in, y=2.309888in, left, base]{\color{textcolor}\rmfamily\fontsize{10.000000}{12.000000}\selectfont \(\displaystyle {\ensuremath{-}100}\)}%
\end{pgfscope}%
\begin{pgfscope}%
\pgfpathrectangle{\pgfqpoint{0.758026in}{2.240123in}}{\pgfqpoint{2.108640in}{0.910803in}}%
\pgfusepath{clip}%
\pgfsetbuttcap%
\pgfsetroundjoin%
\pgfsetlinewidth{0.501875pt}%
\definecolor{currentstroke}{rgb}{0.690196,0.690196,0.690196}%
\pgfsetstrokecolor{currentstroke}%
\pgfsetstrokeopacity{0.700000}%
\pgfsetdash{{1.850000pt}{0.800000pt}}{0.000000pt}%
\pgfpathmoveto{\pgfqpoint{0.758026in}{2.873000in}}%
\pgfpathlineto{\pgfqpoint{2.866666in}{2.873000in}}%
\pgfusepath{stroke}%
\end{pgfscope}%
\begin{pgfscope}%
\pgfsetbuttcap%
\pgfsetroundjoin%
\definecolor{currentfill}{rgb}{0.000000,0.000000,0.000000}%
\pgfsetfillcolor{currentfill}%
\pgfsetlinewidth{0.803000pt}%
\definecolor{currentstroke}{rgb}{0.000000,0.000000,0.000000}%
\pgfsetstrokecolor{currentstroke}%
\pgfsetdash{}{0pt}%
\pgfsys@defobject{currentmarker}{\pgfqpoint{-0.048611in}{0.000000in}}{\pgfqpoint{-0.000000in}{0.000000in}}{%
\pgfpathmoveto{\pgfqpoint{-0.000000in}{0.000000in}}%
\pgfpathlineto{\pgfqpoint{-0.048611in}{0.000000in}}%
\pgfusepath{stroke,fill}%
}%
\begin{pgfscope}%
\pgfsys@transformshift{0.758026in}{2.873000in}%
\pgfsys@useobject{currentmarker}{}%
\end{pgfscope}%
\end{pgfscope}%
\begin{pgfscope}%
\definecolor{textcolor}{rgb}{0.000000,0.000000,0.000000}%
\pgfsetstrokecolor{textcolor}%
\pgfsetfillcolor{textcolor}%
\pgftext[x=0.413889in, y=2.824775in, left, base]{\color{textcolor}\rmfamily\fontsize{10.000000}{12.000000}\selectfont \(\displaystyle {\ensuremath{-}50}\)}%
\end{pgfscope}%
\begin{pgfscope}%
\definecolor{textcolor}{rgb}{0.000000,0.000000,0.000000}%
\pgfsetstrokecolor{textcolor}%
\pgfsetfillcolor{textcolor}%
\pgftext[x=0.288889in,y=2.695525in,,bottom,rotate=90.000000]{\color{textcolor}\rmfamily\fontsize{10.000000}{12.000000}\selectfont Amplitude (dB)}%
\end{pgfscope}%
\begin{pgfscope}%
\pgfpathrectangle{\pgfqpoint{0.758026in}{2.240123in}}{\pgfqpoint{2.108640in}{0.910803in}}%
\pgfusepath{clip}%
\pgfsetrectcap%
\pgfsetroundjoin%
\pgfsetlinewidth{1.003750pt}%
\definecolor{currentstroke}{rgb}{0.000000,0.000000,0.000000}%
\pgfsetstrokecolor{currentstroke}%
\pgfsetdash{}{0pt}%
\pgfpathmoveto{\pgfqpoint{0.757674in}{2.543150in}}%
\pgfpathlineto{\pgfqpoint{0.760837in}{2.538002in}}%
\pgfpathlineto{\pgfqpoint{0.764000in}{2.568573in}}%
\pgfpathlineto{\pgfqpoint{0.767163in}{2.475893in}}%
\pgfpathlineto{\pgfqpoint{0.770326in}{2.527382in}}%
\pgfpathlineto{\pgfqpoint{0.773489in}{2.517084in}}%
\pgfpathlineto{\pgfqpoint{0.776652in}{2.455298in}}%
\pgfpathlineto{\pgfqpoint{0.779815in}{2.542829in}}%
\pgfpathlineto{\pgfqpoint{0.782978in}{2.491340in}}%
\pgfpathlineto{\pgfqpoint{0.786141in}{2.599466in}}%
\pgfpathlineto{\pgfqpoint{0.789304in}{2.552805in}}%
\pgfpathlineto{\pgfqpoint{0.792467in}{2.547656in}}%
\pgfpathlineto{\pgfqpoint{0.795630in}{2.552805in}}%
\pgfpathlineto{\pgfqpoint{0.798793in}{2.573400in}}%
\pgfpathlineto{\pgfqpoint{0.801956in}{2.527060in}}%
\pgfpathlineto{\pgfqpoint{0.805119in}{2.547656in}}%
\pgfpathlineto{\pgfqpoint{0.808282in}{2.542507in}}%
\pgfpathlineto{\pgfqpoint{0.811445in}{2.583698in}}%
\pgfpathlineto{\pgfqpoint{0.814608in}{2.516441in}}%
\pgfpathlineto{\pgfqpoint{0.817771in}{2.552483in}}%
\pgfpathlineto{\pgfqpoint{0.820933in}{2.490696in}}%
\pgfpathlineto{\pgfqpoint{0.824096in}{2.495845in}}%
\pgfpathlineto{\pgfqpoint{0.827259in}{2.583376in}}%
\pgfpathlineto{\pgfqpoint{0.830422in}{2.537036in}}%
\pgfpathlineto{\pgfqpoint{0.833585in}{2.552483in}}%
\pgfpathlineto{\pgfqpoint{0.836748in}{2.588525in}}%
\pgfpathlineto{\pgfqpoint{0.839911in}{2.572756in}}%
\pgfpathlineto{\pgfqpoint{0.843074in}{2.598501in}}%
\pgfpathlineto{\pgfqpoint{0.846237in}{2.541863in}}%
\pgfpathlineto{\pgfqpoint{0.849400in}{2.552161in}}%
\pgfpathlineto{\pgfqpoint{0.852563in}{2.536714in}}%
\pgfpathlineto{\pgfqpoint{0.855726in}{2.557310in}}%
\pgfpathlineto{\pgfqpoint{0.858889in}{2.521268in}}%
\pgfpathlineto{\pgfqpoint{0.862052in}{2.557310in}}%
\pgfpathlineto{\pgfqpoint{0.865215in}{2.505499in}}%
\pgfpathlineto{\pgfqpoint{0.868378in}{2.598179in}}%
\pgfpathlineto{\pgfqpoint{0.871541in}{2.510648in}}%
\pgfpathlineto{\pgfqpoint{0.874704in}{2.536393in}}%
\pgfpathlineto{\pgfqpoint{0.877867in}{2.541541in}}%
\pgfpathlineto{\pgfqpoint{0.881030in}{2.567286in}}%
\pgfpathlineto{\pgfqpoint{0.884193in}{2.515797in}}%
\pgfpathlineto{\pgfqpoint{0.887356in}{2.510648in}}%
\pgfpathlineto{\pgfqpoint{0.890519in}{2.536071in}}%
\pgfpathlineto{\pgfqpoint{0.893682in}{2.474284in}}%
\pgfpathlineto{\pgfqpoint{0.896845in}{2.587559in}}%
\pgfpathlineto{\pgfqpoint{0.900007in}{2.510326in}}%
\pgfpathlineto{\pgfqpoint{0.903170in}{2.587559in}}%
\pgfpathlineto{\pgfqpoint{0.906333in}{2.577262in}}%
\pgfpathlineto{\pgfqpoint{0.909496in}{2.577262in}}%
\pgfpathlineto{\pgfqpoint{0.912659in}{2.489731in}}%
\pgfpathlineto{\pgfqpoint{0.915822in}{2.530600in}}%
\pgfpathlineto{\pgfqpoint{0.918985in}{2.479111in}}%
\pgfpathlineto{\pgfqpoint{0.922148in}{2.551196in}}%
\pgfpathlineto{\pgfqpoint{0.925311in}{2.576940in}}%
\pgfpathlineto{\pgfqpoint{0.928474in}{2.571791in}}%
\pgfpathlineto{\pgfqpoint{0.931637in}{2.618131in}}%
\pgfpathlineto{\pgfqpoint{0.934800in}{2.638726in}}%
\pgfpathlineto{\pgfqpoint{0.937963in}{2.685066in}}%
\pgfpathlineto{\pgfqpoint{0.941126in}{2.638405in}}%
\pgfpathlineto{\pgfqpoint{0.947452in}{2.525129in}}%
\pgfpathlineto{\pgfqpoint{0.950615in}{2.504534in}}%
\pgfpathlineto{\pgfqpoint{0.960104in}{2.602362in}}%
\pgfpathlineto{\pgfqpoint{0.966430in}{2.524808in}}%
\pgfpathlineto{\pgfqpoint{0.969593in}{2.540254in}}%
\pgfpathlineto{\pgfqpoint{0.972756in}{2.545403in}}%
\pgfpathlineto{\pgfqpoint{0.975919in}{2.545403in}}%
\pgfpathlineto{\pgfqpoint{0.979081in}{2.421830in}}%
\pgfpathlineto{\pgfqpoint{0.982244in}{2.565999in}}%
\pgfpathlineto{\pgfqpoint{0.985407in}{2.504212in}}%
\pgfpathlineto{\pgfqpoint{0.988570in}{2.560850in}}%
\pgfpathlineto{\pgfqpoint{0.991733in}{2.545081in}}%
\pgfpathlineto{\pgfqpoint{0.994896in}{2.545081in}}%
\pgfpathlineto{\pgfqpoint{0.998059in}{2.560528in}}%
\pgfpathlineto{\pgfqpoint{1.004385in}{2.514188in}}%
\pgfpathlineto{\pgfqpoint{1.007548in}{2.565677in}}%
\pgfpathlineto{\pgfqpoint{1.010711in}{2.565677in}}%
\pgfpathlineto{\pgfqpoint{1.013874in}{2.545081in}}%
\pgfpathlineto{\pgfqpoint{1.017037in}{2.555057in}}%
\pgfpathlineto{\pgfqpoint{1.020200in}{2.421187in}}%
\pgfpathlineto{\pgfqpoint{1.023363in}{2.601397in}}%
\pgfpathlineto{\pgfqpoint{1.029689in}{2.565355in}}%
\pgfpathlineto{\pgfqpoint{1.032852in}{2.457229in}}%
\pgfpathlineto{\pgfqpoint{1.036015in}{2.519015in}}%
\pgfpathlineto{\pgfqpoint{1.039178in}{2.534462in}}%
\pgfpathlineto{\pgfqpoint{1.042341in}{2.539289in}}%
\pgfpathlineto{\pgfqpoint{1.045504in}{2.559884in}}%
\pgfpathlineto{\pgfqpoint{1.048667in}{2.565033in}}%
\pgfpathlineto{\pgfqpoint{1.051830in}{2.534140in}}%
\pgfpathlineto{\pgfqpoint{1.054993in}{2.549587in}}%
\pgfpathlineto{\pgfqpoint{1.058156in}{2.539289in}}%
\pgfpathlineto{\pgfqpoint{1.061318in}{2.570182in}}%
\pgfpathlineto{\pgfqpoint{1.064481in}{2.492949in}}%
\pgfpathlineto{\pgfqpoint{1.067644in}{2.544116in}}%
\pgfpathlineto{\pgfqpoint{1.070807in}{2.662540in}}%
\pgfpathlineto{\pgfqpoint{1.073970in}{2.734624in}}%
\pgfpathlineto{\pgfqpoint{1.077133in}{2.780964in}}%
\pgfpathlineto{\pgfqpoint{1.083459in}{2.801559in}}%
\pgfpathlineto{\pgfqpoint{1.086622in}{2.780964in}}%
\pgfpathlineto{\pgfqpoint{1.089785in}{2.750071in}}%
\pgfpathlineto{\pgfqpoint{1.096111in}{2.595283in}}%
\pgfpathlineto{\pgfqpoint{1.099274in}{2.518050in}}%
\pgfpathlineto{\pgfqpoint{1.102437in}{2.574687in}}%
\pgfpathlineto{\pgfqpoint{1.105600in}{2.512901in}}%
\pgfpathlineto{\pgfqpoint{1.108763in}{2.579836in}}%
\pgfpathlineto{\pgfqpoint{1.111926in}{2.528347in}}%
\pgfpathlineto{\pgfqpoint{1.115089in}{2.523199in}}%
\pgfpathlineto{\pgfqpoint{1.118252in}{2.491984in}}%
\pgfpathlineto{\pgfqpoint{1.121415in}{2.533174in}}%
\pgfpathlineto{\pgfqpoint{1.124578in}{2.548621in}}%
\pgfpathlineto{\pgfqpoint{1.127741in}{2.491984in}}%
\pgfpathlineto{\pgfqpoint{1.130904in}{2.594961in}}%
\pgfpathlineto{\pgfqpoint{1.134067in}{2.564068in}}%
\pgfpathlineto{\pgfqpoint{1.140392in}{2.600110in}}%
\pgfpathlineto{\pgfqpoint{1.143555in}{2.579193in}}%
\pgfpathlineto{\pgfqpoint{1.146718in}{2.527704in}}%
\pgfpathlineto{\pgfqpoint{1.149881in}{2.517406in}}%
\pgfpathlineto{\pgfqpoint{1.153044in}{2.589490in}}%
\pgfpathlineto{\pgfqpoint{1.156207in}{2.543150in}}%
\pgfpathlineto{\pgfqpoint{1.159370in}{2.568895in}}%
\pgfpathlineto{\pgfqpoint{1.162533in}{2.610086in}}%
\pgfpathlineto{\pgfqpoint{1.165696in}{2.599788in}}%
\pgfpathlineto{\pgfqpoint{1.168859in}{2.573722in}}%
\pgfpathlineto{\pgfqpoint{1.172022in}{2.501638in}}%
\pgfpathlineto{\pgfqpoint{1.175185in}{2.547977in}}%
\pgfpathlineto{\pgfqpoint{1.178348in}{2.568573in}}%
\pgfpathlineto{\pgfqpoint{1.181511in}{2.481042in}}%
\pgfpathlineto{\pgfqpoint{1.184674in}{2.609764in}}%
\pgfpathlineto{\pgfqpoint{1.187837in}{2.522233in}}%
\pgfpathlineto{\pgfqpoint{1.191000in}{2.563424in}}%
\pgfpathlineto{\pgfqpoint{1.194163in}{2.568251in}}%
\pgfpathlineto{\pgfqpoint{1.197326in}{2.588847in}}%
\pgfpathlineto{\pgfqpoint{1.200489in}{2.547656in}}%
\pgfpathlineto{\pgfqpoint{1.203652in}{2.552805in}}%
\pgfpathlineto{\pgfqpoint{1.206815in}{2.552805in}}%
\pgfpathlineto{\pgfqpoint{1.209978in}{2.527060in}}%
\pgfpathlineto{\pgfqpoint{1.219466in}{2.835671in}}%
\pgfpathlineto{\pgfqpoint{1.225792in}{2.907755in}}%
\pgfpathlineto{\pgfqpoint{1.228955in}{2.907755in}}%
\pgfpathlineto{\pgfqpoint{1.232118in}{2.902606in}}%
\pgfpathlineto{\pgfqpoint{1.235281in}{2.866564in}}%
\pgfpathlineto{\pgfqpoint{1.238444in}{2.815075in}}%
\pgfpathlineto{\pgfqpoint{1.241607in}{2.722395in}}%
\pgfpathlineto{\pgfqpoint{1.247933in}{2.469779in}}%
\pgfpathlineto{\pgfqpoint{1.251096in}{2.613947in}}%
\pgfpathlineto{\pgfqpoint{1.254259in}{2.552161in}}%
\pgfpathlineto{\pgfqpoint{1.257422in}{2.567608in}}%
\pgfpathlineto{\pgfqpoint{1.260585in}{2.567608in}}%
\pgfpathlineto{\pgfqpoint{1.263748in}{2.634543in}}%
\pgfpathlineto{\pgfqpoint{1.270074in}{2.546690in}}%
\pgfpathlineto{\pgfqpoint{1.273237in}{2.618774in}}%
\pgfpathlineto{\pgfqpoint{1.276400in}{2.454011in}}%
\pgfpathlineto{\pgfqpoint{1.279563in}{2.551839in}}%
\pgfpathlineto{\pgfqpoint{1.282726in}{2.582732in}}%
\pgfpathlineto{\pgfqpoint{1.285889in}{2.562137in}}%
\pgfpathlineto{\pgfqpoint{1.289052in}{2.613626in}}%
\pgfpathlineto{\pgfqpoint{1.292215in}{2.582732in}}%
\pgfpathlineto{\pgfqpoint{1.295378in}{2.577262in}}%
\pgfpathlineto{\pgfqpoint{1.301703in}{2.525773in}}%
\pgfpathlineto{\pgfqpoint{1.304866in}{2.525773in}}%
\pgfpathlineto{\pgfqpoint{1.308029in}{2.520624in}}%
\pgfpathlineto{\pgfqpoint{1.311192in}{2.546368in}}%
\pgfpathlineto{\pgfqpoint{1.314355in}{2.551517in}}%
\pgfpathlineto{\pgfqpoint{1.317518in}{2.510326in}}%
\pgfpathlineto{\pgfqpoint{1.320681in}{2.510005in}}%
\pgfpathlineto{\pgfqpoint{1.323844in}{2.633577in}}%
\pgfpathlineto{\pgfqpoint{1.327007in}{2.525451in}}%
\pgfpathlineto{\pgfqpoint{1.330170in}{2.525451in}}%
\pgfpathlineto{\pgfqpoint{1.336496in}{2.504856in}}%
\pgfpathlineto{\pgfqpoint{1.339659in}{2.576940in}}%
\pgfpathlineto{\pgfqpoint{1.342822in}{2.592386in}}%
\pgfpathlineto{\pgfqpoint{1.345985in}{2.550874in}}%
\pgfpathlineto{\pgfqpoint{1.349148in}{2.576618in}}%
\pgfpathlineto{\pgfqpoint{1.352311in}{2.556023in}}%
\pgfpathlineto{\pgfqpoint{1.355474in}{2.586916in}}%
\pgfpathlineto{\pgfqpoint{1.364963in}{2.926741in}}%
\pgfpathlineto{\pgfqpoint{1.368126in}{2.983379in}}%
\pgfpathlineto{\pgfqpoint{1.371289in}{3.013950in}}%
\pgfpathlineto{\pgfqpoint{1.374452in}{3.024248in}}%
\pgfpathlineto{\pgfqpoint{1.377615in}{3.019099in}}%
\pgfpathlineto{\pgfqpoint{1.380777in}{2.988206in}}%
\pgfpathlineto{\pgfqpoint{1.383940in}{2.941866in}}%
\pgfpathlineto{\pgfqpoint{1.387103in}{2.864633in}}%
\pgfpathlineto{\pgfqpoint{1.390266in}{2.741060in}}%
\pgfpathlineto{\pgfqpoint{1.393429in}{2.555701in}}%
\pgfpathlineto{\pgfqpoint{1.396592in}{2.601719in}}%
\pgfpathlineto{\pgfqpoint{1.399755in}{2.524486in}}%
\pgfpathlineto{\pgfqpoint{1.402918in}{2.581123in}}%
\pgfpathlineto{\pgfqpoint{1.406081in}{2.570826in}}%
\pgfpathlineto{\pgfqpoint{1.409244in}{2.627463in}}%
\pgfpathlineto{\pgfqpoint{1.412407in}{2.539932in}}%
\pgfpathlineto{\pgfqpoint{1.415570in}{2.514188in}}%
\pgfpathlineto{\pgfqpoint{1.418733in}{2.601719in}}%
\pgfpathlineto{\pgfqpoint{1.421896in}{2.606546in}}%
\pgfpathlineto{\pgfqpoint{1.425059in}{2.519015in}}%
\pgfpathlineto{\pgfqpoint{1.428222in}{2.606546in}}%
\pgfpathlineto{\pgfqpoint{1.431385in}{2.596248in}}%
\pgfpathlineto{\pgfqpoint{1.434548in}{2.601397in}}%
\pgfpathlineto{\pgfqpoint{1.437711in}{2.668332in}}%
\pgfpathlineto{\pgfqpoint{1.440874in}{2.565355in}}%
\pgfpathlineto{\pgfqpoint{1.444037in}{2.519015in}}%
\pgfpathlineto{\pgfqpoint{1.447200in}{2.606224in}}%
\pgfpathlineto{\pgfqpoint{1.450363in}{2.528991in}}%
\pgfpathlineto{\pgfqpoint{1.456689in}{2.647415in}}%
\pgfpathlineto{\pgfqpoint{1.459851in}{2.595926in}}%
\pgfpathlineto{\pgfqpoint{1.463014in}{2.621671in}}%
\pgfpathlineto{\pgfqpoint{1.466177in}{2.595926in}}%
\pgfpathlineto{\pgfqpoint{1.469340in}{2.513544in}}%
\pgfpathlineto{\pgfqpoint{1.472503in}{2.565033in}}%
\pgfpathlineto{\pgfqpoint{1.475666in}{2.580158in}}%
\pgfpathlineto{\pgfqpoint{1.478829in}{2.528669in}}%
\pgfpathlineto{\pgfqpoint{1.481992in}{2.580158in}}%
\pgfpathlineto{\pgfqpoint{1.485155in}{2.487478in}}%
\pgfpathlineto{\pgfqpoint{1.488318in}{2.482329in}}%
\pgfpathlineto{\pgfqpoint{1.494644in}{2.611051in}}%
\pgfpathlineto{\pgfqpoint{1.497807in}{2.564711in}}%
\pgfpathlineto{\pgfqpoint{1.500970in}{2.600432in}}%
\pgfpathlineto{\pgfqpoint{1.510459in}{2.976299in}}%
\pgfpathlineto{\pgfqpoint{1.513622in}{3.038086in}}%
\pgfpathlineto{\pgfqpoint{1.516785in}{3.074128in}}%
\pgfpathlineto{\pgfqpoint{1.519948in}{3.094723in}}%
\pgfpathlineto{\pgfqpoint{1.523111in}{3.089574in}}%
\pgfpathlineto{\pgfqpoint{1.526274in}{3.068657in}}%
\pgfpathlineto{\pgfqpoint{1.529437in}{3.027466in}}%
\pgfpathlineto{\pgfqpoint{1.532600in}{2.960531in}}%
\pgfpathlineto{\pgfqpoint{1.535763in}{2.852404in}}%
\pgfpathlineto{\pgfqpoint{1.538926in}{2.646450in}}%
\pgfpathlineto{\pgfqpoint{1.542088in}{2.610408in}}%
\pgfpathlineto{\pgfqpoint{1.545251in}{2.553770in}}%
\pgfpathlineto{\pgfqpoint{1.548414in}{2.564068in}}%
\pgfpathlineto{\pgfqpoint{1.551577in}{2.563746in}}%
\pgfpathlineto{\pgfqpoint{1.554740in}{2.548299in}}%
\pgfpathlineto{\pgfqpoint{1.557903in}{2.512257in}}%
\pgfpathlineto{\pgfqpoint{1.561066in}{2.574044in}}%
\pgfpathlineto{\pgfqpoint{1.567392in}{2.527704in}}%
\pgfpathlineto{\pgfqpoint{1.570555in}{2.538002in}}%
\pgfpathlineto{\pgfqpoint{1.573718in}{2.517406in}}%
\pgfpathlineto{\pgfqpoint{1.576881in}{2.578871in}}%
\pgfpathlineto{\pgfqpoint{1.580044in}{2.465596in}}%
\pgfpathlineto{\pgfqpoint{1.583207in}{2.573722in}}%
\pgfpathlineto{\pgfqpoint{1.586370in}{2.563424in}}%
\pgfpathlineto{\pgfqpoint{1.589533in}{2.614913in}}%
\pgfpathlineto{\pgfqpoint{1.592696in}{2.517084in}}%
\pgfpathlineto{\pgfqpoint{1.595859in}{2.470744in}}%
\pgfpathlineto{\pgfqpoint{1.599022in}{2.573722in}}%
\pgfpathlineto{\pgfqpoint{1.602185in}{2.547656in}}%
\pgfpathlineto{\pgfqpoint{1.605348in}{2.496167in}}%
\pgfpathlineto{\pgfqpoint{1.608511in}{2.475572in}}%
\pgfpathlineto{\pgfqpoint{1.611674in}{2.578549in}}%
\pgfpathlineto{\pgfqpoint{1.614837in}{2.604293in}}%
\pgfpathlineto{\pgfqpoint{1.618000in}{2.609442in}}%
\pgfpathlineto{\pgfqpoint{1.621162in}{2.537358in}}%
\pgfpathlineto{\pgfqpoint{1.624325in}{2.562780in}}%
\pgfpathlineto{\pgfqpoint{1.627488in}{2.521590in}}%
\pgfpathlineto{\pgfqpoint{1.630651in}{2.542185in}}%
\pgfpathlineto{\pgfqpoint{1.633814in}{2.547334in}}%
\pgfpathlineto{\pgfqpoint{1.636977in}{2.588525in}}%
\pgfpathlineto{\pgfqpoint{1.640140in}{2.495845in}}%
\pgfpathlineto{\pgfqpoint{1.643303in}{2.531887in}}%
\pgfpathlineto{\pgfqpoint{1.649629in}{2.644841in}}%
\pgfpathlineto{\pgfqpoint{1.652792in}{2.861093in}}%
\pgfpathlineto{\pgfqpoint{1.655955in}{2.974368in}}%
\pgfpathlineto{\pgfqpoint{1.662281in}{3.087643in}}%
\pgfpathlineto{\pgfqpoint{1.665444in}{3.108239in}}%
\pgfpathlineto{\pgfqpoint{1.668607in}{3.108239in}}%
\pgfpathlineto{\pgfqpoint{1.671770in}{3.097941in}}%
\pgfpathlineto{\pgfqpoint{1.674933in}{3.061577in}}%
\pgfpathlineto{\pgfqpoint{1.678096in}{2.999791in}}%
\pgfpathlineto{\pgfqpoint{1.681259in}{2.907111in}}%
\pgfpathlineto{\pgfqpoint{1.690748in}{2.448862in}}%
\pgfpathlineto{\pgfqpoint{1.693911in}{2.515797in}}%
\pgfpathlineto{\pgfqpoint{1.697074in}{2.551839in}}%
\pgfpathlineto{\pgfqpoint{1.700236in}{2.525773in}}%
\pgfpathlineto{\pgfqpoint{1.703399in}{2.551517in}}%
\pgfpathlineto{\pgfqpoint{1.706562in}{2.556666in}}%
\pgfpathlineto{\pgfqpoint{1.709725in}{2.505178in}}%
\pgfpathlineto{\pgfqpoint{1.712888in}{2.551517in}}%
\pgfpathlineto{\pgfqpoint{1.716051in}{2.484582in}}%
\pgfpathlineto{\pgfqpoint{1.719214in}{2.566964in}}%
\pgfpathlineto{\pgfqpoint{1.722377in}{2.546368in}}%
\pgfpathlineto{\pgfqpoint{1.725540in}{2.489409in}}%
\pgfpathlineto{\pgfqpoint{1.728703in}{2.571791in}}%
\pgfpathlineto{\pgfqpoint{1.731866in}{2.510005in}}%
\pgfpathlineto{\pgfqpoint{1.735029in}{2.525451in}}%
\pgfpathlineto{\pgfqpoint{1.738192in}{2.510005in}}%
\pgfpathlineto{\pgfqpoint{1.741355in}{2.525451in}}%
\pgfpathlineto{\pgfqpoint{1.744518in}{2.551196in}}%
\pgfpathlineto{\pgfqpoint{1.750844in}{2.540576in}}%
\pgfpathlineto{\pgfqpoint{1.754007in}{2.530278in}}%
\pgfpathlineto{\pgfqpoint{1.757170in}{2.561171in}}%
\pgfpathlineto{\pgfqpoint{1.760333in}{2.530278in}}%
\pgfpathlineto{\pgfqpoint{1.763496in}{2.571469in}}%
\pgfpathlineto{\pgfqpoint{1.766659in}{2.556023in}}%
\pgfpathlineto{\pgfqpoint{1.769822in}{2.483938in}}%
\pgfpathlineto{\pgfqpoint{1.772985in}{2.556023in}}%
\pgfpathlineto{\pgfqpoint{1.776148in}{2.442426in}}%
\pgfpathlineto{\pgfqpoint{1.779311in}{2.514510in}}%
\pgfpathlineto{\pgfqpoint{1.782473in}{2.442426in}}%
\pgfpathlineto{\pgfqpoint{1.785636in}{2.514510in}}%
\pgfpathlineto{\pgfqpoint{1.788799in}{2.519659in}}%
\pgfpathlineto{\pgfqpoint{1.791962in}{2.499063in}}%
\pgfpathlineto{\pgfqpoint{1.795125in}{2.550552in}}%
\pgfpathlineto{\pgfqpoint{1.798288in}{2.473319in}}%
\pgfpathlineto{\pgfqpoint{1.801451in}{2.483295in}}%
\pgfpathlineto{\pgfqpoint{1.804614in}{2.514188in}}%
\pgfpathlineto{\pgfqpoint{1.810940in}{2.452402in}}%
\pgfpathlineto{\pgfqpoint{1.814103in}{2.478146in}}%
\pgfpathlineto{\pgfqpoint{1.817266in}{2.524486in}}%
\pgfpathlineto{\pgfqpoint{1.820429in}{2.534784in}}%
\pgfpathlineto{\pgfqpoint{1.823592in}{2.462699in}}%
\pgfpathlineto{\pgfqpoint{1.826755in}{2.488122in}}%
\pgfpathlineto{\pgfqpoint{1.829918in}{2.524164in}}%
\pgfpathlineto{\pgfqpoint{1.833081in}{2.410889in}}%
\pgfpathlineto{\pgfqpoint{1.836244in}{2.457229in}}%
\pgfpathlineto{\pgfqpoint{1.839407in}{2.585950in}}%
\pgfpathlineto{\pgfqpoint{1.842570in}{2.513866in}}%
\pgfpathlineto{\pgfqpoint{1.845733in}{2.555057in}}%
\pgfpathlineto{\pgfqpoint{1.848896in}{2.493271in}}%
\pgfpathlineto{\pgfqpoint{1.852059in}{2.518693in}}%
\pgfpathlineto{\pgfqpoint{1.855222in}{2.492949in}}%
\pgfpathlineto{\pgfqpoint{1.858385in}{2.446609in}}%
\pgfpathlineto{\pgfqpoint{1.861547in}{2.554735in}}%
\pgfpathlineto{\pgfqpoint{1.864710in}{2.472353in}}%
\pgfpathlineto{\pgfqpoint{1.867873in}{2.487800in}}%
\pgfpathlineto{\pgfqpoint{1.871036in}{2.585629in}}%
\pgfpathlineto{\pgfqpoint{1.874199in}{2.503247in}}%
\pgfpathlineto{\pgfqpoint{1.877362in}{2.544116in}}%
\pgfpathlineto{\pgfqpoint{1.880525in}{2.528669in}}%
\pgfpathlineto{\pgfqpoint{1.883688in}{2.523520in}}%
\pgfpathlineto{\pgfqpoint{1.886851in}{2.544116in}}%
\pgfpathlineto{\pgfqpoint{1.890014in}{2.487478in}}%
\pgfpathlineto{\pgfqpoint{1.893177in}{2.554414in}}%
\pgfpathlineto{\pgfqpoint{1.896340in}{2.528669in}}%
\pgfpathlineto{\pgfqpoint{1.899503in}{2.472032in}}%
\pgfpathlineto{\pgfqpoint{1.902666in}{2.492305in}}%
\pgfpathlineto{\pgfqpoint{1.905829in}{2.466561in}}%
\pgfpathlineto{\pgfqpoint{1.908992in}{2.533496in}}%
\pgfpathlineto{\pgfqpoint{1.912155in}{2.476859in}}%
\pgfpathlineto{\pgfqpoint{1.915318in}{2.528347in}}%
\pgfpathlineto{\pgfqpoint{1.918481in}{2.482008in}}%
\pgfpathlineto{\pgfqpoint{1.921644in}{2.574687in}}%
\pgfpathlineto{\pgfqpoint{1.924807in}{2.487156in}}%
\pgfpathlineto{\pgfqpoint{1.927970in}{2.481686in}}%
\pgfpathlineto{\pgfqpoint{1.931133in}{2.502281in}}%
\pgfpathlineto{\pgfqpoint{1.934296in}{2.507430in}}%
\pgfpathlineto{\pgfqpoint{1.937459in}{2.528026in}}%
\pgfpathlineto{\pgfqpoint{1.940621in}{2.486835in}}%
\pgfpathlineto{\pgfqpoint{1.943784in}{2.775171in}}%
\pgfpathlineto{\pgfqpoint{1.946947in}{2.924489in}}%
\pgfpathlineto{\pgfqpoint{1.950110in}{3.012019in}}%
\pgfpathlineto{\pgfqpoint{1.953273in}{3.068335in}}%
\pgfpathlineto{\pgfqpoint{1.956436in}{3.104377in}}%
\pgfpathlineto{\pgfqpoint{1.959599in}{3.109526in}}%
\pgfpathlineto{\pgfqpoint{1.962762in}{3.109526in}}%
\pgfpathlineto{\pgfqpoint{1.965925in}{3.083782in}}%
\pgfpathlineto{\pgfqpoint{1.969088in}{3.037442in}}%
\pgfpathlineto{\pgfqpoint{1.972251in}{2.960209in}}%
\pgfpathlineto{\pgfqpoint{1.975414in}{2.846934in}}%
\pgfpathlineto{\pgfqpoint{1.978577in}{2.599466in}}%
\pgfpathlineto{\pgfqpoint{1.981740in}{2.537680in}}%
\pgfpathlineto{\pgfqpoint{1.984903in}{2.501638in}}%
\pgfpathlineto{\pgfqpoint{1.988066in}{2.527382in}}%
\pgfpathlineto{\pgfqpoint{1.991229in}{2.527382in}}%
\pgfpathlineto{\pgfqpoint{1.994392in}{2.481042in}}%
\pgfpathlineto{\pgfqpoint{1.997555in}{2.553126in}}%
\pgfpathlineto{\pgfqpoint{2.000718in}{2.491340in}}%
\pgfpathlineto{\pgfqpoint{2.003881in}{2.491018in}}%
\pgfpathlineto{\pgfqpoint{2.010207in}{2.573400in}}%
\pgfpathlineto{\pgfqpoint{2.016533in}{2.511614in}}%
\pgfpathlineto{\pgfqpoint{2.019696in}{2.563102in}}%
\pgfpathlineto{\pgfqpoint{2.022858in}{2.563102in}}%
\pgfpathlineto{\pgfqpoint{2.026021in}{2.557953in}}%
\pgfpathlineto{\pgfqpoint{2.029184in}{2.444356in}}%
\pgfpathlineto{\pgfqpoint{2.032347in}{2.583376in}}%
\pgfpathlineto{\pgfqpoint{2.035510in}{2.470101in}}%
\pgfpathlineto{\pgfqpoint{2.038673in}{2.444356in}}%
\pgfpathlineto{\pgfqpoint{2.041836in}{2.593674in}}%
\pgfpathlineto{\pgfqpoint{2.044999in}{2.567929in}}%
\pgfpathlineto{\pgfqpoint{2.048162in}{2.459803in}}%
\pgfpathlineto{\pgfqpoint{2.051325in}{2.506143in}}%
\pgfpathlineto{\pgfqpoint{2.054488in}{2.583054in}}%
\pgfpathlineto{\pgfqpoint{2.060814in}{2.526417in}}%
\pgfpathlineto{\pgfqpoint{2.063977in}{2.547012in}}%
\pgfpathlineto{\pgfqpoint{2.067140in}{2.552161in}}%
\pgfpathlineto{\pgfqpoint{2.070303in}{2.500672in}}%
\pgfpathlineto{\pgfqpoint{2.073466in}{2.577905in}}%
\pgfpathlineto{\pgfqpoint{2.076629in}{2.510970in}}%
\pgfpathlineto{\pgfqpoint{2.079792in}{2.582732in}}%
\pgfpathlineto{\pgfqpoint{2.082955in}{2.572435in}}%
\pgfpathlineto{\pgfqpoint{2.086118in}{2.531244in}}%
\pgfpathlineto{\pgfqpoint{2.095607in}{2.979195in}}%
\pgfpathlineto{\pgfqpoint{2.101932in}{3.082173in}}%
\pgfpathlineto{\pgfqpoint{2.105095in}{3.097298in}}%
\pgfpathlineto{\pgfqpoint{2.108258in}{3.097298in}}%
\pgfpathlineto{\pgfqpoint{2.111421in}{3.076702in}}%
\pgfpathlineto{\pgfqpoint{2.114584in}{3.035511in}}%
\pgfpathlineto{\pgfqpoint{2.117747in}{2.968576in}}%
\pgfpathlineto{\pgfqpoint{2.120910in}{2.870747in}}%
\pgfpathlineto{\pgfqpoint{2.127236in}{2.505178in}}%
\pgfpathlineto{\pgfqpoint{2.130399in}{2.458516in}}%
\pgfpathlineto{\pgfqpoint{2.133562in}{2.443069in}}%
\pgfpathlineto{\pgfqpoint{2.136725in}{2.571791in}}%
\pgfpathlineto{\pgfqpoint{2.139888in}{2.520302in}}%
\pgfpathlineto{\pgfqpoint{2.143051in}{2.499707in}}%
\pgfpathlineto{\pgfqpoint{2.146214in}{2.520302in}}%
\pgfpathlineto{\pgfqpoint{2.149377in}{2.499707in}}%
\pgfpathlineto{\pgfqpoint{2.152540in}{2.504856in}}%
\pgfpathlineto{\pgfqpoint{2.155703in}{2.494236in}}%
\pgfpathlineto{\pgfqpoint{2.158866in}{2.530278in}}%
\pgfpathlineto{\pgfqpoint{2.162029in}{2.422152in}}%
\pgfpathlineto{\pgfqpoint{2.171518in}{2.525129in}}%
\pgfpathlineto{\pgfqpoint{2.174681in}{2.514832in}}%
\pgfpathlineto{\pgfqpoint{2.177844in}{2.473641in}}%
\pgfpathlineto{\pgfqpoint{2.181006in}{2.560850in}}%
\pgfpathlineto{\pgfqpoint{2.184169in}{2.499063in}}%
\pgfpathlineto{\pgfqpoint{2.187332in}{2.560850in}}%
\pgfpathlineto{\pgfqpoint{2.193658in}{2.514510in}}%
\pgfpathlineto{\pgfqpoint{2.196821in}{2.529956in}}%
\pgfpathlineto{\pgfqpoint{2.199984in}{2.571147in}}%
\pgfpathlineto{\pgfqpoint{2.203147in}{2.524808in}}%
\pgfpathlineto{\pgfqpoint{2.206310in}{2.524486in}}%
\pgfpathlineto{\pgfqpoint{2.209473in}{2.555379in}}%
\pgfpathlineto{\pgfqpoint{2.212636in}{2.524486in}}%
\pgfpathlineto{\pgfqpoint{2.215799in}{2.519337in}}%
\pgfpathlineto{\pgfqpoint{2.218962in}{2.442104in}}%
\pgfpathlineto{\pgfqpoint{2.222125in}{2.534784in}}%
\pgfpathlineto{\pgfqpoint{2.225288in}{2.472997in}}%
\pgfpathlineto{\pgfqpoint{2.228451in}{2.498741in}}%
\pgfpathlineto{\pgfqpoint{2.231614in}{2.472675in}}%
\pgfpathlineto{\pgfqpoint{2.234777in}{2.529313in}}%
\pgfpathlineto{\pgfqpoint{2.237940in}{2.771310in}}%
\pgfpathlineto{\pgfqpoint{2.241103in}{2.884585in}}%
\pgfpathlineto{\pgfqpoint{2.244266in}{2.956669in}}%
\pgfpathlineto{\pgfqpoint{2.247429in}{3.003009in}}%
\pgfpathlineto{\pgfqpoint{2.250592in}{3.023604in}}%
\pgfpathlineto{\pgfqpoint{2.253755in}{3.028753in}}%
\pgfpathlineto{\pgfqpoint{2.256918in}{3.012985in}}%
\pgfpathlineto{\pgfqpoint{2.260081in}{2.976943in}}%
\pgfpathlineto{\pgfqpoint{2.263243in}{2.925454in}}%
\pgfpathlineto{\pgfqpoint{2.266406in}{2.832774in}}%
\pgfpathlineto{\pgfqpoint{2.272732in}{2.472353in}}%
\pgfpathlineto{\pgfqpoint{2.275895in}{2.467205in}}%
\pgfpathlineto{\pgfqpoint{2.279058in}{2.410567in}}%
\pgfpathlineto{\pgfqpoint{2.282221in}{2.420543in}}%
\pgfpathlineto{\pgfqpoint{2.285384in}{2.435990in}}%
\pgfpathlineto{\pgfqpoint{2.288547in}{2.487478in}}%
\pgfpathlineto{\pgfqpoint{2.291710in}{2.446287in}}%
\pgfpathlineto{\pgfqpoint{2.298036in}{2.435990in}}%
\pgfpathlineto{\pgfqpoint{2.301199in}{2.441138in}}%
\pgfpathlineto{\pgfqpoint{2.304362in}{2.451436in}}%
\pgfpathlineto{\pgfqpoint{2.307525in}{2.440817in}}%
\pgfpathlineto{\pgfqpoint{2.310688in}{2.523199in}}%
\pgfpathlineto{\pgfqpoint{2.313851in}{2.487156in}}%
\pgfpathlineto{\pgfqpoint{2.317014in}{2.466561in}}%
\pgfpathlineto{\pgfqpoint{2.320177in}{2.420221in}}%
\pgfpathlineto{\pgfqpoint{2.323340in}{2.466561in}}%
\pgfpathlineto{\pgfqpoint{2.326503in}{2.451114in}}%
\pgfpathlineto{\pgfqpoint{2.329666in}{2.466561in}}%
\pgfpathlineto{\pgfqpoint{2.332829in}{2.425048in}}%
\pgfpathlineto{\pgfqpoint{2.335992in}{2.481686in}}%
\pgfpathlineto{\pgfqpoint{2.339155in}{2.481686in}}%
\pgfpathlineto{\pgfqpoint{2.342317in}{2.404453in}}%
\pgfpathlineto{\pgfqpoint{2.345480in}{2.450793in}}%
\pgfpathlineto{\pgfqpoint{2.348643in}{2.435346in}}%
\pgfpathlineto{\pgfqpoint{2.351806in}{2.445644in}}%
\pgfpathlineto{\pgfqpoint{2.354969in}{2.419899in}}%
\pgfpathlineto{\pgfqpoint{2.358132in}{2.486513in}}%
\pgfpathlineto{\pgfqpoint{2.361295in}{2.476215in}}%
\pgfpathlineto{\pgfqpoint{2.364458in}{2.404131in}}%
\pgfpathlineto{\pgfqpoint{2.367621in}{2.435024in}}%
\pgfpathlineto{\pgfqpoint{2.370784in}{2.435024in}}%
\pgfpathlineto{\pgfqpoint{2.373947in}{2.388684in}}%
\pgfpathlineto{\pgfqpoint{2.377110in}{2.435024in}}%
\pgfpathlineto{\pgfqpoint{2.380273in}{2.424726in}}%
\pgfpathlineto{\pgfqpoint{2.386599in}{2.753932in}}%
\pgfpathlineto{\pgfqpoint{2.392925in}{2.892952in}}%
\pgfpathlineto{\pgfqpoint{2.396088in}{2.918696in}}%
\pgfpathlineto{\pgfqpoint{2.399251in}{2.934143in}}%
\pgfpathlineto{\pgfqpoint{2.402414in}{2.918696in}}%
\pgfpathlineto{\pgfqpoint{2.405577in}{2.892952in}}%
\pgfpathlineto{\pgfqpoint{2.408740in}{2.841141in}}%
\pgfpathlineto{\pgfqpoint{2.411903in}{2.753610in}}%
\pgfpathlineto{\pgfqpoint{2.415066in}{2.619740in}}%
\pgfpathlineto{\pgfqpoint{2.418229in}{2.388041in}}%
\pgfpathlineto{\pgfqpoint{2.424554in}{2.475572in}}%
\pgfpathlineto{\pgfqpoint{2.427717in}{2.460125in}}%
\pgfpathlineto{\pgfqpoint{2.430880in}{2.388041in}}%
\pgfpathlineto{\pgfqpoint{2.437206in}{2.418612in}}%
\pgfpathlineto{\pgfqpoint{2.440369in}{2.418612in}}%
\pgfpathlineto{\pgfqpoint{2.443532in}{2.449505in}}%
\pgfpathlineto{\pgfqpoint{2.446695in}{2.387719in}}%
\pgfpathlineto{\pgfqpoint{2.453021in}{2.449505in}}%
\pgfpathlineto{\pgfqpoint{2.456184in}{2.403166in}}%
\pgfpathlineto{\pgfqpoint{2.459347in}{2.402844in}}%
\pgfpathlineto{\pgfqpoint{2.462510in}{2.407993in}}%
\pgfpathlineto{\pgfqpoint{2.465673in}{2.469779in}}%
\pgfpathlineto{\pgfqpoint{2.468836in}{2.464630in}}%
\pgfpathlineto{\pgfqpoint{2.475162in}{2.387397in}}%
\pgfpathlineto{\pgfqpoint{2.478325in}{2.418290in}}%
\pgfpathlineto{\pgfqpoint{2.481488in}{2.387397in}}%
\pgfpathlineto{\pgfqpoint{2.484651in}{2.402522in}}%
\pgfpathlineto{\pgfqpoint{2.487814in}{2.423117in}}%
\pgfpathlineto{\pgfqpoint{2.490977in}{2.433415in}}%
\pgfpathlineto{\pgfqpoint{2.494140in}{2.417969in}}%
\pgfpathlineto{\pgfqpoint{2.497303in}{2.423117in}}%
\pgfpathlineto{\pgfqpoint{2.503628in}{2.387075in}}%
\pgfpathlineto{\pgfqpoint{2.506791in}{2.402522in}}%
\pgfpathlineto{\pgfqpoint{2.509954in}{2.386753in}}%
\pgfpathlineto{\pgfqpoint{2.513117in}{2.402200in}}%
\pgfpathlineto{\pgfqpoint{2.516280in}{2.402200in}}%
\pgfpathlineto{\pgfqpoint{2.519443in}{2.386753in}}%
\pgfpathlineto{\pgfqpoint{2.522606in}{2.386753in}}%
\pgfpathlineto{\pgfqpoint{2.528932in}{2.448540in}}%
\pgfpathlineto{\pgfqpoint{2.535258in}{2.685388in}}%
\pgfpathlineto{\pgfqpoint{2.538421in}{2.746853in}}%
\pgfpathlineto{\pgfqpoint{2.541584in}{2.782895in}}%
\pgfpathlineto{\pgfqpoint{2.544747in}{2.798341in}}%
\pgfpathlineto{\pgfqpoint{2.547910in}{2.793192in}}%
\pgfpathlineto{\pgfqpoint{2.551073in}{2.772597in}}%
\pgfpathlineto{\pgfqpoint{2.554236in}{2.726257in}}%
\pgfpathlineto{\pgfqpoint{2.557399in}{2.654173in}}%
\pgfpathlineto{\pgfqpoint{2.563725in}{2.386110in}}%
\pgfpathlineto{\pgfqpoint{2.566888in}{2.453045in}}%
\pgfpathlineto{\pgfqpoint{2.570051in}{2.463343in}}%
\pgfpathlineto{\pgfqpoint{2.573214in}{2.442747in}}%
\pgfpathlineto{\pgfqpoint{2.576377in}{2.370663in}}%
\pgfpathlineto{\pgfqpoint{2.579540in}{2.391259in}}%
\pgfpathlineto{\pgfqpoint{2.582702in}{2.386110in}}%
\pgfpathlineto{\pgfqpoint{2.585865in}{2.386110in}}%
\pgfpathlineto{\pgfqpoint{2.589028in}{2.370341in}}%
\pgfpathlineto{\pgfqpoint{2.592191in}{2.385788in}}%
\pgfpathlineto{\pgfqpoint{2.595354in}{2.437277in}}%
\pgfpathlineto{\pgfqpoint{2.598517in}{2.447575in}}%
\pgfpathlineto{\pgfqpoint{2.601680in}{2.401235in}}%
\pgfpathlineto{\pgfqpoint{2.604843in}{2.452723in}}%
\pgfpathlineto{\pgfqpoint{2.608006in}{2.375490in}}%
\pgfpathlineto{\pgfqpoint{2.611169in}{2.385788in}}%
\pgfpathlineto{\pgfqpoint{2.614332in}{2.385466in}}%
\pgfpathlineto{\pgfqpoint{2.617495in}{2.400913in}}%
\pgfpathlineto{\pgfqpoint{2.620658in}{2.385466in}}%
\pgfpathlineto{\pgfqpoint{2.623821in}{2.416359in}}%
\pgfpathlineto{\pgfqpoint{2.626984in}{2.416359in}}%
\pgfpathlineto{\pgfqpoint{2.630147in}{2.385466in}}%
\pgfpathlineto{\pgfqpoint{2.633310in}{2.421508in}}%
\pgfpathlineto{\pgfqpoint{2.636473in}{2.385466in}}%
\pgfpathlineto{\pgfqpoint{2.639636in}{2.385144in}}%
\pgfpathlineto{\pgfqpoint{2.642799in}{2.421187in}}%
\pgfpathlineto{\pgfqpoint{2.645962in}{2.374847in}}%
\pgfpathlineto{\pgfqpoint{2.649125in}{2.400591in}}%
\pgfpathlineto{\pgfqpoint{2.652288in}{2.385144in}}%
\pgfpathlineto{\pgfqpoint{2.655451in}{2.400591in}}%
\pgfpathlineto{\pgfqpoint{2.658614in}{2.400591in}}%
\pgfpathlineto{\pgfqpoint{2.661776in}{2.385144in}}%
\pgfpathlineto{\pgfqpoint{2.664939in}{2.400269in}}%
\pgfpathlineto{\pgfqpoint{2.668102in}{2.353929in}}%
\pgfpathlineto{\pgfqpoint{2.671265in}{2.369376in}}%
\pgfpathlineto{\pgfqpoint{2.674428in}{2.446609in}}%
\pgfpathlineto{\pgfqpoint{2.677591in}{2.431162in}}%
\pgfpathlineto{\pgfqpoint{2.683917in}{2.585629in}}%
\pgfpathlineto{\pgfqpoint{2.690243in}{2.657391in}}%
\pgfpathlineto{\pgfqpoint{2.693406in}{2.657391in}}%
\pgfpathlineto{\pgfqpoint{2.696569in}{2.636795in}}%
\pgfpathlineto{\pgfqpoint{2.702895in}{2.518371in}}%
\pgfpathlineto{\pgfqpoint{2.706058in}{2.399947in}}%
\pgfpathlineto{\pgfqpoint{2.709221in}{2.451436in}}%
\pgfpathlineto{\pgfqpoint{2.712384in}{2.384179in}}%
\pgfpathlineto{\pgfqpoint{2.715547in}{2.445966in}}%
\pgfpathlineto{\pgfqpoint{2.718710in}{2.415072in}}%
\pgfpathlineto{\pgfqpoint{2.721873in}{2.445966in}}%
\pgfpathlineto{\pgfqpoint{2.725036in}{2.384179in}}%
\pgfpathlineto{\pgfqpoint{2.728199in}{2.358435in}}%
\pgfpathlineto{\pgfqpoint{2.731362in}{2.451114in}}%
\pgfpathlineto{\pgfqpoint{2.734525in}{2.337839in}}%
\pgfpathlineto{\pgfqpoint{2.737688in}{2.414750in}}%
\pgfpathlineto{\pgfqpoint{2.740851in}{2.394155in}}%
\pgfpathlineto{\pgfqpoint{2.744013in}{2.435346in}}%
\pgfpathlineto{\pgfqpoint{2.747176in}{2.383857in}}%
\pgfpathlineto{\pgfqpoint{2.750339in}{2.404453in}}%
\pgfpathlineto{\pgfqpoint{2.753502in}{2.389006in}}%
\pgfpathlineto{\pgfqpoint{2.756665in}{2.358113in}}%
\pgfpathlineto{\pgfqpoint{2.759828in}{2.368411in}}%
\pgfpathlineto{\pgfqpoint{2.762991in}{2.373238in}}%
\pgfpathlineto{\pgfqpoint{2.766154in}{2.476215in}}%
\pgfpathlineto{\pgfqpoint{2.772480in}{2.388684in}}%
\pgfpathlineto{\pgfqpoint{2.775643in}{2.383535in}}%
\pgfpathlineto{\pgfqpoint{2.778806in}{2.409280in}}%
\pgfpathlineto{\pgfqpoint{2.781969in}{2.409280in}}%
\pgfpathlineto{\pgfqpoint{2.785132in}{2.383535in}}%
\pgfpathlineto{\pgfqpoint{2.788295in}{2.383214in}}%
\pgfpathlineto{\pgfqpoint{2.791458in}{2.352320in}}%
\pgfpathlineto{\pgfqpoint{2.794621in}{2.357469in}}%
\pgfpathlineto{\pgfqpoint{2.800947in}{2.378065in}}%
\pgfpathlineto{\pgfqpoint{2.804110in}{2.398660in}}%
\pgfpathlineto{\pgfqpoint{2.807273in}{2.336874in}}%
\pgfpathlineto{\pgfqpoint{2.810436in}{2.434702in}}%
\pgfpathlineto{\pgfqpoint{2.813599in}{2.382892in}}%
\pgfpathlineto{\pgfqpoint{2.816762in}{2.351999in}}%
\pgfpathlineto{\pgfqpoint{2.819925in}{2.429232in}}%
\pgfpathlineto{\pgfqpoint{2.823087in}{2.398338in}}%
\pgfpathlineto{\pgfqpoint{2.826250in}{2.398338in}}%
\pgfpathlineto{\pgfqpoint{2.829413in}{2.372594in}}%
\pgfpathlineto{\pgfqpoint{2.832576in}{2.521911in}}%
\pgfpathlineto{\pgfqpoint{2.835739in}{2.506465in}}%
\pgfpathlineto{\pgfqpoint{2.838902in}{2.526738in}}%
\pgfpathlineto{\pgfqpoint{2.842065in}{2.459803in}}%
\pgfpathlineto{\pgfqpoint{2.845228in}{2.506143in}}%
\pgfpathlineto{\pgfqpoint{2.848391in}{2.454654in}}%
\pgfpathlineto{\pgfqpoint{2.851554in}{2.351677in}}%
\pgfpathlineto{\pgfqpoint{2.854717in}{2.387719in}}%
\pgfpathlineto{\pgfqpoint{2.857880in}{2.398017in}}%
\pgfpathlineto{\pgfqpoint{2.861043in}{2.336230in}}%
\pgfpathlineto{\pgfqpoint{2.864206in}{2.335908in}}%
\pgfpathlineto{\pgfqpoint{2.867369in}{2.418290in}}%
\pgfpathlineto{\pgfqpoint{2.867369in}{2.418290in}}%
\pgfusepath{stroke}%
\end{pgfscope}%
\begin{pgfscope}%
\pgfsetrectcap%
\pgfsetmiterjoin%
\pgfsetlinewidth{0.803000pt}%
\definecolor{currentstroke}{rgb}{0.000000,0.000000,0.000000}%
\pgfsetstrokecolor{currentstroke}%
\pgfsetdash{}{0pt}%
\pgfpathmoveto{\pgfqpoint{0.758026in}{2.240123in}}%
\pgfpathlineto{\pgfqpoint{0.758026in}{3.150926in}}%
\pgfusepath{stroke}%
\end{pgfscope}%
\begin{pgfscope}%
\pgfsetrectcap%
\pgfsetmiterjoin%
\pgfsetlinewidth{0.803000pt}%
\definecolor{currentstroke}{rgb}{0.000000,0.000000,0.000000}%
\pgfsetstrokecolor{currentstroke}%
\pgfsetdash{}{0pt}%
\pgfpathmoveto{\pgfqpoint{2.866666in}{2.240123in}}%
\pgfpathlineto{\pgfqpoint{2.866666in}{3.150926in}}%
\pgfusepath{stroke}%
\end{pgfscope}%
\begin{pgfscope}%
\pgfsetrectcap%
\pgfsetmiterjoin%
\pgfsetlinewidth{0.803000pt}%
\definecolor{currentstroke}{rgb}{0.000000,0.000000,0.000000}%
\pgfsetstrokecolor{currentstroke}%
\pgfsetdash{}{0pt}%
\pgfpathmoveto{\pgfqpoint{0.758026in}{2.240123in}}%
\pgfpathlineto{\pgfqpoint{2.866666in}{2.240123in}}%
\pgfusepath{stroke}%
\end{pgfscope}%
\begin{pgfscope}%
\pgfsetrectcap%
\pgfsetmiterjoin%
\pgfsetlinewidth{0.803000pt}%
\definecolor{currentstroke}{rgb}{0.000000,0.000000,0.000000}%
\pgfsetstrokecolor{currentstroke}%
\pgfsetdash{}{0pt}%
\pgfpathmoveto{\pgfqpoint{0.758026in}{3.150926in}}%
\pgfpathlineto{\pgfqpoint{2.866666in}{3.150926in}}%
\pgfusepath{stroke}%
\end{pgfscope}%
\begin{pgfscope}%
\definecolor{textcolor}{rgb}{0.000000,0.000000,0.000000}%
\pgfsetstrokecolor{textcolor}%
\pgfsetfillcolor{textcolor}%
\pgftext[x=1.812346in,y=3.234260in,,base]{\color{textcolor}\rmfamily\fontsize{12.000000}{14.400000}\selectfont 41,597 kHz}%
\end{pgfscope}%
\begin{pgfscope}%
\pgfsetbuttcap%
\pgfsetmiterjoin%
\definecolor{currentfill}{rgb}{1.000000,1.000000,1.000000}%
\pgfsetfillcolor{currentfill}%
\pgfsetlinewidth{0.000000pt}%
\definecolor{currentstroke}{rgb}{0.000000,0.000000,0.000000}%
\pgfsetstrokecolor{currentstroke}%
\pgfsetstrokeopacity{0.000000}%
\pgfsetdash{}{0pt}%
\pgfpathmoveto{\pgfqpoint{3.833026in}{2.240123in}}%
\pgfpathlineto{\pgfqpoint{5.941666in}{2.240123in}}%
\pgfpathlineto{\pgfqpoint{5.941666in}{3.150926in}}%
\pgfpathlineto{\pgfqpoint{3.833026in}{3.150926in}}%
\pgfpathlineto{\pgfqpoint{3.833026in}{2.240123in}}%
\pgfpathclose%
\pgfusepath{fill}%
\end{pgfscope}%
\begin{pgfscope}%
\pgfpathrectangle{\pgfqpoint{3.833026in}{2.240123in}}{\pgfqpoint{2.108640in}{0.910803in}}%
\pgfusepath{clip}%
\pgfsetbuttcap%
\pgfsetroundjoin%
\pgfsetlinewidth{0.501875pt}%
\definecolor{currentstroke}{rgb}{0.690196,0.690196,0.690196}%
\pgfsetstrokecolor{currentstroke}%
\pgfsetstrokeopacity{0.700000}%
\pgfsetdash{{1.850000pt}{0.800000pt}}{0.000000pt}%
\pgfpathmoveto{\pgfqpoint{3.833026in}{2.240123in}}%
\pgfpathlineto{\pgfqpoint{3.833026in}{3.150926in}}%
\pgfusepath{stroke}%
\end{pgfscope}%
\begin{pgfscope}%
\pgfsetbuttcap%
\pgfsetroundjoin%
\definecolor{currentfill}{rgb}{0.000000,0.000000,0.000000}%
\pgfsetfillcolor{currentfill}%
\pgfsetlinewidth{0.803000pt}%
\definecolor{currentstroke}{rgb}{0.000000,0.000000,0.000000}%
\pgfsetstrokecolor{currentstroke}%
\pgfsetdash{}{0pt}%
\pgfsys@defobject{currentmarker}{\pgfqpoint{0.000000in}{-0.048611in}}{\pgfqpoint{0.000000in}{0.000000in}}{%
\pgfpathmoveto{\pgfqpoint{0.000000in}{0.000000in}}%
\pgfpathlineto{\pgfqpoint{0.000000in}{-0.048611in}}%
\pgfusepath{stroke,fill}%
}%
\begin{pgfscope}%
\pgfsys@transformshift{3.833026in}{2.240123in}%
\pgfsys@useobject{currentmarker}{}%
\end{pgfscope}%
\end{pgfscope}%
\begin{pgfscope}%
\definecolor{textcolor}{rgb}{0.000000,0.000000,0.000000}%
\pgfsetstrokecolor{textcolor}%
\pgfsetfillcolor{textcolor}%
\pgftext[x=3.833026in,y=2.142901in,,top]{\color{textcolor}\rmfamily\fontsize{10.000000}{12.000000}\selectfont \(\displaystyle {300000}\)}%
\end{pgfscope}%
\begin{pgfscope}%
\pgfpathrectangle{\pgfqpoint{3.833026in}{2.240123in}}{\pgfqpoint{2.108640in}{0.910803in}}%
\pgfusepath{clip}%
\pgfsetbuttcap%
\pgfsetroundjoin%
\pgfsetlinewidth{0.501875pt}%
\definecolor{currentstroke}{rgb}{0.690196,0.690196,0.690196}%
\pgfsetstrokecolor{currentstroke}%
\pgfsetstrokeopacity{0.700000}%
\pgfsetdash{{1.850000pt}{0.800000pt}}{0.000000pt}%
\pgfpathmoveto{\pgfqpoint{4.360186in}{2.240123in}}%
\pgfpathlineto{\pgfqpoint{4.360186in}{3.150926in}}%
\pgfusepath{stroke}%
\end{pgfscope}%
\begin{pgfscope}%
\pgfsetbuttcap%
\pgfsetroundjoin%
\definecolor{currentfill}{rgb}{0.000000,0.000000,0.000000}%
\pgfsetfillcolor{currentfill}%
\pgfsetlinewidth{0.803000pt}%
\definecolor{currentstroke}{rgb}{0.000000,0.000000,0.000000}%
\pgfsetstrokecolor{currentstroke}%
\pgfsetdash{}{0pt}%
\pgfsys@defobject{currentmarker}{\pgfqpoint{0.000000in}{-0.048611in}}{\pgfqpoint{0.000000in}{0.000000in}}{%
\pgfpathmoveto{\pgfqpoint{0.000000in}{0.000000in}}%
\pgfpathlineto{\pgfqpoint{0.000000in}{-0.048611in}}%
\pgfusepath{stroke,fill}%
}%
\begin{pgfscope}%
\pgfsys@transformshift{4.360186in}{2.240123in}%
\pgfsys@useobject{currentmarker}{}%
\end{pgfscope}%
\end{pgfscope}%
\begin{pgfscope}%
\definecolor{textcolor}{rgb}{0.000000,0.000000,0.000000}%
\pgfsetstrokecolor{textcolor}%
\pgfsetfillcolor{textcolor}%
\pgftext[x=4.360186in,y=2.142901in,,top]{\color{textcolor}\rmfamily\fontsize{10.000000}{12.000000}\selectfont \(\displaystyle {400000}\)}%
\end{pgfscope}%
\begin{pgfscope}%
\pgfpathrectangle{\pgfqpoint{3.833026in}{2.240123in}}{\pgfqpoint{2.108640in}{0.910803in}}%
\pgfusepath{clip}%
\pgfsetbuttcap%
\pgfsetroundjoin%
\pgfsetlinewidth{0.501875pt}%
\definecolor{currentstroke}{rgb}{0.690196,0.690196,0.690196}%
\pgfsetstrokecolor{currentstroke}%
\pgfsetstrokeopacity{0.700000}%
\pgfsetdash{{1.850000pt}{0.800000pt}}{0.000000pt}%
\pgfpathmoveto{\pgfqpoint{4.887346in}{2.240123in}}%
\pgfpathlineto{\pgfqpoint{4.887346in}{3.150926in}}%
\pgfusepath{stroke}%
\end{pgfscope}%
\begin{pgfscope}%
\pgfsetbuttcap%
\pgfsetroundjoin%
\definecolor{currentfill}{rgb}{0.000000,0.000000,0.000000}%
\pgfsetfillcolor{currentfill}%
\pgfsetlinewidth{0.803000pt}%
\definecolor{currentstroke}{rgb}{0.000000,0.000000,0.000000}%
\pgfsetstrokecolor{currentstroke}%
\pgfsetdash{}{0pt}%
\pgfsys@defobject{currentmarker}{\pgfqpoint{0.000000in}{-0.048611in}}{\pgfqpoint{0.000000in}{0.000000in}}{%
\pgfpathmoveto{\pgfqpoint{0.000000in}{0.000000in}}%
\pgfpathlineto{\pgfqpoint{0.000000in}{-0.048611in}}%
\pgfusepath{stroke,fill}%
}%
\begin{pgfscope}%
\pgfsys@transformshift{4.887346in}{2.240123in}%
\pgfsys@useobject{currentmarker}{}%
\end{pgfscope}%
\end{pgfscope}%
\begin{pgfscope}%
\definecolor{textcolor}{rgb}{0.000000,0.000000,0.000000}%
\pgfsetstrokecolor{textcolor}%
\pgfsetfillcolor{textcolor}%
\pgftext[x=4.887346in,y=2.142901in,,top]{\color{textcolor}\rmfamily\fontsize{10.000000}{12.000000}\selectfont \(\displaystyle {500000}\)}%
\end{pgfscope}%
\begin{pgfscope}%
\pgfpathrectangle{\pgfqpoint{3.833026in}{2.240123in}}{\pgfqpoint{2.108640in}{0.910803in}}%
\pgfusepath{clip}%
\pgfsetbuttcap%
\pgfsetroundjoin%
\pgfsetlinewidth{0.501875pt}%
\definecolor{currentstroke}{rgb}{0.690196,0.690196,0.690196}%
\pgfsetstrokecolor{currentstroke}%
\pgfsetstrokeopacity{0.700000}%
\pgfsetdash{{1.850000pt}{0.800000pt}}{0.000000pt}%
\pgfpathmoveto{\pgfqpoint{5.414506in}{2.240123in}}%
\pgfpathlineto{\pgfqpoint{5.414506in}{3.150926in}}%
\pgfusepath{stroke}%
\end{pgfscope}%
\begin{pgfscope}%
\pgfsetbuttcap%
\pgfsetroundjoin%
\definecolor{currentfill}{rgb}{0.000000,0.000000,0.000000}%
\pgfsetfillcolor{currentfill}%
\pgfsetlinewidth{0.803000pt}%
\definecolor{currentstroke}{rgb}{0.000000,0.000000,0.000000}%
\pgfsetstrokecolor{currentstroke}%
\pgfsetdash{}{0pt}%
\pgfsys@defobject{currentmarker}{\pgfqpoint{0.000000in}{-0.048611in}}{\pgfqpoint{0.000000in}{0.000000in}}{%
\pgfpathmoveto{\pgfqpoint{0.000000in}{0.000000in}}%
\pgfpathlineto{\pgfqpoint{0.000000in}{-0.048611in}}%
\pgfusepath{stroke,fill}%
}%
\begin{pgfscope}%
\pgfsys@transformshift{5.414506in}{2.240123in}%
\pgfsys@useobject{currentmarker}{}%
\end{pgfscope}%
\end{pgfscope}%
\begin{pgfscope}%
\definecolor{textcolor}{rgb}{0.000000,0.000000,0.000000}%
\pgfsetstrokecolor{textcolor}%
\pgfsetfillcolor{textcolor}%
\pgftext[x=5.414506in,y=2.142901in,,top]{\color{textcolor}\rmfamily\fontsize{10.000000}{12.000000}\selectfont \(\displaystyle {600000}\)}%
\end{pgfscope}%
\begin{pgfscope}%
\pgfpathrectangle{\pgfqpoint{3.833026in}{2.240123in}}{\pgfqpoint{2.108640in}{0.910803in}}%
\pgfusepath{clip}%
\pgfsetbuttcap%
\pgfsetroundjoin%
\pgfsetlinewidth{0.501875pt}%
\definecolor{currentstroke}{rgb}{0.690196,0.690196,0.690196}%
\pgfsetstrokecolor{currentstroke}%
\pgfsetstrokeopacity{0.700000}%
\pgfsetdash{{1.850000pt}{0.800000pt}}{0.000000pt}%
\pgfpathmoveto{\pgfqpoint{5.941666in}{2.240123in}}%
\pgfpathlineto{\pgfqpoint{5.941666in}{3.150926in}}%
\pgfusepath{stroke}%
\end{pgfscope}%
\begin{pgfscope}%
\pgfsetbuttcap%
\pgfsetroundjoin%
\definecolor{currentfill}{rgb}{0.000000,0.000000,0.000000}%
\pgfsetfillcolor{currentfill}%
\pgfsetlinewidth{0.803000pt}%
\definecolor{currentstroke}{rgb}{0.000000,0.000000,0.000000}%
\pgfsetstrokecolor{currentstroke}%
\pgfsetdash{}{0pt}%
\pgfsys@defobject{currentmarker}{\pgfqpoint{0.000000in}{-0.048611in}}{\pgfqpoint{0.000000in}{0.000000in}}{%
\pgfpathmoveto{\pgfqpoint{0.000000in}{0.000000in}}%
\pgfpathlineto{\pgfqpoint{0.000000in}{-0.048611in}}%
\pgfusepath{stroke,fill}%
}%
\begin{pgfscope}%
\pgfsys@transformshift{5.941666in}{2.240123in}%
\pgfsys@useobject{currentmarker}{}%
\end{pgfscope}%
\end{pgfscope}%
\begin{pgfscope}%
\definecolor{textcolor}{rgb}{0.000000,0.000000,0.000000}%
\pgfsetstrokecolor{textcolor}%
\pgfsetfillcolor{textcolor}%
\pgftext[x=5.941666in,y=2.142901in,,top]{\color{textcolor}\rmfamily\fontsize{10.000000}{12.000000}\selectfont \(\displaystyle {700000}\)}%
\end{pgfscope}%
\begin{pgfscope}%
\definecolor{textcolor}{rgb}{0.000000,0.000000,0.000000}%
\pgfsetstrokecolor{textcolor}%
\pgfsetfillcolor{textcolor}%
\pgftext[x=4.887346in,y=1.963889in,,top]{\color{textcolor}\rmfamily\fontsize{10.000000}{12.000000}\selectfont Frequency (Hz)}%
\end{pgfscope}%
\begin{pgfscope}%
\pgfpathrectangle{\pgfqpoint{3.833026in}{2.240123in}}{\pgfqpoint{2.108640in}{0.910803in}}%
\pgfusepath{clip}%
\pgfsetbuttcap%
\pgfsetroundjoin%
\pgfsetlinewidth{0.501875pt}%
\definecolor{currentstroke}{rgb}{0.690196,0.690196,0.690196}%
\pgfsetstrokecolor{currentstroke}%
\pgfsetstrokeopacity{0.700000}%
\pgfsetdash{{1.850000pt}{0.800000pt}}{0.000000pt}%
\pgfpathmoveto{\pgfqpoint{3.833026in}{2.358532in}}%
\pgfpathlineto{\pgfqpoint{5.941666in}{2.358532in}}%
\pgfusepath{stroke}%
\end{pgfscope}%
\begin{pgfscope}%
\pgfsetbuttcap%
\pgfsetroundjoin%
\definecolor{currentfill}{rgb}{0.000000,0.000000,0.000000}%
\pgfsetfillcolor{currentfill}%
\pgfsetlinewidth{0.803000pt}%
\definecolor{currentstroke}{rgb}{0.000000,0.000000,0.000000}%
\pgfsetstrokecolor{currentstroke}%
\pgfsetdash{}{0pt}%
\pgfsys@defobject{currentmarker}{\pgfqpoint{-0.048611in}{0.000000in}}{\pgfqpoint{-0.000000in}{0.000000in}}{%
\pgfpathmoveto{\pgfqpoint{-0.000000in}{0.000000in}}%
\pgfpathlineto{\pgfqpoint{-0.048611in}{0.000000in}}%
\pgfusepath{stroke,fill}%
}%
\begin{pgfscope}%
\pgfsys@transformshift{3.833026in}{2.358532in}%
\pgfsys@useobject{currentmarker}{}%
\end{pgfscope}%
\end{pgfscope}%
\begin{pgfscope}%
\definecolor{textcolor}{rgb}{0.000000,0.000000,0.000000}%
\pgfsetstrokecolor{textcolor}%
\pgfsetfillcolor{textcolor}%
\pgftext[x=3.419444in, y=2.310306in, left, base]{\color{textcolor}\rmfamily\fontsize{10.000000}{12.000000}\selectfont \(\displaystyle {\ensuremath{-}100}\)}%
\end{pgfscope}%
\begin{pgfscope}%
\pgfpathrectangle{\pgfqpoint{3.833026in}{2.240123in}}{\pgfqpoint{2.108640in}{0.910803in}}%
\pgfusepath{clip}%
\pgfsetbuttcap%
\pgfsetroundjoin%
\pgfsetlinewidth{0.501875pt}%
\definecolor{currentstroke}{rgb}{0.690196,0.690196,0.690196}%
\pgfsetstrokecolor{currentstroke}%
\pgfsetstrokeopacity{0.700000}%
\pgfsetdash{{1.850000pt}{0.800000pt}}{0.000000pt}%
\pgfpathmoveto{\pgfqpoint{3.833026in}{2.885083in}}%
\pgfpathlineto{\pgfqpoint{5.941666in}{2.885083in}}%
\pgfusepath{stroke}%
\end{pgfscope}%
\begin{pgfscope}%
\pgfsetbuttcap%
\pgfsetroundjoin%
\definecolor{currentfill}{rgb}{0.000000,0.000000,0.000000}%
\pgfsetfillcolor{currentfill}%
\pgfsetlinewidth{0.803000pt}%
\definecolor{currentstroke}{rgb}{0.000000,0.000000,0.000000}%
\pgfsetstrokecolor{currentstroke}%
\pgfsetdash{}{0pt}%
\pgfsys@defobject{currentmarker}{\pgfqpoint{-0.048611in}{0.000000in}}{\pgfqpoint{-0.000000in}{0.000000in}}{%
\pgfpathmoveto{\pgfqpoint{-0.000000in}{0.000000in}}%
\pgfpathlineto{\pgfqpoint{-0.048611in}{0.000000in}}%
\pgfusepath{stroke,fill}%
}%
\begin{pgfscope}%
\pgfsys@transformshift{3.833026in}{2.885083in}%
\pgfsys@useobject{currentmarker}{}%
\end{pgfscope}%
\end{pgfscope}%
\begin{pgfscope}%
\definecolor{textcolor}{rgb}{0.000000,0.000000,0.000000}%
\pgfsetstrokecolor{textcolor}%
\pgfsetfillcolor{textcolor}%
\pgftext[x=3.488889in, y=2.836858in, left, base]{\color{textcolor}\rmfamily\fontsize{10.000000}{12.000000}\selectfont \(\displaystyle {\ensuremath{-}50}\)}%
\end{pgfscope}%
\begin{pgfscope}%
\definecolor{textcolor}{rgb}{0.000000,0.000000,0.000000}%
\pgfsetstrokecolor{textcolor}%
\pgfsetfillcolor{textcolor}%
\pgftext[x=3.363889in,y=2.695525in,,bottom,rotate=90.000000]{\color{textcolor}\rmfamily\fontsize{10.000000}{12.000000}\selectfont Amplitude (dB)}%
\end{pgfscope}%
\begin{pgfscope}%
\pgfpathrectangle{\pgfqpoint{3.833026in}{2.240123in}}{\pgfqpoint{2.108640in}{0.910803in}}%
\pgfusepath{clip}%
\pgfsetrectcap%
\pgfsetroundjoin%
\pgfsetlinewidth{1.003750pt}%
\definecolor{currentstroke}{rgb}{0.000000,0.000000,0.000000}%
\pgfsetstrokecolor{currentstroke}%
\pgfsetdash{}{0pt}%
\pgfpathmoveto{\pgfqpoint{3.831971in}{2.527357in}}%
\pgfpathlineto{\pgfqpoint{3.836716in}{2.427312in}}%
\pgfpathlineto{\pgfqpoint{3.841460in}{2.558950in}}%
\pgfpathlineto{\pgfqpoint{3.846205in}{2.537559in}}%
\pgfpathlineto{\pgfqpoint{3.850949in}{2.584949in}}%
\pgfpathlineto{\pgfqpoint{3.855694in}{2.611277in}}%
\pgfpathlineto{\pgfqpoint{3.860438in}{2.590214in}}%
\pgfpathlineto{\pgfqpoint{3.865182in}{2.563887in}}%
\pgfpathlineto{\pgfqpoint{3.869927in}{2.527028in}}%
\pgfpathlineto{\pgfqpoint{3.874671in}{2.590214in}}%
\pgfpathlineto{\pgfqpoint{3.879416in}{2.548090in}}%
\pgfpathlineto{\pgfqpoint{3.884160in}{2.574089in}}%
\pgfpathlineto{\pgfqpoint{3.888905in}{2.568823in}}%
\pgfpathlineto{\pgfqpoint{3.893649in}{2.553027in}}%
\pgfpathlineto{\pgfqpoint{3.898394in}{2.568823in}}%
\pgfpathlineto{\pgfqpoint{3.903138in}{2.637275in}}%
\pgfpathlineto{\pgfqpoint{3.907882in}{2.610947in}}%
\pgfpathlineto{\pgfqpoint{3.912627in}{2.705727in}}%
\pgfpathlineto{\pgfqpoint{3.917371in}{2.710992in}}%
\pgfpathlineto{\pgfqpoint{3.922116in}{2.773849in}}%
\pgfpathlineto{\pgfqpoint{3.926860in}{2.784380in}}%
\pgfpathlineto{\pgfqpoint{3.931605in}{2.763318in}}%
\pgfpathlineto{\pgfqpoint{3.936349in}{2.726460in}}%
\pgfpathlineto{\pgfqpoint{3.945838in}{2.584291in}}%
\pgfpathlineto{\pgfqpoint{3.950582in}{2.584291in}}%
\pgfpathlineto{\pgfqpoint{3.955327in}{2.536901in}}%
\pgfpathlineto{\pgfqpoint{3.960071in}{2.578696in}}%
\pgfpathlineto{\pgfqpoint{3.964816in}{2.547103in}}%
\pgfpathlineto{\pgfqpoint{3.969560in}{2.552369in}}%
\pgfpathlineto{\pgfqpoint{3.974305in}{2.562900in}}%
\pgfpathlineto{\pgfqpoint{3.979049in}{2.531306in}}%
\pgfpathlineto{\pgfqpoint{3.983793in}{2.531306in}}%
\pgfpathlineto{\pgfqpoint{3.988538in}{2.594493in}}%
\pgfpathlineto{\pgfqpoint{3.993282in}{2.599758in}}%
\pgfpathlineto{\pgfqpoint{4.002771in}{2.541508in}}%
\pgfpathlineto{\pgfqpoint{4.007516in}{2.567836in}}%
\pgfpathlineto{\pgfqpoint{4.012260in}{2.530977in}}%
\pgfpathlineto{\pgfqpoint{4.017005in}{2.567836in}}%
\pgfpathlineto{\pgfqpoint{4.021749in}{2.494119in}}%
\pgfpathlineto{\pgfqpoint{4.031238in}{2.604695in}}%
\pgfpathlineto{\pgfqpoint{4.035982in}{2.599100in}}%
\pgfpathlineto{\pgfqpoint{4.040727in}{2.625428in}}%
\pgfpathlineto{\pgfqpoint{4.045471in}{2.736003in}}%
\pgfpathlineto{\pgfqpoint{4.050216in}{2.809721in}}%
\pgfpathlineto{\pgfqpoint{4.054960in}{2.862376in}}%
\pgfpathlineto{\pgfqpoint{4.059705in}{2.883438in}}%
\pgfpathlineto{\pgfqpoint{4.064449in}{2.888703in}}%
\pgfpathlineto{\pgfqpoint{4.069193in}{2.867641in}}%
\pgfpathlineto{\pgfqpoint{4.073938in}{2.835719in}}%
\pgfpathlineto{\pgfqpoint{4.078682in}{2.783064in}}%
\pgfpathlineto{\pgfqpoint{4.083427in}{2.677754in}}%
\pgfpathlineto{\pgfqpoint{4.092916in}{2.546116in}}%
\pgfpathlineto{\pgfqpoint{4.097660in}{2.498726in}}%
\pgfpathlineto{\pgfqpoint{4.102404in}{2.561912in}}%
\pgfpathlineto{\pgfqpoint{4.107149in}{2.572443in}}%
\pgfpathlineto{\pgfqpoint{4.111893in}{2.529990in}}%
\pgfpathlineto{\pgfqpoint{4.116638in}{2.635300in}}%
\pgfpathlineto{\pgfqpoint{4.121382in}{2.582645in}}%
\pgfpathlineto{\pgfqpoint{4.126127in}{2.624769in}}%
\pgfpathlineto{\pgfqpoint{4.130871in}{2.577380in}}%
\pgfpathlineto{\pgfqpoint{4.135616in}{2.551052in}}%
\pgfpathlineto{\pgfqpoint{4.140360in}{2.503662in}}%
\pgfpathlineto{\pgfqpoint{4.145104in}{2.577380in}}%
\pgfpathlineto{\pgfqpoint{4.149849in}{2.508599in}}%
\pgfpathlineto{\pgfqpoint{4.154593in}{2.503333in}}%
\pgfpathlineto{\pgfqpoint{4.159338in}{2.540192in}}%
\pgfpathlineto{\pgfqpoint{4.164082in}{2.540192in}}%
\pgfpathlineto{\pgfqpoint{4.173571in}{2.634971in}}%
\pgfpathlineto{\pgfqpoint{4.178316in}{2.719220in}}%
\pgfpathlineto{\pgfqpoint{4.183060in}{2.829795in}}%
\pgfpathlineto{\pgfqpoint{4.187804in}{2.908449in}}%
\pgfpathlineto{\pgfqpoint{4.192549in}{2.955839in}}%
\pgfpathlineto{\pgfqpoint{4.197293in}{2.982166in}}%
\pgfpathlineto{\pgfqpoint{4.202038in}{2.987432in}}%
\pgfpathlineto{\pgfqpoint{4.206782in}{2.971635in}}%
\pgfpathlineto{\pgfqpoint{4.211527in}{2.934777in}}%
\pgfpathlineto{\pgfqpoint{4.216271in}{2.876856in}}%
\pgfpathlineto{\pgfqpoint{4.221016in}{2.787342in}}%
\pgfpathlineto{\pgfqpoint{4.230504in}{2.481613in}}%
\pgfpathlineto{\pgfqpoint{4.235249in}{2.544799in}}%
\pgfpathlineto{\pgfqpoint{4.239993in}{2.634313in}}%
\pgfpathlineto{\pgfqpoint{4.244738in}{2.539534in}}%
\pgfpathlineto{\pgfqpoint{4.249482in}{2.571127in}}%
\pgfpathlineto{\pgfqpoint{4.254227in}{2.529003in}}%
\pgfpathlineto{\pgfqpoint{4.258971in}{2.586923in}}%
\pgfpathlineto{\pgfqpoint{4.263715in}{2.481284in}}%
\pgfpathlineto{\pgfqpoint{4.268460in}{2.555001in}}%
\pgfpathlineto{\pgfqpoint{4.273204in}{2.597125in}}%
\pgfpathlineto{\pgfqpoint{4.282693in}{2.618188in}}%
\pgfpathlineto{\pgfqpoint{4.287438in}{2.549736in}}%
\pgfpathlineto{\pgfqpoint{4.292182in}{2.612922in}}%
\pgfpathlineto{\pgfqpoint{4.296927in}{2.523408in}}%
\pgfpathlineto{\pgfqpoint{4.301671in}{2.649452in}}%
\pgfpathlineto{\pgfqpoint{4.306415in}{2.596796in}}%
\pgfpathlineto{\pgfqpoint{4.311160in}{2.559938in}}%
\pgfpathlineto{\pgfqpoint{4.315904in}{2.775824in}}%
\pgfpathlineto{\pgfqpoint{4.320649in}{2.902196in}}%
\pgfpathlineto{\pgfqpoint{4.325393in}{2.981179in}}%
\pgfpathlineto{\pgfqpoint{4.330138in}{3.028569in}}%
\pgfpathlineto{\pgfqpoint{4.334882in}{3.054896in}}%
\pgfpathlineto{\pgfqpoint{4.339627in}{3.059833in}}%
\pgfpathlineto{\pgfqpoint{4.344371in}{3.049302in}}%
\pgfpathlineto{\pgfqpoint{4.349115in}{3.012443in}}%
\pgfpathlineto{\pgfqpoint{4.353860in}{2.954522in}}%
\pgfpathlineto{\pgfqpoint{4.358604in}{2.865009in}}%
\pgfpathlineto{\pgfqpoint{4.363349in}{2.722840in}}%
\pgfpathlineto{\pgfqpoint{4.368093in}{2.680716in}}%
\pgfpathlineto{\pgfqpoint{4.372838in}{2.659653in}}%
\pgfpathlineto{\pgfqpoint{4.377582in}{2.559609in}}%
\pgfpathlineto{\pgfqpoint{4.382327in}{2.611935in}}%
\pgfpathlineto{\pgfqpoint{4.387071in}{2.606669in}}%
\pgfpathlineto{\pgfqpoint{4.396560in}{2.522421in}}%
\pgfpathlineto{\pgfqpoint{4.401304in}{2.564545in}}%
\pgfpathlineto{\pgfqpoint{4.406049in}{2.596138in}}%
\pgfpathlineto{\pgfqpoint{4.410793in}{2.564545in}}%
\pgfpathlineto{\pgfqpoint{4.415538in}{2.690917in}}%
\pgfpathlineto{\pgfqpoint{4.420282in}{2.522092in}}%
\pgfpathlineto{\pgfqpoint{4.425026in}{2.553685in}}%
\pgfpathlineto{\pgfqpoint{4.429771in}{2.574747in}}%
\pgfpathlineto{\pgfqpoint{4.434515in}{2.537888in}}%
\pgfpathlineto{\pgfqpoint{4.439260in}{2.580012in}}%
\pgfpathlineto{\pgfqpoint{4.444004in}{2.569481in}}%
\pgfpathlineto{\pgfqpoint{4.448749in}{2.537888in}}%
\pgfpathlineto{\pgfqpoint{4.453493in}{2.822226in}}%
\pgfpathlineto{\pgfqpoint{4.458238in}{2.943004in}}%
\pgfpathlineto{\pgfqpoint{4.462982in}{3.021987in}}%
\pgfpathlineto{\pgfqpoint{4.467726in}{3.069377in}}%
\pgfpathlineto{\pgfqpoint{4.472471in}{3.100970in}}%
\pgfpathlineto{\pgfqpoint{4.477215in}{3.106235in}}%
\pgfpathlineto{\pgfqpoint{4.481960in}{3.095704in}}%
\pgfpathlineto{\pgfqpoint{4.486704in}{3.058846in}}%
\pgfpathlineto{\pgfqpoint{4.491449in}{3.000925in}}%
\pgfpathlineto{\pgfqpoint{4.496193in}{2.916347in}}%
\pgfpathlineto{\pgfqpoint{4.500938in}{2.774178in}}%
\pgfpathlineto{\pgfqpoint{4.505682in}{2.526699in}}%
\pgfpathlineto{\pgfqpoint{4.515171in}{2.537230in}}%
\pgfpathlineto{\pgfqpoint{4.519915in}{2.500372in}}%
\pgfpathlineto{\pgfqpoint{4.524660in}{2.521434in}}%
\pgfpathlineto{\pgfqpoint{4.529404in}{2.516168in}}%
\pgfpathlineto{\pgfqpoint{4.534149in}{2.552698in}}%
\pgfpathlineto{\pgfqpoint{4.538893in}{2.536901in}}%
\pgfpathlineto{\pgfqpoint{4.543637in}{2.526370in}}%
\pgfpathlineto{\pgfqpoint{4.553126in}{2.547432in}}%
\pgfpathlineto{\pgfqpoint{4.557871in}{2.500042in}}%
\pgfpathlineto{\pgfqpoint{4.562615in}{2.515839in}}%
\pgfpathlineto{\pgfqpoint{4.567360in}{2.526370in}}%
\pgfpathlineto{\pgfqpoint{4.572104in}{2.526041in}}%
\pgfpathlineto{\pgfqpoint{4.576849in}{2.536572in}}%
\pgfpathlineto{\pgfqpoint{4.581593in}{2.515510in}}%
\pgfpathlineto{\pgfqpoint{4.586337in}{2.568165in}}%
\pgfpathlineto{\pgfqpoint{4.591082in}{2.815644in}}%
\pgfpathlineto{\pgfqpoint{4.595826in}{2.936751in}}%
\pgfpathlineto{\pgfqpoint{4.600571in}{3.021000in}}%
\pgfpathlineto{\pgfqpoint{4.605315in}{3.068060in}}%
\pgfpathlineto{\pgfqpoint{4.610060in}{3.094388in}}%
\pgfpathlineto{\pgfqpoint{4.614804in}{3.104919in}}%
\pgfpathlineto{\pgfqpoint{4.619549in}{3.089122in}}%
\pgfpathlineto{\pgfqpoint{4.624293in}{3.057529in}}%
\pgfpathlineto{\pgfqpoint{4.629037in}{3.004874in}}%
\pgfpathlineto{\pgfqpoint{4.633782in}{2.915360in}}%
\pgfpathlineto{\pgfqpoint{4.638526in}{2.767926in}}%
\pgfpathlineto{\pgfqpoint{4.643271in}{2.525383in}}%
\pgfpathlineto{\pgfqpoint{4.648015in}{2.556976in}}%
\pgfpathlineto{\pgfqpoint{4.652760in}{2.551710in}}%
\pgfpathlineto{\pgfqpoint{4.657504in}{2.556976in}}%
\pgfpathlineto{\pgfqpoint{4.666993in}{2.504321in}}%
\pgfpathlineto{\pgfqpoint{4.671737in}{2.499055in}}%
\pgfpathlineto{\pgfqpoint{4.676482in}{2.551710in}}%
\pgfpathlineto{\pgfqpoint{4.681226in}{2.556647in}}%
\pgfpathlineto{\pgfqpoint{4.685971in}{2.446071in}}%
\pgfpathlineto{\pgfqpoint{4.690715in}{2.514523in}}%
\pgfpathlineto{\pgfqpoint{4.695460in}{2.488195in}}%
\pgfpathlineto{\pgfqpoint{4.700204in}{2.472398in}}%
\pgfpathlineto{\pgfqpoint{4.704948in}{2.467133in}}%
\pgfpathlineto{\pgfqpoint{4.709693in}{2.519788in}}%
\pgfpathlineto{\pgfqpoint{4.714437in}{2.503992in}}%
\pgfpathlineto{\pgfqpoint{4.719182in}{2.587911in}}%
\pgfpathlineto{\pgfqpoint{4.723926in}{2.566849in}}%
\pgfpathlineto{\pgfqpoint{4.728671in}{2.540521in}}%
\pgfpathlineto{\pgfqpoint{4.733415in}{2.587911in}}%
\pgfpathlineto{\pgfqpoint{4.738160in}{2.451007in}}%
\pgfpathlineto{\pgfqpoint{4.742904in}{2.572114in}}%
\pgfpathlineto{\pgfqpoint{4.747648in}{2.487866in}}%
\pgfpathlineto{\pgfqpoint{4.752393in}{2.566849in}}%
\pgfpathlineto{\pgfqpoint{4.757137in}{2.608644in}}%
\pgfpathlineto{\pgfqpoint{4.761882in}{2.550723in}}%
\pgfpathlineto{\pgfqpoint{4.766626in}{2.571785in}}%
\pgfpathlineto{\pgfqpoint{4.771371in}{2.598113in}}%
\pgfpathlineto{\pgfqpoint{4.776115in}{2.534927in}}%
\pgfpathlineto{\pgfqpoint{4.780860in}{2.587582in}}%
\pgfpathlineto{\pgfqpoint{4.785604in}{2.519130in}}%
\pgfpathlineto{\pgfqpoint{4.790348in}{2.529661in}}%
\pgfpathlineto{\pgfqpoint{4.795093in}{2.566191in}}%
\pgfpathlineto{\pgfqpoint{4.799837in}{2.571456in}}%
\pgfpathlineto{\pgfqpoint{4.804582in}{2.466146in}}%
\pgfpathlineto{\pgfqpoint{4.809326in}{2.529332in}}%
\pgfpathlineto{\pgfqpoint{4.814071in}{2.487208in}}%
\pgfpathlineto{\pgfqpoint{4.818815in}{2.513535in}}%
\pgfpathlineto{\pgfqpoint{4.823559in}{2.513535in}}%
\pgfpathlineto{\pgfqpoint{4.828304in}{2.503004in}}%
\pgfpathlineto{\pgfqpoint{4.833048in}{2.502675in}}%
\pgfpathlineto{\pgfqpoint{4.837793in}{2.544799in}}%
\pgfpathlineto{\pgfqpoint{4.842537in}{2.481613in}}%
\pgfpathlineto{\pgfqpoint{4.847282in}{2.555330in}}%
\pgfpathlineto{\pgfqpoint{4.852026in}{2.476348in}}%
\pgfpathlineto{\pgfqpoint{4.856771in}{2.560596in}}%
\pgfpathlineto{\pgfqpoint{4.861515in}{2.502675in}}%
\pgfpathlineto{\pgfqpoint{4.866259in}{2.802810in}}%
\pgfpathlineto{\pgfqpoint{4.871004in}{2.934119in}}%
\pgfpathlineto{\pgfqpoint{4.875748in}{3.018367in}}%
\pgfpathlineto{\pgfqpoint{4.880493in}{3.060491in}}%
\pgfpathlineto{\pgfqpoint{4.885237in}{3.092084in}}%
\pgfpathlineto{\pgfqpoint{4.889982in}{3.097350in}}%
\pgfpathlineto{\pgfqpoint{4.894726in}{3.092084in}}%
\pgfpathlineto{\pgfqpoint{4.899471in}{3.060491in}}%
\pgfpathlineto{\pgfqpoint{4.904215in}{3.002570in}}%
\pgfpathlineto{\pgfqpoint{4.908959in}{2.917993in}}%
\pgfpathlineto{\pgfqpoint{4.913704in}{2.775824in}}%
\pgfpathlineto{\pgfqpoint{4.923193in}{2.412503in}}%
\pgfpathlineto{\pgfqpoint{4.927937in}{2.502017in}}%
\pgfpathlineto{\pgfqpoint{4.932682in}{2.475689in}}%
\pgfpathlineto{\pgfqpoint{4.937426in}{2.517814in}}%
\pgfpathlineto{\pgfqpoint{4.942171in}{2.412503in}}%
\pgfpathlineto{\pgfqpoint{4.946915in}{2.464829in}}%
\pgfpathlineto{\pgfqpoint{4.951659in}{2.501688in}}%
\pgfpathlineto{\pgfqpoint{4.956404in}{2.501688in}}%
\pgfpathlineto{\pgfqpoint{4.961148in}{2.480626in}}%
\pgfpathlineto{\pgfqpoint{4.965893in}{2.538547in}}%
\pgfpathlineto{\pgfqpoint{4.970637in}{2.512219in}}%
\pgfpathlineto{\pgfqpoint{4.980126in}{2.554343in}}%
\pgfpathlineto{\pgfqpoint{4.984870in}{2.448704in}}%
\pgfpathlineto{\pgfqpoint{4.989615in}{2.464500in}}%
\pgfpathlineto{\pgfqpoint{4.994359in}{2.554014in}}%
\pgfpathlineto{\pgfqpoint{4.999104in}{2.569811in}}%
\pgfpathlineto{\pgfqpoint{5.003848in}{2.464500in}}%
\pgfpathlineto{\pgfqpoint{5.008593in}{2.538217in}}%
\pgfpathlineto{\pgfqpoint{5.013337in}{2.485562in}}%
\pgfpathlineto{\pgfqpoint{5.018082in}{2.527686in}}%
\pgfpathlineto{\pgfqpoint{5.022826in}{2.527357in}}%
\pgfpathlineto{\pgfqpoint{5.027570in}{2.569481in}}%
\pgfpathlineto{\pgfqpoint{5.032315in}{2.532623in}}%
\pgfpathlineto{\pgfqpoint{5.037059in}{2.522092in}}%
\pgfpathlineto{\pgfqpoint{5.041804in}{2.569481in}}%
\pgfpathlineto{\pgfqpoint{5.046548in}{2.527357in}}%
\pgfpathlineto{\pgfqpoint{5.051293in}{2.553685in}}%
\pgfpathlineto{\pgfqpoint{5.056037in}{2.495764in}}%
\pgfpathlineto{\pgfqpoint{5.060782in}{2.532294in}}%
\pgfpathlineto{\pgfqpoint{5.065526in}{2.579683in}}%
\pgfpathlineto{\pgfqpoint{5.070270in}{2.500701in}}%
\pgfpathlineto{\pgfqpoint{5.075015in}{2.584949in}}%
\pgfpathlineto{\pgfqpoint{5.079759in}{2.484904in}}%
\pgfpathlineto{\pgfqpoint{5.084504in}{2.505966in}}%
\pgfpathlineto{\pgfqpoint{5.089248in}{2.511232in}}%
\pgfpathlineto{\pgfqpoint{5.093993in}{2.527028in}}%
\pgfpathlineto{\pgfqpoint{5.098737in}{2.489841in}}%
\pgfpathlineto{\pgfqpoint{5.108226in}{2.531965in}}%
\pgfpathlineto{\pgfqpoint{5.112970in}{2.589885in}}%
\pgfpathlineto{\pgfqpoint{5.117715in}{2.505637in}}%
\pgfpathlineto{\pgfqpoint{5.122459in}{2.521434in}}%
\pgfpathlineto{\pgfqpoint{5.127204in}{2.505637in}}%
\pgfpathlineto{\pgfqpoint{5.131948in}{2.574089in}}%
\pgfpathlineto{\pgfqpoint{5.136693in}{2.484246in}}%
\pgfpathlineto{\pgfqpoint{5.141437in}{2.800177in}}%
\pgfpathlineto{\pgfqpoint{5.146181in}{2.921284in}}%
\pgfpathlineto{\pgfqpoint{5.150926in}{3.010798in}}%
\pgfpathlineto{\pgfqpoint{5.155670in}{3.063453in}}%
\pgfpathlineto{\pgfqpoint{5.160415in}{3.095046in}}%
\pgfpathlineto{\pgfqpoint{5.165159in}{3.105577in}}%
\pgfpathlineto{\pgfqpoint{5.169904in}{3.095046in}}%
\pgfpathlineto{\pgfqpoint{5.174648in}{3.063124in}}%
\pgfpathlineto{\pgfqpoint{5.179393in}{3.010469in}}%
\pgfpathlineto{\pgfqpoint{5.184137in}{2.926220in}}%
\pgfpathlineto{\pgfqpoint{5.188881in}{2.794582in}}%
\pgfpathlineto{\pgfqpoint{5.193626in}{2.510244in}}%
\pgfpathlineto{\pgfqpoint{5.198370in}{2.541837in}}%
\pgfpathlineto{\pgfqpoint{5.203115in}{2.526041in}}%
\pgfpathlineto{\pgfqpoint{5.207859in}{2.478651in}}%
\pgfpathlineto{\pgfqpoint{5.212604in}{2.599429in}}%
\pgfpathlineto{\pgfqpoint{5.217348in}{2.541508in}}%
\pgfpathlineto{\pgfqpoint{5.222093in}{2.504650in}}%
\pgfpathlineto{\pgfqpoint{5.226837in}{2.441464in}}%
\pgfpathlineto{\pgfqpoint{5.231581in}{2.441464in}}%
\pgfpathlineto{\pgfqpoint{5.236326in}{2.557305in}}%
\pgfpathlineto{\pgfqpoint{5.241070in}{2.536243in}}%
\pgfpathlineto{\pgfqpoint{5.245815in}{2.557305in}}%
\pgfpathlineto{\pgfqpoint{5.250559in}{2.509586in}}%
\pgfpathlineto{\pgfqpoint{5.260048in}{2.551710in}}%
\pgfpathlineto{\pgfqpoint{5.264792in}{2.567507in}}%
\pgfpathlineto{\pgfqpoint{5.269537in}{2.546445in}}%
\pgfpathlineto{\pgfqpoint{5.274281in}{2.556976in}}%
\pgfpathlineto{\pgfqpoint{5.279026in}{2.799190in}}%
\pgfpathlineto{\pgfqpoint{5.283770in}{2.925562in}}%
\pgfpathlineto{\pgfqpoint{5.288515in}{3.014747in}}%
\pgfpathlineto{\pgfqpoint{5.293259in}{3.067402in}}%
\pgfpathlineto{\pgfqpoint{5.298004in}{3.098995in}}%
\pgfpathlineto{\pgfqpoint{5.302748in}{3.109526in}}%
\pgfpathlineto{\pgfqpoint{5.307492in}{3.098995in}}%
\pgfpathlineto{\pgfqpoint{5.312237in}{3.067402in}}%
\pgfpathlineto{\pgfqpoint{5.316981in}{3.014747in}}%
\pgfpathlineto{\pgfqpoint{5.321726in}{2.935764in}}%
\pgfpathlineto{\pgfqpoint{5.326470in}{2.798531in}}%
\pgfpathlineto{\pgfqpoint{5.331215in}{2.524725in}}%
\pgfpathlineto{\pgfqpoint{5.335959in}{2.477335in}}%
\pgfpathlineto{\pgfqpoint{5.340704in}{2.477335in}}%
\pgfpathlineto{\pgfqpoint{5.345448in}{2.445742in}}%
\pgfpathlineto{\pgfqpoint{5.354937in}{2.545787in}}%
\pgfpathlineto{\pgfqpoint{5.359681in}{2.493131in}}%
\pgfpathlineto{\pgfqpoint{5.364426in}{2.513864in}}%
\pgfpathlineto{\pgfqpoint{5.369170in}{2.477006in}}%
\pgfpathlineto{\pgfqpoint{5.373915in}{2.508599in}}%
\pgfpathlineto{\pgfqpoint{5.378659in}{2.498068in}}%
\pgfpathlineto{\pgfqpoint{5.383404in}{2.445413in}}%
\pgfpathlineto{\pgfqpoint{5.388148in}{2.440147in}}%
\pgfpathlineto{\pgfqpoint{5.392892in}{2.508599in}}%
\pgfpathlineto{\pgfqpoint{5.397637in}{2.455944in}}%
\pgfpathlineto{\pgfqpoint{5.402381in}{2.503004in}}%
\pgfpathlineto{\pgfqpoint{5.407126in}{2.508270in}}%
\pgfpathlineto{\pgfqpoint{5.411870in}{2.518801in}}%
\pgfpathlineto{\pgfqpoint{5.416615in}{2.739953in}}%
\pgfpathlineto{\pgfqpoint{5.421359in}{2.882122in}}%
\pgfpathlineto{\pgfqpoint{5.426103in}{2.966370in}}%
\pgfpathlineto{\pgfqpoint{5.430848in}{3.019025in}}%
\pgfpathlineto{\pgfqpoint{5.435592in}{3.050618in}}%
\pgfpathlineto{\pgfqpoint{5.440337in}{3.060820in}}%
\pgfpathlineto{\pgfqpoint{5.445081in}{3.055555in}}%
\pgfpathlineto{\pgfqpoint{5.454570in}{2.976572in}}%
\pgfpathlineto{\pgfqpoint{5.459315in}{2.892324in}}%
\pgfpathlineto{\pgfqpoint{5.464059in}{2.760686in}}%
\pgfpathlineto{\pgfqpoint{5.468803in}{2.486879in}}%
\pgfpathlineto{\pgfqpoint{5.473548in}{2.481613in}}%
\pgfpathlineto{\pgfqpoint{5.478292in}{2.523408in}}%
\pgfpathlineto{\pgfqpoint{5.483037in}{2.423363in}}%
\pgfpathlineto{\pgfqpoint{5.487781in}{2.502346in}}%
\pgfpathlineto{\pgfqpoint{5.492526in}{2.491815in}}%
\pgfpathlineto{\pgfqpoint{5.497270in}{2.539205in}}%
\pgfpathlineto{\pgfqpoint{5.502015in}{2.423363in}}%
\pgfpathlineto{\pgfqpoint{5.506759in}{2.407567in}}%
\pgfpathlineto{\pgfqpoint{5.511503in}{2.460222in}}%
\pgfpathlineto{\pgfqpoint{5.516248in}{2.475689in}}%
\pgfpathlineto{\pgfqpoint{5.520992in}{2.507283in}}%
\pgfpathlineto{\pgfqpoint{5.525737in}{2.496752in}}%
\pgfpathlineto{\pgfqpoint{5.530481in}{2.491486in}}%
\pgfpathlineto{\pgfqpoint{5.535226in}{2.459893in}}%
\pgfpathlineto{\pgfqpoint{5.539970in}{2.438831in}}%
\pgfpathlineto{\pgfqpoint{5.544714in}{2.486220in}}%
\pgfpathlineto{\pgfqpoint{5.549459in}{2.444096in}}%
\pgfpathlineto{\pgfqpoint{5.554203in}{2.654388in}}%
\pgfpathlineto{\pgfqpoint{5.558948in}{2.801822in}}%
\pgfpathlineto{\pgfqpoint{5.563692in}{2.886071in}}%
\pgfpathlineto{\pgfqpoint{5.568437in}{2.943991in}}%
\pgfpathlineto{\pgfqpoint{5.573181in}{2.975585in}}%
\pgfpathlineto{\pgfqpoint{5.577926in}{2.991381in}}%
\pgfpathlineto{\pgfqpoint{5.582670in}{2.980850in}}%
\pgfpathlineto{\pgfqpoint{5.587414in}{2.954522in}}%
\pgfpathlineto{\pgfqpoint{5.592159in}{2.896273in}}%
\pgfpathlineto{\pgfqpoint{5.596903in}{2.822555in}}%
\pgfpathlineto{\pgfqpoint{5.601648in}{2.696183in}}%
\pgfpathlineto{\pgfqpoint{5.606392in}{2.480297in}}%
\pgfpathlineto{\pgfqpoint{5.611137in}{2.443438in}}%
\pgfpathlineto{\pgfqpoint{5.615881in}{2.511890in}}%
\pgfpathlineto{\pgfqpoint{5.625370in}{2.406580in}}%
\pgfpathlineto{\pgfqpoint{5.630114in}{2.443109in}}%
\pgfpathlineto{\pgfqpoint{5.634859in}{2.453640in}}%
\pgfpathlineto{\pgfqpoint{5.639603in}{2.406250in}}%
\pgfpathlineto{\pgfqpoint{5.644348in}{2.427312in}}%
\pgfpathlineto{\pgfqpoint{5.649092in}{2.422047in}}%
\pgfpathlineto{\pgfqpoint{5.653837in}{2.427312in}}%
\pgfpathlineto{\pgfqpoint{5.658581in}{2.469437in}}%
\pgfpathlineto{\pgfqpoint{5.663326in}{2.479968in}}%
\pgfpathlineto{\pgfqpoint{5.668070in}{2.405921in}}%
\pgfpathlineto{\pgfqpoint{5.672814in}{2.437514in}}%
\pgfpathlineto{\pgfqpoint{5.677559in}{2.437514in}}%
\pgfpathlineto{\pgfqpoint{5.682303in}{2.405921in}}%
\pgfpathlineto{\pgfqpoint{5.687048in}{2.426983in}}%
\pgfpathlineto{\pgfqpoint{5.696537in}{2.695525in}}%
\pgfpathlineto{\pgfqpoint{5.701281in}{2.790304in}}%
\pgfpathlineto{\pgfqpoint{5.706025in}{2.847896in}}%
\pgfpathlineto{\pgfqpoint{5.710770in}{2.879489in}}%
\pgfpathlineto{\pgfqpoint{5.715514in}{2.895285in}}%
\pgfpathlineto{\pgfqpoint{5.720259in}{2.884754in}}%
\pgfpathlineto{\pgfqpoint{5.725003in}{2.853161in}}%
\pgfpathlineto{\pgfqpoint{5.729748in}{2.805772in}}%
\pgfpathlineto{\pgfqpoint{5.734492in}{2.732054in}}%
\pgfpathlineto{\pgfqpoint{5.739237in}{2.600416in}}%
\pgfpathlineto{\pgfqpoint{5.743981in}{2.389467in}}%
\pgfpathlineto{\pgfqpoint{5.748725in}{2.468449in}}%
\pgfpathlineto{\pgfqpoint{5.753470in}{2.410529in}}%
\pgfpathlineto{\pgfqpoint{5.758214in}{2.389467in}}%
\pgfpathlineto{\pgfqpoint{5.762959in}{2.463184in}}%
\pgfpathlineto{\pgfqpoint{5.767703in}{2.410529in}}%
\pgfpathlineto{\pgfqpoint{5.772448in}{2.436856in}}%
\pgfpathlineto{\pgfqpoint{5.777192in}{2.426325in}}%
\pgfpathlineto{\pgfqpoint{5.781937in}{2.436527in}}%
\pgfpathlineto{\pgfqpoint{5.786681in}{2.389137in}}%
\pgfpathlineto{\pgfqpoint{5.791425in}{2.441793in}}%
\pgfpathlineto{\pgfqpoint{5.796170in}{2.404934in}}%
\pgfpathlineto{\pgfqpoint{5.800914in}{2.436527in}}%
\pgfpathlineto{\pgfqpoint{5.805659in}{2.494448in}}%
\pgfpathlineto{\pgfqpoint{5.810403in}{2.404934in}}%
\pgfpathlineto{\pgfqpoint{5.815148in}{2.389137in}}%
\pgfpathlineto{\pgfqpoint{5.819892in}{2.441464in}}%
\pgfpathlineto{\pgfqpoint{5.824637in}{2.404605in}}%
\pgfpathlineto{\pgfqpoint{5.829381in}{2.436198in}}%
\pgfpathlineto{\pgfqpoint{5.834125in}{2.573102in}}%
\pgfpathlineto{\pgfqpoint{5.838870in}{2.673146in}}%
\pgfpathlineto{\pgfqpoint{5.843614in}{2.731067in}}%
\pgfpathlineto{\pgfqpoint{5.848359in}{2.762660in}}%
\pgfpathlineto{\pgfqpoint{5.853103in}{2.778457in}}%
\pgfpathlineto{\pgfqpoint{5.857848in}{2.772862in}}%
\pgfpathlineto{\pgfqpoint{5.862592in}{2.741269in}}%
\pgfpathlineto{\pgfqpoint{5.867336in}{2.699145in}}%
\pgfpathlineto{\pgfqpoint{5.872081in}{2.614897in}}%
\pgfpathlineto{\pgfqpoint{5.876825in}{2.488524in}}%
\pgfpathlineto{\pgfqpoint{5.881570in}{2.509586in}}%
\pgfpathlineto{\pgfqpoint{5.886314in}{2.425338in}}%
\pgfpathlineto{\pgfqpoint{5.891059in}{2.388479in}}%
\pgfpathlineto{\pgfqpoint{5.895803in}{2.451336in}}%
\pgfpathlineto{\pgfqpoint{5.900548in}{2.440805in}}%
\pgfpathlineto{\pgfqpoint{5.905292in}{2.451336in}}%
\pgfpathlineto{\pgfqpoint{5.910036in}{2.456602in}}%
\pgfpathlineto{\pgfqpoint{5.914781in}{2.419743in}}%
\pgfpathlineto{\pgfqpoint{5.919525in}{2.425009in}}%
\pgfpathlineto{\pgfqpoint{5.924270in}{2.403947in}}%
\pgfpathlineto{\pgfqpoint{5.929014in}{2.440805in}}%
\pgfpathlineto{\pgfqpoint{5.933759in}{2.419414in}}%
\pgfpathlineto{\pgfqpoint{5.938503in}{2.440476in}}%
\pgfpathlineto{\pgfqpoint{5.943248in}{2.387821in}}%
\pgfpathlineto{\pgfqpoint{5.943248in}{2.387821in}}%
\pgfusepath{stroke}%
\end{pgfscope}%
\begin{pgfscope}%
\pgfsetrectcap%
\pgfsetmiterjoin%
\pgfsetlinewidth{0.803000pt}%
\definecolor{currentstroke}{rgb}{0.000000,0.000000,0.000000}%
\pgfsetstrokecolor{currentstroke}%
\pgfsetdash{}{0pt}%
\pgfpathmoveto{\pgfqpoint{3.833026in}{2.240123in}}%
\pgfpathlineto{\pgfqpoint{3.833026in}{3.150926in}}%
\pgfusepath{stroke}%
\end{pgfscope}%
\begin{pgfscope}%
\pgfsetrectcap%
\pgfsetmiterjoin%
\pgfsetlinewidth{0.803000pt}%
\definecolor{currentstroke}{rgb}{0.000000,0.000000,0.000000}%
\pgfsetstrokecolor{currentstroke}%
\pgfsetdash{}{0pt}%
\pgfpathmoveto{\pgfqpoint{5.941666in}{2.240123in}}%
\pgfpathlineto{\pgfqpoint{5.941666in}{3.150926in}}%
\pgfusepath{stroke}%
\end{pgfscope}%
\begin{pgfscope}%
\pgfsetrectcap%
\pgfsetmiterjoin%
\pgfsetlinewidth{0.803000pt}%
\definecolor{currentstroke}{rgb}{0.000000,0.000000,0.000000}%
\pgfsetstrokecolor{currentstroke}%
\pgfsetdash{}{0pt}%
\pgfpathmoveto{\pgfqpoint{3.833026in}{2.240123in}}%
\pgfpathlineto{\pgfqpoint{5.941666in}{2.240123in}}%
\pgfusepath{stroke}%
\end{pgfscope}%
\begin{pgfscope}%
\pgfsetrectcap%
\pgfsetmiterjoin%
\pgfsetlinewidth{0.803000pt}%
\definecolor{currentstroke}{rgb}{0.000000,0.000000,0.000000}%
\pgfsetstrokecolor{currentstroke}%
\pgfsetdash{}{0pt}%
\pgfpathmoveto{\pgfqpoint{3.833026in}{3.150926in}}%
\pgfpathlineto{\pgfqpoint{5.941666in}{3.150926in}}%
\pgfusepath{stroke}%
\end{pgfscope}%
\begin{pgfscope}%
\definecolor{textcolor}{rgb}{0.000000,0.000000,0.000000}%
\pgfsetstrokecolor{textcolor}%
\pgfsetfillcolor{textcolor}%
\pgftext[x=4.887346in,y=3.234260in,,base]{\color{textcolor}\rmfamily\fontsize{12.000000}{14.400000}\selectfont 26,098 kHz}%
\end{pgfscope}%
\begin{pgfscope}%
\pgfsetbuttcap%
\pgfsetmiterjoin%
\definecolor{currentfill}{rgb}{1.000000,1.000000,1.000000}%
\pgfsetfillcolor{currentfill}%
\pgfsetlinewidth{0.000000pt}%
\definecolor{currentstroke}{rgb}{0.000000,0.000000,0.000000}%
\pgfsetstrokecolor{currentstroke}%
\pgfsetstrokeopacity{0.000000}%
\pgfsetdash{}{0pt}%
\pgfpathmoveto{\pgfqpoint{0.758026in}{0.565123in}}%
\pgfpathlineto{\pgfqpoint{2.866666in}{0.565123in}}%
\pgfpathlineto{\pgfqpoint{2.866666in}{1.475926in}}%
\pgfpathlineto{\pgfqpoint{0.758026in}{1.475926in}}%
\pgfpathlineto{\pgfqpoint{0.758026in}{0.565123in}}%
\pgfpathclose%
\pgfusepath{fill}%
\end{pgfscope}%
\begin{pgfscope}%
\pgfpathrectangle{\pgfqpoint{0.758026in}{0.565123in}}{\pgfqpoint{2.108640in}{0.910803in}}%
\pgfusepath{clip}%
\pgfsetbuttcap%
\pgfsetroundjoin%
\pgfsetlinewidth{0.501875pt}%
\definecolor{currentstroke}{rgb}{0.690196,0.690196,0.690196}%
\pgfsetstrokecolor{currentstroke}%
\pgfsetstrokeopacity{0.700000}%
\pgfsetdash{{1.850000pt}{0.800000pt}}{0.000000pt}%
\pgfpathmoveto{\pgfqpoint{0.758026in}{0.565123in}}%
\pgfpathlineto{\pgfqpoint{0.758026in}{1.475926in}}%
\pgfusepath{stroke}%
\end{pgfscope}%
\begin{pgfscope}%
\pgfsetbuttcap%
\pgfsetroundjoin%
\definecolor{currentfill}{rgb}{0.000000,0.000000,0.000000}%
\pgfsetfillcolor{currentfill}%
\pgfsetlinewidth{0.803000pt}%
\definecolor{currentstroke}{rgb}{0.000000,0.000000,0.000000}%
\pgfsetstrokecolor{currentstroke}%
\pgfsetdash{}{0pt}%
\pgfsys@defobject{currentmarker}{\pgfqpoint{0.000000in}{-0.048611in}}{\pgfqpoint{0.000000in}{0.000000in}}{%
\pgfpathmoveto{\pgfqpoint{0.000000in}{0.000000in}}%
\pgfpathlineto{\pgfqpoint{0.000000in}{-0.048611in}}%
\pgfusepath{stroke,fill}%
}%
\begin{pgfscope}%
\pgfsys@transformshift{0.758026in}{0.565123in}%
\pgfsys@useobject{currentmarker}{}%
\end{pgfscope}%
\end{pgfscope}%
\begin{pgfscope}%
\definecolor{textcolor}{rgb}{0.000000,0.000000,0.000000}%
\pgfsetstrokecolor{textcolor}%
\pgfsetfillcolor{textcolor}%
\pgftext[x=0.758026in,y=0.467901in,,top]{\color{textcolor}\rmfamily\fontsize{10.000000}{12.000000}\selectfont \(\displaystyle {300000}\)}%
\end{pgfscope}%
\begin{pgfscope}%
\pgfpathrectangle{\pgfqpoint{0.758026in}{0.565123in}}{\pgfqpoint{2.108640in}{0.910803in}}%
\pgfusepath{clip}%
\pgfsetbuttcap%
\pgfsetroundjoin%
\pgfsetlinewidth{0.501875pt}%
\definecolor{currentstroke}{rgb}{0.690196,0.690196,0.690196}%
\pgfsetstrokecolor{currentstroke}%
\pgfsetstrokeopacity{0.700000}%
\pgfsetdash{{1.850000pt}{0.800000pt}}{0.000000pt}%
\pgfpathmoveto{\pgfqpoint{1.285186in}{0.565123in}}%
\pgfpathlineto{\pgfqpoint{1.285186in}{1.475926in}}%
\pgfusepath{stroke}%
\end{pgfscope}%
\begin{pgfscope}%
\pgfsetbuttcap%
\pgfsetroundjoin%
\definecolor{currentfill}{rgb}{0.000000,0.000000,0.000000}%
\pgfsetfillcolor{currentfill}%
\pgfsetlinewidth{0.803000pt}%
\definecolor{currentstroke}{rgb}{0.000000,0.000000,0.000000}%
\pgfsetstrokecolor{currentstroke}%
\pgfsetdash{}{0pt}%
\pgfsys@defobject{currentmarker}{\pgfqpoint{0.000000in}{-0.048611in}}{\pgfqpoint{0.000000in}{0.000000in}}{%
\pgfpathmoveto{\pgfqpoint{0.000000in}{0.000000in}}%
\pgfpathlineto{\pgfqpoint{0.000000in}{-0.048611in}}%
\pgfusepath{stroke,fill}%
}%
\begin{pgfscope}%
\pgfsys@transformshift{1.285186in}{0.565123in}%
\pgfsys@useobject{currentmarker}{}%
\end{pgfscope}%
\end{pgfscope}%
\begin{pgfscope}%
\definecolor{textcolor}{rgb}{0.000000,0.000000,0.000000}%
\pgfsetstrokecolor{textcolor}%
\pgfsetfillcolor{textcolor}%
\pgftext[x=1.285186in,y=0.467901in,,top]{\color{textcolor}\rmfamily\fontsize{10.000000}{12.000000}\selectfont \(\displaystyle {400000}\)}%
\end{pgfscope}%
\begin{pgfscope}%
\pgfpathrectangle{\pgfqpoint{0.758026in}{0.565123in}}{\pgfqpoint{2.108640in}{0.910803in}}%
\pgfusepath{clip}%
\pgfsetbuttcap%
\pgfsetroundjoin%
\pgfsetlinewidth{0.501875pt}%
\definecolor{currentstroke}{rgb}{0.690196,0.690196,0.690196}%
\pgfsetstrokecolor{currentstroke}%
\pgfsetstrokeopacity{0.700000}%
\pgfsetdash{{1.850000pt}{0.800000pt}}{0.000000pt}%
\pgfpathmoveto{\pgfqpoint{1.812346in}{0.565123in}}%
\pgfpathlineto{\pgfqpoint{1.812346in}{1.475926in}}%
\pgfusepath{stroke}%
\end{pgfscope}%
\begin{pgfscope}%
\pgfsetbuttcap%
\pgfsetroundjoin%
\definecolor{currentfill}{rgb}{0.000000,0.000000,0.000000}%
\pgfsetfillcolor{currentfill}%
\pgfsetlinewidth{0.803000pt}%
\definecolor{currentstroke}{rgb}{0.000000,0.000000,0.000000}%
\pgfsetstrokecolor{currentstroke}%
\pgfsetdash{}{0pt}%
\pgfsys@defobject{currentmarker}{\pgfqpoint{0.000000in}{-0.048611in}}{\pgfqpoint{0.000000in}{0.000000in}}{%
\pgfpathmoveto{\pgfqpoint{0.000000in}{0.000000in}}%
\pgfpathlineto{\pgfqpoint{0.000000in}{-0.048611in}}%
\pgfusepath{stroke,fill}%
}%
\begin{pgfscope}%
\pgfsys@transformshift{1.812346in}{0.565123in}%
\pgfsys@useobject{currentmarker}{}%
\end{pgfscope}%
\end{pgfscope}%
\begin{pgfscope}%
\definecolor{textcolor}{rgb}{0.000000,0.000000,0.000000}%
\pgfsetstrokecolor{textcolor}%
\pgfsetfillcolor{textcolor}%
\pgftext[x=1.812346in,y=0.467901in,,top]{\color{textcolor}\rmfamily\fontsize{10.000000}{12.000000}\selectfont \(\displaystyle {500000}\)}%
\end{pgfscope}%
\begin{pgfscope}%
\pgfpathrectangle{\pgfqpoint{0.758026in}{0.565123in}}{\pgfqpoint{2.108640in}{0.910803in}}%
\pgfusepath{clip}%
\pgfsetbuttcap%
\pgfsetroundjoin%
\pgfsetlinewidth{0.501875pt}%
\definecolor{currentstroke}{rgb}{0.690196,0.690196,0.690196}%
\pgfsetstrokecolor{currentstroke}%
\pgfsetstrokeopacity{0.700000}%
\pgfsetdash{{1.850000pt}{0.800000pt}}{0.000000pt}%
\pgfpathmoveto{\pgfqpoint{2.339506in}{0.565123in}}%
\pgfpathlineto{\pgfqpoint{2.339506in}{1.475926in}}%
\pgfusepath{stroke}%
\end{pgfscope}%
\begin{pgfscope}%
\pgfsetbuttcap%
\pgfsetroundjoin%
\definecolor{currentfill}{rgb}{0.000000,0.000000,0.000000}%
\pgfsetfillcolor{currentfill}%
\pgfsetlinewidth{0.803000pt}%
\definecolor{currentstroke}{rgb}{0.000000,0.000000,0.000000}%
\pgfsetstrokecolor{currentstroke}%
\pgfsetdash{}{0pt}%
\pgfsys@defobject{currentmarker}{\pgfqpoint{0.000000in}{-0.048611in}}{\pgfqpoint{0.000000in}{0.000000in}}{%
\pgfpathmoveto{\pgfqpoint{0.000000in}{0.000000in}}%
\pgfpathlineto{\pgfqpoint{0.000000in}{-0.048611in}}%
\pgfusepath{stroke,fill}%
}%
\begin{pgfscope}%
\pgfsys@transformshift{2.339506in}{0.565123in}%
\pgfsys@useobject{currentmarker}{}%
\end{pgfscope}%
\end{pgfscope}%
\begin{pgfscope}%
\definecolor{textcolor}{rgb}{0.000000,0.000000,0.000000}%
\pgfsetstrokecolor{textcolor}%
\pgfsetfillcolor{textcolor}%
\pgftext[x=2.339506in,y=0.467901in,,top]{\color{textcolor}\rmfamily\fontsize{10.000000}{12.000000}\selectfont \(\displaystyle {600000}\)}%
\end{pgfscope}%
\begin{pgfscope}%
\pgfpathrectangle{\pgfqpoint{0.758026in}{0.565123in}}{\pgfqpoint{2.108640in}{0.910803in}}%
\pgfusepath{clip}%
\pgfsetbuttcap%
\pgfsetroundjoin%
\pgfsetlinewidth{0.501875pt}%
\definecolor{currentstroke}{rgb}{0.690196,0.690196,0.690196}%
\pgfsetstrokecolor{currentstroke}%
\pgfsetstrokeopacity{0.700000}%
\pgfsetdash{{1.850000pt}{0.800000pt}}{0.000000pt}%
\pgfpathmoveto{\pgfqpoint{2.866666in}{0.565123in}}%
\pgfpathlineto{\pgfqpoint{2.866666in}{1.475926in}}%
\pgfusepath{stroke}%
\end{pgfscope}%
\begin{pgfscope}%
\pgfsetbuttcap%
\pgfsetroundjoin%
\definecolor{currentfill}{rgb}{0.000000,0.000000,0.000000}%
\pgfsetfillcolor{currentfill}%
\pgfsetlinewidth{0.803000pt}%
\definecolor{currentstroke}{rgb}{0.000000,0.000000,0.000000}%
\pgfsetstrokecolor{currentstroke}%
\pgfsetdash{}{0pt}%
\pgfsys@defobject{currentmarker}{\pgfqpoint{0.000000in}{-0.048611in}}{\pgfqpoint{0.000000in}{0.000000in}}{%
\pgfpathmoveto{\pgfqpoint{0.000000in}{0.000000in}}%
\pgfpathlineto{\pgfqpoint{0.000000in}{-0.048611in}}%
\pgfusepath{stroke,fill}%
}%
\begin{pgfscope}%
\pgfsys@transformshift{2.866666in}{0.565123in}%
\pgfsys@useobject{currentmarker}{}%
\end{pgfscope}%
\end{pgfscope}%
\begin{pgfscope}%
\definecolor{textcolor}{rgb}{0.000000,0.000000,0.000000}%
\pgfsetstrokecolor{textcolor}%
\pgfsetfillcolor{textcolor}%
\pgftext[x=2.866666in,y=0.467901in,,top]{\color{textcolor}\rmfamily\fontsize{10.000000}{12.000000}\selectfont \(\displaystyle {700000}\)}%
\end{pgfscope}%
\begin{pgfscope}%
\definecolor{textcolor}{rgb}{0.000000,0.000000,0.000000}%
\pgfsetstrokecolor{textcolor}%
\pgfsetfillcolor{textcolor}%
\pgftext[x=1.812346in,y=0.288889in,,top]{\color{textcolor}\rmfamily\fontsize{10.000000}{12.000000}\selectfont Frequency (Hz)}%
\end{pgfscope}%
\begin{pgfscope}%
\pgfpathrectangle{\pgfqpoint{0.758026in}{0.565123in}}{\pgfqpoint{2.108640in}{0.910803in}}%
\pgfusepath{clip}%
\pgfsetbuttcap%
\pgfsetroundjoin%
\pgfsetlinewidth{0.501875pt}%
\definecolor{currentstroke}{rgb}{0.690196,0.690196,0.690196}%
\pgfsetstrokecolor{currentstroke}%
\pgfsetstrokeopacity{0.700000}%
\pgfsetdash{{1.850000pt}{0.800000pt}}{0.000000pt}%
\pgfpathmoveto{\pgfqpoint{0.758026in}{0.683646in}}%
\pgfpathlineto{\pgfqpoint{2.866666in}{0.683646in}}%
\pgfusepath{stroke}%
\end{pgfscope}%
\begin{pgfscope}%
\pgfsetbuttcap%
\pgfsetroundjoin%
\definecolor{currentfill}{rgb}{0.000000,0.000000,0.000000}%
\pgfsetfillcolor{currentfill}%
\pgfsetlinewidth{0.803000pt}%
\definecolor{currentstroke}{rgb}{0.000000,0.000000,0.000000}%
\pgfsetstrokecolor{currentstroke}%
\pgfsetdash{}{0pt}%
\pgfsys@defobject{currentmarker}{\pgfqpoint{-0.048611in}{0.000000in}}{\pgfqpoint{-0.000000in}{0.000000in}}{%
\pgfpathmoveto{\pgfqpoint{-0.000000in}{0.000000in}}%
\pgfpathlineto{\pgfqpoint{-0.048611in}{0.000000in}}%
\pgfusepath{stroke,fill}%
}%
\begin{pgfscope}%
\pgfsys@transformshift{0.758026in}{0.683646in}%
\pgfsys@useobject{currentmarker}{}%
\end{pgfscope}%
\end{pgfscope}%
\begin{pgfscope}%
\definecolor{textcolor}{rgb}{0.000000,0.000000,0.000000}%
\pgfsetstrokecolor{textcolor}%
\pgfsetfillcolor{textcolor}%
\pgftext[x=0.344444in, y=0.635421in, left, base]{\color{textcolor}\rmfamily\fontsize{10.000000}{12.000000}\selectfont \(\displaystyle {\ensuremath{-}100}\)}%
\end{pgfscope}%
\begin{pgfscope}%
\pgfpathrectangle{\pgfqpoint{0.758026in}{0.565123in}}{\pgfqpoint{2.108640in}{0.910803in}}%
\pgfusepath{clip}%
\pgfsetbuttcap%
\pgfsetroundjoin%
\pgfsetlinewidth{0.501875pt}%
\definecolor{currentstroke}{rgb}{0.690196,0.690196,0.690196}%
\pgfsetstrokecolor{currentstroke}%
\pgfsetstrokeopacity{0.700000}%
\pgfsetdash{{1.850000pt}{0.800000pt}}{0.000000pt}%
\pgfpathmoveto{\pgfqpoint{0.758026in}{1.208737in}}%
\pgfpathlineto{\pgfqpoint{2.866666in}{1.208737in}}%
\pgfusepath{stroke}%
\end{pgfscope}%
\begin{pgfscope}%
\pgfsetbuttcap%
\pgfsetroundjoin%
\definecolor{currentfill}{rgb}{0.000000,0.000000,0.000000}%
\pgfsetfillcolor{currentfill}%
\pgfsetlinewidth{0.803000pt}%
\definecolor{currentstroke}{rgb}{0.000000,0.000000,0.000000}%
\pgfsetstrokecolor{currentstroke}%
\pgfsetdash{}{0pt}%
\pgfsys@defobject{currentmarker}{\pgfqpoint{-0.048611in}{0.000000in}}{\pgfqpoint{-0.000000in}{0.000000in}}{%
\pgfpathmoveto{\pgfqpoint{-0.000000in}{0.000000in}}%
\pgfpathlineto{\pgfqpoint{-0.048611in}{0.000000in}}%
\pgfusepath{stroke,fill}%
}%
\begin{pgfscope}%
\pgfsys@transformshift{0.758026in}{1.208737in}%
\pgfsys@useobject{currentmarker}{}%
\end{pgfscope}%
\end{pgfscope}%
\begin{pgfscope}%
\definecolor{textcolor}{rgb}{0.000000,0.000000,0.000000}%
\pgfsetstrokecolor{textcolor}%
\pgfsetfillcolor{textcolor}%
\pgftext[x=0.413889in, y=1.160512in, left, base]{\color{textcolor}\rmfamily\fontsize{10.000000}{12.000000}\selectfont \(\displaystyle {\ensuremath{-}50}\)}%
\end{pgfscope}%
\begin{pgfscope}%
\definecolor{textcolor}{rgb}{0.000000,0.000000,0.000000}%
\pgfsetstrokecolor{textcolor}%
\pgfsetfillcolor{textcolor}%
\pgftext[x=0.288889in,y=1.020525in,,bottom,rotate=90.000000]{\color{textcolor}\rmfamily\fontsize{10.000000}{12.000000}\selectfont Amplitude (dB)}%
\end{pgfscope}%
\begin{pgfscope}%
\pgfpathrectangle{\pgfqpoint{0.758026in}{0.565123in}}{\pgfqpoint{2.108640in}{0.910803in}}%
\pgfusepath{clip}%
\pgfsetrectcap%
\pgfsetroundjoin%
\pgfsetlinewidth{1.003750pt}%
\definecolor{currentstroke}{rgb}{0.000000,0.000000,0.000000}%
\pgfsetstrokecolor{currentstroke}%
\pgfsetdash{}{0pt}%
\pgfpathmoveto{\pgfqpoint{0.756971in}{0.873007in}}%
\pgfpathlineto{\pgfqpoint{0.761716in}{0.883509in}}%
\pgfpathlineto{\pgfqpoint{0.766460in}{0.899262in}}%
\pgfpathlineto{\pgfqpoint{0.771205in}{0.909435in}}%
\pgfpathlineto{\pgfqpoint{0.775949in}{0.904184in}}%
\pgfpathlineto{\pgfqpoint{0.780694in}{0.967195in}}%
\pgfpathlineto{\pgfqpoint{0.785438in}{0.946192in}}%
\pgfpathlineto{\pgfqpoint{0.790182in}{0.919937in}}%
\pgfpathlineto{\pgfqpoint{0.794927in}{0.883181in}}%
\pgfpathlineto{\pgfqpoint{0.799671in}{0.909435in}}%
\pgfpathlineto{\pgfqpoint{0.804416in}{0.830672in}}%
\pgfpathlineto{\pgfqpoint{0.809160in}{0.909107in}}%
\pgfpathlineto{\pgfqpoint{0.813905in}{0.851347in}}%
\pgfpathlineto{\pgfqpoint{0.818649in}{0.914358in}}%
\pgfpathlineto{\pgfqpoint{0.823394in}{0.903856in}}%
\pgfpathlineto{\pgfqpoint{0.828138in}{0.846096in}}%
\pgfpathlineto{\pgfqpoint{0.832882in}{0.919609in}}%
\pgfpathlineto{\pgfqpoint{0.837627in}{0.945863in}}%
\pgfpathlineto{\pgfqpoint{0.842371in}{0.882853in}}%
\pgfpathlineto{\pgfqpoint{0.847116in}{0.887775in}}%
\pgfpathlineto{\pgfqpoint{0.851860in}{0.898277in}}%
\pgfpathlineto{\pgfqpoint{0.856605in}{0.835266in}}%
\pgfpathlineto{\pgfqpoint{0.861349in}{0.861521in}}%
\pgfpathlineto{\pgfqpoint{0.866094in}{0.840517in}}%
\pgfpathlineto{\pgfqpoint{0.870838in}{0.893026in}}%
\pgfpathlineto{\pgfqpoint{0.875582in}{0.982292in}}%
\pgfpathlineto{\pgfqpoint{0.889816in}{1.060727in}}%
\pgfpathlineto{\pgfqpoint{0.894560in}{1.060727in}}%
\pgfpathlineto{\pgfqpoint{0.899305in}{1.023971in}}%
\pgfpathlineto{\pgfqpoint{0.913538in}{0.845440in}}%
\pgfpathlineto{\pgfqpoint{0.918282in}{0.840189in}}%
\pgfpathlineto{\pgfqpoint{0.923027in}{0.818857in}}%
\pgfpathlineto{\pgfqpoint{0.927771in}{0.866115in}}%
\pgfpathlineto{\pgfqpoint{0.932516in}{0.860864in}}%
\pgfpathlineto{\pgfqpoint{0.937260in}{0.923875in}}%
\pgfpathlineto{\pgfqpoint{0.942005in}{0.845112in}}%
\pgfpathlineto{\pgfqpoint{0.946749in}{0.929126in}}%
\pgfpathlineto{\pgfqpoint{0.951493in}{0.913373in}}%
\pgfpathlineto{\pgfqpoint{0.956238in}{0.944879in}}%
\pgfpathlineto{\pgfqpoint{0.960982in}{0.913045in}}%
\pgfpathlineto{\pgfqpoint{0.965727in}{0.928798in}}%
\pgfpathlineto{\pgfqpoint{0.970471in}{0.918296in}}%
\pgfpathlineto{\pgfqpoint{0.975216in}{1.023314in}}%
\pgfpathlineto{\pgfqpoint{0.979960in}{1.091576in}}%
\pgfpathlineto{\pgfqpoint{0.984705in}{1.138834in}}%
\pgfpathlineto{\pgfqpoint{0.989449in}{1.159838in}}%
\pgfpathlineto{\pgfqpoint{0.994193in}{1.175591in}}%
\pgfpathlineto{\pgfqpoint{1.003682in}{1.112252in}}%
\pgfpathlineto{\pgfqpoint{1.008427in}{1.054492in}}%
\pgfpathlineto{\pgfqpoint{1.017916in}{0.828703in}}%
\pgfpathlineto{\pgfqpoint{1.022660in}{0.812950in}}%
\pgfpathlineto{\pgfqpoint{1.027404in}{0.912717in}}%
\pgfpathlineto{\pgfqpoint{1.032149in}{0.938972in}}%
\pgfpathlineto{\pgfqpoint{1.036893in}{0.833625in}}%
\pgfpathlineto{\pgfqpoint{1.041638in}{0.901887in}}%
\pgfpathlineto{\pgfqpoint{1.046382in}{0.912389in}}%
\pgfpathlineto{\pgfqpoint{1.051127in}{0.880883in}}%
\pgfpathlineto{\pgfqpoint{1.055871in}{0.938643in}}%
\pgfpathlineto{\pgfqpoint{1.060616in}{0.896636in}}%
\pgfpathlineto{\pgfqpoint{1.065360in}{0.917640in}}%
\pgfpathlineto{\pgfqpoint{1.070104in}{0.901887in}}%
\pgfpathlineto{\pgfqpoint{1.074849in}{1.048584in}}%
\pgfpathlineto{\pgfqpoint{1.079593in}{1.148352in}}%
\pgfpathlineto{\pgfqpoint{1.084338in}{1.211363in}}%
\pgfpathlineto{\pgfqpoint{1.089082in}{1.242868in}}%
\pgfpathlineto{\pgfqpoint{1.093827in}{1.258621in}}%
\pgfpathlineto{\pgfqpoint{1.098571in}{1.258621in}}%
\pgfpathlineto{\pgfqpoint{1.103316in}{1.232366in}}%
\pgfpathlineto{\pgfqpoint{1.108060in}{1.190359in}}%
\pgfpathlineto{\pgfqpoint{1.112804in}{1.116518in}}%
\pgfpathlineto{\pgfqpoint{1.117549in}{1.000998in}}%
\pgfpathlineto{\pgfqpoint{1.122293in}{0.916983in}}%
\pgfpathlineto{\pgfqpoint{1.127038in}{0.927485in}}%
\pgfpathlineto{\pgfqpoint{1.131782in}{0.869725in}}%
\pgfpathlineto{\pgfqpoint{1.136527in}{0.880227in}}%
\pgfpathlineto{\pgfqpoint{1.141271in}{0.848722in}}%
\pgfpathlineto{\pgfqpoint{1.146016in}{0.874976in}}%
\pgfpathlineto{\pgfqpoint{1.150760in}{0.869397in}}%
\pgfpathlineto{\pgfqpoint{1.155504in}{0.853644in}}%
\pgfpathlineto{\pgfqpoint{1.160249in}{0.916655in}}%
\pgfpathlineto{\pgfqpoint{1.164993in}{0.911404in}}%
\pgfpathlineto{\pgfqpoint{1.169738in}{0.953412in}}%
\pgfpathlineto{\pgfqpoint{1.174482in}{1.037426in}}%
\pgfpathlineto{\pgfqpoint{1.179227in}{1.168699in}}%
\pgfpathlineto{\pgfqpoint{1.183971in}{1.252713in}}%
\pgfpathlineto{\pgfqpoint{1.188715in}{1.304894in}}%
\pgfpathlineto{\pgfqpoint{1.193460in}{1.331149in}}%
\pgfpathlineto{\pgfqpoint{1.198204in}{1.341651in}}%
\pgfpathlineto{\pgfqpoint{1.202949in}{1.331149in}}%
\pgfpathlineto{\pgfqpoint{1.207693in}{1.299643in}}%
\pgfpathlineto{\pgfqpoint{1.212438in}{1.247134in}}%
\pgfpathlineto{\pgfqpoint{1.217182in}{1.157869in}}%
\pgfpathlineto{\pgfqpoint{1.221927in}{1.037098in}}%
\pgfpathlineto{\pgfqpoint{1.226671in}{0.952755in}}%
\pgfpathlineto{\pgfqpoint{1.236160in}{0.905497in}}%
\pgfpathlineto{\pgfqpoint{1.240904in}{0.831984in}}%
\pgfpathlineto{\pgfqpoint{1.245649in}{0.900246in}}%
\pgfpathlineto{\pgfqpoint{1.250393in}{0.810981in}}%
\pgfpathlineto{\pgfqpoint{1.255138in}{0.926501in}}%
\pgfpathlineto{\pgfqpoint{1.259882in}{0.921250in}}%
\pgfpathlineto{\pgfqpoint{1.264627in}{0.894667in}}%
\pgfpathlineto{\pgfqpoint{1.269371in}{0.941925in}}%
\pgfpathlineto{\pgfqpoint{1.274115in}{0.957678in}}%
\pgfpathlineto{\pgfqpoint{1.278860in}{1.151962in}}%
\pgfpathlineto{\pgfqpoint{1.283604in}{1.262231in}}%
\pgfpathlineto{\pgfqpoint{1.288349in}{1.330492in}}%
\pgfpathlineto{\pgfqpoint{1.293093in}{1.377751in}}%
\pgfpathlineto{\pgfqpoint{1.297838in}{1.398754in}}%
\pgfpathlineto{\pgfqpoint{1.302582in}{1.398754in}}%
\pgfpathlineto{\pgfqpoint{1.307327in}{1.387924in}}%
\pgfpathlineto{\pgfqpoint{1.312071in}{1.345917in}}%
\pgfpathlineto{\pgfqpoint{1.316815in}{1.282906in}}%
\pgfpathlineto{\pgfqpoint{1.321560in}{1.188390in}}%
\pgfpathlineto{\pgfqpoint{1.331049in}{0.894339in}}%
\pgfpathlineto{\pgfqpoint{1.335793in}{0.857583in}}%
\pgfpathlineto{\pgfqpoint{1.340538in}{0.894339in}}%
\pgfpathlineto{\pgfqpoint{1.345282in}{0.883509in}}%
\pgfpathlineto{\pgfqpoint{1.350026in}{0.857254in}}%
\pgfpathlineto{\pgfqpoint{1.354771in}{0.888760in}}%
\pgfpathlineto{\pgfqpoint{1.359515in}{0.883509in}}%
\pgfpathlineto{\pgfqpoint{1.364260in}{0.883509in}}%
\pgfpathlineto{\pgfqpoint{1.369004in}{0.888760in}}%
\pgfpathlineto{\pgfqpoint{1.373749in}{0.825749in}}%
\pgfpathlineto{\pgfqpoint{1.378493in}{1.067291in}}%
\pgfpathlineto{\pgfqpoint{1.383238in}{1.229741in}}%
\pgfpathlineto{\pgfqpoint{1.387982in}{1.319006in}}%
\pgfpathlineto{\pgfqpoint{1.392726in}{1.382017in}}%
\pgfpathlineto{\pgfqpoint{1.397471in}{1.418773in}}%
\pgfpathlineto{\pgfqpoint{1.402215in}{1.434526in}}%
\pgfpathlineto{\pgfqpoint{1.406960in}{1.429275in}}%
\pgfpathlineto{\pgfqpoint{1.411704in}{1.408272in}}%
\pgfpathlineto{\pgfqpoint{1.416449in}{1.361013in}}%
\pgfpathlineto{\pgfqpoint{1.421193in}{1.292423in}}%
\pgfpathlineto{\pgfqpoint{1.425938in}{1.176903in}}%
\pgfpathlineto{\pgfqpoint{1.430682in}{0.977369in}}%
\pgfpathlineto{\pgfqpoint{1.435426in}{0.898605in}}%
\pgfpathlineto{\pgfqpoint{1.440171in}{0.882853in}}%
\pgfpathlineto{\pgfqpoint{1.444915in}{0.851347in}}%
\pgfpathlineto{\pgfqpoint{1.449660in}{0.872351in}}%
\pgfpathlineto{\pgfqpoint{1.454404in}{0.825093in}}%
\pgfpathlineto{\pgfqpoint{1.459149in}{0.830015in}}%
\pgfpathlineto{\pgfqpoint{1.463893in}{0.919281in}}%
\pgfpathlineto{\pgfqpoint{1.468637in}{0.914030in}}%
\pgfpathlineto{\pgfqpoint{1.473382in}{0.861521in}}%
\pgfpathlineto{\pgfqpoint{1.478126in}{0.898277in}}%
\pgfpathlineto{\pgfqpoint{1.482871in}{1.124066in}}%
\pgfpathlineto{\pgfqpoint{1.487615in}{1.250088in}}%
\pgfpathlineto{\pgfqpoint{1.492360in}{1.328852in}}%
\pgfpathlineto{\pgfqpoint{1.497104in}{1.381032in}}%
\pgfpathlineto{\pgfqpoint{1.501849in}{1.412538in}}%
\pgfpathlineto{\pgfqpoint{1.506593in}{1.423040in}}%
\pgfpathlineto{\pgfqpoint{1.511337in}{1.407287in}}%
\pgfpathlineto{\pgfqpoint{1.516082in}{1.375782in}}%
\pgfpathlineto{\pgfqpoint{1.520826in}{1.323272in}}%
\pgfpathlineto{\pgfqpoint{1.525571in}{1.244509in}}%
\pgfpathlineto{\pgfqpoint{1.530315in}{1.102406in}}%
\pgfpathlineto{\pgfqpoint{1.535060in}{0.902872in}}%
\pgfpathlineto{\pgfqpoint{1.539804in}{0.876617in}}%
\pgfpathlineto{\pgfqpoint{1.544549in}{0.908123in}}%
\pgfpathlineto{\pgfqpoint{1.549293in}{0.881868in}}%
\pgfpathlineto{\pgfqpoint{1.554037in}{0.908123in}}%
\pgfpathlineto{\pgfqpoint{1.558782in}{0.845112in}}%
\pgfpathlineto{\pgfqpoint{1.563526in}{0.845112in}}%
\pgfpathlineto{\pgfqpoint{1.568271in}{0.823780in}}%
\pgfpathlineto{\pgfqpoint{1.573015in}{0.913045in}}%
\pgfpathlineto{\pgfqpoint{1.577760in}{0.860536in}}%
\pgfpathlineto{\pgfqpoint{1.582504in}{0.844783in}}%
\pgfpathlineto{\pgfqpoint{1.587249in}{0.892042in}}%
\pgfpathlineto{\pgfqpoint{1.591993in}{0.855285in}}%
\pgfpathlineto{\pgfqpoint{1.596737in}{0.886791in}}%
\pgfpathlineto{\pgfqpoint{1.601482in}{0.897293in}}%
\pgfpathlineto{\pgfqpoint{1.606226in}{0.870710in}}%
\pgfpathlineto{\pgfqpoint{1.610971in}{0.860208in}}%
\pgfpathlineto{\pgfqpoint{1.615715in}{0.938972in}}%
\pgfpathlineto{\pgfqpoint{1.620460in}{0.870710in}}%
\pgfpathlineto{\pgfqpoint{1.625204in}{0.886463in}}%
\pgfpathlineto{\pgfqpoint{1.629948in}{0.917968in}}%
\pgfpathlineto{\pgfqpoint{1.634693in}{0.881212in}}%
\pgfpathlineto{\pgfqpoint{1.639437in}{0.854957in}}%
\pgfpathlineto{\pgfqpoint{1.644182in}{0.802120in}}%
\pgfpathlineto{\pgfqpoint{1.648926in}{0.886134in}}%
\pgfpathlineto{\pgfqpoint{1.653671in}{0.896636in}}%
\pgfpathlineto{\pgfqpoint{1.658415in}{0.854629in}}%
\pgfpathlineto{\pgfqpoint{1.663160in}{0.865131in}}%
\pgfpathlineto{\pgfqpoint{1.667904in}{0.896636in}}%
\pgfpathlineto{\pgfqpoint{1.672648in}{0.865131in}}%
\pgfpathlineto{\pgfqpoint{1.677393in}{0.854629in}}%
\pgfpathlineto{\pgfqpoint{1.682137in}{0.854301in}}%
\pgfpathlineto{\pgfqpoint{1.686882in}{1.069588in}}%
\pgfpathlineto{\pgfqpoint{1.691626in}{1.211363in}}%
\pgfpathlineto{\pgfqpoint{1.696371in}{1.311130in}}%
\pgfpathlineto{\pgfqpoint{1.701115in}{1.368890in}}%
\pgfpathlineto{\pgfqpoint{1.705860in}{1.405646in}}%
\pgfpathlineto{\pgfqpoint{1.710604in}{1.421399in}}%
\pgfpathlineto{\pgfqpoint{1.715348in}{1.416148in}}%
\pgfpathlineto{\pgfqpoint{1.720093in}{1.389565in}}%
\pgfpathlineto{\pgfqpoint{1.724837in}{1.347558in}}%
\pgfpathlineto{\pgfqpoint{1.729582in}{1.274045in}}%
\pgfpathlineto{\pgfqpoint{1.734326in}{1.153274in}}%
\pgfpathlineto{\pgfqpoint{1.739071in}{0.937987in}}%
\pgfpathlineto{\pgfqpoint{1.743815in}{0.822467in}}%
\pgfpathlineto{\pgfqpoint{1.748559in}{0.901231in}}%
\pgfpathlineto{\pgfqpoint{1.753304in}{0.806714in}}%
\pgfpathlineto{\pgfqpoint{1.758048in}{0.864146in}}%
\pgfpathlineto{\pgfqpoint{1.762793in}{0.895652in}}%
\pgfpathlineto{\pgfqpoint{1.772282in}{0.874648in}}%
\pgfpathlineto{\pgfqpoint{1.777026in}{0.885150in}}%
\pgfpathlineto{\pgfqpoint{1.781771in}{0.879899in}}%
\pgfpathlineto{\pgfqpoint{1.786515in}{0.900903in}}%
\pgfpathlineto{\pgfqpoint{1.796004in}{1.168371in}}%
\pgfpathlineto{\pgfqpoint{1.800748in}{1.247134in}}%
\pgfpathlineto{\pgfqpoint{1.805493in}{1.299643in}}%
\pgfpathlineto{\pgfqpoint{1.810237in}{1.320647in}}%
\pgfpathlineto{\pgfqpoint{1.814982in}{1.331149in}}%
\pgfpathlineto{\pgfqpoint{1.819726in}{1.320647in}}%
\pgfpathlineto{\pgfqpoint{1.824471in}{1.294392in}}%
\pgfpathlineto{\pgfqpoint{1.829215in}{1.236632in}}%
\pgfpathlineto{\pgfqpoint{1.833959in}{1.147039in}}%
\pgfpathlineto{\pgfqpoint{1.843448in}{0.852988in}}%
\pgfpathlineto{\pgfqpoint{1.848193in}{0.805730in}}%
\pgfpathlineto{\pgfqpoint{1.852937in}{0.852988in}}%
\pgfpathlineto{\pgfqpoint{1.857682in}{0.842486in}}%
\pgfpathlineto{\pgfqpoint{1.862426in}{0.894995in}}%
\pgfpathlineto{\pgfqpoint{1.867171in}{0.816232in}}%
\pgfpathlineto{\pgfqpoint{1.871915in}{0.768645in}}%
\pgfpathlineto{\pgfqpoint{1.876659in}{0.905169in}}%
\pgfpathlineto{\pgfqpoint{1.881404in}{0.868413in}}%
\pgfpathlineto{\pgfqpoint{1.886148in}{0.857911in}}%
\pgfpathlineto{\pgfqpoint{1.890893in}{0.962929in}}%
\pgfpathlineto{\pgfqpoint{1.895637in}{1.172965in}}%
\pgfpathlineto{\pgfqpoint{1.900382in}{1.283234in}}%
\pgfpathlineto{\pgfqpoint{1.905126in}{1.351496in}}%
\pgfpathlineto{\pgfqpoint{1.909870in}{1.393175in}}%
\pgfpathlineto{\pgfqpoint{1.914615in}{1.419430in}}%
\pgfpathlineto{\pgfqpoint{1.919359in}{1.419430in}}%
\pgfpathlineto{\pgfqpoint{1.924104in}{1.403677in}}%
\pgfpathlineto{\pgfqpoint{1.928848in}{1.361670in}}%
\pgfpathlineto{\pgfqpoint{1.933593in}{1.298659in}}%
\pgfpathlineto{\pgfqpoint{1.938337in}{1.204143in}}%
\pgfpathlineto{\pgfqpoint{1.947826in}{0.888760in}}%
\pgfpathlineto{\pgfqpoint{1.952570in}{0.757487in}}%
\pgfpathlineto{\pgfqpoint{1.957315in}{0.836251in}}%
\pgfpathlineto{\pgfqpoint{1.962059in}{0.867756in}}%
\pgfpathlineto{\pgfqpoint{1.966804in}{0.804745in}}%
\pgfpathlineto{\pgfqpoint{1.971548in}{0.846753in}}%
\pgfpathlineto{\pgfqpoint{1.976293in}{0.809996in}}%
\pgfpathlineto{\pgfqpoint{1.981037in}{0.873007in}}%
\pgfpathlineto{\pgfqpoint{1.985782in}{0.799166in}}%
\pgfpathlineto{\pgfqpoint{1.990526in}{0.830672in}}%
\pgfpathlineto{\pgfqpoint{1.995270in}{0.814919in}}%
\pgfpathlineto{\pgfqpoint{2.000015in}{0.893683in}}%
\pgfpathlineto{\pgfqpoint{2.004759in}{0.793915in}}%
\pgfpathlineto{\pgfqpoint{2.009504in}{0.851675in}}%
\pgfpathlineto{\pgfqpoint{2.014248in}{0.967195in}}%
\pgfpathlineto{\pgfqpoint{2.018993in}{0.862177in}}%
\pgfpathlineto{\pgfqpoint{2.023737in}{0.909107in}}%
\pgfpathlineto{\pgfqpoint{2.033226in}{0.851347in}}%
\pgfpathlineto{\pgfqpoint{2.042715in}{0.830343in}}%
\pgfpathlineto{\pgfqpoint{2.047459in}{0.814591in}}%
\pgfpathlineto{\pgfqpoint{2.052204in}{0.951114in}}%
\pgfpathlineto{\pgfqpoint{2.056948in}{0.872351in}}%
\pgfpathlineto{\pgfqpoint{2.061693in}{0.861521in}}%
\pgfpathlineto{\pgfqpoint{2.066437in}{0.893026in}}%
\pgfpathlineto{\pgfqpoint{2.071181in}{0.877273in}}%
\pgfpathlineto{\pgfqpoint{2.075926in}{0.809012in}}%
\pgfpathlineto{\pgfqpoint{2.080670in}{0.819513in}}%
\pgfpathlineto{\pgfqpoint{2.085415in}{0.845768in}}%
\pgfpathlineto{\pgfqpoint{2.090159in}{0.998044in}}%
\pgfpathlineto{\pgfqpoint{2.094904in}{0.809012in}}%
\pgfpathlineto{\pgfqpoint{2.099648in}{1.102734in}}%
\pgfpathlineto{\pgfqpoint{2.109137in}{1.318022in}}%
\pgfpathlineto{\pgfqpoint{2.113881in}{1.365280in}}%
\pgfpathlineto{\pgfqpoint{2.118626in}{1.396785in}}%
\pgfpathlineto{\pgfqpoint{2.123370in}{1.402036in}}%
\pgfpathlineto{\pgfqpoint{2.128115in}{1.391534in}}%
\pgfpathlineto{\pgfqpoint{2.132859in}{1.360029in}}%
\pgfpathlineto{\pgfqpoint{2.137604in}{1.301941in}}%
\pgfpathlineto{\pgfqpoint{2.142348in}{1.223177in}}%
\pgfpathlineto{\pgfqpoint{2.147093in}{1.081403in}}%
\pgfpathlineto{\pgfqpoint{2.151837in}{0.866115in}}%
\pgfpathlineto{\pgfqpoint{2.156581in}{0.803104in}}%
\pgfpathlineto{\pgfqpoint{2.161326in}{0.871366in}}%
\pgfpathlineto{\pgfqpoint{2.166070in}{0.860864in}}%
\pgfpathlineto{\pgfqpoint{2.170815in}{0.808355in}}%
\pgfpathlineto{\pgfqpoint{2.180304in}{0.871038in}}%
\pgfpathlineto{\pgfqpoint{2.185048in}{0.860536in}}%
\pgfpathlineto{\pgfqpoint{2.189792in}{0.797525in}}%
\pgfpathlineto{\pgfqpoint{2.194537in}{0.834282in}}%
\pgfpathlineto{\pgfqpoint{2.199281in}{0.981307in}}%
\pgfpathlineto{\pgfqpoint{2.204026in}{1.180842in}}%
\pgfpathlineto{\pgfqpoint{2.208770in}{1.280609in}}%
\pgfpathlineto{\pgfqpoint{2.213515in}{1.353793in}}%
\pgfpathlineto{\pgfqpoint{2.218259in}{1.395801in}}%
\pgfpathlineto{\pgfqpoint{2.223004in}{1.416804in}}%
\pgfpathlineto{\pgfqpoint{2.227748in}{1.422055in}}%
\pgfpathlineto{\pgfqpoint{2.232492in}{1.401052in}}%
\pgfpathlineto{\pgfqpoint{2.237237in}{1.364295in}}%
\pgfpathlineto{\pgfqpoint{2.241981in}{1.301284in}}%
\pgfpathlineto{\pgfqpoint{2.246726in}{1.196266in}}%
\pgfpathlineto{\pgfqpoint{2.256215in}{0.854629in}}%
\pgfpathlineto{\pgfqpoint{2.260959in}{0.802120in}}%
\pgfpathlineto{\pgfqpoint{2.265704in}{0.854629in}}%
\pgfpathlineto{\pgfqpoint{2.270448in}{0.844127in}}%
\pgfpathlineto{\pgfqpoint{2.275192in}{0.796869in}}%
\pgfpathlineto{\pgfqpoint{2.279937in}{0.865131in}}%
\pgfpathlineto{\pgfqpoint{2.284681in}{0.833625in}}%
\pgfpathlineto{\pgfqpoint{2.289426in}{0.854301in}}%
\pgfpathlineto{\pgfqpoint{2.294170in}{0.849050in}}%
\pgfpathlineto{\pgfqpoint{2.298915in}{0.822795in}}%
\pgfpathlineto{\pgfqpoint{2.303659in}{1.043333in}}%
\pgfpathlineto{\pgfqpoint{2.308404in}{1.190359in}}%
\pgfpathlineto{\pgfqpoint{2.313148in}{1.284875in}}%
\pgfpathlineto{\pgfqpoint{2.317892in}{1.342635in}}%
\pgfpathlineto{\pgfqpoint{2.322637in}{1.374141in}}%
\pgfpathlineto{\pgfqpoint{2.327381in}{1.389565in}}%
\pgfpathlineto{\pgfqpoint{2.332126in}{1.384314in}}%
\pgfpathlineto{\pgfqpoint{2.336870in}{1.363311in}}%
\pgfpathlineto{\pgfqpoint{2.341615in}{1.316052in}}%
\pgfpathlineto{\pgfqpoint{2.346359in}{1.242540in}}%
\pgfpathlineto{\pgfqpoint{2.351103in}{1.116518in}}%
\pgfpathlineto{\pgfqpoint{2.355848in}{0.890729in}}%
\pgfpathlineto{\pgfqpoint{2.360592in}{0.864474in}}%
\pgfpathlineto{\pgfqpoint{2.365337in}{0.806386in}}%
\pgfpathlineto{\pgfqpoint{2.370081in}{0.764379in}}%
\pgfpathlineto{\pgfqpoint{2.374826in}{0.895652in}}%
\pgfpathlineto{\pgfqpoint{2.384315in}{0.806386in}}%
\pgfpathlineto{\pgfqpoint{2.389059in}{0.811637in}}%
\pgfpathlineto{\pgfqpoint{2.393803in}{0.879899in}}%
\pgfpathlineto{\pgfqpoint{2.398548in}{0.806386in}}%
\pgfpathlineto{\pgfqpoint{2.403292in}{0.827062in}}%
\pgfpathlineto{\pgfqpoint{2.408037in}{1.042349in}}%
\pgfpathlineto{\pgfqpoint{2.412781in}{1.168371in}}%
\pgfpathlineto{\pgfqpoint{2.417526in}{1.247134in}}%
\pgfpathlineto{\pgfqpoint{2.422270in}{1.299643in}}%
\pgfpathlineto{\pgfqpoint{2.427015in}{1.325898in}}%
\pgfpathlineto{\pgfqpoint{2.431759in}{1.331149in}}%
\pgfpathlineto{\pgfqpoint{2.436503in}{1.320647in}}%
\pgfpathlineto{\pgfqpoint{2.441248in}{1.288813in}}%
\pgfpathlineto{\pgfqpoint{2.445992in}{1.231053in}}%
\pgfpathlineto{\pgfqpoint{2.450737in}{1.141788in}}%
\pgfpathlineto{\pgfqpoint{2.455481in}{0.994763in}}%
\pgfpathlineto{\pgfqpoint{2.460226in}{0.763723in}}%
\pgfpathlineto{\pgfqpoint{2.464970in}{0.847737in}}%
\pgfpathlineto{\pgfqpoint{2.469714in}{0.810981in}}%
\pgfpathlineto{\pgfqpoint{2.474459in}{0.732217in}}%
\pgfpathlineto{\pgfqpoint{2.479203in}{0.752893in}}%
\pgfpathlineto{\pgfqpoint{2.483948in}{0.794900in}}%
\pgfpathlineto{\pgfqpoint{2.488692in}{0.752893in}}%
\pgfpathlineto{\pgfqpoint{2.493437in}{0.731889in}}%
\pgfpathlineto{\pgfqpoint{2.498181in}{0.815903in}}%
\pgfpathlineto{\pgfqpoint{2.502926in}{0.752893in}}%
\pgfpathlineto{\pgfqpoint{2.507670in}{0.815903in}}%
\pgfpathlineto{\pgfqpoint{2.512414in}{1.015438in}}%
\pgfpathlineto{\pgfqpoint{2.517159in}{1.125379in}}%
\pgfpathlineto{\pgfqpoint{2.521903in}{1.188390in}}%
\pgfpathlineto{\pgfqpoint{2.526648in}{1.230397in}}%
\pgfpathlineto{\pgfqpoint{2.531392in}{1.251401in}}%
\pgfpathlineto{\pgfqpoint{2.536137in}{1.251401in}}%
\pgfpathlineto{\pgfqpoint{2.540881in}{1.235648in}}%
\pgfpathlineto{\pgfqpoint{2.545626in}{1.193641in}}%
\pgfpathlineto{\pgfqpoint{2.550370in}{1.130630in}}%
\pgfpathlineto{\pgfqpoint{2.555114in}{1.030534in}}%
\pgfpathlineto{\pgfqpoint{2.559859in}{0.836251in}}%
\pgfpathlineto{\pgfqpoint{2.564603in}{0.815247in}}%
\pgfpathlineto{\pgfqpoint{2.569348in}{0.773240in}}%
\pgfpathlineto{\pgfqpoint{2.574092in}{0.762738in}}%
\pgfpathlineto{\pgfqpoint{2.578837in}{0.767989in}}%
\pgfpathlineto{\pgfqpoint{2.583581in}{0.715480in}}%
\pgfpathlineto{\pgfqpoint{2.588326in}{0.746985in}}%
\pgfpathlineto{\pgfqpoint{2.593070in}{0.751908in}}%
\pgfpathlineto{\pgfqpoint{2.597814in}{0.778163in}}%
\pgfpathlineto{\pgfqpoint{2.602559in}{0.762410in}}%
\pgfpathlineto{\pgfqpoint{2.607303in}{0.736155in}}%
\pgfpathlineto{\pgfqpoint{2.612048in}{0.793915in}}%
\pgfpathlineto{\pgfqpoint{2.616792in}{0.956693in}}%
\pgfpathlineto{\pgfqpoint{2.621537in}{1.051210in}}%
\pgfpathlineto{\pgfqpoint{2.626281in}{1.108970in}}%
\pgfpathlineto{\pgfqpoint{2.631025in}{1.150649in}}%
\pgfpathlineto{\pgfqpoint{2.635770in}{1.161151in}}%
\pgfpathlineto{\pgfqpoint{2.640514in}{1.155900in}}%
\pgfpathlineto{\pgfqpoint{2.645259in}{1.129645in}}%
\pgfpathlineto{\pgfqpoint{2.650003in}{1.082387in}}%
\pgfpathlineto{\pgfqpoint{2.654748in}{1.003623in}}%
\pgfpathlineto{\pgfqpoint{2.659492in}{0.888103in}}%
\pgfpathlineto{\pgfqpoint{2.664237in}{0.714823in}}%
\pgfpathlineto{\pgfqpoint{2.668981in}{0.730248in}}%
\pgfpathlineto{\pgfqpoint{2.673725in}{0.851019in}}%
\pgfpathlineto{\pgfqpoint{2.678470in}{0.714495in}}%
\pgfpathlineto{\pgfqpoint{2.683214in}{0.714495in}}%
\pgfpathlineto{\pgfqpoint{2.687959in}{0.798510in}}%
\pgfpathlineto{\pgfqpoint{2.692703in}{0.735499in}}%
\pgfpathlineto{\pgfqpoint{2.697448in}{0.877273in}}%
\pgfpathlineto{\pgfqpoint{2.702192in}{0.782757in}}%
\pgfpathlineto{\pgfqpoint{2.706937in}{0.729920in}}%
\pgfpathlineto{\pgfqpoint{2.711681in}{0.792931in}}%
\pgfpathlineto{\pgfqpoint{2.716425in}{0.777178in}}%
\pgfpathlineto{\pgfqpoint{2.721170in}{0.887447in}}%
\pgfpathlineto{\pgfqpoint{2.730659in}{1.018720in}}%
\pgfpathlineto{\pgfqpoint{2.735403in}{1.044974in}}%
\pgfpathlineto{\pgfqpoint{2.740148in}{1.055476in}}%
\pgfpathlineto{\pgfqpoint{2.744892in}{1.044646in}}%
\pgfpathlineto{\pgfqpoint{2.749637in}{1.007890in}}%
\pgfpathlineto{\pgfqpoint{2.754381in}{0.950130in}}%
\pgfpathlineto{\pgfqpoint{2.763870in}{0.750595in}}%
\pgfpathlineto{\pgfqpoint{2.768614in}{0.776850in}}%
\pgfpathlineto{\pgfqpoint{2.773359in}{0.855613in}}%
\pgfpathlineto{\pgfqpoint{2.778103in}{0.713839in}}%
\pgfpathlineto{\pgfqpoint{2.782848in}{0.745016in}}%
\pgfpathlineto{\pgfqpoint{2.792336in}{0.713511in}}%
\pgfpathlineto{\pgfqpoint{2.797081in}{0.713511in}}%
\pgfpathlineto{\pgfqpoint{2.801825in}{0.729263in}}%
\pgfpathlineto{\pgfqpoint{2.806570in}{0.713511in}}%
\pgfpathlineto{\pgfqpoint{2.811314in}{0.713511in}}%
\pgfpathlineto{\pgfqpoint{2.816059in}{0.760769in}}%
\pgfpathlineto{\pgfqpoint{2.820803in}{0.728935in}}%
\pgfpathlineto{\pgfqpoint{2.830292in}{0.833953in}}%
\pgfpathlineto{\pgfqpoint{2.835036in}{0.912717in}}%
\pgfpathlineto{\pgfqpoint{2.839781in}{0.917968in}}%
\pgfpathlineto{\pgfqpoint{2.844525in}{0.944223in}}%
\pgfpathlineto{\pgfqpoint{2.849270in}{0.949473in}}%
\pgfpathlineto{\pgfqpoint{2.854014in}{0.875961in}}%
\pgfpathlineto{\pgfqpoint{2.858759in}{0.781116in}}%
\pgfpathlineto{\pgfqpoint{2.863503in}{0.728607in}}%
\pgfpathlineto{\pgfqpoint{2.868248in}{0.712854in}}%
\pgfpathlineto{\pgfqpoint{2.868248in}{0.712854in}}%
\pgfusepath{stroke}%
\end{pgfscope}%
\begin{pgfscope}%
\pgfsetrectcap%
\pgfsetmiterjoin%
\pgfsetlinewidth{0.803000pt}%
\definecolor{currentstroke}{rgb}{0.000000,0.000000,0.000000}%
\pgfsetstrokecolor{currentstroke}%
\pgfsetdash{}{0pt}%
\pgfpathmoveto{\pgfqpoint{0.758026in}{0.565123in}}%
\pgfpathlineto{\pgfqpoint{0.758026in}{1.475926in}}%
\pgfusepath{stroke}%
\end{pgfscope}%
\begin{pgfscope}%
\pgfsetrectcap%
\pgfsetmiterjoin%
\pgfsetlinewidth{0.803000pt}%
\definecolor{currentstroke}{rgb}{0.000000,0.000000,0.000000}%
\pgfsetstrokecolor{currentstroke}%
\pgfsetdash{}{0pt}%
\pgfpathmoveto{\pgfqpoint{2.866666in}{0.565123in}}%
\pgfpathlineto{\pgfqpoint{2.866666in}{1.475926in}}%
\pgfusepath{stroke}%
\end{pgfscope}%
\begin{pgfscope}%
\pgfsetrectcap%
\pgfsetmiterjoin%
\pgfsetlinewidth{0.803000pt}%
\definecolor{currentstroke}{rgb}{0.000000,0.000000,0.000000}%
\pgfsetstrokecolor{currentstroke}%
\pgfsetdash{}{0pt}%
\pgfpathmoveto{\pgfqpoint{0.758026in}{0.565123in}}%
\pgfpathlineto{\pgfqpoint{2.866666in}{0.565123in}}%
\pgfusepath{stroke}%
\end{pgfscope}%
\begin{pgfscope}%
\pgfsetrectcap%
\pgfsetmiterjoin%
\pgfsetlinewidth{0.803000pt}%
\definecolor{currentstroke}{rgb}{0.000000,0.000000,0.000000}%
\pgfsetstrokecolor{currentstroke}%
\pgfsetdash{}{0pt}%
\pgfpathmoveto{\pgfqpoint{0.758026in}{1.475926in}}%
\pgfpathlineto{\pgfqpoint{2.866666in}{1.475926in}}%
\pgfusepath{stroke}%
\end{pgfscope}%
\begin{pgfscope}%
\definecolor{textcolor}{rgb}{0.000000,0.000000,0.000000}%
\pgfsetstrokecolor{textcolor}%
\pgfsetfillcolor{textcolor}%
\pgftext[x=1.812346in,y=1.559260in,,base]{\color{textcolor}\rmfamily\fontsize{12.000000}{14.400000}\selectfont 19,472 kHz}%
\end{pgfscope}%
\begin{pgfscope}%
\pgfsetbuttcap%
\pgfsetmiterjoin%
\definecolor{currentfill}{rgb}{1.000000,1.000000,1.000000}%
\pgfsetfillcolor{currentfill}%
\pgfsetlinewidth{0.000000pt}%
\definecolor{currentstroke}{rgb}{0.000000,0.000000,0.000000}%
\pgfsetstrokecolor{currentstroke}%
\pgfsetstrokeopacity{0.000000}%
\pgfsetdash{}{0pt}%
\pgfpathmoveto{\pgfqpoint{3.833026in}{0.565123in}}%
\pgfpathlineto{\pgfqpoint{5.941666in}{0.565123in}}%
\pgfpathlineto{\pgfqpoint{5.941666in}{1.475926in}}%
\pgfpathlineto{\pgfqpoint{3.833026in}{1.475926in}}%
\pgfpathlineto{\pgfqpoint{3.833026in}{0.565123in}}%
\pgfpathclose%
\pgfusepath{fill}%
\end{pgfscope}%
\begin{pgfscope}%
\pgfpathrectangle{\pgfqpoint{3.833026in}{0.565123in}}{\pgfqpoint{2.108640in}{0.910803in}}%
\pgfusepath{clip}%
\pgfsetbuttcap%
\pgfsetroundjoin%
\pgfsetlinewidth{0.501875pt}%
\definecolor{currentstroke}{rgb}{0.690196,0.690196,0.690196}%
\pgfsetstrokecolor{currentstroke}%
\pgfsetstrokeopacity{0.700000}%
\pgfsetdash{{1.850000pt}{0.800000pt}}{0.000000pt}%
\pgfpathmoveto{\pgfqpoint{3.833026in}{0.565123in}}%
\pgfpathlineto{\pgfqpoint{3.833026in}{1.475926in}}%
\pgfusepath{stroke}%
\end{pgfscope}%
\begin{pgfscope}%
\pgfsetbuttcap%
\pgfsetroundjoin%
\definecolor{currentfill}{rgb}{0.000000,0.000000,0.000000}%
\pgfsetfillcolor{currentfill}%
\pgfsetlinewidth{0.803000pt}%
\definecolor{currentstroke}{rgb}{0.000000,0.000000,0.000000}%
\pgfsetstrokecolor{currentstroke}%
\pgfsetdash{}{0pt}%
\pgfsys@defobject{currentmarker}{\pgfqpoint{0.000000in}{-0.048611in}}{\pgfqpoint{0.000000in}{0.000000in}}{%
\pgfpathmoveto{\pgfqpoint{0.000000in}{0.000000in}}%
\pgfpathlineto{\pgfqpoint{0.000000in}{-0.048611in}}%
\pgfusepath{stroke,fill}%
}%
\begin{pgfscope}%
\pgfsys@transformshift{3.833026in}{0.565123in}%
\pgfsys@useobject{currentmarker}{}%
\end{pgfscope}%
\end{pgfscope}%
\begin{pgfscope}%
\definecolor{textcolor}{rgb}{0.000000,0.000000,0.000000}%
\pgfsetstrokecolor{textcolor}%
\pgfsetfillcolor{textcolor}%
\pgftext[x=3.833026in,y=0.467901in,,top]{\color{textcolor}\rmfamily\fontsize{10.000000}{12.000000}\selectfont \(\displaystyle {300000}\)}%
\end{pgfscope}%
\begin{pgfscope}%
\pgfpathrectangle{\pgfqpoint{3.833026in}{0.565123in}}{\pgfqpoint{2.108640in}{0.910803in}}%
\pgfusepath{clip}%
\pgfsetbuttcap%
\pgfsetroundjoin%
\pgfsetlinewidth{0.501875pt}%
\definecolor{currentstroke}{rgb}{0.690196,0.690196,0.690196}%
\pgfsetstrokecolor{currentstroke}%
\pgfsetstrokeopacity{0.700000}%
\pgfsetdash{{1.850000pt}{0.800000pt}}{0.000000pt}%
\pgfpathmoveto{\pgfqpoint{4.360186in}{0.565123in}}%
\pgfpathlineto{\pgfqpoint{4.360186in}{1.475926in}}%
\pgfusepath{stroke}%
\end{pgfscope}%
\begin{pgfscope}%
\pgfsetbuttcap%
\pgfsetroundjoin%
\definecolor{currentfill}{rgb}{0.000000,0.000000,0.000000}%
\pgfsetfillcolor{currentfill}%
\pgfsetlinewidth{0.803000pt}%
\definecolor{currentstroke}{rgb}{0.000000,0.000000,0.000000}%
\pgfsetstrokecolor{currentstroke}%
\pgfsetdash{}{0pt}%
\pgfsys@defobject{currentmarker}{\pgfqpoint{0.000000in}{-0.048611in}}{\pgfqpoint{0.000000in}{0.000000in}}{%
\pgfpathmoveto{\pgfqpoint{0.000000in}{0.000000in}}%
\pgfpathlineto{\pgfqpoint{0.000000in}{-0.048611in}}%
\pgfusepath{stroke,fill}%
}%
\begin{pgfscope}%
\pgfsys@transformshift{4.360186in}{0.565123in}%
\pgfsys@useobject{currentmarker}{}%
\end{pgfscope}%
\end{pgfscope}%
\begin{pgfscope}%
\definecolor{textcolor}{rgb}{0.000000,0.000000,0.000000}%
\pgfsetstrokecolor{textcolor}%
\pgfsetfillcolor{textcolor}%
\pgftext[x=4.360186in,y=0.467901in,,top]{\color{textcolor}\rmfamily\fontsize{10.000000}{12.000000}\selectfont \(\displaystyle {400000}\)}%
\end{pgfscope}%
\begin{pgfscope}%
\pgfpathrectangle{\pgfqpoint{3.833026in}{0.565123in}}{\pgfqpoint{2.108640in}{0.910803in}}%
\pgfusepath{clip}%
\pgfsetbuttcap%
\pgfsetroundjoin%
\pgfsetlinewidth{0.501875pt}%
\definecolor{currentstroke}{rgb}{0.690196,0.690196,0.690196}%
\pgfsetstrokecolor{currentstroke}%
\pgfsetstrokeopacity{0.700000}%
\pgfsetdash{{1.850000pt}{0.800000pt}}{0.000000pt}%
\pgfpathmoveto{\pgfqpoint{4.887346in}{0.565123in}}%
\pgfpathlineto{\pgfqpoint{4.887346in}{1.475926in}}%
\pgfusepath{stroke}%
\end{pgfscope}%
\begin{pgfscope}%
\pgfsetbuttcap%
\pgfsetroundjoin%
\definecolor{currentfill}{rgb}{0.000000,0.000000,0.000000}%
\pgfsetfillcolor{currentfill}%
\pgfsetlinewidth{0.803000pt}%
\definecolor{currentstroke}{rgb}{0.000000,0.000000,0.000000}%
\pgfsetstrokecolor{currentstroke}%
\pgfsetdash{}{0pt}%
\pgfsys@defobject{currentmarker}{\pgfqpoint{0.000000in}{-0.048611in}}{\pgfqpoint{0.000000in}{0.000000in}}{%
\pgfpathmoveto{\pgfqpoint{0.000000in}{0.000000in}}%
\pgfpathlineto{\pgfqpoint{0.000000in}{-0.048611in}}%
\pgfusepath{stroke,fill}%
}%
\begin{pgfscope}%
\pgfsys@transformshift{4.887346in}{0.565123in}%
\pgfsys@useobject{currentmarker}{}%
\end{pgfscope}%
\end{pgfscope}%
\begin{pgfscope}%
\definecolor{textcolor}{rgb}{0.000000,0.000000,0.000000}%
\pgfsetstrokecolor{textcolor}%
\pgfsetfillcolor{textcolor}%
\pgftext[x=4.887346in,y=0.467901in,,top]{\color{textcolor}\rmfamily\fontsize{10.000000}{12.000000}\selectfont \(\displaystyle {500000}\)}%
\end{pgfscope}%
\begin{pgfscope}%
\pgfpathrectangle{\pgfqpoint{3.833026in}{0.565123in}}{\pgfqpoint{2.108640in}{0.910803in}}%
\pgfusepath{clip}%
\pgfsetbuttcap%
\pgfsetroundjoin%
\pgfsetlinewidth{0.501875pt}%
\definecolor{currentstroke}{rgb}{0.690196,0.690196,0.690196}%
\pgfsetstrokecolor{currentstroke}%
\pgfsetstrokeopacity{0.700000}%
\pgfsetdash{{1.850000pt}{0.800000pt}}{0.000000pt}%
\pgfpathmoveto{\pgfqpoint{5.414506in}{0.565123in}}%
\pgfpathlineto{\pgfqpoint{5.414506in}{1.475926in}}%
\pgfusepath{stroke}%
\end{pgfscope}%
\begin{pgfscope}%
\pgfsetbuttcap%
\pgfsetroundjoin%
\definecolor{currentfill}{rgb}{0.000000,0.000000,0.000000}%
\pgfsetfillcolor{currentfill}%
\pgfsetlinewidth{0.803000pt}%
\definecolor{currentstroke}{rgb}{0.000000,0.000000,0.000000}%
\pgfsetstrokecolor{currentstroke}%
\pgfsetdash{}{0pt}%
\pgfsys@defobject{currentmarker}{\pgfqpoint{0.000000in}{-0.048611in}}{\pgfqpoint{0.000000in}{0.000000in}}{%
\pgfpathmoveto{\pgfqpoint{0.000000in}{0.000000in}}%
\pgfpathlineto{\pgfqpoint{0.000000in}{-0.048611in}}%
\pgfusepath{stroke,fill}%
}%
\begin{pgfscope}%
\pgfsys@transformshift{5.414506in}{0.565123in}%
\pgfsys@useobject{currentmarker}{}%
\end{pgfscope}%
\end{pgfscope}%
\begin{pgfscope}%
\definecolor{textcolor}{rgb}{0.000000,0.000000,0.000000}%
\pgfsetstrokecolor{textcolor}%
\pgfsetfillcolor{textcolor}%
\pgftext[x=5.414506in,y=0.467901in,,top]{\color{textcolor}\rmfamily\fontsize{10.000000}{12.000000}\selectfont \(\displaystyle {600000}\)}%
\end{pgfscope}%
\begin{pgfscope}%
\pgfpathrectangle{\pgfqpoint{3.833026in}{0.565123in}}{\pgfqpoint{2.108640in}{0.910803in}}%
\pgfusepath{clip}%
\pgfsetbuttcap%
\pgfsetroundjoin%
\pgfsetlinewidth{0.501875pt}%
\definecolor{currentstroke}{rgb}{0.690196,0.690196,0.690196}%
\pgfsetstrokecolor{currentstroke}%
\pgfsetstrokeopacity{0.700000}%
\pgfsetdash{{1.850000pt}{0.800000pt}}{0.000000pt}%
\pgfpathmoveto{\pgfqpoint{5.941666in}{0.565123in}}%
\pgfpathlineto{\pgfqpoint{5.941666in}{1.475926in}}%
\pgfusepath{stroke}%
\end{pgfscope}%
\begin{pgfscope}%
\pgfsetbuttcap%
\pgfsetroundjoin%
\definecolor{currentfill}{rgb}{0.000000,0.000000,0.000000}%
\pgfsetfillcolor{currentfill}%
\pgfsetlinewidth{0.803000pt}%
\definecolor{currentstroke}{rgb}{0.000000,0.000000,0.000000}%
\pgfsetstrokecolor{currentstroke}%
\pgfsetdash{}{0pt}%
\pgfsys@defobject{currentmarker}{\pgfqpoint{0.000000in}{-0.048611in}}{\pgfqpoint{0.000000in}{0.000000in}}{%
\pgfpathmoveto{\pgfqpoint{0.000000in}{0.000000in}}%
\pgfpathlineto{\pgfqpoint{0.000000in}{-0.048611in}}%
\pgfusepath{stroke,fill}%
}%
\begin{pgfscope}%
\pgfsys@transformshift{5.941666in}{0.565123in}%
\pgfsys@useobject{currentmarker}{}%
\end{pgfscope}%
\end{pgfscope}%
\begin{pgfscope}%
\definecolor{textcolor}{rgb}{0.000000,0.000000,0.000000}%
\pgfsetstrokecolor{textcolor}%
\pgfsetfillcolor{textcolor}%
\pgftext[x=5.941666in,y=0.467901in,,top]{\color{textcolor}\rmfamily\fontsize{10.000000}{12.000000}\selectfont \(\displaystyle {700000}\)}%
\end{pgfscope}%
\begin{pgfscope}%
\definecolor{textcolor}{rgb}{0.000000,0.000000,0.000000}%
\pgfsetstrokecolor{textcolor}%
\pgfsetfillcolor{textcolor}%
\pgftext[x=4.887346in,y=0.288889in,,top]{\color{textcolor}\rmfamily\fontsize{10.000000}{12.000000}\selectfont Frequency (Hz)}%
\end{pgfscope}%
\begin{pgfscope}%
\pgfpathrectangle{\pgfqpoint{3.833026in}{0.565123in}}{\pgfqpoint{2.108640in}{0.910803in}}%
\pgfusepath{clip}%
\pgfsetbuttcap%
\pgfsetroundjoin%
\pgfsetlinewidth{0.501875pt}%
\definecolor{currentstroke}{rgb}{0.690196,0.690196,0.690196}%
\pgfsetstrokecolor{currentstroke}%
\pgfsetstrokeopacity{0.700000}%
\pgfsetdash{{1.850000pt}{0.800000pt}}{0.000000pt}%
\pgfpathmoveto{\pgfqpoint{3.833026in}{0.665093in}}%
\pgfpathlineto{\pgfqpoint{5.941666in}{0.665093in}}%
\pgfusepath{stroke}%
\end{pgfscope}%
\begin{pgfscope}%
\pgfsetbuttcap%
\pgfsetroundjoin%
\definecolor{currentfill}{rgb}{0.000000,0.000000,0.000000}%
\pgfsetfillcolor{currentfill}%
\pgfsetlinewidth{0.803000pt}%
\definecolor{currentstroke}{rgb}{0.000000,0.000000,0.000000}%
\pgfsetstrokecolor{currentstroke}%
\pgfsetdash{}{0pt}%
\pgfsys@defobject{currentmarker}{\pgfqpoint{-0.048611in}{0.000000in}}{\pgfqpoint{-0.000000in}{0.000000in}}{%
\pgfpathmoveto{\pgfqpoint{-0.000000in}{0.000000in}}%
\pgfpathlineto{\pgfqpoint{-0.048611in}{0.000000in}}%
\pgfusepath{stroke,fill}%
}%
\begin{pgfscope}%
\pgfsys@transformshift{3.833026in}{0.665093in}%
\pgfsys@useobject{currentmarker}{}%
\end{pgfscope}%
\end{pgfscope}%
\begin{pgfscope}%
\definecolor{textcolor}{rgb}{0.000000,0.000000,0.000000}%
\pgfsetstrokecolor{textcolor}%
\pgfsetfillcolor{textcolor}%
\pgftext[x=3.419444in, y=0.616867in, left, base]{\color{textcolor}\rmfamily\fontsize{10.000000}{12.000000}\selectfont \(\displaystyle {\ensuremath{-}100}\)}%
\end{pgfscope}%
\begin{pgfscope}%
\pgfpathrectangle{\pgfqpoint{3.833026in}{0.565123in}}{\pgfqpoint{2.108640in}{0.910803in}}%
\pgfusepath{clip}%
\pgfsetbuttcap%
\pgfsetroundjoin%
\pgfsetlinewidth{0.501875pt}%
\definecolor{currentstroke}{rgb}{0.690196,0.690196,0.690196}%
\pgfsetstrokecolor{currentstroke}%
\pgfsetstrokeopacity{0.700000}%
\pgfsetdash{{1.850000pt}{0.800000pt}}{0.000000pt}%
\pgfpathmoveto{\pgfqpoint{3.833026in}{1.200584in}}%
\pgfpathlineto{\pgfqpoint{5.941666in}{1.200584in}}%
\pgfusepath{stroke}%
\end{pgfscope}%
\begin{pgfscope}%
\pgfsetbuttcap%
\pgfsetroundjoin%
\definecolor{currentfill}{rgb}{0.000000,0.000000,0.000000}%
\pgfsetfillcolor{currentfill}%
\pgfsetlinewidth{0.803000pt}%
\definecolor{currentstroke}{rgb}{0.000000,0.000000,0.000000}%
\pgfsetstrokecolor{currentstroke}%
\pgfsetdash{}{0pt}%
\pgfsys@defobject{currentmarker}{\pgfqpoint{-0.048611in}{0.000000in}}{\pgfqpoint{-0.000000in}{0.000000in}}{%
\pgfpathmoveto{\pgfqpoint{-0.000000in}{0.000000in}}%
\pgfpathlineto{\pgfqpoint{-0.048611in}{0.000000in}}%
\pgfusepath{stroke,fill}%
}%
\begin{pgfscope}%
\pgfsys@transformshift{3.833026in}{1.200584in}%
\pgfsys@useobject{currentmarker}{}%
\end{pgfscope}%
\end{pgfscope}%
\begin{pgfscope}%
\definecolor{textcolor}{rgb}{0.000000,0.000000,0.000000}%
\pgfsetstrokecolor{textcolor}%
\pgfsetfillcolor{textcolor}%
\pgftext[x=3.488889in, y=1.152358in, left, base]{\color{textcolor}\rmfamily\fontsize{10.000000}{12.000000}\selectfont \(\displaystyle {\ensuremath{-}50}\)}%
\end{pgfscope}%
\begin{pgfscope}%
\definecolor{textcolor}{rgb}{0.000000,0.000000,0.000000}%
\pgfsetstrokecolor{textcolor}%
\pgfsetfillcolor{textcolor}%
\pgftext[x=3.363889in,y=1.020525in,,bottom,rotate=90.000000]{\color{textcolor}\rmfamily\fontsize{10.000000}{12.000000}\selectfont Amplitude (dB)}%
\end{pgfscope}%
\begin{pgfscope}%
\pgfpathrectangle{\pgfqpoint{3.833026in}{0.565123in}}{\pgfqpoint{2.108640in}{0.910803in}}%
\pgfusepath{clip}%
\pgfsetrectcap%
\pgfsetroundjoin%
\pgfsetlinewidth{1.003750pt}%
\definecolor{currentstroke}{rgb}{0.000000,0.000000,0.000000}%
\pgfsetstrokecolor{currentstroke}%
\pgfsetdash{}{0pt}%
\pgfpathmoveto{\pgfqpoint{3.831971in}{0.842139in}}%
\pgfpathlineto{\pgfqpoint{3.836716in}{0.890334in}}%
\pgfpathlineto{\pgfqpoint{3.841460in}{0.858204in}}%
\pgfpathlineto{\pgfqpoint{3.846205in}{0.911419in}}%
\pgfpathlineto{\pgfqpoint{3.850949in}{0.857869in}}%
\pgfpathlineto{\pgfqpoint{3.855694in}{0.857869in}}%
\pgfpathlineto{\pgfqpoint{3.860438in}{0.895354in}}%
\pgfpathlineto{\pgfqpoint{3.865182in}{0.900709in}}%
\pgfpathlineto{\pgfqpoint{3.869927in}{0.793611in}}%
\pgfpathlineto{\pgfqpoint{3.874671in}{0.879289in}}%
\pgfpathlineto{\pgfqpoint{3.879416in}{0.857869in}}%
\pgfpathlineto{\pgfqpoint{3.884160in}{0.712952in}}%
\pgfpathlineto{\pgfqpoint{3.888905in}{0.846825in}}%
\pgfpathlineto{\pgfqpoint{3.893649in}{0.862890in}}%
\pgfpathlineto{\pgfqpoint{3.898394in}{0.814696in}}%
\pgfpathlineto{\pgfqpoint{3.903138in}{0.868245in}}%
\pgfpathlineto{\pgfqpoint{3.907882in}{0.857535in}}%
\pgfpathlineto{\pgfqpoint{3.912627in}{0.836115in}}%
\pgfpathlineto{\pgfqpoint{3.917371in}{0.841470in}}%
\pgfpathlineto{\pgfqpoint{3.922116in}{0.932169in}}%
\pgfpathlineto{\pgfqpoint{3.926860in}{0.980363in}}%
\pgfpathlineto{\pgfqpoint{3.931605in}{1.007138in}}%
\pgfpathlineto{\pgfqpoint{3.936349in}{0.985718in}}%
\pgfpathlineto{\pgfqpoint{3.941094in}{0.975008in}}%
\pgfpathlineto{\pgfqpoint{3.945838in}{0.926814in}}%
\pgfpathlineto{\pgfqpoint{3.950582in}{0.894684in}}%
\pgfpathlineto{\pgfqpoint{3.955327in}{0.841135in}}%
\pgfpathlineto{\pgfqpoint{3.964816in}{0.878285in}}%
\pgfpathlineto{\pgfqpoint{3.969560in}{0.883640in}}%
\pgfpathlineto{\pgfqpoint{3.974305in}{0.835446in}}%
\pgfpathlineto{\pgfqpoint{3.979049in}{0.915769in}}%
\pgfpathlineto{\pgfqpoint{3.983793in}{0.883640in}}%
\pgfpathlineto{\pgfqpoint{3.988538in}{0.915769in}}%
\pgfpathlineto{\pgfqpoint{3.993282in}{0.899705in}}%
\pgfpathlineto{\pgfqpoint{3.998027in}{0.867241in}}%
\pgfpathlineto{\pgfqpoint{4.002771in}{0.926145in}}%
\pgfpathlineto{\pgfqpoint{4.007516in}{0.845821in}}%
\pgfpathlineto{\pgfqpoint{4.012260in}{0.958274in}}%
\pgfpathlineto{\pgfqpoint{4.017005in}{1.033243in}}%
\pgfpathlineto{\pgfqpoint{4.026493in}{1.102857in}}%
\pgfpathlineto{\pgfqpoint{4.031238in}{1.097502in}}%
\pgfpathlineto{\pgfqpoint{4.035982in}{1.081102in}}%
\pgfpathlineto{\pgfqpoint{4.040727in}{1.038263in}}%
\pgfpathlineto{\pgfqpoint{4.050216in}{0.904390in}}%
\pgfpathlineto{\pgfqpoint{4.054960in}{0.845486in}}%
\pgfpathlineto{\pgfqpoint{4.059705in}{0.872261in}}%
\pgfpathlineto{\pgfqpoint{4.064449in}{0.877616in}}%
\pgfpathlineto{\pgfqpoint{4.069193in}{0.877616in}}%
\pgfpathlineto{\pgfqpoint{4.073938in}{0.909410in}}%
\pgfpathlineto{\pgfqpoint{4.083427in}{0.802312in}}%
\pgfpathlineto{\pgfqpoint{4.088171in}{0.748763in}}%
\pgfpathlineto{\pgfqpoint{4.092916in}{0.887991in}}%
\pgfpathlineto{\pgfqpoint{4.097660in}{0.823732in}}%
\pgfpathlineto{\pgfqpoint{4.111893in}{1.133982in}}%
\pgfpathlineto{\pgfqpoint{4.116638in}{1.176821in}}%
\pgfpathlineto{\pgfqpoint{4.121382in}{1.198241in}}%
\pgfpathlineto{\pgfqpoint{4.126127in}{1.208951in}}%
\pgfpathlineto{\pgfqpoint{4.130871in}{1.192886in}}%
\pgfpathlineto{\pgfqpoint{4.135616in}{1.160756in}}%
\pgfpathlineto{\pgfqpoint{4.140360in}{1.085788in}}%
\pgfpathlineto{\pgfqpoint{4.145104in}{0.989399in}}%
\pgfpathlineto{\pgfqpoint{4.149849in}{0.839127in}}%
\pgfpathlineto{\pgfqpoint{4.154593in}{0.930161in}}%
\pgfpathlineto{\pgfqpoint{4.159338in}{0.924806in}}%
\pgfpathlineto{\pgfqpoint{4.164082in}{0.748094in}}%
\pgfpathlineto{\pgfqpoint{4.168827in}{0.892676in}}%
\pgfpathlineto{\pgfqpoint{4.173571in}{0.876612in}}%
\pgfpathlineto{\pgfqpoint{4.178316in}{0.844482in}}%
\pgfpathlineto{\pgfqpoint{4.183060in}{0.865902in}}%
\pgfpathlineto{\pgfqpoint{4.187804in}{0.833438in}}%
\pgfpathlineto{\pgfqpoint{4.192549in}{0.903051in}}%
\pgfpathlineto{\pgfqpoint{4.197293in}{0.999440in}}%
\pgfpathlineto{\pgfqpoint{4.202038in}{1.127958in}}%
\pgfpathlineto{\pgfqpoint{4.206782in}{1.213636in}}%
\pgfpathlineto{\pgfqpoint{4.211527in}{1.256475in}}%
\pgfpathlineto{\pgfqpoint{4.216271in}{1.283250in}}%
\pgfpathlineto{\pgfqpoint{4.221016in}{1.299315in}}%
\pgfpathlineto{\pgfqpoint{4.225760in}{1.282915in}}%
\pgfpathlineto{\pgfqpoint{4.230504in}{1.245431in}}%
\pgfpathlineto{\pgfqpoint{4.235249in}{1.191882in}}%
\pgfpathlineto{\pgfqpoint{4.239993in}{1.100848in}}%
\pgfpathlineto{\pgfqpoint{4.244738in}{0.966976in}}%
\pgfpathlineto{\pgfqpoint{4.254227in}{0.827748in}}%
\pgfpathlineto{\pgfqpoint{4.258971in}{0.892007in}}%
\pgfpathlineto{\pgfqpoint{4.263715in}{0.902382in}}%
\pgfpathlineto{\pgfqpoint{4.268460in}{0.907737in}}%
\pgfpathlineto{\pgfqpoint{4.273204in}{0.795284in}}%
\pgfpathlineto{\pgfqpoint{4.282693in}{0.891672in}}%
\pgfpathlineto{\pgfqpoint{4.287438in}{0.843478in}}%
\pgfpathlineto{\pgfqpoint{4.292182in}{1.052320in}}%
\pgfpathlineto{\pgfqpoint{4.296927in}{1.180837in}}%
\pgfpathlineto{\pgfqpoint{4.301671in}{1.266181in}}%
\pgfpathlineto{\pgfqpoint{4.306415in}{1.325085in}}%
\pgfpathlineto{\pgfqpoint{4.311160in}{1.357215in}}%
\pgfpathlineto{\pgfqpoint{4.315904in}{1.367924in}}%
\pgfpathlineto{\pgfqpoint{4.320649in}{1.357215in}}%
\pgfpathlineto{\pgfqpoint{4.325393in}{1.325085in}}%
\pgfpathlineto{\pgfqpoint{4.330138in}{1.271536in}}%
\pgfpathlineto{\pgfqpoint{4.334882in}{1.191212in}}%
\pgfpathlineto{\pgfqpoint{4.339627in}{1.057005in}}%
\pgfpathlineto{\pgfqpoint{4.344371in}{0.821389in}}%
\pgfpathlineto{\pgfqpoint{4.349115in}{0.891003in}}%
\pgfpathlineto{\pgfqpoint{4.358604in}{0.939197in}}%
\pgfpathlineto{\pgfqpoint{4.363349in}{0.874938in}}%
\pgfpathlineto{\pgfqpoint{4.368093in}{0.928487in}}%
\pgfpathlineto{\pgfqpoint{4.372838in}{0.853519in}}%
\pgfpathlineto{\pgfqpoint{4.377582in}{0.880293in}}%
\pgfpathlineto{\pgfqpoint{4.382327in}{0.869249in}}%
\pgfpathlineto{\pgfqpoint{4.387071in}{1.067380in}}%
\pgfpathlineto{\pgfqpoint{4.391815in}{1.217318in}}%
\pgfpathlineto{\pgfqpoint{4.396560in}{1.308351in}}%
\pgfpathlineto{\pgfqpoint{4.401304in}{1.367255in}}%
\pgfpathlineto{\pgfqpoint{4.406049in}{1.404739in}}%
\pgfpathlineto{\pgfqpoint{4.410793in}{1.415449in}}%
\pgfpathlineto{\pgfqpoint{4.415538in}{1.410094in}}%
\pgfpathlineto{\pgfqpoint{4.420282in}{1.382985in}}%
\pgfpathlineto{\pgfqpoint{4.425026in}{1.334791in}}%
\pgfpathlineto{\pgfqpoint{4.429771in}{1.259822in}}%
\pgfpathlineto{\pgfqpoint{4.434515in}{1.136659in}}%
\pgfpathlineto{\pgfqpoint{4.439260in}{0.884979in}}%
\pgfpathlineto{\pgfqpoint{4.448749in}{0.852849in}}%
\pgfpathlineto{\pgfqpoint{4.458238in}{0.900709in}}%
\pgfpathlineto{\pgfqpoint{4.462982in}{0.820385in}}%
\pgfpathlineto{\pgfqpoint{4.467726in}{0.922128in}}%
\pgfpathlineto{\pgfqpoint{4.477215in}{0.831095in}}%
\pgfpathlineto{\pgfqpoint{4.486704in}{1.222003in}}%
\pgfpathlineto{\pgfqpoint{4.491449in}{1.318392in}}%
\pgfpathlineto{\pgfqpoint{4.496193in}{1.376961in}}%
\pgfpathlineto{\pgfqpoint{4.500938in}{1.419800in}}%
\pgfpathlineto{\pgfqpoint{4.505682in}{1.430510in}}%
\pgfpathlineto{\pgfqpoint{4.510426in}{1.430510in}}%
\pgfpathlineto{\pgfqpoint{4.515171in}{1.409090in}}%
\pgfpathlineto{\pgfqpoint{4.519915in}{1.360896in}}%
\pgfpathlineto{\pgfqpoint{4.524660in}{1.291282in}}%
\pgfpathlineto{\pgfqpoint{4.529404in}{1.178829in}}%
\pgfpathlineto{\pgfqpoint{4.534149in}{0.953588in}}%
\pgfpathlineto{\pgfqpoint{4.538893in}{0.883975in}}%
\pgfpathlineto{\pgfqpoint{4.543637in}{0.851845in}}%
\pgfpathlineto{\pgfqpoint{4.548382in}{0.846490in}}%
\pgfpathlineto{\pgfqpoint{4.553126in}{0.873265in}}%
\pgfpathlineto{\pgfqpoint{4.557871in}{0.819716in}}%
\pgfpathlineto{\pgfqpoint{4.562615in}{0.862555in}}%
\pgfpathlineto{\pgfqpoint{4.567360in}{0.819716in}}%
\pgfpathlineto{\pgfqpoint{4.572104in}{0.872930in}}%
\pgfpathlineto{\pgfqpoint{4.576849in}{0.974673in}}%
\pgfpathlineto{\pgfqpoint{4.581593in}{1.156740in}}%
\pgfpathlineto{\pgfqpoint{4.586337in}{1.258484in}}%
\pgfpathlineto{\pgfqpoint{4.591082in}{1.333452in}}%
\pgfpathlineto{\pgfqpoint{4.595826in}{1.370937in}}%
\pgfpathlineto{\pgfqpoint{4.600571in}{1.387001in}}%
\pgfpathlineto{\pgfqpoint{4.605315in}{1.386667in}}%
\pgfpathlineto{\pgfqpoint{4.610060in}{1.370602in}}%
\pgfpathlineto{\pgfqpoint{4.614804in}{1.327763in}}%
\pgfpathlineto{\pgfqpoint{4.619549in}{1.258149in}}%
\pgfpathlineto{\pgfqpoint{4.624293in}{1.151051in}}%
\pgfpathlineto{\pgfqpoint{4.629037in}{0.963629in}}%
\pgfpathlineto{\pgfqpoint{4.633782in}{0.883305in}}%
\pgfpathlineto{\pgfqpoint{4.643271in}{0.824067in}}%
\pgfpathlineto{\pgfqpoint{4.648015in}{0.931165in}}%
\pgfpathlineto{\pgfqpoint{4.652760in}{0.808002in}}%
\pgfpathlineto{\pgfqpoint{4.657504in}{0.861551in}}%
\pgfpathlineto{\pgfqpoint{4.662249in}{0.808002in}}%
\pgfpathlineto{\pgfqpoint{4.666993in}{0.893680in}}%
\pgfpathlineto{\pgfqpoint{4.671737in}{0.899035in}}%
\pgfpathlineto{\pgfqpoint{4.676482in}{1.070392in}}%
\pgfpathlineto{\pgfqpoint{4.681226in}{1.182511in}}%
\pgfpathlineto{\pgfqpoint{4.685971in}{1.257479in}}%
\pgfpathlineto{\pgfqpoint{4.690715in}{1.300319in}}%
\pgfpathlineto{\pgfqpoint{4.695460in}{1.321738in}}%
\pgfpathlineto{\pgfqpoint{4.700204in}{1.321738in}}%
\pgfpathlineto{\pgfqpoint{4.704948in}{1.305674in}}%
\pgfpathlineto{\pgfqpoint{4.709693in}{1.273544in}}%
\pgfpathlineto{\pgfqpoint{4.714437in}{1.209285in}}%
\pgfpathlineto{\pgfqpoint{4.719182in}{1.107207in}}%
\pgfpathlineto{\pgfqpoint{4.723926in}{0.930495in}}%
\pgfpathlineto{\pgfqpoint{4.728671in}{0.850172in}}%
\pgfpathlineto{\pgfqpoint{4.733415in}{0.909076in}}%
\pgfpathlineto{\pgfqpoint{4.738160in}{0.780558in}}%
\pgfpathlineto{\pgfqpoint{4.742904in}{0.834107in}}%
\pgfpathlineto{\pgfqpoint{4.747648in}{0.850172in}}%
\pgfpathlineto{\pgfqpoint{4.752393in}{0.887656in}}%
\pgfpathlineto{\pgfqpoint{4.757137in}{0.849837in}}%
\pgfpathlineto{\pgfqpoint{4.761882in}{0.881967in}}%
\pgfpathlineto{\pgfqpoint{4.766626in}{0.892676in}}%
\pgfpathlineto{\pgfqpoint{4.771371in}{1.144357in}}%
\pgfpathlineto{\pgfqpoint{4.776115in}{1.267520in}}%
\pgfpathlineto{\pgfqpoint{4.780860in}{1.342489in}}%
\pgfpathlineto{\pgfqpoint{4.785604in}{1.385328in}}%
\pgfpathlineto{\pgfqpoint{4.790348in}{1.417457in}}%
\pgfpathlineto{\pgfqpoint{4.795093in}{1.422478in}}%
\pgfpathlineto{\pgfqpoint{4.799837in}{1.406413in}}%
\pgfpathlineto{\pgfqpoint{4.804582in}{1.368928in}}%
\pgfpathlineto{\pgfqpoint{4.809326in}{1.310025in}}%
\pgfpathlineto{\pgfqpoint{4.814071in}{1.218991in}}%
\pgfpathlineto{\pgfqpoint{4.823559in}{0.881632in}}%
\pgfpathlineto{\pgfqpoint{4.828304in}{0.870922in}}%
\pgfpathlineto{\pgfqpoint{4.833048in}{0.865232in}}%
\pgfpathlineto{\pgfqpoint{4.837793in}{0.870587in}}%
\pgfpathlineto{\pgfqpoint{4.842537in}{0.838458in}}%
\pgfpathlineto{\pgfqpoint{4.847282in}{0.849168in}}%
\pgfpathlineto{\pgfqpoint{4.852026in}{0.854523in}}%
\pgfpathlineto{\pgfqpoint{4.856771in}{0.849168in}}%
\pgfpathlineto{\pgfqpoint{4.861515in}{0.892007in}}%
\pgfpathlineto{\pgfqpoint{4.866259in}{0.779554in}}%
\pgfpathlineto{\pgfqpoint{4.871004in}{0.848833in}}%
\pgfpathlineto{\pgfqpoint{4.875748in}{0.827413in}}%
\pgfpathlineto{\pgfqpoint{4.880493in}{0.848833in}}%
\pgfpathlineto{\pgfqpoint{4.885237in}{0.875608in}}%
\pgfpathlineto{\pgfqpoint{4.889982in}{0.929157in}}%
\pgfpathlineto{\pgfqpoint{4.894726in}{0.880963in}}%
\pgfpathlineto{\pgfqpoint{4.899471in}{0.800639in}}%
\pgfpathlineto{\pgfqpoint{4.908959in}{0.832434in}}%
\pgfpathlineto{\pgfqpoint{4.913704in}{0.859208in}}%
\pgfpathlineto{\pgfqpoint{4.918448in}{0.816369in}}%
\pgfpathlineto{\pgfqpoint{4.923193in}{0.827079in}}%
\pgfpathlineto{\pgfqpoint{4.927937in}{0.869918in}}%
\pgfpathlineto{\pgfqpoint{4.932682in}{0.837789in}}%
\pgfpathlineto{\pgfqpoint{4.937426in}{0.784239in}}%
\pgfpathlineto{\pgfqpoint{4.942171in}{0.864563in}}%
\pgfpathlineto{\pgfqpoint{4.946915in}{0.789260in}}%
\pgfpathlineto{\pgfqpoint{4.951659in}{0.901713in}}%
\pgfpathlineto{\pgfqpoint{4.956404in}{0.741066in}}%
\pgfpathlineto{\pgfqpoint{4.961148in}{1.094489in}}%
\pgfpathlineto{\pgfqpoint{4.965893in}{1.228362in}}%
\pgfpathlineto{\pgfqpoint{4.970637in}{1.319396in}}%
\pgfpathlineto{\pgfqpoint{4.975382in}{1.378300in}}%
\pgfpathlineto{\pgfqpoint{4.980126in}{1.410429in}}%
\pgfpathlineto{\pgfqpoint{4.984870in}{1.420804in}}%
\pgfpathlineto{\pgfqpoint{4.989615in}{1.415449in}}%
\pgfpathlineto{\pgfqpoint{4.994359in}{1.383320in}}%
\pgfpathlineto{\pgfqpoint{4.999104in}{1.335126in}}%
\pgfpathlineto{\pgfqpoint{5.003848in}{1.249447in}}%
\pgfpathlineto{\pgfqpoint{5.008593in}{1.120929in}}%
\pgfpathlineto{\pgfqpoint{5.013337in}{0.842474in}}%
\pgfpathlineto{\pgfqpoint{5.018082in}{0.858539in}}%
\pgfpathlineto{\pgfqpoint{5.022826in}{0.799300in}}%
\pgfpathlineto{\pgfqpoint{5.027570in}{0.826075in}}%
\pgfpathlineto{\pgfqpoint{5.032315in}{0.820720in}}%
\pgfpathlineto{\pgfqpoint{5.037059in}{0.820720in}}%
\pgfpathlineto{\pgfqpoint{5.041804in}{0.836785in}}%
\pgfpathlineto{\pgfqpoint{5.051293in}{0.826075in}}%
\pgfpathlineto{\pgfqpoint{5.065526in}{1.211293in}}%
\pgfpathlineto{\pgfqpoint{5.070270in}{1.275552in}}%
\pgfpathlineto{\pgfqpoint{5.075015in}{1.313037in}}%
\pgfpathlineto{\pgfqpoint{5.079759in}{1.329101in}}%
\pgfpathlineto{\pgfqpoint{5.084504in}{1.318392in}}%
\pgfpathlineto{\pgfqpoint{5.089248in}{1.296972in}}%
\pgfpathlineto{\pgfqpoint{5.093993in}{1.243423in}}%
\pgfpathlineto{\pgfqpoint{5.098737in}{1.173474in}}%
\pgfpathlineto{\pgfqpoint{5.103482in}{1.055666in}}%
\pgfpathlineto{\pgfqpoint{5.108226in}{0.836115in}}%
\pgfpathlineto{\pgfqpoint{5.112970in}{0.862890in}}%
\pgfpathlineto{\pgfqpoint{5.117715in}{0.820050in}}%
\pgfpathlineto{\pgfqpoint{5.122459in}{0.873600in}}%
\pgfpathlineto{\pgfqpoint{5.127204in}{0.820050in}}%
\pgfpathlineto{\pgfqpoint{5.131948in}{0.862890in}}%
\pgfpathlineto{\pgfqpoint{5.136693in}{0.889330in}}%
\pgfpathlineto{\pgfqpoint{5.141437in}{0.803651in}}%
\pgfpathlineto{\pgfqpoint{5.146181in}{0.846490in}}%
\pgfpathlineto{\pgfqpoint{5.160415in}{1.274883in}}%
\pgfpathlineto{\pgfqpoint{5.165159in}{1.333787in}}%
\pgfpathlineto{\pgfqpoint{5.169904in}{1.371271in}}%
\pgfpathlineto{\pgfqpoint{5.174648in}{1.387001in}}%
\pgfpathlineto{\pgfqpoint{5.179393in}{1.387001in}}%
\pgfpathlineto{\pgfqpoint{5.184137in}{1.365582in}}%
\pgfpathlineto{\pgfqpoint{5.188881in}{1.322742in}}%
\pgfpathlineto{\pgfqpoint{5.193626in}{1.247774in}}%
\pgfpathlineto{\pgfqpoint{5.198370in}{1.146030in}}%
\pgfpathlineto{\pgfqpoint{5.203115in}{0.937189in}}%
\pgfpathlineto{\pgfqpoint{5.207859in}{0.915769in}}%
\pgfpathlineto{\pgfqpoint{5.212604in}{0.824401in}}%
\pgfpathlineto{\pgfqpoint{5.217348in}{0.888660in}}%
\pgfpathlineto{\pgfqpoint{5.222093in}{0.861886in}}%
\pgfpathlineto{\pgfqpoint{5.226837in}{0.861886in}}%
\pgfpathlineto{\pgfqpoint{5.231581in}{0.867241in}}%
\pgfpathlineto{\pgfqpoint{5.236326in}{0.824401in}}%
\pgfpathlineto{\pgfqpoint{5.241070in}{0.872595in}}%
\pgfpathlineto{\pgfqpoint{5.245815in}{1.006468in}}%
\pgfpathlineto{\pgfqpoint{5.250559in}{1.188200in}}%
\pgfpathlineto{\pgfqpoint{5.255304in}{1.300653in}}%
\pgfpathlineto{\pgfqpoint{5.260048in}{1.370267in}}%
\pgfpathlineto{\pgfqpoint{5.264792in}{1.413106in}}%
\pgfpathlineto{\pgfqpoint{5.269537in}{1.429171in}}%
\pgfpathlineto{\pgfqpoint{5.274281in}{1.434526in}}%
\pgfpathlineto{\pgfqpoint{5.279026in}{1.413106in}}%
\pgfpathlineto{\pgfqpoint{5.283770in}{1.375622in}}%
\pgfpathlineto{\pgfqpoint{5.288515in}{1.305674in}}%
\pgfpathlineto{\pgfqpoint{5.293259in}{1.203930in}}%
\pgfpathlineto{\pgfqpoint{5.302748in}{0.845152in}}%
\pgfpathlineto{\pgfqpoint{5.307492in}{0.866571in}}%
\pgfpathlineto{\pgfqpoint{5.312237in}{0.845152in}}%
\pgfpathlineto{\pgfqpoint{5.321726in}{0.877281in}}%
\pgfpathlineto{\pgfqpoint{5.326470in}{0.844817in}}%
\pgfpathlineto{\pgfqpoint{5.331215in}{0.844817in}}%
\pgfpathlineto{\pgfqpoint{5.335959in}{0.812687in}}%
\pgfpathlineto{\pgfqpoint{5.340704in}{0.935850in}}%
\pgfpathlineto{\pgfqpoint{5.345448in}{1.155402in}}%
\pgfpathlineto{\pgfqpoint{5.350192in}{1.267855in}}%
\pgfpathlineto{\pgfqpoint{5.354937in}{1.342823in}}%
\pgfpathlineto{\pgfqpoint{5.359681in}{1.391017in}}%
\pgfpathlineto{\pgfqpoint{5.364426in}{1.412102in}}%
\pgfpathlineto{\pgfqpoint{5.369170in}{1.417457in}}%
\pgfpathlineto{\pgfqpoint{5.373915in}{1.401393in}}%
\pgfpathlineto{\pgfqpoint{5.378659in}{1.363908in}}%
\pgfpathlineto{\pgfqpoint{5.383404in}{1.299649in}}%
\pgfpathlineto{\pgfqpoint{5.388148in}{1.203261in}}%
\pgfpathlineto{\pgfqpoint{5.397637in}{0.881967in}}%
\pgfpathlineto{\pgfqpoint{5.402381in}{0.844148in}}%
\pgfpathlineto{\pgfqpoint{5.407126in}{0.795953in}}%
\pgfpathlineto{\pgfqpoint{5.411870in}{0.838793in}}%
\pgfpathlineto{\pgfqpoint{5.416615in}{0.790598in}}%
\pgfpathlineto{\pgfqpoint{5.421359in}{0.828083in}}%
\pgfpathlineto{\pgfqpoint{5.426103in}{0.828083in}}%
\pgfpathlineto{\pgfqpoint{5.430848in}{0.844148in}}%
\pgfpathlineto{\pgfqpoint{5.435592in}{0.854857in}}%
\pgfpathlineto{\pgfqpoint{5.440337in}{1.079429in}}%
\pgfpathlineto{\pgfqpoint{5.445081in}{1.202592in}}%
\pgfpathlineto{\pgfqpoint{5.449826in}{1.288270in}}%
\pgfpathlineto{\pgfqpoint{5.454570in}{1.331109in}}%
\pgfpathlineto{\pgfqpoint{5.459315in}{1.363239in}}%
\pgfpathlineto{\pgfqpoint{5.464059in}{1.363239in}}%
\pgfpathlineto{\pgfqpoint{5.468803in}{1.357884in}}%
\pgfpathlineto{\pgfqpoint{5.473548in}{1.320400in}}%
\pgfpathlineto{\pgfqpoint{5.478292in}{1.261161in}}%
\pgfpathlineto{\pgfqpoint{5.483037in}{1.175482in}}%
\pgfpathlineto{\pgfqpoint{5.487781in}{1.025545in}}%
\pgfpathlineto{\pgfqpoint{5.492526in}{0.784574in}}%
\pgfpathlineto{\pgfqpoint{5.497270in}{0.822059in}}%
\pgfpathlineto{\pgfqpoint{5.502015in}{0.897027in}}%
\pgfpathlineto{\pgfqpoint{5.506759in}{0.752445in}}%
\pgfpathlineto{\pgfqpoint{5.511503in}{0.789929in}}%
\pgfpathlineto{\pgfqpoint{5.516248in}{0.811014in}}%
\pgfpathlineto{\pgfqpoint{5.520992in}{0.821724in}}%
\pgfpathlineto{\pgfqpoint{5.525737in}{0.789594in}}%
\pgfpathlineto{\pgfqpoint{5.530481in}{0.837789in}}%
\pgfpathlineto{\pgfqpoint{5.539970in}{1.116244in}}%
\pgfpathlineto{\pgfqpoint{5.544714in}{1.207277in}}%
\pgfpathlineto{\pgfqpoint{5.549459in}{1.255471in}}%
\pgfpathlineto{\pgfqpoint{5.554203in}{1.287266in}}%
\pgfpathlineto{\pgfqpoint{5.558948in}{1.292621in}}%
\pgfpathlineto{\pgfqpoint{5.563692in}{1.287266in}}%
\pgfpathlineto{\pgfqpoint{5.568437in}{1.255137in}}%
\pgfpathlineto{\pgfqpoint{5.573181in}{1.201588in}}%
\pgfpathlineto{\pgfqpoint{5.577926in}{1.110554in}}%
\pgfpathlineto{\pgfqpoint{5.587414in}{0.832099in}}%
\pgfpathlineto{\pgfqpoint{5.592159in}{0.762150in}}%
\pgfpathlineto{\pgfqpoint{5.596903in}{0.762150in}}%
\pgfpathlineto{\pgfqpoint{5.601648in}{0.810345in}}%
\pgfpathlineto{\pgfqpoint{5.606392in}{0.713956in}}%
\pgfpathlineto{\pgfqpoint{5.611137in}{0.815700in}}%
\pgfpathlineto{\pgfqpoint{5.615881in}{0.751441in}}%
\pgfpathlineto{\pgfqpoint{5.620626in}{0.713956in}}%
\pgfpathlineto{\pgfqpoint{5.625370in}{0.730021in}}%
\pgfpathlineto{\pgfqpoint{5.630114in}{0.852849in}}%
\pgfpathlineto{\pgfqpoint{5.634859in}{1.008142in}}%
\pgfpathlineto{\pgfqpoint{5.639603in}{1.104530in}}%
\pgfpathlineto{\pgfqpoint{5.644348in}{1.163434in}}%
\pgfpathlineto{\pgfqpoint{5.649092in}{1.195563in}}%
\pgfpathlineto{\pgfqpoint{5.653837in}{1.206273in}}%
\pgfpathlineto{\pgfqpoint{5.658581in}{1.200918in}}%
\pgfpathlineto{\pgfqpoint{5.663326in}{1.174144in}}%
\pgfpathlineto{\pgfqpoint{5.668070in}{1.120260in}}%
\pgfpathlineto{\pgfqpoint{5.672814in}{1.039936in}}%
\pgfpathlineto{\pgfqpoint{5.682303in}{0.729352in}}%
\pgfpathlineto{\pgfqpoint{5.687048in}{0.729352in}}%
\pgfpathlineto{\pgfqpoint{5.691792in}{0.734707in}}%
\pgfpathlineto{\pgfqpoint{5.696537in}{0.729352in}}%
\pgfpathlineto{\pgfqpoint{5.706025in}{0.798631in}}%
\pgfpathlineto{\pgfqpoint{5.710770in}{0.761146in}}%
\pgfpathlineto{\pgfqpoint{5.715514in}{0.734372in}}%
\pgfpathlineto{\pgfqpoint{5.720259in}{0.729017in}}%
\pgfpathlineto{\pgfqpoint{5.725003in}{0.761146in}}%
\pgfpathlineto{\pgfqpoint{5.729748in}{0.895019in}}%
\pgfpathlineto{\pgfqpoint{5.734492in}{0.986053in}}%
\pgfpathlineto{\pgfqpoint{5.739237in}{1.050311in}}%
\pgfpathlineto{\pgfqpoint{5.743981in}{1.092816in}}%
\pgfpathlineto{\pgfqpoint{5.748725in}{1.108881in}}%
\pgfpathlineto{\pgfqpoint{5.753470in}{1.103526in}}%
\pgfpathlineto{\pgfqpoint{5.758214in}{1.076751in}}%
\pgfpathlineto{\pgfqpoint{5.762959in}{1.028557in}}%
\pgfpathlineto{\pgfqpoint{5.767703in}{0.948234in}}%
\pgfpathlineto{\pgfqpoint{5.772448in}{0.803651in}}%
\pgfpathlineto{\pgfqpoint{5.777192in}{0.734037in}}%
\pgfpathlineto{\pgfqpoint{5.781937in}{0.696218in}}%
\pgfpathlineto{\pgfqpoint{5.786681in}{0.717638in}}%
\pgfpathlineto{\pgfqpoint{5.791425in}{0.696218in}}%
\pgfpathlineto{\pgfqpoint{5.796170in}{0.712283in}}%
\pgfpathlineto{\pgfqpoint{5.810403in}{0.712283in}}%
\pgfpathlineto{\pgfqpoint{5.815148in}{0.728348in}}%
\pgfpathlineto{\pgfqpoint{5.819892in}{0.765497in}}%
\pgfpathlineto{\pgfqpoint{5.824637in}{0.781562in}}%
\pgfpathlineto{\pgfqpoint{5.838870in}{0.974339in}}%
\pgfpathlineto{\pgfqpoint{5.843614in}{0.995758in}}%
\pgfpathlineto{\pgfqpoint{5.848359in}{0.979694in}}%
\pgfpathlineto{\pgfqpoint{5.853103in}{0.968984in}}%
\pgfpathlineto{\pgfqpoint{5.857848in}{0.920455in}}%
\pgfpathlineto{\pgfqpoint{5.867336in}{0.759808in}}%
\pgfpathlineto{\pgfqpoint{5.872081in}{0.733033in}}%
\pgfpathlineto{\pgfqpoint{5.876825in}{0.727678in}}%
\pgfpathlineto{\pgfqpoint{5.881570in}{0.711614in}}%
\pgfpathlineto{\pgfqpoint{5.886314in}{0.727678in}}%
\pgfpathlineto{\pgfqpoint{5.891059in}{0.695549in}}%
\pgfpathlineto{\pgfqpoint{5.895803in}{0.727344in}}%
\pgfpathlineto{\pgfqpoint{5.900548in}{0.732698in}}%
\pgfpathlineto{\pgfqpoint{5.905292in}{0.748763in}}%
\pgfpathlineto{\pgfqpoint{5.910036in}{0.695214in}}%
\pgfpathlineto{\pgfqpoint{5.919525in}{0.695214in}}%
\pgfpathlineto{\pgfqpoint{5.924270in}{0.727344in}}%
\pgfpathlineto{\pgfqpoint{5.929014in}{0.727344in}}%
\pgfpathlineto{\pgfqpoint{5.933759in}{0.818042in}}%
\pgfpathlineto{\pgfqpoint{5.938503in}{0.801978in}}%
\pgfpathlineto{\pgfqpoint{5.943248in}{0.850172in}}%
\pgfpathlineto{\pgfqpoint{5.943248in}{0.850172in}}%
\pgfusepath{stroke}%
\end{pgfscope}%
\begin{pgfscope}%
\pgfsetrectcap%
\pgfsetmiterjoin%
\pgfsetlinewidth{0.803000pt}%
\definecolor{currentstroke}{rgb}{0.000000,0.000000,0.000000}%
\pgfsetstrokecolor{currentstroke}%
\pgfsetdash{}{0pt}%
\pgfpathmoveto{\pgfqpoint{3.833026in}{0.565123in}}%
\pgfpathlineto{\pgfqpoint{3.833026in}{1.475926in}}%
\pgfusepath{stroke}%
\end{pgfscope}%
\begin{pgfscope}%
\pgfsetrectcap%
\pgfsetmiterjoin%
\pgfsetlinewidth{0.803000pt}%
\definecolor{currentstroke}{rgb}{0.000000,0.000000,0.000000}%
\pgfsetstrokecolor{currentstroke}%
\pgfsetdash{}{0pt}%
\pgfpathmoveto{\pgfqpoint{5.941666in}{0.565123in}}%
\pgfpathlineto{\pgfqpoint{5.941666in}{1.475926in}}%
\pgfusepath{stroke}%
\end{pgfscope}%
\begin{pgfscope}%
\pgfsetrectcap%
\pgfsetmiterjoin%
\pgfsetlinewidth{0.803000pt}%
\definecolor{currentstroke}{rgb}{0.000000,0.000000,0.000000}%
\pgfsetstrokecolor{currentstroke}%
\pgfsetdash{}{0pt}%
\pgfpathmoveto{\pgfqpoint{3.833026in}{0.565123in}}%
\pgfpathlineto{\pgfqpoint{5.941666in}{0.565123in}}%
\pgfusepath{stroke}%
\end{pgfscope}%
\begin{pgfscope}%
\pgfsetrectcap%
\pgfsetmiterjoin%
\pgfsetlinewidth{0.803000pt}%
\definecolor{currentstroke}{rgb}{0.000000,0.000000,0.000000}%
\pgfsetstrokecolor{currentstroke}%
\pgfsetdash{}{0pt}%
\pgfpathmoveto{\pgfqpoint{3.833026in}{1.475926in}}%
\pgfpathlineto{\pgfqpoint{5.941666in}{1.475926in}}%
\pgfusepath{stroke}%
\end{pgfscope}%
\begin{pgfscope}%
\definecolor{textcolor}{rgb}{0.000000,0.000000,0.000000}%
\pgfsetstrokecolor{textcolor}%
\pgfsetfillcolor{textcolor}%
\pgftext[x=4.887346in,y=1.559260in,,base]{\color{textcolor}\rmfamily\fontsize{12.000000}{14.400000}\selectfont 18,116 kHz}%
\end{pgfscope}%
\end{pgfpicture}%
\makeatother%
\endgroup%

		\caption{Measured frequency content of a FM modulated 500 kHz sine wave with modulation frequencies at solutions of the Bessel functions.}
		\label{fig:446}
	\end{figure}

	A sinusoidal carrier with frequency $f_T = 500\,\mathrm{kHz}$ was frequency modulated using a peak frequency deviation of $\Delta f_T = \pm 100\,\mathrm{kHz}$. The modulation was sinusoidal with modulation frequencies
	\[
	f_m = \frac{100\,\mathrm{kHz}}{x_n}, \quad x_n \in \{2.40483,\;3.83171,\;5.13562,\;5.52008\},
	\]
	which correspond to the first zeros of the Bessel function $J_0(\beta)$. The FM modulation index is defined as
	\[
	\beta = \frac{\Delta f_T}{f_m}% = x_n.
	\]
	
	For sinusoidal FM, the spectrum consists of discrete lines at
	\[
	f = f_T \pm n f_m, \quad n = 0,1,\dots,7,
	\]
	with amplitudes proportional to the Bessel functions $J_n(\beta)$. At the selected modulation frequencies, $J_0(\beta)=0$, leading to complete carrier suppression, which is clearly observed in the measured spectra. The amplitudes of the sidebands decrease with increasing order $n$, in agreement with the theoretical behavior of the Bessel functions.
	
	\begin{table}[h]
		\centering
		\caption{Modulation frequencies and qualitative spectral observations}
		\begin{tabular}{c c c l}
			\hline
			$x_n$ & $f_m$ [kHz] & $\beta$ & Observed FM spectrum \\ \hline
			2.40483 & 41.597 & 2.40 & Few dominant sidebands, wide spacing \\
			3.83171 & 26.098 & 3.83 & Increased number of visible sidebands \\
			5.13562 & 19.472 & 5.14 & Dense spectrum, many significant sidebands \\
			5.52008 & 18.116 & 5.52 & Very wide spectrum, strong higher-order bands \\ \hline
		\end{tabular}
	\end{table}
	
	A variation of the modulation frequency from $10\,\mathrm{kHz}$ to $50\,\mathrm{kHz}$ showed that decreasing $f_m$ (increasing $\beta$) leads to a broader spectrum with more sidebands, while increasing $f_m$ reduces the occupied bandwidth. Repeating the measurement with a carrier frequency of $200\,\mathrm{kHz}$ resulted in an identical spectral structure shifted in frequency, confirming that the relative sideband amplitudes depend only on the modulation index $\beta$. Overall, the measured amplitude ratios of the sidebands show good agreement with the values predicted by the Bessel functions.

	\begin{figure}[H]
		\centering
		%% Creator: Matplotlib, PGF backend
%%
%% To include the figure in your LaTeX document, write
%%   \input{<filename>.pgf}
%%
%% Make sure the required packages are loaded in your preamble
%%   \usepackage{pgf}
%%
%% Also ensure that all the required font packages are loaded; for instance,
%% the lmodern package is sometimes necessary when using math font.
%%   \usepackage{lmodern}
%%
%% Figures using additional raster images can only be included by \input if
%% they are in the same directory as the main LaTeX file. For loading figures
%% from other directories you can use the `import` package
%%   \usepackage{import}
%%
%% and then include the figures with
%%   \import{<path to file>}{<filename>.pgf}
%%
%% Matplotlib used the following preamble
%%
\begingroup%
\makeatletter%
\begin{pgfpicture}%
\pgfpathrectangle{\pgfpointorigin}{\pgfqpoint{6.300000in}{3.500000in}}%
\pgfusepath{use as bounding box, clip}%
\begin{pgfscope}%
\pgfsetbuttcap%
\pgfsetmiterjoin%
\definecolor{currentfill}{rgb}{1.000000,1.000000,1.000000}%
\pgfsetfillcolor{currentfill}%
\pgfsetlinewidth{0.000000pt}%
\definecolor{currentstroke}{rgb}{1.000000,1.000000,1.000000}%
\pgfsetstrokecolor{currentstroke}%
\pgfsetdash{}{0pt}%
\pgfpathmoveto{\pgfqpoint{0.000000in}{0.000000in}}%
\pgfpathlineto{\pgfqpoint{6.300000in}{0.000000in}}%
\pgfpathlineto{\pgfqpoint{6.300000in}{3.500000in}}%
\pgfpathlineto{\pgfqpoint{0.000000in}{3.500000in}}%
\pgfpathlineto{\pgfqpoint{0.000000in}{0.000000in}}%
\pgfpathclose%
\pgfusepath{fill}%
\end{pgfscope}%
\begin{pgfscope}%
\pgfsetbuttcap%
\pgfsetmiterjoin%
\definecolor{currentfill}{rgb}{1.000000,1.000000,1.000000}%
\pgfsetfillcolor{currentfill}%
\pgfsetlinewidth{0.000000pt}%
\definecolor{currentstroke}{rgb}{0.000000,0.000000,0.000000}%
\pgfsetstrokecolor{currentstroke}%
\pgfsetstrokeopacity{0.000000}%
\pgfsetdash{}{0pt}%
\pgfpathmoveto{\pgfqpoint{0.758026in}{2.240123in}}%
\pgfpathlineto{\pgfqpoint{2.866666in}{2.240123in}}%
\pgfpathlineto{\pgfqpoint{2.866666in}{3.150926in}}%
\pgfpathlineto{\pgfqpoint{0.758026in}{3.150926in}}%
\pgfpathlineto{\pgfqpoint{0.758026in}{2.240123in}}%
\pgfpathclose%
\pgfusepath{fill}%
\end{pgfscope}%
\begin{pgfscope}%
\pgfpathrectangle{\pgfqpoint{0.758026in}{2.240123in}}{\pgfqpoint{2.108640in}{0.910803in}}%
\pgfusepath{clip}%
\pgfsetbuttcap%
\pgfsetroundjoin%
\pgfsetlinewidth{0.501875pt}%
\definecolor{currentstroke}{rgb}{0.690196,0.690196,0.690196}%
\pgfsetstrokecolor{currentstroke}%
\pgfsetstrokeopacity{0.700000}%
\pgfsetdash{{1.850000pt}{0.800000pt}}{0.000000pt}%
\pgfpathmoveto{\pgfqpoint{0.758026in}{2.240123in}}%
\pgfpathlineto{\pgfqpoint{0.758026in}{3.150926in}}%
\pgfusepath{stroke}%
\end{pgfscope}%
\begin{pgfscope}%
\pgfsetbuttcap%
\pgfsetroundjoin%
\definecolor{currentfill}{rgb}{0.000000,0.000000,0.000000}%
\pgfsetfillcolor{currentfill}%
\pgfsetlinewidth{0.803000pt}%
\definecolor{currentstroke}{rgb}{0.000000,0.000000,0.000000}%
\pgfsetstrokecolor{currentstroke}%
\pgfsetdash{}{0pt}%
\pgfsys@defobject{currentmarker}{\pgfqpoint{0.000000in}{-0.048611in}}{\pgfqpoint{0.000000in}{0.000000in}}{%
\pgfpathmoveto{\pgfqpoint{0.000000in}{0.000000in}}%
\pgfpathlineto{\pgfqpoint{0.000000in}{-0.048611in}}%
\pgfusepath{stroke,fill}%
}%
\begin{pgfscope}%
\pgfsys@transformshift{0.758026in}{2.240123in}%
\pgfsys@useobject{currentmarker}{}%
\end{pgfscope}%
\end{pgfscope}%
\begin{pgfscope}%
\definecolor{textcolor}{rgb}{0.000000,0.000000,0.000000}%
\pgfsetstrokecolor{textcolor}%
\pgfsetfillcolor{textcolor}%
\pgftext[x=0.758026in,y=2.142901in,,top]{\color{textcolor}\rmfamily\fontsize{10.000000}{12.000000}\selectfont \(\displaystyle {100000}\)}%
\end{pgfscope}%
\begin{pgfscope}%
\pgfpathrectangle{\pgfqpoint{0.758026in}{2.240123in}}{\pgfqpoint{2.108640in}{0.910803in}}%
\pgfusepath{clip}%
\pgfsetbuttcap%
\pgfsetroundjoin%
\pgfsetlinewidth{0.501875pt}%
\definecolor{currentstroke}{rgb}{0.690196,0.690196,0.690196}%
\pgfsetstrokecolor{currentstroke}%
\pgfsetstrokeopacity{0.700000}%
\pgfsetdash{{1.850000pt}{0.800000pt}}{0.000000pt}%
\pgfpathmoveto{\pgfqpoint{1.285186in}{2.240123in}}%
\pgfpathlineto{\pgfqpoint{1.285186in}{3.150926in}}%
\pgfusepath{stroke}%
\end{pgfscope}%
\begin{pgfscope}%
\pgfsetbuttcap%
\pgfsetroundjoin%
\definecolor{currentfill}{rgb}{0.000000,0.000000,0.000000}%
\pgfsetfillcolor{currentfill}%
\pgfsetlinewidth{0.803000pt}%
\definecolor{currentstroke}{rgb}{0.000000,0.000000,0.000000}%
\pgfsetstrokecolor{currentstroke}%
\pgfsetdash{}{0pt}%
\pgfsys@defobject{currentmarker}{\pgfqpoint{0.000000in}{-0.048611in}}{\pgfqpoint{0.000000in}{0.000000in}}{%
\pgfpathmoveto{\pgfqpoint{0.000000in}{0.000000in}}%
\pgfpathlineto{\pgfqpoint{0.000000in}{-0.048611in}}%
\pgfusepath{stroke,fill}%
}%
\begin{pgfscope}%
\pgfsys@transformshift{1.285186in}{2.240123in}%
\pgfsys@useobject{currentmarker}{}%
\end{pgfscope}%
\end{pgfscope}%
\begin{pgfscope}%
\definecolor{textcolor}{rgb}{0.000000,0.000000,0.000000}%
\pgfsetstrokecolor{textcolor}%
\pgfsetfillcolor{textcolor}%
\pgftext[x=1.285186in,y=2.142901in,,top]{\color{textcolor}\rmfamily\fontsize{10.000000}{12.000000}\selectfont \(\displaystyle {200000}\)}%
\end{pgfscope}%
\begin{pgfscope}%
\pgfpathrectangle{\pgfqpoint{0.758026in}{2.240123in}}{\pgfqpoint{2.108640in}{0.910803in}}%
\pgfusepath{clip}%
\pgfsetbuttcap%
\pgfsetroundjoin%
\pgfsetlinewidth{0.501875pt}%
\definecolor{currentstroke}{rgb}{0.690196,0.690196,0.690196}%
\pgfsetstrokecolor{currentstroke}%
\pgfsetstrokeopacity{0.700000}%
\pgfsetdash{{1.850000pt}{0.800000pt}}{0.000000pt}%
\pgfpathmoveto{\pgfqpoint{1.812346in}{2.240123in}}%
\pgfpathlineto{\pgfqpoint{1.812346in}{3.150926in}}%
\pgfusepath{stroke}%
\end{pgfscope}%
\begin{pgfscope}%
\pgfsetbuttcap%
\pgfsetroundjoin%
\definecolor{currentfill}{rgb}{0.000000,0.000000,0.000000}%
\pgfsetfillcolor{currentfill}%
\pgfsetlinewidth{0.803000pt}%
\definecolor{currentstroke}{rgb}{0.000000,0.000000,0.000000}%
\pgfsetstrokecolor{currentstroke}%
\pgfsetdash{}{0pt}%
\pgfsys@defobject{currentmarker}{\pgfqpoint{0.000000in}{-0.048611in}}{\pgfqpoint{0.000000in}{0.000000in}}{%
\pgfpathmoveto{\pgfqpoint{0.000000in}{0.000000in}}%
\pgfpathlineto{\pgfqpoint{0.000000in}{-0.048611in}}%
\pgfusepath{stroke,fill}%
}%
\begin{pgfscope}%
\pgfsys@transformshift{1.812346in}{2.240123in}%
\pgfsys@useobject{currentmarker}{}%
\end{pgfscope}%
\end{pgfscope}%
\begin{pgfscope}%
\definecolor{textcolor}{rgb}{0.000000,0.000000,0.000000}%
\pgfsetstrokecolor{textcolor}%
\pgfsetfillcolor{textcolor}%
\pgftext[x=1.812346in,y=2.142901in,,top]{\color{textcolor}\rmfamily\fontsize{10.000000}{12.000000}\selectfont \(\displaystyle {300000}\)}%
\end{pgfscope}%
\begin{pgfscope}%
\pgfpathrectangle{\pgfqpoint{0.758026in}{2.240123in}}{\pgfqpoint{2.108640in}{0.910803in}}%
\pgfusepath{clip}%
\pgfsetbuttcap%
\pgfsetroundjoin%
\pgfsetlinewidth{0.501875pt}%
\definecolor{currentstroke}{rgb}{0.690196,0.690196,0.690196}%
\pgfsetstrokecolor{currentstroke}%
\pgfsetstrokeopacity{0.700000}%
\pgfsetdash{{1.850000pt}{0.800000pt}}{0.000000pt}%
\pgfpathmoveto{\pgfqpoint{2.339506in}{2.240123in}}%
\pgfpathlineto{\pgfqpoint{2.339506in}{3.150926in}}%
\pgfusepath{stroke}%
\end{pgfscope}%
\begin{pgfscope}%
\pgfsetbuttcap%
\pgfsetroundjoin%
\definecolor{currentfill}{rgb}{0.000000,0.000000,0.000000}%
\pgfsetfillcolor{currentfill}%
\pgfsetlinewidth{0.803000pt}%
\definecolor{currentstroke}{rgb}{0.000000,0.000000,0.000000}%
\pgfsetstrokecolor{currentstroke}%
\pgfsetdash{}{0pt}%
\pgfsys@defobject{currentmarker}{\pgfqpoint{0.000000in}{-0.048611in}}{\pgfqpoint{0.000000in}{0.000000in}}{%
\pgfpathmoveto{\pgfqpoint{0.000000in}{0.000000in}}%
\pgfpathlineto{\pgfqpoint{0.000000in}{-0.048611in}}%
\pgfusepath{stroke,fill}%
}%
\begin{pgfscope}%
\pgfsys@transformshift{2.339506in}{2.240123in}%
\pgfsys@useobject{currentmarker}{}%
\end{pgfscope}%
\end{pgfscope}%
\begin{pgfscope}%
\definecolor{textcolor}{rgb}{0.000000,0.000000,0.000000}%
\pgfsetstrokecolor{textcolor}%
\pgfsetfillcolor{textcolor}%
\pgftext[x=2.339506in,y=2.142901in,,top]{\color{textcolor}\rmfamily\fontsize{10.000000}{12.000000}\selectfont \(\displaystyle {400000}\)}%
\end{pgfscope}%
\begin{pgfscope}%
\pgfpathrectangle{\pgfqpoint{0.758026in}{2.240123in}}{\pgfqpoint{2.108640in}{0.910803in}}%
\pgfusepath{clip}%
\pgfsetbuttcap%
\pgfsetroundjoin%
\pgfsetlinewidth{0.501875pt}%
\definecolor{currentstroke}{rgb}{0.690196,0.690196,0.690196}%
\pgfsetstrokecolor{currentstroke}%
\pgfsetstrokeopacity{0.700000}%
\pgfsetdash{{1.850000pt}{0.800000pt}}{0.000000pt}%
\pgfpathmoveto{\pgfqpoint{2.866666in}{2.240123in}}%
\pgfpathlineto{\pgfqpoint{2.866666in}{3.150926in}}%
\pgfusepath{stroke}%
\end{pgfscope}%
\begin{pgfscope}%
\pgfsetbuttcap%
\pgfsetroundjoin%
\definecolor{currentfill}{rgb}{0.000000,0.000000,0.000000}%
\pgfsetfillcolor{currentfill}%
\pgfsetlinewidth{0.803000pt}%
\definecolor{currentstroke}{rgb}{0.000000,0.000000,0.000000}%
\pgfsetstrokecolor{currentstroke}%
\pgfsetdash{}{0pt}%
\pgfsys@defobject{currentmarker}{\pgfqpoint{0.000000in}{-0.048611in}}{\pgfqpoint{0.000000in}{0.000000in}}{%
\pgfpathmoveto{\pgfqpoint{0.000000in}{0.000000in}}%
\pgfpathlineto{\pgfqpoint{0.000000in}{-0.048611in}}%
\pgfusepath{stroke,fill}%
}%
\begin{pgfscope}%
\pgfsys@transformshift{2.866666in}{2.240123in}%
\pgfsys@useobject{currentmarker}{}%
\end{pgfscope}%
\end{pgfscope}%
\begin{pgfscope}%
\definecolor{textcolor}{rgb}{0.000000,0.000000,0.000000}%
\pgfsetstrokecolor{textcolor}%
\pgfsetfillcolor{textcolor}%
\pgftext[x=2.866666in,y=2.142901in,,top]{\color{textcolor}\rmfamily\fontsize{10.000000}{12.000000}\selectfont \(\displaystyle {500000}\)}%
\end{pgfscope}%
\begin{pgfscope}%
\definecolor{textcolor}{rgb}{0.000000,0.000000,0.000000}%
\pgfsetstrokecolor{textcolor}%
\pgfsetfillcolor{textcolor}%
\pgftext[x=1.812346in,y=1.963889in,,top]{\color{textcolor}\rmfamily\fontsize{10.000000}{12.000000}\selectfont Frequency (Hz)}%
\end{pgfscope}%
\begin{pgfscope}%
\pgfpathrectangle{\pgfqpoint{0.758026in}{2.240123in}}{\pgfqpoint{2.108640in}{0.910803in}}%
\pgfusepath{clip}%
\pgfsetbuttcap%
\pgfsetroundjoin%
\pgfsetlinewidth{0.501875pt}%
\definecolor{currentstroke}{rgb}{0.690196,0.690196,0.690196}%
\pgfsetstrokecolor{currentstroke}%
\pgfsetstrokeopacity{0.700000}%
\pgfsetdash{{1.850000pt}{0.800000pt}}{0.000000pt}%
\pgfpathmoveto{\pgfqpoint{0.758026in}{2.405979in}}%
\pgfpathlineto{\pgfqpoint{2.866666in}{2.405979in}}%
\pgfusepath{stroke}%
\end{pgfscope}%
\begin{pgfscope}%
\pgfsetbuttcap%
\pgfsetroundjoin%
\definecolor{currentfill}{rgb}{0.000000,0.000000,0.000000}%
\pgfsetfillcolor{currentfill}%
\pgfsetlinewidth{0.803000pt}%
\definecolor{currentstroke}{rgb}{0.000000,0.000000,0.000000}%
\pgfsetstrokecolor{currentstroke}%
\pgfsetdash{}{0pt}%
\pgfsys@defobject{currentmarker}{\pgfqpoint{-0.048611in}{0.000000in}}{\pgfqpoint{-0.000000in}{0.000000in}}{%
\pgfpathmoveto{\pgfqpoint{-0.000000in}{0.000000in}}%
\pgfpathlineto{\pgfqpoint{-0.048611in}{0.000000in}}%
\pgfusepath{stroke,fill}%
}%
\begin{pgfscope}%
\pgfsys@transformshift{0.758026in}{2.405979in}%
\pgfsys@useobject{currentmarker}{}%
\end{pgfscope}%
\end{pgfscope}%
\begin{pgfscope}%
\definecolor{textcolor}{rgb}{0.000000,0.000000,0.000000}%
\pgfsetstrokecolor{textcolor}%
\pgfsetfillcolor{textcolor}%
\pgftext[x=0.344444in, y=2.357753in, left, base]{\color{textcolor}\rmfamily\fontsize{10.000000}{12.000000}\selectfont \(\displaystyle {\ensuremath{-}100}\)}%
\end{pgfscope}%
\begin{pgfscope}%
\pgfpathrectangle{\pgfqpoint{0.758026in}{2.240123in}}{\pgfqpoint{2.108640in}{0.910803in}}%
\pgfusepath{clip}%
\pgfsetbuttcap%
\pgfsetroundjoin%
\pgfsetlinewidth{0.501875pt}%
\definecolor{currentstroke}{rgb}{0.690196,0.690196,0.690196}%
\pgfsetstrokecolor{currentstroke}%
\pgfsetstrokeopacity{0.700000}%
\pgfsetdash{{1.850000pt}{0.800000pt}}{0.000000pt}%
\pgfpathmoveto{\pgfqpoint{0.758026in}{2.885806in}}%
\pgfpathlineto{\pgfqpoint{2.866666in}{2.885806in}}%
\pgfusepath{stroke}%
\end{pgfscope}%
\begin{pgfscope}%
\pgfsetbuttcap%
\pgfsetroundjoin%
\definecolor{currentfill}{rgb}{0.000000,0.000000,0.000000}%
\pgfsetfillcolor{currentfill}%
\pgfsetlinewidth{0.803000pt}%
\definecolor{currentstroke}{rgb}{0.000000,0.000000,0.000000}%
\pgfsetstrokecolor{currentstroke}%
\pgfsetdash{}{0pt}%
\pgfsys@defobject{currentmarker}{\pgfqpoint{-0.048611in}{0.000000in}}{\pgfqpoint{-0.000000in}{0.000000in}}{%
\pgfpathmoveto{\pgfqpoint{-0.000000in}{0.000000in}}%
\pgfpathlineto{\pgfqpoint{-0.048611in}{0.000000in}}%
\pgfusepath{stroke,fill}%
}%
\begin{pgfscope}%
\pgfsys@transformshift{0.758026in}{2.885806in}%
\pgfsys@useobject{currentmarker}{}%
\end{pgfscope}%
\end{pgfscope}%
\begin{pgfscope}%
\definecolor{textcolor}{rgb}{0.000000,0.000000,0.000000}%
\pgfsetstrokecolor{textcolor}%
\pgfsetfillcolor{textcolor}%
\pgftext[x=0.413889in, y=2.837581in, left, base]{\color{textcolor}\rmfamily\fontsize{10.000000}{12.000000}\selectfont \(\displaystyle {\ensuremath{-}50}\)}%
\end{pgfscope}%
\begin{pgfscope}%
\definecolor{textcolor}{rgb}{0.000000,0.000000,0.000000}%
\pgfsetstrokecolor{textcolor}%
\pgfsetfillcolor{textcolor}%
\pgftext[x=0.288889in,y=2.695525in,,bottom,rotate=90.000000]{\color{textcolor}\rmfamily\fontsize{10.000000}{12.000000}\selectfont Amplitude (dB)}%
\end{pgfscope}%
\begin{pgfscope}%
\pgfpathrectangle{\pgfqpoint{0.758026in}{2.240123in}}{\pgfqpoint{2.108640in}{0.910803in}}%
\pgfusepath{clip}%
\pgfsetrectcap%
\pgfsetroundjoin%
\pgfsetlinewidth{1.003750pt}%
\definecolor{currentstroke}{rgb}{0.000000,0.000000,0.000000}%
\pgfsetstrokecolor{currentstroke}%
\pgfsetdash{}{0pt}%
\pgfpathmoveto{\pgfqpoint{0.758026in}{2.539431in}}%
\pgfpathlineto{\pgfqpoint{0.762770in}{2.592212in}}%
\pgfpathlineto{\pgfqpoint{0.767515in}{2.611405in}}%
\pgfpathlineto{\pgfqpoint{0.772259in}{2.568221in}}%
\pgfpathlineto{\pgfqpoint{0.777003in}{2.568221in}}%
\pgfpathlineto{\pgfqpoint{0.781748in}{2.592212in}}%
\pgfpathlineto{\pgfqpoint{0.786492in}{2.577817in}}%
\pgfpathlineto{\pgfqpoint{0.791237in}{2.606607in}}%
\pgfpathlineto{\pgfqpoint{0.795981in}{2.582615in}}%
\pgfpathlineto{\pgfqpoint{0.800726in}{2.582315in}}%
\pgfpathlineto{\pgfqpoint{0.805470in}{2.596710in}}%
\pgfpathlineto{\pgfqpoint{0.810215in}{2.630298in}}%
\pgfpathlineto{\pgfqpoint{0.814959in}{2.567921in}}%
\pgfpathlineto{\pgfqpoint{0.819703in}{2.601509in}}%
\pgfpathlineto{\pgfqpoint{0.824448in}{2.587114in}}%
\pgfpathlineto{\pgfqpoint{0.829192in}{2.577517in}}%
\pgfpathlineto{\pgfqpoint{0.833937in}{2.615903in}}%
\pgfpathlineto{\pgfqpoint{0.838681in}{2.884307in}}%
\pgfpathlineto{\pgfqpoint{0.843426in}{3.071440in}}%
\pgfpathlineto{\pgfqpoint{0.848170in}{3.090633in}}%
\pgfpathlineto{\pgfqpoint{0.852915in}{2.989869in}}%
\pgfpathlineto{\pgfqpoint{0.857659in}{2.663586in}}%
\pgfpathlineto{\pgfqpoint{0.862403in}{2.629998in}}%
\pgfpathlineto{\pgfqpoint{0.867148in}{2.610805in}}%
\pgfpathlineto{\pgfqpoint{0.871892in}{2.572419in}}%
\pgfpathlineto{\pgfqpoint{0.876637in}{2.582016in}}%
\pgfpathlineto{\pgfqpoint{0.881381in}{2.567321in}}%
\pgfpathlineto{\pgfqpoint{0.886126in}{2.605707in}}%
\pgfpathlineto{\pgfqpoint{0.890870in}{2.610505in}}%
\pgfpathlineto{\pgfqpoint{0.895614in}{2.572119in}}%
\pgfpathlineto{\pgfqpoint{0.900359in}{2.634497in}}%
\pgfpathlineto{\pgfqpoint{0.905103in}{2.653690in}}%
\pgfpathlineto{\pgfqpoint{0.909848in}{2.648892in}}%
\pgfpathlineto{\pgfqpoint{0.914592in}{2.596110in}}%
\pgfpathlineto{\pgfqpoint{0.919337in}{2.624600in}}%
\pgfpathlineto{\pgfqpoint{0.924081in}{2.567021in}}%
\pgfpathlineto{\pgfqpoint{0.928826in}{2.528635in}}%
\pgfpathlineto{\pgfqpoint{0.933570in}{2.576617in}}%
\pgfpathlineto{\pgfqpoint{0.938314in}{2.571819in}}%
\pgfpathlineto{\pgfqpoint{0.943059in}{2.586214in}}%
\pgfpathlineto{\pgfqpoint{0.947803in}{2.605407in}}%
\pgfpathlineto{\pgfqpoint{0.952548in}{2.562223in}}%
\pgfpathlineto{\pgfqpoint{0.957292in}{2.542730in}}%
\pgfpathlineto{\pgfqpoint{0.962037in}{2.595511in}}%
\pgfpathlineto{\pgfqpoint{0.966781in}{2.595511in}}%
\pgfpathlineto{\pgfqpoint{0.971526in}{2.576318in}}%
\pgfpathlineto{\pgfqpoint{0.976270in}{2.595511in}}%
\pgfpathlineto{\pgfqpoint{0.981014in}{2.561923in}}%
\pgfpathlineto{\pgfqpoint{0.985759in}{2.542730in}}%
\pgfpathlineto{\pgfqpoint{0.990503in}{2.576018in}}%
\pgfpathlineto{\pgfqpoint{0.995248in}{2.600009in}}%
\pgfpathlineto{\pgfqpoint{0.999992in}{2.604807in}}%
\pgfpathlineto{\pgfqpoint{1.004737in}{2.647992in}}%
\pgfpathlineto{\pgfqpoint{1.009481in}{2.595211in}}%
\pgfpathlineto{\pgfqpoint{1.014225in}{2.604807in}}%
\pgfpathlineto{\pgfqpoint{1.018970in}{2.580816in}}%
\pgfpathlineto{\pgfqpoint{1.023714in}{2.590413in}}%
\pgfpathlineto{\pgfqpoint{1.028459in}{2.552026in}}%
\pgfpathlineto{\pgfqpoint{1.033203in}{2.590113in}}%
\pgfpathlineto{\pgfqpoint{1.037948in}{2.570920in}}%
\pgfpathlineto{\pgfqpoint{1.042692in}{2.618902in}}%
\pgfpathlineto{\pgfqpoint{1.047437in}{2.537332in}}%
\pgfpathlineto{\pgfqpoint{1.052181in}{2.594911in}}%
\pgfpathlineto{\pgfqpoint{1.056925in}{2.801237in}}%
\pgfpathlineto{\pgfqpoint{1.061670in}{3.065142in}}%
\pgfpathlineto{\pgfqpoint{1.066414in}{3.108327in}}%
\pgfpathlineto{\pgfqpoint{1.071159in}{3.055246in}}%
\pgfpathlineto{\pgfqpoint{1.075903in}{2.752954in}}%
\pgfpathlineto{\pgfqpoint{1.080648in}{2.594611in}}%
\pgfpathlineto{\pgfqpoint{1.085392in}{2.541830in}}%
\pgfpathlineto{\pgfqpoint{1.090137in}{2.541830in}}%
\pgfpathlineto{\pgfqpoint{1.099625in}{2.585014in}}%
\pgfpathlineto{\pgfqpoint{1.104370in}{2.584715in}}%
\pgfpathlineto{\pgfqpoint{1.109114in}{2.608706in}}%
\pgfpathlineto{\pgfqpoint{1.113859in}{2.536732in}}%
\pgfpathlineto{\pgfqpoint{1.118603in}{2.594311in}}%
\pgfpathlineto{\pgfqpoint{1.123348in}{2.584715in}}%
\pgfpathlineto{\pgfqpoint{1.128092in}{2.613504in}}%
\pgfpathlineto{\pgfqpoint{1.132837in}{2.565521in}}%
\pgfpathlineto{\pgfqpoint{1.137581in}{2.560723in}}%
\pgfpathlineto{\pgfqpoint{1.142325in}{2.594311in}}%
\pgfpathlineto{\pgfqpoint{1.147070in}{2.598810in}}%
\pgfpathlineto{\pgfqpoint{1.151814in}{2.608406in}}%
\pgfpathlineto{\pgfqpoint{1.156559in}{2.589213in}}%
\pgfpathlineto{\pgfqpoint{1.161303in}{2.550827in}}%
\pgfpathlineto{\pgfqpoint{1.166048in}{2.618003in}}%
\pgfpathlineto{\pgfqpoint{1.170792in}{2.570020in}}%
\pgfpathlineto{\pgfqpoint{1.175536in}{2.536432in}}%
\pgfpathlineto{\pgfqpoint{1.180281in}{2.608106in}}%
\pgfpathlineto{\pgfqpoint{1.185025in}{2.502544in}}%
\pgfpathlineto{\pgfqpoint{1.189770in}{2.536132in}}%
\pgfpathlineto{\pgfqpoint{1.194514in}{2.521737in}}%
\pgfpathlineto{\pgfqpoint{1.199259in}{2.540930in}}%
\pgfpathlineto{\pgfqpoint{1.204003in}{2.588913in}}%
\pgfpathlineto{\pgfqpoint{1.208748in}{2.536132in}}%
\pgfpathlineto{\pgfqpoint{1.213492in}{2.569720in}}%
\pgfpathlineto{\pgfqpoint{1.218236in}{2.560123in}}%
\pgfpathlineto{\pgfqpoint{1.222981in}{2.555025in}}%
\pgfpathlineto{\pgfqpoint{1.227725in}{2.593411in}}%
\pgfpathlineto{\pgfqpoint{1.232470in}{2.535832in}}%
\pgfpathlineto{\pgfqpoint{1.237214in}{2.516639in}}%
\pgfpathlineto{\pgfqpoint{1.241959in}{2.655789in}}%
\pgfpathlineto{\pgfqpoint{1.246703in}{2.650991in}}%
\pgfpathlineto{\pgfqpoint{1.251448in}{2.579017in}}%
\pgfpathlineto{\pgfqpoint{1.256192in}{2.617103in}}%
\pgfpathlineto{\pgfqpoint{1.260936in}{2.583515in}}%
\pgfpathlineto{\pgfqpoint{1.270425in}{2.583515in}}%
\pgfpathlineto{\pgfqpoint{1.275170in}{2.530734in}}%
\pgfpathlineto{\pgfqpoint{1.279914in}{2.588313in}}%
\pgfpathlineto{\pgfqpoint{1.284659in}{2.617103in}}%
\pgfpathlineto{\pgfqpoint{1.294148in}{2.564022in}}%
\pgfpathlineto{\pgfqpoint{1.298892in}{2.607207in}}%
\pgfpathlineto{\pgfqpoint{1.303636in}{2.583215in}}%
\pgfpathlineto{\pgfqpoint{1.308381in}{2.549627in}}%
\pgfpathlineto{\pgfqpoint{1.322614in}{2.621601in}}%
\pgfpathlineto{\pgfqpoint{1.327359in}{2.520838in}}%
\pgfpathlineto{\pgfqpoint{1.332103in}{2.510941in}}%
\pgfpathlineto{\pgfqpoint{1.336847in}{2.602108in}}%
\pgfpathlineto{\pgfqpoint{1.341592in}{2.534932in}}%
\pgfpathlineto{\pgfqpoint{1.346336in}{2.573319in}}%
\pgfpathlineto{\pgfqpoint{1.351081in}{2.597310in}}%
\pgfpathlineto{\pgfqpoint{1.355825in}{2.602108in}}%
\pgfpathlineto{\pgfqpoint{1.360570in}{2.597310in}}%
\pgfpathlineto{\pgfqpoint{1.365314in}{2.582915in}}%
\pgfpathlineto{\pgfqpoint{1.370059in}{2.582615in}}%
\pgfpathlineto{\pgfqpoint{1.374803in}{2.568221in}}%
\pgfpathlineto{\pgfqpoint{1.384292in}{2.621002in}}%
\pgfpathlineto{\pgfqpoint{1.389036in}{2.563422in}}%
\pgfpathlineto{\pgfqpoint{1.393781in}{2.549027in}}%
\pgfpathlineto{\pgfqpoint{1.398525in}{2.505843in}}%
\pgfpathlineto{\pgfqpoint{1.403270in}{2.597010in}}%
\pgfpathlineto{\pgfqpoint{1.408014in}{2.615903in}}%
\pgfpathlineto{\pgfqpoint{1.412759in}{2.591912in}}%
\pgfpathlineto{\pgfqpoint{1.417503in}{2.596710in}}%
\pgfpathlineto{\pgfqpoint{1.422247in}{2.596710in}}%
\pgfpathlineto{\pgfqpoint{1.426992in}{2.563122in}}%
\pgfpathlineto{\pgfqpoint{1.431736in}{2.572719in}}%
\pgfpathlineto{\pgfqpoint{1.436481in}{2.548728in}}%
\pgfpathlineto{\pgfqpoint{1.441225in}{2.567921in}}%
\pgfpathlineto{\pgfqpoint{1.445970in}{2.562822in}}%
\pgfpathlineto{\pgfqpoint{1.450714in}{2.591612in}}%
\pgfpathlineto{\pgfqpoint{1.455458in}{2.572419in}}%
\pgfpathlineto{\pgfqpoint{1.460203in}{2.682779in}}%
\pgfpathlineto{\pgfqpoint{1.464947in}{2.730762in}}%
\pgfpathlineto{\pgfqpoint{1.469692in}{2.586814in}}%
\pgfpathlineto{\pgfqpoint{1.474436in}{2.596410in}}%
\pgfpathlineto{\pgfqpoint{1.479181in}{2.596410in}}%
\pgfpathlineto{\pgfqpoint{1.483925in}{2.576917in}}%
\pgfpathlineto{\pgfqpoint{1.488670in}{2.581716in}}%
\pgfpathlineto{\pgfqpoint{1.493414in}{2.600909in}}%
\pgfpathlineto{\pgfqpoint{1.498158in}{2.960780in}}%
\pgfpathlineto{\pgfqpoint{1.502903in}{3.109526in}}%
\pgfpathlineto{\pgfqpoint{1.507647in}{3.109526in}}%
\pgfpathlineto{\pgfqpoint{1.512392in}{2.960780in}}%
\pgfpathlineto{\pgfqpoint{1.517136in}{2.615304in}}%
\pgfpathlineto{\pgfqpoint{1.521881in}{2.605407in}}%
\pgfpathlineto{\pgfqpoint{1.526625in}{2.581416in}}%
\pgfpathlineto{\pgfqpoint{1.531370in}{2.605407in}}%
\pgfpathlineto{\pgfqpoint{1.536114in}{2.605407in}}%
\pgfpathlineto{\pgfqpoint{1.540858in}{2.581416in}}%
\pgfpathlineto{\pgfqpoint{1.550347in}{2.610205in}}%
\pgfpathlineto{\pgfqpoint{1.555092in}{2.605407in}}%
\pgfpathlineto{\pgfqpoint{1.559836in}{2.595511in}}%
\pgfpathlineto{\pgfqpoint{1.564581in}{2.576318in}}%
\pgfpathlineto{\pgfqpoint{1.569325in}{2.571519in}}%
\pgfpathlineto{\pgfqpoint{1.574070in}{2.571519in}}%
\pgfpathlineto{\pgfqpoint{1.578814in}{2.581116in}}%
\pgfpathlineto{\pgfqpoint{1.583558in}{2.629099in}}%
\pgfpathlineto{\pgfqpoint{1.588303in}{2.581116in}}%
\pgfpathlineto{\pgfqpoint{1.597792in}{2.547228in}}%
\pgfpathlineto{\pgfqpoint{1.602536in}{2.609606in}}%
\pgfpathlineto{\pgfqpoint{1.607281in}{2.576018in}}%
\pgfpathlineto{\pgfqpoint{1.612025in}{2.585614in}}%
\pgfpathlineto{\pgfqpoint{1.616769in}{2.590413in}}%
\pgfpathlineto{\pgfqpoint{1.621514in}{2.624000in}}%
\pgfpathlineto{\pgfqpoint{1.626258in}{2.590413in}}%
\pgfpathlineto{\pgfqpoint{1.631003in}{2.542430in}}%
\pgfpathlineto{\pgfqpoint{1.635747in}{2.537332in}}%
\pgfpathlineto{\pgfqpoint{1.640492in}{2.585314in}}%
\pgfpathlineto{\pgfqpoint{1.645236in}{2.590113in}}%
\pgfpathlineto{\pgfqpoint{1.649981in}{2.580516in}}%
\pgfpathlineto{\pgfqpoint{1.654725in}{2.609306in}}%
\pgfpathlineto{\pgfqpoint{1.659469in}{2.566121in}}%
\pgfpathlineto{\pgfqpoint{1.664214in}{2.556525in}}%
\pgfpathlineto{\pgfqpoint{1.668958in}{2.551726in}}%
\pgfpathlineto{\pgfqpoint{1.673703in}{2.580216in}}%
\pgfpathlineto{\pgfqpoint{1.683192in}{2.714568in}}%
\pgfpathlineto{\pgfqpoint{1.692681in}{2.522637in}}%
\pgfpathlineto{\pgfqpoint{1.697425in}{2.589813in}}%
\pgfpathlineto{\pgfqpoint{1.702169in}{2.609006in}}%
\pgfpathlineto{\pgfqpoint{1.706914in}{2.580216in}}%
\pgfpathlineto{\pgfqpoint{1.711658in}{2.627899in}}%
\pgfpathlineto{\pgfqpoint{1.721147in}{3.074139in}}%
\pgfpathlineto{\pgfqpoint{1.725892in}{3.102928in}}%
\pgfpathlineto{\pgfqpoint{1.730636in}{3.021358in}}%
\pgfpathlineto{\pgfqpoint{1.735380in}{2.632697in}}%
\pgfpathlineto{\pgfqpoint{1.740125in}{2.575118in}}%
\pgfpathlineto{\pgfqpoint{1.744869in}{2.584715in}}%
\pgfpathlineto{\pgfqpoint{1.749614in}{2.570020in}}%
\pgfpathlineto{\pgfqpoint{1.754358in}{2.594011in}}%
\pgfpathlineto{\pgfqpoint{1.763847in}{2.555625in}}%
\pgfpathlineto{\pgfqpoint{1.768592in}{2.570020in}}%
\pgfpathlineto{\pgfqpoint{1.773336in}{2.560423in}}%
\pgfpathlineto{\pgfqpoint{1.778080in}{2.608406in}}%
\pgfpathlineto{\pgfqpoint{1.782825in}{2.579616in}}%
\pgfpathlineto{\pgfqpoint{1.787569in}{2.588913in}}%
\pgfpathlineto{\pgfqpoint{1.792314in}{2.622501in}}%
\pgfpathlineto{\pgfqpoint{1.801803in}{2.536132in}}%
\pgfpathlineto{\pgfqpoint{1.806547in}{2.603308in}}%
\pgfpathlineto{\pgfqpoint{1.811292in}{2.516939in}}%
\pgfpathlineto{\pgfqpoint{1.816036in}{2.603308in}}%
\pgfpathlineto{\pgfqpoint{1.820780in}{2.593711in}}%
\pgfpathlineto{\pgfqpoint{1.825525in}{2.545429in}}%
\pgfpathlineto{\pgfqpoint{1.830269in}{2.607806in}}%
\pgfpathlineto{\pgfqpoint{1.835014in}{2.588613in}}%
\pgfpathlineto{\pgfqpoint{1.839758in}{2.564622in}}%
\pgfpathlineto{\pgfqpoint{1.844503in}{2.603008in}}%
\pgfpathlineto{\pgfqpoint{1.849247in}{2.559824in}}%
\pgfpathlineto{\pgfqpoint{1.853992in}{2.579017in}}%
\pgfpathlineto{\pgfqpoint{1.858736in}{2.559824in}}%
\pgfpathlineto{\pgfqpoint{1.863480in}{2.564322in}}%
\pgfpathlineto{\pgfqpoint{1.868225in}{2.636296in}}%
\pgfpathlineto{\pgfqpoint{1.872969in}{2.607506in}}%
\pgfpathlineto{\pgfqpoint{1.877714in}{2.569120in}}%
\pgfpathlineto{\pgfqpoint{1.887203in}{2.549927in}}%
\pgfpathlineto{\pgfqpoint{1.891947in}{2.573918in}}%
\pgfpathlineto{\pgfqpoint{1.896691in}{2.564322in}}%
\pgfpathlineto{\pgfqpoint{1.906180in}{2.592812in}}%
\pgfpathlineto{\pgfqpoint{1.910925in}{2.544829in}}%
\pgfpathlineto{\pgfqpoint{1.915669in}{2.607207in}}%
\pgfpathlineto{\pgfqpoint{1.920414in}{2.559224in}}%
\pgfpathlineto{\pgfqpoint{1.925158in}{2.540031in}}%
\pgfpathlineto{\pgfqpoint{1.929903in}{2.626400in}}%
\pgfpathlineto{\pgfqpoint{1.934647in}{2.655189in}}%
\pgfpathlineto{\pgfqpoint{1.939391in}{2.981172in}}%
\pgfpathlineto{\pgfqpoint{1.944136in}{3.033953in}}%
\pgfpathlineto{\pgfqpoint{1.948880in}{2.995567in}}%
\pgfpathlineto{\pgfqpoint{1.953625in}{2.731662in}}%
\pgfpathlineto{\pgfqpoint{1.958369in}{2.611705in}}%
\pgfpathlineto{\pgfqpoint{1.963114in}{2.554126in}}%
\pgfpathlineto{\pgfqpoint{1.967858in}{2.568520in}}%
\pgfpathlineto{\pgfqpoint{1.972603in}{2.578117in}}%
\pgfpathlineto{\pgfqpoint{1.977347in}{2.597010in}}%
\pgfpathlineto{\pgfqpoint{1.982091in}{2.539431in}}%
\pgfpathlineto{\pgfqpoint{1.991580in}{2.558624in}}%
\pgfpathlineto{\pgfqpoint{1.996325in}{2.563422in}}%
\pgfpathlineto{\pgfqpoint{2.001069in}{2.549027in}}%
\pgfpathlineto{\pgfqpoint{2.005814in}{2.592212in}}%
\pgfpathlineto{\pgfqpoint{2.010558in}{2.549027in}}%
\pgfpathlineto{\pgfqpoint{2.015303in}{2.558324in}}%
\pgfpathlineto{\pgfqpoint{2.020047in}{2.596710in}}%
\pgfpathlineto{\pgfqpoint{2.024791in}{2.591912in}}%
\pgfpathlineto{\pgfqpoint{2.039025in}{2.591912in}}%
\pgfpathlineto{\pgfqpoint{2.043769in}{2.587114in}}%
\pgfpathlineto{\pgfqpoint{2.048514in}{2.534333in}}%
\pgfpathlineto{\pgfqpoint{2.053258in}{2.586814in}}%
\pgfpathlineto{\pgfqpoint{2.058002in}{2.514840in}}%
\pgfpathlineto{\pgfqpoint{2.062747in}{2.591612in}}%
\pgfpathlineto{\pgfqpoint{2.067491in}{2.543629in}}%
\pgfpathlineto{\pgfqpoint{2.072236in}{2.529235in}}%
\pgfpathlineto{\pgfqpoint{2.076980in}{2.524436in}}%
\pgfpathlineto{\pgfqpoint{2.081725in}{2.567621in}}%
\pgfpathlineto{\pgfqpoint{2.086469in}{2.543629in}}%
\pgfpathlineto{\pgfqpoint{2.091214in}{2.586514in}}%
\pgfpathlineto{\pgfqpoint{2.095958in}{2.567321in}}%
\pgfpathlineto{\pgfqpoint{2.100702in}{2.528935in}}%
\pgfpathlineto{\pgfqpoint{2.105447in}{2.600909in}}%
\pgfpathlineto{\pgfqpoint{2.110191in}{2.576917in}}%
\pgfpathlineto{\pgfqpoint{2.114936in}{2.600909in}}%
\pgfpathlineto{\pgfqpoint{2.119680in}{2.663286in}}%
\pgfpathlineto{\pgfqpoint{2.124425in}{2.629698in}}%
\pgfpathlineto{\pgfqpoint{2.129169in}{2.547828in}}%
\pgfpathlineto{\pgfqpoint{2.133914in}{2.523836in}}%
\pgfpathlineto{\pgfqpoint{2.138658in}{2.528635in}}%
\pgfpathlineto{\pgfqpoint{2.143402in}{2.571819in}}%
\pgfpathlineto{\pgfqpoint{2.148147in}{2.533433in}}%
\pgfpathlineto{\pgfqpoint{2.152891in}{2.605407in}}%
\pgfpathlineto{\pgfqpoint{2.157636in}{2.845321in}}%
\pgfpathlineto{\pgfqpoint{2.162380in}{2.945785in}}%
\pgfpathlineto{\pgfqpoint{2.167125in}{2.931390in}}%
\pgfpathlineto{\pgfqpoint{2.176613in}{2.552326in}}%
\pgfpathlineto{\pgfqpoint{2.181358in}{2.557124in}}%
\pgfpathlineto{\pgfqpoint{2.186102in}{2.566721in}}%
\pgfpathlineto{\pgfqpoint{2.190847in}{2.523537in}}%
\pgfpathlineto{\pgfqpoint{2.195591in}{2.552326in}}%
\pgfpathlineto{\pgfqpoint{2.200336in}{2.557124in}}%
\pgfpathlineto{\pgfqpoint{2.205080in}{2.537631in}}%
\pgfpathlineto{\pgfqpoint{2.209825in}{2.590413in}}%
\pgfpathlineto{\pgfqpoint{2.219313in}{2.499245in}}%
\pgfpathlineto{\pgfqpoint{2.224058in}{2.542430in}}%
\pgfpathlineto{\pgfqpoint{2.228802in}{2.532833in}}%
\pgfpathlineto{\pgfqpoint{2.233547in}{2.561623in}}%
\pgfpathlineto{\pgfqpoint{2.238291in}{2.552026in}}%
\pgfpathlineto{\pgfqpoint{2.243036in}{2.532533in}}%
\pgfpathlineto{\pgfqpoint{2.247780in}{2.599709in}}%
\pgfpathlineto{\pgfqpoint{2.252525in}{2.566121in}}%
\pgfpathlineto{\pgfqpoint{2.257269in}{2.513340in}}%
\pgfpathlineto{\pgfqpoint{2.262013in}{2.542130in}}%
\pgfpathlineto{\pgfqpoint{2.266758in}{2.561323in}}%
\pgfpathlineto{\pgfqpoint{2.271502in}{2.556525in}}%
\pgfpathlineto{\pgfqpoint{2.276247in}{2.532233in}}%
\pgfpathlineto{\pgfqpoint{2.280991in}{2.556225in}}%
\pgfpathlineto{\pgfqpoint{2.285736in}{2.589813in}}%
\pgfpathlineto{\pgfqpoint{2.290480in}{2.575418in}}%
\pgfpathlineto{\pgfqpoint{2.295225in}{2.575418in}}%
\pgfpathlineto{\pgfqpoint{2.299969in}{2.599409in}}%
\pgfpathlineto{\pgfqpoint{2.304713in}{2.561023in}}%
\pgfpathlineto{\pgfqpoint{2.309458in}{2.513040in}}%
\pgfpathlineto{\pgfqpoint{2.314202in}{2.570620in}}%
\pgfpathlineto{\pgfqpoint{2.318947in}{2.527135in}}%
\pgfpathlineto{\pgfqpoint{2.323691in}{2.579916in}}%
\pgfpathlineto{\pgfqpoint{2.328436in}{2.599109in}}%
\pgfpathlineto{\pgfqpoint{2.333180in}{2.570320in}}%
\pgfpathlineto{\pgfqpoint{2.342669in}{2.632697in}}%
\pgfpathlineto{\pgfqpoint{2.347413in}{2.541530in}}%
\pgfpathlineto{\pgfqpoint{2.352158in}{2.527135in}}%
\pgfpathlineto{\pgfqpoint{2.356902in}{2.575118in}}%
\pgfpathlineto{\pgfqpoint{2.361647in}{2.541230in}}%
\pgfpathlineto{\pgfqpoint{2.366391in}{2.536432in}}%
\pgfpathlineto{\pgfqpoint{2.371136in}{2.498046in}}%
\pgfpathlineto{\pgfqpoint{2.380624in}{2.824329in}}%
\pgfpathlineto{\pgfqpoint{2.385369in}{2.829127in}}%
\pgfpathlineto{\pgfqpoint{2.394858in}{2.574518in}}%
\pgfpathlineto{\pgfqpoint{2.399602in}{2.588913in}}%
\pgfpathlineto{\pgfqpoint{2.404347in}{2.564922in}}%
\pgfpathlineto{\pgfqpoint{2.409091in}{2.507342in}}%
\pgfpathlineto{\pgfqpoint{2.413836in}{2.531334in}}%
\pgfpathlineto{\pgfqpoint{2.418580in}{2.512141in}}%
\pgfpathlineto{\pgfqpoint{2.423324in}{2.555325in}}%
\pgfpathlineto{\pgfqpoint{2.428069in}{2.555325in}}%
\pgfpathlineto{\pgfqpoint{2.432813in}{2.526236in}}%
\pgfpathlineto{\pgfqpoint{2.437558in}{2.555025in}}%
\pgfpathlineto{\pgfqpoint{2.442302in}{2.603008in}}%
\pgfpathlineto{\pgfqpoint{2.447047in}{2.535832in}}%
\pgfpathlineto{\pgfqpoint{2.451791in}{2.569420in}}%
\pgfpathlineto{\pgfqpoint{2.456535in}{2.564622in}}%
\pgfpathlineto{\pgfqpoint{2.461280in}{2.507042in}}%
\pgfpathlineto{\pgfqpoint{2.466024in}{2.531034in}}%
\pgfpathlineto{\pgfqpoint{2.470769in}{2.516639in}}%
\pgfpathlineto{\pgfqpoint{2.475513in}{2.511541in}}%
\pgfpathlineto{\pgfqpoint{2.480258in}{2.549927in}}%
\pgfpathlineto{\pgfqpoint{2.485002in}{2.540331in}}%
\pgfpathlineto{\pgfqpoint{2.489747in}{2.583515in}}%
\pgfpathlineto{\pgfqpoint{2.494491in}{2.559524in}}%
\pgfpathlineto{\pgfqpoint{2.499235in}{2.559524in}}%
\pgfpathlineto{\pgfqpoint{2.503980in}{2.540331in}}%
\pgfpathlineto{\pgfqpoint{2.508724in}{2.497146in}}%
\pgfpathlineto{\pgfqpoint{2.513469in}{2.511241in}}%
\pgfpathlineto{\pgfqpoint{2.518213in}{2.544829in}}%
\pgfpathlineto{\pgfqpoint{2.527702in}{2.492048in}}%
\pgfpathlineto{\pgfqpoint{2.532447in}{2.549627in}}%
\pgfpathlineto{\pgfqpoint{2.537191in}{2.573619in}}%
\pgfpathlineto{\pgfqpoint{2.541935in}{2.554425in}}%
\pgfpathlineto{\pgfqpoint{2.546680in}{2.573319in}}%
\pgfpathlineto{\pgfqpoint{2.551424in}{2.515739in}}%
\pgfpathlineto{\pgfqpoint{2.556169in}{2.602108in}}%
\pgfpathlineto{\pgfqpoint{2.560913in}{2.635696in}}%
\pgfpathlineto{\pgfqpoint{2.565658in}{2.578117in}}%
\pgfpathlineto{\pgfqpoint{2.570402in}{2.549327in}}%
\pgfpathlineto{\pgfqpoint{2.575147in}{2.510941in}}%
\pgfpathlineto{\pgfqpoint{2.579891in}{2.530134in}}%
\pgfpathlineto{\pgfqpoint{2.584635in}{2.525036in}}%
\pgfpathlineto{\pgfqpoint{2.589380in}{2.539431in}}%
\pgfpathlineto{\pgfqpoint{2.594124in}{2.573019in}}%
\pgfpathlineto{\pgfqpoint{2.598869in}{2.668984in}}%
\pgfpathlineto{\pgfqpoint{2.603613in}{2.702572in}}%
\pgfpathlineto{\pgfqpoint{2.608358in}{2.611405in}}%
\pgfpathlineto{\pgfqpoint{2.613102in}{2.549027in}}%
\pgfpathlineto{\pgfqpoint{2.617846in}{2.515439in}}%
\pgfpathlineto{\pgfqpoint{2.627335in}{2.563122in}}%
\pgfpathlineto{\pgfqpoint{2.632080in}{2.543929in}}%
\pgfpathlineto{\pgfqpoint{2.636824in}{2.481552in}}%
\pgfpathlineto{\pgfqpoint{2.641569in}{2.481552in}}%
\pgfpathlineto{\pgfqpoint{2.646313in}{2.534333in}}%
\pgfpathlineto{\pgfqpoint{2.651058in}{2.543929in}}%
\pgfpathlineto{\pgfqpoint{2.655802in}{2.563122in}}%
\pgfpathlineto{\pgfqpoint{2.660546in}{2.534033in}}%
\pgfpathlineto{\pgfqpoint{2.665291in}{2.538831in}}%
\pgfpathlineto{\pgfqpoint{2.670035in}{2.548428in}}%
\pgfpathlineto{\pgfqpoint{2.674780in}{2.505243in}}%
\pgfpathlineto{\pgfqpoint{2.679524in}{2.476453in}}%
\pgfpathlineto{\pgfqpoint{2.684269in}{2.519638in}}%
\pgfpathlineto{\pgfqpoint{2.689013in}{2.514840in}}%
\pgfpathlineto{\pgfqpoint{2.693758in}{2.553226in}}%
\pgfpathlineto{\pgfqpoint{2.698502in}{2.534033in}}%
\pgfpathlineto{\pgfqpoint{2.707991in}{2.490548in}}%
\pgfpathlineto{\pgfqpoint{2.712735in}{2.524136in}}%
\pgfpathlineto{\pgfqpoint{2.717480in}{2.543329in}}%
\pgfpathlineto{\pgfqpoint{2.722224in}{2.476154in}}%
\pgfpathlineto{\pgfqpoint{2.726969in}{2.509742in}}%
\pgfpathlineto{\pgfqpoint{2.731713in}{2.461759in}}%
\pgfpathlineto{\pgfqpoint{2.736458in}{2.562223in}}%
\pgfpathlineto{\pgfqpoint{2.741202in}{2.495047in}}%
\pgfpathlineto{\pgfqpoint{2.745946in}{2.504643in}}%
\pgfpathlineto{\pgfqpoint{2.750691in}{2.562223in}}%
\pgfpathlineto{\pgfqpoint{2.755435in}{2.552626in}}%
\pgfpathlineto{\pgfqpoint{2.760180in}{2.495047in}}%
\pgfpathlineto{\pgfqpoint{2.764924in}{2.552626in}}%
\pgfpathlineto{\pgfqpoint{2.769669in}{2.523836in}}%
\pgfpathlineto{\pgfqpoint{2.774413in}{2.533133in}}%
\pgfpathlineto{\pgfqpoint{2.779157in}{2.590712in}}%
\pgfpathlineto{\pgfqpoint{2.783902in}{2.537931in}}%
\pgfpathlineto{\pgfqpoint{2.788646in}{2.523537in}}%
\pgfpathlineto{\pgfqpoint{2.793391in}{2.504343in}}%
\pgfpathlineto{\pgfqpoint{2.798135in}{2.480352in}}%
\pgfpathlineto{\pgfqpoint{2.802880in}{2.494747in}}%
\pgfpathlineto{\pgfqpoint{2.807624in}{2.504343in}}%
\pgfpathlineto{\pgfqpoint{2.812369in}{2.489949in}}%
\pgfpathlineto{\pgfqpoint{2.817113in}{2.494447in}}%
\pgfpathlineto{\pgfqpoint{2.821857in}{2.552026in}}%
\pgfpathlineto{\pgfqpoint{2.826602in}{2.532833in}}%
\pgfpathlineto{\pgfqpoint{2.836091in}{2.460859in}}%
\pgfpathlineto{\pgfqpoint{2.840835in}{2.518438in}}%
\pgfpathlineto{\pgfqpoint{2.845580in}{2.494447in}}%
\pgfpathlineto{\pgfqpoint{2.850324in}{2.513340in}}%
\pgfpathlineto{\pgfqpoint{2.855069in}{2.489349in}}%
\pgfpathlineto{\pgfqpoint{2.859813in}{2.479752in}}%
\pgfpathlineto{\pgfqpoint{2.864557in}{2.479752in}}%
\pgfpathlineto{\pgfqpoint{2.869302in}{2.537332in}}%
\pgfpathlineto{\pgfqpoint{2.869302in}{2.537332in}}%
\pgfusepath{stroke}%
\end{pgfscope}%
\begin{pgfscope}%
\pgfsetrectcap%
\pgfsetmiterjoin%
\pgfsetlinewidth{0.803000pt}%
\definecolor{currentstroke}{rgb}{0.000000,0.000000,0.000000}%
\pgfsetstrokecolor{currentstroke}%
\pgfsetdash{}{0pt}%
\pgfpathmoveto{\pgfqpoint{0.758026in}{2.240123in}}%
\pgfpathlineto{\pgfqpoint{0.758026in}{3.150926in}}%
\pgfusepath{stroke}%
\end{pgfscope}%
\begin{pgfscope}%
\pgfsetrectcap%
\pgfsetmiterjoin%
\pgfsetlinewidth{0.803000pt}%
\definecolor{currentstroke}{rgb}{0.000000,0.000000,0.000000}%
\pgfsetstrokecolor{currentstroke}%
\pgfsetdash{}{0pt}%
\pgfpathmoveto{\pgfqpoint{2.866666in}{2.240123in}}%
\pgfpathlineto{\pgfqpoint{2.866666in}{3.150926in}}%
\pgfusepath{stroke}%
\end{pgfscope}%
\begin{pgfscope}%
\pgfsetrectcap%
\pgfsetmiterjoin%
\pgfsetlinewidth{0.803000pt}%
\definecolor{currentstroke}{rgb}{0.000000,0.000000,0.000000}%
\pgfsetstrokecolor{currentstroke}%
\pgfsetdash{}{0pt}%
\pgfpathmoveto{\pgfqpoint{0.758026in}{2.240123in}}%
\pgfpathlineto{\pgfqpoint{2.866666in}{2.240123in}}%
\pgfusepath{stroke}%
\end{pgfscope}%
\begin{pgfscope}%
\pgfsetrectcap%
\pgfsetmiterjoin%
\pgfsetlinewidth{0.803000pt}%
\definecolor{currentstroke}{rgb}{0.000000,0.000000,0.000000}%
\pgfsetstrokecolor{currentstroke}%
\pgfsetdash{}{0pt}%
\pgfpathmoveto{\pgfqpoint{0.758026in}{3.150926in}}%
\pgfpathlineto{\pgfqpoint{2.866666in}{3.150926in}}%
\pgfusepath{stroke}%
\end{pgfscope}%
\begin{pgfscope}%
\definecolor{textcolor}{rgb}{0.000000,0.000000,0.000000}%
\pgfsetstrokecolor{textcolor}%
\pgfsetfillcolor{textcolor}%
\pgftext[x=1.812346in,y=3.234260in,,base]{\color{textcolor}\rmfamily\fontsize{12.000000}{14.400000}\selectfont 41,597 kHz}%
\end{pgfscope}%
\begin{pgfscope}%
\pgfsetbuttcap%
\pgfsetmiterjoin%
\definecolor{currentfill}{rgb}{1.000000,1.000000,1.000000}%
\pgfsetfillcolor{currentfill}%
\pgfsetlinewidth{0.000000pt}%
\definecolor{currentstroke}{rgb}{0.000000,0.000000,0.000000}%
\pgfsetstrokecolor{currentstroke}%
\pgfsetstrokeopacity{0.000000}%
\pgfsetdash{}{0pt}%
\pgfpathmoveto{\pgfqpoint{3.833026in}{2.240123in}}%
\pgfpathlineto{\pgfqpoint{5.941666in}{2.240123in}}%
\pgfpathlineto{\pgfqpoint{5.941666in}{3.150926in}}%
\pgfpathlineto{\pgfqpoint{3.833026in}{3.150926in}}%
\pgfpathlineto{\pgfqpoint{3.833026in}{2.240123in}}%
\pgfpathclose%
\pgfusepath{fill}%
\end{pgfscope}%
\begin{pgfscope}%
\pgfpathrectangle{\pgfqpoint{3.833026in}{2.240123in}}{\pgfqpoint{2.108640in}{0.910803in}}%
\pgfusepath{clip}%
\pgfsetbuttcap%
\pgfsetroundjoin%
\pgfsetlinewidth{0.501875pt}%
\definecolor{currentstroke}{rgb}{0.690196,0.690196,0.690196}%
\pgfsetstrokecolor{currentstroke}%
\pgfsetstrokeopacity{0.700000}%
\pgfsetdash{{1.850000pt}{0.800000pt}}{0.000000pt}%
\pgfpathmoveto{\pgfqpoint{3.833026in}{2.240123in}}%
\pgfpathlineto{\pgfqpoint{3.833026in}{3.150926in}}%
\pgfusepath{stroke}%
\end{pgfscope}%
\begin{pgfscope}%
\pgfsetbuttcap%
\pgfsetroundjoin%
\definecolor{currentfill}{rgb}{0.000000,0.000000,0.000000}%
\pgfsetfillcolor{currentfill}%
\pgfsetlinewidth{0.803000pt}%
\definecolor{currentstroke}{rgb}{0.000000,0.000000,0.000000}%
\pgfsetstrokecolor{currentstroke}%
\pgfsetdash{}{0pt}%
\pgfsys@defobject{currentmarker}{\pgfqpoint{0.000000in}{-0.048611in}}{\pgfqpoint{0.000000in}{0.000000in}}{%
\pgfpathmoveto{\pgfqpoint{0.000000in}{0.000000in}}%
\pgfpathlineto{\pgfqpoint{0.000000in}{-0.048611in}}%
\pgfusepath{stroke,fill}%
}%
\begin{pgfscope}%
\pgfsys@transformshift{3.833026in}{2.240123in}%
\pgfsys@useobject{currentmarker}{}%
\end{pgfscope}%
\end{pgfscope}%
\begin{pgfscope}%
\definecolor{textcolor}{rgb}{0.000000,0.000000,0.000000}%
\pgfsetstrokecolor{textcolor}%
\pgfsetfillcolor{textcolor}%
\pgftext[x=3.833026in,y=2.142901in,,top]{\color{textcolor}\rmfamily\fontsize{10.000000}{12.000000}\selectfont \(\displaystyle {100000}\)}%
\end{pgfscope}%
\begin{pgfscope}%
\pgfpathrectangle{\pgfqpoint{3.833026in}{2.240123in}}{\pgfqpoint{2.108640in}{0.910803in}}%
\pgfusepath{clip}%
\pgfsetbuttcap%
\pgfsetroundjoin%
\pgfsetlinewidth{0.501875pt}%
\definecolor{currentstroke}{rgb}{0.690196,0.690196,0.690196}%
\pgfsetstrokecolor{currentstroke}%
\pgfsetstrokeopacity{0.700000}%
\pgfsetdash{{1.850000pt}{0.800000pt}}{0.000000pt}%
\pgfpathmoveto{\pgfqpoint{4.360186in}{2.240123in}}%
\pgfpathlineto{\pgfqpoint{4.360186in}{3.150926in}}%
\pgfusepath{stroke}%
\end{pgfscope}%
\begin{pgfscope}%
\pgfsetbuttcap%
\pgfsetroundjoin%
\definecolor{currentfill}{rgb}{0.000000,0.000000,0.000000}%
\pgfsetfillcolor{currentfill}%
\pgfsetlinewidth{0.803000pt}%
\definecolor{currentstroke}{rgb}{0.000000,0.000000,0.000000}%
\pgfsetstrokecolor{currentstroke}%
\pgfsetdash{}{0pt}%
\pgfsys@defobject{currentmarker}{\pgfqpoint{0.000000in}{-0.048611in}}{\pgfqpoint{0.000000in}{0.000000in}}{%
\pgfpathmoveto{\pgfqpoint{0.000000in}{0.000000in}}%
\pgfpathlineto{\pgfqpoint{0.000000in}{-0.048611in}}%
\pgfusepath{stroke,fill}%
}%
\begin{pgfscope}%
\pgfsys@transformshift{4.360186in}{2.240123in}%
\pgfsys@useobject{currentmarker}{}%
\end{pgfscope}%
\end{pgfscope}%
\begin{pgfscope}%
\definecolor{textcolor}{rgb}{0.000000,0.000000,0.000000}%
\pgfsetstrokecolor{textcolor}%
\pgfsetfillcolor{textcolor}%
\pgftext[x=4.360186in,y=2.142901in,,top]{\color{textcolor}\rmfamily\fontsize{10.000000}{12.000000}\selectfont \(\displaystyle {200000}\)}%
\end{pgfscope}%
\begin{pgfscope}%
\pgfpathrectangle{\pgfqpoint{3.833026in}{2.240123in}}{\pgfqpoint{2.108640in}{0.910803in}}%
\pgfusepath{clip}%
\pgfsetbuttcap%
\pgfsetroundjoin%
\pgfsetlinewidth{0.501875pt}%
\definecolor{currentstroke}{rgb}{0.690196,0.690196,0.690196}%
\pgfsetstrokecolor{currentstroke}%
\pgfsetstrokeopacity{0.700000}%
\pgfsetdash{{1.850000pt}{0.800000pt}}{0.000000pt}%
\pgfpathmoveto{\pgfqpoint{4.887346in}{2.240123in}}%
\pgfpathlineto{\pgfqpoint{4.887346in}{3.150926in}}%
\pgfusepath{stroke}%
\end{pgfscope}%
\begin{pgfscope}%
\pgfsetbuttcap%
\pgfsetroundjoin%
\definecolor{currentfill}{rgb}{0.000000,0.000000,0.000000}%
\pgfsetfillcolor{currentfill}%
\pgfsetlinewidth{0.803000pt}%
\definecolor{currentstroke}{rgb}{0.000000,0.000000,0.000000}%
\pgfsetstrokecolor{currentstroke}%
\pgfsetdash{}{0pt}%
\pgfsys@defobject{currentmarker}{\pgfqpoint{0.000000in}{-0.048611in}}{\pgfqpoint{0.000000in}{0.000000in}}{%
\pgfpathmoveto{\pgfqpoint{0.000000in}{0.000000in}}%
\pgfpathlineto{\pgfqpoint{0.000000in}{-0.048611in}}%
\pgfusepath{stroke,fill}%
}%
\begin{pgfscope}%
\pgfsys@transformshift{4.887346in}{2.240123in}%
\pgfsys@useobject{currentmarker}{}%
\end{pgfscope}%
\end{pgfscope}%
\begin{pgfscope}%
\definecolor{textcolor}{rgb}{0.000000,0.000000,0.000000}%
\pgfsetstrokecolor{textcolor}%
\pgfsetfillcolor{textcolor}%
\pgftext[x=4.887346in,y=2.142901in,,top]{\color{textcolor}\rmfamily\fontsize{10.000000}{12.000000}\selectfont \(\displaystyle {300000}\)}%
\end{pgfscope}%
\begin{pgfscope}%
\pgfpathrectangle{\pgfqpoint{3.833026in}{2.240123in}}{\pgfqpoint{2.108640in}{0.910803in}}%
\pgfusepath{clip}%
\pgfsetbuttcap%
\pgfsetroundjoin%
\pgfsetlinewidth{0.501875pt}%
\definecolor{currentstroke}{rgb}{0.690196,0.690196,0.690196}%
\pgfsetstrokecolor{currentstroke}%
\pgfsetstrokeopacity{0.700000}%
\pgfsetdash{{1.850000pt}{0.800000pt}}{0.000000pt}%
\pgfpathmoveto{\pgfqpoint{5.414506in}{2.240123in}}%
\pgfpathlineto{\pgfqpoint{5.414506in}{3.150926in}}%
\pgfusepath{stroke}%
\end{pgfscope}%
\begin{pgfscope}%
\pgfsetbuttcap%
\pgfsetroundjoin%
\definecolor{currentfill}{rgb}{0.000000,0.000000,0.000000}%
\pgfsetfillcolor{currentfill}%
\pgfsetlinewidth{0.803000pt}%
\definecolor{currentstroke}{rgb}{0.000000,0.000000,0.000000}%
\pgfsetstrokecolor{currentstroke}%
\pgfsetdash{}{0pt}%
\pgfsys@defobject{currentmarker}{\pgfqpoint{0.000000in}{-0.048611in}}{\pgfqpoint{0.000000in}{0.000000in}}{%
\pgfpathmoveto{\pgfqpoint{0.000000in}{0.000000in}}%
\pgfpathlineto{\pgfqpoint{0.000000in}{-0.048611in}}%
\pgfusepath{stroke,fill}%
}%
\begin{pgfscope}%
\pgfsys@transformshift{5.414506in}{2.240123in}%
\pgfsys@useobject{currentmarker}{}%
\end{pgfscope}%
\end{pgfscope}%
\begin{pgfscope}%
\definecolor{textcolor}{rgb}{0.000000,0.000000,0.000000}%
\pgfsetstrokecolor{textcolor}%
\pgfsetfillcolor{textcolor}%
\pgftext[x=5.414506in,y=2.142901in,,top]{\color{textcolor}\rmfamily\fontsize{10.000000}{12.000000}\selectfont \(\displaystyle {400000}\)}%
\end{pgfscope}%
\begin{pgfscope}%
\pgfpathrectangle{\pgfqpoint{3.833026in}{2.240123in}}{\pgfqpoint{2.108640in}{0.910803in}}%
\pgfusepath{clip}%
\pgfsetbuttcap%
\pgfsetroundjoin%
\pgfsetlinewidth{0.501875pt}%
\definecolor{currentstroke}{rgb}{0.690196,0.690196,0.690196}%
\pgfsetstrokecolor{currentstroke}%
\pgfsetstrokeopacity{0.700000}%
\pgfsetdash{{1.850000pt}{0.800000pt}}{0.000000pt}%
\pgfpathmoveto{\pgfqpoint{5.941666in}{2.240123in}}%
\pgfpathlineto{\pgfqpoint{5.941666in}{3.150926in}}%
\pgfusepath{stroke}%
\end{pgfscope}%
\begin{pgfscope}%
\pgfsetbuttcap%
\pgfsetroundjoin%
\definecolor{currentfill}{rgb}{0.000000,0.000000,0.000000}%
\pgfsetfillcolor{currentfill}%
\pgfsetlinewidth{0.803000pt}%
\definecolor{currentstroke}{rgb}{0.000000,0.000000,0.000000}%
\pgfsetstrokecolor{currentstroke}%
\pgfsetdash{}{0pt}%
\pgfsys@defobject{currentmarker}{\pgfqpoint{0.000000in}{-0.048611in}}{\pgfqpoint{0.000000in}{0.000000in}}{%
\pgfpathmoveto{\pgfqpoint{0.000000in}{0.000000in}}%
\pgfpathlineto{\pgfqpoint{0.000000in}{-0.048611in}}%
\pgfusepath{stroke,fill}%
}%
\begin{pgfscope}%
\pgfsys@transformshift{5.941666in}{2.240123in}%
\pgfsys@useobject{currentmarker}{}%
\end{pgfscope}%
\end{pgfscope}%
\begin{pgfscope}%
\definecolor{textcolor}{rgb}{0.000000,0.000000,0.000000}%
\pgfsetstrokecolor{textcolor}%
\pgfsetfillcolor{textcolor}%
\pgftext[x=5.941666in,y=2.142901in,,top]{\color{textcolor}\rmfamily\fontsize{10.000000}{12.000000}\selectfont \(\displaystyle {500000}\)}%
\end{pgfscope}%
\begin{pgfscope}%
\definecolor{textcolor}{rgb}{0.000000,0.000000,0.000000}%
\pgfsetstrokecolor{textcolor}%
\pgfsetfillcolor{textcolor}%
\pgftext[x=4.887346in,y=1.963889in,,top]{\color{textcolor}\rmfamily\fontsize{10.000000}{12.000000}\selectfont Frequency (Hz)}%
\end{pgfscope}%
\begin{pgfscope}%
\pgfpathrectangle{\pgfqpoint{3.833026in}{2.240123in}}{\pgfqpoint{2.108640in}{0.910803in}}%
\pgfusepath{clip}%
\pgfsetbuttcap%
\pgfsetroundjoin%
\pgfsetlinewidth{0.501875pt}%
\definecolor{currentstroke}{rgb}{0.690196,0.690196,0.690196}%
\pgfsetstrokecolor{currentstroke}%
\pgfsetstrokeopacity{0.700000}%
\pgfsetdash{{1.850000pt}{0.800000pt}}{0.000000pt}%
\pgfpathmoveto{\pgfqpoint{3.833026in}{2.407901in}}%
\pgfpathlineto{\pgfqpoint{5.941666in}{2.407901in}}%
\pgfusepath{stroke}%
\end{pgfscope}%
\begin{pgfscope}%
\pgfsetbuttcap%
\pgfsetroundjoin%
\definecolor{currentfill}{rgb}{0.000000,0.000000,0.000000}%
\pgfsetfillcolor{currentfill}%
\pgfsetlinewidth{0.803000pt}%
\definecolor{currentstroke}{rgb}{0.000000,0.000000,0.000000}%
\pgfsetstrokecolor{currentstroke}%
\pgfsetdash{}{0pt}%
\pgfsys@defobject{currentmarker}{\pgfqpoint{-0.048611in}{0.000000in}}{\pgfqpoint{-0.000000in}{0.000000in}}{%
\pgfpathmoveto{\pgfqpoint{-0.000000in}{0.000000in}}%
\pgfpathlineto{\pgfqpoint{-0.048611in}{0.000000in}}%
\pgfusepath{stroke,fill}%
}%
\begin{pgfscope}%
\pgfsys@transformshift{3.833026in}{2.407901in}%
\pgfsys@useobject{currentmarker}{}%
\end{pgfscope}%
\end{pgfscope}%
\begin{pgfscope}%
\definecolor{textcolor}{rgb}{0.000000,0.000000,0.000000}%
\pgfsetstrokecolor{textcolor}%
\pgfsetfillcolor{textcolor}%
\pgftext[x=3.419444in, y=2.359676in, left, base]{\color{textcolor}\rmfamily\fontsize{10.000000}{12.000000}\selectfont \(\displaystyle {\ensuremath{-}100}\)}%
\end{pgfscope}%
\begin{pgfscope}%
\pgfpathrectangle{\pgfqpoint{3.833026in}{2.240123in}}{\pgfqpoint{2.108640in}{0.910803in}}%
\pgfusepath{clip}%
\pgfsetbuttcap%
\pgfsetroundjoin%
\pgfsetlinewidth{0.501875pt}%
\definecolor{currentstroke}{rgb}{0.690196,0.690196,0.690196}%
\pgfsetstrokecolor{currentstroke}%
\pgfsetstrokeopacity{0.700000}%
\pgfsetdash{{1.850000pt}{0.800000pt}}{0.000000pt}%
\pgfpathmoveto{\pgfqpoint{3.833026in}{2.895141in}}%
\pgfpathlineto{\pgfqpoint{5.941666in}{2.895141in}}%
\pgfusepath{stroke}%
\end{pgfscope}%
\begin{pgfscope}%
\pgfsetbuttcap%
\pgfsetroundjoin%
\definecolor{currentfill}{rgb}{0.000000,0.000000,0.000000}%
\pgfsetfillcolor{currentfill}%
\pgfsetlinewidth{0.803000pt}%
\definecolor{currentstroke}{rgb}{0.000000,0.000000,0.000000}%
\pgfsetstrokecolor{currentstroke}%
\pgfsetdash{}{0pt}%
\pgfsys@defobject{currentmarker}{\pgfqpoint{-0.048611in}{0.000000in}}{\pgfqpoint{-0.000000in}{0.000000in}}{%
\pgfpathmoveto{\pgfqpoint{-0.000000in}{0.000000in}}%
\pgfpathlineto{\pgfqpoint{-0.048611in}{0.000000in}}%
\pgfusepath{stroke,fill}%
}%
\begin{pgfscope}%
\pgfsys@transformshift{3.833026in}{2.895141in}%
\pgfsys@useobject{currentmarker}{}%
\end{pgfscope}%
\end{pgfscope}%
\begin{pgfscope}%
\definecolor{textcolor}{rgb}{0.000000,0.000000,0.000000}%
\pgfsetstrokecolor{textcolor}%
\pgfsetfillcolor{textcolor}%
\pgftext[x=3.488889in, y=2.846915in, left, base]{\color{textcolor}\rmfamily\fontsize{10.000000}{12.000000}\selectfont \(\displaystyle {\ensuremath{-}50}\)}%
\end{pgfscope}%
\begin{pgfscope}%
\definecolor{textcolor}{rgb}{0.000000,0.000000,0.000000}%
\pgfsetstrokecolor{textcolor}%
\pgfsetfillcolor{textcolor}%
\pgftext[x=3.363889in,y=2.695525in,,bottom,rotate=90.000000]{\color{textcolor}\rmfamily\fontsize{10.000000}{12.000000}\selectfont Amplitude (dB)}%
\end{pgfscope}%
\begin{pgfscope}%
\pgfpathrectangle{\pgfqpoint{3.833026in}{2.240123in}}{\pgfqpoint{2.108640in}{0.910803in}}%
\pgfusepath{clip}%
\pgfsetrectcap%
\pgfsetroundjoin%
\pgfsetlinewidth{1.003750pt}%
\definecolor{currentstroke}{rgb}{0.000000,0.000000,0.000000}%
\pgfsetstrokecolor{currentstroke}%
\pgfsetdash{}{0pt}%
\pgfpathmoveto{\pgfqpoint{3.833026in}{2.572649in}}%
\pgfpathlineto{\pgfqpoint{3.837770in}{2.597011in}}%
\pgfpathlineto{\pgfqpoint{3.842515in}{2.601883in}}%
\pgfpathlineto{\pgfqpoint{3.847259in}{2.538542in}}%
\pgfpathlineto{\pgfqpoint{3.852003in}{2.543415in}}%
\pgfpathlineto{\pgfqpoint{3.856748in}{2.577521in}}%
\pgfpathlineto{\pgfqpoint{3.861492in}{2.577521in}}%
\pgfpathlineto{\pgfqpoint{3.866237in}{2.606756in}}%
\pgfpathlineto{\pgfqpoint{3.870981in}{2.538542in}}%
\pgfpathlineto{\pgfqpoint{3.875726in}{2.582089in}}%
\pgfpathlineto{\pgfqpoint{3.880470in}{2.557727in}}%
\pgfpathlineto{\pgfqpoint{3.885215in}{2.586962in}}%
\pgfpathlineto{\pgfqpoint{3.889959in}{2.572345in}}%
\pgfpathlineto{\pgfqpoint{3.894703in}{2.538238in}}%
\pgfpathlineto{\pgfqpoint{3.899448in}{2.601579in}}%
\pgfpathlineto{\pgfqpoint{3.904192in}{2.611324in}}%
\pgfpathlineto{\pgfqpoint{3.908937in}{2.616196in}}%
\pgfpathlineto{\pgfqpoint{3.918426in}{2.586657in}}%
\pgfpathlineto{\pgfqpoint{3.923170in}{2.591530in}}%
\pgfpathlineto{\pgfqpoint{3.927915in}{2.645126in}}%
\pgfpathlineto{\pgfqpoint{3.932659in}{2.615892in}}%
\pgfpathlineto{\pgfqpoint{3.937403in}{2.620764in}}%
\pgfpathlineto{\pgfqpoint{3.942148in}{3.005683in}}%
\pgfpathlineto{\pgfqpoint{3.946892in}{3.103131in}}%
\pgfpathlineto{\pgfqpoint{3.951637in}{3.073897in}}%
\pgfpathlineto{\pgfqpoint{3.956381in}{2.873824in}}%
\pgfpathlineto{\pgfqpoint{3.961126in}{2.605842in}}%
\pgfpathlineto{\pgfqpoint{3.965870in}{2.649694in}}%
\pgfpathlineto{\pgfqpoint{3.970614in}{2.630204in}}%
\pgfpathlineto{\pgfqpoint{3.975359in}{2.537629in}}%
\pgfpathlineto{\pgfqpoint{3.984848in}{2.581480in}}%
\pgfpathlineto{\pgfqpoint{3.989592in}{2.600970in}}%
\pgfpathlineto{\pgfqpoint{3.994337in}{2.590921in}}%
\pgfpathlineto{\pgfqpoint{3.999081in}{2.586048in}}%
\pgfpathlineto{\pgfqpoint{4.003826in}{2.571431in}}%
\pgfpathlineto{\pgfqpoint{4.008570in}{2.576303in}}%
\pgfpathlineto{\pgfqpoint{4.013314in}{2.561686in}}%
\pgfpathlineto{\pgfqpoint{4.018059in}{2.590921in}}%
\pgfpathlineto{\pgfqpoint{4.022803in}{2.586048in}}%
\pgfpathlineto{\pgfqpoint{4.027548in}{2.595488in}}%
\pgfpathlineto{\pgfqpoint{4.032292in}{2.546764in}}%
\pgfpathlineto{\pgfqpoint{4.037037in}{2.541892in}}%
\pgfpathlineto{\pgfqpoint{4.041781in}{2.634468in}}%
\pgfpathlineto{\pgfqpoint{4.046526in}{2.624723in}}%
\pgfpathlineto{\pgfqpoint{4.051270in}{2.610106in}}%
\pgfpathlineto{\pgfqpoint{4.056014in}{2.585744in}}%
\pgfpathlineto{\pgfqpoint{4.060759in}{2.580871in}}%
\pgfpathlineto{\pgfqpoint{4.065503in}{2.556205in}}%
\pgfpathlineto{\pgfqpoint{4.070248in}{2.609801in}}%
\pgfpathlineto{\pgfqpoint{4.074992in}{2.575694in}}%
\pgfpathlineto{\pgfqpoint{4.079737in}{3.009337in}}%
\pgfpathlineto{\pgfqpoint{4.084481in}{3.101913in}}%
\pgfpathlineto{\pgfqpoint{4.089225in}{3.077551in}}%
\pgfpathlineto{\pgfqpoint{4.093970in}{2.863166in}}%
\pgfpathlineto{\pgfqpoint{4.098714in}{2.575694in}}%
\pgfpathlineto{\pgfqpoint{4.103459in}{2.653653in}}%
\pgfpathlineto{\pgfqpoint{4.108203in}{2.575390in}}%
\pgfpathlineto{\pgfqpoint{4.112948in}{2.614369in}}%
\pgfpathlineto{\pgfqpoint{4.117692in}{2.526666in}}%
\pgfpathlineto{\pgfqpoint{4.127181in}{2.614369in}}%
\pgfpathlineto{\pgfqpoint{4.131925in}{2.585135in}}%
\pgfpathlineto{\pgfqpoint{4.136670in}{2.570517in}}%
\pgfpathlineto{\pgfqpoint{4.141414in}{2.565645in}}%
\pgfpathlineto{\pgfqpoint{4.146159in}{2.555596in}}%
\pgfpathlineto{\pgfqpoint{4.150903in}{2.594575in}}%
\pgfpathlineto{\pgfqpoint{4.155648in}{2.536106in}}%
\pgfpathlineto{\pgfqpoint{4.160392in}{2.501999in}}%
\pgfpathlineto{\pgfqpoint{4.165137in}{2.604320in}}%
\pgfpathlineto{\pgfqpoint{4.169881in}{2.599447in}}%
\pgfpathlineto{\pgfqpoint{4.174625in}{2.589702in}}%
\pgfpathlineto{\pgfqpoint{4.179370in}{2.599143in}}%
\pgfpathlineto{\pgfqpoint{4.184114in}{2.550419in}}%
\pgfpathlineto{\pgfqpoint{4.193603in}{2.604015in}}%
\pgfpathlineto{\pgfqpoint{4.198348in}{2.589398in}}%
\pgfpathlineto{\pgfqpoint{4.207837in}{2.618632in}}%
\pgfpathlineto{\pgfqpoint{4.212581in}{2.540674in}}%
\pgfpathlineto{\pgfqpoint{4.217325in}{2.564731in}}%
\pgfpathlineto{\pgfqpoint{4.222070in}{2.642690in}}%
\pgfpathlineto{\pgfqpoint{4.226814in}{2.613455in}}%
\pgfpathlineto{\pgfqpoint{4.231559in}{2.598838in}}%
\pgfpathlineto{\pgfqpoint{4.236303in}{2.603711in}}%
\pgfpathlineto{\pgfqpoint{4.241048in}{2.564731in}}%
\pgfpathlineto{\pgfqpoint{4.245792in}{2.593966in}}%
\pgfpathlineto{\pgfqpoint{4.250536in}{2.579349in}}%
\pgfpathlineto{\pgfqpoint{4.255281in}{2.554682in}}%
\pgfpathlineto{\pgfqpoint{4.260025in}{2.598534in}}%
\pgfpathlineto{\pgfqpoint{4.264770in}{2.598534in}}%
\pgfpathlineto{\pgfqpoint{4.269514in}{2.579044in}}%
\pgfpathlineto{\pgfqpoint{4.279003in}{2.598534in}}%
\pgfpathlineto{\pgfqpoint{4.283748in}{2.579044in}}%
\pgfpathlineto{\pgfqpoint{4.288492in}{2.588789in}}%
\pgfpathlineto{\pgfqpoint{4.293236in}{2.612846in}}%
\pgfpathlineto{\pgfqpoint{4.297981in}{2.593357in}}%
\pgfpathlineto{\pgfqpoint{4.302725in}{2.578740in}}%
\pgfpathlineto{\pgfqpoint{4.307470in}{2.554378in}}%
\pgfpathlineto{\pgfqpoint{4.312214in}{2.593357in}}%
\pgfpathlineto{\pgfqpoint{4.316959in}{2.671315in}}%
\pgfpathlineto{\pgfqpoint{4.321703in}{2.598229in}}%
\pgfpathlineto{\pgfqpoint{4.326448in}{2.646953in}}%
\pgfpathlineto{\pgfqpoint{4.331192in}{2.554378in}}%
\pgfpathlineto{\pgfqpoint{4.335936in}{2.568690in}}%
\pgfpathlineto{\pgfqpoint{4.340681in}{2.617414in}}%
\pgfpathlineto{\pgfqpoint{4.345425in}{2.583307in}}%
\pgfpathlineto{\pgfqpoint{4.350170in}{2.578435in}}%
\pgfpathlineto{\pgfqpoint{4.354914in}{3.012078in}}%
\pgfpathlineto{\pgfqpoint{4.359659in}{3.109526in}}%
\pgfpathlineto{\pgfqpoint{4.364403in}{3.080292in}}%
\pgfpathlineto{\pgfqpoint{4.373892in}{2.602492in}}%
\pgfpathlineto{\pgfqpoint{4.378636in}{2.607365in}}%
\pgfpathlineto{\pgfqpoint{4.383381in}{2.587875in}}%
\pgfpathlineto{\pgfqpoint{4.388125in}{2.587875in}}%
\pgfpathlineto{\pgfqpoint{4.392870in}{2.568386in}}%
\pgfpathlineto{\pgfqpoint{4.397614in}{2.534279in}}%
\pgfpathlineto{\pgfqpoint{4.402359in}{2.578131in}}%
\pgfpathlineto{\pgfqpoint{4.407103in}{2.558336in}}%
\pgfpathlineto{\pgfqpoint{4.411847in}{2.568081in}}%
\pgfpathlineto{\pgfqpoint{4.421336in}{2.626550in}}%
\pgfpathlineto{\pgfqpoint{4.426081in}{2.572954in}}%
\pgfpathlineto{\pgfqpoint{4.430825in}{2.577826in}}%
\pgfpathlineto{\pgfqpoint{4.435570in}{2.568081in}}%
\pgfpathlineto{\pgfqpoint{4.440314in}{2.602188in}}%
\pgfpathlineto{\pgfqpoint{4.445059in}{2.567777in}}%
\pgfpathlineto{\pgfqpoint{4.449803in}{2.665225in}}%
\pgfpathlineto{\pgfqpoint{4.454547in}{2.674969in}}%
\pgfpathlineto{\pgfqpoint{4.459292in}{2.577521in}}%
\pgfpathlineto{\pgfqpoint{4.464036in}{2.601883in}}%
\pgfpathlineto{\pgfqpoint{4.468781in}{2.592139in}}%
\pgfpathlineto{\pgfqpoint{4.473525in}{2.567777in}}%
\pgfpathlineto{\pgfqpoint{4.478270in}{2.606756in}}%
\pgfpathlineto{\pgfqpoint{4.483014in}{2.567472in}}%
\pgfpathlineto{\pgfqpoint{4.487759in}{2.625941in}}%
\pgfpathlineto{\pgfqpoint{4.492503in}{2.621069in}}%
\pgfpathlineto{\pgfqpoint{4.497247in}{2.645430in}}%
\pgfpathlineto{\pgfqpoint{4.501992in}{2.621069in}}%
\pgfpathlineto{\pgfqpoint{4.511481in}{2.547983in}}%
\pgfpathlineto{\pgfqpoint{4.516225in}{2.586962in}}%
\pgfpathlineto{\pgfqpoint{4.520970in}{2.601274in}}%
\pgfpathlineto{\pgfqpoint{4.525714in}{2.537933in}}%
\pgfpathlineto{\pgfqpoint{4.530458in}{2.586657in}}%
\pgfpathlineto{\pgfqpoint{4.539947in}{2.586657in}}%
\pgfpathlineto{\pgfqpoint{4.544692in}{2.562295in}}%
\pgfpathlineto{\pgfqpoint{4.549436in}{2.591530in}}%
\pgfpathlineto{\pgfqpoint{4.554181in}{2.586657in}}%
\pgfpathlineto{\pgfqpoint{4.558925in}{2.561991in}}%
\pgfpathlineto{\pgfqpoint{4.563670in}{2.523012in}}%
\pgfpathlineto{\pgfqpoint{4.568414in}{2.591225in}}%
\pgfpathlineto{\pgfqpoint{4.573158in}{2.561991in}}%
\pgfpathlineto{\pgfqpoint{4.577903in}{2.605842in}}%
\pgfpathlineto{\pgfqpoint{4.582647in}{2.625332in}}%
\pgfpathlineto{\pgfqpoint{4.587392in}{2.708163in}}%
\pgfpathlineto{\pgfqpoint{4.592136in}{2.703290in}}%
\pgfpathlineto{\pgfqpoint{4.596881in}{2.625027in}}%
\pgfpathlineto{\pgfqpoint{4.601625in}{2.605538in}}%
\pgfpathlineto{\pgfqpoint{4.606370in}{2.537324in}}%
\pgfpathlineto{\pgfqpoint{4.611114in}{2.586048in}}%
\pgfpathlineto{\pgfqpoint{4.615858in}{2.590921in}}%
\pgfpathlineto{\pgfqpoint{4.620603in}{2.581176in}}%
\pgfpathlineto{\pgfqpoint{4.625347in}{2.610410in}}%
\pgfpathlineto{\pgfqpoint{4.630092in}{3.019691in}}%
\pgfpathlineto{\pgfqpoint{4.634836in}{3.107090in}}%
\pgfpathlineto{\pgfqpoint{4.639581in}{3.077856in}}%
\pgfpathlineto{\pgfqpoint{4.649070in}{2.639340in}}%
\pgfpathlineto{\pgfqpoint{4.653814in}{2.614978in}}%
\pgfpathlineto{\pgfqpoint{4.658558in}{2.614978in}}%
\pgfpathlineto{\pgfqpoint{4.663303in}{2.624723in}}%
\pgfpathlineto{\pgfqpoint{4.668047in}{2.619850in}}%
\pgfpathlineto{\pgfqpoint{4.672792in}{2.585439in}}%
\pgfpathlineto{\pgfqpoint{4.677536in}{2.565950in}}%
\pgfpathlineto{\pgfqpoint{4.682281in}{2.556205in}}%
\pgfpathlineto{\pgfqpoint{4.687025in}{2.565950in}}%
\pgfpathlineto{\pgfqpoint{4.691769in}{2.522098in}}%
\pgfpathlineto{\pgfqpoint{4.696514in}{2.585439in}}%
\pgfpathlineto{\pgfqpoint{4.701258in}{2.551332in}}%
\pgfpathlineto{\pgfqpoint{4.706003in}{2.570822in}}%
\pgfpathlineto{\pgfqpoint{4.710747in}{2.604624in}}%
\pgfpathlineto{\pgfqpoint{4.715492in}{2.546155in}}%
\pgfpathlineto{\pgfqpoint{4.720236in}{2.580262in}}%
\pgfpathlineto{\pgfqpoint{4.724981in}{2.653348in}}%
\pgfpathlineto{\pgfqpoint{4.729725in}{2.692327in}}%
\pgfpathlineto{\pgfqpoint{4.734469in}{2.619241in}}%
\pgfpathlineto{\pgfqpoint{4.739214in}{2.565645in}}%
\pgfpathlineto{\pgfqpoint{4.743958in}{2.614369in}}%
\pgfpathlineto{\pgfqpoint{4.748703in}{2.584830in}}%
\pgfpathlineto{\pgfqpoint{4.753447in}{2.614064in}}%
\pgfpathlineto{\pgfqpoint{4.758192in}{2.599447in}}%
\pgfpathlineto{\pgfqpoint{4.762936in}{2.565340in}}%
\pgfpathlineto{\pgfqpoint{4.767681in}{3.028218in}}%
\pgfpathlineto{\pgfqpoint{4.772425in}{3.106176in}}%
\pgfpathlineto{\pgfqpoint{4.777169in}{3.086687in}}%
\pgfpathlineto{\pgfqpoint{4.786658in}{2.628377in}}%
\pgfpathlineto{\pgfqpoint{4.791403in}{2.613760in}}%
\pgfpathlineto{\pgfqpoint{4.796147in}{2.628377in}}%
\pgfpathlineto{\pgfqpoint{4.800892in}{2.589398in}}%
\pgfpathlineto{\pgfqpoint{4.805636in}{2.569908in}}%
\pgfpathlineto{\pgfqpoint{4.815125in}{2.579653in}}%
\pgfpathlineto{\pgfqpoint{4.819869in}{2.594270in}}%
\pgfpathlineto{\pgfqpoint{4.824614in}{2.579349in}}%
\pgfpathlineto{\pgfqpoint{4.829358in}{2.603711in}}%
\pgfpathlineto{\pgfqpoint{4.834103in}{2.593966in}}%
\pgfpathlineto{\pgfqpoint{4.838847in}{2.554987in}}%
\pgfpathlineto{\pgfqpoint{4.843592in}{2.559859in}}%
\pgfpathlineto{\pgfqpoint{4.848336in}{2.579349in}}%
\pgfpathlineto{\pgfqpoint{4.853080in}{2.530625in}}%
\pgfpathlineto{\pgfqpoint{4.857825in}{2.593966in}}%
\pgfpathlineto{\pgfqpoint{4.862569in}{2.569299in}}%
\pgfpathlineto{\pgfqpoint{4.867314in}{2.574172in}}%
\pgfpathlineto{\pgfqpoint{4.872058in}{2.613151in}}%
\pgfpathlineto{\pgfqpoint{4.876803in}{2.583916in}}%
\pgfpathlineto{\pgfqpoint{4.881547in}{2.603406in}}%
\pgfpathlineto{\pgfqpoint{4.886292in}{2.603406in}}%
\pgfpathlineto{\pgfqpoint{4.891036in}{2.569299in}}%
\pgfpathlineto{\pgfqpoint{4.895780in}{2.593661in}}%
\pgfpathlineto{\pgfqpoint{4.900525in}{2.637208in}}%
\pgfpathlineto{\pgfqpoint{4.905269in}{2.988021in}}%
\pgfpathlineto{\pgfqpoint{4.910014in}{3.065979in}}%
\pgfpathlineto{\pgfqpoint{4.914758in}{3.046490in}}%
\pgfpathlineto{\pgfqpoint{4.924247in}{2.607974in}}%
\pgfpathlineto{\pgfqpoint{4.928992in}{2.627464in}}%
\pgfpathlineto{\pgfqpoint{4.933736in}{2.568995in}}%
\pgfpathlineto{\pgfqpoint{4.938480in}{2.588180in}}%
\pgfpathlineto{\pgfqpoint{4.943225in}{2.568690in}}%
\pgfpathlineto{\pgfqpoint{4.947969in}{2.602797in}}%
\pgfpathlineto{\pgfqpoint{4.952714in}{2.588180in}}%
\pgfpathlineto{\pgfqpoint{4.957458in}{2.563818in}}%
\pgfpathlineto{\pgfqpoint{4.966947in}{2.632031in}}%
\pgfpathlineto{\pgfqpoint{4.971691in}{2.583307in}}%
\pgfpathlineto{\pgfqpoint{4.976436in}{2.621982in}}%
\pgfpathlineto{\pgfqpoint{4.985925in}{2.568386in}}%
\pgfpathlineto{\pgfqpoint{4.990669in}{2.558641in}}%
\pgfpathlineto{\pgfqpoint{4.995414in}{2.568386in}}%
\pgfpathlineto{\pgfqpoint{5.000158in}{2.641472in}}%
\pgfpathlineto{\pgfqpoint{5.004903in}{2.695068in}}%
\pgfpathlineto{\pgfqpoint{5.009647in}{2.583003in}}%
\pgfpathlineto{\pgfqpoint{5.014391in}{2.558336in}}%
\pgfpathlineto{\pgfqpoint{5.019136in}{2.563209in}}%
\pgfpathlineto{\pgfqpoint{5.023880in}{2.558336in}}%
\pgfpathlineto{\pgfqpoint{5.028625in}{2.568081in}}%
\pgfpathlineto{\pgfqpoint{5.033369in}{2.568081in}}%
\pgfpathlineto{\pgfqpoint{5.038114in}{2.621678in}}%
\pgfpathlineto{\pgfqpoint{5.042858in}{2.914021in}}%
\pgfpathlineto{\pgfqpoint{5.047603in}{3.006597in}}%
\pgfpathlineto{\pgfqpoint{5.052347in}{2.977058in}}%
\pgfpathlineto{\pgfqpoint{5.057091in}{2.772417in}}%
\pgfpathlineto{\pgfqpoint{5.061836in}{2.616501in}}%
\pgfpathlineto{\pgfqpoint{5.066580in}{2.582394in}}%
\pgfpathlineto{\pgfqpoint{5.071325in}{2.562904in}}%
\pgfpathlineto{\pgfqpoint{5.076069in}{2.592139in}}%
\pgfpathlineto{\pgfqpoint{5.080814in}{2.582394in}}%
\pgfpathlineto{\pgfqpoint{5.085558in}{2.528798in}}%
\pgfpathlineto{\pgfqpoint{5.090303in}{2.586962in}}%
\pgfpathlineto{\pgfqpoint{5.095047in}{2.582089in}}%
\pgfpathlineto{\pgfqpoint{5.099791in}{2.621069in}}%
\pgfpathlineto{\pgfqpoint{5.104536in}{2.635686in}}%
\pgfpathlineto{\pgfqpoint{5.109280in}{2.601579in}}%
\pgfpathlineto{\pgfqpoint{5.114025in}{2.504131in}}%
\pgfpathlineto{\pgfqpoint{5.118769in}{2.572345in}}%
\pgfpathlineto{\pgfqpoint{5.123514in}{2.557727in}}%
\pgfpathlineto{\pgfqpoint{5.128258in}{2.601274in}}%
\pgfpathlineto{\pgfqpoint{5.133002in}{2.542806in}}%
\pgfpathlineto{\pgfqpoint{5.137747in}{2.630509in}}%
\pgfpathlineto{\pgfqpoint{5.142491in}{2.688978in}}%
\pgfpathlineto{\pgfqpoint{5.147236in}{2.572040in}}%
\pgfpathlineto{\pgfqpoint{5.151980in}{2.562295in}}%
\pgfpathlineto{\pgfqpoint{5.156725in}{2.576912in}}%
\pgfpathlineto{\pgfqpoint{5.161469in}{2.610715in}}%
\pgfpathlineto{\pgfqpoint{5.166214in}{2.571736in}}%
\pgfpathlineto{\pgfqpoint{5.170958in}{2.561991in}}%
\pgfpathlineto{\pgfqpoint{5.175702in}{2.557118in}}%
\pgfpathlineto{\pgfqpoint{5.180447in}{2.829972in}}%
\pgfpathlineto{\pgfqpoint{5.185191in}{2.922548in}}%
\pgfpathlineto{\pgfqpoint{5.189936in}{2.888441in}}%
\pgfpathlineto{\pgfqpoint{5.194680in}{2.674056in}}%
\pgfpathlineto{\pgfqpoint{5.199425in}{2.532452in}}%
\pgfpathlineto{\pgfqpoint{5.204169in}{2.566559in}}%
\pgfpathlineto{\pgfqpoint{5.208914in}{2.581176in}}%
\pgfpathlineto{\pgfqpoint{5.213658in}{2.547069in}}%
\pgfpathlineto{\pgfqpoint{5.218402in}{2.527579in}}%
\pgfpathlineto{\pgfqpoint{5.223147in}{2.586048in}}%
\pgfpathlineto{\pgfqpoint{5.227891in}{2.566559in}}%
\pgfpathlineto{\pgfqpoint{5.232636in}{2.551941in}}%
\pgfpathlineto{\pgfqpoint{5.237380in}{2.585744in}}%
\pgfpathlineto{\pgfqpoint{5.242125in}{2.551637in}}%
\pgfpathlineto{\pgfqpoint{5.246869in}{2.575999in}}%
\pgfpathlineto{\pgfqpoint{5.251613in}{2.532147in}}%
\pgfpathlineto{\pgfqpoint{5.256358in}{2.600361in}}%
\pgfpathlineto{\pgfqpoint{5.261102in}{2.546764in}}%
\pgfpathlineto{\pgfqpoint{5.265847in}{2.590616in}}%
\pgfpathlineto{\pgfqpoint{5.270591in}{2.541892in}}%
\pgfpathlineto{\pgfqpoint{5.275336in}{2.629291in}}%
\pgfpathlineto{\pgfqpoint{5.280080in}{2.624418in}}%
\pgfpathlineto{\pgfqpoint{5.284825in}{2.570822in}}%
\pgfpathlineto{\pgfqpoint{5.289569in}{2.575694in}}%
\pgfpathlineto{\pgfqpoint{5.294313in}{2.522098in}}%
\pgfpathlineto{\pgfqpoint{5.299058in}{2.580567in}}%
\pgfpathlineto{\pgfqpoint{5.303802in}{2.556205in}}%
\pgfpathlineto{\pgfqpoint{5.308547in}{2.580567in}}%
\pgfpathlineto{\pgfqpoint{5.313291in}{2.561077in}}%
\pgfpathlineto{\pgfqpoint{5.318036in}{2.726434in}}%
\pgfpathlineto{\pgfqpoint{5.322780in}{2.809265in}}%
\pgfpathlineto{\pgfqpoint{5.327525in}{2.784903in}}%
\pgfpathlineto{\pgfqpoint{5.332269in}{2.604624in}}%
\pgfpathlineto{\pgfqpoint{5.337013in}{2.570517in}}%
\pgfpathlineto{\pgfqpoint{5.341758in}{2.516921in}}%
\pgfpathlineto{\pgfqpoint{5.346502in}{2.565645in}}%
\pgfpathlineto{\pgfqpoint{5.351247in}{2.526666in}}%
\pgfpathlineto{\pgfqpoint{5.355991in}{2.575085in}}%
\pgfpathlineto{\pgfqpoint{5.365480in}{2.545851in}}%
\pgfpathlineto{\pgfqpoint{5.370225in}{2.560468in}}%
\pgfpathlineto{\pgfqpoint{5.374969in}{2.589702in}}%
\pgfpathlineto{\pgfqpoint{5.379713in}{2.545851in}}%
\pgfpathlineto{\pgfqpoint{5.384458in}{2.545851in}}%
\pgfpathlineto{\pgfqpoint{5.389202in}{2.530929in}}%
\pgfpathlineto{\pgfqpoint{5.393947in}{2.565036in}}%
\pgfpathlineto{\pgfqpoint{5.398691in}{2.516312in}}%
\pgfpathlineto{\pgfqpoint{5.403436in}{2.599143in}}%
\pgfpathlineto{\pgfqpoint{5.408180in}{2.560164in}}%
\pgfpathlineto{\pgfqpoint{5.412924in}{2.633249in}}%
\pgfpathlineto{\pgfqpoint{5.417669in}{2.623505in}}%
\pgfpathlineto{\pgfqpoint{5.422413in}{2.555291in}}%
\pgfpathlineto{\pgfqpoint{5.427158in}{2.569908in}}%
\pgfpathlineto{\pgfqpoint{5.431902in}{2.569604in}}%
\pgfpathlineto{\pgfqpoint{5.436647in}{2.564731in}}%
\pgfpathlineto{\pgfqpoint{5.441391in}{2.516007in}}%
\pgfpathlineto{\pgfqpoint{5.446136in}{2.593966in}}%
\pgfpathlineto{\pgfqpoint{5.450880in}{2.516007in}}%
\pgfpathlineto{\pgfqpoint{5.460369in}{2.706031in}}%
\pgfpathlineto{\pgfqpoint{5.465113in}{2.671924in}}%
\pgfpathlineto{\pgfqpoint{5.469858in}{2.501390in}}%
\pgfpathlineto{\pgfqpoint{5.474602in}{2.554682in}}%
\pgfpathlineto{\pgfqpoint{5.479347in}{2.554682in}}%
\pgfpathlineto{\pgfqpoint{5.488836in}{2.525448in}}%
\pgfpathlineto{\pgfqpoint{5.493580in}{2.564427in}}%
\pgfpathlineto{\pgfqpoint{5.498324in}{2.530320in}}%
\pgfpathlineto{\pgfqpoint{5.503069in}{2.520575in}}%
\pgfpathlineto{\pgfqpoint{5.507813in}{2.525448in}}%
\pgfpathlineto{\pgfqpoint{5.512558in}{2.525143in}}%
\pgfpathlineto{\pgfqpoint{5.517302in}{2.549505in}}%
\pgfpathlineto{\pgfqpoint{5.522047in}{2.544633in}}%
\pgfpathlineto{\pgfqpoint{5.526791in}{2.549505in}}%
\pgfpathlineto{\pgfqpoint{5.531535in}{2.578740in}}%
\pgfpathlineto{\pgfqpoint{5.536280in}{2.559250in}}%
\pgfpathlineto{\pgfqpoint{5.541024in}{2.534888in}}%
\pgfpathlineto{\pgfqpoint{5.545769in}{2.578435in}}%
\pgfpathlineto{\pgfqpoint{5.550513in}{2.573563in}}%
\pgfpathlineto{\pgfqpoint{5.555258in}{2.607669in}}%
\pgfpathlineto{\pgfqpoint{5.560002in}{2.529711in}}%
\pgfpathlineto{\pgfqpoint{5.564747in}{2.549201in}}%
\pgfpathlineto{\pgfqpoint{5.569491in}{2.544328in}}%
\pgfpathlineto{\pgfqpoint{5.574235in}{2.597925in}}%
\pgfpathlineto{\pgfqpoint{5.583724in}{2.534279in}}%
\pgfpathlineto{\pgfqpoint{5.588469in}{2.539151in}}%
\pgfpathlineto{\pgfqpoint{5.593213in}{2.514789in}}%
\pgfpathlineto{\pgfqpoint{5.597958in}{2.602492in}}%
\pgfpathlineto{\pgfqpoint{5.607447in}{2.509917in}}%
\pgfpathlineto{\pgfqpoint{5.612191in}{2.524534in}}%
\pgfpathlineto{\pgfqpoint{5.616935in}{2.568386in}}%
\pgfpathlineto{\pgfqpoint{5.626424in}{2.529102in}}%
\pgfpathlineto{\pgfqpoint{5.631169in}{2.499868in}}%
\pgfpathlineto{\pgfqpoint{5.635913in}{2.509612in}}%
\pgfpathlineto{\pgfqpoint{5.640658in}{2.572954in}}%
\pgfpathlineto{\pgfqpoint{5.645402in}{2.538847in}}%
\pgfpathlineto{\pgfqpoint{5.650147in}{2.533974in}}%
\pgfpathlineto{\pgfqpoint{5.654891in}{2.485250in}}%
\pgfpathlineto{\pgfqpoint{5.659635in}{2.553464in}}%
\pgfpathlineto{\pgfqpoint{5.664380in}{2.543415in}}%
\pgfpathlineto{\pgfqpoint{5.669124in}{2.543415in}}%
\pgfpathlineto{\pgfqpoint{5.673869in}{2.528798in}}%
\pgfpathlineto{\pgfqpoint{5.678613in}{2.562904in}}%
\pgfpathlineto{\pgfqpoint{5.683358in}{2.572649in}}%
\pgfpathlineto{\pgfqpoint{5.688102in}{2.572649in}}%
\pgfpathlineto{\pgfqpoint{5.692846in}{2.606756in}}%
\pgfpathlineto{\pgfqpoint{5.697591in}{2.552855in}}%
\pgfpathlineto{\pgfqpoint{5.702335in}{2.538238in}}%
\pgfpathlineto{\pgfqpoint{5.707080in}{2.509003in}}%
\pgfpathlineto{\pgfqpoint{5.711824in}{2.523621in}}%
\pgfpathlineto{\pgfqpoint{5.716569in}{2.528493in}}%
\pgfpathlineto{\pgfqpoint{5.721313in}{2.562600in}}%
\pgfpathlineto{\pgfqpoint{5.726058in}{2.513876in}}%
\pgfpathlineto{\pgfqpoint{5.730802in}{2.547983in}}%
\pgfpathlineto{\pgfqpoint{5.735546in}{2.513876in}}%
\pgfpathlineto{\pgfqpoint{5.740291in}{2.528188in}}%
\pgfpathlineto{\pgfqpoint{5.745035in}{2.547678in}}%
\pgfpathlineto{\pgfqpoint{5.749780in}{2.508699in}}%
\pgfpathlineto{\pgfqpoint{5.754524in}{2.479465in}}%
\pgfpathlineto{\pgfqpoint{5.759269in}{2.513571in}}%
\pgfpathlineto{\pgfqpoint{5.764013in}{2.513571in}}%
\pgfpathlineto{\pgfqpoint{5.768758in}{2.533061in}}%
\pgfpathlineto{\pgfqpoint{5.773502in}{2.528188in}}%
\pgfpathlineto{\pgfqpoint{5.778246in}{2.484032in}}%
\pgfpathlineto{\pgfqpoint{5.782991in}{2.561991in}}%
\pgfpathlineto{\pgfqpoint{5.787735in}{2.523012in}}%
\pgfpathlineto{\pgfqpoint{5.792480in}{2.498650in}}%
\pgfpathlineto{\pgfqpoint{5.797224in}{2.464543in}}%
\pgfpathlineto{\pgfqpoint{5.801969in}{2.479160in}}%
\pgfpathlineto{\pgfqpoint{5.806713in}{2.547374in}}%
\pgfpathlineto{\pgfqpoint{5.811458in}{2.537324in}}%
\pgfpathlineto{\pgfqpoint{5.816202in}{2.512962in}}%
\pgfpathlineto{\pgfqpoint{5.820946in}{2.503217in}}%
\pgfpathlineto{\pgfqpoint{5.825691in}{2.595793in}}%
\pgfpathlineto{\pgfqpoint{5.830435in}{2.600665in}}%
\pgfpathlineto{\pgfqpoint{5.835180in}{2.498345in}}%
\pgfpathlineto{\pgfqpoint{5.839924in}{2.556814in}}%
\pgfpathlineto{\pgfqpoint{5.844669in}{2.493473in}}%
\pgfpathlineto{\pgfqpoint{5.849413in}{2.508090in}}%
\pgfpathlineto{\pgfqpoint{5.854157in}{2.498041in}}%
\pgfpathlineto{\pgfqpoint{5.858902in}{2.507785in}}%
\pgfpathlineto{\pgfqpoint{5.863646in}{2.507785in}}%
\pgfpathlineto{\pgfqpoint{5.868391in}{2.478551in}}%
\pgfpathlineto{\pgfqpoint{5.873135in}{2.527275in}}%
\pgfpathlineto{\pgfqpoint{5.877880in}{2.512658in}}%
\pgfpathlineto{\pgfqpoint{5.882624in}{2.507785in}}%
\pgfpathlineto{\pgfqpoint{5.892113in}{2.507481in}}%
\pgfpathlineto{\pgfqpoint{5.896857in}{2.497736in}}%
\pgfpathlineto{\pgfqpoint{5.901602in}{2.590312in}}%
\pgfpathlineto{\pgfqpoint{5.906346in}{2.522098in}}%
\pgfpathlineto{\pgfqpoint{5.911091in}{2.541588in}}%
\pgfpathlineto{\pgfqpoint{5.915835in}{2.522098in}}%
\pgfpathlineto{\pgfqpoint{5.920580in}{2.497736in}}%
\pgfpathlineto{\pgfqpoint{5.925324in}{2.541283in}}%
\pgfpathlineto{\pgfqpoint{5.930069in}{2.497431in}}%
\pgfpathlineto{\pgfqpoint{5.939557in}{2.526666in}}%
\pgfpathlineto{\pgfqpoint{5.944302in}{2.526666in}}%
\pgfpathlineto{\pgfqpoint{5.944302in}{2.526666in}}%
\pgfusepath{stroke}%
\end{pgfscope}%
\begin{pgfscope}%
\pgfsetrectcap%
\pgfsetmiterjoin%
\pgfsetlinewidth{0.803000pt}%
\definecolor{currentstroke}{rgb}{0.000000,0.000000,0.000000}%
\pgfsetstrokecolor{currentstroke}%
\pgfsetdash{}{0pt}%
\pgfpathmoveto{\pgfqpoint{3.833026in}{2.240123in}}%
\pgfpathlineto{\pgfqpoint{3.833026in}{3.150926in}}%
\pgfusepath{stroke}%
\end{pgfscope}%
\begin{pgfscope}%
\pgfsetrectcap%
\pgfsetmiterjoin%
\pgfsetlinewidth{0.803000pt}%
\definecolor{currentstroke}{rgb}{0.000000,0.000000,0.000000}%
\pgfsetstrokecolor{currentstroke}%
\pgfsetdash{}{0pt}%
\pgfpathmoveto{\pgfqpoint{5.941666in}{2.240123in}}%
\pgfpathlineto{\pgfqpoint{5.941666in}{3.150926in}}%
\pgfusepath{stroke}%
\end{pgfscope}%
\begin{pgfscope}%
\pgfsetrectcap%
\pgfsetmiterjoin%
\pgfsetlinewidth{0.803000pt}%
\definecolor{currentstroke}{rgb}{0.000000,0.000000,0.000000}%
\pgfsetstrokecolor{currentstroke}%
\pgfsetdash{}{0pt}%
\pgfpathmoveto{\pgfqpoint{3.833026in}{2.240123in}}%
\pgfpathlineto{\pgfqpoint{5.941666in}{2.240123in}}%
\pgfusepath{stroke}%
\end{pgfscope}%
\begin{pgfscope}%
\pgfsetrectcap%
\pgfsetmiterjoin%
\pgfsetlinewidth{0.803000pt}%
\definecolor{currentstroke}{rgb}{0.000000,0.000000,0.000000}%
\pgfsetstrokecolor{currentstroke}%
\pgfsetdash{}{0pt}%
\pgfpathmoveto{\pgfqpoint{3.833026in}{3.150926in}}%
\pgfpathlineto{\pgfqpoint{5.941666in}{3.150926in}}%
\pgfusepath{stroke}%
\end{pgfscope}%
\begin{pgfscope}%
\definecolor{textcolor}{rgb}{0.000000,0.000000,0.000000}%
\pgfsetstrokecolor{textcolor}%
\pgfsetfillcolor{textcolor}%
\pgftext[x=4.887346in,y=3.234260in,,base]{\color{textcolor}\rmfamily\fontsize{12.000000}{14.400000}\selectfont 26,098 kHz}%
\end{pgfscope}%
\begin{pgfscope}%
\pgfsetbuttcap%
\pgfsetmiterjoin%
\definecolor{currentfill}{rgb}{1.000000,1.000000,1.000000}%
\pgfsetfillcolor{currentfill}%
\pgfsetlinewidth{0.000000pt}%
\definecolor{currentstroke}{rgb}{0.000000,0.000000,0.000000}%
\pgfsetstrokecolor{currentstroke}%
\pgfsetstrokeopacity{0.000000}%
\pgfsetdash{}{0pt}%
\pgfpathmoveto{\pgfqpoint{0.758026in}{0.565123in}}%
\pgfpathlineto{\pgfqpoint{2.866666in}{0.565123in}}%
\pgfpathlineto{\pgfqpoint{2.866666in}{1.475926in}}%
\pgfpathlineto{\pgfqpoint{0.758026in}{1.475926in}}%
\pgfpathlineto{\pgfqpoint{0.758026in}{0.565123in}}%
\pgfpathclose%
\pgfusepath{fill}%
\end{pgfscope}%
\begin{pgfscope}%
\pgfpathrectangle{\pgfqpoint{0.758026in}{0.565123in}}{\pgfqpoint{2.108640in}{0.910803in}}%
\pgfusepath{clip}%
\pgfsetbuttcap%
\pgfsetroundjoin%
\pgfsetlinewidth{0.501875pt}%
\definecolor{currentstroke}{rgb}{0.690196,0.690196,0.690196}%
\pgfsetstrokecolor{currentstroke}%
\pgfsetstrokeopacity{0.700000}%
\pgfsetdash{{1.850000pt}{0.800000pt}}{0.000000pt}%
\pgfpathmoveto{\pgfqpoint{0.758026in}{0.565123in}}%
\pgfpathlineto{\pgfqpoint{0.758026in}{1.475926in}}%
\pgfusepath{stroke}%
\end{pgfscope}%
\begin{pgfscope}%
\pgfsetbuttcap%
\pgfsetroundjoin%
\definecolor{currentfill}{rgb}{0.000000,0.000000,0.000000}%
\pgfsetfillcolor{currentfill}%
\pgfsetlinewidth{0.803000pt}%
\definecolor{currentstroke}{rgb}{0.000000,0.000000,0.000000}%
\pgfsetstrokecolor{currentstroke}%
\pgfsetdash{}{0pt}%
\pgfsys@defobject{currentmarker}{\pgfqpoint{0.000000in}{-0.048611in}}{\pgfqpoint{0.000000in}{0.000000in}}{%
\pgfpathmoveto{\pgfqpoint{0.000000in}{0.000000in}}%
\pgfpathlineto{\pgfqpoint{0.000000in}{-0.048611in}}%
\pgfusepath{stroke,fill}%
}%
\begin{pgfscope}%
\pgfsys@transformshift{0.758026in}{0.565123in}%
\pgfsys@useobject{currentmarker}{}%
\end{pgfscope}%
\end{pgfscope}%
\begin{pgfscope}%
\definecolor{textcolor}{rgb}{0.000000,0.000000,0.000000}%
\pgfsetstrokecolor{textcolor}%
\pgfsetfillcolor{textcolor}%
\pgftext[x=0.758026in,y=0.467901in,,top]{\color{textcolor}\rmfamily\fontsize{10.000000}{12.000000}\selectfont \(\displaystyle {100000}\)}%
\end{pgfscope}%
\begin{pgfscope}%
\pgfpathrectangle{\pgfqpoint{0.758026in}{0.565123in}}{\pgfqpoint{2.108640in}{0.910803in}}%
\pgfusepath{clip}%
\pgfsetbuttcap%
\pgfsetroundjoin%
\pgfsetlinewidth{0.501875pt}%
\definecolor{currentstroke}{rgb}{0.690196,0.690196,0.690196}%
\pgfsetstrokecolor{currentstroke}%
\pgfsetstrokeopacity{0.700000}%
\pgfsetdash{{1.850000pt}{0.800000pt}}{0.000000pt}%
\pgfpathmoveto{\pgfqpoint{1.285186in}{0.565123in}}%
\pgfpathlineto{\pgfqpoint{1.285186in}{1.475926in}}%
\pgfusepath{stroke}%
\end{pgfscope}%
\begin{pgfscope}%
\pgfsetbuttcap%
\pgfsetroundjoin%
\definecolor{currentfill}{rgb}{0.000000,0.000000,0.000000}%
\pgfsetfillcolor{currentfill}%
\pgfsetlinewidth{0.803000pt}%
\definecolor{currentstroke}{rgb}{0.000000,0.000000,0.000000}%
\pgfsetstrokecolor{currentstroke}%
\pgfsetdash{}{0pt}%
\pgfsys@defobject{currentmarker}{\pgfqpoint{0.000000in}{-0.048611in}}{\pgfqpoint{0.000000in}{0.000000in}}{%
\pgfpathmoveto{\pgfqpoint{0.000000in}{0.000000in}}%
\pgfpathlineto{\pgfqpoint{0.000000in}{-0.048611in}}%
\pgfusepath{stroke,fill}%
}%
\begin{pgfscope}%
\pgfsys@transformshift{1.285186in}{0.565123in}%
\pgfsys@useobject{currentmarker}{}%
\end{pgfscope}%
\end{pgfscope}%
\begin{pgfscope}%
\definecolor{textcolor}{rgb}{0.000000,0.000000,0.000000}%
\pgfsetstrokecolor{textcolor}%
\pgfsetfillcolor{textcolor}%
\pgftext[x=1.285186in,y=0.467901in,,top]{\color{textcolor}\rmfamily\fontsize{10.000000}{12.000000}\selectfont \(\displaystyle {200000}\)}%
\end{pgfscope}%
\begin{pgfscope}%
\pgfpathrectangle{\pgfqpoint{0.758026in}{0.565123in}}{\pgfqpoint{2.108640in}{0.910803in}}%
\pgfusepath{clip}%
\pgfsetbuttcap%
\pgfsetroundjoin%
\pgfsetlinewidth{0.501875pt}%
\definecolor{currentstroke}{rgb}{0.690196,0.690196,0.690196}%
\pgfsetstrokecolor{currentstroke}%
\pgfsetstrokeopacity{0.700000}%
\pgfsetdash{{1.850000pt}{0.800000pt}}{0.000000pt}%
\pgfpathmoveto{\pgfqpoint{1.812346in}{0.565123in}}%
\pgfpathlineto{\pgfqpoint{1.812346in}{1.475926in}}%
\pgfusepath{stroke}%
\end{pgfscope}%
\begin{pgfscope}%
\pgfsetbuttcap%
\pgfsetroundjoin%
\definecolor{currentfill}{rgb}{0.000000,0.000000,0.000000}%
\pgfsetfillcolor{currentfill}%
\pgfsetlinewidth{0.803000pt}%
\definecolor{currentstroke}{rgb}{0.000000,0.000000,0.000000}%
\pgfsetstrokecolor{currentstroke}%
\pgfsetdash{}{0pt}%
\pgfsys@defobject{currentmarker}{\pgfqpoint{0.000000in}{-0.048611in}}{\pgfqpoint{0.000000in}{0.000000in}}{%
\pgfpathmoveto{\pgfqpoint{0.000000in}{0.000000in}}%
\pgfpathlineto{\pgfqpoint{0.000000in}{-0.048611in}}%
\pgfusepath{stroke,fill}%
}%
\begin{pgfscope}%
\pgfsys@transformshift{1.812346in}{0.565123in}%
\pgfsys@useobject{currentmarker}{}%
\end{pgfscope}%
\end{pgfscope}%
\begin{pgfscope}%
\definecolor{textcolor}{rgb}{0.000000,0.000000,0.000000}%
\pgfsetstrokecolor{textcolor}%
\pgfsetfillcolor{textcolor}%
\pgftext[x=1.812346in,y=0.467901in,,top]{\color{textcolor}\rmfamily\fontsize{10.000000}{12.000000}\selectfont \(\displaystyle {300000}\)}%
\end{pgfscope}%
\begin{pgfscope}%
\pgfpathrectangle{\pgfqpoint{0.758026in}{0.565123in}}{\pgfqpoint{2.108640in}{0.910803in}}%
\pgfusepath{clip}%
\pgfsetbuttcap%
\pgfsetroundjoin%
\pgfsetlinewidth{0.501875pt}%
\definecolor{currentstroke}{rgb}{0.690196,0.690196,0.690196}%
\pgfsetstrokecolor{currentstroke}%
\pgfsetstrokeopacity{0.700000}%
\pgfsetdash{{1.850000pt}{0.800000pt}}{0.000000pt}%
\pgfpathmoveto{\pgfqpoint{2.339506in}{0.565123in}}%
\pgfpathlineto{\pgfqpoint{2.339506in}{1.475926in}}%
\pgfusepath{stroke}%
\end{pgfscope}%
\begin{pgfscope}%
\pgfsetbuttcap%
\pgfsetroundjoin%
\definecolor{currentfill}{rgb}{0.000000,0.000000,0.000000}%
\pgfsetfillcolor{currentfill}%
\pgfsetlinewidth{0.803000pt}%
\definecolor{currentstroke}{rgb}{0.000000,0.000000,0.000000}%
\pgfsetstrokecolor{currentstroke}%
\pgfsetdash{}{0pt}%
\pgfsys@defobject{currentmarker}{\pgfqpoint{0.000000in}{-0.048611in}}{\pgfqpoint{0.000000in}{0.000000in}}{%
\pgfpathmoveto{\pgfqpoint{0.000000in}{0.000000in}}%
\pgfpathlineto{\pgfqpoint{0.000000in}{-0.048611in}}%
\pgfusepath{stroke,fill}%
}%
\begin{pgfscope}%
\pgfsys@transformshift{2.339506in}{0.565123in}%
\pgfsys@useobject{currentmarker}{}%
\end{pgfscope}%
\end{pgfscope}%
\begin{pgfscope}%
\definecolor{textcolor}{rgb}{0.000000,0.000000,0.000000}%
\pgfsetstrokecolor{textcolor}%
\pgfsetfillcolor{textcolor}%
\pgftext[x=2.339506in,y=0.467901in,,top]{\color{textcolor}\rmfamily\fontsize{10.000000}{12.000000}\selectfont \(\displaystyle {400000}\)}%
\end{pgfscope}%
\begin{pgfscope}%
\pgfpathrectangle{\pgfqpoint{0.758026in}{0.565123in}}{\pgfqpoint{2.108640in}{0.910803in}}%
\pgfusepath{clip}%
\pgfsetbuttcap%
\pgfsetroundjoin%
\pgfsetlinewidth{0.501875pt}%
\definecolor{currentstroke}{rgb}{0.690196,0.690196,0.690196}%
\pgfsetstrokecolor{currentstroke}%
\pgfsetstrokeopacity{0.700000}%
\pgfsetdash{{1.850000pt}{0.800000pt}}{0.000000pt}%
\pgfpathmoveto{\pgfqpoint{2.866666in}{0.565123in}}%
\pgfpathlineto{\pgfqpoint{2.866666in}{1.475926in}}%
\pgfusepath{stroke}%
\end{pgfscope}%
\begin{pgfscope}%
\pgfsetbuttcap%
\pgfsetroundjoin%
\definecolor{currentfill}{rgb}{0.000000,0.000000,0.000000}%
\pgfsetfillcolor{currentfill}%
\pgfsetlinewidth{0.803000pt}%
\definecolor{currentstroke}{rgb}{0.000000,0.000000,0.000000}%
\pgfsetstrokecolor{currentstroke}%
\pgfsetdash{}{0pt}%
\pgfsys@defobject{currentmarker}{\pgfqpoint{0.000000in}{-0.048611in}}{\pgfqpoint{0.000000in}{0.000000in}}{%
\pgfpathmoveto{\pgfqpoint{0.000000in}{0.000000in}}%
\pgfpathlineto{\pgfqpoint{0.000000in}{-0.048611in}}%
\pgfusepath{stroke,fill}%
}%
\begin{pgfscope}%
\pgfsys@transformshift{2.866666in}{0.565123in}%
\pgfsys@useobject{currentmarker}{}%
\end{pgfscope}%
\end{pgfscope}%
\begin{pgfscope}%
\definecolor{textcolor}{rgb}{0.000000,0.000000,0.000000}%
\pgfsetstrokecolor{textcolor}%
\pgfsetfillcolor{textcolor}%
\pgftext[x=2.866666in,y=0.467901in,,top]{\color{textcolor}\rmfamily\fontsize{10.000000}{12.000000}\selectfont \(\displaystyle {500000}\)}%
\end{pgfscope}%
\begin{pgfscope}%
\definecolor{textcolor}{rgb}{0.000000,0.000000,0.000000}%
\pgfsetstrokecolor{textcolor}%
\pgfsetfillcolor{textcolor}%
\pgftext[x=1.812346in,y=0.288889in,,top]{\color{textcolor}\rmfamily\fontsize{10.000000}{12.000000}\selectfont Frequency (Hz)}%
\end{pgfscope}%
\begin{pgfscope}%
\pgfpathrectangle{\pgfqpoint{0.758026in}{0.565123in}}{\pgfqpoint{2.108640in}{0.910803in}}%
\pgfusepath{clip}%
\pgfsetbuttcap%
\pgfsetroundjoin%
\pgfsetlinewidth{0.501875pt}%
\definecolor{currentstroke}{rgb}{0.690196,0.690196,0.690196}%
\pgfsetstrokecolor{currentstroke}%
\pgfsetstrokeopacity{0.700000}%
\pgfsetdash{{1.850000pt}{0.800000pt}}{0.000000pt}%
\pgfpathmoveto{\pgfqpoint{0.758026in}{0.733415in}}%
\pgfpathlineto{\pgfqpoint{2.866666in}{0.733415in}}%
\pgfusepath{stroke}%
\end{pgfscope}%
\begin{pgfscope}%
\pgfsetbuttcap%
\pgfsetroundjoin%
\definecolor{currentfill}{rgb}{0.000000,0.000000,0.000000}%
\pgfsetfillcolor{currentfill}%
\pgfsetlinewidth{0.803000pt}%
\definecolor{currentstroke}{rgb}{0.000000,0.000000,0.000000}%
\pgfsetstrokecolor{currentstroke}%
\pgfsetdash{}{0pt}%
\pgfsys@defobject{currentmarker}{\pgfqpoint{-0.048611in}{0.000000in}}{\pgfqpoint{-0.000000in}{0.000000in}}{%
\pgfpathmoveto{\pgfqpoint{-0.000000in}{0.000000in}}%
\pgfpathlineto{\pgfqpoint{-0.048611in}{0.000000in}}%
\pgfusepath{stroke,fill}%
}%
\begin{pgfscope}%
\pgfsys@transformshift{0.758026in}{0.733415in}%
\pgfsys@useobject{currentmarker}{}%
\end{pgfscope}%
\end{pgfscope}%
\begin{pgfscope}%
\definecolor{textcolor}{rgb}{0.000000,0.000000,0.000000}%
\pgfsetstrokecolor{textcolor}%
\pgfsetfillcolor{textcolor}%
\pgftext[x=0.344444in, y=0.685189in, left, base]{\color{textcolor}\rmfamily\fontsize{10.000000}{12.000000}\selectfont \(\displaystyle {\ensuremath{-}100}\)}%
\end{pgfscope}%
\begin{pgfscope}%
\pgfpathrectangle{\pgfqpoint{0.758026in}{0.565123in}}{\pgfqpoint{2.108640in}{0.910803in}}%
\pgfusepath{clip}%
\pgfsetbuttcap%
\pgfsetroundjoin%
\pgfsetlinewidth{0.501875pt}%
\definecolor{currentstroke}{rgb}{0.690196,0.690196,0.690196}%
\pgfsetstrokecolor{currentstroke}%
\pgfsetstrokeopacity{0.700000}%
\pgfsetdash{{1.850000pt}{0.800000pt}}{0.000000pt}%
\pgfpathmoveto{\pgfqpoint{0.758026in}{1.222633in}}%
\pgfpathlineto{\pgfqpoint{2.866666in}{1.222633in}}%
\pgfusepath{stroke}%
\end{pgfscope}%
\begin{pgfscope}%
\pgfsetbuttcap%
\pgfsetroundjoin%
\definecolor{currentfill}{rgb}{0.000000,0.000000,0.000000}%
\pgfsetfillcolor{currentfill}%
\pgfsetlinewidth{0.803000pt}%
\definecolor{currentstroke}{rgb}{0.000000,0.000000,0.000000}%
\pgfsetstrokecolor{currentstroke}%
\pgfsetdash{}{0pt}%
\pgfsys@defobject{currentmarker}{\pgfqpoint{-0.048611in}{0.000000in}}{\pgfqpoint{-0.000000in}{0.000000in}}{%
\pgfpathmoveto{\pgfqpoint{-0.000000in}{0.000000in}}%
\pgfpathlineto{\pgfqpoint{-0.048611in}{0.000000in}}%
\pgfusepath{stroke,fill}%
}%
\begin{pgfscope}%
\pgfsys@transformshift{0.758026in}{1.222633in}%
\pgfsys@useobject{currentmarker}{}%
\end{pgfscope}%
\end{pgfscope}%
\begin{pgfscope}%
\definecolor{textcolor}{rgb}{0.000000,0.000000,0.000000}%
\pgfsetstrokecolor{textcolor}%
\pgfsetfillcolor{textcolor}%
\pgftext[x=0.413889in, y=1.174408in, left, base]{\color{textcolor}\rmfamily\fontsize{10.000000}{12.000000}\selectfont \(\displaystyle {\ensuremath{-}50}\)}%
\end{pgfscope}%
\begin{pgfscope}%
\definecolor{textcolor}{rgb}{0.000000,0.000000,0.000000}%
\pgfsetstrokecolor{textcolor}%
\pgfsetfillcolor{textcolor}%
\pgftext[x=0.288889in,y=1.020525in,,bottom,rotate=90.000000]{\color{textcolor}\rmfamily\fontsize{10.000000}{12.000000}\selectfont Amplitude (dB)}%
\end{pgfscope}%
\begin{pgfscope}%
\pgfpathrectangle{\pgfqpoint{0.758026in}{0.565123in}}{\pgfqpoint{2.108640in}{0.910803in}}%
\pgfusepath{clip}%
\pgfsetrectcap%
\pgfsetroundjoin%
\pgfsetlinewidth{1.003750pt}%
\definecolor{currentstroke}{rgb}{0.000000,0.000000,0.000000}%
\pgfsetstrokecolor{currentstroke}%
\pgfsetdash{}{0pt}%
\pgfpathmoveto{\pgfqpoint{0.758026in}{0.937969in}}%
\pgfpathlineto{\pgfqpoint{0.762770in}{1.099411in}}%
\pgfpathlineto{\pgfqpoint{0.767515in}{1.344021in}}%
\pgfpathlineto{\pgfqpoint{0.772259in}{1.383158in}}%
\pgfpathlineto{\pgfqpoint{0.777003in}{1.324452in}}%
\pgfpathlineto{\pgfqpoint{0.781748in}{1.026029in}}%
\pgfpathlineto{\pgfqpoint{0.786492in}{0.972214in}}%
\pgfpathlineto{\pgfqpoint{0.791237in}{0.903724in}}%
\pgfpathlineto{\pgfqpoint{0.795981in}{0.898832in}}%
\pgfpathlineto{\pgfqpoint{0.800726in}{0.923293in}}%
\pgfpathlineto{\pgfqpoint{0.805470in}{0.903418in}}%
\pgfpathlineto{\pgfqpoint{0.810215in}{0.922987in}}%
\pgfpathlineto{\pgfqpoint{0.814959in}{0.913202in}}%
\pgfpathlineto{\pgfqpoint{0.819703in}{0.952340in}}%
\pgfpathlineto{\pgfqpoint{0.824448in}{0.918095in}}%
\pgfpathlineto{\pgfqpoint{0.829192in}{0.932771in}}%
\pgfpathlineto{\pgfqpoint{0.833937in}{0.927879in}}%
\pgfpathlineto{\pgfqpoint{0.838681in}{0.957232in}}%
\pgfpathlineto{\pgfqpoint{0.843426in}{0.927573in}}%
\pgfpathlineto{\pgfqpoint{0.848170in}{0.937358in}}%
\pgfpathlineto{\pgfqpoint{0.852915in}{0.912897in}}%
\pgfpathlineto{\pgfqpoint{0.857659in}{0.932465in}}%
\pgfpathlineto{\pgfqpoint{0.862403in}{0.903112in}}%
\pgfpathlineto{\pgfqpoint{0.867148in}{1.240673in}}%
\pgfpathlineto{\pgfqpoint{0.871892in}{1.411900in}}%
\pgfpathlineto{\pgfqpoint{0.876637in}{1.421684in}}%
\pgfpathlineto{\pgfqpoint{0.881381in}{1.303966in}}%
\pgfpathlineto{\pgfqpoint{0.886126in}{0.927267in}}%
\pgfpathlineto{\pgfqpoint{0.890870in}{0.966405in}}%
\pgfpathlineto{\pgfqpoint{0.895614in}{0.907699in}}%
\pgfpathlineto{\pgfqpoint{0.900359in}{0.922375in}}%
\pgfpathlineto{\pgfqpoint{0.905103in}{0.902807in}}%
\pgfpathlineto{\pgfqpoint{0.914592in}{0.912285in}}%
\pgfpathlineto{\pgfqpoint{0.919337in}{0.941638in}}%
\pgfpathlineto{\pgfqpoint{0.924081in}{0.912285in}}%
\pgfpathlineto{\pgfqpoint{0.928826in}{0.941638in}}%
\pgfpathlineto{\pgfqpoint{0.933570in}{0.936746in}}%
\pgfpathlineto{\pgfqpoint{0.938314in}{0.951423in}}%
\pgfpathlineto{\pgfqpoint{0.943059in}{0.912285in}}%
\pgfpathlineto{\pgfqpoint{0.947803in}{0.926962in}}%
\pgfpathlineto{\pgfqpoint{0.957292in}{0.970686in}}%
\pgfpathlineto{\pgfqpoint{0.962037in}{0.946225in}}%
\pgfpathlineto{\pgfqpoint{0.966781in}{0.956009in}}%
\pgfpathlineto{\pgfqpoint{0.971526in}{1.322923in}}%
\pgfpathlineto{\pgfqpoint{0.976270in}{1.406090in}}%
\pgfpathlineto{\pgfqpoint{0.981014in}{1.391414in}}%
\pgfpathlineto{\pgfqpoint{0.990503in}{0.950811in}}%
\pgfpathlineto{\pgfqpoint{0.995248in}{0.955703in}}%
\pgfpathlineto{\pgfqpoint{0.999992in}{0.955703in}}%
\pgfpathlineto{\pgfqpoint{1.004737in}{0.941027in}}%
\pgfpathlineto{\pgfqpoint{1.009481in}{0.945919in}}%
\pgfpathlineto{\pgfqpoint{1.014225in}{0.926350in}}%
\pgfpathlineto{\pgfqpoint{1.018970in}{0.867644in}}%
\pgfpathlineto{\pgfqpoint{1.023714in}{0.936135in}}%
\pgfpathlineto{\pgfqpoint{1.028459in}{0.887213in}}%
\pgfpathlineto{\pgfqpoint{1.033203in}{0.940721in}}%
\pgfpathlineto{\pgfqpoint{1.037948in}{0.911368in}}%
\pgfpathlineto{\pgfqpoint{1.042692in}{0.911368in}}%
\pgfpathlineto{\pgfqpoint{1.047437in}{0.916260in}}%
\pgfpathlineto{\pgfqpoint{1.052181in}{0.926044in}}%
\pgfpathlineto{\pgfqpoint{1.056925in}{0.945613in}}%
\pgfpathlineto{\pgfqpoint{1.061670in}{0.911368in}}%
\pgfpathlineto{\pgfqpoint{1.066414in}{0.857248in}}%
\pgfpathlineto{\pgfqpoint{1.071159in}{0.920846in}}%
\pgfpathlineto{\pgfqpoint{1.075903in}{0.886601in}}%
\pgfpathlineto{\pgfqpoint{1.080648in}{0.935523in}}%
\pgfpathlineto{\pgfqpoint{1.085392in}{0.911062in}}%
\pgfpathlineto{\pgfqpoint{1.090137in}{0.847464in}}%
\pgfpathlineto{\pgfqpoint{1.094881in}{0.969768in}}%
\pgfpathlineto{\pgfqpoint{1.099625in}{0.915954in}}%
\pgfpathlineto{\pgfqpoint{1.104370in}{0.920846in}}%
\pgfpathlineto{\pgfqpoint{1.113859in}{0.881403in}}%
\pgfpathlineto{\pgfqpoint{1.118603in}{0.886295in}}%
\pgfpathlineto{\pgfqpoint{1.123348in}{0.842266in}}%
\pgfpathlineto{\pgfqpoint{1.128092in}{0.920541in}}%
\pgfpathlineto{\pgfqpoint{1.132837in}{0.886295in}}%
\pgfpathlineto{\pgfqpoint{1.137581in}{0.910756in}}%
\pgfpathlineto{\pgfqpoint{1.142325in}{0.920235in}}%
\pgfpathlineto{\pgfqpoint{1.147070in}{0.900666in}}%
\pgfpathlineto{\pgfqpoint{1.151814in}{0.876205in}}%
\pgfpathlineto{\pgfqpoint{1.156559in}{0.930019in}}%
\pgfpathlineto{\pgfqpoint{1.161303in}{0.930019in}}%
\pgfpathlineto{\pgfqpoint{1.166048in}{0.885990in}}%
\pgfpathlineto{\pgfqpoint{1.170792in}{0.930019in}}%
\pgfpathlineto{\pgfqpoint{1.175536in}{1.252904in}}%
\pgfpathlineto{\pgfqpoint{1.180281in}{1.404256in}}%
\pgfpathlineto{\pgfqpoint{1.185025in}{1.414040in}}%
\pgfpathlineto{\pgfqpoint{1.189770in}{1.281951in}}%
\pgfpathlineto{\pgfqpoint{1.194514in}{0.954175in}}%
\pgfpathlineto{\pgfqpoint{1.199259in}{0.929714in}}%
\pgfpathlineto{\pgfqpoint{1.204003in}{0.934606in}}%
\pgfpathlineto{\pgfqpoint{1.208748in}{0.993312in}}%
\pgfpathlineto{\pgfqpoint{1.213492in}{0.968851in}}%
\pgfpathlineto{\pgfqpoint{1.218236in}{0.909839in}}%
\pgfpathlineto{\pgfqpoint{1.222981in}{0.914731in}}%
\pgfpathlineto{\pgfqpoint{1.227725in}{0.836456in}}%
\pgfpathlineto{\pgfqpoint{1.232470in}{0.909839in}}%
\pgfpathlineto{\pgfqpoint{1.237214in}{0.909839in}}%
\pgfpathlineto{\pgfqpoint{1.241959in}{0.929408in}}%
\pgfpathlineto{\pgfqpoint{1.246703in}{0.909839in}}%
\pgfpathlineto{\pgfqpoint{1.251448in}{0.880486in}}%
\pgfpathlineto{\pgfqpoint{1.256192in}{0.885072in}}%
\pgfpathlineto{\pgfqpoint{1.260936in}{0.958455in}}%
\pgfpathlineto{\pgfqpoint{1.265681in}{0.904641in}}%
\pgfpathlineto{\pgfqpoint{1.270425in}{0.875288in}}%
\pgfpathlineto{\pgfqpoint{1.275170in}{0.889965in}}%
\pgfpathlineto{\pgfqpoint{1.279914in}{1.247094in}}%
\pgfpathlineto{\pgfqpoint{1.284659in}{1.335154in}}%
\pgfpathlineto{\pgfqpoint{1.289403in}{1.310693in}}%
\pgfpathlineto{\pgfqpoint{1.298892in}{0.923904in}}%
\pgfpathlineto{\pgfqpoint{1.303636in}{0.933688in}}%
\pgfpathlineto{\pgfqpoint{1.308381in}{0.894551in}}%
\pgfpathlineto{\pgfqpoint{1.313125in}{0.992395in}}%
\pgfpathlineto{\pgfqpoint{1.322614in}{0.909228in}}%
\pgfpathlineto{\pgfqpoint{1.327359in}{0.923904in}}%
\pgfpathlineto{\pgfqpoint{1.332103in}{0.904030in}}%
\pgfpathlineto{\pgfqpoint{1.336847in}{0.948059in}}%
\pgfpathlineto{\pgfqpoint{1.341592in}{0.899137in}}%
\pgfpathlineto{\pgfqpoint{1.346336in}{0.933383in}}%
\pgfpathlineto{\pgfqpoint{1.351081in}{0.879569in}}%
\pgfpathlineto{\pgfqpoint{1.355825in}{0.913814in}}%
\pgfpathlineto{\pgfqpoint{1.360570in}{0.879569in}}%
\pgfpathlineto{\pgfqpoint{1.365314in}{0.894245in}}%
\pgfpathlineto{\pgfqpoint{1.370059in}{0.937969in}}%
\pgfpathlineto{\pgfqpoint{1.374803in}{0.937969in}}%
\pgfpathlineto{\pgfqpoint{1.384292in}{1.378266in}}%
\pgfpathlineto{\pgfqpoint{1.389036in}{1.417403in}}%
\pgfpathlineto{\pgfqpoint{1.393781in}{1.348913in}}%
\pgfpathlineto{\pgfqpoint{1.398525in}{1.030921in}}%
\pgfpathlineto{\pgfqpoint{1.403270in}{0.947754in}}%
\pgfpathlineto{\pgfqpoint{1.408014in}{0.898526in}}%
\pgfpathlineto{\pgfqpoint{1.412759in}{0.996370in}}%
\pgfpathlineto{\pgfqpoint{1.417503in}{0.991477in}}%
\pgfpathlineto{\pgfqpoint{1.422247in}{0.922987in}}%
\pgfpathlineto{\pgfqpoint{1.426992in}{0.898526in}}%
\pgfpathlineto{\pgfqpoint{1.431736in}{0.927879in}}%
\pgfpathlineto{\pgfqpoint{1.436481in}{0.874065in}}%
\pgfpathlineto{\pgfqpoint{1.441225in}{0.913202in}}%
\pgfpathlineto{\pgfqpoint{1.445970in}{0.922681in}}%
\pgfpathlineto{\pgfqpoint{1.450714in}{0.873759in}}%
\pgfpathlineto{\pgfqpoint{1.460203in}{0.937358in}}%
\pgfpathlineto{\pgfqpoint{1.464947in}{0.932465in}}%
\pgfpathlineto{\pgfqpoint{1.469692in}{0.893328in}}%
\pgfpathlineto{\pgfqpoint{1.474436in}{0.908004in}}%
\pgfpathlineto{\pgfqpoint{1.479181in}{0.927573in}}%
\pgfpathlineto{\pgfqpoint{1.483925in}{0.888130in}}%
\pgfpathlineto{\pgfqpoint{1.488670in}{0.907699in}}%
\pgfpathlineto{\pgfqpoint{1.493414in}{0.863669in}}%
\pgfpathlineto{\pgfqpoint{1.498158in}{0.878346in}}%
\pgfpathlineto{\pgfqpoint{1.502903in}{0.961513in}}%
\pgfpathlineto{\pgfqpoint{1.507647in}{0.873453in}}%
\pgfpathlineto{\pgfqpoint{1.512392in}{0.937052in}}%
\pgfpathlineto{\pgfqpoint{1.517136in}{1.025111in}}%
\pgfpathlineto{\pgfqpoint{1.521881in}{1.015021in}}%
\pgfpathlineto{\pgfqpoint{1.526625in}{0.887824in}}%
\pgfpathlineto{\pgfqpoint{1.531370in}{0.917177in}}%
\pgfpathlineto{\pgfqpoint{1.536114in}{0.887824in}}%
\pgfpathlineto{\pgfqpoint{1.540858in}{0.887824in}}%
\pgfpathlineto{\pgfqpoint{1.545603in}{0.912285in}}%
\pgfpathlineto{\pgfqpoint{1.550347in}{0.912285in}}%
\pgfpathlineto{\pgfqpoint{1.555092in}{0.956315in}}%
\pgfpathlineto{\pgfqpoint{1.559836in}{0.887518in}}%
\pgfpathlineto{\pgfqpoint{1.564581in}{0.902195in}}%
\pgfpathlineto{\pgfqpoint{1.569325in}{0.931548in}}%
\pgfpathlineto{\pgfqpoint{1.574070in}{0.892411in}}%
\pgfpathlineto{\pgfqpoint{1.578814in}{0.916872in}}%
\pgfpathlineto{\pgfqpoint{1.583558in}{0.897303in}}%
\pgfpathlineto{\pgfqpoint{1.588303in}{1.332707in}}%
\pgfpathlineto{\pgfqpoint{1.593047in}{1.425659in}}%
\pgfpathlineto{\pgfqpoint{1.597792in}{1.391108in}}%
\pgfpathlineto{\pgfqpoint{1.602536in}{1.180744in}}%
\pgfpathlineto{\pgfqpoint{1.607281in}{0.906781in}}%
\pgfpathlineto{\pgfqpoint{1.612025in}{0.887213in}}%
\pgfpathlineto{\pgfqpoint{1.621514in}{0.975272in}}%
\pgfpathlineto{\pgfqpoint{1.626258in}{0.892105in}}%
\pgfpathlineto{\pgfqpoint{1.631003in}{0.926350in}}%
\pgfpathlineto{\pgfqpoint{1.635747in}{0.896691in}}%
\pgfpathlineto{\pgfqpoint{1.640492in}{0.960290in}}%
\pgfpathlineto{\pgfqpoint{1.645236in}{0.906476in}}%
\pgfpathlineto{\pgfqpoint{1.649981in}{0.916260in}}%
\pgfpathlineto{\pgfqpoint{1.654725in}{0.886907in}}%
\pgfpathlineto{\pgfqpoint{1.659469in}{0.950505in}}%
\pgfpathlineto{\pgfqpoint{1.664214in}{0.930937in}}%
\pgfpathlineto{\pgfqpoint{1.668958in}{0.916260in}}%
\pgfpathlineto{\pgfqpoint{1.673703in}{0.955092in}}%
\pgfpathlineto{\pgfqpoint{1.678447in}{0.925739in}}%
\pgfpathlineto{\pgfqpoint{1.683192in}{0.876817in}}%
\pgfpathlineto{\pgfqpoint{1.692681in}{1.400281in}}%
\pgfpathlineto{\pgfqpoint{1.697425in}{1.434526in}}%
\pgfpathlineto{\pgfqpoint{1.702169in}{1.366035in}}%
\pgfpathlineto{\pgfqpoint{1.706914in}{1.013798in}}%
\pgfpathlineto{\pgfqpoint{1.711658in}{0.945002in}}%
\pgfpathlineto{\pgfqpoint{1.716403in}{0.925433in}}%
\pgfpathlineto{\pgfqpoint{1.721147in}{0.930325in}}%
\pgfpathlineto{\pgfqpoint{1.725892in}{0.915649in}}%
\pgfpathlineto{\pgfqpoint{1.730636in}{0.896080in}}%
\pgfpathlineto{\pgfqpoint{1.735380in}{0.900972in}}%
\pgfpathlineto{\pgfqpoint{1.740125in}{0.881403in}}%
\pgfpathlineto{\pgfqpoint{1.744869in}{0.896080in}}%
\pgfpathlineto{\pgfqpoint{1.749614in}{0.871313in}}%
\pgfpathlineto{\pgfqpoint{1.759103in}{0.939804in}}%
\pgfpathlineto{\pgfqpoint{1.763847in}{0.925127in}}%
\pgfpathlineto{\pgfqpoint{1.768592in}{0.925127in}}%
\pgfpathlineto{\pgfqpoint{1.773336in}{0.934912in}}%
\pgfpathlineto{\pgfqpoint{1.778080in}{0.910451in}}%
\pgfpathlineto{\pgfqpoint{1.782825in}{0.895774in}}%
\pgfpathlineto{\pgfqpoint{1.787569in}{0.934606in}}%
\pgfpathlineto{\pgfqpoint{1.792314in}{1.247706in}}%
\pgfpathlineto{\pgfqpoint{1.797058in}{1.394471in}}%
\pgfpathlineto{\pgfqpoint{1.801803in}{1.399364in}}%
\pgfpathlineto{\pgfqpoint{1.806547in}{1.262382in}}%
\pgfpathlineto{\pgfqpoint{1.811292in}{0.915037in}}%
\pgfpathlineto{\pgfqpoint{1.816036in}{0.934606in}}%
\pgfpathlineto{\pgfqpoint{1.820780in}{0.963959in}}%
\pgfpathlineto{\pgfqpoint{1.825525in}{0.973437in}}%
\pgfpathlineto{\pgfqpoint{1.830269in}{0.997898in}}%
\pgfpathlineto{\pgfqpoint{1.835014in}{0.909839in}}%
\pgfpathlineto{\pgfqpoint{1.839758in}{0.939192in}}%
\pgfpathlineto{\pgfqpoint{1.844503in}{0.875594in}}%
\pgfpathlineto{\pgfqpoint{1.853992in}{0.929408in}}%
\pgfpathlineto{\pgfqpoint{1.858736in}{0.909839in}}%
\pgfpathlineto{\pgfqpoint{1.868225in}{0.899749in}}%
\pgfpathlineto{\pgfqpoint{1.872969in}{0.929102in}}%
\pgfpathlineto{\pgfqpoint{1.877714in}{0.880180in}}%
\pgfpathlineto{\pgfqpoint{1.882458in}{0.904641in}}%
\pgfpathlineto{\pgfqpoint{1.887203in}{0.909533in}}%
\pgfpathlineto{\pgfqpoint{1.891947in}{0.889965in}}%
\pgfpathlineto{\pgfqpoint{1.896691in}{1.271555in}}%
\pgfpathlineto{\pgfqpoint{1.901436in}{1.354417in}}%
\pgfpathlineto{\pgfqpoint{1.906180in}{1.325063in}}%
\pgfpathlineto{\pgfqpoint{1.910925in}{1.095131in}}%
\pgfpathlineto{\pgfqpoint{1.915669in}{0.923904in}}%
\pgfpathlineto{\pgfqpoint{1.925158in}{0.904335in}}%
\pgfpathlineto{\pgfqpoint{1.929903in}{0.953257in}}%
\pgfpathlineto{\pgfqpoint{1.934647in}{0.914120in}}%
\pgfpathlineto{\pgfqpoint{1.939391in}{0.889353in}}%
\pgfpathlineto{\pgfqpoint{1.944136in}{0.943167in}}%
\pgfpathlineto{\pgfqpoint{1.948880in}{0.855108in}}%
\pgfpathlineto{\pgfqpoint{1.953625in}{0.894245in}}%
\pgfpathlineto{\pgfqpoint{1.958369in}{0.884461in}}%
\pgfpathlineto{\pgfqpoint{1.963114in}{0.913814in}}%
\pgfpathlineto{\pgfqpoint{1.967858in}{0.904030in}}%
\pgfpathlineto{\pgfqpoint{1.972603in}{0.860000in}}%
\pgfpathlineto{\pgfqpoint{1.977347in}{0.898832in}}%
\pgfpathlineto{\pgfqpoint{1.986836in}{0.928185in}}%
\pgfpathlineto{\pgfqpoint{1.991580in}{0.898832in}}%
\pgfpathlineto{\pgfqpoint{1.996325in}{1.001568in}}%
\pgfpathlineto{\pgfqpoint{2.001069in}{1.246177in}}%
\pgfpathlineto{\pgfqpoint{2.005814in}{1.280422in}}%
\pgfpathlineto{\pgfqpoint{2.010558in}{1.197255in}}%
\pgfpathlineto{\pgfqpoint{2.015303in}{0.908310in}}%
\pgfpathlineto{\pgfqpoint{2.020047in}{0.927879in}}%
\pgfpathlineto{\pgfqpoint{2.024791in}{0.874065in}}%
\pgfpathlineto{\pgfqpoint{2.029536in}{0.903418in}}%
\pgfpathlineto{\pgfqpoint{2.034280in}{0.986585in}}%
\pgfpathlineto{\pgfqpoint{2.039025in}{0.913202in}}%
\pgfpathlineto{\pgfqpoint{2.043769in}{0.893634in}}%
\pgfpathlineto{\pgfqpoint{2.048514in}{0.908310in}}%
\pgfpathlineto{\pgfqpoint{2.053258in}{0.947142in}}%
\pgfpathlineto{\pgfqpoint{2.058002in}{0.893328in}}%
\pgfpathlineto{\pgfqpoint{2.062747in}{0.873759in}}%
\pgfpathlineto{\pgfqpoint{2.067491in}{0.863975in}}%
\pgfpathlineto{\pgfqpoint{2.072236in}{0.868867in}}%
\pgfpathlineto{\pgfqpoint{2.076980in}{0.893328in}}%
\pgfpathlineto{\pgfqpoint{2.081725in}{0.893328in}}%
\pgfpathlineto{\pgfqpoint{2.086469in}{0.844100in}}%
\pgfpathlineto{\pgfqpoint{2.091214in}{0.932160in}}%
\pgfpathlineto{\pgfqpoint{2.095958in}{0.956621in}}%
\pgfpathlineto{\pgfqpoint{2.100702in}{1.030003in}}%
\pgfpathlineto{\pgfqpoint{2.105447in}{1.181661in}}%
\pgfpathlineto{\pgfqpoint{2.110191in}{1.186553in}}%
\pgfpathlineto{\pgfqpoint{2.114936in}{1.030003in}}%
\pgfpathlineto{\pgfqpoint{2.119680in}{0.927267in}}%
\pgfpathlineto{\pgfqpoint{2.124425in}{0.887824in}}%
\pgfpathlineto{\pgfqpoint{2.129169in}{0.878040in}}%
\pgfpathlineto{\pgfqpoint{2.133914in}{0.995452in}}%
\pgfpathlineto{\pgfqpoint{2.138658in}{0.985668in}}%
\pgfpathlineto{\pgfqpoint{2.143402in}{0.873148in}}%
\pgfpathlineto{\pgfqpoint{2.148147in}{0.887824in}}%
\pgfpathlineto{\pgfqpoint{2.152891in}{0.868255in}}%
\pgfpathlineto{\pgfqpoint{2.157636in}{0.858471in}}%
\pgfpathlineto{\pgfqpoint{2.162380in}{0.912285in}}%
\pgfpathlineto{\pgfqpoint{2.171869in}{0.858165in}}%
\pgfpathlineto{\pgfqpoint{2.176613in}{0.941333in}}%
\pgfpathlineto{\pgfqpoint{2.181358in}{0.877734in}}%
\pgfpathlineto{\pgfqpoint{2.186102in}{0.863058in}}%
\pgfpathlineto{\pgfqpoint{2.190847in}{0.936440in}}%
\pgfpathlineto{\pgfqpoint{2.195591in}{0.867950in}}%
\pgfpathlineto{\pgfqpoint{2.200336in}{0.882626in}}%
\pgfpathlineto{\pgfqpoint{2.205080in}{1.004625in}}%
\pgfpathlineto{\pgfqpoint{2.209825in}{1.092685in}}%
\pgfpathlineto{\pgfqpoint{2.214569in}{1.053547in}}%
\pgfpathlineto{\pgfqpoint{2.219313in}{0.877428in}}%
\pgfpathlineto{\pgfqpoint{2.224058in}{0.882321in}}%
\pgfpathlineto{\pgfqpoint{2.228802in}{0.916566in}}%
\pgfpathlineto{\pgfqpoint{2.233547in}{0.862752in}}%
\pgfpathlineto{\pgfqpoint{2.238291in}{0.891799in}}%
\pgfpathlineto{\pgfqpoint{2.243036in}{0.891799in}}%
\pgfpathlineto{\pgfqpoint{2.252525in}{0.872230in}}%
\pgfpathlineto{\pgfqpoint{2.257269in}{0.911368in}}%
\pgfpathlineto{\pgfqpoint{2.262013in}{0.862446in}}%
\pgfpathlineto{\pgfqpoint{2.271502in}{0.921152in}}%
\pgfpathlineto{\pgfqpoint{2.276247in}{0.940721in}}%
\pgfpathlineto{\pgfqpoint{2.280991in}{0.871925in}}%
\pgfpathlineto{\pgfqpoint{2.285736in}{0.901278in}}%
\pgfpathlineto{\pgfqpoint{2.290480in}{0.896386in}}%
\pgfpathlineto{\pgfqpoint{2.295225in}{0.886601in}}%
\pgfpathlineto{\pgfqpoint{2.299969in}{0.920846in}}%
\pgfpathlineto{\pgfqpoint{2.304713in}{0.886601in}}%
\pgfpathlineto{\pgfqpoint{2.309458in}{0.920846in}}%
\pgfpathlineto{\pgfqpoint{2.314202in}{0.993924in}}%
\pgfpathlineto{\pgfqpoint{2.318947in}{0.900972in}}%
\pgfpathlineto{\pgfqpoint{2.323691in}{0.900972in}}%
\pgfpathlineto{\pgfqpoint{2.328436in}{0.852050in}}%
\pgfpathlineto{\pgfqpoint{2.333180in}{0.876511in}}%
\pgfpathlineto{\pgfqpoint{2.337924in}{0.959678in}}%
\pgfpathlineto{\pgfqpoint{2.342669in}{1.013492in}}%
\pgfpathlineto{\pgfqpoint{2.347413in}{0.891188in}}%
\pgfpathlineto{\pgfqpoint{2.352158in}{0.886295in}}%
\pgfpathlineto{\pgfqpoint{2.356902in}{0.900972in}}%
\pgfpathlineto{\pgfqpoint{2.361647in}{0.846852in}}%
\pgfpathlineto{\pgfqpoint{2.366391in}{0.910451in}}%
\pgfpathlineto{\pgfqpoint{2.371136in}{0.890882in}}%
\pgfpathlineto{\pgfqpoint{2.375880in}{0.856637in}}%
\pgfpathlineto{\pgfqpoint{2.380624in}{0.881097in}}%
\pgfpathlineto{\pgfqpoint{2.385369in}{0.881097in}}%
\pgfpathlineto{\pgfqpoint{2.390113in}{0.866421in}}%
\pgfpathlineto{\pgfqpoint{2.394858in}{0.905253in}}%
\pgfpathlineto{\pgfqpoint{2.399602in}{0.836762in}}%
\pgfpathlineto{\pgfqpoint{2.404347in}{0.871007in}}%
\pgfpathlineto{\pgfqpoint{2.413836in}{0.924821in}}%
\pgfpathlineto{\pgfqpoint{2.418580in}{0.826978in}}%
\pgfpathlineto{\pgfqpoint{2.423324in}{0.929714in}}%
\pgfpathlineto{\pgfqpoint{2.428069in}{0.905253in}}%
\pgfpathlineto{\pgfqpoint{2.432813in}{0.909839in}}%
\pgfpathlineto{\pgfqpoint{2.437558in}{0.890270in}}%
\pgfpathlineto{\pgfqpoint{2.442302in}{0.846241in}}%
\pgfpathlineto{\pgfqpoint{2.447047in}{0.929408in}}%
\pgfpathlineto{\pgfqpoint{2.451791in}{0.914731in}}%
\pgfpathlineto{\pgfqpoint{2.456535in}{0.846241in}}%
\pgfpathlineto{\pgfqpoint{2.461280in}{0.836456in}}%
\pgfpathlineto{\pgfqpoint{2.470769in}{0.865809in}}%
\pgfpathlineto{\pgfqpoint{2.475513in}{0.850827in}}%
\pgfpathlineto{\pgfqpoint{2.480258in}{0.855719in}}%
\pgfpathlineto{\pgfqpoint{2.485002in}{0.826366in}}%
\pgfpathlineto{\pgfqpoint{2.494491in}{0.899749in}}%
\pgfpathlineto{\pgfqpoint{2.499235in}{0.865504in}}%
\pgfpathlineto{\pgfqpoint{2.503980in}{0.870396in}}%
\pgfpathlineto{\pgfqpoint{2.508724in}{0.826060in}}%
\pgfpathlineto{\pgfqpoint{2.513469in}{0.899443in}}%
\pgfpathlineto{\pgfqpoint{2.522958in}{0.879874in}}%
\pgfpathlineto{\pgfqpoint{2.527702in}{0.889659in}}%
\pgfpathlineto{\pgfqpoint{2.532447in}{0.840737in}}%
\pgfpathlineto{\pgfqpoint{2.537191in}{0.835845in}}%
\pgfpathlineto{\pgfqpoint{2.541935in}{0.889659in}}%
\pgfpathlineto{\pgfqpoint{2.546680in}{0.967934in}}%
\pgfpathlineto{\pgfqpoint{2.551424in}{0.913814in}}%
\pgfpathlineto{\pgfqpoint{2.556169in}{0.899137in}}%
\pgfpathlineto{\pgfqpoint{2.560913in}{0.889353in}}%
\pgfpathlineto{\pgfqpoint{2.565658in}{0.864892in}}%
\pgfpathlineto{\pgfqpoint{2.570402in}{0.864892in}}%
\pgfpathlineto{\pgfqpoint{2.575147in}{0.879569in}}%
\pgfpathlineto{\pgfqpoint{2.579891in}{0.874676in}}%
\pgfpathlineto{\pgfqpoint{2.584635in}{0.884155in}}%
\pgfpathlineto{\pgfqpoint{2.589380in}{0.845018in}}%
\pgfpathlineto{\pgfqpoint{2.594124in}{0.849910in}}%
\pgfpathlineto{\pgfqpoint{2.598869in}{0.840125in}}%
\pgfpathlineto{\pgfqpoint{2.603613in}{0.864586in}}%
\pgfpathlineto{\pgfqpoint{2.608358in}{0.913508in}}%
\pgfpathlineto{\pgfqpoint{2.613102in}{0.835233in}}%
\pgfpathlineto{\pgfqpoint{2.617846in}{0.840125in}}%
\pgfpathlineto{\pgfqpoint{2.622591in}{0.918400in}}%
\pgfpathlineto{\pgfqpoint{2.627335in}{0.849604in}}%
\pgfpathlineto{\pgfqpoint{2.632080in}{0.825143in}}%
\pgfpathlineto{\pgfqpoint{2.636824in}{0.874065in}}%
\pgfpathlineto{\pgfqpoint{2.646313in}{0.864281in}}%
\pgfpathlineto{\pgfqpoint{2.651058in}{0.883849in}}%
\pgfpathlineto{\pgfqpoint{2.655802in}{0.805574in}}%
\pgfpathlineto{\pgfqpoint{2.660546in}{0.883849in}}%
\pgfpathlineto{\pgfqpoint{2.665291in}{0.868867in}}%
\pgfpathlineto{\pgfqpoint{2.670035in}{0.844406in}}%
\pgfpathlineto{\pgfqpoint{2.679524in}{0.834622in}}%
\pgfpathlineto{\pgfqpoint{2.684269in}{0.883544in}}%
\pgfpathlineto{\pgfqpoint{2.689013in}{0.839514in}}%
\pgfpathlineto{\pgfqpoint{2.693758in}{0.839514in}}%
\pgfpathlineto{\pgfqpoint{2.698502in}{0.893328in}}%
\pgfpathlineto{\pgfqpoint{2.703246in}{0.858777in}}%
\pgfpathlineto{\pgfqpoint{2.707991in}{0.848992in}}%
\pgfpathlineto{\pgfqpoint{2.712735in}{0.834316in}}%
\pgfpathlineto{\pgfqpoint{2.717480in}{0.853885in}}%
\pgfpathlineto{\pgfqpoint{2.722224in}{0.829424in}}%
\pgfpathlineto{\pgfqpoint{2.726969in}{0.853885in}}%
\pgfpathlineto{\pgfqpoint{2.731713in}{0.834316in}}%
\pgfpathlineto{\pgfqpoint{2.736458in}{0.853579in}}%
\pgfpathlineto{\pgfqpoint{2.741202in}{0.853579in}}%
\pgfpathlineto{\pgfqpoint{2.745946in}{0.873148in}}%
\pgfpathlineto{\pgfqpoint{2.750691in}{0.931854in}}%
\pgfpathlineto{\pgfqpoint{2.755435in}{0.907393in}}%
\pgfpathlineto{\pgfqpoint{2.760180in}{0.819334in}}%
\pgfpathlineto{\pgfqpoint{2.764924in}{0.804657in}}%
\pgfpathlineto{\pgfqpoint{2.769669in}{0.819334in}}%
\pgfpathlineto{\pgfqpoint{2.774413in}{0.828812in}}%
\pgfpathlineto{\pgfqpoint{2.779157in}{0.843489in}}%
\pgfpathlineto{\pgfqpoint{2.783902in}{0.838597in}}%
\pgfpathlineto{\pgfqpoint{2.788646in}{0.809243in}}%
\pgfpathlineto{\pgfqpoint{2.793391in}{0.853273in}}%
\pgfpathlineto{\pgfqpoint{2.798135in}{0.804351in}}%
\pgfpathlineto{\pgfqpoint{2.802880in}{0.848381in}}%
\pgfpathlineto{\pgfqpoint{2.807624in}{0.863058in}}%
\pgfpathlineto{\pgfqpoint{2.812369in}{0.848075in}}%
\pgfpathlineto{\pgfqpoint{2.817113in}{0.852967in}}%
\pgfpathlineto{\pgfqpoint{2.821857in}{0.867644in}}%
\pgfpathlineto{\pgfqpoint{2.826602in}{0.818722in}}%
\pgfpathlineto{\pgfqpoint{2.831346in}{0.818722in}}%
\pgfpathlineto{\pgfqpoint{2.836091in}{0.877428in}}%
\pgfpathlineto{\pgfqpoint{2.840835in}{0.823614in}}%
\pgfpathlineto{\pgfqpoint{2.845580in}{0.838291in}}%
\pgfpathlineto{\pgfqpoint{2.850324in}{0.872230in}}%
\pgfpathlineto{\pgfqpoint{2.855069in}{0.926044in}}%
\pgfpathlineto{\pgfqpoint{2.859813in}{0.823309in}}%
\pgfpathlineto{\pgfqpoint{2.869302in}{0.882015in}}%
\pgfpathlineto{\pgfqpoint{2.869302in}{0.882015in}}%
\pgfusepath{stroke}%
\end{pgfscope}%
\begin{pgfscope}%
\pgfsetrectcap%
\pgfsetmiterjoin%
\pgfsetlinewidth{0.803000pt}%
\definecolor{currentstroke}{rgb}{0.000000,0.000000,0.000000}%
\pgfsetstrokecolor{currentstroke}%
\pgfsetdash{}{0pt}%
\pgfpathmoveto{\pgfqpoint{0.758026in}{0.565123in}}%
\pgfpathlineto{\pgfqpoint{0.758026in}{1.475926in}}%
\pgfusepath{stroke}%
\end{pgfscope}%
\begin{pgfscope}%
\pgfsetrectcap%
\pgfsetmiterjoin%
\pgfsetlinewidth{0.803000pt}%
\definecolor{currentstroke}{rgb}{0.000000,0.000000,0.000000}%
\pgfsetstrokecolor{currentstroke}%
\pgfsetdash{}{0pt}%
\pgfpathmoveto{\pgfqpoint{2.866666in}{0.565123in}}%
\pgfpathlineto{\pgfqpoint{2.866666in}{1.475926in}}%
\pgfusepath{stroke}%
\end{pgfscope}%
\begin{pgfscope}%
\pgfsetrectcap%
\pgfsetmiterjoin%
\pgfsetlinewidth{0.803000pt}%
\definecolor{currentstroke}{rgb}{0.000000,0.000000,0.000000}%
\pgfsetstrokecolor{currentstroke}%
\pgfsetdash{}{0pt}%
\pgfpathmoveto{\pgfqpoint{0.758026in}{0.565123in}}%
\pgfpathlineto{\pgfqpoint{2.866666in}{0.565123in}}%
\pgfusepath{stroke}%
\end{pgfscope}%
\begin{pgfscope}%
\pgfsetrectcap%
\pgfsetmiterjoin%
\pgfsetlinewidth{0.803000pt}%
\definecolor{currentstroke}{rgb}{0.000000,0.000000,0.000000}%
\pgfsetstrokecolor{currentstroke}%
\pgfsetdash{}{0pt}%
\pgfpathmoveto{\pgfqpoint{0.758026in}{1.475926in}}%
\pgfpathlineto{\pgfqpoint{2.866666in}{1.475926in}}%
\pgfusepath{stroke}%
\end{pgfscope}%
\begin{pgfscope}%
\definecolor{textcolor}{rgb}{0.000000,0.000000,0.000000}%
\pgfsetstrokecolor{textcolor}%
\pgfsetfillcolor{textcolor}%
\pgftext[x=1.812346in,y=1.559260in,,base]{\color{textcolor}\rmfamily\fontsize{12.000000}{14.400000}\selectfont 19,472 kHz}%
\end{pgfscope}%
\begin{pgfscope}%
\pgfsetbuttcap%
\pgfsetmiterjoin%
\definecolor{currentfill}{rgb}{1.000000,1.000000,1.000000}%
\pgfsetfillcolor{currentfill}%
\pgfsetlinewidth{0.000000pt}%
\definecolor{currentstroke}{rgb}{0.000000,0.000000,0.000000}%
\pgfsetstrokecolor{currentstroke}%
\pgfsetstrokeopacity{0.000000}%
\pgfsetdash{}{0pt}%
\pgfpathmoveto{\pgfqpoint{3.833026in}{0.565123in}}%
\pgfpathlineto{\pgfqpoint{5.941666in}{0.565123in}}%
\pgfpathlineto{\pgfqpoint{5.941666in}{1.475926in}}%
\pgfpathlineto{\pgfqpoint{3.833026in}{1.475926in}}%
\pgfpathlineto{\pgfqpoint{3.833026in}{0.565123in}}%
\pgfpathclose%
\pgfusepath{fill}%
\end{pgfscope}%
\begin{pgfscope}%
\pgfpathrectangle{\pgfqpoint{3.833026in}{0.565123in}}{\pgfqpoint{2.108640in}{0.910803in}}%
\pgfusepath{clip}%
\pgfsetbuttcap%
\pgfsetroundjoin%
\pgfsetlinewidth{0.501875pt}%
\definecolor{currentstroke}{rgb}{0.690196,0.690196,0.690196}%
\pgfsetstrokecolor{currentstroke}%
\pgfsetstrokeopacity{0.700000}%
\pgfsetdash{{1.850000pt}{0.800000pt}}{0.000000pt}%
\pgfpathmoveto{\pgfqpoint{3.833026in}{0.565123in}}%
\pgfpathlineto{\pgfqpoint{3.833026in}{1.475926in}}%
\pgfusepath{stroke}%
\end{pgfscope}%
\begin{pgfscope}%
\pgfsetbuttcap%
\pgfsetroundjoin%
\definecolor{currentfill}{rgb}{0.000000,0.000000,0.000000}%
\pgfsetfillcolor{currentfill}%
\pgfsetlinewidth{0.803000pt}%
\definecolor{currentstroke}{rgb}{0.000000,0.000000,0.000000}%
\pgfsetstrokecolor{currentstroke}%
\pgfsetdash{}{0pt}%
\pgfsys@defobject{currentmarker}{\pgfqpoint{0.000000in}{-0.048611in}}{\pgfqpoint{0.000000in}{0.000000in}}{%
\pgfpathmoveto{\pgfqpoint{0.000000in}{0.000000in}}%
\pgfpathlineto{\pgfqpoint{0.000000in}{-0.048611in}}%
\pgfusepath{stroke,fill}%
}%
\begin{pgfscope}%
\pgfsys@transformshift{3.833026in}{0.565123in}%
\pgfsys@useobject{currentmarker}{}%
\end{pgfscope}%
\end{pgfscope}%
\begin{pgfscope}%
\definecolor{textcolor}{rgb}{0.000000,0.000000,0.000000}%
\pgfsetstrokecolor{textcolor}%
\pgfsetfillcolor{textcolor}%
\pgftext[x=3.833026in,y=0.467901in,,top]{\color{textcolor}\rmfamily\fontsize{10.000000}{12.000000}\selectfont \(\displaystyle {100000}\)}%
\end{pgfscope}%
\begin{pgfscope}%
\pgfpathrectangle{\pgfqpoint{3.833026in}{0.565123in}}{\pgfqpoint{2.108640in}{0.910803in}}%
\pgfusepath{clip}%
\pgfsetbuttcap%
\pgfsetroundjoin%
\pgfsetlinewidth{0.501875pt}%
\definecolor{currentstroke}{rgb}{0.690196,0.690196,0.690196}%
\pgfsetstrokecolor{currentstroke}%
\pgfsetstrokeopacity{0.700000}%
\pgfsetdash{{1.850000pt}{0.800000pt}}{0.000000pt}%
\pgfpathmoveto{\pgfqpoint{4.360186in}{0.565123in}}%
\pgfpathlineto{\pgfqpoint{4.360186in}{1.475926in}}%
\pgfusepath{stroke}%
\end{pgfscope}%
\begin{pgfscope}%
\pgfsetbuttcap%
\pgfsetroundjoin%
\definecolor{currentfill}{rgb}{0.000000,0.000000,0.000000}%
\pgfsetfillcolor{currentfill}%
\pgfsetlinewidth{0.803000pt}%
\definecolor{currentstroke}{rgb}{0.000000,0.000000,0.000000}%
\pgfsetstrokecolor{currentstroke}%
\pgfsetdash{}{0pt}%
\pgfsys@defobject{currentmarker}{\pgfqpoint{0.000000in}{-0.048611in}}{\pgfqpoint{0.000000in}{0.000000in}}{%
\pgfpathmoveto{\pgfqpoint{0.000000in}{0.000000in}}%
\pgfpathlineto{\pgfqpoint{0.000000in}{-0.048611in}}%
\pgfusepath{stroke,fill}%
}%
\begin{pgfscope}%
\pgfsys@transformshift{4.360186in}{0.565123in}%
\pgfsys@useobject{currentmarker}{}%
\end{pgfscope}%
\end{pgfscope}%
\begin{pgfscope}%
\definecolor{textcolor}{rgb}{0.000000,0.000000,0.000000}%
\pgfsetstrokecolor{textcolor}%
\pgfsetfillcolor{textcolor}%
\pgftext[x=4.360186in,y=0.467901in,,top]{\color{textcolor}\rmfamily\fontsize{10.000000}{12.000000}\selectfont \(\displaystyle {200000}\)}%
\end{pgfscope}%
\begin{pgfscope}%
\pgfpathrectangle{\pgfqpoint{3.833026in}{0.565123in}}{\pgfqpoint{2.108640in}{0.910803in}}%
\pgfusepath{clip}%
\pgfsetbuttcap%
\pgfsetroundjoin%
\pgfsetlinewidth{0.501875pt}%
\definecolor{currentstroke}{rgb}{0.690196,0.690196,0.690196}%
\pgfsetstrokecolor{currentstroke}%
\pgfsetstrokeopacity{0.700000}%
\pgfsetdash{{1.850000pt}{0.800000pt}}{0.000000pt}%
\pgfpathmoveto{\pgfqpoint{4.887346in}{0.565123in}}%
\pgfpathlineto{\pgfqpoint{4.887346in}{1.475926in}}%
\pgfusepath{stroke}%
\end{pgfscope}%
\begin{pgfscope}%
\pgfsetbuttcap%
\pgfsetroundjoin%
\definecolor{currentfill}{rgb}{0.000000,0.000000,0.000000}%
\pgfsetfillcolor{currentfill}%
\pgfsetlinewidth{0.803000pt}%
\definecolor{currentstroke}{rgb}{0.000000,0.000000,0.000000}%
\pgfsetstrokecolor{currentstroke}%
\pgfsetdash{}{0pt}%
\pgfsys@defobject{currentmarker}{\pgfqpoint{0.000000in}{-0.048611in}}{\pgfqpoint{0.000000in}{0.000000in}}{%
\pgfpathmoveto{\pgfqpoint{0.000000in}{0.000000in}}%
\pgfpathlineto{\pgfqpoint{0.000000in}{-0.048611in}}%
\pgfusepath{stroke,fill}%
}%
\begin{pgfscope}%
\pgfsys@transformshift{4.887346in}{0.565123in}%
\pgfsys@useobject{currentmarker}{}%
\end{pgfscope}%
\end{pgfscope}%
\begin{pgfscope}%
\definecolor{textcolor}{rgb}{0.000000,0.000000,0.000000}%
\pgfsetstrokecolor{textcolor}%
\pgfsetfillcolor{textcolor}%
\pgftext[x=4.887346in,y=0.467901in,,top]{\color{textcolor}\rmfamily\fontsize{10.000000}{12.000000}\selectfont \(\displaystyle {300000}\)}%
\end{pgfscope}%
\begin{pgfscope}%
\pgfpathrectangle{\pgfqpoint{3.833026in}{0.565123in}}{\pgfqpoint{2.108640in}{0.910803in}}%
\pgfusepath{clip}%
\pgfsetbuttcap%
\pgfsetroundjoin%
\pgfsetlinewidth{0.501875pt}%
\definecolor{currentstroke}{rgb}{0.690196,0.690196,0.690196}%
\pgfsetstrokecolor{currentstroke}%
\pgfsetstrokeopacity{0.700000}%
\pgfsetdash{{1.850000pt}{0.800000pt}}{0.000000pt}%
\pgfpathmoveto{\pgfqpoint{5.414506in}{0.565123in}}%
\pgfpathlineto{\pgfqpoint{5.414506in}{1.475926in}}%
\pgfusepath{stroke}%
\end{pgfscope}%
\begin{pgfscope}%
\pgfsetbuttcap%
\pgfsetroundjoin%
\definecolor{currentfill}{rgb}{0.000000,0.000000,0.000000}%
\pgfsetfillcolor{currentfill}%
\pgfsetlinewidth{0.803000pt}%
\definecolor{currentstroke}{rgb}{0.000000,0.000000,0.000000}%
\pgfsetstrokecolor{currentstroke}%
\pgfsetdash{}{0pt}%
\pgfsys@defobject{currentmarker}{\pgfqpoint{0.000000in}{-0.048611in}}{\pgfqpoint{0.000000in}{0.000000in}}{%
\pgfpathmoveto{\pgfqpoint{0.000000in}{0.000000in}}%
\pgfpathlineto{\pgfqpoint{0.000000in}{-0.048611in}}%
\pgfusepath{stroke,fill}%
}%
\begin{pgfscope}%
\pgfsys@transformshift{5.414506in}{0.565123in}%
\pgfsys@useobject{currentmarker}{}%
\end{pgfscope}%
\end{pgfscope}%
\begin{pgfscope}%
\definecolor{textcolor}{rgb}{0.000000,0.000000,0.000000}%
\pgfsetstrokecolor{textcolor}%
\pgfsetfillcolor{textcolor}%
\pgftext[x=5.414506in,y=0.467901in,,top]{\color{textcolor}\rmfamily\fontsize{10.000000}{12.000000}\selectfont \(\displaystyle {400000}\)}%
\end{pgfscope}%
\begin{pgfscope}%
\pgfpathrectangle{\pgfqpoint{3.833026in}{0.565123in}}{\pgfqpoint{2.108640in}{0.910803in}}%
\pgfusepath{clip}%
\pgfsetbuttcap%
\pgfsetroundjoin%
\pgfsetlinewidth{0.501875pt}%
\definecolor{currentstroke}{rgb}{0.690196,0.690196,0.690196}%
\pgfsetstrokecolor{currentstroke}%
\pgfsetstrokeopacity{0.700000}%
\pgfsetdash{{1.850000pt}{0.800000pt}}{0.000000pt}%
\pgfpathmoveto{\pgfqpoint{5.941666in}{0.565123in}}%
\pgfpathlineto{\pgfqpoint{5.941666in}{1.475926in}}%
\pgfusepath{stroke}%
\end{pgfscope}%
\begin{pgfscope}%
\pgfsetbuttcap%
\pgfsetroundjoin%
\definecolor{currentfill}{rgb}{0.000000,0.000000,0.000000}%
\pgfsetfillcolor{currentfill}%
\pgfsetlinewidth{0.803000pt}%
\definecolor{currentstroke}{rgb}{0.000000,0.000000,0.000000}%
\pgfsetstrokecolor{currentstroke}%
\pgfsetdash{}{0pt}%
\pgfsys@defobject{currentmarker}{\pgfqpoint{0.000000in}{-0.048611in}}{\pgfqpoint{0.000000in}{0.000000in}}{%
\pgfpathmoveto{\pgfqpoint{0.000000in}{0.000000in}}%
\pgfpathlineto{\pgfqpoint{0.000000in}{-0.048611in}}%
\pgfusepath{stroke,fill}%
}%
\begin{pgfscope}%
\pgfsys@transformshift{5.941666in}{0.565123in}%
\pgfsys@useobject{currentmarker}{}%
\end{pgfscope}%
\end{pgfscope}%
\begin{pgfscope}%
\definecolor{textcolor}{rgb}{0.000000,0.000000,0.000000}%
\pgfsetstrokecolor{textcolor}%
\pgfsetfillcolor{textcolor}%
\pgftext[x=5.941666in,y=0.467901in,,top]{\color{textcolor}\rmfamily\fontsize{10.000000}{12.000000}\selectfont \(\displaystyle {500000}\)}%
\end{pgfscope}%
\begin{pgfscope}%
\definecolor{textcolor}{rgb}{0.000000,0.000000,0.000000}%
\pgfsetstrokecolor{textcolor}%
\pgfsetfillcolor{textcolor}%
\pgftext[x=4.887346in,y=0.288889in,,top]{\color{textcolor}\rmfamily\fontsize{10.000000}{12.000000}\selectfont Frequency (Hz)}%
\end{pgfscope}%
\begin{pgfscope}%
\pgfpathrectangle{\pgfqpoint{3.833026in}{0.565123in}}{\pgfqpoint{2.108640in}{0.910803in}}%
\pgfusepath{clip}%
\pgfsetbuttcap%
\pgfsetroundjoin%
\pgfsetlinewidth{0.501875pt}%
\definecolor{currentstroke}{rgb}{0.690196,0.690196,0.690196}%
\pgfsetstrokecolor{currentstroke}%
\pgfsetstrokeopacity{0.700000}%
\pgfsetdash{{1.850000pt}{0.800000pt}}{0.000000pt}%
\pgfpathmoveto{\pgfqpoint{3.833026in}{0.733368in}}%
\pgfpathlineto{\pgfqpoint{5.941666in}{0.733368in}}%
\pgfusepath{stroke}%
\end{pgfscope}%
\begin{pgfscope}%
\pgfsetbuttcap%
\pgfsetroundjoin%
\definecolor{currentfill}{rgb}{0.000000,0.000000,0.000000}%
\pgfsetfillcolor{currentfill}%
\pgfsetlinewidth{0.803000pt}%
\definecolor{currentstroke}{rgb}{0.000000,0.000000,0.000000}%
\pgfsetstrokecolor{currentstroke}%
\pgfsetdash{}{0pt}%
\pgfsys@defobject{currentmarker}{\pgfqpoint{-0.048611in}{0.000000in}}{\pgfqpoint{-0.000000in}{0.000000in}}{%
\pgfpathmoveto{\pgfqpoint{-0.000000in}{0.000000in}}%
\pgfpathlineto{\pgfqpoint{-0.048611in}{0.000000in}}%
\pgfusepath{stroke,fill}%
}%
\begin{pgfscope}%
\pgfsys@transformshift{3.833026in}{0.733368in}%
\pgfsys@useobject{currentmarker}{}%
\end{pgfscope}%
\end{pgfscope}%
\begin{pgfscope}%
\definecolor{textcolor}{rgb}{0.000000,0.000000,0.000000}%
\pgfsetstrokecolor{textcolor}%
\pgfsetfillcolor{textcolor}%
\pgftext[x=3.419444in, y=0.685142in, left, base]{\color{textcolor}\rmfamily\fontsize{10.000000}{12.000000}\selectfont \(\displaystyle {\ensuremath{-}100}\)}%
\end{pgfscope}%
\begin{pgfscope}%
\pgfpathrectangle{\pgfqpoint{3.833026in}{0.565123in}}{\pgfqpoint{2.108640in}{0.910803in}}%
\pgfusepath{clip}%
\pgfsetbuttcap%
\pgfsetroundjoin%
\pgfsetlinewidth{0.501875pt}%
\definecolor{currentstroke}{rgb}{0.690196,0.690196,0.690196}%
\pgfsetstrokecolor{currentstroke}%
\pgfsetstrokeopacity{0.700000}%
\pgfsetdash{{1.850000pt}{0.800000pt}}{0.000000pt}%
\pgfpathmoveto{\pgfqpoint{3.833026in}{1.222406in}}%
\pgfpathlineto{\pgfqpoint{5.941666in}{1.222406in}}%
\pgfusepath{stroke}%
\end{pgfscope}%
\begin{pgfscope}%
\pgfsetbuttcap%
\pgfsetroundjoin%
\definecolor{currentfill}{rgb}{0.000000,0.000000,0.000000}%
\pgfsetfillcolor{currentfill}%
\pgfsetlinewidth{0.803000pt}%
\definecolor{currentstroke}{rgb}{0.000000,0.000000,0.000000}%
\pgfsetstrokecolor{currentstroke}%
\pgfsetdash{}{0pt}%
\pgfsys@defobject{currentmarker}{\pgfqpoint{-0.048611in}{0.000000in}}{\pgfqpoint{-0.000000in}{0.000000in}}{%
\pgfpathmoveto{\pgfqpoint{-0.000000in}{0.000000in}}%
\pgfpathlineto{\pgfqpoint{-0.048611in}{0.000000in}}%
\pgfusepath{stroke,fill}%
}%
\begin{pgfscope}%
\pgfsys@transformshift{3.833026in}{1.222406in}%
\pgfsys@useobject{currentmarker}{}%
\end{pgfscope}%
\end{pgfscope}%
\begin{pgfscope}%
\definecolor{textcolor}{rgb}{0.000000,0.000000,0.000000}%
\pgfsetstrokecolor{textcolor}%
\pgfsetfillcolor{textcolor}%
\pgftext[x=3.488889in, y=1.174181in, left, base]{\color{textcolor}\rmfamily\fontsize{10.000000}{12.000000}\selectfont \(\displaystyle {\ensuremath{-}50}\)}%
\end{pgfscope}%
\begin{pgfscope}%
\definecolor{textcolor}{rgb}{0.000000,0.000000,0.000000}%
\pgfsetstrokecolor{textcolor}%
\pgfsetfillcolor{textcolor}%
\pgftext[x=3.363889in,y=1.020525in,,bottom,rotate=90.000000]{\color{textcolor}\rmfamily\fontsize{10.000000}{12.000000}\selectfont Amplitude (dB)}%
\end{pgfscope}%
\begin{pgfscope}%
\pgfpathrectangle{\pgfqpoint{3.833026in}{0.565123in}}{\pgfqpoint{2.108640in}{0.910803in}}%
\pgfusepath{clip}%
\pgfsetrectcap%
\pgfsetroundjoin%
\pgfsetlinewidth{1.003750pt}%
\definecolor{currentstroke}{rgb}{0.000000,0.000000,0.000000}%
\pgfsetstrokecolor{currentstroke}%
\pgfsetdash{}{0pt}%
\pgfpathmoveto{\pgfqpoint{3.833026in}{0.937847in}}%
\pgfpathlineto{\pgfqpoint{3.837770in}{0.942737in}}%
\pgfpathlineto{\pgfqpoint{3.842515in}{0.913395in}}%
\pgfpathlineto{\pgfqpoint{3.847259in}{0.923176in}}%
\pgfpathlineto{\pgfqpoint{3.852003in}{0.835149in}}%
\pgfpathlineto{\pgfqpoint{3.856748in}{0.884053in}}%
\pgfpathlineto{\pgfqpoint{3.861492in}{0.962299in}}%
\pgfpathlineto{\pgfqpoint{3.870981in}{0.879162in}}%
\pgfpathlineto{\pgfqpoint{3.875726in}{1.255416in}}%
\pgfpathlineto{\pgfqpoint{3.880470in}{1.397237in}}%
\pgfpathlineto{\pgfqpoint{3.885215in}{1.392347in}}%
\pgfpathlineto{\pgfqpoint{3.889959in}{1.255416in}}%
\pgfpathlineto{\pgfqpoint{3.894703in}{0.957103in}}%
\pgfpathlineto{\pgfqpoint{3.899448in}{0.878857in}}%
\pgfpathlineto{\pgfqpoint{3.904192in}{0.903308in}}%
\pgfpathlineto{\pgfqpoint{3.908937in}{0.952212in}}%
\pgfpathlineto{\pgfqpoint{3.913681in}{0.912784in}}%
\pgfpathlineto{\pgfqpoint{3.918426in}{0.937235in}}%
\pgfpathlineto{\pgfqpoint{3.923170in}{0.927455in}}%
\pgfpathlineto{\pgfqpoint{3.927915in}{0.888332in}}%
\pgfpathlineto{\pgfqpoint{3.932659in}{0.878551in}}%
\pgfpathlineto{\pgfqpoint{3.937403in}{0.956797in}}%
\pgfpathlineto{\pgfqpoint{3.942148in}{0.903003in}}%
\pgfpathlineto{\pgfqpoint{3.946892in}{0.903003in}}%
\pgfpathlineto{\pgfqpoint{3.951637in}{0.942126in}}%
\pgfpathlineto{\pgfqpoint{3.956381in}{0.863574in}}%
\pgfpathlineto{\pgfqpoint{3.961126in}{0.946711in}}%
\pgfpathlineto{\pgfqpoint{3.965870in}{0.956491in}}%
\pgfpathlineto{\pgfqpoint{3.970614in}{1.245024in}}%
\pgfpathlineto{\pgfqpoint{3.975359in}{1.406406in}}%
\pgfpathlineto{\pgfqpoint{3.980103in}{1.421078in}}%
\pgfpathlineto{\pgfqpoint{3.984848in}{1.298818in}}%
\pgfpathlineto{\pgfqpoint{3.989592in}{0.956186in}}%
\pgfpathlineto{\pgfqpoint{3.994337in}{0.926843in}}%
\pgfpathlineto{\pgfqpoint{3.999081in}{0.912172in}}%
\pgfpathlineto{\pgfqpoint{4.003826in}{0.912172in}}%
\pgfpathlineto{\pgfqpoint{4.008570in}{0.877940in}}%
\pgfpathlineto{\pgfqpoint{4.013314in}{0.926843in}}%
\pgfpathlineto{\pgfqpoint{4.018059in}{0.961076in}}%
\pgfpathlineto{\pgfqpoint{4.022803in}{0.912172in}}%
\pgfpathlineto{\pgfqpoint{4.027548in}{0.926538in}}%
\pgfpathlineto{\pgfqpoint{4.032292in}{0.897195in}}%
\pgfpathlineto{\pgfqpoint{4.037037in}{0.916757in}}%
\pgfpathlineto{\pgfqpoint{4.041781in}{0.970551in}}%
\pgfpathlineto{\pgfqpoint{4.046526in}{0.916757in}}%
\pgfpathlineto{\pgfqpoint{4.051270in}{0.897195in}}%
\pgfpathlineto{\pgfqpoint{4.056014in}{0.926538in}}%
\pgfpathlineto{\pgfqpoint{4.060759in}{0.902086in}}%
\pgfpathlineto{\pgfqpoint{4.070248in}{1.356586in}}%
\pgfpathlineto{\pgfqpoint{4.074992in}{1.390818in}}%
\pgfpathlineto{\pgfqpoint{4.079737in}{1.297901in}}%
\pgfpathlineto{\pgfqpoint{4.084481in}{0.926232in}}%
\pgfpathlineto{\pgfqpoint{4.089225in}{0.901780in}}%
\pgfpathlineto{\pgfqpoint{4.093970in}{0.847986in}}%
\pgfpathlineto{\pgfqpoint{4.098714in}{0.896890in}}%
\pgfpathlineto{\pgfqpoint{4.103459in}{0.925926in}}%
\pgfpathlineto{\pgfqpoint{4.108203in}{0.921036in}}%
\pgfpathlineto{\pgfqpoint{4.112948in}{0.930817in}}%
\pgfpathlineto{\pgfqpoint{4.117692in}{0.857461in}}%
\pgfpathlineto{\pgfqpoint{4.127181in}{0.950378in}}%
\pgfpathlineto{\pgfqpoint{4.131925in}{0.911255in}}%
\pgfpathlineto{\pgfqpoint{4.136670in}{0.921036in}}%
\pgfpathlineto{\pgfqpoint{4.141414in}{0.925621in}}%
\pgfpathlineto{\pgfqpoint{4.146159in}{0.871827in}}%
\pgfpathlineto{\pgfqpoint{4.150903in}{0.876717in}}%
\pgfpathlineto{\pgfqpoint{4.155648in}{0.959854in}}%
\pgfpathlineto{\pgfqpoint{4.160392in}{1.013648in}}%
\pgfpathlineto{\pgfqpoint{4.165137in}{1.287509in}}%
\pgfpathlineto{\pgfqpoint{4.169881in}{1.326632in}}%
\pgfpathlineto{\pgfqpoint{4.174625in}{1.272838in}}%
\pgfpathlineto{\pgfqpoint{4.179370in}{0.910950in}}%
\pgfpathlineto{\pgfqpoint{4.184114in}{0.905754in}}%
\pgfpathlineto{\pgfqpoint{4.188859in}{0.876411in}}%
\pgfpathlineto{\pgfqpoint{4.193603in}{0.895973in}}%
\pgfpathlineto{\pgfqpoint{4.198348in}{0.900863in}}%
\pgfpathlineto{\pgfqpoint{4.203092in}{0.920425in}}%
\pgfpathlineto{\pgfqpoint{4.207837in}{0.881302in}}%
\pgfpathlineto{\pgfqpoint{4.217325in}{0.944571in}}%
\pgfpathlineto{\pgfqpoint{4.222070in}{0.885886in}}%
\pgfpathlineto{\pgfqpoint{4.226814in}{0.885886in}}%
\pgfpathlineto{\pgfqpoint{4.231559in}{0.944571in}}%
\pgfpathlineto{\pgfqpoint{4.236303in}{0.915229in}}%
\pgfpathlineto{\pgfqpoint{4.241048in}{0.866325in}}%
\pgfpathlineto{\pgfqpoint{4.245792in}{0.856544in}}%
\pgfpathlineto{\pgfqpoint{4.250536in}{0.876106in}}%
\pgfpathlineto{\pgfqpoint{4.255281in}{1.007840in}}%
\pgfpathlineto{\pgfqpoint{4.260025in}{1.350167in}}%
\pgfpathlineto{\pgfqpoint{4.264770in}{1.413742in}}%
\pgfpathlineto{\pgfqpoint{4.269514in}{1.374619in}}%
\pgfpathlineto{\pgfqpoint{4.279003in}{0.919814in}}%
\pgfpathlineto{\pgfqpoint{4.283748in}{0.919814in}}%
\pgfpathlineto{\pgfqpoint{4.288492in}{0.910033in}}%
\pgfpathlineto{\pgfqpoint{4.293236in}{0.895056in}}%
\pgfpathlineto{\pgfqpoint{4.297981in}{0.934179in}}%
\pgfpathlineto{\pgfqpoint{4.302725in}{0.899946in}}%
\pgfpathlineto{\pgfqpoint{4.307470in}{0.914617in}}%
\pgfpathlineto{\pgfqpoint{4.316959in}{0.934179in}}%
\pgfpathlineto{\pgfqpoint{4.321703in}{0.904837in}}%
\pgfpathlineto{\pgfqpoint{4.326448in}{0.895056in}}%
\pgfpathlineto{\pgfqpoint{4.331192in}{0.889860in}}%
\pgfpathlineto{\pgfqpoint{4.335936in}{0.880079in}}%
\pgfpathlineto{\pgfqpoint{4.340681in}{0.938764in}}%
\pgfpathlineto{\pgfqpoint{4.345425in}{0.904531in}}%
\pgfpathlineto{\pgfqpoint{4.350170in}{0.894750in}}%
\pgfpathlineto{\pgfqpoint{4.354914in}{0.860518in}}%
\pgfpathlineto{\pgfqpoint{4.359659in}{0.938764in}}%
\pgfpathlineto{\pgfqpoint{4.364403in}{0.997448in}}%
\pgfpathlineto{\pgfqpoint{4.369148in}{0.972691in}}%
\pgfpathlineto{\pgfqpoint{4.373892in}{0.928677in}}%
\pgfpathlineto{\pgfqpoint{4.378636in}{0.869993in}}%
\pgfpathlineto{\pgfqpoint{4.383381in}{0.889554in}}%
\pgfpathlineto{\pgfqpoint{4.388125in}{0.874883in}}%
\pgfpathlineto{\pgfqpoint{4.392870in}{0.904225in}}%
\pgfpathlineto{\pgfqpoint{4.397614in}{0.923787in}}%
\pgfpathlineto{\pgfqpoint{4.402359in}{0.923787in}}%
\pgfpathlineto{\pgfqpoint{4.407103in}{0.947933in}}%
\pgfpathlineto{\pgfqpoint{4.411847in}{0.923481in}}%
\pgfpathlineto{\pgfqpoint{4.416592in}{0.884358in}}%
\pgfpathlineto{\pgfqpoint{4.421336in}{0.855016in}}%
\pgfpathlineto{\pgfqpoint{4.426081in}{0.928372in}}%
\pgfpathlineto{\pgfqpoint{4.430825in}{0.928372in}}%
\pgfpathlineto{\pgfqpoint{4.435570in}{0.894139in}}%
\pgfpathlineto{\pgfqpoint{4.440314in}{0.962604in}}%
\pgfpathlineto{\pgfqpoint{4.445059in}{0.923176in}}%
\pgfpathlineto{\pgfqpoint{4.449803in}{1.299735in}}%
\pgfpathlineto{\pgfqpoint{4.454547in}{1.417104in}}%
\pgfpathlineto{\pgfqpoint{4.459292in}{1.407323in}}%
\pgfpathlineto{\pgfqpoint{4.464036in}{1.236160in}}%
\pgfpathlineto{\pgfqpoint{4.468781in}{0.937847in}}%
\pgfpathlineto{\pgfqpoint{4.473525in}{0.928066in}}%
\pgfpathlineto{\pgfqpoint{4.478270in}{0.888943in}}%
\pgfpathlineto{\pgfqpoint{4.487759in}{0.869076in}}%
\pgfpathlineto{\pgfqpoint{4.492503in}{0.913089in}}%
\pgfpathlineto{\pgfqpoint{4.497247in}{0.908199in}}%
\pgfpathlineto{\pgfqpoint{4.501992in}{0.932651in}}%
\pgfpathlineto{\pgfqpoint{4.506736in}{0.869076in}}%
\pgfpathlineto{\pgfqpoint{4.511481in}{0.893528in}}%
\pgfpathlineto{\pgfqpoint{4.516225in}{0.888637in}}%
\pgfpathlineto{\pgfqpoint{4.520970in}{0.907893in}}%
\pgfpathlineto{\pgfqpoint{4.525714in}{0.898112in}}%
\pgfpathlineto{\pgfqpoint{4.535203in}{0.927455in}}%
\pgfpathlineto{\pgfqpoint{4.539947in}{0.937235in}}%
\pgfpathlineto{\pgfqpoint{4.544692in}{1.171974in}}%
\pgfpathlineto{\pgfqpoint{4.549436in}{1.328466in}}%
\pgfpathlineto{\pgfqpoint{4.554181in}{1.328466in}}%
\pgfpathlineto{\pgfqpoint{4.558925in}{1.191230in}}%
\pgfpathlineto{\pgfqpoint{4.563670in}{0.912478in}}%
\pgfpathlineto{\pgfqpoint{4.568414in}{0.892916in}}%
\pgfpathlineto{\pgfqpoint{4.573158in}{0.936930in}}%
\pgfpathlineto{\pgfqpoint{4.577903in}{0.936930in}}%
\pgfpathlineto{\pgfqpoint{4.582647in}{0.888026in}}%
\pgfpathlineto{\pgfqpoint{4.587392in}{0.888026in}}%
\pgfpathlineto{\pgfqpoint{4.592136in}{0.927149in}}%
\pgfpathlineto{\pgfqpoint{4.596881in}{0.912172in}}%
\pgfpathlineto{\pgfqpoint{4.601625in}{0.921953in}}%
\pgfpathlineto{\pgfqpoint{4.606370in}{0.936624in}}%
\pgfpathlineto{\pgfqpoint{4.611114in}{0.926843in}}%
\pgfpathlineto{\pgfqpoint{4.615858in}{0.936624in}}%
\pgfpathlineto{\pgfqpoint{4.620603in}{0.921953in}}%
\pgfpathlineto{\pgfqpoint{4.625347in}{0.902392in}}%
\pgfpathlineto{\pgfqpoint{4.630092in}{0.917063in}}%
\pgfpathlineto{\pgfqpoint{4.634836in}{0.955880in}}%
\pgfpathlineto{\pgfqpoint{4.639581in}{1.195509in}}%
\pgfpathlineto{\pgfqpoint{4.644325in}{1.376453in}}%
\pgfpathlineto{\pgfqpoint{4.649070in}{1.396014in}}%
\pgfpathlineto{\pgfqpoint{4.653814in}{1.283536in}}%
\pgfpathlineto{\pgfqpoint{4.658558in}{0.960770in}}%
\pgfpathlineto{\pgfqpoint{4.663303in}{0.872744in}}%
\pgfpathlineto{\pgfqpoint{4.668047in}{0.887415in}}%
\pgfpathlineto{\pgfqpoint{4.672792in}{0.862657in}}%
\pgfpathlineto{\pgfqpoint{4.682281in}{0.911561in}}%
\pgfpathlineto{\pgfqpoint{4.687025in}{0.955574in}}%
\pgfpathlineto{\pgfqpoint{4.691769in}{0.862657in}}%
\pgfpathlineto{\pgfqpoint{4.696514in}{0.916451in}}%
\pgfpathlineto{\pgfqpoint{4.701258in}{0.887109in}}%
\pgfpathlineto{\pgfqpoint{4.706003in}{0.955574in}}%
\pgfpathlineto{\pgfqpoint{4.710747in}{0.925926in}}%
\pgfpathlineto{\pgfqpoint{4.715492in}{0.960159in}}%
\pgfpathlineto{\pgfqpoint{4.720236in}{0.940598in}}%
\pgfpathlineto{\pgfqpoint{4.724981in}{0.925926in}}%
\pgfpathlineto{\pgfqpoint{4.729725in}{0.930817in}}%
\pgfpathlineto{\pgfqpoint{4.739214in}{1.405184in}}%
\pgfpathlineto{\pgfqpoint{4.743958in}{1.434526in}}%
\pgfpathlineto{\pgfqpoint{4.748703in}{1.355974in}}%
\pgfpathlineto{\pgfqpoint{4.753447in}{0.954963in}}%
\pgfpathlineto{\pgfqpoint{4.758192in}{0.989196in}}%
\pgfpathlineto{\pgfqpoint{4.762936in}{0.950073in}}%
\pgfpathlineto{\pgfqpoint{4.767681in}{0.979415in}}%
\pgfpathlineto{\pgfqpoint{4.772425in}{0.906059in}}%
\pgfpathlineto{\pgfqpoint{4.777169in}{0.915840in}}%
\pgfpathlineto{\pgfqpoint{4.781914in}{0.896279in}}%
\pgfpathlineto{\pgfqpoint{4.786658in}{0.856850in}}%
\pgfpathlineto{\pgfqpoint{4.791403in}{0.944877in}}%
\pgfpathlineto{\pgfqpoint{4.796147in}{0.969329in}}%
\pgfpathlineto{\pgfqpoint{4.800892in}{0.910644in}}%
\pgfpathlineto{\pgfqpoint{4.805636in}{0.964438in}}%
\pgfpathlineto{\pgfqpoint{4.810380in}{0.866631in}}%
\pgfpathlineto{\pgfqpoint{4.815125in}{0.925315in}}%
\pgfpathlineto{\pgfqpoint{4.819869in}{0.910644in}}%
\pgfpathlineto{\pgfqpoint{4.824614in}{0.973913in}}%
\pgfpathlineto{\pgfqpoint{4.829358in}{1.076611in}}%
\pgfpathlineto{\pgfqpoint{4.834103in}{1.370034in}}%
\pgfpathlineto{\pgfqpoint{4.838847in}{1.414048in}}%
\pgfpathlineto{\pgfqpoint{4.843592in}{1.365144in}}%
\pgfpathlineto{\pgfqpoint{4.848336in}{1.071721in}}%
\pgfpathlineto{\pgfqpoint{4.853080in}{0.959242in}}%
\pgfpathlineto{\pgfqpoint{4.857825in}{0.915229in}}%
\pgfpathlineto{\pgfqpoint{4.862569in}{0.919814in}}%
\pgfpathlineto{\pgfqpoint{4.867314in}{0.934485in}}%
\pgfpathlineto{\pgfqpoint{4.872058in}{0.875800in}}%
\pgfpathlineto{\pgfqpoint{4.876803in}{0.905142in}}%
\pgfpathlineto{\pgfqpoint{4.881547in}{0.914923in}}%
\pgfpathlineto{\pgfqpoint{4.886292in}{0.939375in}}%
\pgfpathlineto{\pgfqpoint{4.891036in}{0.939375in}}%
\pgfpathlineto{\pgfqpoint{4.895780in}{0.944265in}}%
\pgfpathlineto{\pgfqpoint{4.900525in}{0.904837in}}%
\pgfpathlineto{\pgfqpoint{4.905269in}{0.929289in}}%
\pgfpathlineto{\pgfqpoint{4.914758in}{0.870604in}}%
\pgfpathlineto{\pgfqpoint{4.919503in}{0.924398in}}%
\pgfpathlineto{\pgfqpoint{4.924247in}{0.948850in}}%
\pgfpathlineto{\pgfqpoint{4.928992in}{1.296067in}}%
\pgfpathlineto{\pgfqpoint{4.933736in}{1.374313in}}%
\pgfpathlineto{\pgfqpoint{4.938480in}{1.334885in}}%
\pgfpathlineto{\pgfqpoint{4.943225in}{1.105037in}}%
\pgfpathlineto{\pgfqpoint{4.947969in}{0.958325in}}%
\pgfpathlineto{\pgfqpoint{4.952714in}{0.904531in}}%
\pgfpathlineto{\pgfqpoint{4.957458in}{0.875189in}}%
\pgfpathlineto{\pgfqpoint{4.962203in}{0.914312in}}%
\pgfpathlineto{\pgfqpoint{4.966947in}{0.880079in}}%
\pgfpathlineto{\pgfqpoint{4.971691in}{0.884970in}}%
\pgfpathlineto{\pgfqpoint{4.976436in}{0.943348in}}%
\pgfpathlineto{\pgfqpoint{4.981180in}{0.909116in}}%
\pgfpathlineto{\pgfqpoint{4.985925in}{0.918897in}}%
\pgfpathlineto{\pgfqpoint{4.990669in}{0.933568in}}%
\pgfpathlineto{\pgfqpoint{4.995414in}{0.928677in}}%
\pgfpathlineto{\pgfqpoint{5.000158in}{0.909116in}}%
\pgfpathlineto{\pgfqpoint{5.004903in}{0.928677in}}%
\pgfpathlineto{\pgfqpoint{5.009647in}{0.899335in}}%
\pgfpathlineto{\pgfqpoint{5.014391in}{0.894139in}}%
\pgfpathlineto{\pgfqpoint{5.019136in}{0.894139in}}%
\pgfpathlineto{\pgfqpoint{5.023880in}{1.202233in}}%
\pgfpathlineto{\pgfqpoint{5.028625in}{1.309821in}}%
\pgfpathlineto{\pgfqpoint{5.033369in}{1.285369in}}%
\pgfpathlineto{\pgfqpoint{5.042858in}{0.918591in}}%
\pgfpathlineto{\pgfqpoint{5.047603in}{0.889249in}}%
\pgfpathlineto{\pgfqpoint{5.052347in}{0.937847in}}%
\pgfpathlineto{\pgfqpoint{5.057091in}{0.913395in}}%
\pgfpathlineto{\pgfqpoint{5.061836in}{0.918285in}}%
\pgfpathlineto{\pgfqpoint{5.066580in}{0.874272in}}%
\pgfpathlineto{\pgfqpoint{5.071325in}{0.903614in}}%
\pgfpathlineto{\pgfqpoint{5.076069in}{0.869381in}}%
\pgfpathlineto{\pgfqpoint{5.080814in}{0.942737in}}%
\pgfpathlineto{\pgfqpoint{5.085558in}{0.903614in}}%
\pgfpathlineto{\pgfqpoint{5.090303in}{0.878857in}}%
\pgfpathlineto{\pgfqpoint{5.099791in}{0.917980in}}%
\pgfpathlineto{\pgfqpoint{5.109280in}{0.888637in}}%
\pgfpathlineto{\pgfqpoint{5.114025in}{0.883747in}}%
\pgfpathlineto{\pgfqpoint{5.118769in}{1.098924in}}%
\pgfpathlineto{\pgfqpoint{5.123514in}{1.221183in}}%
\pgfpathlineto{\pgfqpoint{5.128258in}{1.215987in}}%
\pgfpathlineto{\pgfqpoint{5.137747in}{0.893222in}}%
\pgfpathlineto{\pgfqpoint{5.142491in}{0.898112in}}%
\pgfpathlineto{\pgfqpoint{5.147236in}{0.868770in}}%
\pgfpathlineto{\pgfqpoint{5.151980in}{0.917674in}}%
\pgfpathlineto{\pgfqpoint{5.156725in}{0.912784in}}%
\pgfpathlineto{\pgfqpoint{5.161469in}{0.932345in}}%
\pgfpathlineto{\pgfqpoint{5.166214in}{0.868465in}}%
\pgfpathlineto{\pgfqpoint{5.170958in}{0.873355in}}%
\pgfpathlineto{\pgfqpoint{5.175702in}{0.858684in}}%
\pgfpathlineto{\pgfqpoint{5.180447in}{0.883136in}}%
\pgfpathlineto{\pgfqpoint{5.185191in}{0.853793in}}%
\pgfpathlineto{\pgfqpoint{5.189936in}{0.834232in}}%
\pgfpathlineto{\pgfqpoint{5.194680in}{0.873355in}}%
\pgfpathlineto{\pgfqpoint{5.199425in}{0.858378in}}%
\pgfpathlineto{\pgfqpoint{5.204169in}{0.892611in}}%
\pgfpathlineto{\pgfqpoint{5.208914in}{0.897501in}}%
\pgfpathlineto{\pgfqpoint{5.213658in}{0.936624in}}%
\pgfpathlineto{\pgfqpoint{5.218402in}{1.127349in}}%
\pgfpathlineto{\pgfqpoint{5.223147in}{1.137130in}}%
\pgfpathlineto{\pgfqpoint{5.232636in}{0.858378in}}%
\pgfpathlineto{\pgfqpoint{5.237380in}{0.873049in}}%
\pgfpathlineto{\pgfqpoint{5.242125in}{0.858072in}}%
\pgfpathlineto{\pgfqpoint{5.246869in}{0.892305in}}%
\pgfpathlineto{\pgfqpoint{5.251613in}{0.892305in}}%
\pgfpathlineto{\pgfqpoint{5.256358in}{0.906976in}}%
\pgfpathlineto{\pgfqpoint{5.261102in}{0.902086in}}%
\pgfpathlineto{\pgfqpoint{5.265847in}{0.921647in}}%
\pgfpathlineto{\pgfqpoint{5.270591in}{0.872744in}}%
\pgfpathlineto{\pgfqpoint{5.280080in}{0.921342in}}%
\pgfpathlineto{\pgfqpoint{5.284825in}{0.896890in}}%
\pgfpathlineto{\pgfqpoint{5.289569in}{0.906671in}}%
\pgfpathlineto{\pgfqpoint{5.294313in}{0.891999in}}%
\pgfpathlineto{\pgfqpoint{5.299058in}{0.872438in}}%
\pgfpathlineto{\pgfqpoint{5.303802in}{0.828425in}}%
\pgfpathlineto{\pgfqpoint{5.313291in}{0.999282in}}%
\pgfpathlineto{\pgfqpoint{5.318036in}{1.028624in}}%
\pgfpathlineto{\pgfqpoint{5.322780in}{0.921036in}}%
\pgfpathlineto{\pgfqpoint{5.327525in}{0.911255in}}%
\pgfpathlineto{\pgfqpoint{5.332269in}{0.877023in}}%
\pgfpathlineto{\pgfqpoint{5.337013in}{0.916146in}}%
\pgfpathlineto{\pgfqpoint{5.341758in}{0.872132in}}%
\pgfpathlineto{\pgfqpoint{5.346502in}{0.911255in}}%
\pgfpathlineto{\pgfqpoint{5.351247in}{0.891694in}}%
\pgfpathlineto{\pgfqpoint{5.355991in}{0.827813in}}%
\pgfpathlineto{\pgfqpoint{5.360736in}{0.891388in}}%
\pgfpathlineto{\pgfqpoint{5.365480in}{0.915840in}}%
\pgfpathlineto{\pgfqpoint{5.370225in}{0.925621in}}%
\pgfpathlineto{\pgfqpoint{5.374969in}{0.842484in}}%
\pgfpathlineto{\pgfqpoint{5.384458in}{0.891388in}}%
\pgfpathlineto{\pgfqpoint{5.389202in}{0.813142in}}%
\pgfpathlineto{\pgfqpoint{5.393947in}{0.837288in}}%
\pgfpathlineto{\pgfqpoint{5.398691in}{0.939986in}}%
\pgfpathlineto{\pgfqpoint{5.403436in}{0.900863in}}%
\pgfpathlineto{\pgfqpoint{5.408180in}{0.920425in}}%
\pgfpathlineto{\pgfqpoint{5.412924in}{0.959548in}}%
\pgfpathlineto{\pgfqpoint{5.417669in}{0.974219in}}%
\pgfpathlineto{\pgfqpoint{5.422413in}{0.891083in}}%
\pgfpathlineto{\pgfqpoint{5.427158in}{0.886192in}}%
\pgfpathlineto{\pgfqpoint{5.431902in}{0.847069in}}%
\pgfpathlineto{\pgfqpoint{5.436647in}{0.836983in}}%
\pgfpathlineto{\pgfqpoint{5.441391in}{0.915229in}}%
\pgfpathlineto{\pgfqpoint{5.446136in}{0.876106in}}%
\pgfpathlineto{\pgfqpoint{5.450880in}{0.827202in}}%
\pgfpathlineto{\pgfqpoint{5.455624in}{0.871215in}}%
\pgfpathlineto{\pgfqpoint{5.460369in}{0.861435in}}%
\pgfpathlineto{\pgfqpoint{5.465113in}{0.876106in}}%
\pgfpathlineto{\pgfqpoint{5.469858in}{0.836983in}}%
\pgfpathlineto{\pgfqpoint{5.474602in}{0.866019in}}%
\pgfpathlineto{\pgfqpoint{5.479347in}{0.836677in}}%
\pgfpathlineto{\pgfqpoint{5.484091in}{0.861129in}}%
\pgfpathlineto{\pgfqpoint{5.488836in}{0.861129in}}%
\pgfpathlineto{\pgfqpoint{5.493580in}{0.875800in}}%
\pgfpathlineto{\pgfqpoint{5.498324in}{0.900252in}}%
\pgfpathlineto{\pgfqpoint{5.503069in}{0.880690in}}%
\pgfpathlineto{\pgfqpoint{5.507813in}{0.934485in}}%
\pgfpathlineto{\pgfqpoint{5.512558in}{0.963521in}}%
\pgfpathlineto{\pgfqpoint{5.522047in}{0.885275in}}%
\pgfpathlineto{\pgfqpoint{5.526791in}{0.904837in}}%
\pgfpathlineto{\pgfqpoint{5.531535in}{0.846152in}}%
\pgfpathlineto{\pgfqpoint{5.536280in}{0.841262in}}%
\pgfpathlineto{\pgfqpoint{5.541024in}{0.851043in}}%
\pgfpathlineto{\pgfqpoint{5.545769in}{0.836066in}}%
\pgfpathlineto{\pgfqpoint{5.550513in}{0.836066in}}%
\pgfpathlineto{\pgfqpoint{5.555258in}{0.884970in}}%
\pgfpathlineto{\pgfqpoint{5.560002in}{0.904531in}}%
\pgfpathlineto{\pgfqpoint{5.564747in}{0.889860in}}%
\pgfpathlineto{\pgfqpoint{5.569491in}{0.909421in}}%
\pgfpathlineto{\pgfqpoint{5.574235in}{0.884970in}}%
\pgfpathlineto{\pgfqpoint{5.578980in}{0.914312in}}%
\pgfpathlineto{\pgfqpoint{5.583724in}{0.831175in}}%
\pgfpathlineto{\pgfqpoint{5.588469in}{0.835760in}}%
\pgfpathlineto{\pgfqpoint{5.593213in}{0.884664in}}%
\pgfpathlineto{\pgfqpoint{5.597958in}{0.889554in}}%
\pgfpathlineto{\pgfqpoint{5.602702in}{0.879774in}}%
\pgfpathlineto{\pgfqpoint{5.607447in}{0.933568in}}%
\pgfpathlineto{\pgfqpoint{5.612191in}{0.874883in}}%
\pgfpathlineto{\pgfqpoint{5.616935in}{0.899335in}}%
\pgfpathlineto{\pgfqpoint{5.621680in}{0.874577in}}%
\pgfpathlineto{\pgfqpoint{5.626424in}{0.859906in}}%
\pgfpathlineto{\pgfqpoint{5.631169in}{0.825674in}}%
\pgfpathlineto{\pgfqpoint{5.635913in}{0.855016in}}%
\pgfpathlineto{\pgfqpoint{5.640658in}{0.864797in}}%
\pgfpathlineto{\pgfqpoint{5.645402in}{0.825674in}}%
\pgfpathlineto{\pgfqpoint{5.650147in}{0.903920in}}%
\pgfpathlineto{\pgfqpoint{5.654891in}{0.879468in}}%
\pgfpathlineto{\pgfqpoint{5.659635in}{0.811003in}}%
\pgfpathlineto{\pgfqpoint{5.664380in}{0.908504in}}%
\pgfpathlineto{\pgfqpoint{5.669124in}{0.879162in}}%
\pgfpathlineto{\pgfqpoint{5.678613in}{0.840039in}}%
\pgfpathlineto{\pgfqpoint{5.683358in}{0.849820in}}%
\pgfpathlineto{\pgfqpoint{5.688102in}{0.840039in}}%
\pgfpathlineto{\pgfqpoint{5.692846in}{0.840039in}}%
\pgfpathlineto{\pgfqpoint{5.697591in}{0.864185in}}%
\pgfpathlineto{\pgfqpoint{5.702335in}{0.927760in}}%
\pgfpathlineto{\pgfqpoint{5.707080in}{0.834843in}}%
\pgfpathlineto{\pgfqpoint{5.711824in}{0.859295in}}%
\pgfpathlineto{\pgfqpoint{5.716569in}{0.854405in}}%
\pgfpathlineto{\pgfqpoint{5.721313in}{0.834843in}}%
\pgfpathlineto{\pgfqpoint{5.726058in}{0.878857in}}%
\pgfpathlineto{\pgfqpoint{5.730802in}{0.820172in}}%
\pgfpathlineto{\pgfqpoint{5.735546in}{0.873966in}}%
\pgfpathlineto{\pgfqpoint{5.740291in}{0.858989in}}%
\pgfpathlineto{\pgfqpoint{5.745035in}{0.868770in}}%
\pgfpathlineto{\pgfqpoint{5.749780in}{0.805195in}}%
\pgfpathlineto{\pgfqpoint{5.759269in}{0.878551in}}%
\pgfpathlineto{\pgfqpoint{5.764013in}{0.819866in}}%
\pgfpathlineto{\pgfqpoint{5.768758in}{0.824757in}}%
\pgfpathlineto{\pgfqpoint{5.773502in}{0.873661in}}%
\pgfpathlineto{\pgfqpoint{5.782991in}{0.863574in}}%
\pgfpathlineto{\pgfqpoint{5.787735in}{0.873355in}}%
\pgfpathlineto{\pgfqpoint{5.792480in}{0.824451in}}%
\pgfpathlineto{\pgfqpoint{5.797224in}{0.824451in}}%
\pgfpathlineto{\pgfqpoint{5.801969in}{0.873355in}}%
\pgfpathlineto{\pgfqpoint{5.806713in}{0.804890in}}%
\pgfpathlineto{\pgfqpoint{5.811458in}{0.809780in}}%
\pgfpathlineto{\pgfqpoint{5.816202in}{0.853488in}}%
\pgfpathlineto{\pgfqpoint{5.825691in}{0.882830in}}%
\pgfpathlineto{\pgfqpoint{5.830435in}{0.838817in}}%
\pgfpathlineto{\pgfqpoint{5.835180in}{0.819255in}}%
\pgfpathlineto{\pgfqpoint{5.839924in}{0.838817in}}%
\pgfpathlineto{\pgfqpoint{5.844669in}{0.824145in}}%
\pgfpathlineto{\pgfqpoint{5.849413in}{0.789913in}}%
\pgfpathlineto{\pgfqpoint{5.854157in}{0.804278in}}%
\pgfpathlineto{\pgfqpoint{5.858902in}{0.833621in}}%
\pgfpathlineto{\pgfqpoint{5.863646in}{0.843401in}}%
\pgfpathlineto{\pgfqpoint{5.868391in}{0.809169in}}%
\pgfpathlineto{\pgfqpoint{5.873135in}{0.872744in}}%
\pgfpathlineto{\pgfqpoint{5.877880in}{0.809169in}}%
\pgfpathlineto{\pgfqpoint{5.882624in}{0.838511in}}%
\pgfpathlineto{\pgfqpoint{5.887369in}{0.882524in}}%
\pgfpathlineto{\pgfqpoint{5.892113in}{0.901780in}}%
\pgfpathlineto{\pgfqpoint{5.896857in}{0.901780in}}%
\pgfpathlineto{\pgfqpoint{5.901602in}{0.838205in}}%
\pgfpathlineto{\pgfqpoint{5.906346in}{0.803973in}}%
\pgfpathlineto{\pgfqpoint{5.911091in}{0.843096in}}%
\pgfpathlineto{\pgfqpoint{5.915835in}{0.789301in}}%
\pgfpathlineto{\pgfqpoint{5.920580in}{0.794192in}}%
\pgfpathlineto{\pgfqpoint{5.925324in}{0.818644in}}%
\pgfpathlineto{\pgfqpoint{5.930069in}{0.862352in}}%
\pgfpathlineto{\pgfqpoint{5.934813in}{0.862352in}}%
\pgfpathlineto{\pgfqpoint{5.939557in}{0.837900in}}%
\pgfpathlineto{\pgfqpoint{5.944302in}{0.872132in}}%
\pgfpathlineto{\pgfqpoint{5.944302in}{0.872132in}}%
\pgfusepath{stroke}%
\end{pgfscope}%
\begin{pgfscope}%
\pgfsetrectcap%
\pgfsetmiterjoin%
\pgfsetlinewidth{0.803000pt}%
\definecolor{currentstroke}{rgb}{0.000000,0.000000,0.000000}%
\pgfsetstrokecolor{currentstroke}%
\pgfsetdash{}{0pt}%
\pgfpathmoveto{\pgfqpoint{3.833026in}{0.565123in}}%
\pgfpathlineto{\pgfqpoint{3.833026in}{1.475926in}}%
\pgfusepath{stroke}%
\end{pgfscope}%
\begin{pgfscope}%
\pgfsetrectcap%
\pgfsetmiterjoin%
\pgfsetlinewidth{0.803000pt}%
\definecolor{currentstroke}{rgb}{0.000000,0.000000,0.000000}%
\pgfsetstrokecolor{currentstroke}%
\pgfsetdash{}{0pt}%
\pgfpathmoveto{\pgfqpoint{5.941666in}{0.565123in}}%
\pgfpathlineto{\pgfqpoint{5.941666in}{1.475926in}}%
\pgfusepath{stroke}%
\end{pgfscope}%
\begin{pgfscope}%
\pgfsetrectcap%
\pgfsetmiterjoin%
\pgfsetlinewidth{0.803000pt}%
\definecolor{currentstroke}{rgb}{0.000000,0.000000,0.000000}%
\pgfsetstrokecolor{currentstroke}%
\pgfsetdash{}{0pt}%
\pgfpathmoveto{\pgfqpoint{3.833026in}{0.565123in}}%
\pgfpathlineto{\pgfqpoint{5.941666in}{0.565123in}}%
\pgfusepath{stroke}%
\end{pgfscope}%
\begin{pgfscope}%
\pgfsetrectcap%
\pgfsetmiterjoin%
\pgfsetlinewidth{0.803000pt}%
\definecolor{currentstroke}{rgb}{0.000000,0.000000,0.000000}%
\pgfsetstrokecolor{currentstroke}%
\pgfsetdash{}{0pt}%
\pgfpathmoveto{\pgfqpoint{3.833026in}{1.475926in}}%
\pgfpathlineto{\pgfqpoint{5.941666in}{1.475926in}}%
\pgfusepath{stroke}%
\end{pgfscope}%
\begin{pgfscope}%
\definecolor{textcolor}{rgb}{0.000000,0.000000,0.000000}%
\pgfsetstrokecolor{textcolor}%
\pgfsetfillcolor{textcolor}%
\pgftext[x=4.887346in,y=1.559260in,,base]{\color{textcolor}\rmfamily\fontsize{12.000000}{14.400000}\selectfont 18,116 kHz}%
\end{pgfscope}%
\end{pgfpicture}%
\makeatother%
\endgroup%

		\caption{Measured frequency content of a FM modulated 200 kHz sine wave with modulation frequencies at solutions of the Bessel functions.}
		\label{fig:446-200}
	\end{figure}
	
	
	\subsubsection{}
	The DCF77 transmitter uses amplitude shift keying (ASK). The carrier at 77.5 kHz is periodically attenuated to about 15 \% of its amplitude, and the duration of the attenuation encodes the binary time information \cite{wiki:DCF77}.
	
	\subsubsection{}

	A frequency-modulated signal can be demodulated with an envelope detector by converting the frequency variations into amplitude variations using a frequency-dependent network (LC resonant circuit operated on its slope). The resulting amplitude-modulated signal is then demodulated with a conventional AM envelope detector consisting of a diode, resistor, and capacitor.

	\vfill
	
		
	% --- Bibliography ---
	\bibliographystyle{plain}
	\bibliography{referenzen}
	
	
\end{document}